\documentclass[11pt,twoside,openright]{memoir}
\usepackage[localtheorems]{raeez-math-template}
\usepackage{platonic}

% Cross-book notation.  Definitions are idempotent so that every extracted
% chapter remains independently meaningful.
\providecommand{\HHtwo}{\mathbb H_2}
\providecommand{\Lam}{\Lambda}
\providecommand{\Dtheta}{{\Delta^{\theta}_5}}
\providecommand{\Dmon}{{D_5}}
\providecommand{\ChiTen}{{\chi_{10}}}
\providecommand{\phiweak}{{\phi_{0,1}}}
\providecommand{\gIg}{{\mathfrak g_{\Delta_5}}}
\providecommand{\nIg}{{\mathfrak n_{\Delta_5,+}}}
\providecommand{\Gammaeff}{\Gamma_{\mathrm{eff}}}
\providecommand{\Gammaind}{\Gamma_{\mathrm{ind}}}
\providecommand{\Gammphys}{\Gamma_{\mathrm{phys}}}
\providecommand{\PiGram}{\Pi}
\providecommand{\den}{\operatorname{den}}
\providecommand{\DenIg}{\operatorname{Den}_{\Delta}}
\providecommand{\smult}{\operatorname{smult}}
\providecommand{\Pf}{\operatorname{Pf}}
\providecommand{\CY}{\operatorname{CY}}
\providecommand{\CC}{\operatorname{CC}}
\providecommand{\Crit}{\operatorname{Crit}}
\providecommand{\IndCoh}{\operatorname{IndCoh}}
\providecommand{\QCoh}{\operatorname{QCoh}}
\providecommand{\Sh}{\operatorname{Sh}}
\providecommand{\Stab}{\operatorname{Stab}}
\providecommand{\Bun}{\operatorname{Bun}}
\providecommand{\LocSys}{\operatorname{LocSys}}
\providecommand{\Opers}{\operatorname{Op}}
\providecommand{\qgroup}{\mathcal U}
\providecommand{\Yplus}{Y^+}
\providecommand{\Yfull}{Y}
\providecommand{\Mukai}{\widetilde H}
\providecommand{\Dfive}{{D_5}}
\providecommand{\gDelta}{{\mathfrak g_{\Delta_5}}}
\providecommand{\nDelta}{{\mathfrak n_{\Delta_5,+}}}
\providecommand{\Rpi}{\mathcal R_{\Pi}}
\providecommand{\HT}{\mathrm{HT}}
\providecommand{\hCS}{\mathrm{hCS}}
\providecommand{\BV}{\mathrm{BV}}
\providecommand{\QME}{\mathrm{QME}}
\providecommand{\CME}{\mathrm{CME}}
\providecommand{\ev}{\operatorname{ev}}
\providecommand{\supp}{\operatorname{supp}}
\providecommand{\Bimod}{\operatorname{Bimod}}
\providecommand{\Bicomod}{\operatorname{Bicomod}}
\providecommand{\coHH}{\operatorname{coHH}}
\providecommand{\MC}{\operatorname{MC}}
\providecommand{\DR}{\operatorname{DR}}
\providecommand{\Vir}{\operatorname{Vir}}
\providecommand{\Ein}{\operatorname{Ein}}
\providecommand{\BVEnv}{\mathbf{BVEnv}}
\providecommand{\QHdg}{\mathbf{QHdg}}
\providecommand{\EinBV}{\mathbf{EinBV}}
\providecommand{\Coll}{\operatorname{Coll}}
\providecommand{\Def}{\operatorname{Def}}
\providecommand{\BarMix}{\operatorname{Bar}^{\mathrm{mix}}}
\providecommand{\Cent}{\operatorname{Cent}}
\providecommand{\Proj}{\operatorname{Proj}}
\providecommand{\Cal}{\operatorname{Cal}}
\providecommand{\adS}{\mathrm{AdS}}
\providecommand{\Diff}{\operatorname{Diff}}
\providecommand{\FactEnv}{\operatorname{FactEnv}}
\providecommand{\Hall}{\operatorname{Hall}}
\providecommand{\Borch}{\operatorname{Borch}}
\providecommand{\OpenClosed}{\operatorname{OC}}
\providecommand{\Root}{\operatorname{Root}}
\providecommand{\CurE}{\operatorname{Cur}_{E}}
\providecommand{\Quant}{\operatorname{Quant}}
\providecommand{\Char}{\operatorname{Char}}
\providecommand{\FactD}{\operatorname{FactD}}
\providecommand{\FactAlg}{\operatorname{FactAlg}}
\providecommand{\HallBPS}{\operatorname{HallBPS}}
\providecommand{\AutForm}{\operatorname{AutForm}}
\providecommand{\CYQGData}{\mathsf{CYQGData}}
\providecommand{\RealDatum}{\mathsf{Real}_{\Delta_5}}
\providecommand{\GeoDatum}{\mathsf{Geo}_{K3\times E}}
\providecommand{\Uuniv}{\mathcal U_{\Delta_5}^{\mathrm{univ}}}

\title{\texorpdfstring{\textbf{A-infinity Chiral Algebras and 3D Holomorphic--Topological QFT}\\[0.55em]\Large Volume II of the programme: Hamiltonian jets, shifted cotangent fields, universal branes, and the open--closed map}{A-infinity Chiral Algebras and 3D Holomorphic--Topological QFT: Volume II of the programme: Hamiltonian jets, shifted cotangent fields, universal branes, and the open--closed map}}
\author{Raeez Lorgat}
\date{26 July 2026}
\hypersetup{pdftitle={Mixed Holomorphic-Topological Theory},pdfauthor={Raeez Lorgat},pdfsubject={Hamiltonian jets, shifted cotangent fields, universal branes, and open-closed comparison},pdfkeywords={A-infinity chiral algebra, holomorphic-topological field theory, Hamiltonian jet, shifted cotangent, open-closed map}}
\begin{document}
\frontmatter
\maketitle
\thispagestyle{empty}
\cleardoublepage
\chapter*{Abstract}
\addcontentsline{toc}{chapter}{Abstract}

Bulk reconstruction requires coherent actions on boundaries and on the operators between them. This book constructs finite Hamiltonian models, supported collision operations and marked reconstruction maps, while retaining their geometric and physical comparison problems.

Finite Hamiltonian jets give shifted-cotangent models satisfying the quantum master equation by vanishing adjoint traces. For a finite-dimensional complex Lie algebra, relative Hochschild cochains of its Rees envelope correspond to adjoint Chevalley cochains with differential \(\hbarq d_{\mathrm{CE}}\), including finite coefficient quotients and the formal limit. Quadratic Moyal Hamiltonians give the even Weyl algebra through the oscillator Casimir quotient. Separate input and output jets determine the continuous cochain carrier; its operadic refinement remains additional.

Polynomial-root current actions retain a nonzero Massey coset, integral comparison paths and derived specialization. Finite interface centres have brackets and higher compositions beyond their cup products. Tree expansions construct action profiles and multiplicative recipients for finite nilpotent families. Explicit pointed and ordinary multiplicative actions obstruct finality in their specified categories.

Wick fields and mixed forest homotopies give meromorphic coefficients with point-source contacts. Outer principal-part classes survive in every arity at least three. Supported Heisenberg pair maps on regular singleton inputs admit a complete triple homotopy; the common four-input comparison remains open. Finite Gaussian elimination retains cumulants, determinant, conditional differential and sewing. Infinite-mode evolution has its numerical Fock states, formal cubic sewing and a positive quartic stabilization with a sharp continuity threshold.

Hamiltonians act in all arities on the resolved plane and derived matrix spaces. Rank-two relative quantum restriction is an integral quasi-isomorphism; absolute reduction retains its gauge class and higher operations. Cyclic traces retain the central rank, while the full Hamiltonian/cyclic comparison remains open.

Gaussian half-space fields give supported boundary actions and an incoming cap on a specified finite-jet domain. Derived normal support reconstructs on the normal Weyl carrier. Marked Morita reconstruction holds on generated sectors. The Ramond junction has nonzero polynomial right orthogonals, a local residue cone without an original-quartet lift, and a polynomial field prescription that fails derived invariance. Shuffle lattices distinguish integral multiplication from rational cooperation. Tensor reconstruction and gravitational examples specify the additional maps required to preserve local fields, nonlinear equations and physical sewing.

\chapter*{Operations and reconstruction}
\addcontentsline{toc}{chapter}{Operations and reconstruction}
A boundary algebra records compositions. Its derived centre records coherent actions on its module category; the derived end of one selected module instead detects the sector that module generates. Reconstructing bulk fields requires an independently defined bulk, its boundary action and a comparison preserving positions, supports and collisions.

Quantization must deform the map as well as its endpoints. Over \(k[\varepsilon]/(\varepsilon^2)\), the relation \(x^2=\varepsilon\) defines a free rank-two deformation of \(k[x]/(x^2)\). Its augmentation cannot lift to the undeformed coefficient algebra: any lift sends \(x\) to \(\varepsilon c\), whose square is zero. An algebra deformation alone therefore does not provide the required bulk--boundary deformation.

Chapters 1--7 construct the algebraic and supported ingredients of this problem. Morphism homotopies lead to coefficient modules on ordered configurations, finite Hamiltonian jets and continuous cochains. Polynomial-root interfaces retain Massey products, comparison fibres and gauge stabilizers. Action profiles give explicit recipients, while two distinct nonfinality arguments constrain universal pointed and multiplicative actions. Supported pair operations admit a triple comparison; four inputs require a common recipient and the pentagon primitive for the prescribed lower homotopies.

Chapters 8--13 retain quantum and state data. Gaussian elimination transports the differential, odd inclusion, determinant and all cumulants. The cubic matrix trace defect forces higher Hamiltonian components. Integral relative rank-two restriction, absolute gauge cohomology and cyclic trace calculations preserve data lost by degree-zero invariants. Infinite-mode states require a specified completion; a positive quartic stabilization changes the cubic theory and has its own continuity threshold.

Chapters 14--15 begin with native fields and an independent junction. The Gaussian boundary action, corrected cutoff system and incoming cap have explicit coefficient domains. Normal Weyl reconstruction uses derived adic support, with a separate comparison to physical fusion. The Ramond junction's reverse mode algebra is computed, but the junction misses nonzero polynomial boundary objects. Its local residue also produces a cone that cannot carry the original quartet. An acyclic localization shows why unchanged polynomial local fields cannot be the derived mapping objects on every pair.

Chapter 16 compares the remaining algebraic and geometric structures. Ordinary bars, integral shuffle lattices and rational cooperation retain different operations. Matrix coefficients reconstruct a specified tensor category. Gravitational examples distinguish tangent complexes from nonlinear equations and connection coefficients from holonomy. The full reconstruction problem must preserve local mapping objects, both tensor actions, units, preparations and coherent sewing.

\chapter*{Notation and normalization}
\addcontentsline{toc}{chapter}{Notation and normalization}
The ground field has characteristic zero. Ambient stable categories used for colimit constructions are presentable and idempotent-complete, and their tensor products preserve colimits separately. Perfect sectors retain their stated finite-generation conditions. Each completion specifies its coefficient ring, filtration and transition maps; ordinary and derived inverse limits are distinguished. The Igusa normalization is
\[
 \Dtheta=64\Dmon,\qquad \ChiTen=\Dmon^2,
\]
with $\Dmon$ monic.  Physical charge is additive; Gram degree is carried through its polarization cocycle.  Bar, factorization, and Hopf coproducts retain distinct notation and distinct geometric meanings.

\tableofcontents
\mainmatter
\chapter{Boundary actions and their comparison homotopies}
\label{ch:vol2-boundary-actions}

A bulk observable acts on boundary conditions and on the operators
between them. Its action may vanish on each boundary while remaining
nonzero on the category of boundaries. Two cubic boundaries exhibit
this phenomenon. They lead to a complex that retains the actions,
the homotopies comparing them, and the equations imposed by composition.
The resulting ordinary derived centre gives an algebraic model of
coherent actions. In a mixed holomorphic--topological field theory,
which of these actions come from local bulk observables?

\section{A relation detected between two boundaries}
\label{sec:vol2-cubic-boundaries}

Let $k$ be a field of characteristic zero and put $R=k[x]$.
A \emph{matrix factorization} of $W\in R$ is a finite free
parity-graded $R$-module $E=E^{\bar0}\oplus E^{\bar1}$ with an odd
endomorphism $Q_E$ satisfying $Q_E^2=W\id_E$.

For two factorizations of the same potential, set
\begin{equation}\label{eq:vol2-mf-differential}
 \delta f=Q_Ff-(-1)^{|f|}fQ_E,
 \qquad f\in\Hom_R(E,F).
\end{equation}
Here $|f|\in\Z/2$ is its parity. Since $W$ is central,
$\delta^2f=Wf-fW=0$. Composition satisfies the graded Leibniz rule.
Thus the factorizations form a differential category with parity grading.

For the cubic potential, take the factorizations on $R^{1|1}$
\begin{equation}\label{eq:vol2-cubic-pair}
 Q_X=\begin{pmatrix}0&x\\x^2&0\end{pmatrix},
 \qquad
 Q_Y=\begin{pmatrix}0&x^2\\x&0\end{pmatrix},
 \qquad
 f=\begin{pmatrix}x&0\\0&1\end{pmatrix}:X\longrightarrow Y.
\end{equation}
The map $f$ is even and closed. Multiplication by $x$ has an odd
primitive on each object: there are $h_X,h_Y$ with
$\delta h_X=x\id_X$ and $\delta h_Y=x\id_Y$.
Their compatibility along $f$ is measured by
\[
 c(f)=h_Yf-fh_X.
\]
It is closed, because $\delta c(f)=xf-fx=0$.

\begin{proposition}\label{prop:vol2-cubic-defect}
For the objects and map in \eqref{eq:vol2-cubic-pair}, every choice
of $h_X,h_Y$ gives the same nonzero class
\[
 [c(f)]=[J]\in H^{\bar1}\Hom_R(X,Y),
 \qquad J=\begin{pmatrix}0&1\\-1&0\end{pmatrix}.
\]
In particular, the objectwise primitives cannot be extended to a
nullhomotopy compatible with this morphism.
\end{proposition}

\begin{proof}
The most general primitives are
\[
 h_X=\begin{pmatrix}0&b\\1-xb&0\end{pmatrix},
 \qquad
 h_Y=\begin{pmatrix}0&1-xc\\c&0\end{pmatrix},
 \qquad b,c\in R.
\]
Indeed, an odd matrix with upper entry $u$ and lower entry $v$
is a primitive on $X$ precisely when $xv+x^2u=x$, so
$v=1-xu$. The equation on $Y$ is analogous.

Multiplication gives
\begin{equation}\label{eq:vol2-cubic-defect}
 c(f)=(1-x(b+c))J.
\end{equation}
For an even morphism $a=\operatorname{diag}(a_0,a_1)$ and an odd
morphism with entries $b_0,c_0$, the two differentials are
\[
 \delta a=x(xa_1-a_0)J,
 \qquad
 \delta\begin{pmatrix}0&b_0\\c_0&0\end{pmatrix}
       =x(b_0+c_0)f.
\]
Consequently the even cycles and boundaries are $Rf$ and $xRf$,
and the odd cycles and boundaries are $RJ$ and $xRJ$.
The coefficient of $J$ in \eqref{eq:vol2-cubic-defect} is one
modulo $x$, independently of $b,c$.
\end{proof}

For three objects with such primitives and closed even maps
$f:X\to Y$, $g:Y\to Z$, the same definition gives
\[
 c(gf)=c(g)f+gc(f).
\]
Thus even a primitive for each $c(f)$ would leave a compatibility
equation for composition. The complex of coherent actions must contain
components on morphisms and on longer composable strings.

\section{The ordinary derived centre}
\label{sec:vol2-ordinary-centre}

Let $\mathcal E$ be a small unital differential graded category
over $k$. Complexes have cohomological grading,
$V[1]^j=V^{j+1}$. The same constructions apply with degrees reduced
modulo two to the preceding category of matrix factorizations.
All tensors in this section are over $k$, and total Hom complexes
use products. In particular, no boundedness in the cohomological
degrees is implicit.
Write $D(k)$ for the derived category of coefficient complexes.

The \emph{diagonal bimodule} is
$\Delta_{\mathcal E}(a,b)=\mathcal E(a,b)$, with the two actions
given by composition. Write
$\mathcal E^e=\mathcal E^{\mathrm{op}}\otimes_k^{\mathbb L}\mathcal E$
for the derived enveloping category. The
\emph{ordinary derived centre} and its cohomology are
\begin{equation}\label{eq:vol2-ordinary-centre}
 Z_k(\mathcal E)=
 \RHom_{\mathcal E^e}(\Delta_{\mathcal E},\Delta_{\mathcal E}),
 \qquad
 \HH_k^*(\mathcal E)=H^*Z_k(\mathcal E).
\end{equation}
This definition fixes the coefficient tensor. Relative centres over
a larger ring will use that ring's tensor and derived Hom instead.

To compute \eqref{eq:vol2-ordinary-centre}, resolve the diagonal
by its two-sided bar. After bimodule Hom, its cochains have factors
\begin{equation}\label{eq:vol2-all-object-bar}
 \prod_{n\ge0}\prod_{w_0,\ldots,w_n}
 \Hom_k\!\left(
 \mathcal E(w_{n-1},w_n)[1]\otimes\cdots\otimes
 \mathcal E(w_0,w_1)[1],\mathcal E(w_0,w_n)\right).
\end{equation}
At $n=0$ the input is $k$. Every intermediate object is an index
of the product. The normalized cochains vanish when an input is an
identity. Denote their product totalization by
$\CC_k^\bullet(\mathcal E)$.

The differential has an internal part and a part that composes
adjacent inputs or acts at an endpoint. An unsuspended convention
makes their signs explicit. An arity-$n$ cochain $F$ of internal
degree $r$ has total degree $r+n$. Put $\partial F=[\delta,F]$,
the graded Hom differential.

For homogeneous composable inputs,
ordered so that the rightmost acts first, set
\begin{align}\label{eq:vol2-hochschild-signs}
 (b_HF)(a_1,\ldots,a_{n+1})
 &=(-1)^{r|a_1|}a_1F(a_2,\ldots,a_{n+1})\notag\\
 &\quad+\sum_{j=1}^n(-1)^j
 F(a_1,\ldots,a_ja_{j+1},\ldots,a_{n+1})\notag\\
 &\quad+(-1)^{n+1}F(a_1,\ldots,a_n)a_{n+1},\\
 (\mathcal DF)_n&=(-1)^n\partial F_n+b_HF_{n-1},
 \qquad F_{-1}=0.\label{eq:vol2-hochschild-total}
\end{align}
Associativity and the Leibniz rule give
$b_H^2=\partial^2=0$ and $\partial b_H=b_H\partial$.
Hence $\mathcal D^2=0$. Each component involves only two arities,
so the differential is defined on the product totalization.

\begin{proposition}[Bar representation and evaluation]
\label{prop:vol2-centre-evaluation}
The normalized complex $\CC_k^\bullet(\mathcal E)$ represents
$Z_k(\mathcal E)$. For every coefficient complex $P\in D(k)$,
evaluation induces a natural equivalence
\begin{equation}\label{eq:vol2-centre-evaluation}
 \RHom_k(P,Z_k(\mathcal E))
 \simeq
 \RHom_{\mathcal E^e}
 (P\otimes_k^{\mathbb L}\Delta_{\mathcal E},\Delta_{\mathcal E}).
\end{equation}
\end{proposition}

\begin{proof}
Use cofibrant models for the coefficient complexes. The augmented
two-sided bar contracts by insertion of a unit at an outside end.

Its filtration by word length has induced bimodules as successive
quotients. Over the field $k$, their coefficient complexes are
homotopically projective: Hom from them into an acyclic complex
is acyclic. A finite length stage therefore has the same property.

Hom out of the full direct-sum bar is the inverse limit of the
Hom complexes of these stages. The transition maps are surjective
in every degree, since the length filtration splits on graded
modules. The kernels are acyclic by the same coefficient argument,
so primitives can be lifted successively through the tower.
Its inverse limit is therefore acyclic. Thus the bar is a
homotopically projective resolution.

The degeneracy contraction splits off the cochains containing an
identity, leaving the normalized complex. Applying bimodule Hom
to the remaining resolution gives \eqref{eq:vol2-all-object-bar}.
For a projective model of $P$, the tensor--Hom adjunction on each
bar term identifies the two sides of
\eqref{eq:vol2-centre-evaluation}. Products commute with this
adjunction, and the identifications commute with the differential.
\end{proof}

The identity action is the nullary family $\id_w$.
The cup product composes outputs after cutting a composable word.
The brace operations insert the output of one cochain into an input
of another, with the Koszul signs of the suspended inputs.
These are operations on the full cochain complex. Retaining only
the cohomology of the endomorphisms of each object removes the
positive-arity terms where Proposition~\ref{prop:vol2-cubic-defect}
found its obstruction.

\section{The scalar action on all cubic boundaries}
\label{sec:vol2-cubic-scalar-centre}

The cubic calculation extends to every finite free boundary.
It distinguishes a relation in the category from a relation on each
of its objects. Put $R=k[[x]]$. Let $\mathcal C$ be the
parity differential category of all matrix factorizations of $x^3$
on standard finite free parity-graded $R$-modules. Fixing standard
modules makes the objects a set. Its morphisms are $R$-linear,
with differential \eqref{eq:vol2-mf-differential}.

Here the centre is \emph{absolute over $k$}: its cochains are
$k$-multilinear, with no continuity condition on those maps.
They are normalized and product-totalized as above. For $g\in R$,
scalar multiplication defines a cocycle $\rho(g)$ whose nullary
value is $g\id_E$ and whose higher values vanish.

\begin{theorem}\label{thm:vol2-cubic-scalar-kernel}
The kernel of the scalar action
\[
 R\longrightarrow\HH_k^{\bar0}(\mathcal C),
 \qquad g\longmapsto[\rho(g)],
\]
is $(x^2)$. Nevertheless, multiplication by $x$ is nullhomotopic
on every object of $\mathcal C$.
\end{theorem}

\begin{proof}
For a cochain primitive with nullary component $h_E$ and unary
component $u$, equations \eqref{eq:vol2-hochschild-signs} and
\eqref{eq:vol2-hochschild-total} begin with
\begin{align}
 \delta h_E&=g\id_E,\label{eq:vol2-scalar-primitive-zero}\\
 \delta u(f)-u(\delta f)
 &=(-1)^{|f|}fh_E-h_Ff.
 \label{eq:vol2-scalar-primitive-one}
\end{align}
The map $h_E$ is odd and $u$ preserves parity.
If all later components vanish, the binary equation is
$u(vf)=u(v)f+vu(f)$.

Differentiate the matrices in their fixed bases and take
\[
 h_E=\frac13Q'_E,
 \qquad u(f)=\frac13f',
 \qquad F_n=0\quad(n\ge2).
\]
Differentiating $Q_E^2=x^3\id_E$ proves
\eqref{eq:vol2-scalar-primitive-zero} for $g=x^2$.
For a homogeneous morphism,
\[
 (\delta f)'=\delta(f')+Q'_Ff-(-1)^{|f|}fQ'_E,
\]
which proves \eqref{eq:vol2-scalar-primitive-one}.
The derivative kills identities and satisfies the Leibniz rule,
so every remaining component of $\mathcal DF$ vanishes.
Multiplying the components by any $a\in R$ gives a primitive
for $ax^2$, since both cochain differentials commute with the
central scalar action on the output.

Conversely, reduce $g$ modulo $(x^2)$ to $\alpha+\beta x$.
On the object $X$ of \eqref{eq:vol2-cubic-pair}, every odd
matrix has the form
\[
 h=\begin{pmatrix}0&b\\d&0\end{pmatrix},
 \qquad \delta h=x(d+xb)\id_X.
\]
The nullary equation forces $\alpha=0$ and $d+xb=\beta$.

The odd endomorphism
\[
 \psi=\begin{pmatrix}0&1\\-x&0\end{pmatrix}
\]
is closed. The unary equation on $\psi$ would give
\[
 \delta u(\psi)=-(\psi h+h\psi)=-(\beta-2xb)\id_X.
\]
Every matrix boundary on $X$ is divisible by $x$.
Evaluation at $x=0$ forces $\beta=0$. Higher arities do not
enter these two equations, so they cannot remove the obstruction.

It remains to prove the objectwise assertion for arbitrary rank.
Over the fraction field of $R$, the two odd blocks of $Q_E$
are inverse up to $x^3$. The even and odd ranks therefore agree.
Smith normal form over the discrete valuation ring $R$
diagonalizes one block. The equation $Q_E^2=x^3\id_E$ forces
the other block to be diagonal in the complementary bases.
The factorization is consequently a direct sum of blocks
$(x^a,x^{3-a})$, where $0\le a\le3$.

For $a=0,3$ one block is invertible, so the object is contractible.
For $a=1$ the matrix
\[
 H=\begin{pmatrix}0&0\\1&0\end{pmatrix}
 \qquad\text{satisfies}\qquad \delta H=x\id.
\]
Exchanging the two parity summands gives a primitive for $a=2$.
Direct sums and changes of basis preserve these homotopies.
\end{proof}

The scalar image in this absolute centre is therefore
$k[[x]]/(x^2)$. This computes the scalar image, without asserting
that it is the entire centre. The coefficient convention enters
the proof at the unary map:
\[
 u(xf)=\frac13f+xu(f).
\]
It is $k$-linear and is not $R$-linear. A relative centre over
$R$ asks a different question, since it excludes this primitive.

\section{Two scalar endpoints joined by a complex}
\label{sec:vol2-triangular-centre}

The cubic example required an infinite bar of a category.
There is also a finite model in which the entire relative centre,
including its multiplication in a family, can be computed.

Let $R$ be a commutative ring and $M$ a bounded complex of finite
projective $R$-modules. Use integer cohomological degrees and set
\begin{equation}\label{eq:vol2-triangular-algebra}
 T(M)=Re_0\oplus Re_1\oplus M,
 \qquad e_1Me_0=M,
 \qquad M^2=0.
\end{equation}

The idempotents are orthogonal, closed, and sum to the unit.
All other idempotent actions on $M$ vanish. The differential on
$M$ is its given differential. Thus $M$ is the complex of arrows
from object zero to object one, and the endpoint endomorphisms
are scalar. All tensors in this section are over $R$ unless
another ring is displayed.

Put $T=T(M)$ and $T^e=T\otimes_R T^{\mathrm{op}}$, with graded
opposite multiplication. The stated projectivity makes this a
model of the derived enveloping algebra. Define
\[
 Z_R(T)=\RHom_{T^e}(T,T),
 \qquad
 E=\End_R(M),
 \qquad d_EF=d_MF-(-1)^{|F|}Fd_M.
\]
The difference of the two scalar actions is the chain map
\[
 \chi:R^2\longrightarrow E,
 \qquad \chi(a,b)=(a-b)\id_M,
\]
where $a$ acts at the source and $b$ at the target.
Write $D(R)$ for the derived category of $R$-module complexes.

\begin{theorem}\label{thm:vol2-triangular-cone}
The relative centre is represented by the homotopy fibre
\begin{equation}\label{eq:vol2-triangular-cone}
 Z_R(T(M))\simeq\operatorname{Cone}(\chi)[-1].
\end{equation}
With the fibre sign convention, its terms and differential are
\[
 Z^n=(R^2)^n\oplus E^{n-1},
 \qquad d_Z(a,b,F)=(da,db,\chi(a,b)-d_EF).
\]
For every $P\in D(R)$, evaluation gives
\begin{equation}\label{eq:vol2-relative-evaluation}
 \RHom_R(P,Z_R(T))
 \simeq\RHom_{T^e}(P\otimes_R^{\mathbb L}T,T).
\end{equation}
\end{theorem}

\begin{proof}
Let $S=Re_0\oplus Re_1$. The normalized relative bimodule bar
of $T$ over $S$ stops at length one, since $M\otimes_SM=0$:
\begin{equation}\label{eq:vol2-short-bar}
 T\otimes_SM\otimes_ST
\longrightarrow T\otimes_ST\longrightarrow T.
\end{equation}

The first term has the extra degree $-1$ of the bar suspension.
Choose the boundary on an arrow to be
$e_1\otimes m-m\otimes e_0$, with the corresponding suspended
internal differential. Multiplication adds the two copies of
$M$ in the middle term. Their difference is its kernel.
On the idempotents, multiplication is the identity.
This proves exactness of the augmented sequence of complexes.

The terms are induced from the split idempotent bimodules of
$S$, with bounded projective coefficients $R$ or $M$.
They are therefore homotopically projective as $T$-bimodules.
Applying bimodule Hom into $T$ yields $R^2$ at length zero
and $\End_R(M)$ at length one. The two endpoint actions give
$(a-b)\id_M$, and suspension gives $-d_E$.
This is \eqref{eq:vol2-triangular-cone}.
Tensoring the same bar with a projective model of $P$ and using
tensor--Hom adjunction proves \eqref{eq:vol2-relative-evaluation}.
\end{proof}

\begin{example}[Parallel arrows]\label{ex:vol2-parallel-arrows}
For $R=k$ and $M=k^m$ in degree zero, with $m>0$, the centre
complex is
\[
 k^2\longrightarrow\operatorname{Mat}_m(k),
 \qquad (a,b)\longmapsto(a-b)\id,
\]
in degrees zero and one. Hence
\[
 \HH_k^0(T)=k,
 \qquad \HH_k^1(T)=\operatorname{Mat}_m(k)/k\id,
 \qquad \HH_k^j(T)=0\quad(j\ne0,1).
\]
Insertion of one unary cochain into another is composition on
$M$, so the degree-one Gerstenhaber bracket is the matrix
commutator modulo scalars. The product of two positive-degree
classes vanishes: a cup product would require two composable
arrows, and there is no such word. Thus a single arrow has no
positive centre cohomology, whereas three parallel arrows give
eight independent degree-one classes.
\end{example}

\section{A differential varying with a parameter}
\label{sec:vol2-varying-centre}

The joining complex can change while both endpoint algebras
remain fixed. Put $R=k[t]$ and take
\[
 M_t=[Ru\xrightarrow{\ t\ }Rv],
 \qquad |u|=-1,
 \qquad |v|=0.
\]
Its endomorphism complex is
\[
 E^{-1}=R\xrightarrow{s\mapsto(ts,ts)}E^0=R^2
 \xrightarrow{(p,q)\mapsto t(p-q)}E^1=R.
\]
Here $s$ sends $v$ to $su$, the pair $(p,q)$ acts on $(u,v)$,
and the last scalar sends $u$ to a multiple of $v$.
The centre in Theorem~\ref{thm:vol2-triangular-cone} is therefore
\begin{equation}\label{eq:vol2-varying-complex}
 R^3\xrightarrow{D_0}R^2\xrightarrow{D_1}R,
 \qquad
 \begin{aligned}
 D_0(a,b,s)&=(a-b-ts,a-b-ts),\\
 D_1(p,q)&=-t(p-q),
 \end{aligned}
\end{equation}
in degrees zero, one, and two.

For a complex of free $R$-modules, its \emph{Bockstein}
on cohomology modulo $t$ sends $[\bar c]$ to
$[t^{-1}dc\bmod t]$, where $c$ lifts the cocycle $\bar c$.
This is the connecting map for the first-order thickening
$k[t]/(t^2)$, with the submodule $(t)$ identified with $k$.

\begin{theorem}\label{thm:vol2-varying-centre}
The full relative Hochschild cohomology algebra is
\begin{equation}\label{eq:vol2-varying-ring}
 \HH_R^*(T(M_t))
 \cong\frac{R[z,\omega]}
 {(z^2-tz,\ t\omega,\ z\omega,\ \omega^2)},
 \qquad |z|=0,
 \qquad |\omega|=2.
\end{equation}
After inverting $t$, it is $R[t^{-1}]\times R[t^{-1}]$ in
degree zero. The derived fibre
$Z_R(T(M_t))\otimes_R^{\mathbb L}k$ has degree-zero algebra
$k[z]/(z^2)$, together with one class $\epsilon$ in degree one
and one class $\omega$ in degree two. In that fibre,
\[
 z\epsilon=z\omega=\epsilon^2=\epsilon\omega=\omega^2=0.
\]
Its Bockstein is
\begin{equation}\label{eq:vol2-centre-bockstein}
 \beta(\epsilon)=-\omega,
 \qquad \beta(z)=0.
\end{equation}
\end{theorem}

\begin{proof}
A degree-zero cycle obeys $a-b=ts$, so its free coordinates
are $(b,s)$. There are no incoming boundaries.
The two cuts of a one-arrow word give the cup product
\begin{equation}\label{eq:vol2-centre-cup}
 (a,b,s)(a',b',s')=(aa',bb',bs'+sa').
\end{equation}
The unit is $(1,1,0)$, and $z=(t,0,1)$ satisfies $z^2=tz$.
Every degree-zero cycle is uniquely $b\,1+s\,z$.

Since $t$ is a non-zero-divisor, $\ker D_1$ is the diagonal
copy of $R$, which equals $\im D_0$. Thus $H^1Z=0$.
The last cokernel is $R/(t)$, generated by the one-arrow
cochain $\omega:u\mapsto v$. Products of two one-arrow
cochains vanish. The left product $z\omega$ uses the target
scalar zero, while the right product uses the source scalar
$t$ and vanishes in $R/(t)$. These are all relations in
\eqref{eq:vol2-varying-ring}, by the computed cohomology groups.

Every term of \eqref{eq:vol2-varying-complex} is free over $R$,
so reduction modulo $t$ computes the derived fibre.
The reduced differential is
\[
 D_0(a,b,s)=(a-b,a-b),
 \qquad D_1=0.
\]
The cohomology dimensions are $2,1,1$ in degrees $0,1,2$.
The degree-one class is represented by $(1,0)\in Z^1$.
At $t=0$, the representative of $z$ has zero endpoint components,
so its products with $\epsilon$ and $\omega$ contain two
one-arrow cochains and vanish. Formula~\eqref{eq:vol2-centre-cup}
also gives $z^2=0$.

The lift $(1,0)$ of $\epsilon$ has differential $-t$.
The lift $(t,0,1)$ of $z$ is already closed.
This proves \eqref{eq:vol2-centre-bockstein}.
After inverting $t$, the roots $0,t$ of $z(z-t)$ separate
and the torsion class disappears.
\end{proof}

The map $b\,1+s\,z\mapsto(b+ts,b)$ identifies the degree-zero
centre with
\[
 R\times_{R/(t)}R.
\]
Its spectrum consists of two affine lines joined at the origin.
The degree-two torsion has a further effect: derived specialization
produces $\epsilon$, which tensoring the cohomology groups alone
with $k$ would omit. At the special fibre, $z$ itself is a nonzero
central class with zero endpoint components. It is carried entirely
by the homotopy on the joining complex.

\section{Gluing endpoints with their own operations}
\label{sec:vol2-dg-triangular-centre}

For differential graded endpoint algebras, the derived centres
act on the joining bimodule, and their difference gives the same
compatibility problem. Work over the field $k$. Let $A_0,A_1$
be unital differential graded algebras and let $M$ be an
$A_1$--$A_0$ bimodule. Form the two-object category
$\mathcal E_M$ with these endomorphism algebras,
$\mathcal E_M(0,1)=M$, and $\mathcal E_M(1,0)=0$.
Use derived coefficient models throughout.

Tensoring an endomorphism of the diagonal $A_0$-bimodule with
$M$ gives its right characteristic action on $M$.
The left action is defined similarly using $A_1$.
Write $\chi$ for right action minus left action.

\begin{proposition}\label{prop:vol2-dg-triangular}
There is a fibre triangle of complexes
\begin{equation}\label{eq:vol2-dg-triangular}
 \CC_k^\bullet(\mathcal E_M)\longrightarrow
 \CC_k^\bullet(A_0)\oplus\CC_k^\bullet(A_1)
 \xrightarrow{\ \chi\ }
 \RHom_{A_1\otimes_k^{\mathbb L}A_0^{\mathrm{op}}}(M,M).
\end{equation}
The cochains use product totalization over all bar lengths.
\end{proposition}

\begin{proof}
Normalize the two-sided bar relative to the object idempotents.
A composable word stays in one corner or crosses once from
zero to one. Put $\overline A_i=A_i/k1$.

The crossing cochains have factors
\[
 \Hom_k\!\left(
 (\overline A_1[1])^{\otimes p}\otimes M[1]\otimes
 (\overline A_0[1])^{\otimes q},M\right),
 \qquad p,q\ge0.
\]
The double bar resolving $M$ as an $A_1$--$A_0$ bimodule
is the direct sum of
\[
 A_1\otimes(\overline A_1[1])^{\otimes p}\otimes M
\otimes(\overline A_0[1])^{\otimes q}\otimes A_0.
\]

Its differential combines the internal differentials,
adjacent products, and the two actions on $M$.
Unit insertion contracts the augmented bar, first on one side
and then on the other. Its length filtration has induced
bimodules with homotopically projective coefficient complexes
as successive quotients. The surjective inverse-tower argument
of Proposition~\ref{prop:vol2-centre-evaluation} proves that
the whole double bar is homotopically projective, including
when the coefficient complexes are unbounded.

Hom from this resolution into $M$ computes the right-hand term
of \eqref{eq:vol2-dg-triangular}. The crossing factor has one
additional suspension, so the kernel of restriction to the
corners is this complex shifted by $-1$.
The remaining endpoint mergers give the right and left
characteristic actions with opposite signs. The resulting
short exact sequence splits on graded cochains and gives the
stated fibre triangle.
\end{proof}

This is a triangle of complexes. The cup products and braces
remain on the Hochschild cochains and their composable strings.
The last term can be replaced by an ordinary endomorphism
complex only when the chosen bimodule model computes derived Hom.
Likewise, selecting fewer boundary objects computes the centre
of that selected category. Comparison with all boundaries is
a further assertion.

\section{From coefficient actions to local bulk observables}
\label{sec:vol2-centre-geometric-comparison}

The bar evaluation gives a precise universal property over the
declared coefficient ring. In a mixed holomorphic--topological theory,
operators vary on a curve and compose in an ordered topological
direction. Their products may acquire terms supported on collision
divisors. A comparison must preserve these operations and their
coefficient modules as well as the ordinary composition of morphisms.

The physical source must be defined independently. Its fields,
action, boundary conditions, and quantization determine a complex
of bulk observables and their action on boundary insertions.
Identifying that complex with a geometric centre requires a geometric
category of boundary probes and its coefficient tensor.
The evaluation map must be compatible with the supported operations.
The comparison must also retain homotopies on morphisms and their
compatibility with composition, as the cubic example requires.

The ordinary centres above supply complete algebraic models of that
compatibility, with evaluation maps and explicit dependence on
coefficients. Their identification with a native chiral centre
remains a separate construction. A geometric issue is already visible
before quantization. Coefficients with poles along a divisor and
coefficients supported on it do not form one ordinary algebra of
functions. Their respective modules and the operations between them
must be specified.

\chapter{Local coefficients and ordered collisions}
\label{ch:vol2-coefficients}

When insertion points collide, an operation can acquire a pole or a
term supported on the collision. Both depend on the normal coordinate,
but they obey different multiplication rules. A pole permits division
by that coordinate. A supported term can be annihilated by it.
Their relation is an attachment between coefficient modules.
To place the comparison homotopies of Chapter~\ref{ch:vol2-boundary-actions}
on ordered configurations, we first construct these modules and their
derived maps.

\section{Principal parts and their residue currents}
\label{sec:vol2-principal-parts}

Work over $\C$, with one algebraic normal coordinate $z$. Put
\[
 A=\C[z],\qquad L=\C[z,z^{-1}],\qquad E=L/A.
\]
The quotient $E$ is the module of principal parts at the origin.
It has basis $t_n=z^{-n}\bmod A$, $n\geq1$, over $\C$.
Every element is a finite linear combination of these basis vectors.

Multiplication and differentiation give
\begin{equation}\label{eq:vol2-normal-jets}
 zt_1=0,\qquad zt_n=t_{n-1}\quad(n>1),\qquad
 \partial_z t_n=-nt_{n+1}.
\end{equation}
Thus $E$ retains derivatives of every normal order, even though each
of its elements is killed by a power of $z$.

Let $D=A\langle\partial_z\rangle$ be the algebra with relation
$\partial_z z-z\partial_z=1$. A left $D$-module is an $A$-module
with an operator $\nabla$ satisfying
$\nabla(am)=a'm+a\nabla m$. Such an operator is its connection
in the coordinate $z$. The usual derivative makes
\begin{equation}\label{eq:vol2-localization-sequence}
 0\longrightarrow A\longrightarrow L\longrightarrow E\longrightarrow0
\end{equation}
an exact sequence of left $D$-modules.

\begin{proposition}\label{prop:vol2-normal-module}
The map $D/Dz\to E$ sending the class of $1$ to $t_1$ is an
isomorphism of left $D$-modules. There is no unital associative algebra
extending the multiplication on $L$ and containing a nonzero element
$T$ with $zT=0$.
\end{proposition}

\begin{proof}
Using $z\partial_z=\partial_z z-1$, every element of $D/Dz$
has a representative that is a polynomial in $\partial_z$.
Its monomials act on $t_1$ by
$\partial_z^r t_1=(-1)^r r!t_{r+1}$. These images form a basis of $E$,
so the map is both injective and surjective.
In the proposed algebra, associativity and the inverse of $z$ give
$T=z^{-1}(zT)=0$, a contradiction.
\end{proof}

The relation with an analytic contact term fixes the normalization.
Write $z=x+iy$, orient the plane by $dx\wedge dy$, and set
$\partial_z=\tfrac12(\partial_x-i\partial_y)$ and
$\partial_{\bar z}=\tfrac12(\partial_x+i\partial_y)$.

A distribution acts on compactly supported smooth test functions.
Its derivatives are defined by integration by parts. Let $\delta_0$
be the distribution $\langle\delta_0,\phi\rangle=\phi(0)$.
The locally integrable function $1/z$ defines a distribution $P_1$.
For $n\geq1$, define its higher principal values by
\[
 P_n=\frac{(-1)^{n-1}}{(n-1)!}\partial_z^{n-1}P_1.
\]
They agree with $z^{-n}$ off the origin.

\begin{proposition}\label{prop:vol2-residue-current}
With these conventions,
\begin{equation}\label{eq:vol2-residue-normalization}
 \bar\partial P_n
 =\frac{\pi(-1)^{n-1}}{(n-1)!}
       \partial_z^{n-1}\delta_0\,d\bar z.
\end{equation}
The assignment of the right-hand side to $t_n$ identifies $E$ with
the left $D$-module spanned by these residue currents. The coefficient
module is placed in cochain degree zero. Its currents have Dolbeault
degree one when used in a Dolbeault complex.
\end{proposition}

\begin{proof}
For a test function $\phi$, Stokes' formula on the complement of the
disc of radius $\varepsilon$ gives
\[
 -\lim_{\varepsilon\to0}\int_{|z|>\varepsilon}
       \frac{\partial_{\bar z}\phi}{z}\,dx\,dy
 =\lim_{\varepsilon\to0}\frac1{2i}
       \int_{|z|=\varepsilon}\frac{\phi(z)}z\,dz
 =\pi\phi(0).
\]
The circle on the middle expression is counterclockwise. The outer
boundary contributes zero because $\phi$ is compactly supported.
This proves $\bar\partial P_1=\pi\delta_0d\bar z$.
Distributional derivatives commute, giving
\eqref{eq:vol2-residue-normalization} for every $n$.

The identity $z\partial_z^r\delta_0=-r\partial_z^{r-1}\delta_0$
follows by differentiating $z\delta_0=0$. It gives precisely the
multiplication rule in \eqref{eq:vol2-normal-jets}. Differentiation
gives its connection rule. The derivatives of $\delta_0$
are linearly independent: test them against $z^r$ times a smooth
function equal to one near the origin. Thus the assignment is an
isomorphism onto the displayed span.
\end{proof}

The residue is therefore a map between specified coefficient objects.
It supplies no multiplication of a principal value by an arbitrary
supported distribution. To compare operations using these objects,
we need resolutions that retain their connections.

\section{Resolutions and derived coefficient maps}
\label{sec:vol2-coefficient-resolutions}

Complexes have cohomological grading, with $V[1]^j=V^{j+1}$.
A complex of differential modules has a connection commuting with
its differential.

Define a free $A$-complex $Q_L$ by
\begin{equation}\label{eq:vol2-coefficient-telescope}
 Q_L^{-1}=\bigoplus_{n\geq0}A\alpha_n,\qquad
 Q_L^0=\bigoplus_{n\geq0}Ae_n,\qquad
 d\alpha_n=e_n-ze_{n+1}.
\end{equation}
Its augmentation sends $e_n$ to $z^{-n}$. Give it the connection
\begin{equation}\label{eq:vol2-telescope-connection}
 \nabla e_n=-ne_{n+1},\qquad
 \nabla\alpha_n=-(n+1)\alpha_{n+1},
\end{equation}
extended by the Leibniz rule. Put $Q_A=A$ in degree zero, and let
$r:Q_A\to Q_L$ send $1$ to $e_0$. Set $Q_E=\operatorname{Cone}(r)$,
with the convention
\begin{equation}\label{eq:vol2-coefficient-cone}
 \operatorname{Cone}(r)^j=Q_L^j\oplus Q_A^{j+1},\qquad
 d(\zeta,\xi)=(d_{Q_L}\zeta+r\xi,-d_{Q_A}\xi).
\end{equation}

\begin{lemma}\label{lem:vol2-coefficient-telescope}
The augmentations $Q_L\to L$ and $Q_E\to E$ are
quasi-isomorphisms of differential modules. All terms of $Q_A,Q_L,Q_E$
are free over $A$.
\end{lemma}

\begin{proof}
If a finite sum $\sum a_n(e_n-ze_{n+1})$ vanishes, its coefficient
of $e_0$ gives $a_0=0$. The successive coefficients give $a_n=0$
for every $n$. Thus the differential in
\eqref{eq:vol2-coefficient-telescope} is injective.
Its cokernel imposes $e_n=ze_{n+1}$ and is the direct limit
of $A\xrightarrow{z}A\xrightarrow{z}\cdots$, identified with $L$
by $e_n\mapsto z^{-n}$.

The connection commutes with the differential because
\[
 \nabla(e_n-ze_{n+1})
 =-(n+1)e_{n+1}+(n+1)ze_{n+2}
 =d\bigl(-(n+1)\alpha_{n+1}\bigr).
\]
It also commutes with both augmentations and with $r$.
The cone of the injection $A\to L$ has cohomology $E$ in degree
zero and no other cohomology. The cone in
\eqref{eq:vol2-coefficient-cone} consequently resolves $E$.
Its terms are direct sums of the indicated free modules.
\end{proof}

These resolutions are free over $A$. To obtain derived maps of
differential modules, we also resolve the derivative action.
For any one of the three complexes $Q$, use the two Spencer columns
\begin{equation}\label{eq:vol2-spencer}
 D\otimes_A Q\xrightarrow{s}D\otimes_A Q,\qquad
 s(P\otimes q)=P\partial_z\otimes q-P\otimes\nabla q,
\end{equation}
in degrees $-1,0$. In column degree $p$, the total differential is
$s+(-1)^p d_Q$. Denote this total complex by $P_Q$.
The Leibniz rule makes $s$ well defined under $A$-balancing.

\begin{proposition}\label{prop:vol2-spencer-maps}
The complexes $P_{Q_A},P_{Q_L},P_{Q_E}$ are bounded free
$D$-resolutions of $A,L,E$. Their complexes of $D$-linear maps,
with ordinary graded composition, form a differential graded category
computing all derived maps among these three objects.

For any complex $N$ of left $D$-modules, put
$C=\Hom_A(Q,N)$, using product totalization, and define
$\mathcal Df=\nabla_Nf-f\nabla_Q$. Then $\RHom_D(Q,N)$ is represented
by the complex with degree-$j$ term $C^j\oplus C^{j-1}$ and differential
\begin{equation}\label{eq:vol2-connection-hom}
 d(a,b)=(d_Ca,\mathcal Da-d_Cb),\qquad
 d_Cf=d_Nf-(-1)^{|f|}fd_Q.
\end{equation}
\end{proposition}

\begin{proof}
For a differential module $M$, augment the second column of
\eqref{eq:vol2-spencer} by $P\otimes m\mapsto Pm$.
Write $D$ in its right $A$-basis $1,\partial_z,\partial_z^2,\ldots$.
The highest derivative of a nonzero finite sum acquires a nonzero
term of one higher order under $s$, so $s$ is injective.
Its relations successively replace $\partial_z^r\otimes m$ by
$1\otimes\nabla^r m$. Every class in its cokernel therefore has a
unique representative $1\otimes m$, proving exactness.

Apply this resolution in each degree of $Q$. The double complex
is bounded in both degrees, so the total augmentation $P_Q\to Q$
is a quasi-isomorphism. Its terms are free over $D$ because the
terms of $Q$ are free over $A$. A bounded complex of projective
modules computes derived Hom, including when its free terms have
countably many generators.

The adjunction
$\Hom_D(D\otimes_AQ,N)=\Hom_A(Q,N)$ identifies the map induced by
$s$ with $\mathcal D$. With the sign of the second summand chosen as
in \eqref{eq:vol2-connection-hom}, the total Hom differential is
the one displayed there. For every $A$-linear $f$, the map
$\mathcal Df$ is again $A$-linear. The operator $\mathcal D$ is
$\C$-linear, obeys $\mathcal D(af)=a'f+a\mathcal Df$ for $a\in A$,
and commutes with $d_C$. Hence the displayed differential squares
to zero.

The direct sums in $Q$ give products in its Hom complexes.
Using the free $D$-resolutions on both sides supplies associative
composition and the graded Leibniz identity.
\end{proof}

The extra column in \eqref{eq:vol2-connection-hom} records compatibility
with normal differentiation. The underlying $A$-linear Hom complex
alone cannot express that compatibility. A second distinction appears
when two supported coefficients meet: their tensor product must also
be derived.

\section{Supported tensors and redundant equations}
\label{sec:vol2-supported-tensors}

Write $D(A)$ for the derived category of $A$-modules. The symbol
$\otimes_A^{\mathbb L}$ means the tensor computed after a flat
resolution. All tensors in the next proposition are over $A$, rather
than over $D$. Let $k_0=A/(z)$ be the coherent closed fibre.
It has no left $D$-action extending its $A$-action: since $z$ acts
by zero, the commutator $[\partial_z,z]$ could not act by one.
The differential module supported at the origin is instead $E$,
as in Proposition~\ref{prop:vol2-normal-module}.

\begin{proposition}\label{prop:vol2-supported-self-tensor}
In $D(A)$ one has
\begin{equation}\label{eq:vol2-supported-self-tensor}
 E\otimes_AE=0,\qquad E\otimes_A^{\mathbb L}E\simeq E[1],
 \qquad
 k_0\otimes_A^{\mathbb L}k_0\simeq k_0\otimes\Lambda(\epsilon),
 \quad |\epsilon|=-1.
\end{equation}
The last expression has zero differential.
\end{proposition}

\begin{proof}
Every element of $E$ is $z$-power torsion, whereas multiplication by
$z$ on $E$ is surjective. For a simple tensor $e\otimes e'$,
choose $m$ with $z^me=0$ and write $e'=z^mf$. Then
$e\otimes e'=z^me\otimes f=0$.

The flat resolution $[A\to L]$ of $E$ is in degrees $-1,0$.
Since $L\otimes_AE=0$, its tensor with $E$ is $[E\to0]$,
which is $E[1]$.

For the closed fibre, tensor the free resolution
$[A\xrightarrow{z}A]$ with $k_0$. Its differential becomes zero.
The Koszul multiplication gives the exterior generator $\epsilon$.
\end{proof}

Several collision equations introduce a further choice. For a
commutative ring $B$ and an element $f\in B$, let
$B_f=B[f^{-1}]$ and define
\[
 C_f=[B\longrightarrow B_f]
\]
in degrees $0,1$. For a finite list $\mathbf f=(f_1,\ldots,f_s)$,
its \emph{stable support complex} is
\begin{equation}\label{eq:vol2-stable-support}
 C_{\mathbf f}=C_{f_1}\otimes_B\cdots\otimes_B C_{f_s},
\end{equation}
with the usual tensor differential. Its terms are flat. The
projection $C_f\to B$ is the identity in degree zero and zero
in degree one.

\begin{lemma}\label{lem:vol2-redundant-support}
If $g$ lies in the radical of the ideal $(f_1,\ldots,f_s)$, then
the projection
\[
 C_{\mathbf f}\otimes_B C_g\longrightarrow C_{\mathbf f}
\]
is a quasi-isomorphism.
\end{lemma}

\begin{proof}
The complex $C_f$ is the filtered direct limit of
$[B\xrightarrow{f^n}B]$, $n\geq1$. The transition is the identity
in degree zero and multiplication by $f$ in degree one.
Consequently $C_{\mathbf f}$ is the direct limit of
$\bigotimes_i[B\xrightarrow{f_i^n}B]$.
On the $i$th two-term factor, the identity from degree one to
degree zero is a homotopy for multiplication by $f_i^n$.
Its signed tensor with the identities on the other factors gives
the same homotopy on the full tensor product.
Filtered direct limits of modules are exact, so each cohomology
class of $C_{\mathbf f}$ is annihilated by some common powers
$f_1^n,\ldots,f_s^n$.

Choose $N$ with $g^N\in(f_1,\ldots,f_s)$. The ideal
$(f_1,\ldots,f_s)^{s(n-1)+1}$ is contained in
$(f_1^n,\ldots,f_s^n)$, since each generating monomial contains
one $f_i$ at least $n$ times. A power of $g$ therefore annihilates
the given class. It follows that $C_{\mathbf f}\otimes_B B_g$
is acyclic, since localization is exact. Tensoring the triangle
$C_g\to B\to B_g$ with the flat complex $C_{\mathbf f}$ proves
the assertion.
\end{proof}

The finite Koszul complex retains different information. Take
$B=\C[u,v]$ and a third equation $w=u+v$. The finite complex
$K(u,v,w)$ is the exterior algebra over $B$ on degree-$-1$
generators $e_u,e_v,e_w$, with differential
\[
 de_u=u,\qquad de_v=v,\qquad de_w=u+v.
\]
The change of basis $\eta=e_w-e_u-e_v$ gives
\begin{equation}\label{eq:vol2-triangle-excess}
 K(u,v,u+v)=K(u,v)\otimes\Lambda(\eta),\qquad
 d\eta=0,
\end{equation}
and hence $H^*K(u,v,u+v)=\C\otimes\Lambda(\eta)$.
Indeed, $u,v$ is a regular sequence: multiplication by $u$ is
injective in $B$, and multiplication by $v$ is injective in $B/(u)$.
The two successive two-term resolutions prove that $K(u,v)$
resolves $B/(u,v)=\C$.

By contrast, $C_{u,v,u+v}\simeq C_{u,v}$ by
Lemma~\ref{lem:vol2-redundant-support}. The latter has only
degree-two cohomology,
\begin{equation}\label{eq:vol2-two-normal-parts}
 H^2C_{u,v}=
 E_2=\frac{\C[u^{\pm1},v^{\pm1}]}
              {\C[u^{\pm1},v]+\C[u,v^{\pm1}]},
\end{equation}
with basis $u^{-m}v^{-n}$, $m,n\geq1$.
To see the concentration, write $C_{u,v}$ as the tensor over $\C$
of the one-variable complexes for $u$ and $v$. Each has its
principal-parts cohomology in degree one.

The vanishing of $(E_2)_{u+v}$ also has a direct calculation:
\[
 \frac1{u^mv^n}=
 \frac1{u+v}\left(\frac1{u^{m-1}v^n}
                       +\frac1{u^mv^{n-1}}\right).
\]
Repeated application decreases $m+n$ until every term is a power
of $(u+v)^{-1}$ times a fraction regular in one normal coordinate.
That fraction represents zero in $E_2$, so the term vanishes in
$(E_2)_{u+v}$.
The odd class $\eta$ in \eqref{eq:vol2-triangle-excess} belongs
to the finite closed-intersection complex. It does not supply an
extra class in the stable support complex.

\section{Differential modules on configuration spaces}
\label{sec:vol2-geometric-coefficients}

The coordinate calculations extend over a smooth separated complex
algebraic curve $X$. For a smooth scheme $Y$, let $\Dcal_Y$ be its
sheaf of differential operators. Let $\operatorname{DMod}(Y)$
denote the stable enhancement of the unbounded derived category
of right $\Dcal_Y$-modules with quasi-coherent underlying
$\Ocal_Y$-modules. The enhancement retains mapping complexes
and coherent homotopies. We use the enhanced differential-module
formalism of Francis--Gaitsgory \cite{FG12}, \S1.4, including
extraordinary inverse image, direct image and external product,
with their composition coherences.

Right modules give a coordinate-independent form of our local
left modules. If $M$ is a left differential module, the associated
right module is $\omega_Y\otimes_{\Ocal_Y}M$, where $\omega_Y$
is the line of top forms. In one coordinate its action satisfies
\[
 (dz\otimes m)\partial_z=-dz\otimes\nabla m.
\]
For a nonconstant form, the derivative also acts on its coefficient.
This is the transposition that fixes functions, negates coordinate
derivatives and reverses products. It preserves the normal derivatives
in $E$.

Here are the functors and shifts used below. For a map $f:Z\to Y$
of smooth schemes, with dimensions $d_Z,d_Y$, the transfer modules are
\[
 \Dcal_{Z\to Y}
 =\Ocal_Z\otimes_{f^{-1}\Ocal_Y}f^{-1}\Dcal_Y,
 \qquad
 \Dcal_{Y\leftarrow Z}
 =\omega_Z\otimes_{\Ocal_Z}\Dcal_{Z\to Y}
       \otimes_{f^{-1}\Ocal_Y}f^{-1}\omega_Y^{-1}.
\]
The first has a left $\Dcal_Z$-action and a right
$f^{-1}\Dcal_Y$-action. Transposition gives the second a left
$f^{-1}\Dcal_Y$-action and a right $\Dcal_Z$-action.

For right-module complexes $M$ on $Z$ and $N$ on $Y$, set
\begin{align}
 f_+M&=Rf_*^{\Ocal}
       (M\otimes_{\Dcal_Z}^{\mathbb L}\Dcal_{Z\to Y}),
       \label{eq:vol2-dmodule-direct}\\
 f^!N&=(f^{-1}N\otimes_{f^{-1}\Dcal_Y}^{\mathbb L}
                  \Dcal_{Y\leftarrow Z})[d_Z-d_Y].
       \label{eq:vol2-dmodule-inverse}
\end{align}
Here $f^{-1}$ is inverse image of sheaves, and $Rf_*^{\Ocal}$
is derived sheaf direct image. External product, denoted $\boxtimes$,
places two differential modules on independent factors.

For the structure map $p_Y:Y\to\Spec\C$, our normalization is
\begin{equation}\label{eq:vol2-dualizing-shift}
 J_Y=p_Y^!\C=\omega_Y[d_Y],\qquad
 J_Y\boxtimes J_Z=J_{Y\times Z}.
\end{equation}

For an open immersion $j$, write $j_*$ for $j_+$. Its functor
$j^!$ is restriction. For a closed immersion $i$, write $i_*=i_+$.
The functor $i_*$ retains all normal derivatives along the closed
subspace.

Differential modules on a possibly singular reduced
closed subspace are defined by the corresponding supported modules
in smooth embeddings, using Kashiwara's equivalence. Thus the
open--closed recollement identities apply to unions of diagonals:
\begin{equation}\label{eq:vol2-open-recollement}
 i_*i^!N\longrightarrow N\longrightarrow j_*j^!N,
 \qquad i^!j_*=0,\qquad j^!j_*=\id.
\end{equation}
These identities, and the coherent composition laws of the functors,
are part of the specified coefficient formalism.

The cut product will first restrict to a disjoint locus and then
apply $j_*$. Its compatibility with a collision rests on the
following consequence of recollement.

\begin{lemma}\label{lem:vol2-open-basechange}
For a Cartesian square with $j$ an open immersion,
\[
 \begin{tikzcd}
 U'\arrow[r,"f_U"]\arrow[d,"j'"']&U\arrow[d,"j"]\\
 Y'\arrow[r,"f"']&Y,
 \end{tikzcd}
 \qquad
 f^!j_*M\simeq j'_*f_U^!M.
\]
These equivalences are coherent under composition of $f$.
\end{lemma}

\begin{proof}
Let $i:Z\hookrightarrow Y$ and $i':Z'\hookrightarrow Y'$
be the reduced closed complements. The map $fi'$ factors as $ih$.
Transitivity and \eqref{eq:vol2-open-recollement} give
\[
 i'^!f^!j_*M=h^!i^!j_*M=0.
\]
The recollement triangle on $Y'$ therefore identifies $f^!j_*M$
with $j'_*j'^!f^!j_*M$. On the open set $U'$, transitivity
identifies its restriction with $f_U^!M$.

For a composite $f$, all such comparisons lie in the fully faithful
image of $j'_*$. Their restrictions are the given coherent
composition maps for extraordinary inverse image. Full faithfulness
lifts those maps and all their higher homotopies uniquely.
\end{proof}

\section{Ordered words and consecutive cuts}
\label{sec:vol2-word-coefficients}

An ordered insertion retains its place even when its point coincides
with another point. The word $(x,y,x)$, with $x\ne y$, is an elementary
example. Its two consecutive cuts are $(x)\mid(y,x)$ and
$(x,y)\mid(x)$. Neither separates disjoint sets of points.

The intervening operation cannot be deleted. For the matrix units
$E_{ij}$ on $\C^2$, which send the $j$th basis vector to the $i$th,
\[
 E_{12}E_{21}E_{12}=E_{12},\qquad E_{12}^2=0.
\]
We accordingly retain words and allow adjacent positions to collapse
without identifying a word with its unordered set of values.

Let $[n]=\{1,\ldots,n\}$ with its usual order, and let
$\operatorname{OrdSurj}_+$ be the category of these nonempty ordered
sets and monotone surjections. A surjection
$\alpha:[n]\twoheadrightarrow[m]$ defines
\[
 d_\alpha:X^m\longrightarrow X^n,\qquad
 d_\alpha(x_1,\ldots,x_m)_i=x_{\alpha(i)}.
\]
Write $D(\C)$ for complexes of complex vector spaces up to
quasi-isomorphism, with their stable enhancement. The category of
\emph{word coefficients} is
\begin{equation}\label{eq:vol2-word-category}
 \mathcal W_X=D(\C)\times
   \varprojlim_{\operatorname{OrdSurj}_+}\operatorname{DMod}(X^n),
\end{equation}
where the transition from $n$ to $m$ is $d_\alpha^!$.
An object $F$ consists of an empty coefficient $F_0$, coefficients
$F_n$ on $X^n$, equivalences $d_\alpha^!F_n\simeq F_m$, and
their coherent compatibility under composition. The empty
coefficient has no collision arrow to a nonempty one.

For $0<p<n$, define the cut open and its inclusion by
\begin{equation}\label{eq:vol2-cut-open}
 U_{p,n}=\{(x_1,\ldots,x_n):x_i\ne x_j
                         \text{ whenever }i\leq p<j\},
 \qquad j_{p,n}:U_{p,n}\hookrightarrow X^n.
\end{equation}
For $p=0,n$, put $U_{p,n}=X^n$ and let $j_{p,n}$ be the identity.
Define
\begin{equation}\label{eq:vol2-cut-product}
 (F\star_{\mathrm{cut}}G)_n
 =\bigoplus_{p=0}^n j_{p,n,*}j_{p,n}^!
                  (F_p\boxtimes G_{n-p}).
\end{equation}
An empty factor means scalar tensoring. In particular,
$(F\star_{\mathrm{cut}}G)_0=F_0\otimes_\C G_0$.

\begin{theorem}\label{thm:vol2-cut-product}
Formula~\eqref{eq:vol2-cut-product} defines an associative monoidal
product on $\mathcal W_X$. Its unit $K$ has $K_0=\C$ and
$K_n=0$ for $n>0$. Its associativity maps are the identifications
of the iterated products with the sums over consecutive partitions
whose blocks have pairwise disjoint values in $X$.
\end{theorem}

\begin{proof}
First pull a summand of \eqref{eq:vol2-cut-product} back along
$d_\alpha$. If its cut lies inside a fibre of $\alpha$, two
positions on opposite sides of the cut are forced to have the same
value. Its inverse-image open is empty, and
Lemma~\ref{lem:vol2-open-basechange} makes the summand zero.
Otherwise the cut lies between fibres, including either end, hence
determines a unique cut $q$ of $[m]$. Its inverse-image open is
$U_{q,m}$. External-product compatibility, the word equivalences
of $F,G$, and the same lemma identify the pulled-back summand
with the $q$th summand of $(F\star_{\mathrm{cut}}G)_m$.

The coherence in the
lemma supplies the coherence for composite surjections, so the
product is a word coefficient system.

For three factors, either parenthesization is a direct sum over
$0\leq p\leq q\leq n$. Each summand starts with the external
product on the three consecutive blocks. Restrict it to the locus
where values in different blocks are distinct, then apply open
direct image. Successive cut conditions impose precisely this
common open. Base change and composition of open direct images
identify the two expressions. Empty blocks impose no condition.

The same description holds for any number of factors. All
associativity diagrams, including the pentagon, restrict to the
identity of this common expression. The functorial coherences
therefore make them commute. With $K$, only the empty-block
summand survives on the appropriate side, giving the unit maps
and their compatibility.
\end{proof}

No permutation of the blocks has been introduced. In particular,
this construction supplies an associative product without a chosen
symmetry. The empty unit is distinct from a local unit field.
A local vacuum in a boundary algebra is extra data on nonempty
configurations, whereas $K_1=0$.

\begin{example}\label{ex:vol2-outer-cut}
At three positions,
\[
 U_{1,3}=X^3\setminus(\Delta_{12}\cup\Delta_{13}),\qquad
 U_{2,3}=X^3\setminus(\Delta_{13}\cup\Delta_{23}),
\]
where $\Delta_{ij}$ is the locus $x_i=x_j$. Their union has
complement exactly $\Delta_{13}$. If $x_1\ne x_3$ and the first
cut fails, then $x_1=x_2$, so the second cut succeeds.
Conversely, $x_1=x_3$ defeats both cuts. Thus every nonadjacent
equality crossing a cut is required already in the first
three-position calculation.

On an affine coordinate chart,
$u=z_1-z_2$ and $v=z_2-z_3$ give the outer equation $u+v=0$.
The stable and finite complexes of
Section~\ref{sec:vol2-supported-tensors} compute different
coefficient constructions on this same arrangement.
\end{example}

Call $F$ \emph{reduced} if $F_0=0$.
For reduced objects $F^{(1)},\ldots,F^{(r)}$, every nonzero
summand of their iterated cut product at $n$ has $r$ nonempty
consecutive blocks. Hence
\begin{equation}\label{eq:vol2-positive-word-length}
 (F^{(1)}\star_{\mathrm{cut}}\cdots
          \star_{\mathrm{cut}}F^{(r)})_n=0\qquad(r>n).
\end{equation}
This bound concerns word length. It makes tensor-length constructions
finite at a fixed configuration size. It does not bound the order
of a pole or the normal derivatives in a supported coefficient.

\section{Comparing attachments}
\label{sec:vol2-coefficient-attachments}

The localization sequence \eqref{eq:vol2-localization-sequence}
ties regular, meromorphic and supported coefficients together.
A comparison must respect this attachment in addition to the
separate objects. Let $r:U\to Z$ and $r':U'\to Z'$ be cochain
maps of complexes of left $D$-modules, and take $D$-linear Hom
complexes. For derived maps among $A,L,E$, use the bounded free
resolutions of Proposition~\ref{prop:vol2-spencer-maps}.

Give each cone the two-step filtration
$0\subset Z\subset\operatorname{Cone}(r)$, with the convention
of \eqref{eq:vol2-coefficient-cone}. A degree-zero map preserving
these filtrations has the form
\begin{equation}\label{eq:vol2-attachment-map}
 (\zeta,\xi)\longmapsto(g\zeta+H\xi,f\xi),\qquad
 |f|=|g|=0,\quad |H|=-1.
\end{equation}
Here $f:U\to U'$, $g:Z\to Z'$ and $H:U\to Z'$.

\begin{lemma}\label{lem:vol2-attachment-homotopy}
The map \eqref{eq:vol2-attachment-map} is closed if and only if
$f,g$ are cochain maps and
\begin{equation}\label{eq:vol2-attachment-homotopy}
 d_{\Hom}H=gr-r'f.
\end{equation}
The complex of filtration-preserving cone maps is isomorphic,
with the shifted-coordinate signs, to the homotopy fibre of
\[
 \Hom(U,U')\oplus\Hom(Z,Z')\longrightarrow\Hom(U,Z'),
 \qquad(f,g)\longmapsto gr-r'f.
\]
For closed degree-zero maps, composition is
\begin{equation}\label{eq:vol2-attachment-composition}
 (f_2,g_2,H_2)(f_1,g_1,H_1)
 =(f_2f_1,g_2g_1,g_2H_1+H_2f_1).
\end{equation}
\end{lemma}

\begin{proof}
Insert \eqref{eq:vol2-attachment-map} into the cone differential
$(\zeta,\xi)\mapsto(d_Z\zeta+r\xi,-d_U\xi)$. The two diagonal entries in the
chain-map equation say $df=dg=0$. The off-diagonal entry says
$d_{Z'}H+Hd_U=gr-r'f$, which is
\eqref{eq:vol2-attachment-homotopy}.

In general degree $j$, put a factor $(-1)^j$ on the $U[1]\to U'[1]$
entry and a factor $(-1)^j$ on the $U[1]\to Z'$ entry.
In these coordinates the differential is
\[
 (f,g,H)\longmapsto
 (df,dg,gr-r'f-dH).
\]
This is the homotopy-fibre differential, since
$|H|=j-1$. Finally, multiplication of the degree-zero upper triangular
block matrices gives \eqref{eq:vol2-attachment-composition} and
its associativity.
\end{proof}

Equations~\eqref{eq:vol2-connection-hom} and
\eqref{eq:vol2-attachment-homotopy} retain two kinds of compatibility:
with differentiation and with the residue attachment. The cut product
adds the geometry of ordered insertion. Together they specify
coefficient objects and comparison complexes for local operations.

An action on them must still give operations in every arity and
homotopies compatible with substitution and collision restriction.
The construction of a field theory must, independently, supply its
fields, differential, interaction and quantum observables.
It must also give a map from those observables to the geometric actions.
The next calculation
begins this theory-dependent part with finite Hamiltonian jets.

\chapter{Finite Hamiltonian jets and relative quantization}
\label{ch:vol2-hamiltonian-jets}

A Hamiltonian determines an infinitesimal motion of a symplectic space.
To study a boundary fixed at a chosen point, we keep the Hamiltonians
whose vector fields vanish there. Their finite jets give Lie algebras
with an explicit quantization. The resulting associative algebras have
relative Hochschild complexes, so they also give finite models of the
coherent actions considered in Chapter~\ref{ch:vol2-boundary-actions}.

The lowest nonzero jet is already instructive. Quadratic Hamiltonians form
$\mathfrak{sl}_2$; cubic Hamiltonians supply its four-dimensional
representation; the bracket of two cubics first survives at quartic
order. We shall construct these algebras, their Poisson cochains and
their finite Batalin--Vilkovisky actions before comparing the
quantized algebra with its Hochschild cochains.

\section{Hamiltonians fixing a point}

Work over $\C$ on the formal plane with coordinates $z_1,z_2$ and
symplectic form $\omega=dz_1\wedge dz_2$. Set
$\mathfrak m=(z_1,z_2)\subset\C\llbracket z_1,z_2\rrbracket$.
Our Poisson convention is
\begin{equation}\label{eq:vol2-hamiltonian-bracket}
 \{f,g\}=\partial_{z_1}f\,\partial_{z_2}g
             -\partial_{z_2}f\,\partial_{z_1}g.
\end{equation}
Thus the vector field $X_f(g)=\{f,g\}$ vanishes at the origin
exactly when the linear part of $f$ vanishes. Constants give the zero
vector field. Choosing the representative with zero constant term
identifies the pointed Hamiltonian Lie algebra with
$\h=\mathfrak m^2$.

For the quadratic polynomials
\[
 e=\tfrac12z_1^2,\qquad f=-\tfrac12z_2^2,\qquad h=-z_1z_2,
\]
direct differentiation gives
\begin{equation}\label{eq:vol2-quadratic-brackets}
 \{h,e\}=2e,\qquad \{h,f\}=-2f,\qquad \{e,f\}=h.
\end{equation}
They are the Chevalley generators of $\mathfrak{sl}_2$.

Write $R_r$ for the homogeneous polynomials of degree $r$.
Each derivative lowers degree by one, so
$\{R_r,R_s\}\subset R_{r+s-2}$.
The filtration
\[
 F^p\h=\mathfrak m^{p+2}\quad(p\geq0)
\]
is therefore additive:
$\{F^p\h,F^q\h\}\subset F^{p+q}\h$.
It is separated and complete in the formal power-series topology.
For $d\geq2$, define
\begin{equation}\label{eq:vol2-finite-jet}
 \h_{\leq d}=\h/\mathfrak m^{d+1}
             =\bigoplus_{r=2}^{d}R_r.
\end{equation}
The equality on the right specifies the polynomial representatives;
their bracket is projected to degrees at most $d$.

\begin{proposition}[Finite Hamiltonian jets]\label{prop:vol2-finite-jets}
The bracket makes $\h_{\leq d}$ a Lie algebra of dimension
\[
 \dim\h_{\leq d}=\frac{d(d+3)}2-2.
\]
Its quadratic subalgebra is $\mathfrak{sl}_2$, and
$R_r$ is its irreducible representation of highest weight $r$.
\end{proposition}

\begin{proof}
The ideal $\mathfrak m^{d+1}$ is stable under bracketing with
$\mathfrak m^2$. Hence Jacobi descends from the Poisson bracket.
In coordinates, this bracket is induced by the constant bivector
$\partial_{z_1}\wedge\partial_{z_2}$; commuting partial derivatives
cancel the Jacobiator. Its monomial formula is
\begin{equation}\label{eq:vol2-monomial-bracket}
 \{z_1^az_2^b,z_1^cz_2^s\}
 =(as-bc)z_1^{a+c-1}z_2^{b+s-1},
\end{equation}
where a term with zero coefficient is zero, including when a displayed
exponent would be negative. For three exponent vectors
$\alpha,\beta,\gamma$ and $\epsilon=(1,1)$, the Jacobi coefficient is
\[
 \det(\alpha,\beta)\det(\alpha+\beta-\epsilon,\gamma)
 +\text{cyclic permutations}=0.
\]
The terms quadratic in the pairwise determinants cancel. The remaining
terms vanish by the two-dimensional identity
\[
 \det(\alpha,\beta)\gamma+\det(\beta,\gamma)\alpha
 +\det(\gamma,\alpha)\beta=0.
\]
The dimension is $\sum_{r=2}^d(r+1)=d(d+3)/2-2$.

Put $u_j=z_1^{r-j}z_2^j$ for $0\leq j\leq r$. The quadratic action is
\begin{equation}\label{eq:vol2-weight-action}
 \{h,u_j\}=(r-2j)u_j,\qquad
 \{e,u_j\}=ju_{j-1},\qquad
 \{f,u_j\}=(r-j)u_{j+1}.
\end{equation}
Terms outside the index range are zero. A nonzero invariant subspace
contains a weight vector because the eigenvalues of $h$ are distinct.
Repeated application of $e$ reaches $u_0$, and repeated application of
$f$ then reaches every $u_j$. This proves irreducibility.
\end{proof}

The pointed condition is needed for these quotients. The affine algebra
\[
 \h^{\mathrm{aff}}=\C\llbracket z_1,z_2\rrbracket/\C
\]
also contains linear Hamiltonians, which generate translations. In it,
\[
 \{z_1,z_2^{d+1}\}=(d+1)z_2^d.
\]
Consequently $\mathfrak m^{d+1}$ is no longer an ideal. A moving
boundary or a centre-of-mass coordinate requires an affine construction
in addition to the pointed jet quotient.

\section{The first nonlinear brackets}

For $d=2$ the quotient is $\mathfrak{sl}_2$ itself. For $d=3$, set
$v_j=z_1^{3-j}z_2^j$ for $0\leq j\leq3$.
Equation~\eqref{eq:vol2-weight-action} gives the quadratic action on
their span.

\begin{proposition}[The cubic jet]\label{prop:vol2-cubic-jet}
The degree-three quotient is the semidirect product
\[
 \h_{\leq3}=\mathfrak{sl}_2\ltimes R_3,
 \qquad [R_3,R_3]=0.
\]
\end{proposition}

\begin{proof}
A bracket of two cubics has degree four and vanishes in this quotient.
The brackets with quadratics are exactly the action in
\eqref{eq:vol2-weight-action}.
\end{proof}

At $d=4$, put $w_k=z_1^{4-k}z_2^k$ for $0\leq k\leq4$.
The cubic bracket becomes
\begin{equation}\label{eq:vol2-cubic-quartic}
 \{v_i,v_j\}=3(j-i)w_{i+j-1}.
\end{equation}
When the index is outside $0,\ldots,4$, the coefficient is zero.
In particular,
\[
 \{v_0,v_3\}=9w_2,\qquad \{v_1,v_2\}=3w_2.
\]
The first equality differentiates $z_1^3$ and $z_2^3$.
The second is
$4z_1^2z_2^2-z_1^2z_2^2=3z_1^2z_2^2$.
The two cubic pairs reach the same quartic monomial with coefficients
in the ratio $3:1$.

\begin{proposition}[The quartic jet]\label{prop:vol2-quartic-jet}
On the vector space $R_2\oplus R_3\oplus R_4$, all brackets are
determined by \eqref{eq:vol2-quadratic-brackets},
\eqref{eq:vol2-weight-action} for $r=3,4$,
\eqref{eq:vol2-cubic-quartic}, and
\[
 [R_3,R_4]=0,\qquad [R_4,R_4]=0.
\]
This Lie algebra is $\h_{\leq4}$. Its quadratic and cubic elements
generate it; its ideal $R_3\oplus R_4$ is two-step nilpotent with
commutator subalgebra $R_4$.
\end{proposition}

\begin{proof}
The omitted brackets have degrees five and six, so they vanish in the
quotient. Every quartic is a cubic bracket: use the pairs
$(i,j)=(0,1),(0,2),(0,3),(1,3),(2,3)$ to obtain
$w_0,w_1,w_2,w_3,w_4$, respectively, with nonzero coefficients.
These formulas exhaust the brackets of a basis. Jacobi follows from
Proposition~\ref{prop:vol2-finite-jets}. They also give the stated
commutator and its centrality inside the ideal.
\end{proof}

The adjoint trace will measure the divergence of the finite cotangent
action. The homogeneous decomposition already determines this trace.

\begin{proposition}[Unimodularity]\label{prop:vol2-jet-unimodular}
For every $d\geq2$ and $x\in\h_{\leq d}$,
$\operatorname{Tr}(\operatorname{ad}_x)=0$.
\end{proposition}

\begin{proof}
A homogeneous Hamiltonian of degree greater than two strictly raises
polynomial degree. Its adjoint matrix is strictly triangular in the
decomposition \eqref{eq:vol2-finite-jet}. For a quadratic Hamiltonian,
use the basis $e,f,h$. On each $R_r$, the matrices of $e$ and $f$ have
zero diagonal, while the trace of $h$ is
$\sum_{j=0}^r(r-2j)=0$. Linearity proves the assertion.
\end{proof}

\section{Chevalley cochains of the linear Poisson structure}

Fix a finite-dimensional complex Lie algebra $\flie$, with basis
$e_a$ and bracket $[e_a,e_b]=f^c{}_{ab}e_c$; repeated indices are
summed. We shall apply the constructions to $\flie=\h_{\leq d}$.
Use cohomological grading, with $(W[k])^n=W^{n+k}$.
In particular, a vector of degree zero has degree one in $W[-1]$.

Let $\flie^\vee$ carry the coadjoint action
$(\operatorname{ad}_x^*\lambda)(y)=-\lambda([x,y])$.
Define the graded Lie algebra
\[
 \mathfrak q(\flie)=\flie\ltimes\flie^\vee[-1]
\]
with zero internal differential and abelian second summand.
Its elements in the second summand have degree one.
For a finite graded Lie algebra $\mathfrak q$, we use polynomial
Chevalley cochains
$C^*_{\CE}(\mathfrak q)=\Sym(\mathfrak q^\vee[-1])$
with differential dual to the bracket. Here $\Sym$ means the graded
symmetric algebra, without completion in its even generators.

\begin{proposition}[Chevalley--Poisson complex]\label{prop:vol2-ce-poisson}
There is a natural isomorphism of commutative differential graded
algebras
\[
 C^*_{\CE}(\mathfrak q(\flie))
 \cong C^*(\flie,\Sym\flie)
 =\Sym\flie\otimes\Lambda^*\flie^\vee.
\]
On the right the differential is the Chevalley differential for the
adjoint action. It is also the linear Poisson differential on
polynomial polyvector fields on $\flie^*$.
\end{proposition}

\begin{proof}
Dualizing and shifting the degree-one summand gives
$(\flie^\vee[-1])^\vee[-1]=\flie$ in degree zero.
The degree-zero summand gives exterior generators $c^a$ of degree one.
Let $x_a\in\Sym\flie$ denote the linear function
$x_a(\lambda)=\lambda(e_a)$. Dualizing the semidirect bracket gives
\begin{equation}\label{eq:vol2-ce-coordinate}
 d_{\CE}c^c=-\tfrac12 f^c{}_{ab}c^ac^b,
 \qquad d_{\CE}x_a=f^c{}_{ba}c^b x_c.
\end{equation}
These are precisely the bracket and action terms of the Chevalley
differential. On a polynomial $F$ the second formula reads
\[
 d_{\CE}F=c^a\{x_a,F\}_{\mathrm{lin}},
 \qquad \{x_a,x_b\}_{\mathrm{lin}}=f^c{}_{ab}x_c.
\]
Identify $c^a$ with the constant vector field $\partial_{x_a}$.
Use the Schouten bracket extending $[X,F]=X(F)$ by graded
skew-symmetry and the graded derivation rule.
For the bivector
$\pi=\tfrac12 f^c{}_{ab}x_c\partial_{x_a}\wedge\partial_{x_b}$,
the Schouten differential $[\pi,-]$ gives both formulas in
\eqref{eq:vol2-ce-coordinate}. Since both operators are derivations
of the exterior product, they agree on all polyvectors. Jacobi is
$[\pi,\pi]=0$, and hence $d_{\CE}^2=0$.
\end{proof}

\section{The finite BV action}

The preceding complex has degree-zero polynomial coordinates $x_a$.
The BV construction uses the different graded space
$T^*[-1](\flie[1])$. Its polynomial functions are
\[
 \mathcal B(\flie)=\Lambda(c^1,\ldots,c^m)
                  \otimes\C[b_1,\ldots,b_m],
 \qquad |c^a|=1,\quad |b_a|=-2,
\]
where $m=\dim\flie$. Thus the cotangent coordinates $b_a$ are
antifields of degree $-2$. Fix left derivatives in the odd variables
and set
\begin{equation}\label{eq:vol2-bv-conventions}
 \Delta=\sum_a\partial_{b_a}\partial_{c^a},\qquad
 \{F,G\}_{\mathrm{BV}}
 =(-1)^{|F|}\bigl(\Delta(FG)-\Delta(F)G
                    -(-1)^{|F|}F\Delta(G)\bigr).
\end{equation}
The bracket has degree one, with
$\{b_a,c^b\}_{\mathrm{BV}}=\delta_a^b$ and
$\{c^b,b_a\}_{\mathrm{BV}}=-\delta_a^b$.
Commutation of the even derivatives and anticommutation of the odd
ones give $\Delta^2=0$.
These Darboux coordinates specify a constant Berezinian up to a
nonzero scalar; its normalization does not change $\Delta$.
The quantum master equation for this measure is proved in
Theorem~\ref{thm:vol2-finite-qme}.

With these conventions the degree-zero cotangent action is
\begin{equation}\label{eq:vol2-bv-action}
 S=-\tfrac12 b_c f^c{}_{ab}c^ac^b.
\end{equation}
In particular, $\{S,c^a\}_{\mathrm{BV}}=d_{\CE}c^a$.

\begin{theorem}[Finite quantum master equation]\label{thm:vol2-finite-qme}
For every finite-dimensional Lie algebra $\flie$, the action
\eqref{eq:vol2-bv-action} satisfies
\[
 \tfrac12\{S,S\}_{\mathrm{BV}}=0,\qquad
 \Delta S=\sum_b\operatorname{Tr}(\operatorname{ad}_{e_b})c^b.
\]
Consequently, for $\flie=\h_{\leq d}$ and a formal parameter
$\hbarq$ of degree zero,
\[
 \tfrac12\{S,S\}_{\mathrm{BV}}+\hbarq\Delta S=0.
\]
This identity holds in $\mathcal B(\h_{\leq d})\llbracket\hbarq\rrbracket$.
\end{theorem}

\begin{proof}
Put $Q^a=-\tfrac12 f^a{}_{bc}c^bc^c$, so that $S=b_aQ^a$.
The coordinate bracket gives
\[
 \tfrac12\{S,S\}_{\mathrm{BV}}
 =\sum_{a,k}b_kQ^a\partial_{c^a}Q^k.
\]
The coefficient of $b_k$ is $d_{\CE}^2c^k$, which vanishes by
Jacobi. For the second identity, left differentiation gives
\[
 \Delta S=-\sum_{a,b}f^a{}_{ab}c^b
          =\sum_{a,b}f^a{}_{ba}c^b.
\]
The last coefficient is the trace of $\operatorname{ad}_{e_b}$.
Proposition~\ref{prop:vol2-jet-unimodular} makes it zero for every
finite Hamiltonian jet.
\end{proof}

The trace condition is substantive. In the two-dimensional Lie algebra
with basis $u,v$ and $[u,v]=v$, one has
\[
 S=-b_vc^uc^v,\qquad \Delta S=c^u.
\]
The classical master equation still holds, but this constant-measure
action fails the quantum master equation. For Hamiltonian jets,
unimodularity removes this finite-dimensional divergence.
A local field theory has further analytic and renormalization data;
the finite identity alone does not construct its quantum observables.

\section{Rees multiplication and interval operations}

Quantization replaces commuting linear functions on $\flie^*$ by
generators whose commutator is $\hbarq$ times the Lie bracket.
For $N\geq0$, let
\[
 R_N=\C[\hbarq]/(\hbarq^{N+1}),\qquad
 \flie_N=R_N\otimes_\C\flie,
 \qquad [x,y]_N=\hbarq[x,y].
\]
Here $N$ is the coefficient order, independent of the jet order $d$
or a Chan--Paton rank. Define the augmented $R_N$-algebra
\[
 A_N=U_{R_N}(\flie_N)
 =T_{R_N}(\flie_N)/(xy-yx-\hbarq[x,y]),
\]
whose augmentation sends every $x\in\flie$ to zero.
The Poincar\'e--Birkhoff--Witt (PBW) theorem has the following form
over this coefficient ring.

\begin{proposition}[PBW over the finite coefficient ring]\label{prop:vol2-rees-pbw}
Choose an ordered basis of $\flie$. Its ordered monomials form an
$R_N$-basis of $A_N$. If $S_0=\Sym_\C\flie$, symmetrization is a
filtered $R_N$-module isomorphism
\[
 \operatorname{sym}_N:S_0\otimes R_N\longrightarrow A_N,
 \qquad x_1\cdots x_r\longmapsto
 \frac1{r!}\sum_{\sigma\in S_r}x_{\sigma(1)}\cdots x_{\sigma(r)}.
\]
It intertwines the adjoint action through
\begin{equation}\label{eq:vol2-scaled-pbw-action}
 [x,\operatorname{sym}_N(F)]_{A_N}
 =\hbarq\operatorname{sym}_N(\operatorname{ad}_xF).
\end{equation}
\end{proposition}

\begin{proof}
Replace every inverted adjacent pair $yx$ by
$xy+\hbarq[y,x]$. In the first term the number of inversions
decreases; the second has shorter length. Induction on length and
then inversions proves termination. Disjoint replacements commute.
For an overlapping triple $z>y>x$, the difference between the two
reductions, after the shorter words have been ordered, is
$\hbarq^2$ times
\[
 [[z,y],x]+[[y,x],z]+[[x,z],y]=0.
\]
Induction therefore makes the ordered result independent of every
choice of replacement. This defines a normal-form map killing the
relations and equal to the identity on ordered monomials. Those
monomials are consequently a basis, even though $R_N$ has nilpotents.

Give generators PBW filtration degree one. Symmetrization induces
the identity on the associated graded polynomial algebra. Subtracting
the highest-degree symbol and descending in degree constructs its
inverse. This uses only finite degrees and the invertibility of
factorials in $R_N$. Finally, commuting $x$ past each factor of a
symmetrized word gives one term $\hbarq[x,x_i]$ in each position.
Their sum is \eqref{eq:vol2-scaled-pbw-action}.
\end{proof}

Transporting multiplication by $\operatorname{sym}_N$ gives the
associative PBW product on $S_0\otimes R_N$. Its reduction modulo
$\hbarq$ is the ordinary product. For $N\geq1$, the first commutator
is the linear Poisson bracket. Indeed, the coefficient of $\hbarq$ in a
commutator is a biderivation of the commutative reduction; its value
on $x,y\in\flie$ is $[x,y]$. These values determine the biderivation
on the polynomial algebra.

For disjoint intervals $I_1<\cdots<I_r$ contained in an interval $J$,
assign the operation
\[
 A_N^{\otimes_{R_N}r}\longrightarrow A_N,
 \qquad a_1\otimes\cdots\otimes a_r\longmapsto a_1\cdots a_r.
\]
Use the unit for an empty collection and the identity for an inclusion
of one interval. Associativity makes these operations compatible with
nested collections, and their independence of interval lengths makes
them locally constant. They supply the algebraic ordered interval
operations expected of a boundary quantization.
For a cotangent field theory on a half-plane, fix the cotangent field
at the boundary and leave the Lie algebra field free. Identifying the
resulting boundary observables with these interval operations requires
a map from that theory's quantum
observables, compatible with interval inclusions and fusion, and a
proof that the map is a quasi-isomorphism. PBW constructs the target
algebra, while this realization remains a separate comparison.

\section{The relative Hochschild comparison}

The coefficient ring affects the action complex. For the quadratic
jet, $[h,e]=2e$, whereas $[h,e]_{A_N}=2\hbarq e$.
Modulo $\hbarq$, the Hochschild differential of the zero-cochain $e$
vanishes because the algebra is commutative. The ordinary Chevalley
differential has $(d_{\CE}e)(h)=2e$. Thus it is the scaled differential
$\hbarq d_{\CE}$ that can compare with relative Hochschild cochains.

Write $\bar A_N$ for the augmentation ideal.
Use normalized relative Hochschild cochains
\[
 C^n_{\HH}(A_N/R_N,A_N)
 =\operatorname{Hom}_{R_N}(\bar A_N^{\tensor_{R_N}n},A_N)
\]
with differential
\[
\begin{split}
 (bF)(a_1,\ldots,a_{n+1})
 ={}&a_1F(a_2,\ldots,a_{n+1})\\
 &+\sum_{i=1}^n(-1)^iF(a_1,\ldots,a_i a_{i+1},\ldots,a_{n+1})\\
 &+(-1)^{n+1}F(a_1,\ldots,a_n)a_{n+1}.
\end{split}
\]
Let $d_{\CE}$ act on
$C^n(\flie,S_0)=\operatorname{Hom}_\C(\Lambda^n\flie,S_0)$
by the adjoint Chevalley differential.

\begin{theorem}[Relative Rees comparison]\label{thm:vol2-relative-rees}
Antisymmetrization defines a quasi-isomorphism of $R_N$-complexes
\[
 \operatorname{Alt}_N:
 C^\bullet_{\HH}(A_N/R_N,A_N)
 \longrightarrow
 \bigl(C^\bullet(\flie,S_0)\tensor R_N,\hbarq d_{\CE}\bigr),
\]
given by
\[
 (\operatorname{Alt}_NF)(x_1,\ldots,x_n)
 =\operatorname{sym}_N^{-1}
 \sum_{\sigma\in S_n}\operatorname{sgn}(\sigma)
 F(x_{\sigma(1)},\ldots,x_{\sigma(n)}).
\]
These maps commute with coefficient reduction in $N$.
\end{theorem}
\begin{proof}
Write $A=A_N$, $R=R_N$, and $A^e=A\tensor_R A^{\op}$.
In homological degree $n$, set
\[
 K_n=A\tensor_R\Lambda_R^n\flie_N\tensor_R A.
\]
Its augmentation is multiplication.
For $X=x_1\wedge\cdots\wedge x_n$, let $X_{\widehat i}$ and
$X_{\widehat i\widehat j}$ omit the indicated entries without reordering.
Define
\[
\begin{split}
 \partial(a\tensor X\tensor b)
 ={}&\sum_i(-1)^{i+1}
 \bigl(ax_i\tensor X_{\widehat i}\tensor b
       -a\tensor X_{\widehat i}\tensor x_i b\bigr)\\
 &+\hbarq\sum_{i<j}(-1)^{i+j}
 a\tensor[x_i,x_j]\wedge X_{\widehat i\widehat j}\tensor b.
\end{split}
\]
The normalized relative bimodule bar has terms
$\mathcal B_n=A\tensor_R\bar A^{\tensor_R n}\tensor_R A$.
Define
\[
 \iota_n(a\tensor X\tensor b)
 =\sum_{\sigma\in S_n}\operatorname{sgn}(\sigma)
       a[x_{\sigma(1)}|\cdots|x_{\sigma(n)}]b.
\]
The outer bar faces give the first sum in $\partial$.
Pair adjacent transpositions in the internal faces and use
$xy-yx=\hbarq[x,y]$ to obtain the second sum.
For two inputs this reads
\[
\begin{split}
 b\bigl(1[x|y]1-1[y|x]1\bigr)
 ={}&x[y]1-y[x]1+1[x]y-1[y]x\\
 &-\hbarq\,1[[x,y]]1.
\end{split}
\]
Thus $b\iota=\iota\partial$.
PBW makes $\flie_N$ a free summand of $\bar A$.
Since every factorial is invertible in $R$, $\iota$ is injective.
The identity $b^2=0$ then gives $\partial^2=0$.

Give $K_n$ filtration degree equal to its outer PBW degrees plus $n$.
The associated graded is the diagonal Koszul complex of
$\Sym_R(\flie_N)$.
Its differential uses differences of left and right polynomial generators.
A linear change of variables makes these independent polynomial variables.
They form a regular sequence even when $R$ has nilpotents.
For an augmented cycle, lift a primitive of its highest-degree symbol
and subtract its differential.
The filtration degree decreases, so this procedure terminates.
Hence $K_\bullet\to A$ is a free bimodule resolution.

The relative bar is also free over $A^e$, with its usual unit-insertion contraction.
The map $\iota$ lifts the identity of $A$.
To construct a reverse chain map, first lift the augmentation in
degree zero. At each higher degree its required boundary is a cycle;
exactness of $K_\bullet$ and projectivity of the bar term give a lift.
Each composite and the identity induce the same augmentation.
Their difference is therefore a boundary in degree zero, and the same
lifting argument, degree by degree, constructs a chain homotopy.
Applying $\operatorname{Hom}_{A^e}(-,A)$ gives a cochain homotopy equivalence.

The induced differential on
$\operatorname{Hom}_R(\Lambda_R^n\flie_N,A)$ is
\[
\begin{split}
 (\delta_N\phi)(x_1,\ldots,x_{n+1})
 ={}&\sum_i(-1)^{i+1}
 [x_i,\phi(x_1,\ldots,\widehat{x_i},\ldots,x_{n+1})]_A\\
 &+\hbarq\sum_{i<j}(-1)^{i+j}
 \phi([x_i,x_j],x_1,\ldots,\widehat{x_i},
                         \ldots,\widehat{x_j},\ldots,x_{n+1}).
\end{split}
\]
By \eqref{eq:vol2-scaled-pbw-action}, the commutator term also carries
one factor of $\hbarq$. Both sums therefore become $\hbarq$ times the
ordinary Chevalley differential.
Precomposition with $\iota$ gives $\operatorname{Alt}_N$.
All formulas commute with coefficient reduction.
\end{proof}

\section{The formal coefficient limit}
Put $R=\C\llbracket\hbarq\rrbracket$ and $A=\varprojlim_N A_N$.
PBW identifies its underlying module with $S_0\llbracket\hbarq\rrbracket$.
Each coefficient of $\hbarq$ is a polynomial.
We also write $U_\hbar(\flie)=A$ and
obtain the continuous module isomorphism
\[
 \operatorname{sym}_\hbar=\varprojlim_N\operatorname{sym}_N:
 S_0\llbracket\hbarq\rrbracket\xrightarrow{\sim}U_\hbar(\flie).
\]
Define
\[
 A^{\widehat\tensor_R n}=\varprojlim_N A_N^{\tensor_{R_N}n},
 \qquad
 C^\bullet_{\HH,\mathrm{cont}}(A/R,A)
 =\varprojlim_N C^\bullet_{\HH}(A_N/R_N,A_N).
\]
Equivalently, these are continuous $R$-multilinear normalized cochains
on the completed relative tensors.

\begin{corollary}[Formal relative Rees comparison]\label{cor:vol2-formal-relative-rees}
The compatible antisymmetrizations induce a quasi-isomorphism
\[
 C^\bullet_{\HH,\mathrm{cont}}(A/R,A)
 \longrightarrow
 \bigl(C^\bullet(\flie,S_0)\llbracket\hbarq\rrbracket,
                    \hbarq d_{\CE}\bigr).
\]
\end{corollary}
\begin{proof}
Coefficient reduction is surjective in each cochain degree.
For Hochschild cochains, lift outputs on the PBW input basis.
The Chevalley exterior input has finite rank, so its outputs also lift.
The comparison cones are acyclic and their transition maps are degreewise surjective.
The long exact sequence shows that every transition kernel is acyclic.
Choose a primitive of a compatible cycle at level zero.
Lift it one level and correct its differential by a primitive in the transition kernel.
Iteration gives compatible primitives at every order.
The inverse-limit cone is therefore acyclic.
\end{proof}

The coefficient ring is part of this result.
Absolute $\C$-linear cochains also detect variations of the coefficient parameter.
For $\flie=0$, the continuous absolute derivation $\partial_{\hbarq}$
on $R$ has no relative counterpart.
For nonabelian $\flie$, the associated graded for the $\hbarq$ filtration
has zero differential, whereas $d_{\CE}$ need not vanish.
The theorem asserts a quasi-isomorphism of complexes. A comparison of
cup products, braces and their higher coherences requires additional
maps and identities. The inverse
limit above changes only the coefficient order $N$, at a fixed
finite-dimensional Lie algebra. Passing to all Hamiltonian jets also
requires compatible transition maps for the chosen cochain complexes
and a proof that their limit computes the intended continuous carrier.

\chapter{Continuous Hamiltonian cochains}
\label{ch:vol2-continuous-cochains}

A formal Hamiltonian is determined by its finite jets. An operation on
the quantized algebra need not return as many coefficients as it inspects.
For example, extraction of a cubic coefficient can have a scalar
output. A limit of self-valued cochains at equal input and output
jet orders would miss this operation.

We therefore keep the two orders separate. The resulting continuous
Hochschild complex has braces and a relative bar resolution. Its
Chevalley comparison is compatible with both jet orders and with
the quantization parameter. This constructs a continuous algebraic
centre. Relating it to geometric collision operations or to bulk
observables requires maps preserving those additional operations.

\section{Residue duality and opposite directions of variance}

Let $k$ be a field of characteristic zero. The polynomial formulas of
Chapter~\ref{ch:vol2-hamiltonian-jets} apply over $k$. Put
\[
 \mathfrak m=(z_1,z_2),\qquad
 \h=\mathfrak m^2\subset k\llbracket z_1,z_2\rrbracket,
 \qquad \h_{\leq d}=\h/\mathfrak m^{d+1}\quad(d\geq2).
\]
The bracket is \eqref{eq:vol2-hamiltonian-bracket}. Give $\h$ its
order topology and $k$ the discrete topology. A continuous linear
functional on $\h$ vanishes on $\mathfrak m^{d+1}$ for some $d$.
Consequently
\begin{equation}\label{eq:vol2-continuous-dual}
 \h^\vee_{\mathrm{cont}}
   =\varinjlim_d\h_{\leq d}^\vee
   =\bigoplus_{a+b\geq2}k\rho_{ab},
 \qquad
 \rho_{ab}=z_1^{-a-1}z_2^{-b-1}\,dz_1\,dz_2.
\end{equation}
The notation on the right represents the iterated residue functional,
normalized by
\[
 \operatorname{Res}_0(z_1^iz_2^j\rho_{ab})=\delta_{ia}\delta_{jb}.
\]
On the whole power-series ring, the same formula includes all
$a,b\geq0$. The two degree-one modes pair with the translations in
$\h^{\mathrm{aff}}$, whereas the pointed algebra starts in degree two.
The continuous dual has the discrete topology.

\begin{proposition}[Coadjoint currents]\label{prop:vol2-continuous-coadjoint}
For $f=z_1^iz_2^j\in\h$,
\begin{equation}\label{eq:vol2-residue-coadjoint}
 \operatorname{ad}_f^*\rho_{ab}
  =\bigl(j(a+1)-i(b+1)\bigr)\rho_{a-i+1,b-j+1}.
\end{equation}
A mode with a negative index or with total degree less than two is
zero as a functional on $\h$. The formula extends to formal $f$,
and the coadjoint action is continuous.
\end{proposition}

\begin{proof}
By definition,
$(\operatorname{ad}_f^*\rho_{ab})(g)=-\rho_{ab}(\{f,g\})$.
For $g=z_1^pz_2^q$, the monomial bracket has coefficient $iq-jp$
and exponent $(i+p-1,j+q-1)$. The residue can be nonzero only at
$(p,q)=(a-i+1,b-j+1)$, where its negative coefficient is
$j(a+1)-i(b+1)$. This proves the formula.
If $a+b\leq d$, a Hamiltonian of order greater than $d$ acts by
zero: the possible output mode would have degree at most one.
Thus the action on any finite sum of residue modes uses only
finitely many Taylor coefficients of $f$. Its output is again a
finite sum. This also proves continuity for the stated topologies.
\end{proof}

The graded semidirect Lie algebra
\[
 \mathfrak q_{\mathrm{cont}}
   =\h\ltimes\h^\vee_{\mathrm{cont}}[-1]
\]
has degree-zero Hamiltonians and degree-one covectors. The second
summand is abelian. The order topology on the first summand and the
discrete topology on the second make it separated and complete.
Evaluation gives the invariant pairing between the two summands:
its invariance is the defining negative-transpose identity above.

The quotient $\h\to\h_{\leq d}$ induces an inclusion of duals.
Its image is a coadjoint submodule, and there is a surjective Lie map
\begin{equation}\label{eq:vol2-current-subquotient}
 \h\ltimes\h_{\leq d}^\vee[-1]
   \longrightarrow
 \mathfrak q(\h_{\leq d})
   =\h_{\leq d}\ltimes\h_{\leq d}^\vee[-1],
 \qquad (x,\alpha)\longmapsto(x\bmod\mathfrak m^{d+1},\alpha).
\end{equation}
Its kernel is $\mathfrak m^{d+1}$ in the first summand. Indeed,
that ideal acts trivially on the displayed finite dual, and the
coadjoint formula descends to the finite quotient. Thus the finite
current algebra is a quotient of a subalgebra of
$\mathfrak q_{\mathrm{cont}}$.

Simultaneous coordinate truncation on both full summands is not a
Lie map. At $d=2$, it kills both $z_1^3$ and $\rho_{21}$, but
\[
 \operatorname{ad}_{z_1^3}^*\rho_{21}=-6\rho_{02}
\]
has a nonzero quadratic component. The quotient of Hamiltonians and
the inclusion of covectors must retain their opposite directions.

\section{The completed algebra and its continuous operations}

Set $R=k\llbracket\hbarq\rrbracket$ and
$R_N=R/(\hbarq^{N+1})$. For $d\geq2$ and $N\geq0$, define
\begin{equation}\label{eq:vol2-continuous-algebra}
 A_{d,N}=U_{R_N}(\h_{\leq d}\otimes_kR_N,\hbarq\{-,-\}),
 \qquad A=\widehat U_\hbar(\h)=\varprojlim_{d,N}A_{d,N}.
\end{equation}
The defining relation is $xy-yx=\hbarq\{x,y\}$. Choose one ordering
of the Hamiltonian monomials, compatible with the inclusions of
their finite bases. The PBW theorem proved in
Proposition~\ref{prop:vol2-rees-pbw} gives free $R_N$-modules with
the corresponding ordered-monomial bases. Every quotient map in
\eqref{eq:vol2-continuous-algebra} is onto.

Give $A$ the inverse-limit topology with discrete quotients $A_{d,N}$.
At each fixed $(d,N)$, PBW expressions are finite polynomials.
In particular, this construction does not complete the polynomial
degree in the enveloping generators. For $n\geq1$, put
\begin{equation}\label{eq:vol2-quotient-tensor}
 A^{\widehat\otimes_R n}
  =\varprojlim_{d,N}A_{d,N}^{\otimes_{R_N}n},
 \qquad A^{\widehat\otimes_R0}=R.
\end{equation}
Different cutoffs in the tensor factors give the same limit because
a common larger cutoff is cofinal. All continuity statements below
refer to this quotient tensor.

Here is an operation requiring different jet orders. Let $X_{30}$
be the enveloping generator represented by the Hamiltonian $z_1^3$.
Projection to $A_{3,N}$ followed by extraction of the PBW coefficient
of the single monomial $X_{30}$ is an $R_N$-linear scalar-valued map.
The compatible maps in $N$, followed by the scalar inclusion into
$A$, define a continuous $R$-linear map $\lambda_{30}:A\to A$.
Its output in $A_{2,N}$ is nonzero on $X_{30}$, although the
input projection to $A_{2,N}$ kills $X_{30}$. That output therefore
cannot be computed from the quadratic input jet.

\begin{definition}\label{def:vol2-uniform-cochains}
The continuous Hochschild cochains for the quotient tensor are
\begin{equation}\label{eq:vol2-two-index-cochains}
 C^n_{\mathrm{uc}}(A/R,A)
   =\varprojlim_{d,N}\varinjlim_{e\geq d}
       \operatorname{Hom}_{R_N}
       (A_{e,N}^{\otimes_{R_N}n},A_{d,N}).
\end{equation}
The direct maps precompose with input quotients. The inverse maps
reduce the output and the parameter precision. Give each direct
limit the discrete topology and each $C^n_{\mathrm{uc}}$ its
inverse-limit topology. The subscript specifies uniform continuity
for \eqref{eq:vol2-quotient-tensor}.
\end{definition}

\begin{lemma}[An input cutoff for each output]\label{lem:vol2-continuous-maps}
The quotient projections in \eqref{eq:vol2-quotient-tensor} are onto,
and there is a natural identification
\begin{equation}\label{eq:vol2-continuous-hom}
 C^n_{\mathrm{uc}}(A/R,A)
  =\operatorname{Hom}_{R,\mathrm{cont}}
       (A^{\widehat\otimes_R n},A).
\end{equation}
Both the tensor source and the cochain space are separated and complete.
\end{lemma}

\begin{proof}
A finite PBW tensor lifts by the same ordered monomials to all
larger cutoffs; hence every projection is onto. On underlying modules,
symmetrization identifies the tensor source with
\[
 \left(\varprojlim_d
      (\operatorname{Sym}_k\h_{\leq d})^{\otimes_k n}\right)
       \llbracket\hbarq\rrbracket.
\]
Each parameter coefficient is a compatible system of finite
polynomials. Coefficientwise division by $\hbarq^{N+1}$ therefore
identifies its quotient by that parameter power with
$\varprojlim_d A_{d,N}^{\otimes n}$. The inverse-limit description
also gives separatedness and completeness.

Let $f$ be a continuous map on the right of
\eqref{eq:vol2-continuous-hom}. Its composite with the output
$(d,N)$ kills the kernel of some input projection $(e,M)$.
Increase that projection so that $e\geq d$ and $M\geq N$.
Surjectivity makes the induced map $R$-linear. Since the target is
killed by $\hbarq^{N+1}$, it factors through
$A_{e,N}^{\otimes n}$. Two choices give the same map precisely
after a common input refinement. The output compatibilities give
\eqref{eq:vol2-two-index-cochains}.
Conversely, its compatible families define a map into $A$, and their
finite input factorizations at each output make it continuous.
Finally, an inverse limit of discrete modules is a closed submodule
of their product and is complete, proving the assertion for cochains.
\end{proof}

\section{Composition and the relative derived centre}

All enveloping generators have cohomological degree zero, and an
$n$-cochain has cohomological degree $n$. Its Hochschild differential is
\begin{align*}
 (bf)(a_1,\ldots,a_{n+1})={}&a_1f(a_2,\ldots,a_{n+1})\\
 &+\sum_{i=1}^n(-1)^i
       f(a_1,\ldots,a_i a_{i+1},\ldots,a_{n+1})\\
 &+(-1)^{n+1}f(a_1,\ldots,a_n)a_{n+1}.
\end{align*}
For brace signs, use the reduced degree $|f|=n-1$.
If $g_j$ has arity $q_j$, define
\begin{equation}\label{eq:vol2-continuous-braces}
 f\{g_1,\ldots,g_r\}
   =\sum_{\text{ordered disjoint insertions}}
       (-1)^{\sum_j(q_j-1)i_j}
       f(\ldots,g_1(\ldots),\ldots,g_r(\ldots),\ldots),
\end{equation}
where $i_j$ is the number of arguments preceding the block of
$g_j$ in the resulting argument list. Equivalently, these are the
Koszul insertion signs on the suspended tensor coalgebra.
For $f$ of arity $p$ and $g$ of arity $q$, the cup product is
\[
 (f\smile g)(a_1,\ldots,a_{p+q})
   =f(a_1,\ldots,a_p)g(a_{p+1},\ldots,a_{p+q}).
\]

\begin{theorem}[Continuous braces]\label{thm:vol2-continuous-braces}
The differential, cup product and every finite brace preserve
$C^\bullet_{\mathrm{uc}}(A/R,A)$. They are continuous for the
stated cochain topologies and obey the Hochschild brace identities.
\end{theorem}

\begin{proof}
Fix an output quotient $(d,N)$ for the outer cochain $f$. It uses
some input cutoff $e$. Each inner cochain $g_j$ is then needed only
at output $(e,N)$, and that output uses some finite input cutoff
$e_j$. Choose $E$ at least as large as $e,e_1,\ldots,e_r$.
Every summand is now a map from a tensor power of $A_{E,N}$ to
$A_{d,N}$. There are finitely many summands at fixed arities, and
further refinement leaves their class unchanged. This constructs
the brace. Multiplication is defined at each finite quotient, so
the same argument constructs $b$ and the cup product.

To prove continuity at a specified tuple of cochains, fix its outer
component at $(d,N)$ and the finitely many inner components just
used. These equalities describe open neighborhoods, and all outputs
at $(d,N)$ agree throughout them. The argument applies to every
output quotient.
For a brace identity, choose one common quotient for all the
intermediate outputs in its two sides. Both sides enumerate the
same nested or disjoint insertions there. The Koszul sign is the
sign of the same permutation of suspended inputs, so the terms
agree. Thus the identity holds at every output and in the limit.
Associativity of multiplication gives $b^2=0$.
\end{proof}

An individual rectangular complex with coefficients in $A_{d,N}$
and inputs in $A_{e,N}$ is not a self-valued brace algebra when
$e\ne d$. The total cochain space supplies the intermediate output
needed by the next insertion.

We now specify the derived category for this construction. Work with
complete separated linearly topologized $R$-modules and quotient
completed tensors. A map out of a tensor is continuous when each
discrete output factors through discrete quotients of all its inputs.
Use complete $A$-bimodules whose action is continuous for this tensor,
and write $A^e=A\widehat\otimes_R A^{\mathrm{op}}$.
A short exact sequence is \emph{relatively admissible} when it has
a continuous $R$-linear splitting, which need not be $A$-bilinear.
Relative resolutions and derived Hom refer to this exact structure.

For a complete coefficient module $E$ and a complete $A$-bimodule $M$,
restriction to
$1\otimes E\otimes1$ gives the induction adjunction
\begin{equation}\label{eq:vol2-continuous-induction}
 \operatorname{Hom}_{A^e,\mathrm{cont}}
       (A\widehat\otimes_R E\widehat\otimes_R A,M)
   =\operatorname{Hom}_{R,\mathrm{cont}}(E,M).
\end{equation}
The inverse sends $f$ to $a\otimes e\otimes b\mapsto af(e)b$.
First refining the output for the action and then the input for $f$
proves its continuity. If $M'\to M$ has a continuous $R$-linear
section, a map on the right lifts by composition with that section.
The adjunction then lifts the bimodule map. Thus induced modules
are relatively projective.

\begin{theorem}[The completed relative bar]\label{thm:vol2-continuous-bar}
The augmented complex
\[
 B_n(A)=A\widehat\otimes_R A^{\widehat\otimes_R n}
                   \widehat\otimes_R A\quad(n\geq0),
 \qquad B_0(A)\longrightarrow A,
\]
with alternating multiplication boundary and $B_n$ in degree $-n$,
is a relative projective resolution. Hence
\begin{equation}\label{eq:vol2-continuous-centre}
 \operatorname{RHom}^{\mathrm{rel,cont}}_{A^e}(A,A)
    \simeq C^\bullet_{\mathrm{uc}}(A/R,A).
\end{equation}
\end{theorem}

\begin{proof}
Each term is induced, hence relatively projective by
\eqref{eq:vol2-continuous-induction}. On the augmented bar, put
\[
 s_n(a_0\otimes\cdots\otimes a_{n+1})
   =1\otimes a_0\otimes\cdots\otimes a_{n+1},
 \qquad s_{-1}(a)=1\otimes a.
\]
These maps are continuous and $R$-linear. In $\partial s+s\partial$,
every multiplication except the newly inserted first one occurs
twice with opposite signs. The remaining term is the identity.
On cycles, $s$ is therefore a continuous $R$-linear section of the
boundary onto those cycles. This proves admissible exactness.
Continuous bimodule Hom identifies the term in degree $n$ with
\eqref{eq:vol2-continuous-hom}. Its induced differential is exactly
$b$, proving \eqref{eq:vol2-continuous-centre}.
\end{proof}

This is the derived centre in the stated relative continuous theory.
An absolute derived category of topological modules or an analytic
tensor product requires a comparison with this resolution.

\section{The Chevalley comparison and the inverse limit}

For $e\geq d$, let $A_{d,N}$ be an $A_{e,N}$-bimodule through the
quotient. Define
\begin{equation}\label{eq:vol2-rectangular-chevalley}
 P^n_{e,d,N}
  =\operatorname{Hom}_k\bigl(\Lambda^n\h_{\leq e},
                  \operatorname{Sym}_k\h_{\leq d}\bigr)\otimes_kR_N.
\end{equation}
The coefficient action of $\h_{\leq e}$ factors through
$\h_{\leq d}$ and extends its adjoint action as a derivation of
the symmetric algebra. Let $d_{\mathrm{CE}}$ be the Chevalley
differential with this action and with the Poisson Lie bracket.
For example,
\[
 (d_{\mathrm{CE}}f)(x,y)=x\cdot f(y)-y\cdot f(x)-f(\{x,y\})
\]
for a one-cochain. In general the action term indexed by $i$ has
sign $(-1)^i$ on inputs $x_0,\ldots,x_n$, and the bracket term
indexed by $i<j$ has sign $(-1)^{i+j}$, with $\{x_i,x_j\}$ put
first and the remaining inputs in their original order.
The differential on $P^\bullet_{e,d,N}$ is
$\hbarq d_{\mathrm{CE}}$.

\begin{theorem}[Chevalley comparison with two jet orders]
\label{thm:vol2-continuous-chevalley}
Output PBW symmetrization and input antisymmetrization define
natural quasi-isomorphisms
\begin{align}
 C^\bullet_{\HH}(A_{e,N}/R_N,A_{d,N})
   &\longrightarrow (P^\bullet_{e,d,N},\hbarq d_{\mathrm{CE}}),
       \label{eq:vol2-rectangular-map}\\
 f&\longmapsto\left[(x_1,\ldots,x_n)\longmapsto
     \operatorname{PBW}^{-1}\sum_{\sigma\in S_n}
       \operatorname{sgn}(\sigma)
       f(x_{\sigma(1)},\ldots,x_{\sigma(n)})\right].\notag
\end{align}
They commute with input refinement, output quotient and parameter
reduction. Consequently they induce a quasi-isomorphism
\begin{equation}\label{eq:vol2-continuous-chevalley-map}
 C^\bullet_{\mathrm{uc}}(A/R,A)
   \longrightarrow
 \mathcal P_{\mathrm{uc}}
   :=\varprojlim_{d,N}\varinjlim_{e\geq d}
                   (P^\bullet_{e,d,N},\hbarq d_{\mathrm{CE}}).
\end{equation}
For both systems in this formula, the ordinary inverse limit
computes the derived inverse limit.
\end{theorem}

\begin{proof}
Use the finite Chevalley bimodule resolution from the proof of
Theorem~\ref{thm:vol2-relative-rees}, now for $A_{e,N}$ and with
coefficient bimodule $A_{d,N}$. Its terms are
\[
 A_{e,N}\otimes_{R_N}\Lambda^n(\h_{\leq e}\otimes_kR_N)
          \otimes_{R_N}A_{e,N}.
\]
They are free bimodules. Filtering by PBW degree plus exterior
degree gives the Koszul resolution of the diagonal of a polynomial
algebra over $R_N$. The diagonal differences form a regular sequence
even over this nilpotent coefficient ring. A cycle has finite
filtration degree, and lifting an associated-graded primitive lowers
that degree. This proves exactness without a convergence assumption.

Antisymmetrization of the exterior factors is a map into the bar
resolution: internal bar products combine in pairs to give the
scaled Lie commutator, and the two endpoint terms give the left and
right actions. It is therefore a chain homotopy equivalence between
projective resolutions. Applying bimodule Hom into $A_{d,N}$ gives
\eqref{eq:vol2-rectangular-map}. Under PBW, the commutator action of
$x$ on the coefficient algebra is $\hbarq$ times its adjoint Poisson
action; the bracket term also carries $\hbarq$. This proves the
stated differential. All formulas commute with the three quotient
maps, proving naturality.

Filtered direct limits of modules are exact. For the output inverse
limit, use the countable cofinal sequence $(d,N)=(r+2,r)$.
Every cochain at one output factors through a finite input order
$e$. For the next output, increase $e$ if necessary and precompose
with the input quotient. Lift the value of each PBW tensor basis
element to the next output and parameter precision, then extend
linearly over the new coefficient ring. Each value is a finite
polynomial, so this defines an allowed cochain. The exterior basis
gives the same construction on $P^n_{e,d,N}$. These are degreewise
lifts, not chain sections; no $R_N$-linear section of
$R_{N+1}\to R_N$ is used.

Thus both cochain towers are degreewise surjective. For a countable
tower of complexes with transitions $p_r:K_{r+1}\to K_r$, the
derived inverse limit is the homotopy fibre of
\[
 1-\mathrm{shift}:\prod_rK_r\longrightarrow\prod_rK_r.
\]
Degreewise surjectivity of the transition maps makes this product
map degreewise onto: choose the next coordinate recursively to
solve $x_r-p_rx_{r+1}=y_r$. Its kernel is $\varprojlim_rK_r$.
The homotopy fibre is consequently quasi-isomorphic to that kernel.
Applying the derived inverse limit to the levelwise comparisons
therefore gives \eqref{eq:vol2-continuous-chevalley-map}.
\end{proof}

At $\hbarq=0$, the Chevalley differential vanishes, as required by
the commutative polynomial special fibre. At positive parameter
precision the two Chevalley terms remain scaled together.
The theorem identifies complexes. It does not identify a chosen
polyvector wedge with the Hochschild cup product or make
antisymmetrization a strict brace map.

Finite comparison homotopies can also be computed. At parameter
zero, write $t_i=x_i-y_i$ for the diagonal differences in the
Koszul resolution and $\epsilon_i$ for its exterior generators.
On positive total $t,\epsilon$ degree, the operators
\[
 d_0=\sum_i t_i\partial_{\epsilon_i},\qquad
 h_0=\frac{\sum_i\epsilon_i\partial_{t_i}}
                    {\deg_t+\deg_\epsilon}
\]
satisfy $d_0h_0+h_0d_0=1$; on $t,\epsilon$-constant terms the
homotopy is zero. Let $i_0$ include this constant summand and let
$p_0$ project to it. The anticommutator before division is the
total Euler operator, so
\[
 d_0h_0+h_0d_0=1-i_0p_0,\qquad
 h_0^2=h_0i_0=p_0h_0=0.
\]
After PBW identification the quantum differential is
$d_K=d_0+\delta$, where $\delta$ has order $\hbarq$.
Modulo $\hbarq^{N+1}$, put
\[
 T=\sum_{j=0}^{N}(-\delta h_0)^j,\qquad
 h_K=h_0T,\qquad p_K=p_0T,\qquad
 i_K=\sum_{j=0}^{N}(-h_0\delta)^j i_0.
\]
These sums are finite. Multiplying them and using
$d_K^2=0$ and the preceding identities gives
\[
 d_Kh_K+h_Kd_K=1-i_Kp_K,\qquad p_Ki_K=1.
\]
The differential transferred to the constant summand is zero by
degree. That summand represents the diagonal quotient $A_{e,N}$,
so $p_K$ is its augmentation in PBW coordinates. This supplies
the contraction on the augmented resolution.

Starting with the augmentation lift, define the bar-to-Chevalley
comparison on a free middle generator by
\[
 \psi_n(1[a_1|\cdots|a_n]1)
   =h_K\psi_{n-1}\partial(1[a_1|\cdots|a_n]1),
\]
and extend bimodule-linearly. Its argument is a cycle by the previous
stage, so the contraction identity proves the chain-map equation.
If $\alpha$ is the antisymmetrization map into the bar, define on
its free middle generators
\[
 H_n=s_n\bigl((1-\alpha\psi)_n-H_{n-1}\partial\bigr),
 \qquad H_{-1}=0,
\]
and again extend bimodule-linearly. The expression in parentheses
is a cycle by induction, so $\partial H+H\partial=1-\alpha\psi$.
These are finite algorithms at every fixed cutoff.
Transporting higher operations through a simultaneous continuous
choice requires the compatibility of those homotopies as well.

\section{A localization test for limits and adjoints}

The surjectivity used above is a substantive hypothesis. A scalar
example shows both what can fail without it and why changing the
completion changes the problem.

Let $S=k[t]$, $L=k[t,t^{-1}]$, and write $D(S)$ and $D(L)$ for
the derived categories of modules. Restriction of scalars
$i:D(L)\to D(S)$ is fully faithful, since localization is flat and
$L\otimes_S L=L$. Its left adjoint is
$L\otimes_S^{\mathbb L}(-)$; its right adjoint is
$\operatorname{RHom}_S(L,-)$. The right adjoint represents maps
from localized objects into a given object, whereas the left
adjoint represents maps out of that object into localized ones.

\begin{proposition}[The two localization adjoints]
\label{prop:vol2-localization-adjoints}
Their values at $S$ are
\[
 L\otimes_S^{\mathbb L}S=L,
 \qquad
 \operatorname{RHom}_S(L,S)
     \simeq(k\llbracket t\rrbracket/k[t])[-1].
\]
The latter equivalence is one of $S$-complexes. Multiplication by
$t$ is invertible on its displayed cohomology and gives the
$L$-module structure. For $\widehat S=k\llbracket t\rrbracket$
and $\widehat L=k((t))$, the right adjoint takes $\widehat S$ to
zero, while the left adjoint takes it to $\widehat L$.
\end{proposition}

\begin{proof}
The direct system $S\xrightarrow{t}S\xrightarrow{t}\cdots$ has
colimit $L$, where the $n$th copy maps by multiplication by $t^{-n}$.
Its free telescope resolution has terms $\bigoplus_{n\geq0}S$
and differential $e_n\mapsto e_n-te_{n+1}$. Injectivity follows
by inspecting the least nonzero coordinate of a finitely supported
vector. Its cokernel imposes exactly the relations of that colimit.
Applying $\operatorname{Hom}_S(-,S)$ gives the complex in degrees
zero and one
\begin{equation}\label{eq:vol2-polynomial-telescope}
 \prod_{n\geq0}S\xrightarrow{D}\prod_{n\geq0}S,
 \qquad D(a)_n=a_n-ta_{n+1}.
\end{equation}
If $D(a)=0$, each $a_n$ is divisible by every power of $t$, so
$a_n=0$. Define
\[
 \Theta(b)=\left[\sum_{n\geq0}t^nb_n\right]
               \in k\llbracket t\rrbracket/k[t].
\]
The series converges coefficientwise. For $b=D(a)$ it telescopes
to the polynomial $a_0$, so $\Theta$ annihilates the image of $D$.
It is onto, by taking the $b_n$ to be scalar coefficients of a series.
If its value is zero, the series equals a polynomial $c$. Put
\[
 a_n=t^{-n}\left(c-\sum_{j<n}t^jb_j\right).
\]
The expression in parentheses is a polynomial divisible by $t^n$,
so $a_n\in S$, and $D(a)=b$. This proves the asserted cokernel.
On that cokernel, the inverse of multiplication by $t$ sends
$[f]$ to $[(f-f(0))/t]$.

Over $\widehat S$, the same telescope differential has the inverse
\[
 a_n=\sum_{j\geq0}t^j b_{n+j}.
\]
Each series converges $t$-adically, and the homogeneous equation
still has only the zero solution. Thus the derived right-adjoint
value is zero. Flat localization computes both left-adjoint values.
\end{proof}

For $b_n=1$, the class $\Theta(b)=[(1-t)^{-1}]$ is nonzero.
Every finite subsystem nevertheless has a polynomial solution:
with $a_{M+1}=0$, take
\[
 a_n=1+t+\cdots+t^{M-n}\qquad(0\leq n\leq M).
\]
These degrees grow with $M$ and give no polynomial solution of the
infinite system. Completion supplies such a solution by admitting
a different class of coefficients. The continuous Hamiltonian
comparison avoids this inference by proving surjectivity of its
actual cochain transitions.

\section{Compatible models and geometric actions}

The algebra $A$ of \eqref{eq:vol2-continuous-algebra} is the formal
Hamiltonian boundary algebra, and \eqref{eq:vol2-continuous-centre}
constructs its relative continuous derived centre. The explicit
model $\mathcal P_{\mathrm{uc}}$ records the corresponding
Hamiltonian cochains with their two variances. A different
polyvector realization must be compared with this complex in its
own topology.

To compare another realization obtained from finite models, one
must also specify its transitions. At a fixed jet order let
\[
 A_d=\varprojlim_NA_{d,N},\qquad
 \mathcal P_d=
 \bigl(C^\bullet(\h_{\leq d},\operatorname{Sym}\h_{\leq d})
              \llbracket\hbarq\rrbracket,\hbarq d_{\mathrm{CE}}\bigr).
\]
A \emph{compatible jet comparison} consists of complete $R$-linear
cochain models $P_d,H_d$, continuous transition maps, levelwise comparisons
$P_d\simeq\mathcal P_d$ and
$H_d\simeq C^\bullet_{\HH,\mathrm{cont}}(A_d/R,A_d)$, and
compatible quasi-isomorphisms $P_d\to H_d$. All comparison maps
and limit identifications are continuous. Require degreewise
surjective transitions or use derived inverse limits. Specify the
identifications of their chosen limits with a proposed complex
$\mathcal P_\infty$ and with $C^\bullet_{\mathrm{uc}}(A/R,A)$.
An ordinary limit may be used only when its comparison with the
derived limit is supplied. All identifications use the quotient
topology fixed above.

\begin{theorem}[Continuous comparison from compatible jets]
\label{thm:vol2-compatible-jets}
A compatible jet comparison induces a continuous quasi-isomorphism
\[
 \Phi_\infty:\mathcal P_\infty\xrightarrow{\sim}
                C^\bullet_{\mathrm{uc}}(A/R,A).
\]
Compatible homotopy-$\Etwo$ and Swiss-cheese refinements, whose
operations converge in the chosen topology, give the corresponding
complete open--closed operations.
\end{theorem}

\begin{proof}
The derived inverse limit preserves levelwise equivalences. For
degreewise surjective towers it agrees with the ordinary limit by
the product calculation in the proof of
Theorem~\ref{thm:vol2-continuous-chevalley}. Apply the specified
limit identifications. For the final assertion, the supplied
operadic comparisons identify each operation and its coherence
homotopies; convergence defines them on the limits.
\end{proof}

The two-index construction already gives the algebraic cochain
comparison \eqref{eq:vol2-continuous-chevalley-map}. It supplies
neither a chosen geometric polyvector product nor its equivalence
with the brace operations. In particular, an invariant continuous
Duflo map requires a multiplicative comparison and a specified
prescription for its wheel coefficients.

Finite matrix branes are representation-theoretic images of the
formal Hamiltonian boundary algebra whenever their representations
are continuous for this completion. Its relative derived centre is
then an algebraic target for coherent bulk actions. A field-theory
identification starts instead from that theory's observables and
constructs their action map. It must preserve the insertion supports,
collisions, vacuum and higher substitutions. The continuous algebraic
resolution provides the target of this comparison; it does not
construct its geometric or physical source.

\chapter{Current actions and derived interface centres}
\label{ch:vol2-currents}

A boundary action has to commute with fields inserted near the boundary.
When this commutation holds up to homotopy, two fields require a
compatibility between the two homotopies and their collision. The
compatibility is additional data. We first recover its cochains from
supported maps and construct them for a resolved current action.

The Heisenberg algebra makes the resulting distinctions explicit. At a
polynomial root, forgetting the current action retains many more
operators than the current cohomology sees. Their derived composition
has a nonzero Massey product. Over the level ring, the path in the
forgetful fibre turns a torsion current class into a free class. Finite
categories of interfaces then give centres in which cup products,
insertion brackets and parameter fibres can all be computed. These
calculations specify operations that a geometric bulk--boundary
comparison must preserve.

Throughout the chapter, $k$ is a field of characteristic zero. Complexes
are cohomological, with $K[1]^n=K^{n+1}$. For a homogeneous map $f$ the
Hom differential is $d_{\mathrm{Hom}}f=d f-(-1)^{|f|}f d$.
Polynomial states have finite support in their variables. A completion
will always be stated separately.

\section{Current homotopies from a supported map}
\label{sec:vol2-current-homotopies}

Let $O=k[y]$ and $D_y=O\langle\partial_y\rangle$, with
$[\partial_y,y]=1$. The Heisenberg vacuum at level $\kappa\in k^\times$
is
\begin{equation}\label{eq:vol2-fock-currents}
 V_\kappa=k[x_1,x_2,\ldots],\qquad
 J_{-n}=L_{x_n},\quad J_n=\kappa n\partial_{x_n}\quad(n>0),
 \qquad J_0=0.
\end{equation}
Here $L_f$ denotes multiplication by $f$. Its vacuum is $1$, its
current state is $J=x_1$, and its translation is
\[
 T=\sum_{n\geq1}n x_{n+1}\partial_{x_n},\qquad
 \mathcal V=O\otimes_kV_\kappa,\qquad \nabla=\partial_y+T.
\]
For a formal variable $u$, the field
$Y(J,u)=\sum_{m\in\mathbb Z}J_m u^{-m-1}$ determines the other
fields by normal products and translation. Direct calculation
gives $[J_m,J_n]=\kappa m\delta_{m+n,0}$ and
$[\nabla,J_m]=-mJ_{m-1}$. Each field is lower truncated on each state.

Take a second coordinate $z$ and put $u=z-y$. A map supported on $u=0$
takes values in principal parts, with all finite normal derivatives
retained as in Section~\ref{sec:vol2-principal-parts}. The current
generates a free $D_y$-submodule: $T^nJ=n!x_{n+1}$. Denote it by
$L_J=D_yJ$.

A semi-free complex has an increasing filtration by subcomplexes,
split on underlying graded modules, whose successive quotients
are direct sums of shifts of $D_y$.

\begin{lemma}[Laurent coefficients of supported maps]
\label{lem:vol2-laurent-current}
Let $Q$ be a semi-free complex of left $D_y$-modules. Let $j$ include
the complement of the diagonal and let $\Delta$ include the diagonal.
Then
\[
 R\operatorname{Hom}_{D_{z,y}}
   \bigl(j_*j^*(L_J\boxtimes Q),\Delta_+\mathcal V\bigr)
\]
is computed by ordinary differential-module maps. Taking the residue
identifies this complex with families of $O$-linear maps
$H_m:Q\to\mathcal V$, $m\in\mathbb Z$, satisfying
\begin{equation}\label{eq:vol2-laurent-horizontal}
 [\nabla,H_m]+mH_{m-1}=0.
\end{equation}
For each input, only finitely many $H_m$ with $m\geq0$ are nonzero.
There is no bound required uniformly over all inputs.
\end{lemma}
\begin{proof}
For one free generator the localized source has the telescope
presentation
\[
 D_{z,y}[u^{-1}]=\varinjlim_k D_{z,y}u^{-k},
\]
whose maps are right multiplication by $u$. The Ore localization
exists because $\operatorname{ad}u$ is locally nilpotent. Applying
Hom gives the inverse system of principal parts under multiplication
by $u$. These maps are onto: $u^{-r}$ has preimage $u^{-r-1}$.
Thus its ordinary limit computes its derived limit, by the product
model in the proof of Theorem~\ref{thm:vol2-continuous-chevalley}. Direct sums of
free generators become products of these exact calculations.

Use the semi-free filtration of $Q$, split on underlying graded
modules. At each finite stage the assertion follows by extensions.
The restriction maps between the ordinary Hom complexes are
degreewise onto, so the assertion passes to the entire filtration.
The resulting totalization uses products.

The map with residue family $(H_m)$ is reconstructed by
\begin{equation}\label{eq:vol2-laurent-reconstruction}
 \varphi(u^mJ\boxtimes q)
   =\sum_{r\geq0}H_{m+r}(q)u^{-r-1}.
\end{equation}
The stated finiteness makes this a principal part for each $m,q$.
Multiplication by $u$ shifts the formula. Linearity in $\partial_z$
then determines the values on every $T^nJ$. Finally
$\partial_y|_z=\partial_y|_u-\partial_u$ and
$\operatorname{Res}_u\partial_u=0$ give
\[
 \nabla H_m(q)=-mH_{m-1}(q)+H_m(\partial_yq),
\]
which is precisely the remaining differential-module relation.
These constructions are inverse and commute with the Hom differential.
\end{proof}

To extract these maps from a collision, retain its attaching map.
Write a one-current Cousin complex as $A[1]\oplus B$, where $A$ is
the open term, $B$ is the supported term, and $\mu:A\to B$ is the
unshifted residue map. For the suspension $s:A\to A[1]$ of degree
$-1$, its differential is
\[
 d(sa,b)=(-s\,d_Aa,\mu a+d_Bb).
\]
A map to a second such complex has open and closed components
$f_U,f_\Delta$ and a degree $-1$ cross component $h$. Its equation is
$d_{\mathrm{Hom}}h=f_\Delta\mu-\mu'f_U$.

Suppose that the prescribed unary operation is
$a:Q\to\mathcal V$, and that the open and closed restrictions are
identified with $1\boxtimes a$ and $a$ by specified homotopies
$\gamma,\delta$. Orient them by
$d\gamma=f_U-1\boxtimes a$ and $d\delta=f_\Delta-a$.
Writing $\partial_P,\partial_C$ for the source and target residue
maps, set
\begin{equation}\label{eq:vol2-corrected-current-face}
 \widetilde h=h+\partial_C\gamma-\delta\partial_P,
 \qquad
 d\widetilde h=a\partial_P-\partial_C(1\boxtimes a).
\end{equation}
For a semi-free $P$, take $Q=P\otimes_O\mathcal V$ with its
prescribed boundary current action. On a free $D_y$-summand,
the change from $\partial_y\otimes1$ to the diagonal connection
$\partial_y\otimes1+1\otimes\nabla$ is triangular in PBW degree.
It identifies the tensor with the module induced from the
$O$-free module $\mathcal V$, so the tensor is $D_y$-free.
Applying this to the filtration quotients makes $Q$ semi-free.
The residue map $\partial_P$
acts on the factor $\mathcal V$. The identity
\[
 \operatorname{Res}_u\operatorname{pp}_u(u^mY(J,u)w)=J_mw
\]
therefore gives
\begin{equation}\label{eq:vol2-current-first}
 H_m=-\operatorname{Res}_u\widetilde h(u^mJ\boxtimes-),
 \qquad d_{\mathrm{Hom}}H_m=J_ma-aJ_m.
\end{equation}
The source action is the rightmost $J_m$ in this formula. In particular,
the zero and negative modes are part of the comparison. Equality of
the restrictions in cohomology would not specify $\gamma,\delta$ or
the corrected supported coordinate.

\subsection{The second equation and all higher equations}

Let $\mathfrak a$ be an ordinary Lie algebra and let $M,V$ be dg
$\mathfrak a$-modules. Put $E=\operatorname{Hom}_k(M,V)$ and
$T_xf=\rho_V(x)f-f\rho_M(x)$. For an alternating $p$-cochain set
\begin{align}\label{eq:vol2-current-ce}
 (\delta f)(x_0,\ldots,x_p)
 ={}&\sum_i(-1)^iT_{x_i}f(x_0,\ldots,\widehat x_i,\ldots,x_p)
 \notag\\
 &+\sum_{i<j}(-1)^{i+j}
 f([x_i,x_j],x_0,\ldots,\widehat x_i,\ldots,\widehat x_j,\ldots,x_p).
\end{align}
The bracket occupies the first slot. The representation identity and
Jacobi identity give $\delta^2=0$; moreover $\delta d_E=d_E\delta$.
Use the product totalization with differential
$\delta+(-1)^p d_E$ on Chevalley degree $p$.

A coherent current-linear map is a family
$H_p:\Lambda^p\mathfrak a\to E^{-p}$ satisfying
\begin{equation}\label{eq:vol2-current-hierarchy}
 H_0=a,\qquad d_Ea=0,\qquad d_EH_p=\delta H_{p-1}\quad(p\geq1).
\end{equation}
Indeed, $\alpha_p=(-1)^{p(p-1)/2}H_p$ is then a total degree-zero
cocycle. At two inputs the equation is
\begin{equation}\label{eq:vol2-current-second}
 d_EH_2(x,y)=T_xH_1(y)-T_yH_1(x)-H_1([x,y]).
\end{equation}
Its right side is closed once the first equation holds. Its class in
$\operatorname{Hom}(\Lambda^2\mathfrak a,H^{-1}E)$ must vanish
before this chosen first homotopy extends. The three terms keep the
two orders of current--root collision and the current--current
collision. The last is the third channel.

\begin{theorem}[A resolved current action]\label{thm:vol2-resolved-current}
Suppose $I$ is a dg $\mathfrak a$-module and $i:V\to I$ is equivariant.
Let $(p,i,h)$ be a contraction of complexes with
\[
 pi=1,\qquad d_Ih+hd_I=1-ip,\qquad hi=ph=h^2=0.
\]
For every equivariant cochain map $j:M\to I$ there is a coherent
current-linear map with unary part $pj$. It is given in all arities
by a series finite in each output arity. Neither $p$ nor $h$ need
be equivariant.
\end{theorem}
\begin{proof}
Postcomposition gives the same contraction on the Hom complexes.
On Chevalley degree $p$ its homotopy is $h_C=(-1)^ph$. Treat
$\delta$ as a perturbation of $(-1)^pd_{\mathrm{Hom}}$ and define
\begin{equation}\label{eq:vol2-current-transfer}
 p'=p(1+\delta h_C)^{-1}
    =\sum_{r\geq0}(-1)^rp(\delta h_C)^r.
\end{equation}
Each perturbation raises Chevalley degree, so this defines a map on
the product totalization. Since $i$ is equivariant and $hi=0$,
$h_C\delta i=0$. The transferred differential is consequently
the ordinary total Chevalley differential on $E$.
Multiplying the geometric series and using the contraction identity
shows that $p'$ intertwines the total differentials. The map $j$ is
a total cocycle in arity zero. Set
$H_p=(-1)^{p(p-1)/2}(p'j)_p$ to obtain
\eqref{eq:vol2-current-hierarchy}.

For the first two coefficients the formulas reduce to
\begin{align}\label{eq:vol2-current-transfer-first}
 H_1(x)&=-p\rho_I(x)hj,\\
 H_2(x,y)&=p\bigl(\rho_I(x)h\rho_I(y)h
                      -\rho_I(y)h\rho_I(x)h\bigr)j.
 \label{eq:vol2-current-transfer-second}
\end{align}
For example, applying $d_E$ to the first gives
$-p\rho_I(x)(1-ip)j=\rho_V(x)pj-pj\rho_M(x)$.
In the second, differentiation of the two copies of $h$ leaves
the two action terms and their commutator, giving
\eqref{eq:vol2-current-second}. The terms containing $ph$, $hi$
or $h^2$ vanish by the side conditions.
\end{proof}

\begin{example}\label{ex:vol2-second-current}
The second homotopy may be nonzero for commuting currents. Let $I$
have basis $z,c$ in degree $-2$, $d,a$ in degree $-1$, and $b$ in
degree zero, with $d_Ic=d$ and $d_Ia=b$. Put $V=kz$, let $p$ be
projection, and take $h(d)=c$, $h(b)=a$. Define commuting cochain
operators by $\rho_x(c)=z$, $\rho_y(a)=d$, with all other values
zero. For $M=I$ and $j=1$ the formulas give
\[
 H_1(x)(d)=-z,\qquad H_1(y)=0,\qquad H_2(x,y)(b)=z.
\]
All other values vanish. Thus $d_EH_2(x,y)(a)=-z$, which is
the nonzero right side of \eqref{eq:vol2-current-second} on $a$.
\end{example}

For an infinite current algebra the same construction requires either
pointwise finiteness of the operations or a specified completed
Chevalley domain. The theorem begins with an equivariant resolved
presentation. Applying it to a geometric boundary action still
requires a map from that action's supported collision complex to
this presentation, preserving its unary path and all three channels.

\section{Several roots and upward extensions}
\label{sec:vol2-root-support}

Return to the Heisenberg state space $V=V_\kappa$ of
\eqref{eq:vol2-fock-currents}. The acting current adds a root to
the source of a module cochain.
This explains why a bound on the roots of the original module is not
a bound on every output of its derived endomorphisms. We make the
coefficient category explicit before proving the bound that does hold.

For a finite set $I$, take a chart $X^I_\sigma$ for each total order
$\sigma$ of $I$, where $X=\mathbb A^1_k$. Glue two charts along the
open set obtained by removing $x_i=x_j$ whenever their orders of
$i,j$ disagree. Denote the result by $\mathscr R_I$. The triple
overlap condition holds: disagreement of two orders on a pair forces
one of them to disagree with any third order on that pair. Thus these
opens define a smooth, generally nonseparated scheme. At a tuple
with equality blocks $B$, its fibre over $X^I$ has
$\prod_B|B|!$ points, one order for each block.

A surjection $\alpha:I\twoheadrightarrow J$ with an order on each
fibre defines $\delta_\alpha:\mathscr R_J\to\mathscr R_I$ by
$x_i=y_{\alpha(i)}$: expand every ordered label of $J$ into its
ordered fibre. Composition concatenates these orders. The maps
glue because the preimage of a chart overlap is the corresponding
overlap on $J$. Let $\mathrm{OSurj}_+$ be the category of nonempty
finite sets and these ordered-fibre surjections. Define
\begin{equation}\label{eq:vol2-full-coefficient-category}
 \mathcal C_+=\varprojlim_{\mathrm{OSurj}_+}
                 D\operatorname{Mod}(\mathscr R_I),
 \qquad \mathcal C=D(k)\times\mathcal C_+,
\end{equation}
using extraordinary inverse images $\delta_\alpha^!$ on right
differential modules. Differential modules on the atlas are defined
by descent from its charts. An object retains the equivalences
$\delta_\alpha^!F_I\simeq F_J$ and all their coherences.
The factor $D(k)$ is the empty coefficient; a positive object has
zero empty coefficient.

For a decomposition $I=A\sqcup B$ let $U_{A,B}$ be the locus on
which every cross pair is distinct, and let $j_{A,B}$ be its open
inclusion. Restriction of orders identifies this open with the
disjoint locus in $\mathscr R_A\times\mathscr R_B$. Set
\begin{equation}\label{eq:vol2-full-convolution}
 (F\star G)_I=
 \bigoplus_{I=A\sqcup B}j_{A,B,*}
             (F_A\boxtimes G_B)|_{U_{A,B}}.
\end{equation}
An empty subset contributes scalar tensor product with its coefficient
in $D(k)$; at $I=\varnothing$ the value is
$F_\varnothing\otimes_kG_\varnothing$.
Both subsets may be empty. This is the all-subset convolution; the
cut product of Section~\ref{sec:vol2-word-coefficients} used only
consecutive decompositions of a fixed word.

\begin{proposition}\label{prop:vol2-full-convolution}
The product \eqref{eq:vol2-full-convolution} is symmetric monoidal.
Its unit $K$ is $k$ at the empty coefficient and zero elsewhere.
Every parenthesized product of $r$ positive objects has the common
formula
\begin{equation}\label{eq:vol2-full-partitions}
 (F_1\star\cdots\star F_r)_I
 =\bigoplus_{I=I_1\sqcup\cdots\sqcup I_r}
      j_{P,*}(F_{1,I_1}\boxtimes\cdots\boxtimes F_{r,I_r})|_{U_P},
\end{equation}
where $P=(I_1,\ldots,I_r)$ and $U_P$ removes all cross-block
diagonals. In particular this coefficient is zero for $r>|I|$.
\end{proposition}
\begin{proof}
Under an ordered collision, a summand whose blocks meet the same
collapsed fibre has zero extraordinary inverse image: that fibre
misses its open locus. Every remaining block is the inverse image
of a unique block downstairs. Open base change
(Lemma~\ref{lem:vol2-open-basechange}) identifies these remaining
summands with the convolution downstairs. This also proves
compatibility with composed collisions.

Nested open conditions impose exactly the cross-block inequalities
of $U_P$. Localization commutes with external products, so all
parenthesizations have \eqref{eq:vol2-full-partitions} as their
common expression. The coherent base-change identifications give
the associativity diagrams. Permutation of the blocks gives the
symmetry with the usual Koszul signs. Only an empty block can
contribute the object $K$, proving the unit assertion. Positivity
forces every $I_j$ to be nonempty, proving the last assertion.
\end{proof}

For a right differential module $M$ on $X$, its main-diagonal family
$iM$ is $\Delta_{I,*}M$ on each ordered chart. The branch maps are
the extraordinary restrictions of this diagonal direct image.
Let $V^r=\omega_X\otimes_O\mathcal V$. The residue pairings of the
Heisenberg fields give the degree-zero Lie object
\[
 \mathfrak g=L_{V_\kappa}=iV^r
\]
for $\star$. Explicitly, on a two-block summand its bracket is the
principal part of $Y(v,z-y)w$, with its right transfer-module
normal derivatives. The field skew symmetry and Jacobi identity
are the antisymmetry and Jacobi identities of this pairing, on
every ordered branch. The local vacuum is the state $1\in V$;
it is distinct from the empty convolution unit $K$.

Write
\[
 N_a=\operatorname{Sym}^a_\star\mathfrak g,\qquad
 A=U_\star\mathfrak g.
\]
Symmetric powers use coinvariants, identified with invariants by
averaging. We use the full derived category of $A$-modules and
the adjoint action on $N_a$ and $A$, without imposing preservation
of conformal weight. PBW symmetrization gives
$A=\bigoplus_{a\geq0}N_a$ as adjoint modules. The usual multilinear
PBW identities apply here: convolution preserves colimits, and
finite symmetric-group coinvariants are exact in characteristic
zero. Since $N_a(I)=0$ for $a>|I|$, this sum is also the product
in compatible coefficient families. On each chart, localization
and diagonal direct image are exact, so each $N_a$ has
differential-module coefficients in degree zero.

\begin{theorem}[The degree-dependent support bound]
\label{thm:vol2-root-bound}
For $w,a\geq1$,
\begin{equation}\label{eq:vol2-root-bound}
 \operatorname{Ext}^q_{\mathfrak g}(N_w,N_a)=0\qquad(q<a-w).
\end{equation}
Consequently, for $w\geq1$ and $q\geq0$,
\begin{equation}\label{eq:vol2-root-window}
 \tau_{\leq q}R\operatorname{Hom}_{\mathfrak g}(N_w,A)
 \simeq\tau_{\leq q}\prod_{0\leq a\leq w+q}
              R\operatorname{Hom}_{\mathfrak g}(N_w,N_a).
\end{equation}
The truncation is applied to the full complexes. In a connective
coefficient resolution, degree-$q$ cycle equations retain Chevalley
columns through $q+1$.
\end{theorem}
\begin{proof}
The coefficient family $N_a$ is a sum of open direct images from
exactly $a$ distinct roots, with diagonal support within each
block. If a coefficient family $B$ is supported on at most
$a-1$ roots, ordinary open adjunction gives
$R\operatorname{Hom}_{\mathcal C}(B,N_a)=0$.
This is an adjunction before imposing module equivariance.

The module Chevalley resolution has columns
\[
 R\operatorname{Hom}_{\mathcal C}
       (\Lambda^n_\star\mathfrak g\star N_w,N_a)[-n].
\]
Their sources have at most $n+w$ roots, so a column vanishes
when $n+w<a$. Otherwise it is $n$-connective: its coefficient
objects are differential modules in degree zero and their
derived Homs have no negative cohomology. The compatible-family
limits preserve this bound. In each total degree only finitely
many columns contribute. This proves \eqref{eq:vol2-root-bound}.

Derived Hom into $A$ is the product of derived Homs into the
$N_a$. Every omitted factor in \eqref{eq:vol2-root-window} starts
above degree $q$; products preserve that lower bound. Finally
column $q+1$ can receive the outgoing differential of a
degree-$q$ cochain, whereas later columns cannot enter its
cycle equation.
\end{proof}

The boundary of this vanishing range is attained. Let
$\mathsf D=\partial_{x_1}$ on polynomial states. It satisfies
$\mathsf D1=0$, $\mathsf DJ=1$ and $[\mathsf D,T]=0$.
Moreover $[\mathsf D,Y(J,u)]=1=Y(\mathsf DJ,u)$.
Normal products and translation extend this derivation identity
to every state. Thus $\phi_\tau=e^{\tau\mathsf D}$ is a
vertex-algebra automorphism over $k[\tau]$, finite on each state,
and induces a Lie derivation of $\mathfrak g$. This is the current
translation subgroup described by Matsuo--Nagatomo
\cite[\S3.3, Theorem 3.6]{MatsuoNagatomo1997}.

\begin{proposition}\label{prop:vol2-upward-current}
Let $m_w:\mathfrak g\star N_w\to N_{w+1}$ be symmetric
multiplication. For every $w\geq1$,
\[
 \gamma_w(x;m)=m_w(\mathsf Dx,m)
\]
defines a nonsplit extension of $N_w$ by $N_{w+1}$ as
$\mathfrak g$-modules, preserving the vacuum action.
\end{proposition}
\begin{proof}
The Chevalley differential of $\gamma_w$ is
\[
 m_w([x,\mathsf Dy]+[\mathsf Dx,y]-\mathsf D[x,y],m)=0.
\]
Hence the action on $N_{w+1}\oplus N_w$ is
$x(n,m)=(xn+\gamma_w(x;m),xm)$. Its off-diagonal vacuum action
vanishes because $\mathsf D1=0$.

Ordinary coefficient adjunction gives
$\operatorname{Hom}_{\mathcal C}(N_w,N_{w+1})=0$, so the underlying
split extension has exactly one coefficient splitting. This
splitting is not equivariant: at $w+1$ distinct roots,
$\gamma_w$ sends $J$ at the acting root and vacua at the other
roots to a nonzero external vacuum. Degree-zero maps between
these degree-zero modules are ordinary equivariant maps by the
nonnegative Chevalley resolution. The extension therefore remains
nonzero in the derived category.
\end{proof}

An unrestricted derived open-support adjunction for these modules
would force this extension to vanish. Such an adjunction cannot be
inferred from the corresponding adjunction for their underlying
coefficient families. For instance, input weight two has a genuine
degree-one output in $N_3$.

\subsection{The extension integrates to an interface}

For any degree-zero Lie object with a natural derivation $\mathsf D$,
the formula
\begin{equation}\label{eq:vol2-triangular-current-action}
 \rho_{0,\tau}(x)m=\operatorname{ad}_x(m)+\tau(\mathsf Dx)m
\end{equation}
defines a module structure on $\operatorname{Sym}_\star\mathfrak g[\tau]$.
Indeed, if $M_z$ is multiplication by $z$, then
\[
 [\operatorname{ad}_x,M_{\mathsf Dy}]=M_{[x,\mathsf Dy]},\qquad
 [M_{\mathsf Dx},\operatorname{ad}_y]=M_{[\mathsf Dx,y]},\qquad
 [M_{\mathsf Dx},M_{\mathsf Dy}]=0.
\]
Their sum is the representation identity. The added term raises
symmetric degree by one. For $w\geq1$ and $r\geq0$, quotienting
the degrees at least $w$ by
those above $w+r$ therefore gives a flat family of modules on
$\bigoplus_{a=w}^{w+r}N_a$. For $r=1$ its extension coordinate is
$\tau\gamma_w$.

Suppose $\mathsf D$ is locally nilpotent and the Rees envelope
$A_\nu=U_{\nu,\star}\mathfrak g$, with
$xy-yx=\nu[x,y]$, is PBW flat. Write $\phi_t=e^{t\mathsf D}$.
PBW gives a unique integral operator $\operatorname{ad}_\nu(x)$
with $\nu\operatorname{ad}_\nu(x)m=xm-mx$. Then
\begin{equation}\label{eq:vol2-rees-current-action}
 \rho_{\nu,\tau}(x)m
 =\operatorname{ad}_\nu(x)m+
   \sum_{j\geq1}\frac{\nu^{j-1}\tau^j}{j!}(\mathsf D^jx)m
\end{equation}
is integral over $k[\nu,\tau]$. After inverting $\nu$ it equals
$\nu^{-1}(\phi_{\nu\tau}(x)m-mx)$.
Commuting left and right multiplications gives
\[
 [L_{\phi x}-R_x,L_{\phi y}-R_y]
   =\nu(L_{\phi[x,y]}-R_{[x,y]}).
\]
Torsion-freeness then proves the representation identity before
inversion. At $\nu=0$ only the $j=1$ shifted term survives, giving
\eqref{eq:vol2-triangular-current-action}.

For an associative algebra $A$ and its automorphisms $\phi_\tau$,
let $B_\tau={}_{\phi_\tau}A$ have ordinary right action and left
action $x\cdot b=\phi_\tau(x)b$. These twisted regular bimodules
compose by
\begin{equation}\label{eq:vol2-interface-composition}
 B_\sigma\otimes_A^L B_\tau\xrightarrow{\sim}B_{\sigma+\tau},
 \qquad a\otimes b\longmapsto\phi_\tau(a)b.
\end{equation}
They are free on either side, so ordinary tensor computes the
derived tensor. The formula is balanced and left equivariant;
its inverse is $c\mapsto1\otimes c$. Both associations of three
factors give $\phi_{\tau+\upsilon}(a)\phi_\upsilon(b)c$.
Thus $B_0=A$ is the unit and $B_{-\tau}$ is inverse to $B_\tau$.
Equation~\eqref{eq:vol2-rees-current-action} is the normalized
left-minus-right action on $B_{\nu\tau}$. This constructs the
algebraic interface underlying the upward extension. A field
realization must also supply its collision and interchange maps.

Choose semi-free resolutions $P_w\to N_w$ and form the dg category
$\mathcal E(w,a)=\operatorname{Hom}_{\mathfrak g}(P_w,P_a)$.
Composition is composition of resolution maps. The dg Yoneda
functor identifies the subcategory generated by the $P_w$ under
finite cones, shifts and retracts with $\operatorname{Perf}(\mathcal E)$:
it is fully faithful on the generators, this property extends by
those operations in each variable, and represented modules generate
the perfect modules. This is the precise generated-sector conclusion
of dg Morita theory \cite[Theorem 3.8(b)]{KellerDG2006}.
It supplies no generation statement for all physical boundaries.

\section{Operators at a polynomial root}
\label{sec:vol2-polynomial-operators}

Forgetting the current action leaves an endomorphism of the moving
differential module. The allowed endomorphisms must be specified
before this forgetful map can be computed. Put
\begin{equation}\label{eq:vol2-affine-operators}
 N=\operatorname{End}_O(\mathcal V),\qquad
 \nabla_NA=\nabla A-A\nabla,\qquad
 G(N)=[N\xrightarrow{\nabla_N}N],
\end{equation}
where $G(N)$ is a complex of $k$-vector spaces in degrees zero and
one. The identity $\nabla_N(fA)=f'A+f\nabla_NA$ for $f\in O$
makes $N$ a differential module. In a monomial
state basis, $N$ consists of column-finite matrices over $O$: each
input has a finite polynomial output, with no common bound on all
outputs. For example, $v_j\mapsto y^jv_j$ lies in $N$ but not in
$O\otimes_k\operatorname{End}_k(V)$.

\begin{lemma}\label{lem:vol2-internal-operator}
For an $O$-free left $D_y$-module $M$ in degree zero, the internal
derived endomorphism for $\otimes_O^L$ is
$\operatorname{End}_O(M)$ in degree zero, with its commutator
connection. This holds without a finite-rank hypothesis.
\end{lemma}
\begin{proof}
Evaluation gives
\[
 \operatorname{Hom}_{D_y}(P,\operatorname{Hom}_O(M,Q))
 =\operatorname{Hom}_{D_y}(P\otimes_OM,Q).
\]
The connection on the Hom is forced by this equality. Tensoring
with $M$ is exact. Its right adjoint is exact as well: after choosing
an $O$-basis, its underlying module is a product of copies of $Q$,
and products of modules are exact. Both functors preserve
quasi-isomorphisms, so their adjunction descends to derived
categories. Taking $Q=M$ proves the assertion.
\end{proof}

This affine internal Hom need not commute with localization of an
infinite-rank source. The supported calculation below uses the
restriction of a specified coefficient object, not a new internal
Hom formed on the open set.

\subsection{The supported degrees}

Use right densities $dy$ on $X$ and $dy\wedge du$ on $Y=X^2$.
The monoidal left-to-right normalization for the diagonal tensor is
$M\mapsto M^r[1]$, with $M^r=\omega_X\otimes_OM$; its units are
$J_X=\omega_X[1]$ and $J_Y=\omega_Y[2]$.
Subscripts $\mathcal D$ on Hom denote right differential modules.
For any left $D_y$-module $L$ with connection $\nabla_L$, write
$G(L)=[L\xrightarrow{\nabla_L}L]$ in degrees zero and one.
Its normal direct image has the model
\[
 \Delta_+L=L[\eta],\qquad \partial_u=\eta,\qquad
 u=-\partial_\eta.
\]
Its normal Spencer complex is
$[L[\eta]\xrightarrow{\eta}L[\eta]]$ in degrees zero and one.
Project the degree-one term to its constant coefficient and include
that coefficient back. The homotopy sends $\eta^m\ell$ in degree
one to $\eta^{m-1}\ell$ for $m\geq1$ and to zero for $m=0$.
It satisfies $dh+hd=1-ip$ on every monomial. In tangential form
degree $a$, give both the normal differential and this homotopy
the sign $(-1)^a$. The tangential terms then cancel, proving
\begin{equation}\label{eq:vol2-affine-supported-hom}
 R\operatorname{Hom}_{\mathcal D_Y}
       (J_Y,\Delta_*(L^r[1]))\simeq G(L)[-2].
\end{equation}
One shift comes from the normal Spencer complex, and the other is
the density shift $[1-2]$.

In particular, suppose a coefficient object $E_2$ has
$\Delta^!E_2=N^r[1]$. Recollement and open adjunction give
\begin{equation}\label{eq:vol2-affine-restriction-fibre}
 \operatorname{Fib}\left(
 R\operatorname{Hom}_{\mathcal D_Y}(J_Y,E_2)
 \longrightarrow R\operatorname{Hom}_{\mathcal D_U}(J_U,j^!E_2)
 \right)\simeq G(N)[-2].
\end{equation}
The fibre has cohomology only in degrees two and three. Restriction
therefore induces an equivalence on the unshifted mapping spaces.
For a source $J_Y[-d]$, the fibre instead is $G(N)[d-2]$.
Thus shifted sources detect supported operators invisible to the
unshifted comparison.

For example, let $P_0$ be projection to the constant state.
Both $P_0$ and $\mathsf D=\partial_{x_1}$ commute with $T$ and are
horizontal endomorphisms. Their commutator is nonzero, since
$(P_0\mathsf D-\mathsf DP_0)J=1$. They give classes in the
degree-zero supported fibre for $d=2$. The inclusion of that fibre
need not be injective on cohomology; realizing the operators by
full boundary actions is an additional question.

\subsection{Translation towers}

Give $x_n$ energy $n$ and let $V_+$ be the positive-energy states.
If $x_N$ is the largest variable occurring in a polynomial $f$,
the coefficient of $x_{N+1}$ in $Tf$ is
$N\partial_{x_N}f$. Descent on $N$ proves $\ker T=k1$.
Choose a graded complement $U$ of $TV_+$ in $V_+$, containing
$J$ and $J^2$. Induction on energy gives
\begin{equation}\label{eq:vol2-root-towers}
 V=k1\oplus k[T]\otimes_kU,\qquad
 \mathcal V\cong O\oplus P,\qquad P=D_y\otimes_kU.
\end{equation}
The second map sends $f(y)\partial_y^n\otimes u$ to $f(y)T^nu$;
PBW and the connection identify the two sides. If $p(E)$ is the
partition number, $\dim U_1=1$ and
$\dim U_E=p(E)-p(E-1)$ for $E\geq2$.

On one tower write $t$ for $T$, so its connection is
$\partial_y+t$ on $k[y,t]$. Define
\begin{align}\label{eq:vol2-tower-reduction}
 \operatorname{red}\left(\sum_n f_n(y)t^n\right)
    &=\sum_n(-1)^nf_n^{(n)}(y),\\
 H_P(f(y)t^n)&=\sum_{j=0}^{n-1}(-1)^jf^{(j)}(y)t^{n-1-j},
 \qquad H_P(f)=0.\label{eq:vol2-tower-homotopy}
\end{align}
If $j$ includes the constant tower coefficient, cancellation gives
\[
 \operatorname{red}\nabla=0,\qquad
 \nabla H_P=1-j\operatorname{red},\qquad H_P\nabla=1.
\]
On the vacuum use evaluation at $y=0$ and
$H_O(f)=\int_0^yf(s)\,ds$. Together these maps define a contraction
\begin{equation}\label{eq:vol2-state-contraction}
 G(\mathcal V)\rightleftarrows k\oplus W[-1],
 \qquad W=k[y]\otimes_kU.
\end{equation}
Write $H_{\mathcal V}$ for its homotopy and
$\operatorname{red}_+$ for reduction of its positive-state summand.

Evaluation at the vacuum splits $N$ as differential modules:
\[
 N=\sigma(\mathcal V)\oplus N_0,
 \qquad \sigma(v)=vP_0,\qquad N_0=\{A:A1=0\}.
\]
For $A\in N_0$ define an $O$-linear operator $R(A)$ by
\begin{equation}\label{eq:vol2-column-recursion}
 R(A)1=R(A)u=0,\qquad
 R(A)(T^{n+1}u)=\nabla R(A)(T^nu)-A(T^nu).
\end{equation}
Every column requires finitely many operations, so $R(A)$ is an
allowed endomorphism. If $F\in N_0$, let $\Phi(F)$ be the
horizontal operator with $\Phi(F)(T^nu)=\nabla^nF(u)$ and
$\Phi(F)1=0$. Recursion then gives
\[
 \nabla_NR(A)=A,\qquad R(\nabla_NF)=F-\Phi(F).
\]

\begin{theorem}\label{thm:vol2-polynomial-operators}
The operator complex has a contraction
\[
 G(N)\rightleftarrows B_0\oplus W[-1],
 \qquad B_0=\operatorname{End}_{D_y}(\mathcal V).
\]
Its degree-one projection is $A\mapsto\operatorname{red}_+(A1)$,
and its homotopy is
\begin{equation}\label{eq:vol2-operator-homotopy}
 H_N(A)=\sigma(H_{\mathcal V}(A1))+R(A-\sigma(A1)).
\end{equation}
For a basis $\mathcal U$ of $U$, in the order $(O,P)$ one has
\begin{equation}\label{eq:vol2-horizontal-operator-ring}
 B_0=\begin{pmatrix}
 k&\displaystyle\prod_{u\in\mathcal U}k[y]\\
 0&\operatorname{Mat}_{\mathrm{col},\mathcal U}(D_y^{\mathrm{op}})
 \end{pmatrix}.
\end{equation}
The complex computes $R\operatorname{Hom}_{D_y}(\mathcal V,\mathcal V)$.
\end{theorem}
\begin{proof}
The preceding contractions apply to the two summands of $N$.
The degree-zero projection is the horizontal operator whose value
on the vacuum is the constant-state coefficient of $A1$ at $y=0$
and whose value on each primitive $u$ is $A(u)$. The degree-one
section is $w\mapsto\sigma(jw)$. The displayed identities prove
$dH_N+H_Nd=1-ip$, $pi=1$, $H_Ni=pH_N=H_N^2=0$.

The free module $P$ has arbitrary images on its basis $U$.
The connection $\partial_y+t$ is injective on a polynomial tower,
so $\operatorname{Hom}_{D_y}(O,P)=0$. A map $O\to O$ is a
scalar; a map $P\to O$ chooses a polynomial for each generator.
Endomorphisms of the direct sum of copies of $D_y$ are the
column-finite matrices over $D_y^{\mathrm{op}}$. This proves
\eqref{eq:vol2-horizontal-operator-ring}.

Finally $P$ is projective and
$[D_y\xrightarrow{(-)\partial_y}D_y]$ resolves $O$ in degrees
$-1,0$. Equivalently, the one-dimensional relative Spencer
resolution for the $O$-free module $\mathcal V$ gives the
commutator complex \eqref{eq:vol2-affine-operators}. Thus it is
a model of the stated derived Hom.
\end{proof}

Its cohomology product has $W^2=0$. A matrix
$b=\left(\begin{smallmatrix}a&v\\0&C\end{smallmatrix}\right)$
acts on the left of $W$ through $C$ and on the right by $a$.
On $D_y/\partial_yD_y=k[y]$, right multiplication by
$\partial_y$ acts as $f\mapsto-f'$. This identifies the
cohomology algebra. The contraction was a contraction of complexes;
it has made no assertion about higher compositions.

\subsection{The charge line and the comparison homotopy}

Define the polynomial current cochain algebra
\[
 Q_\kappa=k[y,x_1,x_2,\ldots]\otimes
                  \Lambda(\xi_0,\xi_1,\ldots),\qquad
 dx_n=\kappa n\xi_n,\qquad d\xi_n=0.
\]
The ghosts $\xi_n$ have degree one. Translation extends by
$T\xi_n=(n+1)\xi_{n+1}$. Take its root de Rham complex
$\mathcal C_\kappa=Q_\kappa\oplus dy\,Q_\kappa$, with total
differential $d+dy(\partial_y+T)$ and forms written on the left;
in particular, $d(dy\,q)=-dy\,dq$. Ghost deletion and multiplication
define the cochain map
\begin{equation}\label{eq:vol2-current-forgetting}
 \rho_\kappa:\mathcal C_\kappa\longrightarrow G(N),
 \qquad f\longmapsto L_f,
\end{equation}
with every positive ghost monomial sent to zero. This is the
coefficient map that forgets the regular-current equivariance.
It keeps the differential-module connection.

\begin{proposition}\label{prop:vol2-charge-image}
The current complex has cohomology $k$ in degree zero and $k\gamma$
in degree one, where
\[
 \gamma=\xi_0+dy\,J/\kappa.
\]
The maps induced by $\rho_\kappa$ are
$k\hookrightarrow B_0$, $a\mapsto a\operatorname{id}$, and
$k\gamma\hookrightarrow W$, $\gamma\mapsto J/\kappa$.
\end{proposition}
\begin{proof}
Contract every Koszul pair $(x_n,\xi_n)$, $n>0$, by the Euler
homotopy: on positive total polynomial-plus-ghost degree divide
$\sum_{n>0}x_n(\kappa n)^{-1}\partial_{\xi_n}$ by that degree.
The remaining complex is $k[y]\otimes\Lambda(\xi_0)$.
Perturb by root translation. The perturbation series terminates
after one root form, since $dy^2=0$, and the induced translation
on the zero ghost vanishes because $T\xi_0=\xi_1$.
The transferred differential is $dy\partial_y$, while the
corrected ghost representative is $\gamma$. Directly, its two
differentials are $dy\xi_1$ and $-dy\xi_1$.
Polynomial integration in $y$ leaves the two stated classes.
Finally $\rho_\kappa\gamma=dy\,L_J/\kappa$, whose normal form
is $J/\kappa$ by Theorem~\ref{thm:vol2-polynomial-operators}.
\end{proof}

\begin{corollary}\label{cor:vol2-charge-lift}
For $A\in N$, the class of $dy\,A$ lifts through
$\rho_\kappa$ exactly when $\operatorname{red}_+(A1)\in kJ$.
If this normal form is $cJ/\kappa$, a lift is $c\gamma$, with
comparison homotopy $H_N(A-cL_J/\kappa)$. The fibre satisfies
\[
 H^1\operatorname{Fib}(\rho_\kappa)=B_0/k\operatorname{id},
 \qquad H^2\operatorname{Fib}(\rho_\kappa)=W/kJ,
\]
and has no other cohomology.
\end{corollary}
\begin{proof}
The image criterion follows from Proposition~\ref{prop:vol2-charge-image}.
When it holds, the degree-one projection of $A-cL_J/\kappa$
vanishes, so the contraction identity gives
$\nabla_NH_N(A-cL_J/\kappa)=A-cL_J/\kappa$.
The two injections in that proposition give the fibre groups.
\end{proof}

For example, $L_{J^2}$ and $yL_J$ do not lift: their normal forms
are $J^2$ and $yJ$. A scalar charge $q$ does integrate, with
$J_0=q\operatorname{id}$ and
$\nabla_q=\partial_y+T+(q/\kappa)L_J$. Only the $n=1$ moving-mode
equation changes, by
$[(q/\kappa)L_J,\kappa\partial_{x_1}]=-q$.
When $q\ne0$ the vector $1$ is no longer horizontal.

There are two distinct completions. For a new parameter $g$,
suppose $A(g)\in gN[[g]]$ kills $1$. The map
\[
 U_g1=1,\qquad U_g(T^nu)=(\partial_y+T+A(g))^nu
\]
extends $k[y][[g]]$-linearly and intertwines the two connections.
It is the identity modulo $g$, so it has a $g$-adic inverse.
Every coefficient of every column is a polynomial. This trivializes
the connection with a chosen horizontal vacuum, but leaves its
horizontal automorphisms as stabilizers.

Completing the root instead gives $\widehat{\mathcal V}=V[[y]]$.
The gauge $e^{-yT}$ turns its connection into $\partial_y$; on
continuous endomorphisms it acts by conjugation. Formal integration
is onto, hence
\begin{equation}\label{eq:vol2-completed-root-operators}
 H^0G(\widehat N)=\operatorname{End}_k(V),\qquad
 H^1G(\widehat N)=0,
 \qquad
 \widehat N=\operatorname{End}_{k[[y]],\mathrm{cont}}(V[[y]]).
\end{equation}
In this category the charge has the primitive
\[
 g_\kappa(y)=\frac1\kappa\sum_{n\geq0}
          \frac{(-1)^ny^{n+1}}{(n+1)!}T^nJ,
 \qquad (\partial_y+T)g_\kappa=J/\kappa.
\]
Its infinitely many nonzero state coefficients prevent it from
being a polynomial-root primitive. Completion has changed the
coefficient category and the cohomology.

\section{A nonformal derived operator algebra}
\label{sec:vol2-root-composition}

One free translation tower and the vacuum already detect the
composition missing from the linear contraction. Write $D=D_y$
and choose the tower generated by $J$. Let $P=D$ and
$O=D/D\partial$. Resolve $O$ by
\[
 Q=[D\xrightarrow{R_\partial}D],\qquad
 Q^{-1}=D,\quad Q^0=D,\quad R_\partial(a)=a\partial.
\]
The Weyl PBW basis proves injectivity of $R_\partial$ and identifies
its cokernel with $O$. Put
\[
 B=D^{\mathrm{op}}=k\langle y,p\rangle/(py-yp+1),\qquad
 \mathscr E=\operatorname{End}_D(P\oplus Q).
\]
Thus $p$ represents right multiplication by $\partial$, with
$[p,y]=-1$. Products in $\mathscr E$ are matrix compositions, the
rightmost map acting first. The free cells $P,Q^0,Q^{-1}$ have
degrees $(0,0,-1)$; the only nonzero entry of their differential
is $(Q^0,Q^{-1})=p$. Hence
\[
 \delta X=dX-(-1)^{|X|}Xd
\]
makes $\mathscr E$ a dg algebra of ranks $2,5,2$ over $B$ in
degrees $-1,0,1$. It models the marked derived endomorphisms of
$P\oplus O$. The splitting \eqref{eq:vol2-root-towers} embeds it
as a corner in a projective model of the full Heisenberg module.

\subsection{Polynomial division gives a contraction}

Use normal monomials $y^ap^b$. Define
\begin{align}\label{eq:vol2-weyl-remainders}
 r_L(y^ap^b)&=\partial_y^b(y^a),&
 \ell(y^ap^b)&=\sum_{j=0}^{\min(a,b-1)}
                    \frac{a!}{(a-j)!}y^{a-j}p^{b-1-j},\\
 r_R\left(\sum_bf_b(y)p^b\right)&=f_0(y),&
 r\left(\sum_bf_b(y)p^b\right)&=\sum_{b\geq1}f_b(y)p^{b-1}.
 \notag
\end{align}
Empty sums are zero. The relation $yp=py+1$ gives
\begin{equation}\label{eq:vol2-weyl-divisions}
 F=p\ell(F)+r_L(F)=r(F)p+r_R(F).
\end{equation}
In particular $pB=\ker r_L$ and $Bp=\ker r_R$, and neither
contains $1$.

The block $\operatorname{Hom}(P,Q)$ is
$[B\xrightarrow{p(-)}B]$ in degrees $-1,0$. Its projection is
$r_L$, its section includes polynomials, and its homotopy is
$\ell$. The block $\operatorname{Hom}(Q,P)$ is
$[B\xrightarrow{-(-)p}B]$ in degrees $0,1$; its maps are
$r_R$, polynomial inclusion and $-r$.

On $\operatorname{End}(Q)$ write a degree-zero map as $(a,c)$,
acting on $Q^0,Q^{-1}$, and a degree-one map as $z:Q^{-1}\to Q^0$.
Then
\begin{equation}\label{eq:vol2-vacuum-block}
 \delta h=(ph,hp),\qquad \delta(a,c)=pc-ap.
\end{equation}
With $I(f)=\int_0^yf(t)dt$, its projection to $k$ is
$r_L(a)(0)$, the section is $s\mapsto(s,s)$, and its homotopy is
\begin{equation}\label{eq:vol2-vacuum-block-homotopy}
 H(a,c)=\ell(a),\qquad
 H(z)=(-I(r_Lz),\ell(z)-I(r_Lz)).
\end{equation}
It is zero on degree $-1$.

\begin{proposition}\label{prop:vol2-marked-contraction}
These block maps, and the identity on $\operatorname{End}(P)=B$,
form a contraction of $\mathscr E$ onto its cohomology with
$\pi i=1$, $H^2=Hi=\pi H=0$. In the order $(P,O)$ the result is
\[
 H^0\mathscr E=\begin{pmatrix}B&0\\B/pB&k\end{pmatrix},
 \qquad
 H^1\mathscr E=\begin{pmatrix}0&B/Bp\\0&0\end{pmatrix},
 \qquad H^{\ne0,1}\mathscr E=0.
\]
Every map in this contraction is polynomial on each input.
\end{proposition}
\begin{proof}
The off-diagonal identities are the two divisions in
\eqref{eq:vol2-weyl-divisions}. For the vacuum block, the maps
$(a,c)\mapsto r_L(a)$ and $z\mapsto r_L(z)$ map the complex to
$[k[y]\xrightarrow{-\partial_y}k[y]]$.
A section sends $f$ to $(f,f)$ and $g$ to $g$, since
$pf-fp=-f'$. In the kernel, write $a=pa_1$, $z=pz_1$.
The differential becomes
\[
 B\longrightarrow B\oplus B\longrightarrow B,
 \qquad h\longmapsto(h,hp),\quad(a_1,c)\longmapsto c-a_1p.
\]
The homotopy sends $(a_1,c)$ to $a_1$ and $z_1$ to $(0,z_1)$.
Contract the remaining derivative by $-I$. Composition gives
\eqref{eq:vol2-vacuum-block-homotopy} and
$\delta H+H\delta=1-i\pi$. The identity
$\ell(ap)=\ell(a)p+r_L(a)$ verifies the middle degree directly.
Integration vanishes at zero and $\ell$ kills polynomials, giving
the side conditions. All maps preserve the two object idempotents.
\end{proof}

\subsection{The full Massey coset}

Consider the closed maps
\[
 a:P\to Q,\quad a^0=1;\qquad
 b:Q\to P[1],\quad b^{-1}=1;\qquad
 c:P\to P,\quad c=p.
\]
Their degrees are $0,1,0$, and $ba=0$, while
$cb=\delta h$ for $h:Q^0\to P$, $h=-1$.
For primitives $\delta h=cb$, $\delta k=ba$, use the defining
Massey representative $ha-ck$. Taking $k=0$ gives $-1_P$.

\begin{theorem}\label{thm:vol2-root-massey}
The exact composable Massey coset is
\begin{equation}\label{eq:vol2-root-massey}
 \langle c,b,a\rangle=-1+pB\ \subset H^0\operatorname{End}(P)=B.
\end{equation}
It excludes zero. Both $\mathscr E$ and the full polynomial-root
derived endomorphism algebra of $\mathcal V$ are nonformal.
\end{theorem}
\begin{proof}
A second primitive of $cb$ differs by a degree-zero cocycle in
$\operatorname{Hom}(Q,P)$. There are no such nonzero cocycles,
since its differential is right multiplication by $-p$.
A primitive of $ba=0$ can differ by any element of $B$.
The indeterminacy is therefore exactly $pB$. The remainder map
$r_L$ annihilates this ideal and sends $1$ to $1$, so zero is
excluded.

In a dg algebra with zero differential, a defined triple Massey
product contains zero, by choosing zero primitives. A dg
quasi-isomorphism preserves the Massey coset: it sends defining
systems to defining systems, and surjectivity on cohomology
recovers all changes of primitives, whose images are the two
indeterminacy subspaces. Applying this in both directions along
a quasi-isomorphism zigzag proves nonformality; equivalently one
may use the minimal-model criterion of
\cite[\S\S3.3--3.4]{KellerAInfinity2001}.

For the full module, sandwich a defining system between the
idempotents of the chosen $P\oplus O$ summand. Its $P\to P$
projection is a defining system above. Additional matrix blocks
do not change the projected indeterminacy: the degree-zero
$O\to P$ cohomology is still zero, and the $P\to P$ corner
is $B$. A zero total Massey class would thus give zero in
\eqref{eq:vol2-root-massey}, which is impossible.
\end{proof}

Replacing $a$ by the representative $a_f=f(y)$ gives
$\langle c,b,a_f\rangle=-f(y)+pB$, nonzero for every nonzero
polynomial $f$. The obstruction concerns the whole coset, not a
particular transferred ternary coefficient. At level $\kappa$,
$a$ is the restriction of $J_1/\kappa=\partial_{x_1}$ to the
$J$-tower, $b$ is its charge class $[J]$, and $c$ is its translation.
Thus this composition uses three of the operators already present
in the current comparison.

\subsection{Higher compositions and the ordinary centre}

The polynomial contraction gives all higher operations explicitly.
Let $\mathcal H=H^*\mathscr E$, and let
$s:\mathscr E\to\mathscr E[1]$ have degree $-1$. On the tensor
coalgebra of suspended matrices use
\[
 b_1(sx)=-s\delta x,\qquad
 b_2(sx,sy)=(-1)^{|x|}s(xy),\qquad h_s=-sHs^{-1}.
\]
Extend $b_1,b_2$ as coderivations, and write $i_s,\pi_s$ for
the suspended inclusion and projection. The tensor contraction is
\[
 \mathbf H|_{\mathscr E[1]^{\otimes n}}
 =\sum_{j=1}^n(i_s\pi_s)^{\otimes(j-1)}\otimes h_s
                                \otimes1^{\otimes(n-j)},
\]
with graded tensor signs. Let $i_T,\pi_T$ be the tensor extensions
of $i_s,\pi_s$.
Homological perturbation gives
\begin{align}\label{eq:vol2-root-ainfinity}
 b_{\mathcal H}&=\pi_Tb_2(1+\mathbf Hb_2)^{-1}i_T,\notag\\
 \mathbf I&=(1+\mathbf Hb_2)^{-1}i_T,&
 \mathbf P&=\pi_T(1+b_2\mathbf H)^{-1},&
 \mathbf H'&=\mathbf H(1+b_2\mathbf H)^{-1}.
\end{align}
Each binary merger lowers length, so at input length $n$ these
inverse series have at most $n-1$ merger stages. Expanding the
contraction identity yields $b_{\mathcal H}^2=0$,
$(b_1+b_2)\mathbf I=\mathbf I b_{\mathcal H}$ and
$\mathbf P\mathbf I=1$, with comparison homotopy $\mathbf H'$.
Deconcatenation and the tensor homotopy give the coalgebra
identities. Tensoring over the two scalar idempotents retains
composable strings; the side conditions give strict units.
This constructs a minimal $A_\infty$ category in every arity.
Its deciding value is
\[
 (b_{\mathcal H})_3(sc,sb,sa)=-s1_P:
\]
the merger $ba$ vanishes, while the other merger has homotopy
$h=-1$. The Massey coset prevents removal of all higher
operations by a different minimal model.

Its ordinary bimodule centre has a different answer. The free
module $P=D$ generates $\operatorname{Perf}(D)$ and $Q$ is a
two-term cone of free modules. Thus $P\oplus Q$ is another
generator. The functors
\[
 R\operatorname{Hom}_D(P\oplus Q,-),\qquad
 -\otimes_{\mathscr E}^L(P\oplus Q)
\]
are inverse on the generated perfect categories: evaluation
is an equivalence on the generator and extends by cones,
shifts and retracts. Derived Morita invariance of Hochschild
cochains \cite[\S5.4]{KellerDG2006} gives
\begin{equation}\label{eq:vol2-marked-root-centre}
 \operatorname{HH}^*(\mathscr E)
   \cong\operatorname{HH}^*(D^{\mathrm{op}})
   \cong\operatorname{HH}^*(D)=k.
\end{equation}
For the last equality, resolve the Weyl algebra as a bimodule
by its two-generator PBW Koszul complex. Its cochains have
differential given by $[\partial,-]=\partial_y$ and
$[y,-]=-\partial_\partial$, hence are the polynomial de Rham
complex after a change of one exterior generator. Exactness of
the bimodule resolution follows from its associated graded
polynomial Koszul resolution and the bounded-below PBW filtration.
Euler contraction, divided by positive polynomial-plus-form
degree, leaves only the constant class. This proves
\eqref{eq:vol2-marked-root-centre} without a root completion.

Nonformality of a selected derived operator algebra is compatible
with a scalar ordinary centre of its perfect module sector.
The centre of the full geometric boundary theory requires its
own category and its own action comparison.

\section{The level ring and the forgetful fibre}
\label{sec:vol2-integral-current-fibre}

The nonzero-level calculation divided by $\kappa$. To retain the
singular level, work over $R=k[\lambda]$ and use a separate symbol
for the Koszul differential:
\[
 \mathcal Q=k[x_1,x_2,\ldots]\otimes\Lambda(\xi_0,\xi_1,\ldots),
 \qquad d_0x_n=n\xi_n,\quad d_0\xi_n=0.
\]
Translation is $Tx_n=nx_{n+1}$ and
$T\xi_n=(n+1)\xi_{n+1}$. It commutes with $d_0$ on generators
and hence everywhere. With forms on the left define
\begin{equation}\label{eq:vol2-integral-current-complex}
 \begin{split}
 \mathcal C_\lambda&=\mathcal Q_R[y]\oplus dy\,\mathcal Q_R[y],\qquad
        \mathcal Q_R=R\otimes_k\mathcal Q,\\
 \mathsf D_\lambda(q+dy\,r)
   &=\lambda d_0q+dy\bigl((\partial_y+T)q-\lambda d_0r\bigr).
 \end{split}
\end{equation}
This is the polynomial-root current complex at the universal level.
There are only two root-form degrees, so its totalization needs
no additional convention. The parameter $\lambda$ has degree zero.

\subsection{Translation coinvariants before specialization}

Let $\mathcal Q_+$ be the augmentation ideal. Translation is
injective on each $\mathcal Q_+^q$. To see this, choose the largest
occupied index of an even or an odd variable. Its translate
introduces a new variable of that type, with coefficient its
corresponding partial derivative. No other term can cancel that
new variable. Descending on the occupied indices removes every
variable, including $\xi_0$; only the unit in the full algebra
can be translation-invariant.

Put $K^q=k[y]\otimes\mathcal Q_+^q$. The operator
$\nabla=\partial_y+T$ is also injective on $K^q$: the highest
power of $y$ in a nonzero kernel element would have a nonzero
coefficient killed by $T$. Define the vector spaces
\begin{equation}\label{eq:vol2-integral-coinvariants}
 \mathsf W^q=K^q/\nabla K^q,\qquad
 \overline d[q]=[d_0q],\qquad
 \mathsf B^q=\operatorname{im}(\overline d:
                   \mathsf W^{q-1}\to\mathsf W^q)\quad(q\geq1).
\end{equation}
They are $k$-vector spaces. The quotient is not an $O$-module
quotient, since $[\nabla,y]=1$.

For computation, choose translation primitives in each state energy
and exterior degree. These pieces are finite dimensional, including
the energy-zero ghost. The same tower division as
\eqref{eq:vol2-tower-reduction} gives the finite normal form
\[
 [f(y)T^nu]=(-1)^nf^{(n)}(y)[u].
\]
Apply $d_0$, decompose its result into towers, and reduce to compute
$\overline d$. In degree zero $\mathsf W^0=W$ of
\eqref{eq:vol2-state-contraction}, as vector spaces.

\begin{lemma}\label{lem:vol2-coinvariant-cohomology}
The cohomology of $(\mathsf W,\overline d)$ is the
one-dimensional span $k\cdot[J]$ in degree zero and zero in
positive degrees.
\end{lemma}
\begin{proof}
The Euler homotopy for the positive Koszul pairs is
$\sum_{n>0}(x_n/n)\partial_{\xi_n}$ divided by positive total
degree. It leaves only $\xi_0$ in the augmentation ideal, so
$H^*(K,d_0)=k[y]\xi_0$ in degree one. Translation acts trivially
there, because $T\xi_0=\xi_1=d_0J$. Thus $\nabla$ acts on this
cohomology as $\partial_y$, which is onto with kernel $k\xi_0$.
The exact sequence of complexes
\[
 0\longrightarrow K\xrightarrow{\nabla}K
       \longrightarrow\mathsf W\longrightarrow0
\]
now proves the assertion. The connecting preimage of $\xi_0$ is
$[J]$, since $d_0J=\nabla\xi_0$.
\end{proof}

\begin{theorem}\label{thm:vol2-integral-current}
There is an integral quasi-isomorphism of complexes
\begin{equation}\label{eq:vol2-integral-current-reduction}
 \mathcal C_\lambda\simeq
 R\oplus(\mathsf W\otimes_kR,\lambda\overline d)[-1].
\end{equation}
Consequently
\begin{align}\label{eq:vol2-integral-current-cohomology}
 H^0\mathcal C_\lambda&=R,&
 H^1\mathcal C_\lambda&=R\Gamma_\lambda,
       &\Gamma_\lambda&=\lambda\xi_0+dy\,J,\notag\\
 H^{q+1}\mathcal C_\lambda&=\mathsf B^q\otimes_kR/(\lambda)
       &&(q\geq1).
\end{align}
\end{theorem}
\begin{proof}
The constant-state root complex contracts to $R$. On the positive
part the map $q+dy\,r\mapsto[r]$ has target differential
$-\lambda\overline d$, the sign of the shift. In its kernel
$r=\nabla s$ has a unique polynomial solution $s$. The homotopy
sends $dy\nabla s$ to $s$ and the root-degree-zero part to zero.
The terms $\lambda d_0s$ cancel in $\mathsf D_\lambda H+
H\mathsf D_\lambda$, leaving the identity. This proves the
quasi-isomorphism without inverting the level.

Multiplication by $\lambda$ is injective on each extended vector
space. Thus the cycles of $\lambda\overline d$ are the cycles of
$\overline d$, and its boundaries are $\lambda\mathsf B^q$.
Lemma~\ref{lem:vol2-coinvariant-cohomology} gives the groups.
Finally $\Gamma_\lambda$ is closed, since its two nonzero
differentials are $\lambda dy\xi_1$ and $-\lambda dy\xi_1$,
and its image in the quotient is $[J]$.
\end{proof}

This is a linear quotient calculation. Multiplication by $y$ does
not descend to $\mathsf W$, and no multiplicative quotient of
$\mathcal C_\lambda$ has been asserted.

\subsection{An equivalence of arrows}

Over the level ring use the actual full operator modules
\[
 \mathcal N_R=\operatorname{End}_{R[y]}(R[y]\otimes_kV),
 \qquad
 \mathcal B_R=\operatorname{End}_{D_y\otimes R}(R[y]\otimes_kV).
\]
These permit different parameter degrees in different columns;
they are not replaced by tensoring $N$ or $B_0$ with $R$.
The tower and column recurrences work directly over $R$, giving
\[
 G(\mathcal N_R)\simeq
 \mathcal B_R\oplus(\mathsf W^0\otimes_kR)[-1].
\]
Let $\rho:\mathcal C_\lambda\to G(\mathcal N_R)$ be ghost deletion
followed by multiplication. Define its fibre by
\begin{equation}\label{eq:vol2-integral-fibre-convention}
 \mathcal F_\lambda^n=\mathcal C_\lambda^n\oplus
                  G(\mathcal N_R)^{n-1},\qquad
 d(q,A)=(\mathsf D_\lambda q,\rho q-d_GA).
\end{equation}
Its second coordinate is the retained comparison path.

The contractions of source and target alone would not identify
this fibre. We also need their comparison on the arrow $\rho$.
Vacuum evaluation splits
$\mathcal N_R=\sigma(\mathcal V_R)\oplus\mathcal N_{0,R}$,
with $\sigma(v)=vP_0$ and $\mathcal N_{0,R}1=0$.
On the first summand, evaluation of $\rho$ is exactly ghost
deletion, so the translation quotients commute.

On the second summand, column recursion contracts positive root
degree. Its remaining horizontal component is homotopic to
$a\mapsto a(1-P_0)$ on constants and zero on the positive
quotient. Here is the homotopy. On $dy\,v$ with no ghosts, take
the horizontal operator whose vacuum value is zero and whose
primitive-column value on $u$ is $H_{\mathcal V}(v)u$; extend by
$T^nu\mapsto\nabla^n(H_{\mathcal V}(v)u)$. On positive ghost
inputs it is zero. The identity
\[
 H_{\mathcal V}\nabla f=f-\epsilon(f)(0)
\]
shows that its composite with $\mathsf D_\lambda$ is exactly
the horizontal component of multiplication by $f$ minus
$\epsilon(f)(0)(1-P_0)$. Here $\epsilon$ extracts the constant
state. The first summand contributes $aP_0$. Thus the arrow is
equivalent, with its homotopy supplied, to
\begin{equation}\label{eq:vol2-integral-reduced-arrow}
 R\oplus(\mathsf W_R,\lambda\overline d)[-1]
   \longrightarrow\mathcal B_R\oplus\mathsf W_R^0[-1],
 \qquad (a,w)\longmapsto(a\operatorname{id},\pi_0w),
\end{equation}
where $\mathsf W_R=\mathsf W\otimes_kR$ and $\pi_0$ is the
brutal degree-zero quotient.

\begin{theorem}\label{thm:vol2-integral-fibre}
There is a constructed equivalence of complexes
\begin{equation}\label{eq:vol2-integral-fibre-model}
 \mathcal F_\lambda\simeq
 (\mathcal B_R/R\operatorname{id})[-1]\oplus
       (\mathsf W_R^{\geq1},\lambda\overline d)[-1],
\end{equation}
where $\mathsf W^{\geq1}$ is the brutal subcomplex. Its cohomology is
\begin{align}\label{eq:vol2-integral-fibre-cohomology}
 H^1\mathcal F_\lambda&=\mathcal B_R/R\operatorname{id},\notag\\
 H^2\mathcal F_\lambda&=\mathsf B^1\otimes_kR,\notag\\
 H^n\mathcal F_\lambda&=\mathsf B^{n-1}\otimes_kR/(\lambda)
                              \quad(n\geq3),
\end{align}
and all other groups vanish.
\end{theorem}
\begin{proof}
In \eqref{eq:vol2-integral-reduced-arrow}, the scalar inclusion
splits by vacuum constant evaluation. The second map is
degreewise onto, so its fibre is its kernel
$\mathsf W_R^{\geq1}[-1]$. Explicitly, the redundant two terms
in its cone are the graph
$w\mapsto(-\lambda\overline dw,w)$.
Projection $(v,u)\mapsto v+\lambda\overline du$ and homotopy
$(v,u)\mapsto u$ cancel that graph without dividing by $\lambda$.
This proves the equivalence. In degree two the incoming
$\lambda\mathsf B^1$ disappeared with the cancelled term,
leaving the free module $\mathsf B^1_R$. In higher degrees the
boundaries remain $\lambda\mathsf B^{n-1}$, giving the formula.
\end{proof}

The exact sequence in degree two is therefore
\begin{equation}\label{eq:vol2-integral-path-extension}
 0\longrightarrow(\mathsf W^0/(k\cdot[J]))\otimes_kR
 \xrightarrow{\lambda\overline d}\mathsf B^1\otimes_kR
 \longrightarrow\mathsf B^1\otimes_kR/(\lambda)
 \longrightarrow0.
\end{equation}
Under the isomorphism
$\overline d:\mathsf W^0/(k\cdot[J])\to\mathsf B^1$, the
first map is multiplication by $\lambda$. A concrete instance
uses $F=J^2/2$, for which $d_0F=x_1\xi_1$. The fibre differential
gives
\[
 d(-dy\,F,0)=(\lambda dy\,x_1\xi_1,-dy\,L_F).
\]
Hence
\begin{equation}\label{eq:vol2-integral-path-witness}
 \boxed{\lambda[(dy\,x_1\xi_1,0)]
                 =[(0,dy\,L_{J^2/2})].}
\end{equation}
After its path is forgotten, the current class on the left is
$\lambda$-torsion. In the fibre its lift is a nonzero free class.
Indeed, $J^2$ is a translation primitive independent of $J$,
and the kernel of $\overline d$ in degree zero is only $k\cdot[J]$.

The integral charge-lifting test follows from the same arrow:
\begin{equation}\label{eq:vol2-integral-lift-test}
 [dy\,A]\text{ lifts through }\rho
 \quad\Longleftrightarrow\quad
 \operatorname{red}_+(A1)\in RJ.
\end{equation}
If the normal form is $a(\lambda)J$, the lift is
$a\Gamma_\lambda$ and the comparison homotopy is
$H_N(A-aL_J)$. No inverse level is used. Once a base lift for
a horizontal source shifted by $s$ has been chosen, its lift
space has homotopy groups $H^{s-j}\mathcal F_\lambda$ in
degrees $j\geq0$. Without a base lift the fibre may be empty.

\subsection{The singular level}

The terms of $\mathsf W_R$ are free over $R$.
The quotient $\mathcal B_R/R\operatorname{id}$ is a direct
summand of the torsion-free module $\mathcal B_R$ and hence is
flat over the PID $R$. Derived specialization is therefore
computed by reducing the displayed fibre complex. At $\lambda=0$
its positive differential vanishes, giving
\begin{equation}\label{eq:vol2-zero-level-fibre}
 H^{q+1}(\mathcal F_\lambda\otimes_R^Lk)=\mathsf W^q
                              \qquad(q\geq1).
\end{equation}
In degree two the new classes occur in
\[
 0\longrightarrow\mathsf B^1\longrightarrow\mathsf W^1
      \xrightarrow{\overline d}\mathsf B^2\longrightarrow0.
\]
They are the Tor contribution of integral $H^3$.

For example, $[\xi_0\xi_1]\ne0$ in $\mathsf W^2$.
A polynomial preimage under $\nabla$ would have root degree
zero, by injectivity of $T$ on $\mathcal Q_+$. It would then
give a two-ghost translation preimage in energy zero, where
there is no such monomial. Since
$d_0(x_1\xi_0)=-\xi_0\xi_1$, this is a nonzero element of
$\mathsf B^2$. Thus $dy\,x_1\xi_0$ becomes a degree-two
cocycle at zero level without being an integral cocycle.
The same torsion prevents an $R$-linear deformation retraction
of this fibre complex onto its cohomology: its nonzero
$\lambda$-torsion could not have a cochain section into the
torsion-free terms. This excludes that linear retraction,
not every possible derived replacement.

\subsection{The strict mode-preserving gauge groupoid}

Complete only a new deformation parameter $g$. On
$\mathcal V_R[[g]]$ fix
\[
 J_{-n}=L_{x_n},\qquad J_n=\lambda n\partial_{x_n}\quad(n>0),
 \qquad J_0=q(g)\operatorname{id},\quad q(g)\in gR[[g]].
\]
Seek connections $\nabla+A(g)$ with
$A\in g\mathcal N_R[[g]]$ satisfying
$[\nabla+A,J_n]=-nJ_{n-1}$ for every integer $n$.
A mode-preserving gauge is an element of
$1+g\mathcal N_R[[g]]$ commuting with every nonzero mode.

\begin{theorem}\label{thm:vol2-integral-gauge}
Such a connection exists exactly when $q(g)=\lambda a(g)$
for $a(g)\in gR[[g]]$. Every one is
\begin{equation}\label{eq:vol2-integral-gauge-normal-form}
 \partial_y+T+a(g)L_J+c(y,g)\operatorname{id},
                 \qquad c(y,g)\in gR[y][[g]].
\end{equation}
It is mode-preserving gauge equivalent to $\partial_y+T+a(g)L_J$.
The stabilizer of that normal form is $1+gR[[g]]$.
Thus the strict formal groupoid is the disjoint union, indexed
by $a(g)$, of the classifying groupoids of these scalar units.
\end{theorem}
\begin{proof}
The negative-mode equations force $A$ to commute with all
multiplications by $x_n$. These generate the state module
from $1$, so $A=L_f$. The equations for $n\geq2$ give
$\lambda\partial_{x_n}f=0$. Torsion-freeness makes $f$
independent of those variables. The $n=1$ equation is
$\lambda\partial_{x_1}f=q$. Polynomial coefficient comparison
forces $q=\lambda a$ and $f=aJ+c$. Conversely this formula
satisfies every mode relation.

A gauge commuting with the negative modes is multiplication
by a formal unit. Commutation with the positive modes and
torsion-freeness make it independent of the states. The unit
\[
 U=\exp\left(-\int_0^yc(t,g)\,dt\right)
\]
has polynomial coefficients in every $g$-degree and satisfies
$(\nabla+aL_J+c)U=U(\nabla+aL_J)$.
Between two normal forms a state-independent gauge cannot
change $a$, and its derivative must be zero. These are exactly
the stated scalar stabilizers.
\end{proof}

This groupoid fixes all nonzero modes. It is not the full
derived deformation groupoid, nor the quotient by gauges that
change those modes. At zero level the positive modes vanish;
then arbitrary multiplication corrections are permitted when
$q=0$. Most are not specializations of
\eqref{eq:vol2-integral-gauge-normal-form}.

Root completion also changes this problem. The series
$g_0=\sum_{n\geq0}(-1)^ny^{n+1}T^nJ/(n+1)!$ satisfies
$\nabla g_0=J$, so the integral charge image $dy\,L_J$ becomes
exact. In the completed $J$-tower the section
$1\mapsto e^{-yT}J=J-Tg_0$ splits the quotient by state
translation. This is the completion of the polynomial
extension $0\to P\xrightarrow{R_\partial}P\to O\to0$.
The polynomial projective model need not remain projective in
the continuous category. Its nonzero Massey coset and its
integral path relation therefore cannot be transported by
an undeclared equivalence across this change of coefficients.

\section{Finite categories of interfaces}
\label{sec:vol2-finite-interface-centres}

An invertible bimodule composes by derived tensor product, as in
\eqref{eq:vol2-interface-composition}. A collection of interfaces
with specified junction maps has more data: the centre must act
compatibly with every composition of those maps. For two objects,
Proposition~\ref{prop:vol2-dg-triangular} expresses this as the fibre
of the difference between the two endpoint characteristic actions
on the joining bimodule. With more objects, the longer composable
strings contribute further terms.

Let $R$ be a commutative $k$-algebra and let $\mathcal A$ be a dg
category with ordered objects $0,\ldots,n-1$, scalar endomorphisms
$\mathcal A(i,i)=R$, and $\mathcal A(i,j)=0$ for $i>j$.
Suppose each forward complex $M_{ij}$ is bounded and finite
projective over $R$, and all dg compositions are specified and
associative. Let $A$ be the sum algebra, $S=\bigoplus_iRe_i$,
and $\overline A=\ker(A\to S)$. The normalized relative bimodule
bar is
\begin{equation}\label{eq:vol2-directed-bar}
 P_\ell=A\otimes_S(\overline A[1])^{\otimes_S\ell}\otimes_SA,
                     \qquad0\leq\ell\leq n-1.
\end{equation}
For an unsuspended tensor of internal degree $q$ and bar length
$\ell$, the total degree is $q-\ell$. One sign convention is
$d_{\mathrm{int}}+(-1)^q\sum_{j=0}^{\ell}(-1)^jd_j$, where
$d_j$ multiplies adjacent factors, including the two endpoint
actions. Suspension translates this to
\eqref{eq:vol2-directed-bar}.

\begin{theorem}\label{thm:vol2-directed-bar}
The total complex $P$ is a finite homotopically projective
$A$-bimodule resolution of $A$. Thus the relative derived centre is
\[
 Z_R(A)=R\operatorname{Hom}_{A\otimes_RA^{\mathrm{op}}}(A,A)
          \simeq\operatorname{Hom}_{A^e}(P,A).
\]
Its evaluation gives, for every complex $F$ of $R$-modules,
\begin{equation}\label{eq:vol2-directed-evaluation}
 R\operatorname{Hom}_R(F,Z_R(A))
 \simeq R\operatorname{Hom}_{A^e}(F\otimes_R^LA,A).
\end{equation}
For every derived commutative scalar extension $R\to R'$,
\begin{equation}\label{eq:vol2-directed-base-change}
 Z_R(A)\otimes_R^LR'
       \simeq Z_{R'}(A\otimes_R^LR'),
\end{equation}
compatibly with evaluation, cup products and Hochschild insertions.
\end{theorem}
\begin{proof}
A nonzero string of $\ell$ forward arrows visits $\ell+1$
strictly increasing objects, so $\ell\leq n-1$.
The unit-insertion contraction of the augmented unnormalized
bar proves its exactness after forgetting an outer action.
The degenerate subcomplex splits off, leaving the normalized
resolution.

Each remaining term is induced from a tensor of bounded finite
projective coefficient complexes and two representable
idempotent modules. Hom into an acyclic bimodule reduces to
Hom from that coefficient complex into an acyclic idempotent
component, hence is acyclic. The finite totalization is therefore
homotopically projective. Tensor--Hom adjunction applied to
$F\otimes_RP$ proves \eqref{eq:vol2-directed-evaluation}.

Every cochain coefficient is a Hom between bounded finite
projective $R$-complexes. These Homs and the finite sums commute
with derived scalar extension. The bar differential, cuts and
insertions use the same composition maps after extension.
This proves \eqref{eq:vol2-directed-base-change} with its operations.
\end{proof}

Strictly increasing object indices made the bar finite.
With arbitrary dg endpoint algebras the repeated corner inputs
of Proposition~\ref{prop:vol2-dg-triangular} instead require
product totalization. We now use the finite case to see what
the operations and the parameter fibres retain.

\section{Koszul parameters and contraction brackets}
\label{sec:vol2-koszul-interface}

Let $r\geq1$, $R=k[h_1,\ldots,h_r]$, $I=(h_1,\ldots,h_r)$,
and $S_0=R/I$. Let $U_0=k^r$ have degree-one basis
$u_1,\ldots,u_r$, and write $L_i$ for exterior multiplication
by $u_i$. Take the arrow complex
\begin{equation}\label{eq:vol2-koszul-arrow}
 \begin{gathered}
 M=R\otimes_k\Lambda U_0,\qquad d_M=\sum_ih_iL_i,\\
 T(M)=Re_0\oplus Re_1\oplus M,\qquad e_1Me_0=M,\quad M^2=0.
 \end{gathered}
\end{equation}
The differential is a differential of the arrow module; it is
not a derivation of the auxiliary exterior multiplication.
Anticommutation gives $d_M^2=0$, and the Koszul complex has
cohomology $S_0$ in degree $r$.

Put $E=\operatorname{End}_R(M)$ with
$\delta f=d_Mf-(-1)^{|f|}fd_M$. Reverse the one-arrow coordinate
in Theorem~\ref{thm:vol2-triangular-cone}. The centre then has
the finite model
\begin{equation}\label{eq:vol2-koszul-centre}
 Z^n=(R^2)^n\oplus E^{n-1},\qquad
 d(a,b;f)=(0,0;(b-a)\operatorname{id}-\delta f).
\end{equation}
The endpoint terms occur only in degree zero. Use the
source-first opposite-cup convention
\begin{equation}\label{eq:vol2-koszul-cup}
 (a,b;f)(a',b';f')=(aa',bb';fb'+af').
\end{equation}
In this convention, for $\sigma f=(0,0;f)$ of degree $|f|+1$,
the insertion bracket is
\begin{equation}\label{eq:vol2-koszul-bracket}
 [\sigma f,\sigma g]
       =\sigma\bigl(fg-(-1)^{|f||g|}gf\bigr).
\end{equation}
Endpoint scalars have zero bracket with the mixed terms. The
bracket has degree $-1$.

These operations follow directly from the one-arrow bar. Its
two cuts give \eqref{eq:vol2-koszul-cup}, while insertion into
the only arrow input is composition in $E$. Two mixed cup
factors would require two composable forward arrows, which do
not exist. The differential identity on endpoints is
$bb'-aa'=(b-a)b'+a(b'-a')$; the other identities reduce to
$R$-linearity of $\delta$. On $Z[1]$ the mixed differential
is $\delta$, so the bracket differential and Jacobi identities
are the graded-commutator identities. This fixes the signs
of the comparison as well as its underlying complex.

\subsection{The integral ring and its bracket}

Let $\iota_i$ contract by the dual of $u_i$. The relations are
\begin{equation}\label{eq:vol2-koszul-clifford}
 \{L_i,L_j\}=0,\qquad \{\iota_i,\iota_j\}=0,
 \qquad \{\iota_i,L_j\}=\delta_{ij}\operatorname{id}.
\end{equation}
The operators $L_A\iota_B$, with increasing subsets $A,B$,
form a basis of $E$: for one generator $1,L,\iota,L\iota$
span all matrices, and graded tensor products give the general
case. Moreover $\delta L_i=0$ and
$\delta\iota_i=h_i\operatorname{id}$. At fixed creation
monomial $L_A$, the contraction variables are the Koszul
resolution of $(h_i)$. Therefore
\[
 H^q(E)=S_0\otimes_k\Lambda^qU_0\quad(0\leq q\leq r),
 \qquad H^q(E)=0\text{ otherwise}.
\]

\begin{theorem}\label{thm:vol2-koszul-centre-ring}
The relative centre of \eqref{eq:vol2-koszul-arrow} has
\begin{align*}
 \operatorname{HH}_R^0(T)&=R\times_{S_0}R,&
 \operatorname{HH}_R^1(T)&=0,\\
 \operatorname{HH}_R^{q+1}(T)&=S_0\otimes_k\Lambda^qU_0
                         &&(1\leq q\leq r),
\end{align*}
and no other cohomology. Its positive-degree cup ideal is
square-zero; a pair in degree zero acts through its common
residue in $S_0$. For $t_i=[(0,h_i;\iota_i)]$,
\begin{equation}\label{eq:vol2-koszul-centre-presentation}
 \operatorname{HH}_R^0(T)
 =R[t_1,\ldots,t_r]/(h_it_j-h_jt_i,\ t_it_j-h_it_j)_{i,j}.
\end{equation}
\end{theorem}
\begin{proof}
The map $R^2\to H^0(E)=S_0$ is $(a,b)\mapsto b-a$.
It is onto with kernel $R\times_{S_0}R$, so the fibre exact
sequence gives the groups. The mixed cup product is zero;
the endpoint residue gives its scalar action.

The representative of $t_i$ is closed because
$\delta\iota_i=h_i$. It maps to $(0,h_i)$ in the fibre product.
Every pair is a diagonal scalar plus $(0,c)$ with $c\in I$,
so these classes generate. The quadratic relations reduce
each polynomial to $a+\sum_ib_it_i$. It maps to zero exactly
when $a=0$ and $\sum_ib_ih_i=0$. The first Koszul syzygies
$h_it_j-h_jt_i$ generate this kernel, proving the presentation.
\end{proof}

For $A=\{a_1<\cdots<a_q\}$ set
$\alpha_A=[\sigma L_{a_1}\cdots L_{a_q}]$, of degree $q+1$.
Then for $q\geq2$,
\begin{equation}\label{eq:vol2-koszul-contraction-brackets}
 [t_i,\alpha_A]=
 \begin{cases}
 (-1)^{j-1}\alpha_{A\setminus\{i\}},&i=a_j,\\
 0,&i\notin A.
 \end{cases}
\end{equation}
Indeed, the graded commutator of $\iota_i$ with the creation
product is precisely this contraction. When $q=1$ the result
is the scalar identity, which is the boundary of an endpoint
difference in \eqref{eq:vol2-koszul-centre}, so the bracket
is zero. All brackets between two positive classes or two
degree-zero classes vanish; the generators supercommute in
their respective creation or contraction factors.

For $r=2$, with $\omega=\alpha_{\{1,2\}}$, this gives
\[
 [t_1,\omega]=\alpha_{\{2\}},\qquad
 [t_2,\omega]=-\alpha_{\{1\}}.
\]
The same $t_i$ annihilate every positive class under cup
product. Thus the bracket detects a composition invisible
to that product. The proof of the integral calculation also
works for any regular sequence $(h_i)$ in a commutative
$\mathbb Q$-algebra: regularity is used in the Koszul
cohomology and in its first syzygies.

\subsection{The full singular fibre}

The finite bar commutes with derived scalar extension by
Theorem~\ref{thm:vol2-directed-bar}. If some $h_i$ is invertible,
$h_i^{-1}\iota_i$ contracts the arrow $M$ and the centre becomes
the two scalar endpoint sectors. At the common zero, put
$V_0=\Lambda U_0$. The arrow differential vanishes, and the
only centre differential sends the endpoint difference to
$\operatorname{id}\in E^0$, in centre degree one. Hence
\begin{equation}\label{eq:vol2-koszul-singular-fibre}
 \operatorname{HH}_k^*(T\otimes_Rk)
   =k\ltimes(\operatorname{End}_k(V_0)/k\operatorname{id})[-1].
\end{equation}
Its displayed ideal has square-zero cup product and shifted
graded-commutator bracket. The unit has zero bracket.
Since
\[
 \dim_k\operatorname{End}^p(\Lambda U_0)
    =\sum_j\binom rj\binom r{j+p}=\binom{2r}{r+p},
\]
the dimensions are
\begin{equation}\label{eq:vol2-koszul-fibre-dimensions}
 \dim_k\operatorname{HH}^n_k(T\otimes_Rk)=
 \begin{cases}
 1+\binom{2r}{r-1},&n=0,\\
 \binom{2r}{r}-1,&n=1,\\
 \binom{2r}{r+n-1},&n\ne0,1.
 \end{cases}
\end{equation}
Out-of-range binomial coefficients are zero. The total dimension
is $4^r$. For $r=2$, degrees $-1,0,1,2,3$ have dimensions
$1,5,5,4,1$. In contrast, ordinary specialization of integral
degree-zero cohomology has basis $1,t_1,t_2$ and dimension three.
Specializing the complex retains the missing extension data.

There is no $R$-linear deformation retraction of
\eqref{eq:vol2-koszul-centre} onto its integral cohomology.
The nonzero $\operatorname{HH}_R^2(T)$ is annihilated by $h_1$,
whereas multiplication by $h_1$ is injective on each free
cochain term. A cochain section of this torsion group must be
zero and cannot have a left inverse. This argument concerns
a linear retraction over $R$; it is not a classification of
all formality zigzags.

\subsection{A second-order obstruction}

At $r=2$ and zero parameters, use the basis
$v_0=1,v_1=u_1,v_2=u_2,v_3=u_1u_2$ of $V_0$. Define the
degree-one endomorphism $f$ by
\[
 f(v_0)=v_1,\qquad f(v_1)=v_3,\qquad f(v_2)=f(v_3)=0.
\]
Over $k[\epsilon]/(\epsilon^2)$, the arrow differential
$\epsilon f$ defines a first-order deformation of $T_0$.
But $f^2(v_0)=v_3$, and
\begin{equation}\label{eq:vol2-koszul-obstruction}
 \tfrac12[\sigma f,\sigma f]=\sigma f^2\ne0
                         \quad\text{in }\operatorname{HH}^3_k(T_0).
\end{equation}
It is nonzero in \eqref{eq:vol2-koszul-singular-fibre} because
an internal degree-two map is not a scalar identity.
Thus the second-order Hochschild equation
$d_{\mathrm{Hoch}}g+\tfrac12[\sigma f,\sigma f]=0$ has no
solution. This excludes even curvature-free $A_\infty$
corrections on the same graded module with the same
first-order class, although the cup square of $\sigma f$
is zero. By contrast, the original Koszul directions
$xL_1+yL_2$ square to zero and integrate over $k[x,y]$.

\section{Four composable objects}
\label{sec:vol2-four-interfaces}

The preceding category had one arrow direction and no composable
pair of nonidentity arrows. Four ordered objects provide the
first three-arrow string. Its two composition orders give an
associativity equation, and the centre retains that equation
in its own differential.

Use chronological products in this section: $a_{ij}a_{jk}$ means
first $i\to j$, then $j\to k$, so $e_ia_{ij}=a_{ij}=a_{ij}e_j$.
Take objects $0<1<2<3$, scalar identities, and a degree-zero
basis arrow $a_{ij}$ for each $i<j$. Set
\begin{equation}\label{eq:vol2-four-products}
 a_{01}a_{12}=pa_{02},\quad a_{12}a_{23}=qa_{13},\quad
 a_{02}a_{23}=ra_{03},\quad a_{01}a_{13}=sa_{03}.
\end{equation}
All other products are forced by identities or are zero.
The only associator of nonidentity arrows is
\[
 (a_{01}a_{12})a_{23}-a_{01}(a_{12}a_{23})
       =Fa_{03},\qquad F=pr-qs.
\]
Thus the associative family is the rank-ten algebra $A_B$ over
$B=k[p,q,r,s]/(pr-qs)$.

\subsection{The centre differential and its operations}

Order the edges as $01,02,03,12,13,23$ and the triangles as
$012,123,023,013$. A normalized cochain assigns a scalar to
each increasing string, multiplying its endpoint arrow.
The finite bar gives
\begin{equation}\label{eq:vol2-four-centre}
 Z_B(A_B)\simeq
 [B^4\xrightarrow{d_0}B^6\xrightarrow{d_1}B^4\xrightarrow{d_2}B]
\end{equation}
in degrees zero through three, with
\begin{align}\label{eq:vol2-four-differentials}
 (d_0\lambda)_{ij}&=\lambda_j-\lambda_i,\notag\\
 d_1f&=\bigl(p(f_{01}+f_{12}-f_{02}),
             q(f_{12}+f_{23}-f_{13}),\notag\\
 &\hspace{20mm}r(f_{02}+f_{23}-f_{03}),
             s(f_{01}+f_{13}-f_{03})\bigr),\notag\\
 d_2(\alpha,\beta,\gamma,\delta)&=-r\alpha+s\beta-p\gamma+q\delta.
\end{align}
Indeed, in arity zero the two Hochschild endpoint terms give
$d_0$. On a triangle the three faces have the same composition
scalar, giving $d_1$. On $0123$ the four terms give $d_2$.
Before imposing the relation, direct substitution gives
\begin{equation}\label{eq:vol2-four-square}
 d_1d_0=0,\qquad
 d_2d_1f=-F(f_{01}+f_{12}+f_{23}-f_{03}).
\end{equation}
The square-zero condition remembers the original associator.

Write $m_{012}=p,m_{123}=q,m_{023}=r,m_{013}=s$.
For positive arities $n,m$, the cup product is
\begin{equation}\label{eq:vol2-four-cup}
 (f\smile g)_{i_0\cdots i_{n+m}}
   =m_{i_0i_ni_{n+m}}f_{i_0\cdots i_n}g_{i_n\cdots i_{n+m}}.
\end{equation}
An arity-zero factor acts by the appropriate first or last
endpoint scalar. A product with no increasing string is zero.
The bar cut proves this formula; two successive cuts prove
associativity using $pr=qs$, and the faces give the Leibniz
identity. Its unit is $(1,1,1,1)$.

For insertion, use
$f\circ g=\sum_{j=0}^{n-1}(-1)^{j(m-1)}f\circ_jg$ and
$[f,g]=f\circ g-(-1)^{(n-1)(m-1)}g\circ f$.
For unary $f$ and binary $g$ this gives
\[
 [f,g]_{ijk}=(f_{ik}-f_{ij}-f_{jk})g_{ijk}.
\]
For binary cochains it gives
\begin{equation}\label{eq:vol2-four-bracket}
 [f,g]_{0123}=f_{023}g_{012}-f_{013}g_{123}
                   +g_{023}f_{012}-g_{013}f_{123}.
\end{equation}
These are the two insertion slots on the unique three-arrow
string. With the Hochschild signs of
\eqref{eq:vol2-four-differentials}, a positive arity-$n$
cochain has $df=(-1)^{n-1}[\mu,f]$, where $\mu$ is the binary
product. Rephasing arities gives the dg Lie convention
$[\mu,-]$. In particular $[\mu,\mu]/2=Fa_{03}$.

\subsection{The parameter fibres and the integral module}

Let $B\to K$ be a map to a field, and let $c$ count the nonzero
parameters among $p,q,r,s$. Theorem~\ref{thm:vol2-directed-bar}
makes the specialized complex the full centre over $K$.
Its cohomology dimensions are
\begin{equation}\label{eq:vol2-four-fibres}
 \begin{array}{c|rrrr}
 c&\operatorname{HH}^0&\operatorname{HH}^1&
          \operatorname{HH}^2&\operatorname{HH}^3\\\hline
 0&1&3&4&1\\
 1&1&2&2&0\\
 2&1&1&1&0\\
 4&1&0&0&0
 \end{array}
\end{equation}
The proof is a rank calculation. The incidence matrix $d_0$
has rank three. The four unweighted triangle rows of $d_1$
have one relation with all four coefficients nonzero, and
any three are independent. Thus $\operatorname{rank}d_1$
is $\min(c,3)$. The last row has rank one off the origin
and zero at the origin. Subtracting adjacent ranks gives
the table. The equation $pr=qs$ excludes $c=3$ and excludes
the supports $\{p,r\}$ and $\{q,s\}$ when $c=2$.

Integrally $B$ is a domain: $pr-qs$ is primitive and
irreducible as a polynomial in $r$ over $k[p,q,s]$.
The answer retains a module supported at the origin:
\begin{align}\label{eq:vol2-four-integral}
 \operatorname{HH}_B^0(A_B)&=B,\qquad
 \operatorname{HH}_B^1(A_B)=0,\notag\\
 \operatorname{HH}_B^3(A_B)&=B/(p,q,r,s),\notag\\
 \operatorname{HH}_B^2(A_B)&=
   \ker(-r,s,-p,q)/\operatorname{im}L.
\end{align}
Here $L(x,y,z)=(-px,-qy,r(x-z),s(y-z))$.
For the degree-one assertion, nonzero parameters in the domain
force all four triangle additivity equations. Setting
$\lambda_0=0$ and $\lambda_i=f_{0i}$ gives
$f_{ij}=\lambda_j-\lambda_i$. More generally, modulo a
vertex coboundary an edge cochain has a unique representative
with $f_{01}=f_{12}=f_{23}=0$. Its other coordinates
$(f_{02},f_{13},f_{03})=(x,y,z)$ give the displayed image
$L$. The cokernel of $d_2$ gives degree three.

There is no $B$-linear deformation retraction onto these
cohomology groups. Its nonzero degree-three group is killed
by $p$, and therefore cannot have a section into the free
degree-three cochain term. Scalar extension must act on
\eqref{eq:vol2-four-centre}, not merely on
\eqref{eq:vol2-four-integral}.

At the origin, let the binary cochain $f$ have
$f_{012}=f_{023}=1$ and $f_{123}=f_{013}=0$.
The first-order product $\epsilon f$ is associative modulo
$\epsilon^2$, but
\[
 \tfrac12[f,f]_{0123}=1\in\operatorname{HH}^3=k.
\]
The positive-arity differential is zero there, so no
second-order correction removes this obstruction. Equivalently,
$p=r=\epsilon$, $q=s=0$ gives $F=\epsilon^2$.
A ternary operation cannot help on the unchanged degree-zero
arrow module over an ordinary ring: it would have degree $-1$
and no target. This is the deformation problem with the
objects, arrow bases and units fixed.

\subsection{A primitive and a finite coherent resolution}

Introduce a coefficient of degree $-1$ with the required
differential:
\begin{equation}\label{eq:vol2-four-coherent-base}
 R_{\mathrm{coh}}=k[p,q,r,s]\otimes\Lambda(\chi),
                  \qquad d\chi=pr-qs.
\end{equation}
Its cohomology is $B$ in degree zero, since $F$ is a
non-zero-divisor. A map to a derived test algebra nevertheless
retains the image of $\chi$ as a specified primitive.
On the rank-ten arrow module, the products
\eqref{eq:vol2-four-products} and
\begin{equation}\label{eq:vol2-four-ternary}
 m_3(a_{01},a_{12},a_{23})=-\chi a_{03}
\end{equation}
define a strictly unital $A_\infty$ category, with convention
$dm_3=-\operatorname{Assoc}(m_2)$ on closed degree-zero
arrows. There is no increasing string of four nonidentity
inputs, so all higher normalized identities vanish.

The binary homotopies and their compatibility can also be
retained in a strict dg algebra. Let $\mathcal Q$ be the
free directed path algebra over $R_{\mathrm{coh}}$ on
\[
 X_{ij}:i\to j\ (|X_{ij}|=0),\qquad
 H_{ijk}:i\to k\ (|H_{ijk}|=-1),\qquad
 K_{03}:0\to3\ (|K_{03}|=-2),
\]
for increasing indices. Extend the coefficient differential
by the graded Leibniz rule and
\begin{align}\label{eq:vol2-four-resolution}
 dX_{ij}&=0,\notag\\
 dH_{ijk}&=X_{ij}X_{jk}-m_{ijk}X_{ik},\notag\\
 dK_{03}&=H_{012}X_{23}-X_{01}H_{123}
                 +pH_{023}-qH_{013}+\chi X_{03}.
\end{align}
Odd coefficients are moved to the left with their Koszul sign.
In $d^2K_{03}$ the first four terms cancel all paths except
$-FX_{03}$; the last gives $FX_{03}$. Thus $d^2=0$.
The free basis has fifteen degree-zero paths (four identities,
six arrows, four length-two paths and one length-three path),
six degree-$-1$ paths (the four $H$ and their two composable
extensions), and the single $K_{03}$. Its rank is $22$.

Let $V_{\mathrm{ar}}$ be the submodule spanned by the identities
and the six $X_{ij}$, with its coefficient differential.
Let $i$ include it, and let $\pi$ be the identity there and
zero on the $H_{ijk},K_{03}$. Specify its remaining values by
\begin{align}\label{eq:vol2-four-projection}
 \pi(X_{ij}X_{jk})&=m_{ijk}X_{ik},\notag\\
 \pi(X_{01}X_{12}X_{23})&=prX_{03},&
 \pi(X_{01}H_{123})&=\chi X_{03},&
 \pi(H_{012}X_{23})&=0.
\end{align}
The degree-$-1$ homotopy is
\begin{align}\label{eq:vol2-four-homotopy}
 h(X_{ij}X_{jk})&=H_{ijk},\notag\\
 h(X_{01}X_{12}X_{23})&=H_{012}X_{23}+pH_{023},&
 h(X_{01}H_{123})&=-K_{03},
\end{align}
and is zero on all other basis paths. Extend $\pi$ linearly
and $h$ graded-linearly, so $h(\chi a)=-\chi h(a)$.

\begin{theorem}\label{thm:vol2-four-contraction}
Over all of $R_{\mathrm{coh}}$ these maps satisfy
\[
 \pi i=1,\qquad dh+hd=1-i\pi,\qquad
 h^2=hi=\pi h=0,\qquad d\pi=\pi d.
\]
The transferred product on $V_{\mathrm{ar}}$ is
\eqref{eq:vol2-four-products}, with the ternary value
\eqref{eq:vol2-four-ternary}; every normalized higher
operation vanishes. No parameter is inverted.
\end{theorem}
\begin{proof}
The intervals $0\to2$ and $1\to3$ contain the pairs
$H_{012}\mapsto X_{01}X_{12}-pX_{02}$ and
$H_{123}\mapsto X_{12}X_{23}-qX_{13}$.
On $0\to3$ put
\begin{align*}
 v_1&=X_{02}X_{23}-rX_{03},& h_1&=H_{023},\\
 v_2&=X_{01}X_{13}-sX_{03},& h_2&=H_{013},\\
 v_3&=X_{01}X_{12}X_{23}-pX_{02}X_{23},& h_3&=H_{012}X_{23},\\
 \ell&=H_{012}X_{23}-X_{01}H_{123}
                   +pH_{023}-qH_{013}+\chi X_{03}.
\end{align*}
Together with $K_{03},X_{03}$ these are another free basis,
obtained by a triangular change with constant invertible
pivots. Its differential has $dh_j=v_j$, $dK_{03}=\ell$,
$d\ell=0$, leaving $X_{03}$. Contracting these four pairs
and returning to the original basis gives
\eqref{eq:vol2-four-projection}--\eqref{eq:vol2-four-homotopy}.
In particular $\pi(\ell)=0$ forces the nonzero value
$\pi(X_{01}H_{123})=\chi X_{03}$. Graded linearity of $h$
cancels the coefficient-differential terms.

The transferred binary operation is $\pi\mu(i\otimes i)$.
Its two ternary trees give
\[
 \pi\bigl(h(X_{01}X_{12})X_{23}
                   -X_{01}h(X_{12}X_{23})\bigr)=-\chi X_{03}.
\]
There are no longer composable nonidentity strings, and the
side conditions kill higher trees with an identity input.
This proves the strictly unital structure.
\end{proof}

After an ordinary specialization $F=\chi=0$, the augmentation
$X_{ij}\mapsto a_{ij}$, $H_{ijk},K_{03}\mapsto0$ is a strict
dg algebra quasi-isomorphism, by this contraction. Over
$R_{\mathrm{coh}}$ the projection is a cochain map to the
$A_\infty$ category, not a strict algebra map.
The primitive has not removed the ordinary obstruction:
a dg map to $k[\epsilon]/(\epsilon^3)$ with $p=r=\epsilon$
and $q=s=0$ would need a degree-$-1$ primitive of
$\epsilon^2$, which that ring does not contain.

The finite categories in this chapter specify algebraic
compositions of interfaces. Their object labels alone supply
no field-theoretic realization. Such a realization must
identify the arrows, junctions, collision maps and local
vacuum operations. Likewise, a native centre comparison must
carry the full current hierarchy, the polynomial Massey
coset and the integral path relation, together with derived
specialization. These are tests on an actual evaluation map;
they do not construct a common recipient for arbitrary
geometric boundary actions.

\chapter{Action profiles and universal recipients}
\label{chap:vol2-recipients}

A universal boundary action must receive actions together with their
comparison homotopies. A multiplication on the underlying coefficient
algebra does not suffice. For an action whose first nonzero term is
quadratic in its parameter, doubling that parameter multiplies the term
by four, whereas composing two copies adds their first terms and
multiplies it by two. The standard additive multiplication is therefore
incompatible with such an action whenever its first class is nonzero.
Nevertheless, a nonlinear change of actor can place the action inside
a multiplicative one.

The higher terms of these actions have a simple scalar example. Over a
characteristic-zero field, the formal solution of
\[
 p=q+\frac{z^2}{2}p^2,\qquad p\equiv q\pmod z,
\]
is
\[
 p=\sum_{n\geq0}\frac{1}{n+1}\binom{2n}{n}
                 \left(\frac{z^2}{2}\right)^n q^{n+1}.
\]
Writing $p=qC(z^2q/2)$ reduces the equation to $C=1+uC^2$,
whose solution with constant term one is
$(1-\sqrt{1-4u})/(2u)$. For noncommuting coefficients the same
implicit construction uses ordered trees. Its terms define actual
actions on the collision coefficients. The universal question then
concerns all maps between these actions, including their higher paths.

\begingroup
\newcommand{\vCE}{\operatorname{CE}_{\star}}
\newcommand{\vU}{U_{\star}}
\newcommand{\vB}{\operatorname{B}_{\star}}
\newcommand{\vSym}{\operatorname{Sym}_{\star}}
\newcommand{\vRH}{\operatorname{RHom}}
\newcommand{\vMap}{\operatorname{Map}}
\newcommand{\vot}{\mathbin{\otimes^!}}
\newcommand{\vM}{\mathscr M}
\newcommand{\vF}{\mathscr F}
\newcommand{\vAct}{\mathscr A_C}
\newcommand{\vMon}{\mathscr B_C}
\newcommand{\vEnd}{\mathcal E_C}
\newcommand{\vid}{\mathrm{id}}

\section{The coefficients of an action}
\label{sec:vol2-action-carrier}

Fix a field $k$ of characteristic zero and $X=\mathbb A^1_k$.
At a collision, the order of the coinciding labels is retained.
For a nonempty finite set $I$, take a copy $X^I_\sigma$ of $X^I$
for each total order $\sigma$ on $I$. Glue two copies by the identity
on the open
\[
 V_{\sigma\rho}=X^I\setminus
 \bigcup_{\substack{\{i,j\}\subset I\\
                  \sigma,\rho\text{ order }i,j\text{ oppositely}}}
                  \{z_i=z_j\}.
\]
The resulting scheme, denoted by $\mathscr R_I$, is the ordered
configuration scheme of \cite{AKY2024}*{\S4.1}. The gluing condition
holds because, whenever two orders disagree on a pair, at least one
disagrees with any third order on that pair. Thus
$V_{\sigma\rho}\cap V_{\sigma\eta}
 =V_{\sigma\rho}\cap V_{\rho\eta}$.
The scheme is smooth, with a finite affine-chart cover and
quasi-compact intersections. It can be nonseparated. Over a tuple with
equality blocks $B$, the map $\mathscr R_I\to X^I$ has
$\prod_B|B|!$ points, one for each order inside each block.

A surjection $\alpha:I\twoheadrightarrow J$ equipped with an order
on each fibre defines
\[
 \delta_\alpha:\mathscr R_J\longrightarrow\mathscr R_I,
 \qquad z_i=z_{\alpha(i)}.
\]
On charts it expands each label into its ordered fibre. Fibre orders
compose by concatenation. The chart maps glue, since the inverse image
of an overlap is the overlap of the corresponding orders on $J$.
Their compositions are the maps for the composite ordered surjections.
In particular, a nonconsecutive collision uses its actual map into
$\mathscr R_I$; it is not replaced by the small diagonal of one chart.

Let $\mathscr C$ be the enhanced category of compatible complexes of
right differential modules on these schemes:
\begin{equation}\label{eq:vol2-action-category}
 \mathscr C=D(k)\times
     \varprojlim_{I,\alpha}D\!\operatorname{Mod}(\mathscr R_I),
 \qquad \delta_\alpha^!F_I\simeq F_J.
\end{equation}
Differential modules on $\mathscr R_I$ are defined by descent on its
smooth chart cover. The displayed equivalences include all composition
homotopies. The first factor is the isolated empty coefficient
$F_\varnothing$. An object is \emph{positive} if this coefficient is
zero. Complexes are cohomological, suspension $s=[1]$ has degree $-1$,
and every permutation uses the Koszul sign. Extraordinary inverse
images and the diagonal tensor retain their canonical density lines
and shifts. Passing to left modules for a coordinate computation uses
the canonical-bundle equivalence.

For an ordered decomposition $I=A\sqcup B$, let $U_{A,B}$ be the open
where every point in $A$ is distinct from every point in $B$.
Collisions inside either part remain. Restriction of orders identifies
this open with the disjointness open in
$\mathscr R_A\times\mathscr R_B$. With $j_{A,B}$ its inclusion, set
\begin{equation}\label{eq:vol2-action-convolution}
 (F\star G)_I=
 \bigoplus_{I=A\sqcup B}
 j_{A,B,*}\bigl((F_A\boxtimes G_B)|_{U_{A,B}}\bigr).
\end{equation}
An empty part contributes scalar tensor product. The unit $K$ has
$K_\varnothing=k$ and zero nonempty coefficients. Iterating
\eqref{eq:vol2-action-convolution} sums over ordered partitions into
the required number of parts. All parenthesizations impose the same
cross-part inequalities. This proves associativity and its coherences;
permuting the parts gives symmetry. Under a collision, a summand
vanishes if one collision fibre meets two parts. Every surviving
partition descends uniquely to the quotient labels. Localization and
closed-support base change therefore give the collision coherences.
These operations preserve colimits separately.

The second tensor product uses coefficients at the same configuration:
\[
 (F\vot G)_I=\Delta_{\mathscr R_I}^!(F_I\boxtimes G_I).
\]
Its unit $J$ is the family of dualizing objects
$J_I=(\mathscr R_I\to\operatorname{Spec}k)^!k$, with empty
coefficient $k$. The objects $J$ and $K$ have different nonempty
coefficients. The forward interchange
\begin{equation}\label{eq:vol2-action-interchange}
 (F\star G)\vot(H\star L)\longrightarrow
       (F\vot H)\star(G\vot L)
\end{equation}
keeps summands with the same ordered partition and applies the
localization comparison there. In a rectangular iterated product,
either order of interchange keeps exactly the common partitions.
This proves its coherence and gives a diagonal tensor on
coassociative $\star$-coalgebras.

Write $\vM$ for connected coaugmented counital coherent
coassociative coalgebras in $\mathscr C$. Connected means that the
empty coefficient and its unit maps are $k$. A coalgebra $F$ is
\emph{factorizing} if every ordered partition
$P=(I_1|\cdots|I_r)$ into nonempty parts has an equivalence
\begin{equation}\label{eq:vol2-action-factorization}
 \delta_{P,F}:j_P^!F_I\xrightarrow{\sim}
       (F_{I_1}\boxtimes\cdots\boxtimes F_{I_r})|_{U_P}.
\end{equation}
The coproduct supplies this map and its refinement homotopies.
Let $\vF\subset\vM$ be the full subcategory of such coalgebras.
The diagonal tensor preserves it by
\eqref{eq:vol2-action-interchange}.

Fix a boundary $C\in\vF$. A \emph{pointed action} consists of
$F\in\vF$, a point $u:J\to F$, an action
$\alpha:F\vot C\to C$ in $\vM$, and a specified path
$\alpha(u\vot\vid_C)\simeq\vid_C$.
A morphism $(F,\alpha)\to(G,\beta)$ includes a pointed coalgebra
map $f:F\to G$, a path
$\alpha\simeq\beta(f\vot\vid_C)$, and its compatibility on $J$.
Their higher homotopies are retained. Denote this category by $\vAct$.
The functor it classifies is
\begin{equation}\label{eq:vol2-action-functor}
 \Phi_C(F,u)=\vMap^{\vM}_{C/}(F\vot C,C),
\end{equation}
the homotopy fibre of restriction over the identity of $C$.
Representability of $\Phi_C$ is equivalent to a final object of
$\vAct$.

The ordered product of two actions is their composition, with actor
$F\vot G$. Denote it by $a\circledast b$; the first action is
outermost. The unit action on $J$ is initial. A \emph{multiplicative
action} has an associative unital monoid actor for $\vot$ and a
coherent module structure on $C$. Its binary path compares acting
after multiplying the two actors with their composed action. All
associativity and unit homotopies are part of the structure.
Write $\vMon$ for the category of multiplicative actions, with monoid
maps and coherent action paths. A local centre, in this categorical
sense, is a final object of $\vMon$. A final pointed action has a
unique coherent multiplicative structure. The converse requires a
comparison with pointed actions whose actors have no specified
multiplication.

\section{Trees integrating a derivation}
\label{sec:vol2-action-profiles}

Let $\mathfrak g$ be a positive Lie object in $\mathscr C$ whose
coefficient modules lie in degree zero, and put $A=\vU\mathfrak g$.
Fix a degree-zero derivation $D$ of the full coefficient family.
Its extension to $A$ satisfies Leibniz and commutes with all
collision and normal-transfer maps. Products below are ordered products in $A$.
PBW symmetrization uses averaging in characteristic zero.

Give an odd symbol $e$ degree one and zero differential, and form
$\mathfrak h=\mathfrak g\ltimes e\mathfrak g$, with
\[
 [x,em]=e[x,m],\qquad [em,en]=0.
\]
The Chevalley coalgebra is
$\vCE(\mathfrak h)=K\oplus\bigoplus_{r\geq1}
\operatorname{Sym}^r_\star(s\mathfrak h)$, with unshuffle coproduct.
Its coderivation has
$Q_1(sx)=-s(dx)$ and
$Q_2(sx\,sy)=(-1)^{|x|}s[x,y]$.
The \emph{actor weight} counts the factors $s(em)$, which have degree
zero. It is independent of PBW output number.

A degree-one twisting cochain to $A$ can be nonzero only on one
ordinary suspended input. Write its value as
$\theta_w(x;m_1,\ldots,m_w)$ and put $\theta_0(x)=x$.
For $m_S=\mathop{\circ}_{i\in S}m_i$ in
$P_{|S|}=\vSym^{|S|}\mathfrak g$, the twisting equation is
\begin{equation}\label{eq:vol2-profile-MC}
\begin{aligned}
0={}&[x,\theta_w(y;m)]-\theta_w(y;[x,m])
       -[y,\theta_w(x;m)]+\theta_w(x;[y,m])
       -\theta_w([x,y];m)\\
 &+\sum_{\substack{S\sqcup T=[w]\\S,T\ne\varnothing}}
 \bigl(\theta_{|S|}(x;m_S)\theta_{|T|}(y;m_T)
       -\theta_{|S|}(y;m_S)\theta_{|T|}(x;m_T)\bigr).
\end{aligned}
\end{equation}
Here $[x,m]$ is the sum of brackets with each module input, and the
splitting $(S,T)$ is ordered. To verify the signs, use convolution
\[
 (f*g)(c)=\sum_{(c)}(-1)^{|g||c'|}f(c')g(c'').
\]
On two ordinary suspended inputs, the tensor sign of two degree-one
cochains is minus. The terms with no actor input on one side give
the four adjoint terms above; the Chevalley bracket gives the fifth.
All other components vanish by degree. Thus
\eqref{eq:vol2-profile-MC} is both necessary and sufficient.

For a nonempty set $B$ of actor labels, define
\[
 u_B=\frac1{|B|!}\sum_{\sigma\in S_B}
                   m_{\sigma(1)}\cdots m_{\sigma(|B|)}.
\]
Let $\mathcal T_r$ consist of nonnegative tuples
$(a,b_1,\ldots,b_r)$ such that
\[
 a+\sum_i b_i=r,\qquad
 a+\sum_{i=1}^j b_i\geq j+1\quad(0\leq j<r).
\]
These are the child counts of plane rooted trees in preorder, with
the root first. For envelope inputs $v_1,\ldots,v_r$, set
\begin{equation}\label{eq:vol2-profile-trees}
 T_r^D(x;v_1,\ldots,v_r)=
 \sum_{(a,b)\in\mathcal T_r}\frac{D^ax}{a!}
 \sum_{\sigma\in S_r}
 \frac{D^{b_1}v_{\sigma(1)}\cdots D^{b_r}v_{\sigma(r)}}
      {b_1!\cdots b_r!},\qquad T_0^D(x)=x.
\end{equation}
The slots are symmetric; each displayed product remains ordered.
For a scalar sequence $\mathbf t=(t_1,t_2,\ldots)$, define
\begin{equation}\label{eq:vol2-profile-partitions}
 \theta_w^{\mathbf t}(x;m_1,\ldots,m_w)=
 \sum_{\pi\in\operatorname{Part}([w])}
       \left(\prod_{B\in\pi}t_{|B|}\right)
       T_{|\pi|}^D(x;u_B:B\in\pi).
\end{equation}
The partitions are unordered, with nonempty blocks. For example,
\begin{align*}
 \theta_1^{\mathbf t}(x;m)&=t_1Dx\,m,\\
 \theta_2^{\mathbf t}(x;m,n)
 &=t_2Dx\,\frac{mn+nm}{2}
   +t_1^2\left(D^2x\,\frac{mn+nm}{2}
                   +Dx(Dm\,n+Dn\,m)\right).
\end{align*}
In weight three the term proportional to $t_1t_2$ is the sum of
$T_2^D(x;m_i,u_B)$ over $\{i\}\sqcup B=[3]$.

\begin{theorem}[Simultaneous profiles]\label{thm:vol2-profiles}
The maps \eqref{eq:vol2-profile-partitions} satisfy
\eqref{eq:vol2-profile-MC}. They define a map of coassociative
coalgebras
\[
 \vCE(\mathfrak g\ltimes e\mathfrak g)
       \longrightarrow\vCE(\mathfrak g)
\]
under $\vCE(\mathfrak g)$, in the derived category of coalgebras.
Every nonconstant term on a configuration $I$ has actor weight
$w<|I|$. Thus each scalar sequence defines coefficient maps without
completing states or normal derivatives.
\end{theorem}

\begin{proof}
Use a central formal variable $z$ and define right substitution by
\[
 f\{R\}=\sum_{j\geq0}\frac{D^jf}{j!}R^j,
 \qquad U=\sum_{r\geq1}\frac{t_rz^r}{r!}m^r,
 \qquad R=U\{R\},\qquad G_x=x\{R\}.
\]
Every coefficient of $R\in zA[[z]]$ is determined by the preceding
ones. Let $\partial_n$ vary $m$ in direction $n$, including its
$D$-derivatives. Adjoint equivariance gives
$\partial_{[x,m]}U=[x,U]$. Put $U_j=D^jU/j!$ and $X_i=D^ix/i!$.
Differentiating the fixed-point equation gives a linear filtered
equation for $\partial_{[x,m]}R$. The same equation is satisfied by
$[G_x,R]$: the difference of its expanded sides is
\[
 \sum_{i,j}X_i[R^i,U_j]R^j
 =\sum_iX_i\sum_{p+q=i-1}R^p
                 \left(\sum_j[R,U_j]R^j\right)R^q=0,
\]
because the inner sum is $[R,R]$. The linear term of the equation
raises the $z$-filtration, so uniqueness gives
$\partial_{[x,m]}R=[G_x,R]$.

Leibniz for $D$ now implies
\begin{equation}\label{eq:vol2-profile-transport}
 [G_x,G_y]-G_{[x,y]}
       -\partial_{[x,m]}G_y+\partial_{[y,m]}G_x=0.
\end{equation}
Indeed, with $Y_j=D^jy/j!$, the first two terms differ by
$\sum_{i,j}X_i[R^i,Y_j]R^j-\sum_{i,j}Y_j[R^j,X_i]R^i$.
The two derivative terms cancel precisely these sums. Expanding the
root and its ordered subtrees gives
\eqref{eq:vol2-profile-trees}; allocating the original inputs to the
subtree slots gives \eqref{eq:vol2-profile-partitions}. Its
normalization is
$G_x=x+\sum_{w\geq1}z^w\theta_w^{\mathbf t}(x;m,\ldots,m)/w!$.
Coefficient extraction in \eqref{eq:vol2-profile-transport}, followed
by polarization, proves the twisting equation.

A twisting cochain defines a map into $\vB A$. The Chevalley--bar
comparison identifies $\vB A$ with $\vCE(\mathfrak g)$.
Explicitly, this comparison sends a symmetric word in suspended Lie
inputs to the sum of its signed bar words, without a factorial divisor.
It preserves the unshuffle and deconcatenation coproducts. PBW reduces
its associated graded to the Koszul comparison for a symmetric algebra.
The usual Koszul homotopy, divided by positive polynomial weight,
contracts its augmentation complex. At $I$, both filtrations stop
at $|I|$, so the comparison is an equivalence. Its inverse in the
pointed derived category gives the asserted map.

Each nonconstant term consumes one positive ordinary input and $w$
positive actor inputs. Hence $w+1\leq|I|$. All operations in each
surviving term are full coefficient maps, so the finite formulas
retain collision and normal-transfer compatibility.
\end{proof}

For $D=\partial_q$, $x=m=q$ and $t_2=1$ with the other parameters
zero, right substitution becomes $f(q+R)$. Thus $p=q+R$ satisfies
the scalar equation at the chapter's opening. The noncommutative
proof explains why the order of the powers of $R$ is fixed.

The same construction supplies mixed homotopies. In the relative
convolution algebra with differential $d_0$ at the scalar action,
write the positive cochain as
\[
 \theta_{\mathbf t}=\sum_rt_ru_r+
             \sum_{p\leq q}t_pt_q\eta_{pq}+\cdots.
\]
Its equation gives $d_0\eta_{pq}=-[u_p,u_q]$ for $p<q$ and
$d_0\eta_{pp}=-u_p^2$. These are compatible primitives because they
are coefficients of one Maurer--Cartan element. More explicitly,
give a free graded Lie generator $e_\alpha$ degree one for every
nonzero finite multi-index $\alpha$, with weight $\sum r\alpha_r$,
and set
\begin{equation}\label{eq:vol2-profile-cells}
 de_\alpha=-\frac12
   \sum_{\substack{\beta+\gamma=\alpha\\\beta,\gamma\ne0}}
                  [e_\beta,e_\gamma].
\end{equation}
This is a differential. In the tensor envelope it is
$de_\alpha=-\sum_{\beta+\gamma=\alpha}e_\beta e_\gamma$;
applying $d$ twice gives opposite sums over ordered three-part
splittings. Coefficient extraction from the profile equation shows
that $e_\alpha\mapsto[\mathbf t^\alpha]\theta_{\mathbf t}$ is a
dg Lie map. Only the positive-weight completion is used.

\section{A change of actor and an actual bar composite}
\label{sec:vol2-actor-composite}

For a vertex algebra $V$, let $L_V$ be its main-diagonal Lie family
on \eqref{eq:vol2-action-category}, with singleton the right
differential module of states and connection induced by translation.
In each ordered chart its higher coefficient is the full normal
direct image of this module on the small diagonal. The chiral bracket
uses principal parts. Put $C_V=\vCE(L_V)$.
For the vacuum vertex algebra $k$, this is the dualizing coalgebra
$J$. For $E=\Lambda(e)$, $|e|=1$, with zero differential and
translation, write $C_E$ for its factorizing coalgebra.

The following tensor comparison specifies how the algebraic profiles
act on the coefficient family.

\begin{lemma}\label{lem:vol2-strong-boundary-tensor}
If $C$ is a cocommutative factorizing coalgebra with empty
coefficient $k$, then $F\mapsto F\vot C$ is a colimit-preserving
strong symmetric monoidal functor for $\star$. For $C=C_V$ it
sends $L_k$ to $L_V$ and $L_E$ to $L_V\ltimes eL_V$, with their
Lie structures. Consequently
\begin{equation}\label{eq:vol2-exterior-comparison}
 C_E\vot C_V\simeq\vCE(L_V\ltimes eL_V).
\end{equation}
The comparison preserves the vacuum, actor grading, finite scalar
extensions, and all collision maps and normal derivatives.
\end{lemma}

\begin{proof}
For each ordered partition, tensor its convolution summand with
$C_I$ and use \eqref{eq:vol2-action-factorization} for $C$.
The projection formula for localization identifies the result with
the corresponding summand of
$(F\vot C)\star(G\vot C)$. At the empty coefficient the map is
the scalar identity. All iterated comparisons use the same ordered
partition and the same iterated coproduct of $C$. This gives
associativity and units; cocommutativity gives symmetry. Collision
pullback keeps precisely the saturated partitions on both sides,
with the same base-change maps. Thus the functor is strong monoidal
on the enhanced diagram and commutes with its Chevalley complexes,
envelopes and bars.

On the main diagonal, the projection formula identifies
$L_k\vot C_V$ with $L_V$. The principal-part boundary of the vacuum
line, tensored with the factorization identification of $C_V$ on the
punctured two-point locus, is the chiral bracket of $V$. This checks
the Lie identity on both ordered branches with the full normal
direct image. The even vacuum and odd $e$ summands give the stated
semidirect product. These coefficient formulas preserve actor
scaling and tensor powers. Finite free scalar extensions preserve
their equivalences, which proves the remaining assertions.
\end{proof}

For $\mathfrak g=L_V$ and this comparison, write
$\alpha_{\mathbf t}:C_E\vot C_V\to C_V$ for the profile action.
The sequence $t_1=1$, $t_r=0$ for $r>1$ gives $\alpha_D$.
The actor maps receiving the other profiles are nonlinear even
though a map between the exterior coefficient algebras is linear
on its degree-one generator.

Let $\mathfrak l=L_k$. In $\vU(\mathfrak l\ltimes e\mathfrak l)$
define
\begin{equation}\label{eq:vol2-marked-PBW}
 \rho_r(m_1,\ldots,m_r)=\frac1{r!}\sum_{\sigma\in S_r}
       m_{\sigma(1)}\cdots m_{\sigma(r-1)}(em_{\sigma(r)}).
\end{equation}
This is the averaged PBW product with one marked odd factor.
In particular its diagonal value is $m^{r-1}em$. The cochain
\[
 \tau_{\mathbf t}(sx)=x,\qquad
 \tau_{\mathbf t}(s(em_1)\cdots s(em_r))=t_r\rho_r(m_1,\ldots,m_r)
\]
has all other reduced components zero.

\begin{proposition}\label{prop:vol2-profile-actor}
The cochain $\tau_{\mathbf t}$ is twisting and defines a map
$\Sigma_{\mathbf t}:C_E\to C_E$ under and over $J$.
Its formal coordinate series is
$\sum_{r\geq1}t_rp^r/r!$. There is a pointed action path
\begin{equation}\label{eq:vol2-profile-comparison}
 \alpha_D(\Sigma_{\mathbf t}\vot\vid_{C_V})
                    \simeq\alpha_{\mathbf t}.
\end{equation}
It preserves every ordered cut, collision, normal transfer and unit
homotopy. In the associative bar its composite corestriction is
exactly \eqref{eq:vol2-profile-partitions}.
\end{proposition}

\begin{proof}
The maps $\rho_r$ are adjoint-equivariant, giving the twisting
equation with one ordinary input. With two ordinary inputs it is
the equation for the canonical cochain. On pure actor inputs,
polarization reduces the equation to the square of
\[
 \sum_{r\geq1}\frac{t_r}{r!}m^{r-1}em.
\]
It is zero because $[m,em]=0$ and $(em)^2=0$. The induced map
to the associative bar is the finite sum of tensor powers of
$s\tau_{\mathbf t}$ applied to the iterated reduced coproduct.
The inverse Chevalley--bar comparison gives $\Sigma_{\mathbf t}$.
Its restriction to $J$ is the identity and its augmentation is
fixed. Removing the marked odd factor in
\eqref{eq:vol2-marked-PBW} gives the scalar coordinate coefficient
$t_r/r!$, with exponential polarization.

To calculate the composite, put
$B=\vU(\mathfrak g\ltimes e\mathfrak g)$ and let $\pi:B\to A$
kill the odd ideal. There is a degree-$-1$ projection derivation
\begin{equation}\label{eq:vol2-actor-contraction}
 q(x)=0,\quad q(em)=m,\quad q(1)=0,\qquad
 q(uv)=q(u)\pi(v)+(-1)^{|u|}\pi(u)q(v).
\end{equation}
It respects the envelope relations: the mixed relation reduces to
the same bracket identity, while a relation with two odd inputs
maps to zero. In $\operatorname{Hom}(\vB B,A)$ let $a$ be the
canonical length-one twisting cochain induced by $\pi$ and $p$
the length-one cochain induced by $q$. They have degrees one and
zero and satisfy
\[
 da+a^2=0,\qquad dp+[a,p]=0.
\]
Apply $D$ on the output and solve
\[
 R=\sum_{j\geq0}\frac{D^jp}{j!}R^j,
 \qquad A_R=\sum_{j\geq0}\frac{D^ja}{j!}R^j
\]
in positive bar length. These equations give a twisting cochain.
For completeness, put $a_j=D^ja/j!$, $p_j=D^jp/j!$.
The differentiated identities are
$da_j=-\sum_{u+v=j}a_ua_v$ and
$dp_j=-\sum_{u+v=j}[a_u,p_v]$. Moreover,
\[
 \sum_{r,s}[a_r,p_s]R^{r+s}
 =[A_R,R]-\sum_s p_s[A_R,R^s],
\]
since their difference is
$\sum_r a_r[R^r,\sum_s p_sR^s]=0$.
The differentiated fixed-point equation therefore has the unique
solution $dR=-[A_R,R]$. Using that each $a_j$ is odd gives
\[
 dA_R=-\sum_{r,s}a_ra_sR^{r+s}
                   +\sum_r a_r[A_R,R^r],\qquad
 A_R^2-\sum_{r,s}a_ra_sR^{r+s}
                   =\sum_r a_r[R^r,A_R].
\]
Thus $dA_R+A_R^2=0$.

The associated bar map $\mathcal F_D:\vB B\to\vB A$ has
length-$(w+1)$ corestriction on $[v_0|\cdots|v_w]$
\begin{equation}\label{eq:vol2-profile-bar-tree}
 \sum_{(a,b)\in\mathcal T_w}
 \frac{D^a\pi(v_0)}{a!}
 \frac{D^{b_1}q(v_1)\cdots D^{b_w}q(v_w)}{b_1!\cdots b_w!}.
\end{equation}
It vanishes unless the first input has exterior degree zero and the
others have exterior degree one. Restriction to the Chevalley coalgebra
is the tree action $\alpha_D$.

Lemma~\ref{lem:vol2-strong-boundary-tensor} carries the actor map
from $\mathfrak l$ to $\mathfrak g$. Its bar representative
partitions the actor inputs into marked blocks and places the one
ordinary input in a singleton. Formula
\eqref{eq:vol2-profile-bar-tree} forces this singleton first.
Equation~\eqref{eq:vol2-actor-contraction} gives $q\rho_r=u_{[r]}$.
For each unordered partition, the sum over orders of its blocks
is therefore $T_{|\pi|}^D$, with coefficient
$\prod_{B\in\pi}t_{|B|}$. This proves equality of the actual
corestrictions. The pointed Chevalley--bar inverses and the strong
tensor comparison transport that equality to
\eqref{eq:vol2-profile-comparison}, including all unit paths.
Each calculation is finite on every configuration.
\end{proof}

The quadratic map with $t_2=1$ has coordinate series $p^2/2$.
The map with coordinate series $p^2$ has $t_2=2$. This normalization
is independent of the unnormalized permutation sum in the
Chevalley--bar comparison.

\section{Natural homotopies and multiplication}
\label{sec:vol2-natural-actions}

The profile comparison gives a map to $\alpha_D$. To receive
compositions, that target must carry its own coherent multiplication.
The obstruction groups can be computed before evaluating on a
particular boundary.

Let $W$ be finite-dimensional in degree zero and set
$E_W=\Lambda(W[-1])$, whose generators have degree one.
Let $\mathfrak d$ be a Lie algebra of degree-zero derivations of an
ordinary Lie algebra $\mathfrak g$. A \emph{natural operation} uses
each Lie input once, applies elements of $U(\mathfrak d)$ to these
inputs, and multiplies the results in a specified order in
$U(\mathfrak g)$. Brackets are commutators. Operations on
$\operatorname{CE}(\mathfrak g\otimes E_W)$ include signed
symmetrization of suspended inputs. Two formulas are identified if
they agree for every Lie algebra with such a derivation action.
This restriction defines a universal algebra of operations; it does
not define all maps of geometric coefficient families.

Put $P=W^*$. Denote the natural convolution algebra by
$\mathcal N(P;\mathfrak d)$, with differential twisted by the scalar
cochain $\tau$. Thus
\[
 \delta f=-(-1)^{|f|}fQ+\tau*f-(-1)^{|f|}f*\tau.
\]
Complete in actor weight and allow only finite sums of differential
operations at each weight and cochain degree. The ideal of positive
actor weight is $\mathcal N_+$. Natural actions and homotopies mean
its Maurer--Cartan space, using polynomial differential forms on
simplices as parameters.

For commuting derivations $D_1,\ldots,D_r$, the unary cochain
$x_{S,a}$ extracts an exterior coefficient $e_S$ and applies
$D^a/a!$. It has degree $1-|S|$. Write $t=x_{\varnothing,0}$.
If $\epsilon(B,C)$ is the exterior shuffle sign, the unary
differential is
\begin{equation}\label{eq:vol2-natural-differential}
\begin{aligned}
 \delta x_{S,a}
 ={}&-\sum_{\substack{B\sqcup C=S\\b+c=a}}
       (-1)^{|B|(|C|+1)}\epsilon(B,C)x_{B,b}x_{C,c}\\
    &+tx_{S,a}-(-1)^{1-|S|}x_{S,a}t.
\end{aligned}
\end{equation}
It follows from
$Q_2(s(e_Bx)s(e_Cy))=(-1)^{|B|}s(e_Be_C[x,y])$
and divided Leibniz. The two terms containing $t$ cancel when
$x_{S,a}\ne t$, while $\delta t=t^2$.
For one actor and one derivation,
\[
 \delta x_{\{1\},1}=x_{\{1\},0}x_{\varnothing,1}
                      -x_{\varnothing,1}x_{\{1\},0},
 \qquad \delta x_{\varnothing,2}=-x_{\varnothing,1}^2.
\]
These identities already show the polynomial actor coordinate and
the exterior derivation class in cohomology.

\begin{theorem}\label{thm:vol2-natural-cohomology}
There is a surjective filtered dg algebra map
\begin{equation}\label{eq:vol2-natural-projection}
 \pi:\mathcal N(P;\mathfrak d)\longrightarrow
          \prod_{q\geq0}\operatorname{Sym}^q(P)
                     \otimes\Lambda(\mathfrak d[-1])
\end{equation}
with acyclic kernel in each actor weight. The target differential
is zero; its actor coordinates have degree zero and its derivation
generators degree one. On positive actor weights it induces an
equivalence of Maurer--Cartan spaces. Their components are
\begin{equation}\label{eq:vol2-natural-classes}
 \prod_{q\geq1}\operatorname{Sym}^q(W^*)\otimes\mathfrak d.
\end{equation}
At each point the fundamental group is the additive group
$\mathfrak m_W=\prod_{q\geq1}\operatorname{Sym}^q(W^*)$,
and all higher homotopy groups vanish.
\end{theorem}

\begin{proof}
Unary operations freely generate the convolution algebra. To prove
this, evaluate on the free Lie algebra on named copies of the free
$\mathfrak d$-module $U(\mathfrak d)$. In its tensor envelope,
an ordered monomial records the input, its exterior label and the
operator applied to it. Distinct such words are independent.
Polarization recovers repeated exterior types with the suspended
signs. A cochain of actor weight $q$ and degree $p$ uses $q+p$
inputs.

The Hopf coproduct of $U(\mathfrak d)$ expresses its Leibniz rule.
Rescaling the exterior-label generator by
$(-1)^{\binom{|S|}{2}}$ identifies the differential with the cobar
differential of $\Lambda(P[1])\otimes U(\mathfrak d)$, apart from
the scalar generator $t$. Thus the natural algebra is
\begin{equation}\label{eq:vol2-natural-cobar}
 (k\langle t\rangle,\delta t=t^2)
       *\Omega\bigl(\Lambda(P[1])\otimes U(\mathfrak d)\bigr).
\end{equation}
Here $*$ is free product, with the stated weight completion.
The augmentation ideal of the first factor contracts by
$h(t^{2j})=t^{2j-1}$ and $h(t^{2j-1})=0$. In the free product,
contract the first such factor of each alternating word, with
the sign of its prefix. This removes $t$ without changing
cohomology.

PBW symmetrization identifies $U(\mathfrak d)$ with the symmetric
coalgebra on $\mathfrak d$. This is a coalgebra identification,
not an identification of products. The cobar projection sends
primitive actor and derivation generators to their namesakes and
all other generators to zero. The cobar differential maps to
graded commutators or odd squares, so the projection is a dg
algebra map.

Its cohomology is computed in finite auxiliary weights. The dual
algebra of $\Lambda(P[1])\otimes\operatorname{Sym}^c(\mathfrak d)$
is exterior on degree-one variables $\alpha_i$ and polynomial on
degree-zero variables $z_j$. A free resolution of its scalar
module is
\[
 \Lambda(\alpha_i)\otimes k[z_j]
       \otimes k[v_i]\otimes\Lambda(\eta_j),\qquad
 |v_i|=0,\quad |\eta_j|=-1,\quad
 dv_i=\alpha_i,\quad d\eta_j=z_j.
\]
For an odd pair, $d(v^{a+1})=(a+1)\alpha v^a$; for an even
pair this is the two-term Koszul resolution. Each pair has scalar
cohomology, so their tensor product resolves $k$. Applying Hom
into $k$ gives zero differential and the algebra on the right
of \eqref{eq:vol2-natural-projection}. The bar resolution computes
the same Ext. The primitive projection is surjective on
cohomology and both sides have the same finite dimension in each
auxiliary weight, proving the quasi-isomorphism. For an
infinite-dimensional derivation algebra take the filtered union
of finite-dimensional symmetric-coalgebra computations.

To lift a Maurer--Cartan element, solve successively in actor
weight. At the first unsolved weight, its curvature is a closed
degree-two element of the kernel of $\pi$. A degree-one primitive
removes it. The same argument with forms relative to a simplex
boundary lifts every compatible simplex. Tensoring a weightwise
acyclic complex with these forms preserves acyclicity over $k$.
The resulting fibres are contractible. In the target, the
commutator bracket and differential vanish. Its degree-one
elements give \eqref{eq:vol2-natural-classes}, its degree-zero
elements give the stated loops, and its negative degrees vanish.
\end{proof}

For one derivation, retain also the derivative order. The sector
where actor weight and derivative order both equal $q$ has
cohomology
\[
 \operatorname{Sym}^q(P)\otimes\Lambda^q(kD)
 \quad\text{in degree }q.
\]
It is therefore acyclic for $q\geq2$. For several commuting
derivations, the same formula has their span in place of $kD$;
its degree-two term need not vanish.

\begin{lemma}\label{lem:vol2-natural-bar-lifting}
Restriction along
$\operatorname{CE}(\mathfrak h)\to\operatorname B U(\mathfrak h)$
is a surjective quasi-isomorphism on natural convolution complexes
twisted by the scalar action. Its kernel contracts naturally,
preserving actor weights and derivative orders. It has the same
property relative to any union of actor unit faces.
\end{lemma}

\begin{proof}
For abelian $\mathfrak h$, multilinear PBW weight $d$ of the bar
is the finite complex of ordered partitions of $d$ named inputs,
with differential merging adjacent blocks. The Koszul resolution
used above identifies its homology with suspended symmetrization.
Choose rational complements to cycles and boundaries and average
over $S_d$. Inverting the differential on the boundary complement
gives a permutation-equivariant contraction, normalized to
\[
 pj=1,\quad hj=ph=h^2=0,\qquad dh+hd=1-jp.
\]
Evaluation on tensor powers gives a natural contraction for every
graded $\mathfrak h$, compatible with its internal differential.

For nonabelian $\mathfrak h$, PBW identifies the change in the
bar differential with a perturbation $b$ that lowers primitive
degree. The formulas
$p'=p(1+bh)^{-1}$ and $h'=h(1+bh)^{-1}$ are finite on each
input. The inclusion stays $j$, since $bj=jQ_2$ and $hj=0$.
The transferred differential is the Chevalley differential.
Precomposition by $p'$ splits restriction, and
$f\mapsto(-1)^{|f|}fh'$ contracts its untwisted kernel.

The scalar twist preserves that kernel because $j$ is a
coalgebra map. Its commutator raises the least primitive input
degree on which a cochain is nonzero. Perturbing the kernel
contraction again therefore gives finite sums on each input.
All these maps use permutations, PBW operations, products and
brackets. They commute with exterior face projections and
preserve derivative orders. Restriction to cochains vanishing
on the chosen faces proves the relative assertion.
\end{proof}

Natural identities evaluate on the coefficient category because
they are identities of these multilinear operations. On a
configuration of $m$ positions with $n$ exterior generators,
actor weight is at most $mn$. Their formal weight products
introduce no completion of a state or of a normal derivative.

\begin{theorem}\label{thm:vol2-derivation-multiplicative}
For a vertex derivation $D$ commuting with translation, the
action $\alpha_D$ admits a coherent unital multiplicative
structure for the ordinary exterior product of $C_E$.
Within natural operations of equal actor and derivative weight,
the higher fillings are contractible after fixing the unary
action and a compatible binary bar path.
Every profile $\alpha_{\mathbf t}$ has a map of pointed actions
to this multiplicative action. In particular its two unit
insertions have the common recipient
\begin{equation}\label{eq:vol2-profile-recipient}
 \mu(\Sigma_{\mathbf t}\vot\Sigma_{\mathbf t}):
 (C_E,\alpha_{\mathbf t})^{\circledast2}
                         \longrightarrow(C_E,\alpha_D).
\end{equation}
\end{theorem}

\begin{proof}
Use the natural bar complex for $n$ exterior actors. Restriction
to their proper unit faces is surjective: the coordinate
projections commute, so inclusion--exclusion extends compatible
face data. Its kernel is the fully mixed part. For $n\geq2$
that part has actor weight at least two and is acyclic in the
equal derivative-weight sector, by
Theorem~\ref{thm:vol2-natural-cohomology} and
Lemma~\ref{lem:vol2-natural-bar-lifting}.

Consequently compatible Maurer--Cartan data on the actor faces
and the boundary of a parameter cube extend across the cube.
Lift the data linearly, then kill the first remaining curvature
in each actor weight by a primitive in the fully mixed complex
tensored with the forms relative to the cube boundary. The
Bianchi identity makes the curvature closed. Further parameter
simplices give the same lifting, proving contractibility of
the relative choices.

The component for $n$ actors uses the cube $[0,1]^{n-1}$.
At $t_i=0$ multiply the adjacent actors; at $t_i=1$ compose
the two shorter bar maps. An end unit deletes its adjacent
parameter. An interior unit combines the two parameters by
\[
 a\vee b=a+b-ab.
\]
This polynomial operation is associative, with identity zero
and absorbing element one. It makes all repeated unit
restrictions agree with the endpoint faces. The preceding
lifting constructs the binary path and then every higher
component, or extends a specified compatible binary path.

The endpoint compositions are the geometric ones. Indeed,
$\vid+hD$ over $k[h]/(h^2)$ is a vertex automorphism, and
Lemma~\ref{lem:vol2-strong-boundary-tensor} carries $D$ to
$D\otimes1$ on exterior coefficient extensions. Its naturality
identifies every bar operation and contraction with the
corresponding coefficient map. Thus the constructed families
retain collisions, normal transfers and the vacuum.

Proposition~\ref{prop:vol2-profile-actor} supplies each incoming
map. Composing its tensor square with $\mu$ gives
\eqref{eq:vol2-profile-recipient}. On either unit axis it is
$\Sigma_{\mathbf t}$ with its action path. The unital coherence
identifies these restrictions on their common vacuum, supplying
the required relative two-simplex. The same higher product
coherences receive any finite ordered product of profile maps.
\end{proof}

The contractibility is in the specified natural-operation
complex. The space of all geometric multiplicative structures
can have additional components or homotopy groups.

\section{First-weight classes and actor size}
\label{sec:vol2-first-weight}

Fix $\kappa\in k^\times$ and the Heisenberg vertex algebra
\[
 V_\kappa=k[x_1,x_2,\ldots],\qquad
 \mathsf j=x_1,\qquad x_n=\mathsf j_{-n}\mathbf1.
\]
Its states are polynomials. Translation and the positive current
modes are
\[
 Tx_n=nx_{n+1},\qquad
 \mathsf j_n=\kappa n\partial_{x_n}\quad(n>0).
\]
Let $L_0x_n=nx_n$, extended as a derivation of the polynomial
algebra. Put $\mathfrak g=L_{V_\kappa}$, $A=\vU\mathfrak g$,
$P_w=\vSym^w\mathfrak g$ and $C=C_{V_\kappa}$. Lie-module
actions on $P_w$ and $A$ are adjoint actions. A coefficient
derivation is a degree-zero map of the full family $\mathfrak g$
which obeys Leibniz for its Lie bracket. Their Lie algebra is
$\mathfrak d_V=\operatorname{Der}_{\mathscr C}(\mathfrak g)$.

In left-module coordinates, the singleton coefficient has
connection $\partial_y+T$. A coefficient derivation is therefore
represented there by a $k[y]$-linear map $B(y)$ with
\begin{equation}\label{eq:vol2-horizontal-derivation}
 \partial_yB+[T,B]=0,
\end{equation}
together with compatibility with the chiral bracket. Each input
state has a finite polynomial image. A common polynomial-degree
bound on all inputs is not imposed. For example,
$B(y)=L_0+yT$ is horizontal because $[T,L_0]=-T$. The conformal
identity for $Y(a,z)$ gives
\[
 B(q)Y(a,p-q)b
   =Y(B(p)a,p-q)b+Y(a,p-q)B(q)b.
\]
Taking arbitrary Laurent principal parts proves its coefficient
derivation identity, including the normal derivatives.

The current shift $D=\partial_{x_1}$ is another such derivation:
it is the derivative of $\mathsf j\mapsto\mathsf j+t\mathbf1$.
Every state zero mode $a_{(0)}$ is a vertex derivation commuting
with $T$, by the mode commutator identity. Put
$Z_r=(\mathsf j^r)_{(0)}$ for $r\geq2$. Normal ordering gives
\begin{equation}\label{eq:vol2-Heisenberg-derivations}
 Z_2=2\kappa T,\qquad
 Z_r\mathsf j=r(r-1)\kappa x_1^{r-2}x_2,
 \qquad (\mathsf j^2)_{(1)}=2\kappa L_0.
\end{equation}
Indeed the contraction of $\mathsf j^r(z)$ with one current is
$r\kappa z^{-2}{:}\mathsf j(z)^{r-1}{:}$, whose residue is
$r\kappa T(\mathsf j^{r-1})$. There is no multiple contraction
with the single current on the right. The quadratic zero and
first modes give the displayed formulas for $T$ and $L_0$ on
every polynomial state. The images of $\mathsf j$ show that
$D,Z_2,Z_3,\ldots$ are linearly independent.

\subsection{An additional output root}

The family $P_w$ is the open direct image of its coefficient
on $w$ distinct roots, supported on configurations with at most
$w$ distinct positions. On the complement of the locus with
fewer roots, each ordered partial diagonal is a copy of
$\operatorname{Conf}_w(X)$. Its coefficient is the tensor of
$w$ state modules, equivariant under $S_w$. Reconstruction is
exact: use closed direct image on each partial diagonal and
localization along its cross-root divisors. This also retains
every normal derivative.

\begin{lemma}\label{lem:vol2-root-support-bound}
For $a,w\geq1$,
\[
 \operatorname{Ext}^q_{\mathfrak g}(P_w,P_a)=0
                       \quad\text{if }q<a-w.
\]
For every coefficient derivation $B$, the formula
\begin{equation}\label{eq:vol2-extra-root}
 \gamma_B^{(w)}(x;v_1\cdots v_w)
        =m_{w+1}(Bx,v_1,\ldots,v_w)
\end{equation}
defines a class in $\operatorname{Ext}^1_{\mathfrak g}(P_w,P_{w+1})$,
where $m_{w+1}$ is symmetric multiplication. The map from
$\mathfrak d_V$ to this extension group is injective.
\end{lemma}

\begin{proof}
A coefficient family supported on at most $a-1$ roots has zero
derived Hom to $P_a$, by open adjunction. The $n$th Chevalley
term computing the desired module maps is
\[
 \vRH_{\mathscr C}(\Lambda^n_\star\mathfrak g\star P_w,P_a)[-n].
\]
It vanishes for $w+n<a$. Before the shift it has no negative
cohomology, since its inputs are coefficient modules in degree
zero. The first-quadrant totalization proves the bound.

Give $P_{w+1}\oplus P_w$ the triangular action with
off-diagonal term \eqref{eq:vol2-extra-root}. Its module defect
is symmetric multiplication by
$[x,By]+[Bx,y]-B[x,y]$, hence is zero on the full coefficient
diagram. Open adjunction also gives
$\operatorname{Hom}_{\mathscr C}(P_w,P_{w+1})=0$.
The underlying split extension has only its canonical
coefficient splitting. If $B\ne0$, this splitting is not
equivariant: on $w+1$ distinct roots put a state on which $B$
is nonzero at the acting root and the vacuum at all other
roots. Its image under \eqref{eq:vol2-extra-root} is nonzero.
The same argument for differences proves injectivity.
\end{proof}

To compare homotopies, resolve the pointed arrow before taking
relative convolution cochains. Bar--cobar duality gives maps
under $A$ from $U_\star(\mathfrak g\ltimes e\mathfrak g)$ to
$A$. At the scalar action, the weight-$w$ complex is
\begin{equation}\label{eq:vol2-action-tangent}
 T_w=\vRH_{\mathfrak g}(P_w,A).
\end{equation}
Indeed the underlying relative Chevalley comodule is
$\vCE(\mathfrak g)\star\operatorname{Sym}^c_
 \star(e\mathfrak g[1])$, and $e\mathfrak g[1]=\mathfrak g$.
The coefficient resolutions are taken in the full enhanced
diagram. This computes the homotopy fibre specifying the
pointing, rather than the kernel of an unresolved restriction.

Scaling $e$ by a parameter over $k[t]/(t^2)$ defines
\[
 \nu:\pi_0\Phi_C(C_E)\longrightarrow H^1(T_1).
\]
The first coefficient is a cocycle; a pointed homotopy changes
it by a boundary. More generally a relative homotopy has a
normalized convolution gauge $G=1+g_1+g_2+\cdots$ satisfying
$d_0G=G\theta-\theta'G$. This follows by integrating the
polynomial interval equation successively in positive weight.
Every inverse and exponential is finite on a fixed
configuration. If both actions vanish in weights below $r$,
the equations give
\begin{equation}\label{eq:vol2-leading-gauge}
 d_0g_j=0\quad(j<r),\qquad
 d_0g_r=\theta_r-\theta'_r.
\end{equation}
The lower gauge coefficients need not be zero.

\begin{theorem}\label{thm:vol2-first-weight-classification}
There are vector-space identifications
\begin{equation}\label{eq:vol2-first-weight-classification}
 \operatorname{Ext}^1_{\mathfrak g}(P_1,P_2)
       \cong\mathfrak d_V,
 \qquad H^1(T_1)=k\xi\oplus\mathfrak d_V.
\end{equation}
The first identification is $B\mapsto[\gamma_B^{(1)}]$.
The class $\xi$ changes the zero-mode charge at one root.
For every $B\in\mathfrak d_V$, the tree action $\alpha_B$
has first class $[\gamma_B^{(1)}]$.
\end{theorem}

\begin{proof}
In the Chevalley extension spectral sequence
\[
 E_1^{n,q}=\operatorname{Ext}^q_{\mathscr C}
      (\Lambda^n_\star\mathfrak g\star\mathfrak g,P_2),
 \qquad n,q\geq0,
\]
open adjunction kills the $n=0$ column. Only $(n,q)=(1,0)$
can contribute in total degree one. There are no incoming
boundaries, and every higher differential has a negative index.
Thus the extension group consists precisely of module cocycles
$\Gamma:\mathfrak g\star\mathfrak g\to P_2$.

On the two-root locus write $R=k[p,q,(p-q)^{-1}]$. Such a
map is an $R$-linear endomorphism of
$R\otimes(V_\kappa\otimes V_\kappa)$ with
\[
 \partial_p\Gamma+[T\otimes1,\Gamma]=0,
 \qquad\partial_q\Gamma+[1\otimes T,\Gamma]=0.
\]
Its individual input images are finite; no uniform pole bound
is assumed. In the module cocycle equation give the two acting
inputs positions $r,p$ and the module input position $q$.
Away from the triple diagonal the target has three disjoint
pieces $r=q$, $p=q$, $r=p$, whose open direct images detect
equality.

On $r=q$, comparison of the zeroth normal coefficients after
testing $(r-q)^n$ gives
\[
 \Gamma(y\otimes x_{(n)}v)=(1\otimes x_{(n)})\Gamma(y\otimes v)
 \quad(n\geq0).
\]
The full normal transfer is retained; this projection supplies
necessary commutation identities. Take $x=\mathsf j$, $n>0$,
and $x=\mathsf j^2/(2\kappa)$, $n=1$.
Then $\Gamma$ commutes with all partial derivatives and $L_0$
in the second state factor. Its restriction to the weight-zero
line of that factor defines $B$. By induction on second-state
weight, $\Gamma(a\otimes v)-Ba\otimes v$ is killed by all
second-factor partial derivatives. It is therefore constant
in that factor, while having positive weight, and is zero.
Hence $\Gamma=B(p,q)\otimes1$.

The second horizontality equation gives $\partial_qB=0$.
A $q$-invariant element of $R$ lies in $k[p]$: it depends only
on $p$ in the fraction field, and substitution $q=p-1$ makes
it polynomial. This applies to each matrix coefficient, so
$B=B(p)$ and it obeys \eqref{eq:vol2-horizontal-derivation}.
On $r=p$ the cocycle equation is the derivation identity for
$B$, tensored with the unchanged spectator at $q$. The normal
transfer differentiates its coordinate dependence; $B(r)$ is
not replaced by $B(p)$ before that transfer. Tensoring with a
spectator vacuum and inverting $q-p$ is faithful, so this
identity holds in the full coefficient family. This proves
surjectivity; Lemma~\ref{lem:vol2-root-support-bound} proves
injectivity.

Equivariant PBW decomposes $A$ by output number. It is both a
sum and a product, because only finitely many output numbers
occur at any configuration. Lemma~\ref{lem:vol2-root-support-bound}
leaves only one and two outputs in $H^1(T_1)$. At one root,
extensions are extensions of the unital moving Heisenberg
vacuum. The positive modes give the polynomial Koszul
contraction for $\kappa n\partial_{x_n}$. Its nonconstant
monomials contract by the Euler homotopy, leaving the zero-mode
class. The de Rham complex of the affine root coordinate has
only its scalar class, so this contribution is $k\xi$.
Concretely, over $k[\epsilon]/(\epsilon^2)$ take
$\mathsf j_0=\epsilon$ and connection
$\partial_y+T+\epsilon x_1/\kappa$. Its commutator with
$\mathsf j_1$ is $-\epsilon$, and the other current translation
relations remain unchanged. This is the charge extension;
its zero-mode action prevents a split. These calculations
identify the one-output summand.

Finally the first tree term is $Bx\,v$. PBW writes it as
$m_2(Bx,v)+[Bx,v]/2$. For the zero-cochain $B$,
$(dB)(x;v)=[x,Bv]-B[x,v]=-[Bx,v]$.
The output-one term is exact and the first class is the
asserted output-two extension.
\end{proof}

Horizontality in the proof does not mean that an arbitrary
two-root map has bounded poles on the entire state space.
For $z=p-q$, the operator
\begin{equation}\label{eq:vol2-horizontal-poles}
 \mathcal K_z=\kappa\sum_{m,n\geq1}
  \frac{(-1)^{m-1}(m+n-1)!}{(m-1)!(n-1)!}
      z^{-m-n}\partial_{x_m}\otimes\partial_{x_n}
\end{equation}
is horizontal and finite on each polynomial input. For its
scalar coefficients $c_{m,n}$,
$m c_{m+1,n}=-(m+n)c_{m,n}$ and
$n c_{m,n+1}=(m+n)c_{m,n}$; together with
$[T,\partial_{x_m}]=-(m-1)\partial_{x_{m-1}}$ these prove
the two connection equations. Its pole order on
$x_m\otimes x_n$ grows with $m+n$. The commutant condition,
not a uniform pole assumption, is what removes the poles in
Theorem~\ref{thm:vol2-first-weight-classification}.

\subsection{Bounds on an actor}

\begin{corollary}\label{cor:vol2-unbounded-arity-actions}
For every $N\geq1$ there is a nontrivial pointed Heisenberg
action which is the trivial action on every configuration with
at most $N$ positions, including all its collisions and normal
derivatives.
\end{corollary}
\begin{proof}
Take $t_N=1$ and the other profile parameters zero. Its first
nonzero cochain is $Dx\,u_{[N]}$. Its PBW output-$N+1$
component is the nonzero extension
$[\gamma_D^{(N)}]$ of
Lemma~\ref{lem:vol2-root-support-bound}.
Equation~\eqref{eq:vol2-leading-gauge} prevents a homotopy to
the trivial action. Every nonconstant term uses at least
$N+1$ positive supports, so vanishes on the stated smaller
configurations.
\end{proof}

\begin{proposition}\label{prop:vol2-higher-weight-product}
Suppose a pointed action of $C_E$ has zero positive weights
below $r>1$ and nonzero first class $[\theta_r]$. It cannot
be multiplicative for the standard exterior-algebra product.
\end{proposition}
\begin{proof}
Pull the two candidate binary actions back along the primitive
diagonal induced by $e\mapsto e_1+e_2$. Multiplication followed
by the action has first class $2^r[\theta_r]$. Composition has
first class $2[\theta_r]$. To see the latter with its paths,
scale the two actors independently by $s,t$ and truncate at
total degree $r$. The first nonzero jet of composition is
the sum of the pure jets, since a product of two perturbations
has degree at least $2r$. Unit homotopies identify the pure
terms. A pointed binary path would therefore force
$(2^r-2)[\theta_r]=0$ by
\eqref{eq:vol2-leading-gauge}, contrary to characteristic zero.
\end{proof}

The profile actions satisfy this obstruction when their first
parameter has index greater than one. They still have the
multiplicative recipient \eqref{eq:vol2-profile-recipient}.

For any pointed factorizing actor $F$, let $B_F$ be a cobar
model under $U_\star L_k$, with its induced adjoint $L_k$-action.
At the augmented map $C_E\to J\to F$, the first-weight
mapping complex is
\[
 T_1^{E,F}=\vRH_{L_k}(L_k,B_F).
\]
The relative Chevalley argument of
\eqref{eq:vol2-action-tangent} gives this identification,
without a degree bound on $B_F$. A fixed action $\beta$ of
$F$ induces a linear map
\begin{equation}\label{eq:vol2-action-chain-rule}
 b_\beta:H^1(T_1^{E,F})\longrightarrow H^1(T_1),\qquad
 \nu\bigl(\beta(f\vot\vid_C)\bigr)=b_\beta(\nu_F(f)).
\end{equation}
Indeed, terms with one perturbation in a compatible bar
composite give a chain map; terms with two perturbations
vanish over $k[t]/(t^2)$. Resolving the pointed arrow retains
its prescribed homotopy. This finite free scalar extension
commutes with derived coefficient maps and diagram limits.

\begin{lemma}\label{lem:vol2-finite-exterior-tangent}
For $W$ finite-dimensional in degree zero and
$F=C_{E_W}$, one has $H^1(T_1^{E,F})=W$.
Its vectors are represented by the linear exterior maps
$e\mapsto w$.
\end{lemma}
\begin{proof}
In PBW for $U_\star L_{E_W}$, let $a$ count output roots
and $b$ their total exterior degree. The support bound
puts Chevalley degree at least $a-1$. Total degree one
allows only $(a,b)=(1,0),(1,1),(2,0)$. The first is
$H^1_{\mathrm{dR}}(\mathbb A^1;k)=0$. The second is
$\operatorname{Hom}_{L_k}(L_k,L_k)\otimes W=W$.
For $(2,0)$, open adjunction leaves only the scalar
multiple $m$ of symmetric multiplication in Chevalley
degree one. Its differential is
$(dm)(x,y;v)=m([x,y],v)$, nonzero on a simple pole between
the acting points with the module root distinct. The vacuum
principal-part bracket has residue one there. Thus this
last summand supplies no class.
\end{proof}

\begin{theorem}\label{thm:vol2-finite-actor-obstruction}
If a pointed action $(F,\beta)$ represents $\Phi_C$, then
$H^1(T_1^{E,F})$ is infinite-dimensional. The same conclusion
holds if $(F,\beta)$ is a local centre in $\vMon$.
Consequently no coalgebra equivalent as a pointed object
to $C_{E_W}$, with $W$ finite-dimensional and $d=T=0$
on its exterior generators, can underlie either object.
\end{theorem}
\begin{proof}
Representability supplies an incoming map from every
$\alpha_{Z_r}$. Equation~\eqref{eq:vol2-action-chain-rule}
puts every $[\gamma_{Z_r}^{(1)}]$ in the image of $b_\beta$.
These classes are linearly independent by
\eqref{eq:vol2-Heisenberg-derivations} and
Theorem~\ref{thm:vol2-first-weight-classification}.
If only finality in $\vMon$ is assumed, the same argument
applies because each $\alpha_{Z_r}$ is multiplicative by
Theorem~\ref{thm:vol2-derivation-multiplicative}.
Lemma~\ref{lem:vol2-finite-exterior-tangent} excludes the
stated finite exterior actors in both cases.
\end{proof}

A finite exterior actor can still receive a specified finite
family. The theorem excludes its universal property for all
actions. It does not exclude a larger factorizing centre or
identify the whole formal deformation functor of a recipient.

\section{Recipients for nilpotent families}
\label{sec:vol2-nilpotent-recipients}

Several derivations need not commute. Their action classes compose
by the Baker--Campbell--Hausdorff law. An exterior actor can carry
that law, but its multiplication must retain which input supplies
each output generator. We first describe these marked coefficients.

\subsection{Natural actor maps}

In this subsection all actor maps are under and over $J$, with both
comparison paths. Let $W,V$ be finite-dimensional vector spaces
in degree zero, and write $C_W=C_{E_W}$ and $C_V=C_{E_V}$.
Natural operations between $C_W$
and $C_V$ use products, brackets, permutations and linear maps of
exterior coefficients. Each ordinary Lie input is used once.
No additional derivation acts on it. The relative convolution
cochains vanish on purely ordinary inputs and take values in the
augmentation ideal of $U_\star L_{E_V}\to U_\star L_k$.

\begin{theorem}\label{thm:vol2-natural-actor-maps}
The space of natural maps $C_W\to C_V$ is equivalent to the
discrete space of formal matrices
\[
 A(t)\in\operatorname{Hom}(W,V)[[W^*]].
\]
The coefficient of polynomial degree $r-1$ has source actor
weight $r$. Its coordinate contraction is $\Phi_A(t)=A(t)t$.
Composition and diagonal tensor are
\begin{equation}\label{eq:vol2-marked-composition}
 (B\circ A)(t)=B(A(t)t)A(t),\qquad
 (A\oplus B)(t,s)=
 \begin{pmatrix}A(t)&0\\0&B(s)\end{pmatrix}.
\end{equation}
The identity is the constant identity matrix.
\end{theorem}

\begin{proof}
Let $K_W=(E_W^+)^\vee$ and let $\mathfrak q_W$ be its Lie
cobar. For a basis of $W$, a generator $y_I$ indexed by a
nonempty exterior subset has degree $1-|I|$ and differential
\[
 dy_I=\frac12\sum_{\substack{B\sqcup D=I\\B,D\ne\varnothing}}
       (-1)^{|B|-1}\epsilon(B,D)[y_B,y_D].
\]
The sign is obtained by the exterior-label rescaling used in
\eqref{eq:vol2-natural-cobar}. In particular
$dy_{\{i,j\}}=[y_i,y_j]$ for $i<j$.
Adjoin a degree-one Lie generator $z$ with
$dz=[z,z]/2$ by Lie coproduct, writing
$\mathfrak h_W=\operatorname{Lie}(z)\amalg\mathfrak q_W$.

The natural convolution algebra identifies with the summand
of $\widehat{U(\mathfrak h_W\otimes E_V)}$ having positive
source weight and positive target exterior degree. To check
this assertion, polarize the inputs and evaluate on the free
Lie algebra on their names. A Lie word specifies an output
root, and a PBW product specifies a forest of roots, each with
its exterior label. Distinct forests give independent
multilinear operations. Evaluation on the group-like exponential
of the suspended inputs turns convolution into multiplication.
Transposing exterior multiplication gives the displayed Lie
cobar differential. The scalar twist cancels the couplings
between $z$ and the nonempty labels by equivariance; its
remaining equation is $dz=z^2$.

Killing $z$ is a quasi-isomorphism in each source weight.
Before tensoring with $E_V$, the associative envelope is
\[
 (k\langle z\rangle,dz=z^2)*U(\mathfrak q_W).
\]
Contract the first positive $z$ factor in every alternating
word by $h(z^{2j})=z^{2j-1}$, $h(z^{2j-1})=0$, with the
prefix sign. PBW identifies the Lie algebra with a direct
summand of its envelope, so the Lie projection is also a
quasi-isomorphism. Tensoring with $E_V$ and taking envelopes
preserve it: tensor powers and symmetric coinvariants are
exact over $k$.

The projection $\mathfrak q_W\to W^*$ sends $y_i$ to its
namesake in an abelian Lie algebra and kills longer exterior
labels. The resolution
$E_W\otimes k[v_i]$, $dv_i=e_i$, contracts in positive
combined weight $N$ by $N^{-1}\sum_i v_i\partial_{e_i}$.
Dualizing the bar comparison in each finite weight proves
$H(U\mathfrak q_W)=\operatorname{Sym}(W^*)$, with this
projection. The PBW length-one summand again gives the Lie
quasi-isomorphism. Consequently the relative natural algebra
has a surjective dg algebra projection, with weightwise
acyclic kernel, to
\[
 \mathcal B(W,V)=
 \left[\widehat{\operatorname{Sym}}(W^*\otimes E_V)
                         \right]_{E_V^+}.
\]
The subscript denotes the ideal generated by the positive
target exterior labels. Its differential is zero. It has
no terms in degrees at most zero and
\[
 \mathcal B^1(W,V)=
 \widehat{\operatorname{Sym}}(W^*)\otimes W^*\otimes V.
\]
Relative Maurer--Cartan lifting, as in
Theorem~\ref{thm:vol2-natural-cohomology}, has contractible
fibres. In the target a degree-one simplex is a polynomial
function with zero differential, hence constant. This proves
the claimed discreteness.

The surviving $W^*\otimes V$ factor marks the input supplying
the target exterior generator. The other input factors are
unmarked. Exponential polarization supplies their factorials.
In a bar composite, unmarked coefficients substitute
$t\mapsto A(t)t$, while the marked factor is transported by
$A(t)$ and then by $B(A(t)t)$. All joined-root terms vanish
under the projection. This proves the first formula in
\eqref{eq:vol2-marked-composition}; disjoint actor groups give
the block-diagonal formula. The relative bar contraction
transports these calculations to actual maps and paths.
At $m$ positions the source actor weight is at most
$m\dim W$, so every formula is finite there.
\end{proof}

The theorem classifies the specified natural maps. Full
geometric mapping spaces can contain coefficient extensions.
For example, with left differential modules and
$D_x=k\langle x,\partial_x\rangle$, put
$M=k[x,x^{-1}]$ and $\Delta=D_x/D_xx$. The presentation
$M=D_x/D_x(x\partial_x+1)$ gives a two-term free resolution.
On $\partial_x^n\delta\in\Delta$, the resulting Hom
differential $x\partial_x+1$ is multiplication by $-n$.
Thus $\operatorname{Hom}_{D_x}(M,\Delta)=k$ and
$\operatorname{Ext}^1_{D_x}(M,\Delta)=k$. This normal-chart
extension alone does not construct a compatible operation on
the whole collision diagram.

\subsection{Multiplication with a chosen frame}

An associative multiplication on $C_W$ has a matrix
$M(x,y)=(L(x,y)\mid R(x,y))$ and coordinate contraction
$F(x,y)=L(x,y)x+R(x,y)y$. Its unit conditions are
$L(x,0)=I$ and $R(0,y)=I$. By
\eqref{eq:vol2-marked-composition}, associativity is precisely
\begin{equation}\label{eq:vol2-multiplication-blocks}
\begin{aligned}
 L(F(x,y),z)L(x,y)&=L(x,F(y,z)),\\
 L(F(x,y),z)R(x,y)&=R(x,F(y,z))L(y,z),\\
 R(F(x,y),z)&=R(x,F(y,z))R(y,z).
\end{aligned}
\end{equation}
It gives the fibrewise-linear formal group law
\[
 (x,u)(y,v)=(F(x,y),L(x,y)u+R(x,y)v).
\]
Both its zero section and its diagonal are homomorphisms.
Conversely, such a law gives the displayed matrix equations.
Discreteness of the natural mapping spaces supplies a
contractible space of higher monoid coherences.

Let $\mathfrak a$ be the tangent Lie algebra of $F$.
A compatible left-symmetric product on $\mathfrak a$ is
$x\cdot y=\ell_xy$ such that
\[
 \ell_xy-\ell_yx=[x,y],\qquad
 [\ell_x,\ell_y]=\ell_{[x,y]}.
\]
These products describe the fibre laws up to changes of frame
which fix the diagonal and are the identity on the fibre at zero.
Here is the construction, including the frame itself.

The affine Lie map
\[
 x\longmapsto\begin{pmatrix}\ell_x&x\\0&0\end{pmatrix}
\]
integrates formally to a representation $\rho_0$ and a cocycle
$b$, with $db_0=I$ and
$b(F(x,y))=b(x)+\rho_0(x)b(y)$.
In exponential coordinates these are
$\rho_0(x)=e^{\ell_x}$ and
$b(x)=\int_0^1e^{t\ell_x}x\,dt$. The formal matrix
exponential and its Baker--Campbell--Hausdorff identity prove
both equations. Choose an invertible frame $M(x)$ with
$M(0)=I$ and $M(x)x=b(x)$. One always exists: if
$b=\sum_{r\geq1}b_r$ is its homogeneous expansion, take
$M=\sum_{r\geq1}Db_r/r$ and use Euler's identity.
Then
\begin{equation}\label{eq:vol2-graded-frame}
 L(x,y)=M(F(x,y))^{-1}M(x),\qquad
 R(x,y)=M(F(x,y))^{-1}\rho_0(x)M(y)
\end{equation}
satisfies \eqref{eq:vol2-multiplication-blocks} and has
coordinate contraction $F$.

For the converse, the zero section splits the fibre group.
Multiplication $(0,w)(x,0)=(x,L(0,x)w)$ trivializes it.
The frame $M(x)=L(0,x)^{-1}$ transforms the law to a
semidirect product, and the diagonal becomes a cocycle
$b(x)=M(x)x$. Differentiating its identities gives exactly
the two identities for $\ell$. Another frame $M'$ changes
the original law by $S(x)=M'(x)^{-1}M(x)$, with $S(x)x=x$
and $S(0)=I$. Conversely such an isomorphism leaves the
infinitesimal product unchanged. This proves the assertion
about changes of frame, while retaining the distinct
multiplication matrices on a fixed actor.

For instance, on $W=k^2$ with additive $F$, put
$N(t)=t(-t_2,t_1)$ and $S(t)=I+N(t)$. Then $N^2=0$,
$S^{-1}=I-N$ and $S(t)t=t$. The frame gives
$L(x,y)=S(x+y)^{-1}S(x)$ and
$R(x,y)=S(x+y)^{-1}S(y)$. Although its fibre conjugation
is trivial, its multiplication is not the constant additive
matrix: $L((1,0)^T,(0,1)^T)=
\left(\begin{smallmatrix}2&1\\1&1\end{smallmatrix}\right)$.

For a finite-dimensional positively graded Lie algebra
$\mathfrak a=\bigoplus_{j>0}\mathfrak a_j$, let
$G|_{\mathfrak a_j}=j$. The product
\begin{equation}\label{eq:vol2-positive-product}
 \ell_xy=G^{-1}[x,Gy]
\end{equation}
is compatible: the derivation identity for $G$ proves its
antisymmetrization, and
$\ell_x=G^{-1}\operatorname{ad}_xG$ proves its representation
identity. In Baker--Campbell--Hausdorff coordinates choose
$M(x)=\int_0^1e^{t\ell_x}\,dt$. Every $\ell_x$ raises
positive grading. The exponentials, $M^{-1}$ and $F$ are
therefore polynomials. Formula~\eqref{eq:vol2-graded-frame}
constructs a natural coherent monoid on $C_{\mathfrak a}$.

\subsection{Composition of action classes}

For a natural action with full twisting cochain $\tau$,
adjoin independent central degree-zero test elements $z_i$
annihilated by the derivations and put $\zeta=\sum_i e_i z_i$.
The operator
\[
 T_\tau(z)x=\tau\bigl(sx\exp(s\zeta)\bigr)
\]
has constant term $x$. It defines a series in
$\prod_{q\geq0}\operatorname{Sym}^q(W^*)\otimes U(\mathfrak d)$.
The twisting equation on two ordinary inputs gives
$T_\tau[x,y]=[T_\tau x,T_\tau y]$. Evaluate on the free
Lie algebra on the named free modules
$U(\mathfrak d)x\oplus U(\mathfrak d)y$. The map
$a\otimes b\mapsto[ax,by]$ is injective, by projection to
the corresponding ordered word in its tensor envelope.
The preceding identity therefore says
$\Delta T_\tau=T_\tau\otimes T_\tau$.

Its logarithm is primitive, hence belongs to the completed
$\mathfrak d$-valued series. PBW projection to degree one
identifies this logarithm with the action class $\phi$ of
\eqref{eq:vol2-natural-classes}: powers of the one element
$\phi$ are already symmetric powers. Thus $T_\tau=e^\phi$.
The actual tensor extension of an inner action carries
$D$ to $D\otimes1$, by
Lemma~\ref{lem:vol2-strong-boundary-tensor}. It fixes the
outer central test parameters. Composition with an outer
action of class $\phi$ and an inner action of class $\psi$
therefore has operator $e^\phi e^\psi$ and class
\begin{equation}\label{eq:vol2-action-BCH}
 \operatorname{BCH}(\phi,\psi)=\log(e^\phi e^\psi).
\end{equation}
The bracket is $[fD,gB]=fg[D,B]$ and raises actor weight.
For example $[D,Z_3]=6\kappa T$ and $T$ is central, so
\[
 \operatorname{BCH}(pD,qZ_3)=pD+qZ_3+3\kappa pqT.
\]
Reversing the action order reverses the last sign.

Pullback along a natural actor map of matrix $A(t)$ sends
an action class $\psi$ and a scalar loop $\sigma$ to
\begin{equation}\label{eq:vol2-class-substitution}
 \psi(A(t)t),\qquad \sigma(A(t)t).
\end{equation}
Indeed removing the marked generator gives $q_j\mapsto
\sum_iA_{ji}(t)t_i$. Precomposition on the natural bar
preserves products; its scalar derivation generators are
unchanged. This proves the formula on the whole cohomology
algebra. Postcomposition by a fixed outer action preserves
the scalar loop of the inner action: the outer operator
fixes its central, derivation-inert test coefficients.

For a finite-dimensional nilpotent Lie algebra $\mathfrak a$,
let $F(x,y)=\operatorname{BCH}(x,y)$. Its normalized formal
$q$-cochains with trivial scalar coefficients are formal series
$c(x_1,\ldots,x_q)$ which vanish whenever any argument is zero.
The completion is in total coordinate degree. Their differential is
\begin{equation}\label{eq:vol2-formal-group-differential}
\begin{aligned}
 (\delta c)(x_1,\ldots,x_{q+1})
 ={}&c(x_2,\ldots,x_{q+1})\\
 &+\sum_{i=1}^{q}(-1)^i
 c(x_1,\ldots,F(x_i,x_{i+1}),\ldots,x_{q+1})\\
 &+(-1)^{q+1}c(x_1,\ldots,x_q).
\end{aligned}
\end{equation}
Associativity and the unit give $\delta^2=0$. Degree zero is
$k$ with zero differential. Denote the cohomology by
$H^\bullet_{\mathrm{for}}(F;k)$. These are functions near the
formal identity, rather than cochains on the discrete group of
$k$-points.

\begin{theorem}\label{thm:vol2-graded-multiplicative-action}
Let $\mathfrak a$ be finite-dimensional and positively graded,
and let $\rho:\mathfrak a\to\mathfrak d_V$ be a Lie map.
There is a natural action of $C_{\mathfrak a}$ on the
Heisenberg boundary with class $x\mapsto\rho(x)$, compatible
with the monoid \eqref{eq:vol2-graded-frame}. Compatibility
includes every unit and associativity homotopy. Every natural
action in this class admits such a structure.
\end{theorem}

\begin{proof}
First work universally with derivation algebra $\mathfrak a$
and class $\theta(x)=x$. The natural convolution projection
constructs an action in this class. Its two binary composites
have the same class, since
\[
 \theta(F(x,y))=\operatorname{BCH}(\theta(x),\theta(y)).
\]
Choose a comparison path $h(x,y)$ relative to both unit axes.
The fully mixed summand is a direct summand of the acyclic
kernel of that projection, so relative Maurer--Cartan lifting
retains these unit restrictions.

Let $K_n$ be the associahedron of bracketings of $n$ ordered
inputs; its dimension is $n-2$. For a module the last input
is $C$ and the others are actors. The path $h$ occupies
$K_3$. The boundary of $K_4$ is now a scalar loop
$o(x,y,z)$. It vanishes on each unit face. Replacing $h$ by
a normalized scalar loop $\sigma(x,y)$ changes it by
\[
 (\delta\sigma)(x,y,z)=\sigma(y,z)-\sigma(F(x,y),z)
                +\sigma(x,F(y,z))-\sigma(x,y).
\]
These transports are \eqref{eq:vol2-class-substitution} and
the unchanged inner scalar loop. The boundary in the next
associahedron gives $\delta o=0$: each transported edge
occurs twice with opposite orientation, the square faces
express functoriality, and the pure actor face is supplied
by the monoid coherence. Hence $o$ defines a class in the
normalized formal-group cohomology $H^3_{\mathrm{for}}(F;k)$,
with trivial scalar coefficients.

This cohomology equals $H^3(\mathfrak a;k)$ equivariantly
for the grading torus. To see this, finite-order distributions
at the identity form $U(\mathfrak a)$, with PBW associated
graded. Formal functions are its full linear dual. The
normalized group cochain complex is the dual bar resolution
of the scalar module. The Chevalley resolution computes the
same derived Hom; its exactness follows from the polynomial
Koszul resolution under PBW. Antisymmetrization into the bar
dualizes to differentiation at the identity. This gives the
asserted comparison in every degree; see also
\cite{Lechner2012}*{Theorem~4.2 and Corollary~4.3}.

The obstruction class is independent of the chosen lift and
binary path. The grading torus preserves $\theta$,
\eqref{eq:vol2-positive-product}, its frame and the natural
bar comparisons, so this class is invariant. But every
weight of $\Lambda^q\mathfrak a^*$ is strictly negative
for $q>0$. Thus its invariant cohomology vanishes. Choose
a normalized $\sigma$ with $\delta\sigma=-o$ and correct
the binary path. The resulting boundary fills.

All later attaching spheres have dimension at least two.
Their action spaces have zero homotopy groups in these
dimensions by Theorem~\ref{thm:vol2-natural-cohomology}.
Successive module associahedra therefore fill, with their
specified unit faces. Positive grading has finite pieces,
so the formal corrections respect the completion. Evaluation
is finite at every configuration. Finally applying $U(\rho)$
to the universal operations preserves all their equations,
paths and units. A natural homotopy to another action in
the same class transports the constructed module structure.
\end{proof}

The binary path in this proof is chosen and corrected. The
theorem does not require it to equal an independently prescribed
path. Nor does the invariant vanishing assert that all third
cohomology vanishes. For $[a,b]=c$ with weights $1,1,2$,
the Chevalley differential is $dc^*=-a^*\wedge b^*$.
All degree-two cochains are closed, so
$[a^*\wedge b^*\wedge c^*]$ is a nonzero top class,
of grading weight $-4$.

\begin{corollary}\label{cor:vol2-nilpotent-recipient}
A finite family of natural exterior-actor actions whose classes
take values in one finite-dimensional nilpotent Lie subalgebra
of $\mathfrak d_V$ has a common coherent multiplicative
recipient on a finite-dimensional exterior actor. Their
classes may be formal series.
\end{corollary}
\begin{proof}
If the subalgebra is zero, use $J$. Otherwise choose its
generators and nilpotence class $c$. The free nilpotent Lie
algebra $\mathfrak a$ of class $c$ on these generators is
finite-dimensional, positively graded by bracket length, and
surjects onto the subalgebra. A linear section lifts each
formal action class to a series $\Phi(t)\in\mathfrak a[[W^*]]$
with zero constant term. If $\Phi=\sum_{r\geq1}\Phi_r$,
the matrix
\[
 A_\Phi(t)=\sum_{r\geq1}\frac1rD\Phi_r(t)
\]
satisfies $A_\Phi(t)t=\Phi(t)$. Theorem~\ref{thm:vol2-natural-actor-maps}
realizes it. Pullback of the action in
Theorem~\ref{thm:vol2-graded-multiplicative-action} has the
given class by \eqref{eq:vol2-class-substitution}, so equality
of natural classes supplies the required pointed action path.
\end{proof}

The incoming maps in this corollary need not preserve
multiplication. For additive one-dimensional actors, the
map $\Phi(t)=t^2$ has matrix $A(t)=t$, but
$\Phi(s+t)-\Phi(s)-\Phi(t)=2st$.

There is also a precise limitation on paths within this
natural augmented category. Fix an actor map $f:C_W\to C_V$
and an action path $p$ over it. For a nonzero scalar loop
$\sigma\in\mathfrak m_W$, the paths $p$ and $p\cdot\sigma$
define morphisms in different components. Their difference
persists after postcomposition through any natural exterior
actor. Indeed the path fibre has components forming a torsor
under $\mathfrak m_W$, and the actor mapping spaces are
discrete, so no actor loop identifies these components.
Postcomposition concatenates both paths with the same
pulled-back path. This obstruction concerns the specified
natural maps preserving augmentation. A geometric recipient
can have additional actor loops, whose images must be included
in the full question.

\section{The full density object}
\label{sec:vol2-action-density}

Fix an arbitrary factorizing boundary $C\in\vF$.
Consider all pointed factorizing actions on $C$, with their geometric
mapping spaces. Their ambient colimit is an action, but factorization
requires comparing one common actor index with independent indices
in the separated blocks. This comparison has a section. Its kernel
is the remaining obstruction to pointed representability.

\subsection{Ambient endomorphisms and a small index}

The enhanced coefficient category $\mathscr C$ is presentable.
Indeed its differential-module categories are presentable, its
extraordinary pullbacks preserve colimits, and the small limit
defining \eqref{eq:vol2-action-category} is a limit of presentable
categories and colimit-preserving functors. Colimits are computed
on the coefficients. The same assertions hold for its positive
subcategory $\mathscr C_+$.

\begin{lemma}\label{lem:vol2-ambient-endomorphisms}
The category $\vM$ is presentable. Its colimits are computed
componentwise on reduced coefficients, and $\vot$ preserves
them separately. Thus $-\vot C$ has a right adjoint and
\[
 \vEnd=[C,C]_{\vot}
\]
is an ambient associative endomorphism monoid, with unit
$\eta:J\to\vEnd$ and evaluation $\vEnd\vot C\to C$.
Maps of monoids to $\vEnd$ classify coherent actions on $C$.
\end{lemma}
\begin{proof}
For positive coefficient families, the $n$th convolution power
vanishes on configurations with fewer than $n$ positions.
Bar and cobar therefore use finite filtrations at each
configuration. Normalized bar terms are $a^{\star n}[n]$
for the positive augmentation ideal $a$, and normalized cobar
terms are $c^{\star n}[-n]$ for the positive reduced coalgebra
$c$. Simplicial realization and cosimplicial totalization
give their coherent structures. Collision pullbacks preserve
these finite normalized constructions.

In $\Omega BA$, flatten a term to its $n$ original
$a$-factors and filter increasingly by $n$. Each of the
$n-1$ gaps is either inside a bar block or between cobar
factors. On the associated graded, multiplication of
original inputs disappears and the differential changes
one inside gap to an outside gap. After the suspension
signs, this is the tensor product of $n-1$ two-term identity
complexes, tensored with the coefficient word. Opposite
signs around each square give the usual tensor differential.
For $n>1$ a contraction in its first gap contracts the
product. For $n=1$ the surviving coefficient is $a$;
the empty term is $K$. The comparison is the identity
on these surviving terms.

For $B\Omega C$, use the decreasing filtration by original
$c$-factors. Coproduct splitting increases that filtration.
Its associated graded has the same gap complexes with
the direction of the cuts reversed, again leaving only
$K$ and $c$. Both filtrations end at $|I|$ on configuration
$I$, so their comparisons are equivalences. Higher coherent
products lower the first filtration and higher coproducts
raise the second; the same associated graded proof applies.
The bar and cobar adjunction retains the triangle coherences.
Thus connected coalgebras and positive augmented algebras
are equivalent as enhanced categories.

Such algebras are nonunital associative algebras in
$\mathscr C_+$. The tensor-word monad is accessible and
preserves sifted colimits, giving a presentable algebra
category \cite{LurieHA}*{Corollary~3.2.3.5}. This proves
presentability of $\vM$.
On the coalgebra side, the reduced objects are nonunital
algebras in $\mathscr C_+^{\mathrm{op}}$. Limits of algebras
are computed on their underlying objects: the operations
and their coherences are obtained by composing with each
projection. Reversing arrows proves the assertion about
coalgebra colimits.

The reduced part of a diagonal tensor is the diagonal tensor
of the reduced parts. Its pullback and external product
preserve colimits. The adjoint functor theorem
\cite{LurieHTT}*{Corollary~5.5.2.9} gives the asserted
right adjoint. Composing two evaluations and
adjointing defines multiplication on $\vEnd$. Associativity
and units are those of composition, including their higher
coherences. Adjunction transports precisely the monoid-map
identities to the coherent action identities.
\end{proof}

This lemma constructs an ambient object; it imposes no
factorization on $\vEnd$. Write
\[
 \mathscr V_C=(\vM_{/\vEnd})_{(J,\eta)/}.
\]
Then $\vAct$ is its full subcategory with factorizing
underlying object. An action $a$ will be written
$(F_a,u_a,\beta_a)$, where $\beta_a:F_a\to\vEnd$ includes
the path $\beta_a u_a\simeq\eta$.

The category $\vAct$ is accessible. For each finite ordered
partition, factorization is the condition that a natural
arrow between accessible coefficient functors be an
equivalence. The category of equivalences in an arrow
category is accessible. Pulling it back for the small set
of partitions, and then adding the pointing and slice,
gives the assertion. Choose one sufficiently large regular
cardinal $\lambda$ for these functors. The inclusion
$\vAct\to\mathscr V_C$ preserves $\lambda$-filtered
presentations: on a filtered index the diagonal to any
finite power is cofinal, so finite external products of
coefficients preserve the factorization comparisons.

Choose a small skeleton $\mathscr A_0$ of the
$\lambda$-compact generators, and let $\mathscr T_C$ be
the small full monoidal subcategory generated by it and
the unit action. Fullness retains all mapping spaces
between the selected objects. Define
\begin{equation}\label{eq:vol2-density-object}
 L_C=\operatorname*{colim}_{a\in\mathscr T_C}F_a
            \quad\text{in }\mathscr V_C,
 \qquad \lambda_a:F_a\to L_C.
\end{equation}
The unit action is initial, so this index has contractible
nerve. Its underlying colimit is consequently the colimit
in $\vM$, with the specified point and map
$\beta_L:L_C\to\vEnd$.

\begin{proposition}\label{prop:vol2-density-cocone}
The density object is a unital ambient monoid over $\vEnd$.
Its cocone extends naturally to all of $\vAct$, retaining
every mapping-space coherence, and is multiplicative for
the ordered action product $\circledast$.
\end{proposition}
\begin{proof}
Separate preservation of colimits defines multiplication by
\[
 m_L(\lambda_a\vot\lambda_b)=\lambda_{a\circledast b}.
\]
The iterated products of the index supply associativity;
its initial unit supplies both unit laws. The maps $\beta_a$
make $\beta_L$ multiplicative.

For any action $a$, the filtered slice $\mathscr A_0/a$
presents $F_a$. Thus the inclusion into $\mathscr V_C$
is the left Kan extension of its restriction to
$\mathscr A_0$. The same holds after restricting it to
$\mathscr T_C$, because that subcategory is full and
has the same slices at its objects. The universal property
of left Kan extension therefore extends the cocone,
including its higher naturality, to all actions. Passing
to compact filtered presentations in each variable extends
the multiplication identity and its coherences as well.
\end{proof}

\subsection{The section and its kernel}

For an ordered partition $P=(I_1|\cdots|I_r)$, set
\[
 X_P(F)=j_P^!F_I,\qquad
 Y_P(F)=\left(\boxtimes_{j=1}^rF_{I_j}\right)|_{U_P}.
\]
These objects lie in the stable differential-module category
on $U_P$. Internal block collisions remain in that open.
The coproduct is $\delta_{P,F}:X_P(F)\to Y_P(F)$.
For factorizing $F$ it is an equivalence; its inverse is
taken naturally in the functor category.

\begin{proposition}\label{prop:vol2-density-section}
The map $\delta_{P,L_C}$ has a section $s_P$. If
\begin{equation}\label{eq:vol2-density-kernel}
 K_{C,P}=\operatorname{fib}(\delta_{P,L_C}),\qquad
 p_P=1-s_P\delta_{P,L_C},
\end{equation}
then
\[
 X_P(L_C)\simeq Y_P(L_C)\oplus K_{C,P},\qquad
 \operatorname{cofib}(\delta_{P,L_C})\simeq K_{C,P}[1].
\]
For $h=a^{\circledast r}$ and its unit insertions $v_j:a\to h$,
put
\begin{equation}\label{eq:vol2-density-generator}
 r_{a,P}=X_P(v_1)-\delta_{P,F_h}^{-1}
             \left(\boxtimes_j(v_j)_{I_j}\right)\delta_{P,F_a}.
\end{equation}
Then $p_PX_P(\lambda_a)=X_P(\lambda_h)r_{a,P}$,
with its natural homotopy.
\end{proposition}
\begin{proof}
For independently indexed actions $a_1,\ldots,a_r$, let
$h=a_1\circledast\cdots\circledast a_r$ and let
$v_j:a_j\to h$ insert the other units. On this tuple define
\begin{equation}\label{eq:vol2-section-formula}
 s_P\big|_{(a_j)}=X_P(\lambda_h)\delta_{P,F_h}^{-1}
                   \left(\boxtimes_j(v_j)_{I_j}\right).
\end{equation}
This is a natural cocone on $\mathscr T_C^r$. The external
product preserves its separate colimits, giving a map
$Y_P(L_C)\to X_P(L_C)$. Naturality of $\delta$ makes its
composite with $\delta_{P,L_C}$ equal to the target cocone,
since $\lambda_hv_j\simeq\lambda_{a_j}$. Hence $\delta s\simeq1$.
The split fibre sequence gives the two decompositions.
Using the tuple $(a,\ldots,a)$ in
\eqref{eq:vol2-section-formula} gives
\eqref{eq:vol2-density-generator}. Compact presentations
extend the formula to every action. Internal collisions
and refinements use the same iterated coproduct coherences.
\end{proof}

\begin{theorem}\label{thm:vol2-density-criterion}
The following are equivalent: the pointed-action functor
$\Phi_C$ is representable; $L_C$ factorizes; and all
$K_{C,P}$ vanish. It suffices to test every ordered binary
block partition, with its full block coefficients.
When these conditions hold, $L_C$ is also the final
factorizing monoid action.
\end{theorem}
\begin{proof}
Binary factorization iterated inside blocks gives all
partitions. The split fibre sequence identifies
factorization with kernel vanishing.

If $L_C$ factorizes, it belongs to $\vAct$. Naturality of
the extended cocone gives
$\lambda_{L_C}\lambda_a\simeq\lambda_a$ on the small
index, hence $\lambda_{L_C}\simeq1$ by its colimit
property. Naturality on the whole space
$\vMap_{\vAct}(a,L_C)$ now homotopes its identity map
to the constant map with value $\lambda_a$. This space
is contractible, so $L_C$ is final and represents $\Phi_C$.

Conversely a final pointed action $Z$ gives a cocone
and thus a map $\rho:L_C\to Z$. The extended cocone
gives $\lambda_Z:Z\to L_C$. Naturality on the full
small index gives $\lambda_Z\rho\simeq1$. The reverse
composite is the identity by finality of $Z$. Thus
$L_C\simeq Z$ and factorizes. Finally all operation
spaces from finite products to a final object are
contractible. They give its unique coherent monoid
structure and a contractible space of monoid-action
maps into it from every factorizing monoid actor.
\end{proof}

This proof uses finality among all pointed actions.
Finality among multiplicative actions alone does not
supply the cocone from arbitrary pointed actors.

\subsection{Specified monoid extensions}

A monoid extension is $(a,M,f)$ with $M\in\vMon$ and
a pointed-action morphism $f:a\to\mathsf U M$, where
$\mathsf U$ forgets multiplication. Its morphisms include
an action map on $a$, a monoid-action map on $M$, and
a specified homotopy making their square commute.
Denote this category, with its higher coherences, by
$\mathscr E_C=\vAct^{\Delta^1}\times_{\vAct}\vMon$,
using target evaluation on the arrow category and $\mathsf U$
on the monoid category.

\begin{theorem}\label{thm:vol2-extension-cancellation}
On $\mathscr E_C$ there is a natural nullhomotopy
$p_PX_P(\lambda_a)\simeq0$. It retains the vacuum paths,
ordered coproducts and all internal collision comparisons.
\end{theorem}
\begin{proof}
For $h=a^{\circledast r}$ set
$q=\mu_M^{(r)}f^{\vot r}:h\to\mathsf U M$.
The monoid-action identity makes this a pointed-action
map in the prescribed order. Its unit paths give
$qv_j\simeq f$ for every $j$. Naturality of factorization
then gives
\[
 X_P(q)\delta_{P,F_h}^{-1}
       \left(\boxtimes_j(v_j)_{I_j}\right)\delta_{P,F_a}
 \simeq\delta_{P,M}^{-1}Y_P(f)\delta_{P,F_a}
 \simeq X_P(f).
\]
The other term of $X_P(q)r_{a,P}$ is also $X_P(f)$.
Thus $X_P(q)r_{a,P}\simeq0$. The density cocone has
$\lambda_Mq\simeq\lambda_h$, proving the assertion.

These are equations of natural transformations on
$\mathscr E_C$: the unit paths come from its coherent
monoid structures, and the factorization inverses are
natural inverses there. Their composite is therefore
a homotopy on every simplex, not a collection of
unrelated objectwise homotopies. Restriction to either
ordered collision branch and to internal block collisions
preserves the calculation. The projector's domain remains
$U_P$; no extension through cross-block collisions is asserted.
\end{proof}

The full subcategory of pointed actions which happen to
admit multiplicative structures also has maps that do
not preserve those structures. It is not the extension
category just defined. This distinction matters when
passing the nullhomotopy to a colimit.

For a small diagram $\mathscr E\to\mathscr E_C$, define
\[
 A_P^{\mathscr E}=\operatorname*{colim}_{e\in\mathscr E}
                       X_P(F_{a_e}),\qquad
 \chi_P:A_P^{\mathscr E}\longrightarrow X_P(L_C).
\]
The map is the extended density cocone. Theorem
\ref{thm:vol2-extension-cancellation} gives a nullhomotopy
of $p_P\chi_P$, so
\begin{equation}\label{eq:vol2-relative-kernel}
 \operatorname{cofib}(\chi_P)\simeq
 \operatorname{cofib}(\delta_{P,L_C}\chi_P)\oplus K_{C,P}.
\end{equation}
Indeed, put $A=A_P^{\mathscr E}$, $X=X_P(L_C)$,
$Y=Y_P(L_C)$ and $K=K_{C,P}$. In the splitting $X=Y\oplus K$, write
$\chi=(\chi_Y,\chi_K)$ and represent the specified
homotopy by $\chi_K=d_Kh+hd_A$. Our cone convention is
\[
\begin{gathered}
 \operatorname{Cone}(\chi)^n=Y^n\oplus K^n\oplus A^{n+1},\\
 d(y,k,a)=(d_Yy+\chi_Ya,d_Kk+\chi_Ka,-d_Aa).
\end{gathered}
\]
The cochain isomorphism is the shear
\begin{equation}\label{eq:vol2-cone-shear}
 (y,k,a)\longmapsto((y,a),k+h(a)).
\end{equation}
Its last differential is
$d_Kk+\chi_Ka-hd_Aa=d_K(k+ha)$, and its inverse subtracts
$ha$. In particular the retraction to the kernel includes
the chosen homotopy.

For example take $A=(k\xrightarrow{5}k)$,
$Y=(k\xrightarrow{2}k)$ and
$K'=(k\xrightarrow{(3,0)}k^2)$ in degrees $-1,0$.
Let $\chi^{-1}=(10,35)$ and $\chi^0=(4,21,0)$.
Its kernel component has homotopy $h^0=7$.
Projection to $k$ alone is not a chain map on the cone,
because of the term $21a$; the corrected projection is
$k+7a$. The nonzero class in $H^0(K')$ remains a retract.

\begin{corollary}\label{cor:vol2-extension-density}
Pointed representability holds if and only if a small
diagram of specified monoid extensions makes $\chi_P$
an equivalence for every ordered binary partition.
Equivalently, one may use a different such diagram for
each partition.
\end{corollary}
\begin{proof}
If the maps are equivalences, their cofibres vanish.
Equation~\eqref{eq:vol2-relative-kernel} makes every
$K_{C,P}$ a retract of zero, and
Theorem~\ref{thm:vol2-density-criterion} applies.
Conversely a final pointed action $Z$ has its coherent
monoid structure. Its final cocone supplies a functor
$\mathscr T_C\to\mathscr E_C$ with constant target $Z$.
The resulting $\chi_P$ is the original density colimit
on every block-disjoint open, hence an equivalence.
\end{proof}

Nonzero relative cofibres alone do not obstruct this
criterion: they may come from the first summand in
\eqref{eq:vol2-relative-kernel}. Nor is a defect before
the density colimit a class of its kernel. Take an exterior
actor with a multiplicative action, such as $(C_E,\alpha_D)$.
At two separated positions let $c_{L,i},c_{R,i}$ be the
degree-one generators in the exterior state factors of
two copies, ordered first by copy and then by position.
They give the nonzero expression
\[
 c_{L,1}c_{L,2}-c_{L,1}c_{R,2}.
\]
Multiplication sends both monomials to $c_1c_2$ and
annihilates their difference. When this multiplication
is a pointed-action map, the full density cocone
annihilates the same difference. Its action paths are
part of this conclusion.

\section{A pair surviving every pointed factorizing action}
\label{sec:vol2-pointed-sphere}

A common recipient of objects need not equalize their maps.
For the polynomial Heisenberg boundary, two prescribed maps remain
distinct after every receiving pointed factorizing action.
The vacuum imposes the restriction responsible for this obstruction:
a degree-zero operation coming from one root of a factorizing actor
must commute with the boundary-module action.

Retain $C=C_{V_\kappa}$ with $\kappa\ne0$, and put
$\mathfrak l=L_k$, $\mathfrak g=L_{V_\kappa}$ and
$A=\vU\mathfrak g$.
Write $\ell=\mathfrak l_{\{1\}}=\omega_X$ and
$G=\mathfrak g_{\{1\}}$ for the right differential modules on $X$.
Thus $J_{\{1\}}=\ell[1]$ and $C_{\{1\}}=G[1]$.
All state spaces and tensors at a fixed configuration are algebraic.

For a coalgebra $Z$, let $\overline Z=\operatorname{cofib}(K\to Z)$.
With $t=[-1]$ of degree one, the normalized cobar is
\[
 \Omega Z=K\oplus\bigoplus_{r\ge1}(t\overline Z)^{\star r},
 \qquad
 d(tc)=-t(d_Zc)+\sum_{\overline\Delta c=\sum c'\otimes c''}
                  (-1)^{|c'|+1}(tc')(tc'').
\]
Its product is concatenation. The sums are finite on every
configuration. These are the cobar and suspension conventions used
in the bar--cobar comparison above.

\subsection{An antisymmetric actor map}

Put $F_1=C_{k}$ and $F_2=C_{k^2}$ in the notation
$C_W=C_{\Lambda(W[-1])}$ of
Section~\ref{sec:vol2-nilpotent-recipients}.
Denote their exterior generators by $\varepsilon$ and
$\varepsilon_1,\varepsilon_2$, respectively.
The subscript $k$ here denotes the one-dimensional actor space;
the vacuum coalgebra remains $J=C_0$.

\begin{lemma}\label{lem:vol2-antisymmetric-actor}
There is a natural actor map $f:F_2\to F_1$, under and over $J$,
whose marked matrix and coordinate contraction are
\[
 A_f(t_1,t_2)=\tfrac12(t_2,-t_1),\qquad \Phi_f=0.
\]
Its pure linear components vanish. In a cobar representative its
pure mixed binary component is
\begin{equation}\label{eq:vol2-antisymmetric-lift}
 \tfrac12(a_mM_n-a_nM_m),\qquad
 a_m=t(s(\varepsilon m)),\quad M_n=t(sn),\quad m,n\in\mathfrak l.
\end{equation}
Every positive-weight natural action pulled back along $f$ is
naturally homotopic to the trivial action.
\end{lemma}

\begin{proof}
The matrix classification in
Theorem~\ref{thm:vol2-natural-actor-maps} already constructs an
actual map with this matrix in every weight.
An explicit composite fixes its binary representative.
Take two auxiliary odd generators $z_1,z_2$, and let $D_z$ be the
even derivation with $D_zz_1=z_2$ and $D_zz_2=0$.
The tree action extends to these graded coefficients by inserting
the Koszul sign in each permutation of the actor coefficients.

Here is a check of that extension. A nonzero term in the twisting
equation has exactly two ordinary inputs: a source bracket joins
them or joins one to an actor input, and a product uses one in each
factor. For each homogeneous coefficient input of nonnegative
degree $d$, multiply it by a product of $d$ distinct auxiliary
odd variables of degree $-1$. Use disjoint variables for different
inputs. The modified inputs have degree zero.
Extend $D_z$ trivially over these variables and apply the
fixed-point proof of Theorem~\ref{thm:vol2-profiles}.
Extracting their ordered product gives the graded twisting
identity. Permuting actor inputs gives their Koszul sign; moving
the ordinary suspension uses its degree one lower.
Thus the graded action is an actual coefficient map, with the
same finite-configuration bound.

On one $z$ input and $w$ pure auxiliary actor inputs its value is
$(D_z^wz)n\,\operatorname{pbw}(m_1\cdots m_w)$.
Act on $z_1,z_2$ by the auxiliary generator, then project
$z_1\mapsto0$, $z_2\mapsto\varepsilon$.
The resulting map $b:C_{k^3}\to F_1$ has marked matrix
\[
 B(y_1,y_2,y_3)=(y_2,0,1).
\]
Choose the inner map $a:F_2\to C_{k^3}$ with matrix
\[
 A_a(t)=
 \begin{pmatrix}
  1&0\\0&1\\-t_2/2&-t_1/2
 \end{pmatrix}.
\]
Its coordinate contraction is
$(t_1,t_2,-t_1t_2)$.
The composition law \eqref{eq:vol2-marked-composition}
gives $A_{ba}=(t_2,-t_1)/2$ and $\Phi_{ba}=0$.

For the binary calculation, relabel the intermediate generators
as $z_1,\eta,z_2$, with $\eta$ the auxiliary actor.
Write $a_z(m)=t(s(zm))$.
A lift of the mixed inner component is
\[
 L_{12}(m,n)=t\bigl(s(z_1m)s(\eta n)\bigr)
 -\tfrac12\bigl(a_{z_2}(m)M_n+a_{z_2}(n)M_m\bigr).
\]
The source input $s(\varepsilon_1\varepsilon_2p)$ lifts to
$-t(s(\eta z_1p))$. The cobar differential gives
\[
 dL_{12}
 =-t\bigl(s(\eta z_1[m,n])\bigr)
       -a_{z_1}(m)a_\eta(n)-a_\eta(n)a_{z_1}(m).
\]
This is the image of the source differential.
The primitive product correction is closed.
The relative bar contraction extends these compatible components
and keeps their binary values.
Applying $b$ sends the first term of $L_{12}$ to
$(\varepsilon m)n$ and its correction to
$-\bigl((\varepsilon m)n+(\varepsilon n)m\bigr)/2$.
This proves \eqref{eq:vol2-antisymmetric-lift}.
Both pure diagonal pairs and both pure linear components vanish.

Finally, \eqref{eq:vol2-class-substitution} identifies pullback of a
natural action class $\psi$ with $\psi(\Phi_f)$.
Its constant term is zero, so this is zero in every weight.
The Maurer--Cartan equivalence of
Theorem~\ref{thm:vol2-natural-cohomology} supplies the homotopy.
Evaluation and the relative bar comparison preserve its vacuum
path and all collision coefficients.
\end{proof}

Let $f_0:F_2\to J\to F_1$ be the augmented map.
For $D=\partial_{x_1}$ the action $\alpha_D$ has natural class $pD$.
Choose the all-weight homotopy from the preceding lemma, with the
constant homotopy for $f_0$. They give two actual morphisms
\begin{equation}\label{eq:vol2-pointed-sphere}
 (f,p_f),(f_0,p_0):
 (F_2,\mathrm{triv})\rightrightarrows(F_1,\alpha_D)
 \quad\text{in }\vAct.
\end{equation}
Equality of their first action classes alone would not have supplied
these paths.

\subsection{The operation at one root}

For a right differential module $N$ on $X$, let
$\iota_{\mathrm{diag}}N$ be its positive coefficient family, with
the direct image of $N$ on every small-diagonal branch.
For any positive family $H$ there is an equivalence
\begin{equation}\label{eq:vol2-main-diagonal-adjunction}
 \vRH_{\mathscr C}(\iota_{\mathrm{diag}}N,H)
       \simeq\vRH_X(N,H_{\{1\}}).
\end{equation}
Indeed the union of small-diagonal branches is closed on each
$\mathscr R_I$. Local cohomology there is the sum
$\bigoplus_\sigma\delta_{\sigma,*}\delta_\sigma^!H_I$.
Collision pullback preserves these unions and their local
cohomology. The resulting coefficient family is
$\iota_{\mathrm{diag}}H_{\{1\}}$, by the collision identifications.
Its counit induces an equivalence on maps from a family supported
on the branches. Restriction to one position is inverse to
$\iota_{\mathrm{diag}}$ on such supported families, proving the
formula.

Fix a pointed factorizing action $(Z,u,\gamma)$ and set
\[
 B=\Omega Z,\qquad
 \mathcal T_C=(-)\vot C,\qquad
 \Phi:\mathcal T_C(B)\longrightarrow A.
\]
Here $\Phi$ is the cobar of $\gamma$, followed by the
Chevalley--enveloping comparison.
Lemma~\ref{lem:vol2-strong-boundary-tensor} identifies
$\mathcal T_C(\Omega Z)$ with $\Omega(Z\vot C)$.
The pointing gives $\iota:\mathfrak l\to B$, and its action path
identifies $\Phi\mathcal T_C(\iota)$ with $\mathfrak g\to A$.

For an affine open $Y\subset X$ and a class
$a\in H^0\vRH_Y(\ell,B_{\{1\}})$, use
\eqref{eq:vol2-main-diagonal-adjunction} to regard $a$ as a
one-root coefficient map. Define
\[
 \Psi(a)=\Phi\mathcal T_C(a).
\]
Since $A_{\{1\}}=G$, the same adjunction identifies $\Psi(a)$
with an endomorphism of the full family $\mathfrak g$ over $Y$.
No cohomological bound on $B$ is imposed.

\begin{lemma}\label{lem:vol2-cobar-product-exchange}
On the open $U$ of two distinct roots, both ordered products
\[
 \mu_{12},\mu_{21}:
 (B_{\{1\}}\boxtimes B_{\{2\}})|_U\longrightarrow B_{\{1,2\}}|_U
\]
are equivalences, with the source order understood in each case.
For the external symmetry $\sigma$, put
$R_B=\mu_{12}^{-1}\mu_{21}\sigma$.
Every algebra map $F:B\to B'$ between such algebras satisfies
\[
 (F_1\boxtimes F_1)R_B=R_{B'}(F_1\boxtimes F_1).
\]
For $A=\vU\mathfrak g$, the exchange is the identity.
\end{lemma}

\begin{proof}
On two roots the reduced cobar has one generator
$t\overline Z_2$ and two product summands
$t\overline Z_1\boxtimes t\overline Z_1$.
Its two differential components are the coproducts of $Z$,
with the displayed suspension signs.
Both are equivalences on $U$.
Quotienting by either product summand leaves the cone of the other
coproduct equivalence. Thus both product inclusions are equivalences,
even for unbounded complexes.
The exchange identity follows after composition with $\mu'_{12}$
from preservation of the two products.
In $A$, the difference of the two generator products is the
chiral bracket, supported on the diagonal. Hence its open exchange
is the identity.
\end{proof}

Write $\mu_{\mathfrak g}:\mathfrak g\star\mathfrak g\to\mathfrak g$
for the coefficient bracket.
An endomorphism $M$ of the module $\mathfrak g$ obeys
$M\mu_{\mathfrak g}=\mu_{\mathfrak g}(\vid\star M)$.

\begin{lemma}\label{lem:vol2-vacuum-module-map}
Every one-root operation $M=\Psi(a)$ satisfies
\begin{equation}\label{eq:vol2-vacuum-module-map}
 M\mu_{\mathfrak g}
 =\mu_{\mathfrak g}(\vid\star M)
 =\mu_{\mathfrak g}(M\star\vid)
\end{equation}
on both ordered branches. If it is also a coefficient derivation,
then $M=0$.
\end{lemma}

\begin{proof}
All maps are derived coefficient maps over $Y$.
On the two-root scheme write $j:U\hookrightarrow\mathscr R_2$
for the disjoint open and $i$ for the union of the two closed
diagonal branches. Put $R=\ell\boxtimes\ell$, pulled back from
$Y^2$. Both convolution summands have source $j_*j^*R$.
The difference $\kappa_a=[\iota,a]$ of their ordered products obeys
\[
 \Phi\mathcal T_C(\kappa_a)
       =\mu_{\mathfrak g}(\vid\star M)
\]
by the envelope relation and the chosen unit path.

The commutator may have an open part. Define on $U$
\[
 w=(R_B-\vid)(\iota\boxtimes a),\qquad
 P(w)=\mu_{12}j_*(w):j_*j^*R\longrightarrow B_2.
\]
The second formula uses the global product from the first
convolution summand. Thus $\kappa_a+P(w)$ restricts to zero on $U$.
The exchange identity for $\Phi$, and the strong symmetric
monoidal comparison for $\mathcal T_C$, give
\[
 \Phi\mathcal T_C(P(w))=0.
\]
Indeed the image of $w$ is killed by $R_A-\vid$.
The open projection formula carries this zero map and its
nullhomotopy through $j_*$ and multiplication.
No primitive is extended from the open.

Put $Q=\operatorname{cofib}(j_!j^*R\to j_*j^*R)$.
Open adjunction lifts $\kappa_a+P(w)$ through $Q$.
The source $R$ is regular holonomic, so $j_!$ is defined for this
source regardless of the unbounded target.
The two localization triangles give
\[
 i_*i^*R\longrightarrow Q\longrightarrow i_*i^!R[1].
\]
On each branch $i^!R=\ell[-1]$ and $i^*R=\ell[1]$.
The exact sequence of principal parts
$\omega\to\omega(*0)\to\delta_0$ fixes the first shift;
duality fixes the second.
These are the right-module conventions of
\cite{BD04}*{\S\S2.1.1--2.1.3}.
The map $j_*j^*R\to Q\to i_*\ell$ is the vacuum
principal-part bracket.

The closed projection formula and $i^!C_2=C_1$ give
\[
 \mathcal T_C(i_*i^*R)=i_*G[1],\qquad
 \mathcal T_C(i_*i^!R[1])=i_*G.
\]
The finite PBW filtration puts $A_2$ in the right-module heart.
Its quotients are open and closed direct images of degree-zero
external products; these images are exact on the affine collision
charts. Also $i^!A_2=G$.
Consequently
\[
 \operatorname{Hom}(i_*G[1],A_2)
   =\operatorname{Ext}^{-1}(G,G)=0.
\]
After applying $\mathcal T_C$ and $\Phi$, the lifted map from $Q$
therefore factors on each branch through an endomorphism $N$ of $G$.
It follows that
\[
 \mu_{\mathfrak g}(\vid\star M)=N\mu_{\mathfrak g}.
\]

Insert the vacuum in the first slot. The bracket becomes
$\mathrm{can}:j_*j^*(\ell\boxtimes G)\to i_*G$,
and the last equality becomes $(N-M)\mathrm{can}=0$.
Precomposition by $\mathrm{can}$ is injective on degree-zero
maps to $A_2$: its kernel is an image of
$\operatorname{Hom}((\ell\boxtimes G)[1],A_2)=0$,
again by the right-module heart.
Thus $N=M$. Skew symmetry gives the identity with $M$ in the
first slot, proving \eqref{eq:vol2-vacuum-module-map}.

If $M$ is a derivation, subtract its second-slot module identity
from the derivation identity.
The result is $\mu_{\mathfrak g}(M\star\vid)=0$.
Inserting the vacuum in the second slot and using the same
principal-part injectivity gives $M=0$.
\end{proof}

\begin{corollary}\label{cor:vol2-one-root-scalars}
Every $\Psi(a)$ over a nonempty affine open is a constant scalar
multiple of the identity.
\end{corollary}

\begin{proof}
The module identity implies commutation with every state mode.
Normal transfer can be written with the normal derivative in the
first coordinate; the matrix of $M$ on the second coordinate then
acquires no derivatives.
Put $c=M\mathbf1$. The regular coefficient of $J(z)$ multiplied by $(z-w)^{-n}$
detects the creation mode $J_{-n}$ in the principal-part module
identity. Commutation with these modes gives
$M(f)=fc$ for every polynomial state.
Commutation with $J_n=\kappa n\partial_{x_n}$ for $n>0$ gives
$\partial_{x_n}c=0$. Thus $c$ is a scalar function of the root.
Horizontality makes it constant.
\end{proof}

Both kinds of modes are necessary. The derivation
$D=\partial_{x_1}$ commutes with all annihilation modes, but
\[
 (DJ_{-1}-J_{-1}D)\mathbf1=\mathbf1.
\]
It is not a module map.

\subsection{A singleton class that no recipient kills}

Normalize the fibre of a right differential module at $q\in X(k)$ by
\[
 N_q=i_q^!N[1].
\]
Then $\ell_q=k$ and $G_q=V_\kappa$ lie in degree zero.
This is a derived restriction, so its induced maps on coefficient
complexes preserve boundaries.

\begin{lemma}\label{lem:vol2-affine-kunneth}
For a nonempty affine open $Y\subset X$, a complex $N$ of
differential modules on $Y$ and a complex $W$ of $k$-vector spaces,
\[
 \vRH_Y(\ell,N\otimes_kW)
 \simeq\vRH_Y(\ell,N)\otimes_kW.
\]
No bound on either complex is required. Every cohomology class
has finite support in the algebraic K\"unneth direct sum.
\end{lemma}

\begin{proof}
The unshifted left--right conversion takes $\ell$ to $\mathcal O_Y$.
For the affine coordinate $z$, the Spencer resolution is
\[
 0\longrightarrow D_Y
 \xrightarrow{P\mapsto P\partial_z}D_Y
 \xrightarrow{P\mapsto P(1)}\mathcal O_Y\longrightarrow0.
\]
Affineness makes it a two-term resolution by finite free modules
for the derived Hom in question. Its Hom complex commutes with
tensoring by any complex $W$. Over the field $k$, the cohomology
of an algebraic tensor product is the direct sum of the tensor
products of cohomology groups. Every element has finite support.
\end{proof}

Let $D_0\in\mathfrak d_V$ be a coefficient derivation, and let
$\alpha$ be any action of $F_1$ with first-weight class
$[\gamma_{D_0}^{(1)}]$ in
Theorem~\ref{thm:vol2-first-weight-classification}.
This includes the natural action $\alpha_{D_0}$.
For a receiving morphism
\[
 h:(F_1,\alpha)\longrightarrow(Z,\gamma),
\]
let $h_0$ be its twisting component on the pure singleton
$s(\varepsilon m)$, $m\in\mathfrak l$.
It defines a class
\[
 [h_0]\in H^1\vRH_X(\ell,B_{\{1\}}).
\]
This class forgets equivariance by $\mathfrak l$.
Nonvanishing of an equivariant extension would not by itself
prove its nonvanishing here.

\begin{theorem}\label{thm:vol2-plain-singleton}
If $D_0\ne0$, then $[h_0]|_Y\ne0$ on every nonempty affine
open $Y\subset X$, for every receiving pointed factorizing action.
\end{theorem}

\begin{proof}
Suppose $[h_0]|_Y=0$. Resolve the pointed arrow before taking
its relative convolution, retaining the full coefficient
differential. A primitive of $h_0$ defines a positive-actor-weight
gauge. Its weight-one equation replaces $h_0$ by zero exactly;
higher weights change by the usual gauge series, finite at each
configuration. This gives $h'\simeq h$ under $J$ with $h'_0=0$.
Composition transports the specified action path.

The first mixed component of $h'$ has one ordinary input and
one actor input. It is a closed degree-zero cochain $c_1$.
It need not be supported on collisions.
The normalized bar formula for $\mathcal T_C(h')$ on these inputs
has the term $\mathcal T_C(c_1)$ and products involving
$\mathcal T_C(h'_0)$. Those products vanish. Its composite
with the action therefore has mixed component
$\Phi\mathcal T_C(c_1)$.

Choose $q\in Y(k)$ and put $Y'=Y\setminus\{q\}$.
Such a choice exists because $k$ is infinite and $Y$ is a nonempty
open of $\mathbb A^1_k$.
Keep the actor root at $q$ and let the ordinary root vary over $Y'$.
The first ordered product and normalized fibre identify the target
on $Y'\times\{q\}$ with $B_{\{1\}}|_{Y'}\otimes_kB_q$.
Lemma~\ref{lem:vol2-affine-kunneth} writes the class of $c_1$
as a finite sum with bidegrees $(r,-r)$.
After the action, its second factors lie in
\[
 H^{-r}\vRH_k(V_\kappa,V_\kappa)=0\qquad(r\ne0).
\]
Thus the resulting operation has the form
\[
 \sum_{i=1}^m M_i\otimes N_i,\qquad
 M_i=\Psi(a_i),\quad N_i=\Psi_q(b_i),
\]
where $a_i$ and $b_i$ have degree zero and $\Psi_q$ is the
normalized point action. The module maps $M_i$ act on the moving
ordinary-root factor.

The action path identifies the first-weight class of this composite
with $[\gamma_{D_0}^{(1)}]$.
The two-output part of the proof of
Theorem~\ref{thm:vol2-first-weight-classification} identifies such
classes with actual module cocycles, without boundaries.
Indeed open adjunction kills the zero Chevalley column in every
degree, and total degree one leaves only its first column in
coefficient degree zero. This argument also holds on $Y$.
Restricting the resulting equality of maps gives
\begin{equation}\label{eq:vol2-singleton-finite-sum}
 D_0|_{Y'}\otimes\vid_{V_\kappa}
       =\sum_{i=1}^m M_i\otimes N_i.
\end{equation}

Algebraic tensors of endomorphisms inject into external
endomorphisms. For the finite list of second factors, finitely
many vector evaluations and linear functionals separate every
linearly independent sublist; applying them recovers each first
factor. Choose a functional $\lambda$ on
$\operatorname{span}_k\{\vid,N_1,\ldots,N_m\}$ with
$\lambda(\vid)=1$. Equation~\eqref{eq:vol2-singleton-finite-sum}
then gives
\[
 D_0|_{Y'}=\sum_i\lambda(N_i)M_i.
\]
Lemma~\ref{lem:vol2-vacuum-module-map} makes the right side a
module map. Since $D_0$ is a derivation, that lemma forces
$D_0|_{Y'}=0$.
The coefficient module $\mathcal O_X\otimes_kV_\kappa$ is
torsion-free, and every state has an algebraic image.
Localization is faithful on each such image, so $D_0=0$ on $X$,
a contradiction.
\end{proof}

\subsection{Homotopy invariance on two roots}

The obstruction must survive homotopies of coalgebra maps,
including those with nonzero intermediate singleton components.
Use compatible differential graded models of the full coefficient
mapping complexes, with the pointed arrow resolved, and write
$\partial$ for their differential.
The two pure inputs $s(\varepsilon_1m)$ and
$s(\varepsilon_2n)$ have degree zero. Their bidegrees record the
numbers of $\varepsilon_1$ and $\varepsilon_2$.
Let $j$ denote restriction to the two-root disjoint open and put
$P_{1,1}=\mathfrak l\star\mathfrak l$, keeping the labels.

\begin{lemma}\label{lem:vol2-binary-homotopy}
Suppose maps $\phi_0,\phi_1:F_2\to Z$ under $J$ have twisting
representatives with zero pure singleton actor components.
If they are homotopic under $J$, their restricted pure binary
components give the same class in
\[
 H^1\vRH(j^*P_{1,1},j^*B).
\]
\end{lemma}

\begin{proof}
Represent the relative homotopy by $\tau(u)+c(u)\,du$,
where $\tau$ and $c$ have degrees one and zero.
Polynomial interval forms suffice in the positive filtration,
which is finite on each configuration.
On a pure singleton input the source differential and reduced
coproduct vanish. Hence
\[
 \tau_{1,0}=\partial E,\quad \tau_{0,1}=\partial F,
 \qquad E(u)=\int_0^u c_{1,0}(v)\,dv,
 \quad F(u)=\int_0^u c_{0,1}(v)\,dv.
\]
We have $E(0)=F(0)=0$ and
$\partial E(1)=\partial F(1)=0$.
On the disjoint open the source bracket vanishes. The mixed
interval equation, with a prime denoting $u$-derivative, is
\[
 \tau'_{1,1}=\partial c_{1,1}
             +[\partial E,F']+[\partial F,E'].
\]
The brackets use both labelled external products followed by
the target multiplication. Since $E$ and $F$ have degree zero,
the graded Leibniz rule gives
\begin{equation}\label{eq:vol2-binary-boundary}
 [\partial E,F']+[\partial F,E']
       =\partial[E,F']-\frac{d}{du}[E,\partial F].
\end{equation}
Therefore
\[
 \tau_{1,1}(1)-\tau_{1,1}(0)
 =\partial\int_0^1\bigl(c_{1,1}+[E,F']\bigr)\,du
       -[E,\partial F]\big|_0^1.
\]
The endpoint term is zero. At both endpoints the
Maurer--Cartan equation makes the restricted binary term closed,
because the singleton terms vanish. Their difference is thus a
boundary in the stated full coefficient complex.
\end{proof}

The commutator correction need not be zero as a cochain.
For $E=u(1-u)E_0$ and $F=u^2(1-u)F_0$, with $E_0,F_0$
of degree zero in a differential graded algebra,
\[
 \int_0^1[E,F']\,du=\frac1{60}[E_0,F_0].
\]
Its differential can be nonzero; the lemma uses its exactness.

\begin{theorem}\label{thm:vol2-pointed-sphere}
For every morphism
$h:(F_1,\alpha_D)\to(Z,\gamma)$ in $\vAct$, the composites
$hf$ and $hf_0$ from \eqref{eq:vol2-pointed-sphere} are not
homotopic under $J$.
Thus $\vAct$ is not filtered, and $\Phi_C$ is not representable
on pointed factorizing coalgebras for the polynomial Heisenberg
boundary at nonzero level.
\end{theorem}

\begin{proof}
Both composites have zero pure singleton actor components.
On two distinct roots the source cobar differential gives
\[
 d\,t\bigl(s(\varepsilon n)s(m)\bigr)
 =-t\bigl(s(\varepsilon n)\bigr)t\bigl(s(m)\bigr)
   +t\bigl(s(m)\bigr)t\bigl(s(\varepsilon n)\bigr).
\]
Applying $\Omega h$ reorders $h_0(n)\iota(m)$ into
$\iota(m)h_0(n)$ modulo a boundary.
The boundary comes from the cocommutative source; no centrality
of the vacuum in $Z$ is assumed.

Fix $q\in X(k)$ and put $Y=X\setminus\{q\}$.
Restrict to $Y\times\{q\}$ by the normalized derived fibre.
In the first product order, the binary class of $hf$ is
\begin{equation}\label{eq:vol2-sphere-binary-class}
 \frac12\bigl([h_0]|_Y\otimes[\iota_q]
                  -[\iota]|_Y\otimes[h_0]_q\bigr).
\end{equation}
Lemma~\ref{lem:vol2-affine-kunneth} puts the two terms in the
distinct K\"unneth summands of bidegrees $(1,0)$ and $(0,1)$.
The first is nonzero by Theorem~\ref{thm:vol2-plain-singleton}
and $\Psi_q(\iota_q)=\vid_{V_\kappa}$.
The second cannot cancel it and need not itself be nonzero.
The composite $hf_0$ has zero binary class.
Lemma~\ref{lem:vol2-binary-homotopy} excludes a homotopy under
$J$, and hence a homotopy of pointed actions with its action
two-cell.

A filtered infinity-category equalizes every parallel pair
after some receiving morphism. The pair
\eqref{eq:vol2-pointed-sphere} violates that condition.
A representing object for $\Phi_C$ would be a final object of
$\vAct$; its mapping space from $(F_2,\mathrm{triv})$ would
be contractible. That too would identify the two composites,
which is impossible.
\end{proof}

\begin{corollary}\label{cor:vol2-independent-unit-products}
For every such $h$, the two classes represented on the disjoint
open by $h_0(m)\iota(n)$ and $h_0(n)\iota(m)$ are linearly
independent in $H^1\vRH(j^*P_{1,1},j^*B)$.
\end{corollary}

\begin{proof}
Write the classes as $c_1,c_2$, with the root labels fixed.
Fix $q\in X(k)$ and put $Y=X\setminus\{q\}$.
Suppose $a c_1+b c_2=0$. On $Y\times\{q\}$ the
$(1,0)$ K\"unneth component is
$a[h_0]|_Y\otimes[\iota_q]$, which is nonzero unless $a=0$.
Now restrict to $\{q\}\times Y$ and take the moving root first
in the product order. The same component is
$b[h_0]|_Y\otimes[\iota_q]$, so $b=0$.
The source cobar boundary used in the theorem permits both
reorderings. Thus neither centrality of the recipient vacuum nor
nonvanishing of $[h_0]_q$ is required.
\end{proof}

Multiplying the marked matrix of $f$ by $2\lambda$, for
$\lambda\in k^\times$, replaces $1/2$ in
\eqref{eq:vol2-sphere-binary-class} by $\lambda$.
The coordinate contraction remains zero, so the same construction
and obstruction apply.

The theorem concerns pointed actions with all pointed maps.
For the additive exterior multiplications, $f$ is not a monoid map.
Its first marked form is $(t_2\,dt_1-t_1\,dt_2)/2$.
Pulling it back by addition and subtracting its values on the two
factors leaves
\[
 \tfrac12(s_2\,dt_1-s_1\,dt_2+t_2\,ds_1-t_1\,ds_2),
\]
which is nonzero in the first marked-matrix invariant.
Thus the displayed pair does not decide finality in $\vMon$,
where comparison maps must preserve multiplication.

\subsection{How the ambient endomorphism object kills the class}

The object $\vEnd$ of Lemma~\ref{lem:vol2-ambient-endomorphisms}
receives every action, without a factorization requirement.
Its failure to factorize can be seen directly in its cochains.
Put $c=\overline C$ and $a=\overline{\Omega C}$, and let
$\underline{\vRH}^{!}$ be the internal Hom for $\vot$ on positive
coefficient families. There is an equivalence
\begin{equation}\label{eq:vol2-internal-action-cochains}
 \Omega\vEnd\simeq K\oplus\underline{\vRH}^{!}(c,a).
\end{equation}
The differential on the right is the internal mapping
differential. Its cup product uses the coproduct of $C$,
the interchange map, evaluation, and multiplication in $\Omega C$.

To prove the underlying equivalence, test against the free augmented
algebra $T_\star(V)$ on a positive coefficient family $V$.
Its bar contracts to $K\oplus sV$, with zero reduced coproduct:
the same cut contraction as in
Lemma~\ref{lem:vol2-ambient-endomorphisms} leaves only the
terms of tensor length one. Tensoring with $C$ leaves the primitive reduced
coefficient $sV\vot c$. The bar--cobar and internal-Hom
adjunctions identify maps successively with
\[
 \vMap_{\mathscr C}(V,\overline{\Omega\vEnd})
 \simeq\vMap_{\vM}\bigl((K\oplus sV)\vot C,C\bigr)
 \simeq\vMap_{\mathscr C}\bigl(V,
                    \underline{\vRH}^{!}(c,a)\bigr).
\]
Yoneda gives the equivalence and its differential.
For the product, take strict representatives and put
$H=K\oplus\underline{\vRH}^{!}(c,a)$ with the stated cup product.
Let $\pi:\overline{\vB H}\to\underline{\vRH}^{!}(c,a)$ be the
bar twisting cochain. Evaluation gives
\[
 \tau=\operatorname{ev}(\pi\vot1):
                 \overline{\vB H}\vot c\longrightarrow a.
\]
Naturality of interchange identifies the quadratic term in the
twisting equation for $\tau$ with the evaluation of the cup-product
term for $\pi$. Evaluation is a chain map, so the linear terms
agree as well. Thus $\tau$ defines $\vB H\vot C\to C$.
Adjunction and cobar give $H\to\Omega\vEnd$.
On a map $V\to\underline{\vRH}^{!}(c,a)$, its adjoint is exactly
evaluation of that map. The underlying ideal map is therefore
the identity under the equivalence just proved. This proves the
algebra equivalence in \eqref{eq:vol2-internal-action-cochains}.

\begin{proposition}\label{prop:vol2-ambient-kills-singleton}
For every coefficient derivation $D_0\in\mathfrak d_V$, the
classifying map $\beta_{\alpha_{D_0}}:F_1\to\vEnd$ has zero
singleton actor component. For every receiving morphism
$h:(F_1,\alpha_{D_0})\to(Z,\gamma)$, the classifying map
$\beta_\gamma:Z\to\vEnd$ sends $[h_0]$ to zero.
The ambient endomorphism object $\vEnd$ is not factorizing.
\end{proposition}

\begin{proof}
The singleton classifying map is adjoint to the singleton action.
The twisting cochain for $\alpha_{D_0}$ needs one ordinary input;
a pure singleton actor input has none. Its image is therefore
zero. The specified action path gives
$\beta_\gamma h\simeq\beta_{\alpha_{D_0}}$ under $J$.
Restriction to one position proves the assertion about $[h_0]$.

Equation~\eqref{eq:vol2-internal-action-cochains} supplies a
degree-zero internal cochain $a_D:\mathfrak l\to\Omega\vEnd$
whose one-root action is $D$.
At a singleton, its adjunction is precisely that between
$\ell\vot C_1=G$ and $D:G\to G$;
\eqref{eq:vol2-main-diagonal-adjunction} extends it to the full
coefficient family. If $\vEnd$ factorized,
Lemma~\ref{lem:vol2-vacuum-module-map} would make $D$ a module
map, contradicting $(DJ_{-1}-J_{-1}D)\mathbf1=\mathbf1$.
\end{proof}

For any nonzero coefficient derivation $D_0$, the same adjunction
gives $a_{D_0}:\mathfrak l\to\Omega\vEnd$.
Its vacuum commutator $[\iota,a_{D_0}]$ is nonzero on the
two-root disjoint open.
If it vanished there, the commutator itself would lift through
$Q$ in the proof of Lemma~\ref{lem:vol2-vacuum-module-map}.
The supported part of that proof would give
$\mu_{\mathfrak g}(\vid\star D_0)=N\mu_{\mathfrak g}$.
The first vacuum slot forces $N=D_0$; the derivation identity
and the second vacuum slot then force $D_0=0$.
Once open vanishing is assumed, this argument uses no
factorization of the target.
The nonzero open commutator becomes zero after the action,
whose bracket image is supported on the diagonal.
This explains how $\vEnd$ can carry the mixed derivation
while killing the singleton class.

Finally, Theorems~\ref{thm:vol2-density-criterion} and
\ref{thm:vol2-pointed-sphere} imply that $L_C$ is not factorizing:
at least one of its two-block kernel coefficients $K_{C,A|B}$ is
nonzero. This consequence does not yet exhibit a representative
of such a coefficient in the full density bar. That bar also
retains maps capable of killing defects of individual objects.

\section{Multiplication and a surviving commutator}
\label{sec:vol2-monoid-obstruction}

The pointed sphere did not use a multiplication on its source
actor. To obstruct a final multiplicative action, we need two
monoid-action maps. A trivially acting submonoid supplies them:
its inclusion and the map through the unit have the same action.
The multiplication of a larger actor can nevertheless distinguish
these maps after every receiving action.

Work over $\mathbb C$. Let $Q$ be an algebraic unital vertex algebra
in cohomological degree zero, with translation $T$, and let
$C_Q=\vCE(L_Q)$ on the full ordered collision carrier.
A vertex derivation is a linear map $D:Q\to Q$ satisfying
\[
 D\mathbf1=0,\qquad DT=TD,\qquad
 D(x_{(n)}y)=(Dx)_{(n)}y+x_{(n)}Dy\quad(n\in\mathbb Z).
\]
We retain the category of multiplicative actions defined in
Section~\ref{sec:vol2-action-carrier}. Receiving actors may be
unbounded. They need not have an augmentation to $J$.

\subsection{An exterior monoid with a trivial direction}

Let $A_0=\mathbb C a\oplus\mathbb C b$ have multiplication
\[
 a^2=a,\qquad ab=ba=b,\qquad b^2=0.
\]
The unit $a$ belongs to this auxiliary algebra; its exterior
generator is distinct from the vacuum. Set
\[
 \mathsf F=C_{\Lambda(A_0[-1])},\qquad
 \mathsf G=C_{\Lambda(\mathbb C a[-1])},\qquad
 \mathsf H=C_{\Lambda(\mathbb C b[-1])}.
\]
The exterior generators have degree one, zero differential and
zero translation. For $x\in A_0$ let $L_x$ be multiplication by
$x$ and put
\[
 R(x)=e^{L_x},\qquad
 M(x)=\sum_{n\geq0}\frac{L_x^n}{(n+1)!}.
\]
The constant term of $M$ is the identity. Define the marked
multiplication matrix
\begin{equation}\label{eq:vol2-dual-number-product}
 (\mathcal L(x,y)\mid\mathcal R(x,y))
 =\bigl(M(x+y)^{-1}M(x)\mid
        M(x+y)^{-1}R(x)M(y)\bigr).
\end{equation}

\begin{lemma}\label{lem:vol2-dual-number-monoid}
This matrix gives a coherent unital monoid on $\mathsf F$ under
and over $J$, with coordinate contraction $x+y$.
The two inclusions
$i:\mathsf G\to\mathsf F$ and $j:\mathsf H\to\mathsf F$
are monoid maps. The induced multiplication on $\mathsf H$ is
additive. All these maps preserve the auxiliary grading
$\deg a=0$, $\deg b=1$.
\end{lemma}
\begin{proof}
The law
\[
 (x,U)(y,V)=(x+y,U+R(x)V)
\]
is associative because $R(x+y)=R(x)R(y)$.
Changing fibre coordinates by $U=M(x)u$ gives
\eqref{eq:vol2-dual-number-product}. Its three associativity
identities are \eqref{eq:vol2-multiplication-blocks}, and its
unit identities follow from $M(0)=I$. In $A_0$,
\[
 M(x)x=e^x-a,\qquad
 e^{x+y}-a=(e^x-a)+e^x(e^y-a).
\]
Thus the diagonal $u=x$ is a homomorphism, proving the coordinate
contraction. The $a$- and $b$-lines are preserved; on the $b$-line,
$L_{sb}b=0$, so both multiplication blocks are the identity.
The algebra product preserves auxiliary degree, hence so do
$R$, $M$ and $M^{-1}$.

Theorem~\ref{thm:vol2-natural-actor-maps} realizes these matrices
and their identities, with a contractible space of coherent
choices. The multigraded contractions in its proof preserve the
auxiliary grading. At $m$ positions, exterior degree bounds the
source actor weight by $2m$. The formal series therefore evaluate
as finite sums on the original coefficient diagrams.
\end{proof}

The zero auxiliary degree of $a$ prevents the use of positive
grading alone to remove action loops. Prescribing the full
$\mathsf G$-action removes those loops instead.

\begin{lemma}\label{lem:vol2-relative-dual-number-action}
For every vertex derivation $D$, there is a coherent unital
$\mathsf F$-action $\alpha$ on $C_Q$ whose restriction to
$\mathsf G$ is trivial and whose restriction to $\mathsf H$ is
the single-derivation action $\alpha_D$.
With these comparisons, $i$ and $j$ are monoid-action maps.
\end{lemma}
\begin{proof}
First use one formal derivation. For $r$ actor copies let
$t_\nu,s_\nu$ be coordinates dual to $a,b$. Give them auxiliary
weights $0,-1$ and give $D$ weight $1$. Take natural cochains
of weight zero which vanish when all $b$-inputs are absent.
The scalar-bar and Koszul projection of
Theorem~\ref{thm:vol2-natural-cohomology} has target
\begin{equation}\label{eq:vol2-relative-monoid-model}
 \left[(s_1,\ldots,s_r)
 \mathbb C[[t_1,s_1,\ldots,t_r,s_r]]
       \otimes\Lambda(\mathbb C D[-1])\right]^{\mathbb G_m}.
\end{equation}
Completion is in total actor weight. Restriction to the
$a$-subspace has a splitting, and taking auxiliary weight zero
is exact in each finite total weight and derivative order.
The projection is consequently surjective with weightwise
acyclic kernel on this relative complex.

The target has no nonpositive degrees: a scalar monomial in
the displayed ideal has negative auxiliary weight. Its degree-one
part consists of $\sum_\nu s_\nu f_\nu(t_1,\ldots,t_r)D$.
Relative Maurer--Cartan lifting therefore makes each corresponding
component contractible. The same assertion holds with prescribed
unit faces and restrictions to the $b$-lines. Indeed the deletion
projections for actor copies and the two exterior directions
commute and split. Inclusion--exclusion extends compatible data
on their union; the residual kernel is a sum of the same acyclic
multigraded kernels. Polynomial forms relative to parameter
boundaries then supply the required fillings.

Choose the class $sD$ on one copy. Pullback by multiplication
has class $(s_1+s_2)D$, since the coordinate contraction is
addition. The composite action has the same class by
\eqref{eq:vol2-action-BCH}, or directly by
$e^{s_1D}e^{s_2D}=e^{(s_1+s_2)D}$.
Contractibility supplies its multiplication path, relative to
the prescribed line actions and units. Each subsequent module
associahedron has boundary in the component with class
$(s_1+\cdots+s_r)D$. Relative contractibility fills it and all
parameterized comparison boundaries.

Finally evaluate on $Q$. The strong symmetric monoidal functor
$\mathcal T_Q=(-)\vot C_Q$ of
Lemma~\ref{lem:vol2-strong-boundary-tensor} sends the vacuum
Lie object to $L_Q$, with its bracket. It therefore identifies
$\mathsf F\vot C_Q$ with
$\vCE(L_Q\otimes\Lambda(A_0[-1]))$ and commutes with envelopes
and bars. The infinitesimal vertex automorphism $1+uD$ over
$\mathbb C[u]/(u^2)$ shows that this comparison carries $D$ to
the coefficient derivation, retaining collisions and normal
transfers. It carries the universal polynomials and their
parameter homotopies as well. At each configuration, actor
weight is bounded and derivative order equals the $b$-count.
Every evaluated operation is finite.
\end{proof}

\subsection{A commutator detected on two roots}

Let $\mathsf S=C_{\Lambda(\varepsilon_1,\varepsilon_2)}$.
Map it to $\mathsf F^{\vot4}$ by
\[
 \varepsilon_1\longmapsto a^{(1)}-a^{(3)},\qquad
 \varepsilon_2\longmapsto b^{(2)}-b^{(4)},
\]
and multiply the four factors in that order. Call the resulting
pointed map $w:\mathsf S\to\mathsf F$. Define the formal series
$\chi(t)=\sum_{n\geq0}t^n/(n+1)!$.

\begin{lemma}\label{lem:vol2-monoid-commutator}
There is a natural pointed map $f:\mathsf S\to\mathsf H$ with
\begin{equation}\label{eq:vol2-monoid-commutator}
 w\simeq jf,\qquad A_f(t,s)=(-s\chi(t)\mid t\chi(t)).
\end{equation}
Its coordinate contraction and singleton actor components vanish.
Its first marked form is $t\,ds-s\,dt$.
\end{lemma}
\begin{proof}
In the affine coordinates used above, an $a$-line element with
coordinate $t$ and fibre $ra$ has linear part $e^t$ and
translation $\chi(t)ra$. A $b$-line element with coordinate $s$
and fibre $qb$ has linear part multiplication by $a+sb$ and
translation $qb$. Negating coordinate and fibre gives inversion,
since $M(-x)=R(-x)M(x)$. The linear parts commute. Their affine
commutator has linear part one and translation
\[
 -s\chi(t)rb+(e^t-1)qb=\chi(t)(tq-sr)b.
\]
Its coordinate is zero, where the frame is the identity.
The marked-map classification now gives \eqref{eq:vol2-monoid-commutator}
and its pointed homotopy. Finally $A_f(t,s)(t,s)^T=0$,
$A_f(0,0)=0$ and $\chi(0)=1$ prove the remaining assertions.
\end{proof}

The singleton detector extends from the Heisenberg algebra to
$Q$ without a scalar-endomorphism classification.
Let $h:(\mathsf H,\alpha_D)\to(Z,\beta)$ be any receiving
pointed factorizing action and put $B=\Omega Z$.
Write $h_0$ for its pure singleton twisting component and
$\iota:L_{\mathbb C}\to B$ for its pointing.

\begin{lemma}\label{lem:vol2-vertex-singleton}
If $D\ne0$, then $[h_0]|_Y\ne0$ in
$H^1\vRH_Y(\ell,B_{\{1\}})$ for every nonempty affine
open $Y\subset X$.
\end{lemma}
\begin{proof}
The proof of Lemma~\ref{lem:vol2-vacuum-module-map} applies
with $G=\omega_X\otimes Q$ and $\mathfrak g=L_Q$.
Here are the hypotheses it uses. Factorization makes both
ordered cobar products invertible on distinct roots. The open
commutator correction has zero action by the exchange identity
and open projection. Its remaining supported part factors
through the vacuum localization cofiber, whose terms on each
branch are $\ell[1]$ and $\ell$. Tensoring with $C_Q$ replaces
them by $G[1]$ and $G$. The PBW envelope of $L_Q$ lies in the
right-module heart on each affine collision chart, with
extraordinary diagonal restriction $G$. Negative Ext therefore
kills the first term and makes the principal-part vacuum map
injective on degree-zero maps. Inserting the vacuum in both
slots gives
\[
 M\mu_{L_Q}=\mu_{L_Q}(1\star M)=\mu_{L_Q}(M\star1)
\]
for every plain degree-zero coefficient operation $M$.
Such an operation that is also a derivation is zero.

If $[h_0]|_Y=0$, the relative gauge used in
Theorem~\ref{thm:vol2-plain-singleton} removes $h_0$.
Separate the moving ordinary root on $Y'=Y\setminus\{q\}$
from the actor root at $q\in Y$. The finite Spencer comparison
of Lemma~\ref{lem:vol2-affine-kunneth} applies to arbitrary
receiving complexes. Since $Q$ is in degree zero, only the
bidegree $(0,0)$ survives after its action. The action path gives
\[
 D|_{Y'}\otimes1_Q=\sum_{\nu=1}^m M_\nu\otimes N_\nu.
\]
This is equality of two-output cocycles. Resolve the coefficient
maps in nonnegative degrees before taking the Chevalley complex.
Open adjunction annihilates its zero column; total degree one
then has only column one in coefficient degree zero, with no
incoming boundary. Each $M_\nu$ satisfies the module
identities just proved. Finite vector evaluations separate the
span of $1_Q,N_1,\ldots,N_m$. Applying a functional taking
$1_Q$ to $1$ makes $D|_{Y'}$ a linear combination of these
module maps. It is a derivation as well, so it vanishes.
Localization is faithful on the free algebraic state module
$\mathcal O_X\otimes Q$. Hence $D=0$, a contradiction.
\end{proof}

\begin{theorem}\label{thm:vol2-monoid-nonfinality}
If $Q$ has a nonzero vertex derivation, its category of coherent
unital multiplicative actions on ordinary factorizing coalgebras
has no final object. There are parallel monoid-action maps
which no receiving monoid action equalizes up to homotopy.
\end{theorem}
\begin{proof}
Let $i_0:\mathsf G\to J\to\mathsf F$ be augmentation followed
by the unit. Lemma~\ref{lem:vol2-relative-dual-number-action}
makes $i$ and $i_0$ morphisms from the trivial $\mathsf G$-action
to $(\mathsf F,\alpha)$.

For every receiving pointed map $h$ from $(\mathsf H,\alpha_D)$,
the binary class of $hf$ on $(X\setminus\{q\})\times\{q\}$ is
\begin{equation}\label{eq:vol2-monoid-binary-class}
 [\iota]|_{X\setminus\{q\}}\otimes[h_0]_q
       -[h_0]|_{X\setminus\{q\}}\otimes[\iota]_q.
\end{equation}
The two terms lie in distinct K\"unneth bidegrees. The second
is nonzero by Lemma~\ref{lem:vol2-vertex-singleton} and the unit
identity: $[\iota]_q$ acts as $1_Q$.
Changing the ordered product changes the representative by the
cobar boundary used in Theorem~\ref{thm:vol2-pointed-sphere}.
Lemma~\ref{lem:vol2-binary-homotopy} makes this class invariant
under every pointed geometric homotopy. Thus $hf$ cannot be
homotopic to the map through $J$.

Now suppose a receiving monoid-action map
$g:(\mathsf F,\alpha)\to(Z,\beta)$ equalized $i$ and $i_0$.
Forgetting the action retains a homotopy under $J$. In the
four-factor word $gw$, multiplicativity of $g$ and that
homotopy replace the first and third factors by units.
The second and fourth factors cancel, since $j$ is a monoid map
and the additive $\mathsf H$ has inverse $b\mapsto-b$.
Tensoring the comparison homotopies and applying the monoid unit
coherences therefore gives $gw\simeq\mathsf S\to J\to Z$.
But $w\simeq jf$, and $gj$ receives $\alpha_D$, contradicting
\eqref{eq:vol2-monoid-binary-class}.

A final action would receive $(\mathsf F,\alpha)$ and identify
the two composites from the trivial $\mathsf G$-action, because
its mapping space from that action would be contractible.
The contradiction excludes finality.
\end{proof}

For the polynomial Heisenberg algebra at nonzero level,
$D=\partial_{x_1}$ is a nonzero vertex derivation by
\eqref{eq:vol2-Heisenberg-derivations}. The theorem therefore
excludes a final multiplicative action on this boundary over
$\mathbb C$, independently of the pointed-sphere argument.
For $Q=\mathbb C$, by contrast, the action category is the
slice of transverse monoids over $J$: an action
$F\vot J\to J$ is a monoid map $F\to J$.
Its final object is $1_J$, with evaluation $J\vot J\simeq J$.
The nonzero-derivation hypothesis thus distinguishes an actual
positive case. The ambient endomorphism object remains defined
in both cases, with its different universal property.

\section{Recipient bars and the role of maps}
\label{sec:vol2-recipient-bars}

The density comparison depends on the maps of its index,
not just on the available actors. A weighted bar records
which two actors meet in a common recipient, and how
different choices of their incoming maps are related.

Let $\mathcal T$ be a small infinity-category, let
$Q:\mathcal T^2\to D(k)$ be a coefficient diagram, and
write $\Delta:\mathcal T\to\mathcal T^2$ for the diagonal.
For $d=(a,b)$ the category $d\downarrow\Delta$ consists
of recipients $t$ with maps $a\to t$ and $b\to t$,
including their mapping spaces and coherences. Put
\[
 W(d)=C_*(|d\downarrow\Delta|;k),\qquad
 \widetilde W(d)=\operatorname{fib}(W(d)\to k).
\]
Chains have nonpositive cohomological degrees and the map
is augmentation. An empty recipient category gives the
same formula with $W(d)=0$.

\begin{proposition}\label{prop:vol2-full-recipient-coend}
For the full coefficient diagram, there is an equivalence
\begin{equation}\label{eq:vol2-full-recipient-coend}
 \operatorname{fib}\left(
   \operatorname*{hocolim}_aQ(a,a)
       \longrightarrow\operatorname*{hocolim}_{a,b}Q(a,b)
                  \right)
 \simeq\int^{d\in\mathcal T^2}_{\mathbb L}
                   Q(d)\otimes_k\widetilde W(d).
\end{equation}
No splitting of a vacuum coefficient is required.
\end{proposition}
\begin{proof}
The weight $W(d)$ is the homotopy colimit over $t$ of
the chains on $\operatorname{Map}_{\mathcal T^2}(d,\Delta t)$.
Derived co-Yoneda identifies its coend with
$\operatorname*{hocolim}_tQ(\Delta t)$.
The constant weight $k$ gives
$\operatorname*{hocolim}_dQ(d)$, and its augmentation is
the comparison displayed above. Tensor product and coend
preserve colimits. In a stable category they therefore
preserve finite fibres, proving the formula.
\end{proof}

The same proof works for coefficients in the stable
differential-module category on $U_P$, with its tensoring
over $D(k)$. For the binary density comparison use
\[
 Q(a,b)=(F_{a,I_1}\boxtimes F_{b,I_2})|_{U_P}.
\]
Factorization identifies its diagonal value with $X_P(F_a)$.
Thus \eqref{eq:vol2-full-recipient-coend} computes $K_{C,P}$
on its full index, before any coefficient reduction.

For an ordinary category $\mathcal D$ and strict resolved
diagrams, the normalized two-sided bar has length-$n$ term
\begin{equation}\label{eq:vol2-weighted-bar}
 \bigoplus_{d_0\to\cdots\to d_n}
          \bigl(Q(d_0)\otimes_k\widetilde W(d_n)\bigr)[n].
\end{equation}
Degenerate strings are removed. The first face transports
$Q$, the interior faces compose arrows, and the last pulls
the weight back. The differential is
$\partial+(-1)^nd_{\mathrm{int}}$, with
$\partial=\sum_{i=0}^n(-1)^id_i$ for $n>0$ and no face
term at length zero. Total degree is internal degree
minus $n$, and realization uses direct sums. For an
infinity-category the bar uses chains on mapping spaces,
not just sets of arrows.

If $\mathcal D$ has a terminal object $t$, the bar of
a covariant diagram $V$ contracts to $V(t)$. Let
$\epsilon$ transport length-zero coefficients to $t$
and let $i$ include that summand. On a length-$n$ string,
append the unique arrow to $t$ with sign $(-1)^{n+1}$,
and set the result to zero if that arrow is an identity.
The resulting degree-minus-one map satisfies
\begin{equation}\label{eq:vol2-terminal-contraction}
 dh+hd=1-i\epsilon.
\end{equation}
All faces except the last cancel against the corresponding
faces of the original string. The last gives the identity;
at length zero the other face gives $-i\epsilon$.
The length change reverses the internal differential sign,
so the internal terms cancel too.

\subsection{Equality patterns}

Let $\mathrm{FinInj}$ be a skeleton of finite sets and
injections, with $V(S)=k[S]$. Its monoidal product is
disjoint union and its initial unit is the empty set.
For $r\geq1$, consider
\[
 \operatorname*{hocolim}_S V(S)^{\otimes r}
       \longrightarrow
       \left(\operatorname*{hocolim}_SV(S)\right)^{\otimes r}.
\]
If $B_r$ is the number of partitions of $[r]$, this is
the summation map $k^{B_r}\to k$ in degree zero.
Its fibre is $k^{B_r-1}$.

Indeed the source bar is the linearized nerve of the
category of pairs $(S,t)$ with $t\in S^r$. An injection
preserves the equality partition of $t$. Each such
component has an initial object: the set of blocks,
with its canonical tuple. Its unique injection to any
tuple of that pattern determines every element. Thus
each partition contributes one contractible component.
The case $r=1$ computes the target as $k$, and every
pattern maps to its generator. The discrete partition
gives a section. If all finite-set maps are allowed,
the singleton is terminal, so both sides contract to
$k$ and the comparison is an equivalence.

The first fibre dimensions are $0,1,4,14,51,202,876,4139$.
For two singleton inputs, injections retain separate
same-image and distinct-image recipient components.
A quotient map joins them. Common upper bounds and an
initial unit therefore do not imply diagonal cofinality.

A pointed variant gives the split map explicitly.
Let $\mathsf I$ be finite pointed sets and
point-preserving injections, and put
\[
 V=\operatorname*{hocolim}_{\mathsf I}k[S],\qquad
 W=\operatorname*{hocolim}_{\mathsf I}k[S^2].
\]
The basepoint is a nonzero basis vector. There are two
one-entry patterns, and five two-entry patterns
\[
 (0,0),\quad(0,x),\quad(x,0),\quad(x,x),\quad(x,y),
       \qquad 0,x,y\text{ distinct}.
\]
Their minimal pointed sets are initial in their
components. Hence $V=k^2$ and $W=k^5$, in degree zero.
The comparison $d:W\to V\otimes V$ and its section are
\begin{equation}\label{eq:vol2-pointed-matrices}
 d=\begin{pmatrix}
 1&0&0&0&0\\0&1&0&0&0\\0&0&1&0&0\\0&0&0&1&1
 \end{pmatrix},\qquad
 s=\begin{pmatrix}
 1&0&0&0\\0&1&0&0\\0&0&1&0\\0&0&0&0\\0&0&0&1
 \end{pmatrix}.
\end{equation}
The section comes from Cartesian products of pointed
sets: $(x,y)$ is sent to $((x,0),(0,y))$.
When both entries are nonzero this selects distinct
images. Thus $ds=1$, but the kernel is generated by
$(0,0,0,1,-1)^T$, on which $1-sd$ is the identity.
The cofiber is $k[1]$. This is a computation on the
stated finite-set category, not on $\vAct$.

\subsection{Parallel maps}

Let $\mathcal U$ have two objects $a,b$ and two distinct
nonidentity arrows $f,g:a\to b$. For the constant
diagram $Q=k$ on $\mathcal U^2$, the fibre of its
diagonal comparison has
\[
 H^{-1}=k,\qquad H^0=k,
\]
and no other cohomology. The nerve of $\mathcal U$ is
a circle, and that of $\mathcal U^2$ is a torus.
The diagonal is an isomorphism on $H_0$ and the
diagonal inclusion $k\to k^2$ on $H_1$. Its fibre
therefore receives the target $H_2$ in cohomological
degree $-1$ and the quotient of $H_1$ in degree zero.
Adjoining $h:b\to c$ with $hf=hg$ makes $c$ terminal.
The comparison then contracts to the identity of $k$.

Even the recipient components do not suffice by
themselves: for $(a,a)$ there are three components.
The choices $(f,f)$ and $(g,g)$ connect through
$(a,\vid,\vid)$, whereas $(f,g)$ and $(g,f)$ are
isolated. The coend also retains transport in the
source variables and every bar face. A functional
detecting a density kernel class must be a cocycle
on this whole bar. In degree zero its components
must obey $\ell_bV(f)=\ell_a$ for every arrow;
its higher components obey the dual bar equations.

\section{Evaluation and the remaining homotopies}
\label{sec:vol2-density-homotopies}

The density kernel acts trivially on the boundary.
Its own vanishing is a stronger assertion. We give
the action nullhomotopy first, then write the equations
for a nullhomotopy of the projector itself.

Fix a binary partition $P$. On $U_P$ abbreviate
\[
 Z=X_P(L_C),\quad W=Y_P(L_C),\quad
 A=X_P(C),\quad A'=Y_P(C),\quad
 \delta=\delta_{P,L_C},\quad t=\delta_{P,C}.
\]
The map $t$ is an equivalence. Let $s$ be the density
section and $p=1-s\delta$. The action supplies
$\alpha:Z\vot A\to A$ and
$\alpha^Y:W\vot A'\to A'$.
Choose resolved models of their full diagram. For
homogeneous maps use $\partial u=du-(-1)^{|u|}ud$.
The prescribed paths give degree-minus-one maps
$e,n,h$ and a chain map $g:A'\to A$ with
\[
 \partial e=\delta s-1,\qquad
 \partial n=t\alpha-\alpha^Y(\delta\vot t),\qquad
 \partial h=gt-1.
\]
Put
\begin{equation}\label{eq:vol2-density-evaluation-homotopy}
 \sigma=-e\delta,\qquad
 N=\alpha^Y(\sigma\vot t)+n(p\vot1),\qquad
 H_{\mathrm{ev}}=gN-h\alpha(p\vot1).
\end{equation}
Then $\partial\sigma=\delta p$,
$\partial N=t\alpha(p\vot1)$, and hence
\[
 \partial H_{\mathrm{ev}}=\alpha(p\vot1).
\]
These identities prove the asserted nullhomotopy.
All its paths come from the same enhanced diagram,
so it restricts coherently along the density cocone
and every internal collision. Evaluation on a closed
boundary state, followed by a closed functional,
therefore cannot detect this kernel on cohomology.

\subsection{A map retaining the fibre homotopy}

Set $Z_E=X_P(\vEnd)$, $W_E=Y_P(\vEnd)$ and
$v=\delta_{P,\vEnd}$. Define its factorization defect
$\operatorname{Def}_P(\vEnd)=\operatorname{Fib}(v)$,
using the cochain convention
\begin{equation}\label{eq:vol2-density-fibre-convention}
 \operatorname{Fib}(u)^n=U^n\oplus V^{n-1},\qquad
 d(x,y)=(dx,ux-dy),\qquad u:U\to V.
\end{equation}
Write $\beta_0=X_P(\beta_L)$ and
$\beta_1=Y_P(\beta_L)$. Their coproduct square gives
$\partial\eta=v\beta_0-\beta_1\delta$.
It induces
\[
 T_\eta:K_{C,P}\to\operatorname{Def}_P(\vEnd),
 \qquad T_\eta(x,y)=(\beta_0x,\beta_1y+\eta x).
\]
The map
\begin{equation}\label{eq:vol2-density-detector}
 q_P=T_\eta\pi:Z\to\operatorname{Def}_P(\vEnd),\qquad
 \pi=(p,-e\delta),\qquad
 q_P=(\beta_0p,\eta p-\beta_1e\delta)
\end{equation}
is a cochain map. Indeed $\partial(-e\delta)=\delta p$
is exactly the fibre equation. If $i:K_{C,P}\to Z$
projects to the first component, $i\pi=p$, while
\[
 R(x,y)=(sy,ey),\qquad dR+Rd=1-\pi i.
\]
Thus $\pi i\simeq1$, and $q_P$ is nullhomotopic if
and only if $T_\eta$ is nullhomotopic.

The two components in \eqref{eq:vol2-density-detector}
record an insertion discrepancy. More explicitly,
let $\ell=\beta_0s$ and $b=\eta s+\beta_1e$.
Then $\partial b=v\ell-\beta_1$, and
\[
 (\beta_0-\ell\delta,\eta-b\delta)=q_P.
\]
For an action $a$, restrict this formula along its
density map and inverse factorization, defining
\begin{equation}\label{eq:vol2-insertion-class}
 \Gamma_a=
 [q_PX_P(\lambda_a)\delta_{P,F_a}^{-1}].
\end{equation}
Its lift $\ell Y_P(\lambda_a)$ is precisely the
ordered insertion lift: on the tuple $(a,a)$,
\eqref{eq:vol2-section-formula} uses the two unit
insertions into $a\circledast a$, followed by
composition in $\vEnd$. The path $b$ is their
coproduct and unit path. Composition of homotopy
squares retains the term
$\eta'\beta_0+\beta'_1\eta$, so this identification
includes the second fibre component.

A nonzero $\Gamma_a$ would force $K_{C,P}\ne0$.
Projection to $Z_E$ alone can lose such a class.
Evaluation loses more: its coproduct square sends
$\operatorname{Def}_P(\vEnd)\vot A$ to the fibre
of $t$, which is zero. Conversely, vanishing of
$q_P$ would leave any part of the density kernel
annihilated by $T_\eta$ undetermined.

\subsection{Equations on the full action bar}

Choose a simplicial-category presentation of
$\mathscr T_C$ and a resolved enriched diagram
$X_a=X_P(F_a)$. Let
$Q(a,b)=C_*(\operatorname{Map}_{\mathscr T_C}(a,b);k)$,
in nonpositive cohomological degrees. The constant
right module $k$ uses its augmentation. The two-sided
bar has terms
\begin{equation}\label{eq:vol2-action-bar}
 B_n=\bigoplus_{a_0,\ldots,a_n}
 Q(a_{n-1},a_n)\otimes\cdots\otimes Q(a_0,a_1)
                                      \otimes X_{a_0}.
\end{equation}
Normalize in the simplicial direction, denoting the
result by $NB_n$, and put
\[
 B=\bigoplus_{n\geq0}NB_n[n],\qquad
 D|_{NB_n[n]}=(-1)^nd_{\mathrm{int}}
                        +\sum_{j=0}^n(-1)^jd_j.
\]
The face sum is absent for $n=0$. Simplicial identities
and compatibility with the internal differential give
$D^2=0$. This bar computes $Z$: resolve the constant
weight by represented right modules, whose augmented
bar has the extra identity degeneracy, then tensor
with the coefficient diagram. All tensors are derived.

Let $p_B:B\to B$ represent $p$, and let
$\iota_n:NB_n\to B$ be the inclusion, of degree $-n$.
A degree-minus-one map $H$ with $DH+HD=p_B$ is
equivalent to maps $h_n=H\iota_n$, of degree $-n-1$,
such that
\begin{equation}\label{eq:vol2-full-primitive-equations}
 \partial h_n=p_B\iota_n-
              \sum_{j=0}^n(-1)^jh_{n-1}d_j,
 \qquad
 \partial h_n=Dh_n+(-1)^nh_nd_{\mathrm{int}}.
\end{equation}
The graded Leibniz rule proves the forward implication;
the same equations on every summand define $H$ on
the direct sum and prove the converse.

For an object $a$ and arrow $f:a\to b$, the first
equations are
\begin{equation}\label{eq:vol2-object-arrow-primitives}
 \partial h_a=p_B\iota_a,\qquad
 \partial h_f=p_B\iota_f-h_bX(f)+h_a.
\end{equation}
The cocone term $p_B\iota_f$ remains after the object
equations are solved. A mapping chain $\gamma$ of
degree $-1$, with $d\gamma=g-f$, gives
\[
 \partial\iota_\gamma=\iota_bX(\gamma)-\iota_g+\iota_f,
 \qquad
 \partial h_\gamma=p_B\iota_\gamma-h_bX(\gamma)+h_g-h_f.
\]
Thus paths between action maps impose their own
degree-minus-three primitives.

Unit insertion supplies part of these equations.
Let $D_2a=a\circledast a$ and put
\[
 w_a=\delta_{P,F_{D_2a}}^{-1}
       ((v_1)_{I_1}\boxtimes(v_2)_{I_2})\delta_{P,F_a},
 \qquad r_a=X_P(v_1)-w_a.
\]
The map $w_a$ is a coefficient map; it has not thereby
become an action morphism $a\to D_2a$. Let $V$ and
$W_0$ be the induced maps on the bar from $D_2$ with
$X(v_1)$ and $w$, retaining their coherences. The
section formula gives $W_0\simeq s\delta$; choose
$p_B=1-W_0$. The natural transformation $v_1:1\to D_2$
has a prism $P$ with $DP+PD=V-1$. On an ordinary
string it inserts $a_j\to D_2a_j$ in slot $j$ with
sign $(-1)^j$ and applies $D_2$ thereafter. Interior
faces cancel in pairs and the end faces give $V-1$.
For mapping-space chains the naturality simplices
supply the same cancellations. Consequently
\begin{equation}\label{eq:vol2-unit-residual}
 p_B=\partial(-P)+R_0,\qquad
 R_0=V-W_0,\qquad
 R_0\iota_a=\iota_{D_2a}r_a.
\end{equation}
The projector is nullhomotopic exactly when $R_0$ is.

The residual equations therefore begin with
\begin{equation}\label{eq:vol2-residual-primitives}
 \partial g_a=\iota_{D_2a}r_a,\qquad
 \partial g_f=R_0\iota_f-g_bX(f)+g_a,
\end{equation}
and continue by \eqref{eq:vol2-full-primitive-equations}.
A specified monoid extension supplies an object
primitive: let $q=\mu_M(f\circledast f)$ and let its
unit and action paths give $\partial\tau=X_P(q)r_a$.
Using the extended density cocone if necessary,
\[
 g_a=\iota_M\tau-\iota_qr_a
\]
has the required differential, since
$\partial\iota_q=\iota_MX_P(q)-\iota_{D_2a}$.
The extension category carries its higher coherences.
Extending them to all action maps is an additional
requirement.

\subsection{Objectwise and infinite compatibility}

Consider a category with initial object $j$ and two
parallel arrows $f,g:a\to b$, whose composites from
$j$ agree. Take $V_j=V_b=0$ and $V_a=k$. Its
normalized bar is
\[
 T=[k\{f,g\}\xrightarrow{(-1,-1)}k\{a\}],
\]
in degrees $-1,0$. Every object inclusion is
nullhomotopic, since the one at $a$ sends $1$ to
$-f$ under a primitive. Yet $H^{-1}(T)=k\{f-g\}$,
so $1_T$ is not nullhomotopic. Adding a constant
diagram $k$ gives a vacuum summand. The projector
onto $T$ annihilates that vacuum and every object
inclusion up to homotopy, while preserving the
parallel-arrow class. This finite example isolates
the missing arrow equations.

There is also an inverse-limit condition. Let
$F_mB$ be the subcomplex of bar length at most $m$.
Let $M_m=\vRH(F_mB,B)$ be the derived $k$-linear mapping
complex in the differential-module category on $U_P$.
Since $B=\operatorname*{colim}_mF_mB$,
\[
 \vRH(B,B)\simeq\operatorname*{Rlim}_mM_m.
\]
The fibre model of $1-\mathrm{shift}$ on
$\prod_mM_m$, and exactness of products of vector
spaces, give
\begin{equation}\label{eq:vol2-density-Milnor}
 0\longrightarrow\lim_m{}^1 H^{-1}(M_m)
 \longrightarrow H^0\vRH(B,B)
 \longrightarrow\lim_mH^0(M_m)\longrightarrow0.
\end{equation}
No vanishing of the first term has been proved here.
Nullhomotopies on finite skeleta must be compatible
before they define one on the whole bar.

For the split projector itself there is a stronger criterion.
The finite skeleta in this criterion retain the full coefficient
complexes and every mapping-space chain; only bar length is bounded.

\begin{proposition}\label{prop:vol2-projector-skeleta}
Let $B=\operatorname*{colim}_m B_m$ be an exhaustive filtered
presentation in the derived category of differential modules on
a block-disjoint open. Suppose $B\simeq Y\oplus K$ and
$p\simeq\operatorname{diag}(0,1_K)$.
Then $K=0$ if and only if $H^q(p i_m)=0$ for every $m,q$,
where $i_m:B_m\to B$ is the colimit map.
In particular, $p i_m\simeq0$ for every $m$ implies $p\simeq0$.
\end{proposition}
\begin{proof}
On each affine ordered chart, filtered colimits of differential
modules are exact, so
\[
 H^q(B)=\operatorname*{colim}_mH^q(B_m).
\]
The stated vanishing makes $H^q(p)$ zero on the image of every
stage, hence zero everywhere. But its restriction to $H^q(K)$
is the identity. Thus $K$ is acyclic on every affine chart.
Those restrictions are jointly conservative, so $K=0$.
The converse and the last assertion follow from the splitting.
\end{proof}

Apply the proposition to $p=1-s\delta$ and $B_m=F_mB$.
A nonzero density kernel is therefore detected on some full
finite bar skeleton. No bound on that skeleton follows, and no
finite truncation of states, normal derivatives or mapping chains
is asserted. The proposition also does not make the whole
$\lim^1$ group in \eqref{eq:vol2-density-Milnor} vanish.

Prescribed paths retain a different obstruction. If
$r:M\to N$ restricts the complex of natural transformations
to a specified extension diagram, its fibre has differential
$d(x,y)=(dx,rx-dy)$. A prescribed cancellation $dh=rp$ gives
the cycle $(p,h)$ in this fibre. It need not vanish when the
absolute class $[p]$ vanishes: if $p=dk$, the remaining class
is represented by the closed element $h-rk$ of $N^{-1}$,
modulo the image of $H^{-1}(M)$. This relative cycle is not a
split projector on a coefficient object. The proposition does
not supply its required path.

\begin{problem}\label{qn:vol2-geometric-recipient}
For the Heisenberg boundary $C_{V_\kappa}$ at $\kappa\ne0$,
exhibit a nonzero class of a binary kernel $K_{C,A|B}$ in
the full action bar, with its object, arrow and higher components.
Determine whether such a class can be an insertion class
\eqref{eq:vol2-insertion-class}.
Over $\mathbb C$, determine whether there are two objects of
$\vMon$ with no common receiving multiplicative factorizing
action. The parallel maps of
Theorem~\ref{thm:vol2-monoid-nonfinality} have no receiving
equalizer; this does not decide the two-object question.
\end{problem}

The tree and nilpotent constructions give actual
families of incoming actions. The density section,
evaluation nullhomotopy and unit prism describe their
images and the remaining equations on the full bar.
For this Heisenberg boundary,
Theorem~\ref{thm:vol2-pointed-sphere} proves that these equations
cannot all have a compatible solution: the pointed-action
functor is not representable, and some density kernel is nonzero.
The kernel criterion and primitive equations still apply to
other boundaries. Over $\mathbb C$,
Theorem~\ref{thm:vol2-monoid-nonfinality} also excludes finality
among multiplicative actions, using the multiplication itself.
Common recipients for two specified action objects remain a
distinct comparison problem.

\endgroup

\chapter{Supported compositions and outer collision classes}
\label{ch:vol2-supported-compositions}

Three quadratic current fields have a connected Gaussian correlation
with a pole on each pair diagonal. Its triangle of contractions already
lies beyond a forest calculation. A second operator product must act on
these poles and on their residue attachments, as well as on the states
left by the first product. The coefficient modules of
Chapter~\ref{ch:vol2-coefficients} provide the domains for those maps.

We begin with the finite Wick formula and its triangle current. Mixed
Green kernels give a second construction, in which each edge has a point
source and a homotopy must retain contact terms. Ordered collision
complexes then show which classes disappear when only consecutive
cuts are kept. The final comparison concerns compositions on supported
field coefficients, with the lower homotopies fixed.

Unless another base is specified, $k$ is a field of characteristic
zero. Polynomial states contain finitely many jets; infinite Taylor
series occur only in the indicated formal variables. Analytic currents
are taken over $\mathbb C$. Complexes are cohomological, with
$K[1]^n=K^{n+1}$, and the differential of a homogeneous map is
$\delta f=d_{\mathrm{out}}f-(-1)^{|f|}fd_{\mathrm{in}}$.

\section{Finite fields and the Gaussian triangle}
\label{sec:vol2-wick-fields}

Allow the Heisenberg level to be the indeterminate $\lambda$. The
polynomial representation of \eqref{eq:vol2-fock-currents} is then
defined over $k[\lambda]$, including at $\lambda=0$:
\[
 V=k[\lambda][x_1,x_2,\ldots],\qquad
 J_{-n}=L_{x_n},\quad J_n=\lambda n\partial_{x_n}\quad(n>0),
 \qquad J_0=0.
\]
Set $J=x_1$, $T=\sum_{n\geq1}n x_{n+1}\partial_{x_n}$ and
$j_r=T^rJ=r!x_{r+1}$. These are raw, undivided jets.
Write $J(z)=\sum_{m\in\mathbb Z}J_mz^{-m-1}$.
For a polynomial $F(j_0,j_1,\ldots)$, define its field $Y(F,z)$
by replacing $j_r$ with $\partial_z^rJ(z)$ and putting every
creation operator $J_{-n}$ to the left of every annihilation
operator $J_n$. The vacuum field is $Y(1,z)=1$.

The commutator of an annihilation part with a creation part is
the expansion of $\lambda/(z-w)^2$ in nonnegative powers of $w/z$.
Consequently the contraction between raw jets is
\begin{equation}\label{eq:vol2-jet-contraction}
 K_{rs}(z,w)
 =\partial_z^r\partial_w^s\frac{\lambda}{(z-w)^2}
 =\frac{\lambda(-1)^r(r+s+1)!}{(z-w)^{r+s+2}}.
\end{equation}
It is a meromorphic coefficient. A chosen ordering of the formal
variables specifies its Laurent expansion.

For polynomial inputs $F_i$ at distinct labels $i=1,\ldots,n$, let
$j_r^{(i)}$ be independent copies of the jet variables and put
\[
 \Gamma_{ij}
 =\sum_{r,s\geq0}K_{rs}(z_i,z_j)
       \partial_{j_r^{(i)}}\partial_{j_s^{(j)}}\qquad(i<j).
\]
After these operators act, colons mean that the surviving variables
are replaced by their fields at the stated positions and all resulting
creation operators are moved to the left, without further contractions.

\begin{proposition}[Polynomial Wick evaluator]
\label{prop:vol2-wick-evaluator}
On each state, use the formal recipient
$V((z_1))\cdots((z_n))$ prescribed by
$|z_1|>\cdots>|z_n|$. Then
\begin{equation}\label{eq:vol2-wick-evaluator}
 Y(F_1,z_1)\cdots Y(F_n,z_n)
 ={:}\exp\left(\sum_{i<j}\Gamma_{ij}\right)
                       \prod_iF_i(j^{(i)}){:}.
\end{equation}
The exponential is finite on each list of polynomial inputs. Before
Laurent expansion, every coefficient is meromorphic with a finite pole
order on each pair diagonal. The formula commutes with differentiation
in the positions and with insertion of the vacuum. Grouping successive
inputs gives the same operator product in a common iterated Laurent
recipient.
\end{proposition}

\begin{proof}
Move each annihilation operator past the creation operators to its
right. Their commutators are scalars, so a term is specified by a set
of disjoint pairs of field slots belonging to different inputs. Each
pair contributes \eqref{eq:vol2-jet-contraction}; the unpaired slots
remain normally ordered. A product of $q$ contraction operators counts
the same unordered set of pairs $q!$ times, cancelled by the
exponential coefficient. Repeated derivatives count the choices among
identical slots. This proves the formula on monomials and hence on
polynomials.

Each contraction removes two slots. There are only finitely many
slots, so the exponential terminates, and each scalar coefficient is
a finite sum of products of the displayed rational functions.
Differentiating the contraction identity differentiates the corresponding
slot; the product rule accounts for all choices. A vacuum slot has no
jet variables, so it has no contractions and contributes the identity.
Finally, the operators $\Gamma_{ij}$ commute. Grouping the factors
partitions the same set of pairs into internal and cross pairs, each
used once. Their exponentials multiply to the single exponential in
\eqref{eq:vol2-wick-evaluator}. Expanding every rational coefficient in
one common order proves the last assertion.
\end{proof}

The meromorphic expression and its Laurent expansions have different
roles. At a collision with separation $u$, the differential-module
operation applied to a Laurent test $f(u)$ is
$\operatorname{pp}_u(f(u)Y(F,u)G)$.
A pole in the test can detect a regular field term. For example,
$\operatorname{pp}_u(u^{-1}Y(1,u)G)=u^{-1}G$.
Thus a vacuum insertion remains visible even though its field has no
singular part before multiplication by the test.

\begin{proposition}[All current modes]
\label{prop:vol2-polynomial-current-modes}
For every polynomial state $F$ and integer $m$,
\begin{equation}\label{eq:vol2-polynomial-current-modes}
 [J_m,Y(F,w)]
 =\lambda\sum_{r\geq0}(m)_{r+1}w^{m-r-1}
                 Y(\partial_{j_r}F,w),
 \qquad (m)_{r+1}=m(m-1)\cdots(m-r).
\end{equation}
The sum is finite, also for negative $m$.
\end{proposition}

\begin{proof}
A single current can contract with only one slot of $F$.
The contraction with $j_r(w)$ is
$\lambda(r+1)!(z-w)^{-r-2}$.
The coefficient of $(z-w)^{r+1}$ in the formal expansion of
$z^m$ about $w$ is $\binom m{r+1}w^{m-r-1}$, for every integer
$m$. Taking the residue in $z$ therefore gives the displayed falling
power. Only the finitely many jet variables occurring in $F$ contribute.
\end{proof}

Write $\langle-\rangle$ for the Gaussian vacuum functional, defined on a product of fields by
applying the product to $1\in V$ and then taking its constant
polynomial coefficient. In \eqref{eq:vol2-wick-evaluator} this retains
exactly the complete pairings. For fields $A_1,\ldots,A_n$, their
connected vacuum coefficient is
\[
 \kappa_n^{\mathrm{vac}}(A_1,\ldots,A_n)
 =\sum_{\pi\in\operatorname{Part}([n])}
   (-1)^{|\pi|-1}(|\pi|-1)!
   \prod_{B\in\pi}\left\langle\prod_{i\in B}A_i\right\rangle.
\]
Products within a block use the inherited order. This definition
specifies an algebraic vacuum functional; a numerical state or a
field-theoretic preparation requires its own construction.

Put $Q(z)={:}J(z)^2{:}$ and $K_{ij}=\lambda/(z_i-z_j)^2$.

\begin{proposition}[Connected quadratic fields]
\label{prop:vol2-wick-cycles}
For $n\geq3$,
\begin{equation}\label{eq:vol2-wick-cycles}
 \kappa_n^{\mathrm{vac}}(Q(z_1),\ldots,Q(z_n))
 =2^{n-1}\sum_{\sigma\in S_{\{2,\ldots,n\}}}
   K_{1,\sigma(2)}K_{\sigma(2),\sigma(3)}\cdots K_{\sigma(n),1}.
\end{equation}
For $n=2$ it is $2K_{12}^2$. In particular,
\begin{equation}\label{eq:vol2-wick-triangle}
 \kappa_3^{\mathrm{vac}}(Q(z_1),Q(z_2),Q(z_3))
 =\frac{8\lambda^3}{z_{12}^2z_{23}^2z_{13}^2},
 \qquad z_{ij}=z_i-z_j.
\end{equation}
These identities hold as meromorphic coefficients.
\end{proposition}

\begin{proof}
The partition formula cancels every disconnected pairing. Indeed,
for a fixed graph with $c$ connected components, its coefficient is
the sum of $(-1)^{|\pi|-1}(|\pi|-1)!$ over partitions of those
components. This sum is $1$ for $c=1$ and $0$ for $c>1$, as follows
by comparing coefficients in $\log(\exp t)=t$.

Normal ordering forbids a contraction within one input. A connected
pairing of two slots at each of at least three vertices is therefore
a cycle. Each unoriented cycle has $2^n$ labelled matchings, since
each vertex can assign its two slots to its neighbours in two ways.
Fixing vertex $1$ and permuting the other vertices counts that cycle
twice, once in each orientation. This gives the factor $2^{n-1}$.
At two vertices there are two bijections between the slots. At three
vertices there is one unoriented cycle and eight labelled matchings,
giving \eqref{eq:vol2-wick-triangle}.
\end{proof}

This triangle has no interaction vertex. It is already a correlation
of composite fields in a free theory. Its nonzero coefficient must be
retained by any comparison covering all polynomial boundary fields.
A calculation restricted to forests does not contain it.

\subsection{Changing the covariance}

The normal-order convention can itself be changed algebraically.
Let $B$ be a commutative $\mathbb Q$-algebra and
$A=B[q_1,\ldots,q_N]$. For symmetric matrices $K,R$ over $B$,
independent of the variables $q_i$, set
\[
 \Gamma_K=\sum_{i,j}K_{ij}\partial_{q_i}\otimes\partial_{q_j},
 \qquad m_K=m\exp(\Gamma_K),\qquad
 \Delta_R=\sum_{i,j}R_{ij}\partial_{q_i}\partial_{q_j},
 \qquad \alpha_R=\exp(\Delta_R/2),
\]
where $m$ is ordinary multiplication. All exponentials terminate on
polynomials. The variables may include separate copies at several
positions; the output operator then includes the cross derivatives
between those copies.

\begin{proposition}\label{prop:vol2-covariance-intertwiner}
For another symmetric constant matrix $S$,
\[
 \alpha_Rm_K=m_{K+R}(\alpha_R\otimes\alpha_R),
 \qquad \alpha_R\alpha_S=\alpha_{R+S}.
\]
\end{proposition}

\begin{proof}
The second-order product rule is
\[
 \tfrac12\Delta_Rm
 =m\bigl(\tfrac12\Delta_R\otimes1+
          1\otimes\tfrac12\Delta_R+\Gamma_R\bigr).
\]
All operators on the right commute. Exponentiating and composing
with $\exp(\Gamma_K)$ proves the first formula. The second follows
from $\Delta_R+\Delta_S=\Delta_{R+S}$.
\end{proof}

An intertwiner of these products does not specify a change of the
vacuum functional. Comparing states requires the functional as well
as the multiplication.

\section{Residue attachments and closed coincidences}
\label{sec:vol2-triangle-attachments}

The triangle has three dependent differences. Nevertheless its
meromorphic coefficient admits an explicit current extension using
only products in independent normal coordinates.
Work over $\mathbb C$, put $u=z_{12}$ and $v=z_{23}$, and use
the distributions $P_j$ of Proposition~\ref{prop:vol2-residue-current}.
Write $R_j=\bar\partial P_j$. Thus
\[
 P_j=\frac{(-1)^{j-1}}{(j-1)!}\partial_u^{j-1}P_1,
 \qquad
 R_j=\frac{\pi(-1)^{j-1}}{(j-1)!}
              \partial_u^{j-1}\delta_0\,d\bar u.
\]

The rational identity
\begin{equation}\label{eq:vol2-triangle-fractions}
 \frac1{u^2v^2(u+v)^2}
 =\frac1{v^4u^2}+\frac1{v^4(u+v)^2}
   -\frac2{v^5u}+\frac2{v^5(u+v)}
\end{equation}
follows by multiplying by the common denominator. It defines the
following extension of the triangle coefficient:
\begin{equation}\label{eq:vol2-triangle-current}
 \mathcal P_\triangle
 =P_4(v)\bigl(P_2(u)+P_2(u+v)\bigr)
       -2P_5(v)\bigl(P_1(u)-P_1(u+v)\bigr).
\end{equation}
Each summand is an external product in either $(u,v)$ or $(u+v,v)$.
Both pairs are linear coordinates, so every product in this formula
is defined.

Every meromorphic Wick coefficient also has a uniform principal-value
realization. For $F=G/Q$ on a smooth complex chart, it is
$[F]=\left.|Q|^{2s}G/Q\right|_{s=0}$, defined by analytic
continuation. This assignment is linear, independent of the
holomorphic presentation, and compatible with holomorphic derivatives;
see Andersson--Wulcan (2018)~\cite[\S2, Example~4.6, Corollary~4.26]{AnderssonWulcan2018}.
It applies to the finite rational coefficients of
Proposition~\ref{prop:vol2-wick-evaluator}, including cycles and
repeated edges. In independent coordinates it agrees with the
external products above, so $\mathcal P_\triangle$ is the
principal value of \eqref{eq:vol2-triangle-fractions}.
Restriction off the poles makes the assignment injective.
On its image, $[F]\diamond[G]=[FG]$ transports associative
meromorphic multiplication. This specified product does not extend
to arbitrary supported currents.

\begin{proposition}\label{prop:vol2-triangle-current}
The current $8\lambda^3\mathcal P_\triangle$ extends
\eqref{eq:vol2-wick-triangle}. Its complete Dolbeault attachment is
\begin{align}\label{eq:vol2-triangle-attachment}
 \bar\partial\mathcal P_\triangle
 &=R_4(v)\bigl(P_2(u)+P_2(u+v)\bigr)
   +P_4(v)\bigl(R_2(u)+R_2(u+v)\bigr)\notag\\
 &\quad-2R_5(v)\bigl(P_1(u)-P_1(u+v)\bigr)
   -2P_5(v)\bigl(R_1(u)-R_1(u+v)\bigr).
\end{align}
Holomorphic differentiation gives the attachments of its jets.
\end{proposition}

\begin{proof}
Away from the three divisors, the principal values are their
ordinary rational functions, so \eqref{eq:vol2-triangle-fractions}
proves the extension assertion. The distributional differential
satisfies the product rule on external products. Apply it separately
in the two coordinate systems to obtain all eight terms in
\eqref{eq:vol2-triangle-attachment}. Holomorphic derivatives commute
with $\bar\partial$ on distributions.
\end{proof}

The displayed extension fixes its contact terms. The rational
identity by itself would not exclude adding a current supported on
the collision. Nor does this construction define the ordinary product
of arbitrary singular distributions. The localization and contact
objects remain separate, as required by
Section~\ref{sec:vol2-supported-tensors}.

\subsection{The finite Koszul complex of a graph}

A closed coincidence gives a different complex. Let $\Gamma$ be
a finite oriented multigraph with vertex set $I$ and edge set $E$;
loops and repeated edges are allowed. Put
\begin{equation}\label{eq:vol2-graph-koszul}
 S=k[z_i:i\in I],\qquad
 K_\Gamma=S\otimes\Lambda(e_a:a\in E),\qquad
 |e_a|=-1,\quad de_{(i,j)}=z_i-z_j.
\end{equation}
No difference is inverted. A loop has zero differential.

Choose a spanning forest $T$ and one root in each component.
For each nonforest edge $a$, express its difference as a signed
sum along the unique forest path and set
\[
 c_a=e_a-\sum_{t\in\operatorname{path}(a)}\epsilon_t e_t.
\]
For a loop the sum is empty. The vectors $c_a$ form a basis of the
cycle space of the graph and satisfy $dc_a=0$.
Root coordinates and the forest differences $w_t$ form a linear
coordinate system: removing a leaf successively reconstructs every
vertex coordinate. With $B=k[z_{\mathrm{roots}}]$, the complex is
therefore
\[
 K_\Gamma=B[w_t:t\in T]\otimes\Lambda(e_t:t\in T)
                          \otimes\Lambda(c_a:a\notin T),
 \qquad d=\sum_tw_t\partial_{e_t}.
\]
Let $p$ set all $w_t,e_t$ to zero, and let $i$ include the resulting
root-and-cycle algebra.

\begin{proposition}[Forest contraction]
\label{prop:vol2-graph-contraction}
Let $N$ count total polynomial degree in the $w_t$ plus exterior
degree in the $e_t$, excluding the $c_a$. Define
\[
 H=\frac1N\sum_te_t\partial_{w_t}\quad(N>0),
 \qquad H=0\quad(N=0).
\]
Then
\[
 dH+Hd=1-ip,\qquad H^2=Hi=pH=0,
 \qquad H^*(K_\Gamma)=B\otimes\Lambda(c_a:a\notin T).
\]
The homotopy and the chosen section are $B$-linear. The augmentation
$p$ is $S$-linear for the quotient action on its target.
\end{proposition}

\begin{proof}
The graded commutator of $\sum_tw_t\partial_{e_t}$ and
$\sum_te_t\partial_{w_t}$ is the Euler operator in exactly the
$w_t,e_t$ variables. Both preserve $N$. Dividing by $N$ proves the
contraction identity on positive degree; on degree zero $ip$ is the
identity. The numerator of $H$ squares to zero because the exterior
multiplications anticommute and the ordinary derivatives commute.
Evaluation gives $pH=0$, and $Hi=0$ follows from the absence of
forest variables in the image of $i$. The target has zero
differential, proving the cohomology formula.

The formulas leave root coordinates unchanged. In contrast,
$H(w_t)=e_t$ whereas $w_tH(1)=0$, so $H$ is generally not
$S$-linear. The augmentation respects $S$ because setting all
forest differences to zero is precisely its quotient action.
\end{proof}

For the triangle the surviving generator is
$e_{13}-e_{12}-e_{23}$. Identifying the endpoints of the path
$12,23$ instead gives the closed generator $e_{12}+e_{23}$.
Keeping the original edge generators makes vertex identifications
maps of Koszul complexes. The choice of spanning forest calculates
their cohomology without removing those cycles.

The stable support complex of
Lemma~\ref{lem:vol2-redundant-support} is unchanged by a redundant
support equation. The finite Koszul complex retains its excess
cycle. Neither construction identifies that cycle with the
meromorphic triangle current. A graph loop, for example, gives a
closed exterior generator in \eqref{eq:vol2-graph-koszul}, but no
rational function $1/(z_i-z_i)$.

\section{Mixed kernels and their forest homotopies}
\label{sec:vol2-mixed-forest-homotopies}

To compare a mixed propagator with a boundary coefficient, its
point source must remain in the differential. Work on
$\mathbb R_t\times\mathbb C_z$, where $z=x+iy$, with orientation
$dt\wedge dx\wedge dy$. Partial forms are generated by $dt,d\bar z$,
and
\[
 D=dt\,\partial_t+d\bar z\,\partial_{\bar z},\qquad
 \partial_{\bar z}=\tfrac12(\partial_x+i\partial_y).
\]
Distribution coefficients are measured against $dt\,dx\,dy$.
Put $r=(t^2+|z|^2)^{1/2}$ and define
\begin{equation}\label{eq:vol2-mixed-axial-kernels}
 P=\frac{t\,d\bar z-2\bar z\,dt}{4\pi r^3},\qquad
 A=-\frac{\delta(t)\,dt}{\pi z},\qquad
 \Delta=\delta^{(3)}(t,z)\,dt\wedge d\bar z.
\end{equation}
The coefficients of $P$ are locally integrable. The coefficient of
$A$ is the external product of the one-dimensional delta and the
locally integrable Cauchy kernel.

\begin{lemma}\label{lem:vol2-kernel-source}
One has $DP=DA=\Delta$ as currents.
\end{lemma}

\begin{proof}
The coefficient of $dt\wedge d\bar z$ in $DP$ is
\[
 \frac1{4\pi}\left(\partial_t\frac{t}{r^3}
               +2\partial_{\bar z}\frac{\bar z}{r^3}\right).
\]
The imaginary angular terms cancel. The remainder is $1/(4\pi)$
times the divergence of $(t,x,y)/r^3$, whose flux through a sphere
is $4\pi$. It is therefore $\delta^{(3)}$ distributionally.
For $A$, use $\partial_{\bar z}(1/(\pi z))=\delta^{(2)}(z)$.
The minus sign in $A$ cancels the sign from moving $d\bar z$
past $dt$.
\end{proof}

Let $H(t)$ be the Heaviside distribution. Its value at zero is
irrelevant to the current. Define
\begin{equation}\label{eq:vol2-kernel-primitive}
 S(t,z)=\frac{H(t)-\tfrac12-\tfrac{t}{2r}}{\pi z}.
\end{equation}

\begin{proposition}\label{prop:vol2-kernel-primitive}
The function $S$ is locally integrable and satisfies
\[
 DS=P-A,\qquad
 \int_{\mathbb R}P=\int_{\mathbb R}A=-\frac1{\pi z}.
\]
At fixed $z$, its two tails are $O(|\bar z|/t^2)$.
\end{proposition}

\begin{proof}
On $t>0$ and $t<0$ respectively,
\[
 S_+=\frac{1-t/r}{2\pi z},\qquad
 S_-=\frac{-1-t/r}{2\pi z}.
\]
Both have derivatives
$\partial_tS_\pm=-\bar z/(2\pi r^3)$ and
$\partial_{\bar z}S_\pm=t/(4\pi r^3)$.
Their jump across $t=0$ is $1/(\pi z)$, supplying the
$\delta(t)/(\pi z)$ term in $P-A$. For fixed $t\ne0$ the
numerator vanishes at $z=0$, so there is no further axial delta.
The function is homogeneous of degree $-1$ with bounded angular
factor near the origin. Its first derivatives have degree $-2$.
A distribution supported only at the origin is a finite sum of
delta derivatives, whose degrees are at most $-3$. Thus no
additional point-supported term occurs, and the identity holds
on the whole space.

Integration of the $dt$ coefficient gives
\[
 -\frac{\bar z}{2\pi}\int_{-\infty}^{\infty}
          (t^2+|z|^2)^{-3/2}\,dt=-\frac1{\pi z}.
\]
Its absolute integral is locally integrable in $z$, which justifies
the distributional pushforward. The axial integral has the same
value. Expanding $t/r$ at either end gives the tail estimate.
\end{proof}

Multiplication by $-\pi$ converts this pushforward to the Cauchy
coefficient $1/z$, whose contour residue is one under
$(2\pi i)^{-1}\oint dz$. A physical coupling or pairing must
include that normalization. Every fixed finite holomorphic or
normal derivative commutes with the current identity.

\subsection{Cutoffs and boundary conditions}

For a smooth cutoff $\chi$ equal to one near the origin,
\begin{equation}\label{eq:vol2-kernel-cutoff}
 D(\chi P)=\Delta+D\chi\wedge P,\qquad
 D(\chi S)=\chi(P-A)+D\chi\,S.
\end{equation}
The annular terms belong to the supported comparison. If a
composition needs a proper projection, choose cutoffs giving that
property and retain these terms in its localization complex.

A smooth regulator follows from
$D^*=-\partial_t\iota_{dt}-4\partial_z\iota_{d\bar z}$.
For $0<\varepsilon<L$, set
\[
 K_\ell=(4\pi\ell)^{-3/2}e^{-r^2/(4\ell)},\qquad
 P_{\varepsilon,L}
 =D^*\left(\int_\varepsilon^L K_\ell\,d\ell
                           \;dt\wedge d\bar z\right).
\]
Since $DD^*+D^*D$ is the negative Euclidean Laplacian,
the heat equation gives
\[
 DP_{\varepsilon,L}=(K_\varepsilon-K_L)\,dt\wedge d\bar z.
\]
Both heat endpoints occur in its differential.

On a half-line the relative de Rham complex has Dirichlet
zero-forms and Neumann one-forms. Writing
$G_\ell(a)=(4\pi\ell)^{-1/2}e^{-a^2/(4\ell)}$, their kernels are
\[
 K^0_\ell(t,s)=G_\ell(t-s)-G_\ell(t+s),\qquad
 K^1_\ell(t,s)=G_\ell(t-s)+G_\ell(t+s).
\]
Differentiation gives $\partial_tK^0_\ell=-\partial_sK^1_\ell$.
Integration by parts, using the relative boundary value, proves
cochain compatibility. Tensoring with the tangential Dolbeault
complex preserves the identity. The polarization of a particular
field theory determines whether its normal complex is relative
or absolute.

\subsection{Forest cubes and contact terms}

Let $F$ be an oriented forest on a finite labelled set. For
$e=(i,j)$ put $t_e=t_i-t_j$ and $z_e=z_i-z_j$.
Choose one root per connected component. The root coordinates
and edge differences are a linear coordinate system with
determinant $\pm1$, by successive leaf deletion. Hence the
pullbacks $P_e,A_e,S_e,\Delta_e$ for different forest edges
have independent normal variables. Their ordered products,
including any fixed finite normal derivatives, are external
products of distributions in these coordinates.

\begin{theorem}\label{thm:vol2-forest-cubes}
On every labelled forest, the mixed and axial decorations are
joined by homotopy cubes compatible with contact contraction,
spectators, and products or graftings that remain forests.
This includes the edge $13$ with spectator $2$ and the crossed
pair of edges $13,24$.
\end{theorem}

\begin{proof}
Introduce $s_e\in[0,1]$ for each edge, with parameter forms
written before spatial forms, and set
\[
 \mathbb P_e=A_e+s_e(P_e-A_e)+ds_e\,S_e.
\]
Then $(D+d_s)\mathbb P_e=\Delta_e$: the term
$ds_e(P_e-A_e)$ cancels $-ds_eDS_e$.
Independence of the edge variables defines every forest product.
Its differential is the signed sum replacing one edge by its
contact current. At $s_e=0$ or $1$ the edge becomes $A_e$ or
$P_e$. Restrictions to two faces commute, giving the higher
face identities. The diagonal interval varies all edges together.

Contact contraction identifies the two endpoints of an edge;
the remaining differences are coordinates of the contracted
forest. The distributional equations persist under finite normal
differentiation and retain its contact jets. Adding an isolated
vertex adds a spectator, and the empty product is the coefficient
unit. Edge reorderings have the usual form signs.
\end{proof}

For two edges define $H_{e,f}=S_eP_f-A_eS_f$. Since
$|P|=|A|=1$ and $|S|=0$, the product rule gives
\begin{equation}\label{eq:vol2-two-edge-contact}
 P_eP_f-A_eA_f
 =DH_{e,f}+\Delta_eS_f-S_e\Delta_f.
\end{equation}
The last two terms express the common point sources of the
kernels. They remain even though $P-A$ is exact.

At three positions put $z_{12}=u$, $z_{23}=v$, $z_{13}=u+v$.
The axial products have the same time delta and oriented time
two-form. Their scalar coefficients satisfy
\[
 \frac1{uv}-\frac1{u(u+v)}-\frac1{v(u+v)}=0.
\]
Each term is locally integrable on $\mathbb C^2$ in its
independent linear coordinates. Thus the rational identity is
also a distributional identity. Define
$\mathcal K=H_{12,23}-H_{12,13}-H_{13,23}$. Applying
\eqref{eq:vol2-two-edge-contact} to the three pairs gives
\begin{align}\label{eq:vol2-three-channel-current}
 P_{12}P_{23}-P_{12}P_{13}-P_{13}P_{23}
 &=D\mathcal K+\mathcal C,\notag\\
 \mathcal C
 &=\Delta_{12}S_{23}-S_{12}\Delta_{23}
   -\Delta_{12}S_{13}+S_{12}\Delta_{13}\notag\\
 &\quad-\Delta_{13}S_{23}+S_{13}\Delta_{23}.
\end{align}
Every product here has independent normal variables.

These are scalar coefficient currents on the mixed space.
Decorating them by fields requires a map preserving their
contact terms, translations, partial vacua and supported
attachments. A cycle has dependent edge differences, so its
mixed current product needs an additional extension. The
meromorphic triangle of \eqref{eq:vol2-triangle-current} supplies
a specified boundary coefficient, without identifying it with
such a mixed product.

\section{Classes discarded by consecutive cuts}
\label{sec:vol2-discarded-collisions}

Consecutive cuts preserve an order on the labels. A collision can
join labels that are separated in that order. To measure this
difference, we compare two coefficient complexes before introducing
a map from coefficients to fields.

Let $R_n=k[z_1,\ldots,z_n]$ and
$L_n=R_n[(\prod_{i<j}z_{ij})^{-1}]$.
An ordered partition $P=(B_1|\cdots|B_r)$ has coefficient module
\begin{equation}\label{eq:vol2-partition-coefficient}
 M_P=R_n[z_{ij}^{-1}:i,j\text{ lie in different blocks of }P]
       \subset L_n.
\end{equation}
Equal subrings attached to different ordered partitions remain
different summands. Tensor every term with the common right
density $\omega_n=\omega_{\mathbb A^n}$; for a constant density
$\nu$, $(f\nu)\partial_i=-(\partial_i f)\nu$.
We suppress this common factor in the coordinate formulas.

Write $e_P=\xi_{B_1}\cdots\xi_{B_r}$, where
$|\xi_B|=1-|B|$. The full complex $C_n^{\mathrm{all}}$ has
term $M_Pe_P$ in degree $r-n$. Its differential splits a block
$B_i=A\sqcup C$ into two nonempty ordered parts, with coefficient
the localization inclusion and sign
\begin{equation}\label{eq:vol2-partition-sign}
 (-1)^{\sum_{j<i}(1-|B_j|)+|A|+1}\epsilon(A,C).
\end{equation}
Here $\epsilon(A,C)$ is the sign of the concatenated increasing
lists of $A,C$ relative to the increasing list of $B_i$.
A partition is \emph{consecutive} when its blocks are intervals
in increasing order. Projection to those summands gives the
consecutive-cut complex $C_n^{\mathrm{cut}}$. Its kernel is
\begin{equation}\label{eq:vol2-discarded-complex}
 F_n^{r-n}=\bigoplus_{\substack{P\text{ has }r\text{ blocks}\\
                                  P\text{ is not consecutive}}}M_Pe_P.
\end{equation}
There is no one-block term in $F_n$.

\begin{lemma}\label{lem:vol2-partition-complex}
The discarded summands form a subcomplex and $d^2=0$.
After extending every coefficient to $L_n$, the full complex is
the oriented cellular chain complex of the permutahedron of
dimension $n-1$, with chain degree negated. For $n\geq3$, the
map from discarded two-block terms to three-block terms is
injective, both after localization and on the actual modules.
\end{lemma}

\begin{proof}
If a refinement is consecutive, merging its two adjacent new
blocks recovers a consecutive partition. Thus a nonconsecutive
partition has only nonconsecutive refinements. Two splits in
different blocks cancel by the odd Leibniz sign. Two successive
splits of the same block have the same density permutation and
opposite cobar signs. Hence $d^2=0$.

The permutahedron is the convex hull of the permutations of
$(1,\ldots,n)$. A face is specified by an ordered partition;
its dimension is $n-r$, and splitting one block gives a boundary
face. Orient the block factor by the increasing labels, then
orient the product in block order. Its boundary orientation is
exactly \eqref{eq:vol2-partition-sign}. This identifies the
fully localized incidence complex.

The permutahedron is a ball. In its cellular complex the kernel
on facets is the image of the unique top cell, spanned by
$d\xi_{[n]}$. Every facet has coefficient $1$ or $-1$ in that
boundary. The discarded facet space omits all $n-1$ consecutive
facets, so it contains no nonzero multiple of the top boundary.
This proves injectivity over $k$, hence over $L_n$. Each actual
coefficient module embeds in $L_n$, giving the last assertion.
\end{proof}

\subsection{The outer principal-part module}

Put $u=z_{1n}$. Retain in $d\xi_{[n]}$ only those facets which
keep $1,n$ together:
\begin{equation}\label{eq:vol2-outer-numerator}
 T_n=\sum_{\substack{A\sqcup C=[n],\ A,C\ne\varnothing\\
                        1,n\text{ lie in the same part}}}
          (-1)^{|A|+1}\epsilon(A,C)\xi_A\xi_C,
 \qquad b_n=dT_n.
\end{equation}
These facets are all discarded. The degrees of $T_n$ and $b_n$
are $2-n$ and $3-n$, respectively.

\begin{theorem}[Outer classes in every arity]
\label{thm:vol2-outer-injection}
For every $n\geq3$, the formula
\begin{equation}\label{eq:vol2-outer-injection}
 \omega_n\otimes_{R_n}(R_n[u^{-1}]/R_n)
       \lhook\joinrel\longrightarrow H^{3-n}(F_n),
 \qquad [f]\longmapsto[fb_n]
\end{equation}
defines an injective map of right differential modules.
Thus no arity-independent lower bound exists for the degrees
of nonzero cohomology of the discarded complexes.
\end{theorem}

\begin{proof}
Consider a three-block partition which still keeps $1,n$
together. Its two facet coarsenings both occur in $T_n$, so
their contributions to $dT_n$ cancel. Every nonzero coordinate
of $b_n$ therefore separates $1,n$ and has a coefficient module
in which $u$ is invertible. Consequently $fb_n$ is an allowed
cochain for every $f\in R_n[u^{-1}]$, and it is closed.
If $f\in R_n$, the primitive $fT_n$ is allowed, so the formula
factors through principal parts.

Conversely, suppose $fb_n=dv$ for an actual discarded facet
cochain $v$. Over $L_n$, the injectivity in
Lemma~\ref{lem:vol2-partition-complex} forces $v=fT_n$.
Choose a facet $P$ occurring in $T_n$. Its module does not
invert $u$, and unique factorization gives
$M_P\cap R_n[u^{-1}]=R_n$. The corresponding coordinate of
$v$ is $\pm f$, so $f$ is regular. Such a facet exists for
$n\geq3$, and the same injectivity shows $b_n\ne0$.

The coefficients of $b_n$ are constants. Differentiating
$fb_n$ therefore differentiates only $f$. Restoring the common
right density proves differential-module linearity. Every normal
derivative of the simple principal part is retained.
\end{proof}

For four positions the numerator is
\begin{equation}\label{eq:vol2-outer-four-numerator}
 \begin{split}
 T_4={}&-\xi_{124}\xi_3+\xi_{134}\xi_2-\xi_{14}\xi_{23}\\
       &-\xi_2\xi_{134}-\xi_{23}\xi_{14}+\xi_3\xi_{124}.
 \end{split}
\end{equation}
Thus $z_{14}^{-1}dT_4$ is a nonzero class in $H^{-1}(F_4)$.
Its fully localized primitive $z_{14}^{-1}T_4$ has a forbidden
pole in each displayed facet coefficient. The obstruction is
precisely that failure to lie in the specified module.

The relation with the cut class also fixes a sign and a degree.
The degreewise exact sequence is
\[
 0\longrightarrow F_n\longrightarrow C_n^{\mathrm{all}}
                  \longrightarrow C_n^{\mathrm{cut}}\longrightarrow0.
\]
Its cut top boundary is
\[
 U_n=d_{\mathrm{cut}}\xi_{[n]}
    =\sum_{p=1}^{n-1}(-1)^{p+1}\xi_{[1,p]}\xi_{[p+1,n]}.
\]
Every term separates $1,n$, so $fU_n$ is an allowed cut
cocycle for every $f\in R_n[u^{-1}]$.

\begin{proposition}\label{prop:vol2-outer-connecting}
The connecting homomorphism sends
\[
 H^{2-n}(C_n^{\mathrm{cut}})\longrightarrow H^{3-n}(F_n),
 \qquad [fU_n]\longmapsto-[fb_n].
\]
\end{proposition}

\begin{proof}
Write $d\xi_{[n]}=U_n+T_n+V_n$, where $V_n$ consists of
discarded facets separating $1,n$. Lift $fU_n$ to the same
cochain in the full complex. Its differential is
$d(fU_n)=-fb_n-d(fV_n)$. The latter term is a legal boundary
in $F_n$, since every facet in $V_n$ inverts $u$.
\end{proof}

This comparison concerns restriction of all-subset cobar
coefficients to a chosen order and cobar after restriction to
cuts. Its kernel does not refute bar--cobar duality on either
of its stated categories. A map from these coefficient classes
to a field-action complex remains additional data.

\subsection{Global cohomology at four positions}

The full four-position permutahedron has $14$ facets, $36$
edges and $24$ vertices. Removing consecutive partitions removes
three facets, three edges and the ascending vertex. Hence
\begin{equation}\label{eq:vol2-four-discarded-matrices}
 F_4:\qquad \bigoplus_{11}M_P\xrightarrow{A}
                 \bigoplus_{33}M_Q\xrightarrow{B}L^{23},
 \qquad L=L_4,
\end{equation}
in degrees $-2,-1,0$. The fully localized ranks are $11,22$.
For the second rank, the edges connect all remaining vertices:
adjacent swaps increasing inversion number connect every
nonascending permutation to the reverse permutation without
passing through the removed vertex. Thus over $L$,
$\ker B=\operatorname{im}A$ and $\ker A=0$.

For an unordered pair $h$, write
\[
 L_h=R_4[z_{ab}^{-1}:\{a,b\}\ne h],\qquad W_h=L/L_h.
\]
A three-block term has a unique unsplit pair $h$ and coefficient
$L_h$. Let $I_h$ be the discarded facets keeping $h$ together
and $J_h$ the three-block terms keeping it together. Denote the
corresponding submatrix of $A$ by $K_h$.

\begin{lemma}\label{lem:vol2-four-pair-ranks}
The sizes and ranks are
\[
 \begin{array}{c|rrrrrr}
 h&12&13&14&23&24&34\\\hline
 |I_h|&4&5&6&4&5&4\\
 |J_h|&5&6&6&5&6&5\\
 \operatorname{rank}K_h&4&5&5&4&5&4.
 \end{array}
\]
Only $K_{14}$ has a kernel. In the order of the six terms in
\eqref{eq:vol2-outer-four-numerator}, that kernel is generated by
$(-1,1,-1,-1,-1,1)$.
\end{lemma}

\begin{proof}
Use \eqref{eq:vol2-partition-sign}. In the displayed order,
\[
 K_{14}=\begin{pmatrix}
 1&0&-1&0&0&0\\0&1&1&0&0&0\\-1&0&0&1&0&0\\
 0&0&0&-1&1&0\\0&-1&0&0&0&1\\0&0&0&0&-1&-1
 \end{pmatrix}.
\]
Its equations determine every coordinate from the first. For
the other pairs, order the partitions lexicographically by
their increasing block lists. The rows for $12,23,34$ give
respectively
\[
 \begin{array}{c|ccccc}
 12&-v_1&-v_2&v_2+v_3&v_1-v_4&-v_3+v_4\\
 23&v_1&-v_2&v_2-v_3&-v_1+v_4&v_3-v_4\\
 34&v_1&v_2&-v_2+v_3&-v_1-v_4&-v_3+v_4.
 \end{array}
\]
Each zero system forces all four coordinates to vanish. For
$13$ the rows are
$-v_1,v_1-v_2,-v_3,v_3+v_4,v_2+v_5,-v_4-v_5$;
for $24$ they are
$v_1,v_2,-v_1-v_4,v_3+v_4,-v_2+v_5,-v_3-v_5$.
These force all five coordinates to vanish. Every elimination
uses only $1,-1$, so the same conclusions hold with coefficients
in $W_h$.
\end{proof}

\begin{theorem}\label{thm:vol2-four-outer-cohomology}
Globally on the standard affine chart,
\[
 H^{-2}(F_4)=0,\qquad
 H^{-1}(F_4)=\omega_4\otimes_{R_4}
                    (R_4[z_{14}^{-1}]/R_4).
\]
The second identification sends $[f]$ to $[fdT_4]$.
\end{theorem}

\begin{proof}
The first assertion follows from facet injectivity. Suppress
the density for the second. A degree-minus-one cocycle $w$ has
a unique fully localized primitive $v\in L^{11}$, with $Av=w$.
Requiring the $J_h$ coordinates of $w$ to lie in $L_h$ gives
\[
 K_h(v_i\bmod L_h)_{i\in I_h}=0.
\]
For $h\ne14$, Lemma~\ref{lem:vol2-four-pair-ranks} forces each
indicated coordinate of $v$ to be regular along $h$. Unique
factorization gives, inside $L$,
\[
 M_P=\bigcap_{h\text{ internal to }P}L_h.
\]
Putting $u=z_{14}$, we obtain $v_P\in M_P[u^{-1}]$.
If $P$ separates $1,4$, this is already $M_P$, and its
coordinate is a legal boundary. The six remaining principal
parts are proportional to the coefficients of $T_4$. Thus
$H^{-1}$ is the intersection, inside $W_{14}$, of the images
of $M_P[u^{-1}]/M_P$ for these six facets.

To compute that intersection at the other collision divisors
as well, use
\[
 u=z_{14},\quad a=z_{12},\quad b=z_{23},\quad
 c=a+b=z_{13},\quad y=z_4,
 \qquad D=k[y,a,b].
\]
The remaining differences are $z_{24}=u-a$, $z_{34}=u-c$.
Finite negative Laurent expansion at $u=0$ identifies
\[
 W_{14}=\bigoplus_{j\geq1}D[(abc)^{-1}]u^{-j}.
\]
Indeed, expand $(u-a)^{-1},(u-c)^{-1}$ only to the finite
order required by the pole in $u$. A zero negative part means
that the numerator is divisible by that power of $u$, hence
the original fraction belongs to $L_{14}$.

The three distinct facet images, from $124|3$, $134|2$ and
$14|23$, are respectively
\[
 \bigoplus_{j\geq1}D[(bc)^{-1}]u^{-j},\quad
 \bigoplus_{j\geq1}D[(ab)^{-1}]u^{-j},\quad
 \bigoplus_{j\geq1}D[(ac)^{-1}]u^{-j}.
\]
The reversed facets repeat them. For example, $M_{124|3}$
inverts $b,c,u-c$, whose negative Laurent coefficients invert
only $b,c$; all such coefficients occur. The irreducibles
$a,b,c=a+b$ are pairwise nonassociated in $D$, so the
intersection of the three coefficient rings is $D$.
The principal-part intersection is therefore
$\bigoplus_{j\geq1}Du^{-j}=R_4[u^{-1}]/R_4$.
The identification is $[f]\mapsto[fdT_4]$, already proved
differential-linear in Theorem~\ref{thm:vol2-outer-injection}.
\end{proof}

Degree zero has a finite global presentation. Let $G$ be the
permutahedron graph on $S_4\setminus\{1234\}$ and choose
$\sigma_0=4321$. Label each edge by its swapped pair $h$ and
define
\[
 \mathcal G_{R_4}=
 \frac{\displaystyle\bigoplus_{\sigma\ne\sigma_0}L e_\sigma}
      {\displaystyle\sum_{\{\sigma,\tau\}\in E(G)}
            L_{h(\sigma,\tau)}(e_\sigma-e_\tau)},
 \qquad e_{\sigma_0}=0.
\]
Then
\begin{equation}\label{eq:vol2-four-H0}
 H^0(F_4)=\omega_4\otimes_{R_4}(L\oplus\mathcal G_{R_4}).
\end{equation}
To see this, orient the singleton density basis so that $B$
is the graph incidence map. Its image lies in the kernel of
the sum of vertex coordinates. Inclusion at $\sigma_0$ splits
that sum. In its kernel, the basis
$e_\sigma-e_{\sigma_0}$ gives exactly the displayed relations.
All basis changes are constant and preserve the differential
module structure. This presents the global module by $22$
generator modules and $33$ relation modules; it does not
decompose its extensions across all diagonal intersections.

\subsection{The crossed four-position chart}

On the open set separating $\{1,3\}$ from $\{2,4\}$, put
\[
 \begin{gathered}
 u=z_{13},\quad v=z_{24},\quad
 A=R_4[(z_{12}z_{14}z_{23}z_{34})^{-1}],\\
 A_u=A[u^{-1}],\quad A_v=A[v^{-1}],\quad L=A[(uv)^{-1}].
 \end{gathered}
\]
The two remaining collision coordinates $u,v$ are independent.
The coefficient types in \eqref{eq:vol2-four-discarded-matrices}
become
\[
 A^2\oplus A_u^3\oplus A_v^3\oplus L^3
 \longrightarrow A_u^6\oplus A_v^6\oplus L^{21}
 \longrightarrow L^{23}.
\]

A nonzero scalar entry $a$ between two identical coefficient
modules can be cancelled. Isolating that entry by triangular
row and column operations gives a contractible two-term
summand with homotopy $a^{-1}$. The remaining differential is
the Schur complement
$D'_{rc}=D_{rc}-D_{rj}a^{-1}D_{ic}$.
Every entry still composes allowed localization inclusions.
These changes commute with differential operators.

Order the partitions lexicographically by increasing block
lists. Cancel the first available scalar between identical
modules, using lower degree first, then row, then column.
The nine cancellations from degree $-2$ have types
\[
 A_u,L,A_v,A_v,L,A_u,A_u,L,A_v
\]
and pivots $1,1,-1,-1,1,1,1,1,1$.
The next eighteen cancellations, from degree $-1$, have type
$L$ and pivot $1$. The remaining complex is
\[
 A^2\xrightarrow{M_0}A_u^3\oplus A_v^3
                              \xrightarrow{M_1}L^5.
\]
In its retained middle order $(A_u,A_v,A_v,A_u,A_u,A_v)$,
\[
 M_0=\begin{pmatrix}1&0\\-1&0\\0&1\\0&-1\\-1&1\\1&-1\end{pmatrix},
 \qquad
 M_1=\begin{pmatrix}
 1&1&0&0&0&0\\0&0&1&1&0&0\\0&-1&-1&0&1&0\\
 -1&0&0&-1&0&1\\0&0&0&0&-1&-1
 \end{pmatrix}.
\]
Use the basis matrices
\[
 S=\begin{pmatrix}
 1&0&1&0&0&0\\0&0&0&1&0&1\\0&0&0&0&-1&0\\
 0&-1&0&0&0&0\\-1&1&0&0&0&0\\0&0&0&-1&1&0
 \end{pmatrix},\qquad
 W=\begin{pmatrix}
 1&0&1&1&1\\0&-1&0&0&0\\-1&1&0&-1&0\\
 -1&1&-1&0&0\\1&-1&0&0&0
 \end{pmatrix}.
\]
Both are invertible. The first three columns of $S$ have
coefficient $A_u$, the last three $A_v$, so this basis change
respects the modules. Matrix multiplication gives
\begin{equation}\label{eq:vol2-crossed-canonical}
 S^{-1}M_0=
 \begin{pmatrix}1&0\\0&1\\0&0\\-1&0\\0&-1\\0&0\end{pmatrix},
 \qquad
 W^{-1}M_1S=
 \begin{pmatrix}
 1&0&0&1&0&0\\0&1&0&0&1&0\\0&0&1&0&0&0\\
 0&0&0&0&0&1\\0&0&0&0&0&0
 \end{pmatrix}.
\end{equation}

\begin{theorem}[Crossed local form]
\label{thm:vol2-crossed-collision}
On this open set, $F_4$ is chain-homotopy equivalent by right
differential-module maps to
\begin{equation}\label{eq:vol2-crossed-splitting}
 L[0]\oplus[A_u\to L]\oplus[A_v\to L]
       \oplus[A\to A_u\oplus A_v\to L]^{\oplus2}.
\end{equation}
The two-term complexes have degrees $-1,0$; the three-term
complexes have degrees $-2,-1,0$. In particular,
\[
 H^{<0}(F_4)=0,\qquad
 H^0(F_4)=\omega_4\otimes_A
 \left(L\oplus L/A_u\oplus L/A_v
              \oplus\bigl(L/(A_u+A_v)\bigr)^{\oplus2}\right).
\]
\end{theorem}

\begin{proof}
The first two columns and their fourth and fifth partners
in \eqref{eq:vol2-crossed-canonical} give two complexes with
maps $a\mapsto(a,-a)$ and $(f,g)\mapsto f+g$.
Columns three and six give the two remaining two-term
complexes. The last target row gives $L[0]$.
All preceding cancellations were chain homotopy equivalences.

The maps $A_u\to L$ and $A_v\to L$ are injective. In a
three-term summand the middle kernel is $(f,-f)$ with
$f\in A_u\cap A_v=A$. The last equality follows by unique
factorization in the localization of $R_4$: the independent
prime factors $u,v$ remain nonassociated and neither has been
inverted in $A$. Only the displayed cokernels remain.
\end{proof}

The codimension-two module $L/(A_u+A_v)$ is nonzero. Near
$u=v=0$, its class $1/(uv)$ has a negative power in both
independent normals, which cannot occur in a sum of terms
with poles in only one of them. The crossed chart has retained
these two copies of the collision module. It has inverted
$z_{14}$, so its vanishing negative cohomology agrees with the
global outer class, whose support is $z_{14}=0$.

At a generic point of that outer divisor, where all other
differences are invertible, set
$u=z_{14}$ and $A=R_4[(z_{12}z_{13}z_{23}z_{24}z_{34})^{-1}]$.
The same cancellations give
\[
 [A\to A[u^{-1}]]_{[-2,-1]}
 \oplus[A\xrightarrow{(1,-1)}A[u^{-1}]^2]_{[-1,0]}.
\]
Thus $H^{-1}=A[u^{-1}]/A$ and
$H^0=A[u^{-1}]\oplus A[u^{-1}]/A$ there. These are full
principal-part modules, including every normal derivative.

\section{Independent normals and the full diagonal}
\label{sec:vol2-full-diagonal-support}

The crossed calculation retains proper collision strata. The
full diagonal has a simpler normal calculation, whose shift
depends on the number of independent equations. We compute it
before imposing any identification with a field coefficient.

Let $X=\mathbb A^b_k$ have coordinates $y_1,\ldots,y_b$, let
$Y=X\times\mathbb A^c_k$ have normal coordinates
$u_1,\ldots,u_c$, and let $i:X\hookrightarrow Y$ be the zero
section. For a left $D_X$-module $L$, its normal direct image is
\[
 i_+L=L[\eta_1,\ldots,\eta_c],\qquad
 \partial_{u_j}=\eta_j,\qquad u_j=-\partial_{\eta_j}.
\]
Tangential operators act on $L$. All polynomial normal jets are
retained. Write $\operatorname{DR}_X(L)$ for the unshifted
algebraic de Rham complex in degrees $0,\ldots,b$.

\begin{proposition}\label{prop:vol2-independent-normal-support}
There is an explicit contraction
\[
 R\operatorname{Hom}_{D_Y}(\mathcal O_Y,i_+L)
       \simeq\operatorname{DR}_X(L)[-c].
\]
For right differential modules, put
$L^r=\omega_X\otimes_{\mathcal O_X}L$ and
$\mathbb J_Y=\omega_Y[b+c]$. Then
\begin{equation}\label{eq:vol2-independent-density-shift}
 R\operatorname{Hom}_{D_Y}
       \bigl(\mathbb J_Y,i_*(L^r[b])\bigr)
       \simeq\operatorname{DR}_X(L)[-2c],
\end{equation}
where $i_*$ denotes right differential-module direct image.
\end{proposition}

\begin{proof}
Apply Hom from the Spencer resolution of $\mathcal O_Y$.
One normal direction contributes the cochain complex
\[
 K_j=[k[\eta_j]\xrightarrow{\eta_j}k[\eta_j]]
\]
in degrees zero and one. Let $p_j$ extract the constant
coefficient in degree one and let $i_j$ insert it. Let $h_j$
vanish on degree zero and send $\eta_j^a$ in degree one to
$\eta_j^{a-1}$ for $a>0$, and $1$ to zero. Then
\[
 dh_j+h_jd=1-i_jp_j,\qquad p_ji_j=1,
 \qquad h_j^2=h_ji_j=p_jh_j=0.
\]
On the tensor product of the normal complexes set
\[
 H=\sum_{j=1}^c(i_1p_1)\otimes\cdots\otimes(i_{j-1}p_{j-1})
                         \otimes h_j\otimes1\otimes\cdots\otimes1,
\]
using the Koszul signs for tensor products of graded maps.
The identity follows by telescoping
$1-q_1\cdots q_c=\sum_jq_1\cdots q_{j-1}(1-q_j)$, with
$q_j=i_jp_j$. The cross terms in $H^2$ contain $h_ji_j$ or
$p_jh_j$ and vanish. Tensoring with the tangential de Rham
complex contracts the whole Spencer complex onto that complex
in normal degree $c$. Every polynomial jet has finite degree;
the argument also applies to an infinite-rank module $L$.

Left-to-right conversion identifies the source and target of
\eqref{eq:vol2-independent-density-shift} with
$\mathcal O_Y[b+c]$ and $i_+L[b]$. Their relative shift is
$-c$. Combining it with the normal Spencer degree gives
$-2c$.
\end{proof}

For an ordered $n$-point affine chart, the first coordinate
and its differences with the other points give one tangential
coordinate and $n-1$ independent normals. Let
$N=\operatorname{End}_{k[y]}(\mathcal V)$, with its commutator
connection, be the polynomial-root endomorphism module of
Theorem~\ref{thm:vol2-polynomial-operators}. Its de Rham
complex is $G(N)$, with cohomology $B_0$ in degree zero and
$W$ in degree one.

Suppose a specified collision comparison identifies the actual
full-diagonal coefficient of an ambient object with $N^r[1]$.
The supported linear mapping complex from
$\mathbb J_{\mathbb A^n}$ is then
\begin{equation}\label{eq:vol2-full-diagonal-groups}
 G(N)[-2(n-1)],\qquad
 H^{2n-2}=B_0,\quad H^{2n-1}=W,
\end{equation}
and all other cohomology vanishes. This follows from
Proposition~\ref{prop:vol2-independent-normal-support} with
$b=1,c=n-1$ and the stated endomorphism contraction. Replacing
the source by $\mathbb J_{\mathbb A^n}[-d]$ changes the complex
to $G(N)[d-2(n-1)]$.

The collision identification is a hypothesis about an actual
coefficient object and its maps. Matching a rank does not supply
it. For $n\geq3$, the complement of the full diagonal still
contains proper collision strata; their simultaneous support
requires their attachment maps. The localization triangle also
controls the map from supported to ambient cohomology, which
need not be injective. Finally, dependent graph equations have
codimension equal to their rank, and this calculation concerns
linear Hom. Coalgebra and substitution compatibility require
their respective operation maps.

\section{Heisenberg compositions on the pair support}
\label{sec:vol2-supported-heisenberg}

We now construct an outer operation on a complete pair-support
complex. Its source includes poles in the remaining separation
and the triple-diagonal term. A correction to the binary path
will make the operation closed on this larger domain.

Fix $\kappa\in k^\times$ and $X=\mathbb A^1_k$. The dg vertex
algebra is $\mathcal Z=V_\kappa\otimes\Lambda[c]$, with
$|J|=0$, $|c|=1$, $dJ=\kappa Tc$ and $dc=0$.
Its state complex is
\begin{equation}\label{eq:vol2-supported-state-complex}
 C=k[u_0,u_1,\ldots]\otimes\Lambda(c_0,c_1,\ldots),
 \quad u_n=T^nJ,\quad c_n=T^nc,\quad
 du_n=\kappa c_{n+1},\quad dc_n=0.
\end{equation}
All states and state tensor products are algebraic. The current
bracket is $[J_\nu J]=\kappa\nu\mathbf1$; brackets with $c$
vanish. The Fourier variable $\nu$ records normal derivatives:
the contact of $v^{-n-1}$ is the $n$th normal derivative divided
by $n!$. On regular state inputs this is the usual bracket
$[a_\nu b]=\operatorname{Res}_v e^{\nu v}Y(a,v)b\,dv$.

\subsection{A common resolution and its state contraction}

Let $E$ be the free quartet vertex algebra on fields of degrees
$(0,0,1,-1)$. Form $\mathcal W=V_\kappa\otimes E^{\otimes3}$,
with quartet generators $(t_s,p_s,\eta_s,\iota_s)$ and
\[
 Qt_s=\eta_s,\qquad Q\iota_s=p_s,
 \qquad [p_s{}_\nu t_l]=[\iota_s{}_\nu\eta_l]=\delta_{sl}.
\]
The other generating brackets vanish, apart from those forced
by graded skew symmetry and translation. In this ambient
algebra take $QJ=0$; its current commutes with the quartets.
Put
\[
 K_3=J+p_1+p_2+p_3,\qquad
 \mathcal B=\operatorname{Com}(K_3,\mathcal W),\qquad
 U_s=J+\sum_{l<s}p_l-\kappa Tt_s.
\]
The commutant means the common kernel of the nonnegative modes
of $K_3$. There are vertex chain maps
\begin{equation}\label{eq:vol2-supported-slots}
 j_s:\mathcal Z\longrightarrow\mathcal B,
 \qquad j_s(J)=U_s,\quad j_s(c)=-\eta_s.
\end{equation}
Distinct images $j_s(\mathcal Z)$ commute.
Indeed $QU_s=-\kappa T\eta_s$ and
$[U_r{}_\nu U_s]=\delta_{rs}\kappa\nu$.
For $r<s$, the mixed contribution
$[-\kappa Tt_r{}_\nu p_r]=-\kappa\nu$ cancels the
current bracket. All current--ghost and ghost--ghost
brackets are zero, and these images commute with $K_3$.
The assignments therefore extend to the stated vertex maps.

Here is a direct description of the state space of $\mathcal B$.
Write $\tau^{(s)}_n=T^nt_s/n!$ for the divided quartet jets,
with $n\geq0$, and use the current modes $x_n=J_{-n}1$.
In these polynomial coordinates, the nonnegative modes of $K_3$
are
\[
 \sum_s\partial_{\tau^{(s)}_0},\qquad
 \kappa n\partial_{x_n}+\sum_s\partial_{\tau^{(s)}_n}\quad(n>0).
\]
Changing coordinates to
$x_n-\kappa n\tau^{(1)}_n$ and
$\tau^{(s)}_n-\tau^{(1)}_n$ for $s=2,3$ leaves these as
derivatives in the remaining $\tau^{(1)}_n$ coordinates.
Their common kernel is therefore freely spanned by the PBW
monomials in the translates of
\[
 U_1,\eta_1;\qquad
 a_2,a_3,e_2,e_3;\qquad
 \iota_1,\iota_2,\iota_3,p_1,p_2,p_3,
 \quad a_s=t_s-t_1,\quad e_s=\eta_s-\eta_1.
\]
Their differentials are
$QU_1=-\kappa T\eta_1$, $Qa_s=e_s$, $Q\iota_s=p_s$,
with the remaining images zero. The fields after the first
semicolon are auxiliary.

Let $r:\mathcal B\to C$ set every auxiliary coordinate to
zero, send $U_1$ to $J$ and $\eta_1$ to $-c$. It is a map of
state complexes. Let $N$ count total auxiliary number and let
$s$ be the odd polynomial derivation with
\[
 s(T^me_j)=T^ma_j\quad(j=2,3),\qquad
 s(T^mp_j)=T^m\iota_j\quad(j=1,2,3),
\]
zero on the other generators. For $q$ tensor factors use the
graded sum $s_q$ and total number $N_q$, and put
$h_q^{\mathrm{aux}}=s_q/N_q$ on positive auxiliary number,
zero on number zero.

\begin{lemma}\label{lem:vol2-supported-state-contraction}
With $i=j_1$, these maps satisfy
\[
 dh_q^{\mathrm{aux}}+h_q^{\mathrm{aux}}d
       =1-i^{\otimes q}r^{\otimes q},\qquad
 r^{\otimes q}i^{\otimes q}=1,
\]
and $(h_q^{\mathrm{aux}})^2=r^{\otimes q}h_q^{\mathrm{aux}}
=h_q^{\mathrm{aux}}i^{\otimes q}=0$.
They commute with translation in each factor and with graded
permutations. Moreover $rj_s=1$ for every slot.
\end{lemma}

\begin{proof}
On generators, $Qs+sQ=N$ and $s^2=0$. The graded Leibniz
rule extends both identities to all PBW polynomials. Their
graded sums give the same identities on tensor powers, since
cross terms cancel. Translation preserves each auxiliary pair
and its number. Division by the positive integer $N_q$ gives
the contraction; on number zero, $i^{\otimes q}r^{\otimes q}$
is the identity. A nonzero value of $h_q^{\mathrm{aux}}$ has
positive auxiliary number, proving the side conditions.
The sums are invariant under graded permutations. Finally all
$U_s$ project to $J$ and all $-\eta_s$ to $c$.
\end{proof}

Write $h=h_1^{\mathrm{aux}}$. The maps $r,h$ are polynomial
state maps, not vertex morphisms. For example, auxiliary
fields killed by $r$ can have a singular product equal to the
vacuum. The collision-coefficient correction will account for
that product.

\subsection{Adjacent paths and all six pair lifts}

Let $A_r\subset\mathcal B$, for $r=1,2$, be generated by
\[
 U_r,\quad D_r=t_{r+1}-t_r,\quad
 p_r,p_{r+1},\eta_r,\eta_{r+1},\iota_r,\iota_{r+1}.
\]
Put $A=\operatorname{Com}(J+p_1+p_2,V_\kappa\otimes E^{\otimes2})$.
Each $A_r$ is a copy of this two-quartet commutant. The first
commutes with $j_3(\mathcal Z)$: it commutes with
$J+p_1+p_2$ and with the third quartet. The second is obtained
by sending $J$ to $J+p_1$ and shifting the two quartet indices;
it commutes with $j_1(\mathcal Z)$. The substitutions preserve
the generating brackets, the differential and the PBW basis.
Thus multiplication gives vertex maps
$A_1\otimes\mathcal Z\to\mathcal B$ and
$\mathcal Z\otimes A_2\to\mathcal B$.

Put polynomial forms $k[s,ds]$ before fields. In $A_r$ define
two paths on the generating states by
\begin{equation}\label{eq:vol2-supported-paths}
 \begin{aligned}
 F_r(J)&=U_r-s\kappa TD_r,&
 F_r(c)&=-\eta_r-s(\eta_{r+1}-\eta_r)-ds\,D_r,\\
 G_r(J)&=U_r-\kappa TD_r+sp_r+ds\,\iota_r,&
 G_r(c)&=-\eta_{r+1}.
 \end{aligned}
\end{equation}
Extend them by the vertex operations. Their images have current
self-bracket $\kappa\nu$, zero current--ghost bracket, and
closed ghost. The tensor differential gives
$dF_r(J)=\kappa TF_r(c)$ and $dG_r(J)=\kappa TG_r(c)$;
the parameter derivatives cancel the differentials of the
$ds$ terms. Hence they are vertex chain maps. They commute
with the fixed spectator throughout, by the commutant
calculation above. Their concatenation runs from $j_r$ to
$j_{r+1}$.

Let $h_r$ be the sum of their integrated $ds$ coefficients.
Each polynomial input has a finite polynomial integrand, so
integration is defined over $k$. Stokes' formula gives
\begin{equation}\label{eq:vol2-supported-path-integral}
 Qh_r+h_rd=j_{r+1}-j_r,
 \qquad h_r(J)=\iota_r,\quad h_r(c)=-D_r,\quad h_r(1)=0.
\end{equation}
The maps commute with translation and satisfy $rh_1=rh_2=0$.
Indeed every form coefficient contains an auxiliary factor
which survives the polynomial state projection until it is
set to zero.

Let $q_I:\mathscr R_I\to X^3$ be the ordered atlas of
Section~\ref{sec:vol2-action-carrier}. Put
$N_C=\omega_X\otimes C$ with
$(fa\,dz)\partial_z=-(f'a+fTa)\,dz$, and let
$P=q_I^!(N_C\boxtimes N_C\boxtimes N_C)$ be the regular
singleton source. For a partition $\pi$, let $\mathcal A_\pi$
be the external product of $\mathcal B^r=\omega_X\otimes\mathcal B$
over its block roots, localized where distinct roots are
separated and pushed onto the corresponding diagonal. This
direct image retains every normal derivative. Order the blocks
by their largest labels and transport their right densities
with that order. The three-position coefficient is
\begin{equation}\label{eq:vol2-supported-three-complex}
 \begin{split}
 \mathcal D_3&=\mathcal A_{1|2|3}
       \oplus\bigoplus_{ab}\mathcal A_{ab|c}[-1]
       \oplus\mathcal A_{123}[-2],\\
 D(x,y,z)&=(dx,b_1x-dy,b_2y+dz).
 \end{split}
\end{equation}
Here $b_1,b_2$ are full vertex contact maps, including
normal transfers and the Koszul signs. On regular states,
$b_1=\varrho_{12}+\varrho_{13}-\varrho_{23}$, where
\[
 \begin{aligned}
 \varrho_{12}(X,Y,W)&=[X_\lambda Y]\boxtimes W,\\
 \varrho_{13}(X,Y,W)&=(-1)^{|X||Y|}Y\boxtimes[X_\lambda W],\\
 \varrho_{23}(X,Y,W)&=X\boxtimes[Y_\mu W].
 \end{aligned}
\]
The second contacts are $[M_{\lambda+\mu}W]$ on $12|3$,
$[Y_\mu M]$ on $13|2$, and $[X_\lambda M]$ on $23|1$.
The vertex Jacobi identity gives
\[
 b_2b_1(X,Y,W)=[[X_\lambda Y]_{\lambda+\mu}W]
 +(-1)^{|X||Y|}[Y_\mu[X_\lambda W]]-[X_\lambda[Y_\mu W]]=0.
\]
These contact formulas extend by the differential-module
relations and retain their Laurent tests.

For $D_{ab}=q_I^{-1}\Delta_{ab}$ and
$D_{123}=q_I^{-1}\Delta_{123}$, including every ordered branch,
the support complexes are
\begin{equation}\label{eq:vol2-supported-pair-complexes}
 M_{ab}=R\Gamma_{D_{ab}}\mathcal D_3
       =\mathcal A_{ab|c}[-1]\oplus\mathcal A_{123}[-2],
 \qquad T_3=R\Gamma_{D_{123}}\mathcal D_3
       =\mathcal A_{123}[-2].
\end{equation}
To verify this description, apply the localization triangle to
the finite filtration by number of blocks. The separated term
is a direct image from the complement of $D_{ab}$ and has zero
extraordinary restriction there. The other pair terms meet
$D_{ab}$ only at the triple diagonal, away from which they are
localized. Only the two displayed terms remain. Their
inclusion is the support counit; localization agrees on chart
overlaps, so the description holds on the atlas.

An order $\sigma$ assigns its labels to slots $1,2,3$ and gives
a singleton insertion $E_\sigma:P\to\mathcal D_3$ through
$j_1,j_2,j_3$. The commuting slot algebras make its pair and
triple components zero. For a binary pair set
\[
 H_r(a,b)=-h_r(a)\boxtimes j_{r+1}(b)
                 +(-1)^{|a|}j_{r+1}(a)\boxtimes h_r(b).
\]
Its Hom differential is
$j_r\boxtimes j_{r+1}-j_{r+1}\boxtimes j_r$.
Use $H_1\boxtimes j_3$ for slots $1,2$ and
$j_1\boxtimes H_2$ for slots $2,3$, with the graded tensor
sign. Relabel source, target, coordinates and densities
together; call the resulting map $H_{\sigma,r}$.
The six edges are
\begin{equation}\label{eq:vol2-supported-hexagon-edges}
 \begin{array}{c|c|c|c}
 \text{edge}&\text{pair}&r&(\text{spectator},\text{slot})\\\hline
 123\to213&12&1&(3,3)\\
 213\to231&13&2&(2,1)\\
 231\to321&23&1&(1,3)\\
 321\to312&12&2&(3,1)\\
 312\to132&13&1&(2,3)\\
 132\to123&23&2&(1,1).
 \end{array}
\end{equation}
Both occurrences of $13$ exchange adjacent slots; neither
moves directly through a fixed middle slot.

\begin{proposition}\label{prop:vol2-supported-pair-lifts}
For an edge $e=(\sigma,r)$ exchanging labels $a,b$, put
\[
 \beta_e=-(b_1H_{\sigma,r})_{ab|c},\qquad
 \widehat\beta_e=(\beta_e,0):P\longrightarrow M_{ab}.
\]
This is a cochain map. Its composite with the support counit
is homotopic to $E_\sigma-E_{\sigma s_r}$ by
$(H_{\sigma,r},0,0)$. It retains the spectator and every normal
derivative. It vanishes on a vacuum in a paired input and
recovers the binary lift on a vacuum spectator.
\end{proposition}

\begin{proof}
The path equation gives the separated component of the
endpoint difference. The image of the paired homotopy lies
in $A_r$, which commutes with the spectator. Thus the two
unwanted first contacts vanish. The remaining contact is
\begin{equation}\label{eq:vol2-supported-binary-lift}
 \beta_r(a,b)=[h_r(a)_\nu j_{r+1}(b)]
       -(-1)^{|a|}[j_{r+1}(a)_\nu h_r(b)],
\end{equation}
tensored with the spectator and the prescribed sign. Its
image lies in $A_r$, so its second contact with that spectator
also vanishes. In the total complex this says
\[
 (E_\sigma-E_{\sigma s_r})
          -\delta(H_{\sigma,r},0,0)=(0,\beta_e,0).
\]
The right side is closed and factors through $M_{ab}$.

All maps commute with translation and use full contacts, so
their differential-module extensions retain the normal
derivatives. Vertex locality bounds the normal order on each
polynomial input. Relabelling transports the normal coordinate
and density once; the same formulas therefore agree on chart
overlaps and separation restrictions. Finally $h_r(1)=0$ and
the regular singular bracket with $1$ is zero, giving the
vacuum assertions.
\end{proof}

\subsection{The two-colour coefficient and its endpoints}

For free dg vertex algebras $D,E$, put
\[
 V(D,E)=j_*j^!(D^r\boxtimes E^r),\qquad
 W(D,E)=\Delta_*(D\otimes E)^r,
\]
where $j$ is the complement of the diagonal in the two-root
space and $D^r=\omega_X\otimes D$. If $v$ is the first root
minus the second, the tensor-product contact is
\begin{equation}\label{eq:vol2-coloured-contact}
 c_{DE}\bigl(f(v)(a\boxtimes b)\bigr)
 =\operatorname{pp}_v\bigl(f(v)(e^{vT_D}a)\otimes b\bigr).
\end{equation}
Only finitely many Taylor terms meet any fixed principal part.
The diagonal term retains the whole tensor $D\otimes E$.
This gives the two-colour complex
\[
 \mathcal Q(D,E)=V(D,E)\oplus W(D,E)[-1],\qquad
 D(x,z)=(dx,c_{DE}x-dz).
\]
Translation makes the contact differential-linear, and the
internal differential commutes with it.

For the resolving algebra itself let
$W(\mathcal B)=\Delta_*\mathcal B^r$ and put
\[
 \mathcal F_{\mathcal B}
   =V(\mathcal B,\mathcal B)\oplus W(\mathcal B)[-1],
 \qquad D(x,z)=(dx,b_{\mathcal B}x-dz),
\]
with its full vertex contact. Denote the positive diagonal
inclusion by $\iota$.

There are two orientations: $R$ for $(A,\mathcal Z)$ and
$L$ for $(\mathcal Z,A)$. Write $\alpha:A\to A_1$ and
$\gamma:A\to A_2$ for the two embeddings. Their placements
in the common resolution are
\begin{equation}\label{eq:vol2-coloured-placements}
 \begin{array}{c|cc}
 &u_+|_V&u_-|_V\\\hline
 R&\alpha\boxtimes j_3&\gamma\boxtimes j_1\\
 L&j_1\boxtimes\gamma&j_3\boxtimes\alpha.
 \end{array}
\end{equation}
On the diagonal use the corresponding normally ordered
products. The images commute, so these are vertex morphisms
and hence coefficient chain maps to $\mathcal F_{\mathcal B}$.

To use them at a pair $ab$, apply right differential-module
direct image along $i_{ab}:X^{\{ab,c\}}\hookrightarrow X^3$.
It includes the entire inner normal transfer and its density.
The subsequent diagonal direct image is $i_{123,+}$.
After the pair-support suspension, the placements have target
$M_{ab}$, and $\iota$ becomes the counit $T_3\to M_{ab}$.

Let $r_A:A\to C$ be the two-quartet polynomial retraction.
Projection in the $A$ factor gives maps $P_R=r_A\otimes1$
and $P_L=1\otimes r_A$, with components $P_V,P_W$.
Since $r_A$ commutes with translation,
\begin{equation}\label{eq:vol2-coloured-projection}
 P_Wc_\eta=c_{\mathcal Z\mathcal Z}P_V,
 \qquad \eta=R,L.
\end{equation}
This uses the tensor-product contact, not multiplicativity
of $r_A$ as a vertex map. The first-slot inclusion and the
binary auxiliary homotopy give a contraction $I,P,H_A$ of
these coloured complexes. The homotopy on the diagonal has
the negative of its state sign; equation
\eqref{eq:vol2-coloured-contact} makes the off-diagonal terms
cancel, so $\delta H_A=1-IP$.

On $V(\mathcal Z,\mathcal Z)$ define
\[
 H_1=-h_1\boxtimes j_2+j_2\boxtimes h_1,
 \qquad K_R=H_1+j_1\boxtimes h_2,
 \qquad K_L=H_1-h_2\boxtimes j_1.
\]
These tensor formulas include the graded signs. Extend them
over every finite meromorphic coefficient and make them zero
on the diagonal source. Their internal differentials are
$j_1\boxtimes j_3-j_2\boxtimes j_1$ and
$j_1\boxtimes j_2-j_3\boxtimes j_1$, respectively.
With $f_{st}(a,b)={:}j_sa\,j_tb{:}$, put
\begin{equation}\label{eq:vol2-coloured-defect}
 F_R=f_{13}-f_{21},\qquad F_L=f_{12}-f_{31},\qquad
 b_\eta=-b_{\mathcal B}K_\eta.
\end{equation}
Here $b_\eta$ has internal degree $-1$, while $F_\eta$ has
degree zero on the diagonal coefficient.

\begin{lemma}\label{lem:vol2-coloured-lift}
The full supported map is
\[
 \mathbb B^0_\eta(x,z)=b_\eta P_Vx+F_\eta P_Wz:
       \mathcal Q_\eta\longrightarrow W(\mathcal B)[-1].
\]
It is closed. If $E_A=u_+-u_-$, then
\begin{equation}\label{eq:vol2-coloured-endpoint-homotopy}
 E_A-\iota\mathbb B^0_\eta=\delta\mathbb K_A,
 \qquad \mathbb K_A=K_\eta P+E_AH_A.
\end{equation}
\end{lemma}

\begin{proof}
The path equation and the chain equations of the placements
give, in unshifted state degrees,
\begin{equation}\label{eq:vol2-coloured-chain-defect}
 \delta b_\eta=-F_\eta c_{\mathcal Z\mathcal Z},
 \qquad \delta F_\eta=0.
\end{equation}
Together with \eqref{eq:vol2-coloured-projection}, these are
the two components of closure into $W(\mathcal B)[-1]$.
On the two-singleton source, subtracting $\delta K_\eta$
from the endpoint difference leaves exactly its supported
components $b_\eta x+F_\eta z$. On arbitrary auxiliary
states, use $E_AI$ for that singleton difference and
$1-IP=\delta H_A$. This gives
\eqref{eq:vol2-coloured-endpoint-homotopy}, including the
term $E_AH_A$.
\end{proof}

The state projection has two useful identities:
\begin{equation}\label{eq:vol2-coloured-endpoint-projection}
 rF_R=rF_L=0,
 \qquad (r\boxtimes r)u_+|_V=(r\boxtimes r)u_-|_V=P_V.
\end{equation}
Indeed each placement is linear on free generators, and the
normally ordered product of its two commuting images is their
concatenated PBW polynomial. Both endpoints have the same
polynomial projection. This does not identify that polynomial
product with the vertex $(-1)$-product in $C$.
The second equality follows from $r\alpha=r\gamma=r_A$
and $rj_s=1$. Translation extends both to the complete
coefficient modules.

Define
\begin{equation}\label{eq:vol2-corrected-punctured-operation}
 s_\eta=hF_\eta,\qquad
 N_\eta=b_\eta+s_\eta c_{\mathcal Z\mathcal Z}.
\end{equation}
The auxiliary contraction and
\eqref{eq:vol2-coloured-endpoint-projection} give
$\delta s_\eta=F_\eta$; thus
\eqref{eq:vol2-coloured-chain-defect} gives $\delta N_\eta=0$.
On regular singleton inputs the tensor-product contact
vanishes. The additional $h_2$ term in $K_\eta$ has zero
regular contact with $j_1$, by the commutant property. Hence
$N_R=N_L=\beta_1$ there. On pole coefficients those terms
and the correction remain part of the operation.

\subsection{The outer map on the actual support}

Use the suspended coloured complex
\begin{equation}\label{eq:vol2-suspended-colours}
 \widetilde{\mathcal Q}_\eta
   =V_\eta[-1]\oplus W_\eta[-2],\qquad
 D(y,z)=(-dy,c_\eta y+dz).
\end{equation}
It is $\mathcal Q_\eta[-1]$ with its diagonal coordinate
multiplied by $-1$. Under the positive support counit, the
coloured lift becomes
\[
 \mathbb B_\eta(y,z)=-b_\eta P_Vy+F_\eta P_Wz.
\]

\begin{theorem}[Outer extension]
\label{thm:vol2-supported-outer-extension}
On the complete pair support define
\begin{equation}\label{eq:vol2-supported-outer-map}
 \mathcal C_\eta(y,z)=-N_\eta(r\boxtimes r)y,
 \qquad L_\eta(y,z)=-s_\eta P_Wz.
\end{equation}
The first is a degree-zero map $M_{ab}\to T_3$ and the second
has total degree $-1$ on the coloured source. Both placements
satisfy
\begin{equation}\label{eq:vol2-supported-placement-equation}
 \delta\mathcal C_\eta=0,
 \qquad \mathcal C_\eta u_\pm-\mathbb B_\eta=\delta L_\eta.
\end{equation}
Use $R$ for $12|3$ and $L$ for $23|1,13|2$, with the stated
root orders and normal direct images. If
$\mathbb K_A^{\mathrm{supp}}$ is the suspended homotopy of
\eqref{eq:vol2-coloured-endpoint-homotopy}, then
\[
 E_A-\iota\mathcal C_\eta u_\pm
 =\delta(\mathbb K_A^{\mathrm{supp}}-\iota L_\eta).
\]
Thus the extension retains the prescribed endpoint products.
\end{theorem}

\begin{proof}
The internal degree of $N_\eta$ is $-1$; the difference
between the pair and triple support shifts makes
$\mathcal C_\eta$ degree zero. It kills the diagonal component
of $M_{ab}$, so its chain equation is precisely
$\delta N_\eta=0$. This proves closure on arbitrary
polynomial auxiliary states and scalar poles.

On either placement, equations
\eqref{eq:vol2-coloured-projection} and
\eqref{eq:vol2-coloured-endpoint-projection} give the difference
\[
 -s_\eta P_Wc_\eta y-F_\eta P_Wz.
\]
The signs $(-,+)$ in
\eqref{eq:vol2-suspended-colours} and the odd degree of
$L_\eta$ give
\[
 \delta L_\eta(y,z)
 =-s_\eta P_Wc_\eta y-(Qs_\eta+s_\eta d)P_Wz,
\]
which is that difference. The final identity follows by
subtracting $\iota\delta L_\eta$ from the suspended endpoint
identity. Every equality is on the complete coloured source,
including its diagonal term.
\end{proof}

Neither placement needs to be inverted. In fact it is not an
equivalence. On diagonal state cohomology its map is
\[
 \Lambda(\alpha,\beta)\longrightarrow\Lambda(\varepsilon),
 \qquad \alpha,\beta\longmapsto\varepsilon.
\]
To compute these groups, contract the pairs
$(u_n,\kappa c_{n+1})$ in \eqref{eq:vol2-supported-state-complex};
only $c_0$ remains. The auxiliary contraction gives the same
cohomology for each resolving algebra. The displayed diagonal
map has kernel containing $\alpha-\beta$ and $\alpha\beta$.
Its separated part is a quasi-isomorphism, so the triangle of
the two-term coefficients identifies its cone with the shifted
cone of this diagonal map. Thus an inverse-placement method
would fail, while \eqref{eq:vol2-supported-outer-map} supplies
the outer extension directly.

The pole correction has a small explicit diagnostic. In the
$R$ orientation put
\[
 a=\iota_1+\iota_2,\quad p=p_1+p_2,\quad
 t=T(t_3-t_1),\quad e=T(\eta_3-\eta_1).
\]
For $x_J=v^{-1}(1\boxtimes J)$ and
$x_{Tc}=v^{-1}(1\boxtimes Tc)$, let $z_J,z_{Tc}$ denote the
corresponding diagonal transfers. The defining formulas give
\[
 \begin{gathered}
 Dx_J=\kappa x_{Tc}+z_J,\quad Dx_{Tc}=z_{Tc},\quad
 Dz_J=-\kappa z_{Tc},\\
 b_Rx_J=-a,\quad b_Rx_{Tc}=t,\quad
 F_Rz_J=p-\kappa t,\quad F_Rz_{Tc}=-e.
 \end{gathered}
\]
Since $h(p-\kappa t)=a$ and $h(-e)=-t$, the corrected
values are $N_Rx_J=-a+a=0$ and $N_Rx_{Tc}=t-t=0$.
Omitting the correction leaves chain defect $-(p-\kappa t)$
on $x_J$; reversing its sign leaves $-2(p-\kappa t)$.
These are nonzero states. They explain the correction on
poles without asserting an obstruction to the corrected map.

\section{The triple homotopy and its coefficient transfer}
\label{sec:vol2-supported-triple-homotopy}

The six pair lifts have two occurrences of each labelled pair.
A cube filling compares these occurrences. Its triple contact
currents sum to zero, whereas the actual iterated outer maps
need a further comparison. We construct both maps before
identifying their difference as a boundary.

\subsection{The nine-square filling}

Let $\mathsf I$ be the cellular chain complex of the interval
$[1,3]$, with vertices $v_1,v_2,v_3$ in degree zero, edges
$e_1,e_2$ in degree one and
$\partial e_r=v_{r+1}-v_r$. A cell of $\mathsf I^{\otimes3}$
assigns a height or a moving edge to each label, in label
order. Define
\[
 \Gamma(c_1\otimes c_2\otimes c_3)
 =\phi(c_1)\boxtimes\phi(c_2)\boxtimes\phi(c_3),
 \qquad \phi(v_s)=j_s,\quad\phi(e_r)=h_r.
\]
Its target is the separated term of
\eqref{eq:vol2-supported-three-complex}; a cell of dimension
$d$ gives an operation of degree $-d$. The graded tensor
formula and \eqref{eq:vol2-supported-path-integral} give
$\delta\Gamma(c)=\Gamma(\partial c)$ for the internal
differential.

For an edge $\sigma\to\sigma s_r$, let $x,y$ be its two
exchanged labels and $z$ the spectator at height $q$.
The corresponding chain is
\[
 c_{\sigma,r}=-(x:e_r,y:v_{r+1},z:v_q)
                  +(x:v_{r+1},y:e_r,z:v_q),
\]
with factors put in label order. Its boundary is the first
permutation vertex minus the second, and
$\Gamma(c_{\sigma,r})=H_{\sigma,r}$.
The sum around the hexagon is the boundary of
\begin{equation}\label{eq:vol2-nine-square-filling}
 \Sigma=\sum_{i=1}^3(-1)^i
       \bigl((e_1,e_2)_i+(e_2,e_1)_i+(e_2,e_2)_i\bigr),
\end{equation}
where $(c,c')_i$ fixes label $i$ at $v_3$ and assigns $c,c'$
to the other two labels in increasing order. To verify this,
the hexagon chain has twelve edges: it passes through the six
permutation vertices and through
\[
 (2,2,3),(3,1,3),(3,2,2),(3,3,1),(2,3,2),(1,3,3).
\]
On each of the three faces through $(3,3,3)$, it removes the
square at the opposite $(1,1)$ corner. The remaining three
squares are those in \eqref{eq:vol2-nine-square-filling}.
Their twelve interior edges cancel: two within each face and
two along each of the three shared edges. The exterior twelve
edges have precisely the hexagon orientation.

Put $\Phi=\Gamma(\Sigma)$ and define its contact maps
\begin{equation}\label{eq:vol2-cube-contacts}
 K_{ab}=(b_1\Phi)_{ab|c},\qquad \gamma_{ab}=-b_2K_{ab}.
\end{equation}
Their internal state degree is $-2$. Thus $(0,K_{ab},0)$
has total degree $-1$ and $(0,0,\gamma_{ab})$ has total
degree zero in \eqref{eq:vol2-supported-three-complex}.

\begin{proposition}\label{prop:vol2-cube-contacts}
If $e_+,e_-$ are the two hexagon occurrences of the pair $ab$,
then
\begin{equation}\label{eq:vol2-cube-pair-identity}
 \widehat\beta_{e_+}+\widehat\beta_{e_-}
       =\delta(0,K_{ab},0)+(0,0,\gamma_{ab}),
 \qquad \delta\gamma_{ab}=0,
\end{equation}
and $\gamma_{12}+\gamma_{23}+\gamma_{13}=0$.
The whole supported hexagon bounds:
\[
 \sum_e(H_e,0,0)+\sum_{ab}(0,K_{ab},0)=\delta(\Phi,0,0).
\]
\end{proposition}

\begin{proof}
The cube boundary gives
$\delta(\Phi,0,0)=(\sum_eH_e,b_1\Phi,0)$.
Differentiate it. The endpoint differences telescope, and
Proposition~\ref{prop:vol2-supported-pair-lifts} identifies
the remaining terms as the six supported lifts. Comparing
each pair component gives
$\beta_{e_+}+\beta_{e_-}=-\delta K_{ab}$ internally.
Its triple component gives the first identity with
$\gamma_{ab}=-b_2K_{ab}$. Differentiating that identity shows
$\gamma_{ab}$ is closed. Finally
$\sum\gamma_{ab}=-b_2b_1\Phi=0$ by vertex Jacobi.
\end{proof}

For the comparison with iterated brackets, set $g=h_1+h_2$
and define the merged-pair transport
\[
 k(a,b)=(-1)^{|a|}
   \bigl([h_1a_\nu h_2b]+[h_2a_\nu h_1b]+[h_2a_\nu h_2b]\bigr).
\]
Its internal differential is $\delta k=\beta_1-\beta_2$.
This follows by differentiating the two homotopies; the
contacts of $h_1$ with $j_3$ and of $h_2$ with $j_1$
vanish by the commutant property. Taking the contacts of the
nine squares gives
\begin{align}\label{eq:vol2-cube-pair-formulas}
 K_{12}(a,b,c)&=-k(a,b)\boxtimes j_3c
          +(-1)^{|a|+|b|}\beta_2(a,b)\boxtimes gc,\notag\\
 K_{23}(a,b,c)&=(-1)^{|a|}ga\boxtimes\beta_2(b,c)
                                      +j_3a\boxtimes k(b,c),\\
 K_{13}(a,b,c)&=(-1)^{|b|(|a|+1)}gb\boxtimes\beta_2(a,c)
                         +(-1)^{|a||b|}j_3b\boxtimes k(a,c).
                         \notag
\end{align}
For the first formula, the face fixing the spectator at
$v_3$ contributes the three terms of $-k$. On the other two
faces the spectator moves along $g$; the $h_1$ contacts with
$j_3$ vanish and the rest give $\beta_2$. The same calculation
gives the second formula. Moving label $2$ in front of
label $1$ gives the third formula and its Koszul signs.

\subsection{The three actual compositions}

Define the bracket on states by
\[
 B_\nu(a,b)=(-1)^{|a|}r\beta_1(a,b),
 \qquad x(a,b)=r\beta_1(a,b).
\]
It has degree $-1$ and is the operation fixed by the two
polynomial paths and the PBW retraction. Its three channels
are
\begin{equation}\label{eq:vol2-supported-three-channels}
 \begin{split}
 L(a,b,c)&=B_{\lambda+\mu}(B_\lambda(a,b),c),\\
 R_B(a,b,c)&=B_\lambda(a,B_\mu(b,c)),\\
 N_B(a,b,c)&=B_\mu(b,B_\lambda(a,c)),
 \qquad \epsilon=(-1)^{(|a|+1)(|b|+1)}.
 \end{split}
\end{equation}
The Fourier variables belong to positions $1,2$ relative to
position $3$. Sesquilinearity sends their translations to
$T_1=-\lambda$, $T_2=-\mu$, $T_3=T+\lambda+\mu$.

Choose the edges
$e_{12}:123\to213$, $e_{23}:231\to321$ and
$e_{13}:213\to231$. Their placements are respectively
$u_{R,+}$, $u_{L,-}$ and $u_{L,+}$. Define literal compositions
through the pair supports by
\begin{equation}\label{eq:vol2-supported-iterated-maps}
 G_{12}=\mathcal C_R\widehat\beta_{e_{12}},\qquad
 G_{23}=-\mathcal C_L\widehat\beta_{e_{23}},\qquad
 G_{13}=\mathcal C_L\widehat\beta_{e_{13}}.
\end{equation}
Their source remains $P$. To compare with the coloured map,
let $a_e:P\to\widetilde{\mathcal Q}_\eta$ have the indicated
binary lift in $A$ tensored with the spectator, and zero
diagonal component. Its outer scalar coefficient is regular,
so $c_\eta a_e=0$, including after any inner normal transfer.
It is therefore closed, $u_ea_e=\widehat\beta_e$ and
$L_\eta a_e=0$. Equation
\eqref{eq:vol2-supported-placement-equation} becomes
\[
 \mathcal C_\eta\widehat\beta_e=\mathbb B_\eta a_e.
\]
This is equality of operations on the whole regular singleton
source.

Projection of the preferred pair components gives
\[
 \begin{array}{c|c|c}
 \text{pair}&(r\boxtimes r)\beta_e(a,b,c)&\text{outer variable}\\\hline
 12&x(a,b)\boxtimes c&\lambda+\mu\\
 23&-(-1)^{|a|}a\boxtimes x(b,c)&\lambda\\
 13&(-1)^{|a||b|+|b|}b\boxtimes x(a,c)&\mu.
 \end{array}
\]
The minus sign in the second row is its first-contact sign;
the third row includes moving the spectator to the first root.
Since $N_R=N_L=\beta_1$ on a regular outer coefficient,
\begin{align}\label{eq:vol2-supported-tree-formulas}
 G_{12}(a,b,c)&=-\beta_1(x(a,b),c),\notag\\
 G_{23}(a,b,c)&=-(-1)^{|a|}\beta_1(a,x(b,c)),\\
 G_{13}(a,b,c)&=-(-1)^{|a||b|+|b|}\beta_1(b,x(a,c)).\notag
\end{align}
The degree of $x(a,b)$ is $|a|+|b|-1$, so
\begin{equation}\label{eq:vol2-supported-channel-projections}
 rG_{12}=(-1)^{|b|}L,
 \qquad rG_{23}=-(-1)^{|b|}R_B,
 \qquad rG_{13}=(-1)^{|b|}\epsilon N_B.
\end{equation}
For example, the first sign is
$-(-1)^{|a|+|b|-1}(-1)^{|a|}=(-1)^{|b|}$;
the third uses
$-(-1)^{|a||b|+|a|}=(-1)^{|b|}\epsilon$.

\subsection{Comparing cube contacts with the bracket}

The comparison uses a polynomial projection $\pi$ retaining
$U_1,\eta_1,p_1,\iota_1$ and killing every other auxiliary
coordinate and its translates.

\begin{lemma}\label{lem:vol2-supported-pbw-reduction}
For $Z\in\mathcal B$ and $a,b,y\in C$,
\begin{gather*}
 r[Z_\nu j_1y]=r[\pi Z_\nu j_1y],\qquad
 r[j_1y_\nu Z]=r[j_1y_\nu\pi Z],\\
 \pi h=h\pi,\qquad \pi Y=\pi j_2(rY)\quad(Y\in A_2),
 \qquad \pi hj_2=\pi h_1,\\
 r[h_1a_\nu h_1b]=0,\\
 r\beta_1(a,b)=r[h_1a_\nu j_1b]
                   -(-1)^{|a|}r[j_1a_\nu h_1b].
\end{gather*}
\end{lemma}

\begin{proof}
Only $U_1,p_1$ contract with $U_1$, and only $\iota_1$
contracts with $\eta_1$. A coordinate killed by $\pi$ thus
survives every contraction with $j_1y$ and is killed by $r$.
This proves the first two equalities. The auxiliary reverse
derivation carries the killed coordinates among themselves,
and their absence leaves auxiliary number unchanged, proving
$\pi h=h\pi$. On the generators of $A_2$, both $\pi$ and
$\pi j_2r$ send $U_2$ to $U_1+p_1$, send the ghosts to
$\eta_1$, and kill the remaining auxiliary fields. Hence
they agree on all PBW polynomials.

For $\pi hj_2=\pi h_1$, the first path contributes a factor
of $a_2$, which $\pi$ kills. In the second path, a monomial
with $m$ current factors choosing $p_1$ has form coefficient
obtained by replacing one of them by $\iota_1$ and weighting
the other $m-1$ by $s^{m-1}$. Integration gives $1/m$,
exactly the auxiliary contraction of the same monomial in
$j_2$. The graded sign of the odd replacement is the same.

To prove $r[h_1a_\nu h_1b]=0$, split $h_1=h^F+h^G$ by
the two paths. Every monomial of $h^F$ contains one translate
of $D=t_2-t_1$; every monomial of $h^G$ contains one translate
of $\iota_1$. Two $F$ terms leave both $D$ factors uncontracted,
and two $G$ terms leave both $\iota_1$ factors. Their
projections vanish. For a mixed term, use the generating
polynomials represented coefficientwise by
\[
 a=\exp\Bigl(\sum_n\alpha_nu_n+\sum_n\theta_nc_n\Bigr),
 \quad
 b=\exp\Bigl(\sum_n\alpha'_nu_n+\sum_n\theta'_nc_n\Bigr),
\]
where $\alpha,\alpha'$ are even and $\theta,\theta'$ odd.
Only finitely many parameters are needed for any chosen
coefficient. If $A_s,B_t$ are the zero-form exponents of the
two paths, their form coefficients are
\[
 h^Fa=\int_0^1\Bigl(\sum_n\theta_nT^nD\Bigr)e^{A_s}ds,
 \quad
 h^Gb=\int_0^1\Bigl(\sum_m\alpha'_mT^m\iota_1\Bigr)e^{B_t}dt.
\]
For the projection to survive, $D$ must contract with a
$tp_1$ factor of $B_t$ and $\iota_1$ with an
$-(1-s)\eta_1$ factor of $A_s$. Since
$[D_\nu p_1]=[\eta_1{}_\nu\iota_1]=1$, both contractions
have jet kernel
\[
 P_{nm}(z-w)=(-1)^n(n+m)!(z-w)^{-n-m-1}.
\]
Their product is $-t(1-s)$ times the square of the odd
element $\sum_{n,m}\theta_n\alpha'_mP_{nm}$, hence zero.
All other contractions multiply this zero. Reversing the
two paths gives the same argument with $\theta'$.
Taking polynomial coefficients proves the identity on all
states.

Finally, the contact of $h_1\otimes h_1$ followed by $r$
is zero. Its differential, using
$\delta h_1=j_2-j_1$, is the difference between
$r\beta_1$ and the last displayed bracket expression.
That difference is therefore zero, proving the final formula.
\end{proof}

Put
\begin{equation}\label{eq:vol2-cube-transfer-homotopy}
 T_{ab}=-r b_2(h\text{ on the merged factor})K_{ab},
 \qquad \widehat H=-(T_{12}+T_{23}+T_{13}).
\end{equation}
These maps have internal degree $-3$. The operation on the
merged factor uses its graded tensor sign.

\begin{proposition}\label{prop:vol2-cube-state-comparison}
On all polynomial states,
\[
 r\gamma_{12}=(-1)^{|b|}L+\delta T_{12},\qquad
 r\gamma_{23}=-(-1)^{|b|}R_B+\delta T_{23},\qquad
 r\gamma_{13}=(-1)^{|b|}\epsilon N_B+\delta T_{13}.
\]
Thus $\delta\widehat H=(-1)^{|b|}(L-R_B+\epsilon N_B)$.
\end{proposition}

\begin{proof}
For $12$, use $\delta(h\otimes1)=(1-j_1r)\otimes1$
and the fact that $rb_2$ commutes with the internal
differential. Equation \eqref{eq:vol2-cube-pair-identity}
gives
\[
 r\gamma_{12}-\delta T_{12}
 =-rb_2(j_1r\otimes1)K_{12}
   +rb_2(h\otimes1)(\beta_{e_+}+\beta_{e_-}).
\]
Here $\beta_{e_+}+\beta_{e_-}=
\beta_1\otimes j_3-\beta_2\otimes j_1$, and
$r\beta_1=r\beta_2=x$, since the slot shift has the same
polynomial projection in both binary copies. In the first
term, the part with $j_3$ has zero contact with $j_1$;
in the part with $g$, only $h_1$ remains. In the second term,
$h\beta_1\in A_1$ commutes with $j_3$.
Lemma~\ref{lem:vol2-supported-pbw-reduction} replaces the
remaining contact of $h\beta_2$ with $j_1$ by that of
$h_1x$. The result is
\[
 -rb_2(j_1x\otimes h_1+h_1x\otimes j_1)
       =-r\beta_1\circ(x\otimes1)=(-1)^{|b|}L.
\]

For $23$, the tensor formulas are
$K_{23}=g\otimes\beta_2+j_3\otimes k$ and
$\beta_{e_+}+\beta_{e_-}=j_1\otimes\beta_2-j_3\otimes\beta_1$.
Apply $1\otimes h$ to the merged factor. The same commutant
vanishings and reduction give
\[
 r\gamma_{23}-\delta T_{23}
 =rb_2(j_1\otimes h_1x-h_1\otimes j_1x)
 =-r\beta_1\circ(1\otimes x)=-(-1)^{|b|}R_B.
\]
The tensor $1\otimes x$ contributes $(-1)^{|a|}$.
For $13$, transpose $a,b$ and $\lambda,\mu$ with input sign
$(-1)^{|a||b|}$. This is a chain map. The result is
$-(-1)^{|a||b|+|a|}N_B=(-1)^{|b|}\epsilon N_B$.
Summing the three identities and using $\sum\gamma_{ab}=0$
proves the assertion about $\widehat H$.
\end{proof}

\begin{theorem}[Supported triple homotopy]
\label{thm:vol2-supported-triple-homotopy}
The maps
\begin{equation}\label{eq:vol2-supported-comparison-primitive}
 U_{ab}=j_1T_{ab}+h(\gamma_{ab}-G_{ab})
\end{equation}
satisfy $\delta U_{ab}=\gamma_{ab}-G_{ab}$. Consequently
\begin{equation}\label{eq:vol2-supported-triple-primitive}
 \mathcal H=-\sum_{ab}U_{ab}
       =j_1\widehat H+h(G_{12}+G_{23}+G_{13})
\end{equation}
has total degree $-1$ and
\[
 \delta\mathcal H=G_{12}+G_{23}+G_{13},\qquad
 r\mathcal H=\widehat H.
\]
These are maps into the full triple-supported coefficient
on every ordered branch, with all normal derivatives retained.
\end{theorem}

\begin{proof}
Both $\gamma_{ab}$ and $G_{ab}$ are closed. Their projection
difference is $\delta T_{ab}$ by
\eqref{eq:vol2-supported-channel-projections} and
Proposition~\ref{prop:vol2-cube-state-comparison}. The
auxiliary contraction gives
\[
 \delta\bigl(h(\gamma_{ab}-G_{ab})\bigr)
 =\gamma_{ab}-G_{ab}-j_1\delta T_{ab}.
\]
Adding $\delta(j_1T_{ab})$ proves the first equation.
The sum of the cube contacts is zero, so summation gives
\eqref{eq:vol2-supported-triple-primitive} and its differential.
The side condition $rh=0$ gives its projection.
Every constituent commutes with translation and extends over
the full normal transfer; all are polynomial on each input.
Thus these are identities of supported differential-module
operations, rather than only identities of state values.
\end{proof}

Changing a preferred occurrence of a pair changes its channel
by a specified boundary. Applying the closed outer map to
\eqref{eq:vol2-cube-pair-identity}, whose triple component it
kills, gives
\[
 \mathcal C_{ab}\widehat\beta_{e_+}
 +\mathcal C_{ab}\widehat\beta_{e_-}
 =\delta\bigl(\mathcal C_{ab}(0,K_{ab},0)\bigr).
\]
Change $\mathcal H$ by the same signed map. Changes for
different pairs add and therefore commute. This retains all
eight choices of one occurrence per pair, including both
$13|2$ edges.

\subsection{Normal transfers, substitutions and vacua}

The diagram just constructed has objects $P$, the coloured
complexes $\widetilde{\mathcal Q}_\eta$, the three pair
supports $M_{ab}$, the triple support $T_3$ and $\mathcal D_3$.
Its arrows include the first lifts $a_e$, both placements,
the outer maps and the support counits. Equation
\eqref{eq:vol2-supported-placement-equation} supplies the
placement homotopy, and
Theorem~\ref{thm:vol2-supported-triple-homotopy} supplies its
three-channel comparison.

If a specified collision substitution gives a chain map
$v:D'\to\widetilde{\mathcal Q}_\eta$, precomposition gives
\[
 \mathcal C_\eta u_\pm v-\mathbb B_\eta v
          =\delta(L_\eta v).
\]
For a further substitution $w$, the homotopies agree
literally: $(L_\eta v)w=L_\eta(vw)$. At three positions the
maps in this diagram are the vertex placements and the first
lifts already constructed. This statement does not assume
an unspecified system of higher substitution maps.

Normal differentiation acts on coefficients as well as
states. With $\nu$ the first-root Fourier variable and $T$
output translation, the corrected operation satisfies, for
every integer pole order $m$,
\begin{align}\label{eq:vol2-outer-normal-relations}
 N_\eta(Ta,b;v^{-m})
   &=-\nu N_\eta(a,b;v^{-m})
                       +mN_\eta(a,b;v^{-m-1}),\notag\\
 N_\eta(a,Tb;v^{-m})
   &=(T+\nu)N_\eta(a,b;v^{-m})
                       -mN_\eta(a,b;v^{-m-1}).
\end{align}
Differentiate $v^{-m}Y(a,v)b$ and then take principal parts
to obtain these identities for the contact part. The same
calculation with \eqref{eq:vol2-coloured-contact} gives them
for the correction. The initial closed embedding retains
the inner collision, the two-term coloured differential
retains the outer one, and the triple comparison retains their
intersection. Open restrictions localize these identities.
Relabelling transports states, normals and densities together.
Different ordered branches of the triple diagonal have no
overlap, so the formulas glue branchwise.

On a pole coefficient with either colour equal to the vacuum,
$rb_\eta=0$: the contact is a principal part of a translate of
$h_s(a)$, and $rh_s=0$. Also $rs_\eta=0$. Thus $N_\eta$
is closed with zero projection on that face, and $hN_\eta$
is its explicit null-homotopy. The canonical outer map has
the negative of this homotopy. On the common all-vacuum face
both are zero. The endpoint-relative homotopy remains the
one of Theorem~\ref{thm:vol2-supported-outer-extension}.

For the ternary maps, a vacuum in the paired inputs makes
the first lift zero. A vacuum spectator makes the regular
outer bracket zero. Hence every $G_{ab}$ vanishes on a
vacuum input. The formulas for $K_{ab}$ show that each
$\gamma_{ab}$ and $T_{ab}$ also vanishes there, so $U_{ab}$
and $\mathcal H$ do as well. A vacuum spectator still leaves
$K_{12}(a,b,1)=-k(a,b)\boxtimes1$; its outer contact is zero,
while its lower binary transport is retained.

\subsection{Removing the auxiliary fields on coefficients}

The state retraction must be corrected before it acts on
collision coefficients. For a finite set $I$ of $n$ positions,
write $\mathcal D_{A,I}=C_{A,I}[-n]$, for $A=\mathcal B$
or $\mathcal Z$. A partition with $q$ blocks has term
$\mathcal A_\pi[q-n]$. Split its differential as
\[
 D_A=D_{0,A}+b_A,\qquad
 D_{0,A}|_\pi=(-1)^{n-q}d_q,
\]
where $b_A$ merges blocks by full contacts.
Extend $i^{\otimes q},r^{\otimes q},h_q^{\mathrm{aux}}$
over the scalar coefficients and normal direct images.
On that partition put
\[
 I_\pi=i^{\otimes q},\qquad P_\pi=r^{\otimes q},\qquad
 H_\pi=(-1)^{n-q}h_q^{\mathrm{aux}}.
\]
Then $D_{0,\mathcal B}H+HD_{0,\mathcal B}=1-IP$.
The inclusion is a vertex map, so $b_{\mathcal B}I=Ib_{\mathcal Z}$.
Define finite operators
\begin{equation}\label{eq:vol2-supported-coefficient-transfer}
 R_I=\sum_{j=0}^{n-1}(-b_{\mathcal B}H)^j,
 \qquad \Pi_I=PR_I,\qquad \mathcal H_I=HR_I.
\end{equation}
Every contact lowers the block number. Thus
$(b_{\mathcal B}H)^n=0$ and $R_I=(1+b_{\mathcal B}H)^{-1}$.

\begin{proposition}\label{prop:vol2-supported-coefficient-transfer}
The maps in \eqref{eq:vol2-supported-coefficient-transfer}
give a contraction of the actual coefficient complexes:
\[
 \Pi_ID_{\mathcal B}=D_{\mathcal Z}\Pi_I,\qquad
 \Pi_II=1,\qquad
 D_{\mathcal B}\mathcal H_I+\mathcal H_ID_{\mathcal B}=1-I\Pi_I,
\]
with $\mathcal H_I^2=\Pi_I\mathcal H_I=\mathcal H_II=0$.
They preserve every fixed collision support and its counit,
commute with relabelling and open restriction, and preserve
the scalar vacuum maps. Their values have finite normal order
on each polynomial input.
\end{proposition}

\begin{proof}
Split the internal complex as $I\mathcal D_{\mathcal Z,I}
\oplus K$, where $K=\ker P$. On $K$ let the internal
differential be $d_K$ and the homotopy be $h_K$. Then
$d_Kh_K+h_Kd_K=1$, $h_K^2=0$. Since $I$ preserves contacts,
the full differential has the block form
\[
 D_{\mathcal B}=\begin{pmatrix}D_{\mathcal Z}&a\\0&D_K\end{pmatrix},
 \qquad D_K=d_K+b_K,\qquad D_{\mathcal Z}a+aD_K=0.
\]
The contact identities give $b_K^2=0$ and
$d_Kb_K+b_Kd_K=0$.
Put $A_K=1+b_Kh_K$. Direct multiplication gives
$D_KA_K=A_Kd_K$, since their difference is
$b_K(1-d_Kh_K-h_Kd_K)=0$.
Therefore $h'_K=h_KA_K^{-1}$ satisfies
$D_Kh'_K+h'_KD_K=1$.
Also $A_K^{-1}h_K=h_K$, so $(h'_K)^2=0$.

The displayed finite formulas are exactly
\[
 \Pi_I(z,k)=z-ah'_Kk,\qquad
 \mathcal H_I(z,k)=(0,h'_Kk).
\]
The asserted equations follow by matrix multiplication,
using $D_{\mathcal Z}a+aD_K=0$. In particular the target
differential is the original one with its vertex contacts.

The state maps commute with translation in each block and
are symmetric under graded permutations. The contacts have
the same equivariance with their Chevalley signs. Thus every
term agrees on ordered chart overlaps and open restrictions.
A partition lies in a fixed collision support when it is
coarser than its defining partition. State maps preserve
partitions and contacts only coarsen them, proving support
compatibility. The localization description of
\eqref{eq:vol2-supported-pair-complexes} identifies its
inclusion with the counit. The vacuum map factors through
$I$, on which $H$ vanishes. Finally there are at most $n-1$
contacts, each of finite normal order on a polynomial input.
\end{proof}

At three positions, $\Pi_3=P-PbH+PbHbH$.
For $w=\iota_1\boxtimes\eta_1\boxtimes1$ and
$x=p_1\boxtimes\eta_1\boxtimes1$, the differential is
$D_{\mathcal B}w=x+v_{\mathcal B}$, where
$v_{\mathcal B}$ is the simple pair-contact vacuum.
One has $Hx=w$, $Hw=Hv_{\mathcal B}=0$, and
$Pw=Px=0$, $Pv_{\mathcal B}=v_{\mathcal Z}$.
Consequently $\Pi_3(x)=-v_{\mathcal Z}$ and
$\Pi_3(v_{\mathcal B})=v_{\mathcal Z}$.
The uncorrected projection has defect $v_{\mathcal Z}\ne0$
on $w$.

On a pair support the formula becomes
\begin{equation}\label{eq:vol2-supported-pair-transfer}
 \Pi_M(y,z)=\bigl((r\boxtimes r)y,
                       rz+rb_2h_2^{\mathrm{aux}}y\bigr).
\end{equation}
There is at most one contact, and the pair homotopy has sign
$-h_2^{\mathrm{aux}}$, giving the plus sign. On the triple
support, $\Pi_T=r$ and $I_T=j_1$.
Let $I_M$ be its pair inclusion and put
\[
 \beta_e^C=\Pi_M\widehat\beta_e,\qquad
 \mathcal C_{ab}^C=\Pi_T\mathcal C_{ab}I_M.
\]
Then
\begin{equation}\label{eq:vol2-supported-composition-transfer}
 \Pi_T\mathcal C_{ab}\widehat\beta_e
                  =\mathcal C_{ab}^C\beta_e^C,\qquad
 \Pi_T\mathcal H=\widehat H.
\end{equation}
Indeed $\mathcal C_{ab}$ kills the triple component of the
pair homotopy, and $(r\boxtimes r)h_2^{\mathrm{aux}}=0$.
Hence $\mathcal C_{ab}\mathcal H_M=0$. Apply the closed
$\mathcal C_{ab}$ to $1-I_M\Pi_M=\delta\mathcal H_M$ to
obtain $\mathcal C_{ab}=\mathcal C_{ab}I_M\Pi_M$.
This proves the strict composition identity. The pair-to-triple
term of \eqref{eq:vol2-supported-pair-transfer} still occurs in
each transferred first lift and its endpoint homotopy.

\subsection{A nonzero state Jacobiator}

The homotopy is necessary even though the cube-contact sum
vanishes. Write $u=u_0$, $c=c_0$. The bracket fixed above
satisfies
\begin{align}\label{eq:vol2-composite-bracket-values}
 B_\nu(u^2,cTc)&=-2uTc+2\nu uc+2(Tu)c,\notag\\
 B_\nu(u^2,uc)&=-2u^2-\tfrac\kappa2\nu^2,\notag\\
 B_\nu(u^2,uTc)&=-2\nu u^2-2uTu-\tfrac\kappa6\nu^3,\\
 B_\nu(u^2,(Tu)c)&=-2uTu-\tfrac\kappa3\nu^3,\qquad
 B_\nu(u^2,u^2)=0.\notag
\end{align}
For the scalar terms, the parts of $h_1(uc)$, $h_1(uTc)$
and $h_1((Tu)c)$ without $\iota_1$ are respectively
\[
 -{:}U_1D_1{:}+\tfrac\kappa2{:}(TD_1)D_1{:},\quad
 -{:}U_1TD_1{:}+\tfrac\kappa2{:}(TD_1)^2{:},\quad
 -{:}(TU_1)D_1{:}+\tfrac\kappa2{:}(T^2D_1)D_1{:}.
\]
The factor $1/2$ is the integral of the first path parameter.
Every omitted term contains $\iota_1$, which cannot contract
with the second embedding and is killed by $r$.
The contractions of $U_2$ with $D_1,TD_1,T^2D_1$ are
$-v^{-1},-v^{-2},-2v^{-3}$. Their double contractions with
$U_2^2$, before the bracket's minus sign, are
$\kappa v^{-3},\kappa v^{-4},2\kappa v^{-4}$.
Converting $v^{-m-1}$ to $\nu^m/m!$ gives the three scalar
coefficients. Single contractions and the Taylor terms give
the displayed state terms. The two ordered ghost contractions
give the first line, and every term in the last bracket
retains an auxiliary field.

Substitution in the three channels gives
\[
 L(u^2,u^2,cTc)-R_B(u^2,u^2,cTc)
                   +\epsilon N_B(u^2,u^2,cTc)
                  =\frac\kappa3(\lambda+\mu)^3\mathbf1.
\]
The state terms cancel and $\epsilon=-1$.
Set $H(a,b,c)=(-1)^{|b|}\widehat H(a,b,c)$.
The ordinary Hom equation becomes
\begin{align}\label{eq:vol2-state-jacobi-homotopy}
 L-R_B+\epsilon N_B={}&dH(a,b,c)+H(da,b,c)\notag\\
 &-(-1)^{|a|}H(a,db,c)
             +(-1)^{|a|+|b|}H(a,b,dc).
\end{align}
Since $d(u^2)=2\kappa uTc$ and $d(cTc)=0$, degree and
the displayed value force
\[
 H(uTc,u^2,cTc)-H(u^2,uTc,cTc)=\frac{(\lambda+\mu)^3}{6}.
\]
The scalar vacuum need not be exact in $C$: the input
differential contributes the boundary in this operation
complex. The noncocycle example therefore does not define
an obstruction class from three cohomology classes.

The construction supplies outer operations on the complete
pair-support complexes and their triple homotopy for the
regular singleton source. Extending these fixed maps to
higher merged collision diagrams remains a further equation.
It must preserve the endpoint homotopies, proper collision
strata and vacuum faces already specified here.

\section{Directed composition beyond the associator}
\label{sec:vol2-directed-higher-composition}

A finite directed category isolates the next coherence
equation. Four arrows have five parenthesizations, so their
associativity homotopies must satisfy a pentagon relation.
This calculation fixes a carrier of rank one on every forward
arrow. Its solution will not identify that carrier with a
supported field-operation complex.

Fix objects $0,\ldots,4$, degree-zero arrows $a_{ij}$ for
$i<j$, scalar identities and no backward arrows. Products
are written chronologically. Introduce ten degree-zero
coefficients and their associativity equations:
\[
 a_{ij}a_{jk}=m_{ijk}a_{ik},\qquad
 F_{ijkl}=m_{ijk}m_{ikl}-m_{jkl}m_{ijl}.
\]
The dg coefficient algebra
\[
 K=k[m_{ijk}]\otimes\Lambda(\chi_{ijkl}),\qquad
 |\chi_{ijkl}|=-1,\quad d\chi_{ijkl}=F_{ijkl},\quad dm_{ijk}=0
\]
retains a chosen homotopy for each binary associativity
equation. All index sequences are increasing.

\begin{proposition}[The relative pentagon class]
\label{prop:vol2-directed-pentagon}
The element
\begin{align}\label{eq:vol2-directed-pentagon}
 S={}&m_{034}\chi_{0123}+m_{123}\chi_{0134}
           +m_{014}\chi_{1234}\notag\\
    &-m_{234}\chi_{0124}-m_{012}\chi_{0234}
\end{align}
is a nonzero class in $H^{-1}(K)$. Adjoining an even generator
$\zeta$ of degree $-2$ with $d\zeta=S$ therefore fills a
relation not already filled in $K$.
\end{proposition}

\begin{proof}
The five parenthesizations of the four-arrow path have
coefficients
\[
 \begin{array}{lll}
 P_1=m_{012}m_{023}m_{034},&P_2=m_{123}m_{013}m_{034},&
 P_3=m_{123}m_{134}m_{014},\\
 P_4=m_{234}m_{124}m_{014},&P_5=m_{012}m_{234}m_{024}.&
 \end{array}
\]
The differentials of the five displayed terms in $S$ are
$P_1-P_2,P_2-P_3,P_3-P_4,P_4-P_5,P_5-P_1$, so $dS=0$.
Every element of $K^{-2}$ is a polynomial combination of
products $\chi_I\chi_J$. Its differential has coefficients
of $m$-degree at least two. The nonzero coefficients of $S$
have $m$-degree one, so $S$ cannot be a boundary.
Finally $d^2\zeta=dS=0$; its even degree requires the
polynomial algebra $k[\zeta]$.
\end{proof}

The obstruction fixes the lower homotopies. In
$A=\Lambda(\varepsilon)$ with $|\varepsilon|=-1$ and $d=0$,
set every $m_{ijk}=1$, set $\chi_{0123}=\varepsilon$ and
all other $\chi$ to zero. All $F_{ijkl}$ vanish, giving a
map $K\to A$. It cannot extend to $K[\zeta]$, since $S$
maps to $\varepsilon$ and $A^{-2}=0$. Setting every $\chi$
to zero would extend. Thus the example obstructs prescribed
homotopies, not all coherent structures with those binary
products.

\subsection{The coefficient algebra in every arity}

For $N$ ordered objects $0,\ldots,N-1$, introduce a variable
$\beta_I$ of degree $2-n$ for every increasing sequence
$I=(i_0<\cdots<i_n)$, $n\geq2$. Let $R_N$ be their free
graded-commutative algebra, with
\begin{equation}\label{eq:vol2-directed-all-arity}
 d\beta_I=-\sum_{s=2}^{n-1}\sum_{r=0}^{n-s}
 (-1)^{r+(r+1)(2-s)}
 \beta_{(i_r,\ldots,i_{r+s})}
 \beta_{(i_0,\ldots,i_r,i_{r+s},\ldots,i_n)}.
\end{equation}
The inner coefficient is written first, fixing the sign
convention. The sum is empty on binary coefficients.
There are $2^N-1-N-\binom N2$ generators.

\begin{theorem}\label{thm:vol2-directed-all-arity}
Equation \eqref{eq:vol2-directed-all-arity} defines a
differential. For every commutative dg algebra $A$, maps
$R_N\to A$ specify strictly unital directed $A_\infty$
structures on the fixed labelled rank-one carrier, with
unary differential that of $A$. For $N=5$, the substitutions
\[
 \beta_{ijk}=m_{ijk},\qquad
 \beta_{ijkl}=-\chi_{ijkl},\qquad
 \beta_{01234}=-\zeta
\]
identify it with the preceding pentagon algebra.
\end{theorem}

\begin{proof}
Suspend each forward arrow to degree $-1$, denoting it by
$\bar a_{ij}$. Define the degree-one operations
\[
 b_n(\bar a_{i_0i_1}|\cdots|\bar a_{i_{n-1}i_n})
                  =\beta_I\bar a_{i_0i_n}.
\]
Keep coefficients on the left and extend them as a
coderivation of the tensor coalgebra over the object
idempotents. Its unary part differentiates coefficients.
Inserting an $s$-ary operation after $r$ inputs contributes
$(-1)^r$. Moving its coefficient of degree $2-s$ past those
inputs and then through the outer degree-one operation
contributes $(-1)^{(r+1)(2-s)}$. Thus the component of
$b^2$ with output length one is $d\beta_I$ plus the
insertion sum in \eqref{eq:vol2-directed-all-arity}.

Apply the prescribed derivation once more. Each resulting
term is a planar substitution with three operation vertices.
Its two internal edges can be contracted in exactly two
orders. For nested substitutions these are the two
associative orders of composing graded maps; for disjoint
substitutions they interchange two odd insertions. The
graded Leibniz rule gives opposite signs and the same
coefficient permutation in each pair. Hence every term
cancels and $d^2\beta_I=0$.

A map $R_N\to A$ therefore gives exactly the coefficients
with $b^2=0$. No nonidentity composable string has length
greater than $N-1$. Normalize away strings with an identity
input and impose the ordinary unit products; these give
the strictly unital structure. Conversely, its coefficients
give the map. Substitution at arities three and four gives
$d\chi=F$ and $d\zeta=S$ with the displayed signs.
\end{proof}

Polynomial differential forms on simplices can be used as
test algebras to retain paths and higher simplices of these
structures. The carrier and its bases stay fixed. Taking
a quotient by changes of basis or by general $A_\infty$
equivalences is a separate moduli problem. Nor does the
presentation assert that $R_N$ resolves the ordinary
associative-parameter ring at every singular point.

\subsection{An algebraic contraction of the invertible chart}

Assume $N\geq3$ and invert every binary coefficient:
\[
 R_N^\times=R_N[m_{ijk}^{-1}:i<j<k],\qquad
 q_{ij}=m_{0ij},\qquad
 L_N=k[q_{ij}^{\pm1}:1\leq i<j\leq N-1].
\]
These are algebraic localizations; no completion is used.
The first object supplies a contracting partner for every
coefficient on the other objects.

\begin{theorem}\label{thm:vol2-directed-invertible-chart}
There are dg algebra maps
$L_N\xrightarrow{\iota}R_N^\times\xrightarrow{p}L_N$,
with $p\iota=1$, and an $L_N$-linear degree-minus-one
homotopy satisfying $dH+Hd=1-\iota p$.
The inclusion sends $q_{ij}$ to $m_{0ij}$, and
\begin{equation}\label{eq:vol2-directed-retraction}
 p(m_{ijk})=\frac{q_{ij}q_{jk}}{q_{ik}}\quad(0<i<j<k),
 \qquad p(\beta_I)=0\quad(|I|\geq4).
\end{equation}
Thus $H^0(R_N^\times)=L_N$ and all other cohomology
vanishes, for every finite $N$.
\end{theorem}

\begin{proof}
For each increasing $I=(i_1,\ldots,i_\ell)$ with $i_1>0$
and $\ell\geq3$, set
\begin{equation}\label{eq:vol2-directed-contracting-pairs}
 \eta_I=q_{i_1i_\ell}^{-1}\beta_{0I},\qquad
 z_I=q_{i_1i_\ell}^{-1}d\beta_{0I}.
\end{equation}
Then $d\eta_I=z_I$ and $dz_I=0$.
In the insertion formula for $d\beta_{0I}$, the term
$r=1,s=\ell-1$ is $+\beta_Iq_{i_1i_\ell}$: the exponent
$1+2(3-\ell)$ is odd and the formula has an initial minus
sign. Every other coefficient with no zero index has fewer
than $\ell$ indices. Coefficients containing zero are
already $q$'s or unit multiples of $\eta$'s. Thus
$z_I=\beta_I+$ terms of smaller nonzero index size.
Increasing induction on $\ell$ inverts this triangular
polynomial change of generators.

For $I=(i,j,k)$,
$z_I=m_{ijk}-c_I$, with $c_I=q_{ij}q_{jk}/q_{ik}$ a
unit of $L_N$. Consequently $R_N^\times$ is the tensor
product over $L_N$ of the factors
\[
 A_I=
 \begin{cases}
 L_N[z_I,(z_I+c_I)^{-1}]\otimes\Lambda(\eta_I),&|I|=3,\\
 \operatorname{Sym}_{L_N}(L_Nz_I\oplus L_N\eta_I),&|I|\geq4,
 \end{cases}
\]
with $|z_I|=3-|I|$, $|\eta_I|=2-|I|$.
Evaluation at both variables zero is defined in every
factor, since each localized $c_I$ is a unit.

For $|I|=3$, a contraction of evaluation is
\[
 h(f(z))=\frac{f(z)-f(0)}z\eta,\qquad h(f(z)\eta)=0.
\]
The quotient belongs to the same localized ring because
the evaluation kernel is the ideal $(z)$.
For odd $|I|\geq5$, use
$h(z^a)=z^{a-1}\eta$ for $a\geq1$, $h(1)=0$ and
$h(z^a\eta)=0$. For even $|I|\geq4$, use
\[
 h(\eta^az)=\frac{\eta^{a+1}}{a+1},\qquad
 h(\eta^a)=0\quad(a\geq0).
\]
The graded Leibniz rule gives $dh+hd=1-\iota p$ in
each case, with $h\iota=ph=h^2=0$.
The denominators are invertible in characteristic zero.

Order the finite set of factors and take the signed
tensor homotopy
$H=\sum_j(\iota_1p_1)\otimes\cdots\otimes
(\iota_{j-1}p_{j-1})\otimes h_j\otimes1\otimes\cdots$.
The contraction equation telescopes, as in
Proposition~\ref{prop:vol2-independent-normal-support}.
After the binary substitutions, induction through the
triangular coordinates sends every higher coefficient
to zero. Evaluation is therefore exactly
\eqref{eq:vol2-directed-retraction}.
\end{proof}

The inclusion and projection are algebra maps; the
displayed homotopy is a homotopy of complexes. It also
determines the tangent complex before changes of arrow
basis. At the point with every $m_{ijk}=1$ and every higher
coefficient zero, let $T^\bullet$ be the dual of the
coefficient cotangent complex. Then
\begin{equation}\label{eq:vol2-directed-tangent-ranks}
 \dim T^j=\binom N{j+3},\qquad
 \operatorname{rank}(d:T^j\to T^{j+1})=\binom{N-1}{j+3}
 \quad(j\geq0),
\end{equation}
with out-of-range binomial coefficients zero. Indeed the
cotangent complex is obtained by differentiating the free
localized presentation. Each pair
\eqref{eq:vol2-directed-contracting-pairs} gives an identity
from tangent degree $|I|-3$ to $|I|-2$, and the
$\binom{N-1}{2}$ variables $q_{ij}$ remain unpaired.
Thus its cohomology is $k^{\binom{N-1}{2}}$ in degree zero
and zero otherwise. For $N=5$, the dimensions are
$(10,5,1)$ and ranks $(4,1)$; for $N=6$, they are
$(20,15,6,1)$ and $(10,5,1)$.

\subsection{Ordinary centres on two loci}

The tangent complex of the coefficient space differs from
the ordinary Hochschild complex of a category at a chosen
point. Let $C$ be the directed category on these $N$ objects,
with scalar identities and a single degree-zero basis arrow
$a_{ij}$ for each $i<j$.

\begin{proposition}\label{prop:vol2-directed-centre-loci}
At an associative point with every $m_{ijk}\ne0$,
$HH^0(C)=k$ and $HH^{>0}(C)=0$.
At the zero-product point,
\[
 \begin{gathered}
 \dim HH^0(C)=1,\qquad \dim HH^1(C)=\binom{N-1}{2},\\
 \dim HH^j(C)=\binom N{j+1}\quad(2\leq j\leq N-1),
 \end{gathered}
\]
and the higher groups vanish.
\end{proposition}

\begin{proof}
Associativity on $(0,i,j,k)$ gives
$m_{ijk}=q_{ij}q_{jk}/q_{ik}$. Replace $a_{ij}$ by
$q_{ij}^{-1}a_{ij}$ when $i>0$ and leave $a_{0i}$ unchanged.
Every composable product then has coefficient one.
The relative bar of Theorem~\ref{thm:vol2-directed-bar}
identifies its normalized Hochschild cochains with the
simplicial cochains of the ordered $(N-1)$-simplex.
For a positive-degree cochain put
\[
 (hf)(i_0,\ldots,i_{n-1})=
 \begin{cases}
 f(0,i_0,\ldots,i_{n-1}),&i_0>0,\\
 0,&i_0=0.
 \end{cases}
\]
In $dh+hd$, terms deleting any vertex other than the
inserted zero cancel. The remaining term is $f$ in positive
degree and $f-f(0)$ in degree zero. This contracts onto
the constants.

At zero products, every differential from degree one
upward is zero: each merger and endpoint action multiplies
two forward arrows. The degree-zero map is the difference
of vertex values on every forward edge. Its kernel is
$k$ and its rank $N-1$. Degree $j$ has one coordinate for
each increasing $(j+1)$-tuple. Subtracting that rank in
degree one gives the stated dimensions.
\end{proof}

At the zero-product point, the coefficient differential
has no linear part, so its tangent complex retains every
generator. Neither this tangent complex nor its ordinary
centre is obtained by substituting zero into the localized
contraction.

\section{The four-input comparison}
\label{sec:vol2-four-input-comparison}

The preceding constructions solve different composition
problems. The discarded coefficient complex retains outer
and crossed collision classes. The Heisenberg construction
gives actual pair-support outer maps and a triple homotopy
on regular singleton inputs. The directed coefficient
algebra retains all arities on its fixed rank-one carrier
and contracts on its invertible chart. A comparison with
field operations must supply the maps relating these
domains.

At four inputs, first construct a common operation complex
which retains the local vacuum and every nonconsecutive
collision channel. Transport the five parenthesizations
and their chosen associativity homotopies into it, including
the operations composing a homotopy with the remaining
inputs. Their pentagon then gives a specified cycle.
The unresolved equation is to construct a degree-minus-two
primitive of that cycle, or to prove that its class is
nonzero, with those lower homotopies fixed.

The directed algebra includes its higher generators and
their prescribed differentials. Theorem~\ref{thm:vol2-directed-invertible-chart}
computes its localized cohomology; it does not assert that
every choice of lower homotopies in another coefficient
algebra extends. The common supported operation complex and
the comparison into it remain to be constructed. Likewise
the three-channel scalar
current identity \eqref{eq:vol2-three-channel-current}
retains the contact terms needed by a field comparison,
but its residue relation alone is not that operation
homotopy. These missing maps preserve the distinction
between a coefficient calculation and a coherent action
on fields.

\chapter{Gaussian elimination and higher state operations}
\label{ch:vol2-gaussian-elimination}

Let $p$ be expectation for a centred real Gaussian variable of
variance $\sigma>0$. Then $p(x)=0$ and $p(x^2)=\sigma$.
Eliminating $x$ therefore leaves information in products that its
mean does not contain. The connected contractions of
Chapter~\ref{ch:vol2-supported-compositions} exhibit the same effect
for fields. A reduction that preserves their operations must retain
these fluctuations together with the differential and the
normalization of the state.

In a finite polynomial model, an explicit homotopy identifies
each observable with its expectation. The failure to preserve
products becomes part of the comparison map. Formal interactions
change both the expectation and the differential of the variables
that remain. Successive integrations must preserve this changed
differential and the state on every product. We construct these
maps and compare them under changes of coordinates and of the
integration data.

The normal direction of a boundary comparison has an additional
datum: its endpoint condition. We first isolate that elementary
support question. The finite quantum construction that follows
does not require a choice of continuum propagator.

\section{A supported normal contraction}
\label{sec:vol2-normal-contraction}

Put $I=[0,\infty)$ and work over $\mathbb C$. The relative compactly
supported de Rham complex
\[
 C=\Omega_c^\bullet(I,\{0\})
\]
has $C^0=\{f\in C_c^\infty(I):f(0)=0\}$ and
$C^1=C_c^\infty(I)\,dt$, with $df=f'\,dt$.
Thus the relative condition constrains functions at the endpoint;
it imposes no restriction on the coefficient of a one-form there.
Choose $\rho\in C_c^\infty(I)$ with
$\int_0^\infty\rho(t)\,dt=1$.

Define maps to and from $\mathbb C[-1]$, placed in cohomological
degree one, by
\begin{equation}\label{eq:vol2-normal-maps}
 \begin{aligned}
 p(\beta\,dt)&=\int_0^\infty\beta(t)\,dt,& p(f)&=0,\\
 i(a)&=a\rho(t)\,dt,& H(f)&=0,\\
 H(\beta\,dt)(t)
   &=\int_0^t\left(\beta(s)-\rho(s)
                    \int_0^\infty\beta(r)\,dr\right)ds.
 \end{aligned}
\end{equation}
The subtraction in the primitive removes its constant tail.
It is required to keep the homotopy in the stated support class.

\begin{proposition}\label{prop:vol2-normal-retract}
The maps in \eqref{eq:vol2-normal-maps} are continuous and satisfy
\[
 pi=1,\qquad dH+Hd=1-ip,\qquad
 H^2=Hi=pH=0.
\]
The normal support of $H\beta$ lies in the compact interval spanned
by $0$, $\operatorname{supp}\beta$ and $\operatorname{supp}\rho$.
Tensoring this contraction with the full Dolbeault complex on an
open set $U\subset\mathbb C^d$ retains every tangential form degree.
\end{proposition}
\begin{proof}
The integrand defining $H\beta$ has total integral zero. Its
primitive vanishes at zero and again beyond the stated interval.
Differentiation gives $dH(\beta\,dt)=\beta\,dt-ip(\beta\,dt)$.
For $f\in C^0$, compact support and $f(0)=0$ give
\[
 p(df)=f(\infty)-f(0)=0,\qquad H(df)=f.
\]
These are the homotopy identities in the two degrees. The equation
$pi=1$ is the normalization of $\rho$; substituting $i(a)$ in $H$
gives $Hi=0$. The remaining side identities follow from degrees.

On functions or forms supported in a fixed compact interval,
integration is bounded by a finite smooth seminorm. The same
estimate for the primitive and its derivatives proves continuity
into the enlarged compact interval. These estimates give
continuity for the usual inductive-limit topology of compact
supports. On
$C\widehat\otimes_\pi\Omega^{0,\bullet}(U)$, the completed
projective tensor product, the maps extend by continuity.
The tensor differential is
\[
 d(\alpha\otimes\gamma)
 =d\alpha\otimes\gamma+(-1)^{|\alpha|}
                         \alpha\otimes\bar\partial\gamma.
\]
Its cross terms with $H\otimes1$ cancel. The resulting target is
$\Omega^{0,\bullet}(U)[-1]$, with differential $-\bar\partial$,
as required by the shift.
\end{proof}

Changing the endpoint condition changes the answer. In the
absolute complex $\Omega_c^\bullet(I)$, the primitive
\[
 H_{\mathrm{abs}}(\beta\,dt)(t)
   =-\int_t^\infty\beta(s)\,ds
\]
is allowed even when it is nonzero at $t=0$. It is compactly
supported and satisfies $dH_{\mathrm{abs}}=1$ on one-forms and
$H_{\mathrm{abs}}d=1$ on functions. The absolute complex is
acyclic. In a field theory, the polarization must determine which
normal complex occurs before either contraction is used.

\section{The finite Gaussian Ward complex}
\label{sec:vol2-gaussian-ward}

Let $R_0$ be a nonzero commutative $\mathbb Q$-algebra and put
$R=R_0[\hbar]$, where $\hbar$ has degree zero. Fix a symmetric
invertible matrix $A\in\operatorname{Mat}_m(R_0)$ and write
$C=A^{-1}$. Introduce even coordinates $q_1,\ldots,q_m$ of
degree zero and odd coordinates $\eta_1,\ldots,\eta_m$ of
degree $-1$. Our observable algebra is
\[
 \mathcal G_A=R[q_1,\ldots,q_m]\otimes
                         \Lambda(\eta_1,\ldots,\eta_m).
\]
All odd derivatives act on the left. Tensor products and
reorderings use the cohomological Koszul sign.

The quadratic potential $S_0(q)=\tfrac12q^{\mathrm t}Aq$
defines the \emph{Ward differential}
\begin{equation}\label{eq:vol2-gaussian-ward}
 D_A=\sum_i\bigl((Aq)_i-\hbar\partial_{q_i}\bigr)
                                             \partial_{\eta_i}.
\end{equation}
This is the finite-dimensional homological form of Gaussian
integration by parts; see Gwilliam--Johnson-Freyd
\cite{GwilliamJohnsonFreyd2019}*{\S\S1--3}.
The even operators $(Aq)_i-\hbar\partial_{q_i}$ commute, because
$A$ is symmetric. The odd derivatives anticommute, so $D_A^2=0$.
In one variable, with $A=(a)$,
\begin{equation}\label{eq:vol2-ward-square}
 D_A\eta=ax,\qquad
 D_A(\eta x)=ax^2-\hbar.
\end{equation}
The second equality already records the fluctuation that survives
elimination.

We will construct a \emph{contraction} of this complex onto $R$:
cochain maps $i:R\to\mathcal G_A$, $p:\mathcal G_A\to R$ and
a degree-minus-one map $H$ satisfying
\[
 pi=1,\qquad D_AH+HD_A=1-ip,
 \qquad H^2=Hi=pH=0.
\]
The last three equations are the side conditions of a strong
deformation retract. The projection $p$, normalized by $p(1)=1$,
will be the formal Gaussian state.

\subsection{Heat conjugation and the Euler homotopy}

Let
\[
 d_0=\sum_i(Aq)_i\partial_{\eta_i},\qquad
 T_A=\exp\!\left(\frac{\hbar}{2}
                 \sum_{i,j}C_{ij}\partial_{q_i}\partial_{q_j}\right).
\]
The heat series and its inverse terminate on every polynomial.
The commutator of its exponent with $q_i$ is
$\hbar\sum_jC_{ij}\partial_{q_j}$, and the next commutator is zero.
Consequently
\begin{equation}\label{eq:vol2-ward-conjugation}
 T_A^{-1}q_iT_A=q_i-\hbar\sum_jC_{ij}\partial_{q_j},
 \qquad D_A=T_A^{-1}d_0T_A.
\end{equation}

Let $N$ count total polynomial degree in the $q_i$ plus exterior
degree in the $\eta_i$, and put
\[
 R_A=\sum_{i,j}\eta_iC_{ij}\partial_{q_j},\qquad
 h_0=N^{-1}R_A\quad(N>0),\qquad h_0(1)=0.
\]
The notation $N^{-1}$ divides a homogeneous monomial by its
positive integer degree. These divisions are valid over $R_0$.
Let $\epsilon$ set every $q_i,\eta_i$ equal to zero, and let
$i$ include constants.

\begin{lemma}\label{lem:vol2-euler-contraction}
The maps $(\epsilon,i,h_0)$ contract $(\mathcal G_A,d_0)$ onto $R$.
\end{lemma}
\begin{proof}
The anticommutator of the odd derivations $d_0$ and $R_A$ is the
Euler operator:
\[
 d_0R_A+R_Ad_0
   =\sum_iq_i\partial_{q_i}
      +\sum_i\eta_i\partial_{\eta_i}=N.
\]
Indeed, its value on each $q_i$ and $\eta_i$ is that generator;
an even derivation is determined by these values. Both odd
operators preserve $N$. Division by $N$ proves
$d_0h_0+h_0d_0=1$ in positive total degree, while both sides of
$1-i\epsilon$ vanish on constants. The constant-coefficient even
derivatives in $R_A$ commute and its odd coefficients
anticommute, giving $R_A^2=0$. Hence $h_0^2=0$.
The equations $h_0i=\epsilon h_0=0$ follow from the same formulas.
\end{proof}

\begin{theorem}\label{thm:vol2-gaussian-retract}
The maps
\begin{equation}\label{eq:vol2-gaussian-retract}
 p_A=\epsilon T_A,\qquad i(a)=a,\qquad
 H_A=T_A^{-1}h_0T_A
\end{equation}
give a contraction of $(\mathcal G_A,D_A)$ onto $R$.
On even polynomials, $p_A$ is the Wick functional of covariance
$\hbar A^{-1}$; it kills every term of negative cohomological degree.
All maps are defined over $R_0[\hbar]$.
\end{theorem}
\begin{proof}
Conjugate every identity of Lemma~\ref{lem:vol2-euler-contraction}
by $T_A$. Equation~\eqref{eq:vol2-ward-conjugation} gives the
new differential, and $T_Ai=i$ keeps the inclusion fixed.
Thus the homotopy and all side identities follow on every
polynomial, without a convergence argument.

To identify the state, introduce independent degree-zero formal
variables $J_1,\ldots,J_m$. Applying $\epsilon T_A$
coefficientwise gives
\begin{equation}\label{eq:vol2-wick-generating}
 p_A(e^{J^{\mathrm t}q})
       =\exp\!\left(\frac{\hbar}{2}J^{\mathrm t}CJ\right).
\end{equation}
Differentiation in the $J_i$ pairs all input coordinates, with
$\hbar C_{ij}$ for each pair. This proves Wick's formula and
determines $p_A$ on every polynomial. Heat operators preserve
exterior degree, so evaluation kills negative cohomological degrees.
\end{proof}

For $A=(a)$, the retract gives the concrete primitive
\begin{equation}\label{eq:vol2-quadratic-primitive}
 p_A(x^2)=\frac{\hbar}{a},\qquad
 H_A(x^2)=\frac{\eta x}{a},\qquad
 D_A\left(\frac{\eta x}{a}\right)=x^2-\frac{\hbar}{a}.
\end{equation}
At a positive real specialization $A>0$, $\hbar>0$,
equation~\eqref{eq:vol2-wick-generating} is also the generating
function of the normalized convergent Gaussian integral.
Positivity is needed for that real integral, and is not a
hypothesis of the algebraic contraction.

\subsection{Ordinary and transported products}

The ordinary product on $\mathcal G_A$ is not a cochain operation:
\eqref{eq:vol2-ward-square} gives
\[
 D_A(\eta x)-(D_A\eta)x=-\hbar.
\]
One can transport a product along the heat isomorphism. It is
useful to do this explicitly, since the transported square has a
different expectation.

\begin{proposition}\label{prop:vol2-transported-product}
In one variable define
\begin{equation}\label{eq:vol2-transported-product}
 f\diamond g=T_A^{-1}\bigl((T_Af)(T_Ag)\bigr)
 =\sum_{r\geq0}\frac{(-\hbar/a)^r}{r!}
                  (\partial_x^rf)(\partial_x^rg).
\end{equation}
This is a unital graded-commutative associative product for which
$D_A$ is a derivation. The map $p_A$ is an augmentation for
$\diamond$ and a quasi-isomorphism. In particular,
\[
 x\diamond x=x^2-\hbar/a,\qquad
 p_A(x\diamond x)=0,\qquad p_A(x^2)=\hbar/a.
\]
\end{proposition}
\begin{proof}
The differential $d_0=ax\partial_\eta$ is a derivation for the
ordinary product. Transporting that differential and product by
$T_A^{-1}$ proves the algebra and differential identities.
Evaluation at zero is an augmentation of the classical Koszul
complex, so its transported map is $p_A$.

To obtain the sum in \eqref{eq:vol2-transported-product}, apply
the heat operators to two independent variables $u,v$ before
setting $u=v=x$. The second derivative of a product is induced
by $(\partial_u+\partial_v)^2$. Its two individual square terms
cancel those of $T_A$ on the factors, leaving
$\exp(-\hbar\partial_u\partial_v/a)$. All these series terminate
on polynomials. The displayed square follows at once, and the
quasi-isomorphism follows from Theorem~\ref{thm:vol2-gaussian-retract}.
\end{proof}

For the ordinary observables, the variance in
\eqref{eq:vol2-quadratic-primitive} still has to be retained.
Its higher analogues give a comparison of products without
changing which polynomial is being observed.

\section{Cumulants as operation maps}
\label{sec:vol2-state-cumulants}

Let $B$ be a unital graded-commutative algebra over a commutative
$\mathbb Q$-algebra. Write
\[
 S^c(B)=\bigoplus_{n\geq0}\operatorname{Sym}^n(B)
\]
for the coaugmented conilpotent symmetric coalgebra of its
underlying graded module, with the unshuffle coproduct.
A coaugmentation-preserving coalgebra map out of $S^c(B)$ is
determined by its \emph{Taylor coefficients}: its composites
from each symmetric length to the length-one component of the
target. The scalar in length zero and the algebra unit
$1_B\in B$, which lies in length one, are distinct.

The products of $B$ determine a coalgebra automorphism $\tau_B$
whose $n$th Taylor coefficient is
\[
 (\tau_B)_n(b_1\odot\cdots\odot b_n)=b_1\cdots b_n,
 \qquad n\geq1.
\]
On an input of length $n$, it sums over partitions of
$\{1,\ldots,n\}$: multiply within each block and take the
symmetric tensor of the block products. Within a block the
indices are increasing; blocks are ordered by their least
indices. Reordering homogeneous inputs into this order gives
the Koszul sign. The term with singleton blocks is the identity,
and every other term lowers length. Thus the inverse is finite
on each input.

\begin{lemma}\label{lem:vol2-cumulant-inverse}
The $n$th Taylor coefficient of $\tau_B^{-1}$ is
\[
 (\tau_B^{-1})_n=(-1)^{n-1}(n-1)!\,\mu_n,
\]
where $\mu_n$ is multiplication of $n$ elements.
\end{lemma}
\begin{proof}
Corestrict the proposed composite $\tau_B^{-1}\tau_B$ to length
one. All partitions into $k$ blocks have the same coefficient
$(-1)^{k-1}(k-1)!$ after their products are multiplied.
The resulting scalar is
\[
 \sum_{k=1}^n
 \left\{\begin{matrix}n\\k\end{matrix}\right\}
                       (-1)^{k-1}(k-1)! =
 \begin{cases}1,&n=1,\\0,&n>1,\end{cases}
\]
where the braces count set partitions into $k$ nonempty blocks.
This identity is the coefficient identity of
$\log(e^z)=z$. The same regrouping supplies the Koszul signs
of the graded products. A coalgebra map is determined by these
corestrictions, proving the claimed inverse.
\end{proof}

\subsection{The deviations of a differential}

Suppose $D:B\to B$ has degree one, $D^2=0$ and $D1_B=0$.
It need not be a derivation. Let $\widetilde D$ be the unary
coderivation on $S^c(B)$: it applies $D$ to one factor at a
time, with the tensor signs. Put
\begin{equation}\label{eq:vol2-deviation-coderivation}
 \mathscr D_B=\tau_B^{-1}\widetilde D\tau_B.
\end{equation}
This is again a degree-one square-zero coderivation. Such a
coderivation with zero constant Taylor coefficient is an
$L_\infty[1]$ structure; its operations are graded symmetric
and have degree one.

\begin{proposition}\label{prop:vol2-multiplicative-deviations}
The first Taylor coefficients of \eqref{eq:vol2-deviation-coderivation}
are
\[
 q_1(b)=Db,\qquad
 q_2(b,c)=D(bc)-(Db)c-(-1)^{|b|}b(Dc).
\]
If $D$ has differential-operator order at most two, all Taylor
coefficients of arity at least three vanish.
\end{proposition}
\begin{proof}
In length two, the product term in $\tau_B$ followed by $D$
contributes $D(bc)$. The length-two part of $\tau_B^{-1}$
subtracts the two unary terms, giving $q_2$.
More generally, the Taylor coefficient is the iterated graded
commutator of $D$ with the multiplication operators of the
inputs, evaluated at $1_B$. By definition, an operator of order
at most two has zero commutators with three or more
multiplication operators. The square-zero identities for the
remaining operations follow from conjugation in
\eqref{eq:vol2-deviation-coderivation}.
\end{proof}

For $D_A$, the binary deviation is divisible by $\hbar$.
Equation~\eqref{eq:vol2-ward-square} says
$q_2(\eta,x)=-\hbar$ in one variable. This second-order
operation is part of the Ward complex.

\subsection{The higher components of a normalized map}

Let $P:(B,D_B)\to(C,D_C)$ be a degree-zero cochain map between
unital graded-commutative algebras, with $P1_B=1_C$.
The linear map need not preserve products. Define
\begin{equation}\label{eq:vol2-cumulant-coalgebra}
 \mathscr K(P)=\tau_C^{-1}S^c(P)\tau_B.
\end{equation}
Its $n$th Taylor coefficient is the \emph{cumulant}
\begin{equation}\label{eq:vol2-state-cumulant}
 \kappa_n^P(b_1,\ldots,b_n)=
 \sum_{\pi\in\Pi_n}(-1)^{|\pi|-1}(|\pi|-1)!
 \epsilon(\pi;b)
 \prod_{E\in\pi}P\!\left(\prod_{j\in E}b_j\right).
\end{equation}
Here $\Pi_n$ is the set of partitions into nonempty blocks and
$\epsilon(\pi;b)$ is the sign of the grouping specified above.
In particular,
\[
 \kappa_1^P=P,\qquad
 \kappa_2^P(b,c)=P(bc)-P(b)P(c).
\]
This coalgebra construction and its compatibility with
Koszul brackets are described by Bandiera
\cite{Bandiera2020}*{\S1.2, Lemma 1.16 and Proposition 1.21}.

\begin{theorem}\label{thm:vol2-cumulant-comparison}
The map $\mathscr K(P)$ is an $L_\infty[1]$ morphism between
the multiplicative deviations of $D_B$ and $D_C$:
\[
 \mathscr D_C\mathscr K(P)=\mathscr K(P)\mathscr D_B.
\]
If $P$ is a quasi-isomorphism, this is an $L_\infty[1]$
quasi-isomorphism. Every cumulant of arity greater than one
vanishes on an algebra-unit input. For a further normalized
cochain map $Q:C\to E$,
\begin{equation}\label{eq:vol2-cumulant-sewing}
 \mathscr K(QP)=\mathscr K(Q)\mathscr K(P).
\end{equation}
\end{theorem}
\begin{proof}
Inserting \eqref{eq:vol2-cumulant-coalgebra} and cancelling
$\tau_C\tau_C^{-1}$ reduces the morphism equation to
$\widetilde D_CS^c(P)=S^c(P)\widetilde D_B$.
This is precisely the cochain equation for $P$, applied to
each tensor factor. The linear Taylor coefficient is $P$;
being a quasi-isomorphism in that coefficient is the definition
of an $L_\infty$ quasi-isomorphism.

The unit assertion can be seen by adding independent marked
variables of the degrees that make a sum $F$ of marked inputs
have total degree zero. The cumulants are the coefficients of
$\log P(e^F)$. Adding a scalar mark $z1_B$ changes this expression
to $z+\log P(e^F)$ because $P1_B=1_C$.
All mixed coefficients involving $z$ therefore vanish.
This is a finite coefficient identity in any specified arity.

Finally the two adjacent multiplication automorphisms in
$\mathscr K(Q)\mathscr K(P)$ cancel. The remaining composite
is \eqref{eq:vol2-cumulant-coalgebra} for $QP$, proving
\eqref{eq:vol2-cumulant-sewing} in every arity.
\end{proof}

At arity two, \eqref{eq:vol2-cumulant-sewing} is the law of
total covariance:
\begin{equation}\label{eq:vol2-total-covariance}
 \kappa_2^{QP}(b,c)
   =Q\bigl(\kappa_2^P(b,c)\bigr)
                       +\kappa_2^Q(Pb,Pc).
\end{equation}
The second term measures fluctuations of the variables left
after the first elimination.

\subsection{A Gaussian state has higher operations on composites}

Take one variable of stiffness $a$ and write $p=p_A$.
For even polynomials $f,g_2,\ldots,g_n$, the morphism equation
of Theorem~\ref{thm:vol2-cumulant-comparison} on
$\eta f,g_2,\ldots,g_n$ gives, for $n\geq2$,
\begin{equation}\label{eq:vol2-higher-ward}
 \begin{split}
 \kappa_n^p(axf-\hbar f',g_2,\ldots,g_n)
  =\hbar\sum_{j=2}^n
    \kappa_{n-1}^p(fg_j',g_2,\ldots,\widehat g_j,\ldots,g_n).
 \end{split}
\end{equation}
Indeed the only binary input terms are
$q_2(\eta f,g_j)=-\hbar fg_j'$, and the target differential
and its deviations are zero. For $n=2$, $f=1$, $g_2=x$,
the equation forces $a\kappa_2^p(x,x)=\hbar$.
The variance is required by the cochain equation itself.

Although the higher cumulants of the linear variable vanish
above arity two, those of its square do not.
Wick's formula gives the formal generating function
\[
 p(e^{tx^2})=(1-2t\hbar/a)^{-1/2}.
\]
Taking its logarithm and comparing coefficients of $t^n/n!$
proves
\begin{equation}\label{eq:vol2-square-cumulants}
 \kappa_n^p(x^2,\ldots,x^2)
   =2^{n-1}(n-1)!\left(\frac{\hbar}{a}\right)^n,
 \qquad n\geq1.
\end{equation}
Thus the state on the full polynomial algebra has nonzero
components in arbitrarily high arity. Its Gaussian character
does not truncate the operations of composite observables.

\section{Formal interactions and conditional elimination}
\label{sec:vol2-interacting-elimination}

Introduce a degree-zero coupling parameter $g$. We work in
\[
 \bigl(R_0[\hbar,q]\otimes\Lambda\eta\bigr)[[g]].
\]
This means a polynomial in the even and odd coordinates at
each coupling order. Polynomial degree may grow with that order;
there is no completion in the field coordinates themselves.
The positive $g$-adic filtration makes the following perturbation
construction well defined.

\subsection{A filtered perturbation identity}

\begin{lemma}\label{lem:vol2-filtered-perturbation}
Let $(i,p,h)$ be a contraction of a complex $(C,d)$, complete
and separated for a descending filtration, onto $(B,d_B)$, with
\[
 dh+hd=1-ip,\qquad pi=1,\qquad h^2=hi=ph=0.
\]
Assume all these maps preserve filtration. Let $\delta$ have
degree one, strictly raise filtration, and satisfy
$(d+\delta)^2=0$. Put
\[
 F=(1+\delta h)^{-1}\delta=\delta(1+h\delta)^{-1}.
\]
Then a contraction for $d+\delta$ is given by
\begin{equation}\label{eq:vol2-perturbation-formulas}
 \begin{aligned}
 d_B'&=d_B+pFi,&
 i'&=i-hFi=(1+h\delta)^{-1}i,\\
 p'&=p-pFh=p(1+\delta h)^{-1},&
 h'&=h-hFh=h(1+\delta h)^{-1}.
 \end{aligned}
\end{equation}
In particular, $(d_B')^2=0$, $p'i'=1$ and
\[
 (d+\delta)h'+h'(d+\delta)=1-i'p',\qquad
 (h')^2=h'i'=p'h'=0.
\]
\end{lemma}
\begin{proof}
All inverses are geometric series. At any fixed filtration
precision only finitely many perturbation factors occur.
On endomorphisms write
$\partial T=dT-(-1)^{|T|}Td$.
The hypotheses give $\partial\delta=-\delta^2$ and
$\partial h=1-ip$. Apply $\partial$ to
$F=\delta-\delta hF$. The graded product rule gives
\[
 (1+\delta h)\partial F=-\delta ipF,
 \qquad\text{hence}\qquad
 dF+Fd=-FipF.
\]

It follows that
\[
 (d_B')^2=p(dF+Fd+FipF)i=0.
\]
Substitution in the inclusion equation gives
\[
 (d+\delta)i'-i'd_B'=h(dF+Fd+FipF)i=0.
\]
Using both $F=\delta-\delta hF$ and
$F=\delta-Fh\delta$, the homotopy equation expands to
\begin{align*}
 (d+\delta)h'+h'(d+\delta)
 &=1-ip+ipFh+hFip+h(dF+Fd)h\\
 &=1-(i-hFi)(p-pFh).
\end{align*}
The original side conditions give $p'i'=1$ and
$(h')^2=h'i'=p'h'=0$ directly. Commuting the homotopy
equation with $d+\delta$ and using the inclusion equation
gives $i'(d_B'p'-p'(d+\delta))=0$. Apply $p'$ to obtain
the projection equation. This proves all the contraction
identities.
\end{proof}

\subsection{An interacting state over the integral quantum ring}

Let
\[
 S(q)=\tfrac12q^{\mathrm t}Aq+V(q),\qquad
 V\in gR_0[\hbar,q][[g]].
\]
The potential is independent of the odd variables. Its Ward
differential is
\begin{equation}\label{eq:vol2-potential-ward}
 D_S=\sum_i(S_i-\hbar\partial_{q_i})\partial_{\eta_i},
 \qquad S_i=\partial_{q_i}S.
\end{equation}
The operators in parentheses commute, since
\[
 [S_i-\hbar\partial_{q_i},S_j-\hbar\partial_{q_j}]
     =\hbar(\partial_{q_j}S_i-\partial_{q_i}S_j)=0.
\]
Thus $D_S^2=0$. The perturbation
$\delta_V=\sum_iV_i\partial_{\eta_i}$ has positive coupling
order and kills constants.

\begin{theorem}\label{thm:vol2-formal-state-retract}
The Gaussian maps of Theorem~\ref{thm:vol2-gaussian-retract}
give an interacting contraction with
\begin{equation}\label{eq:vol2-formal-state-retract}
 p_V=p_A(1+\delta_VH_A)^{-1},\qquad
 i_V=i,\qquad
 H_V=H_A(1+\delta_VH_A)^{-1}.
\end{equation}
Its target is $R_0[\hbar][[g]]$ with zero differential.
All coefficients are polynomial in $\hbar$ at each coupling
order. The normalized exponential formula
\begin{equation}\label{eq:vol2-interacting-expectation}
 p_V(f)=\frac{p_A(fe^{-V/\hbar})}{p_A(e^{-V/\hbar})}
\end{equation}
agrees with this projection, with negative cohomological degrees
sent to zero. The quotient is first interpreted in
$R_0[\hbar,\hbar^{-1},q][[g]]$; its normalized value lies in
the integral quantum ring.
\end{theorem}
\begin{proof}
Apply Lemma~\ref{lem:vol2-filtered-perturbation} to the
Gaussian contraction. Since $\delta_Vi=0$, the inclusion and
target differential remain unchanged. The other formulas are
\eqref{eq:vol2-formal-state-retract}. Each coupling coefficient
contains finitely many operations on polynomials. None of these
operations divides by $\hbar$.

In the Laurent coefficient ring, the denominator of
\eqref{eq:vol2-interacting-expectation} is $1+O(g)$ and
therefore a unit. Gaussian integration by parts says
\[
 p_A\bigl((Aq)_if-\hbar\partial_{q_i}f\bigr)=0.
\]
Apply this identity to $fe^{-V/\hbar}$ at each coupling order.
The normalized quotient $L$ then satisfies
\[
 L\bigl(S_if-\hbar\partial_{q_i}f\bigr)=0,
 \qquad L(1)=1.
\]
These equations give the cochain identity on inputs with
one odd generator. With two or more odd generators,
the image under $D_S$ still has negative cohomological degree
and is killed by $L$.

The normalized cochain functional is unique. For an even
polynomial $f$, the free contraction gives
$f-ip_Af=D_AH_Af$. Hence any such functional satisfies
\[
 L(f)=p_Af-L(\delta_VH_Af).
\]
Solving recursively in coupling order gives
$L=p_A(1+\delta_VH_A)^{-1}$. This identifies the quotient
with the already integral projection.
\end{proof}

For the formal cubic potential $S=ax^2/2+gx^3/6$, the first
moments are
\begin{equation}\label{eq:vol2-cubic-state-moments}
 \begin{aligned}
 p_V(x)&=-\frac{g\hbar}{2a^2}
             -\frac{5g^3\hbar^2}{8a^5}+O(g^5),\\
 p_V(x^2)&=\frac{\hbar}{a}
             +\frac{5g^2\hbar^2}{4a^4}+O(g^4).
 \end{aligned}
\end{equation}
These follow by expanding \eqref{eq:vol2-interacting-expectation}
and using $p_A(x^{2n})=(2n-1)!!(\hbar/a)^n$.
They obey $a p_V(x)+(g/2)p_V(x^2)=0$, the Ward equation on
the odd generator. All series in this example are formal in $g$.

\subsection{The differential of the retained variables}

Split the even variables into $x=(x_1,\ldots,x_r)$ and
$y=(y_1,\ldots,y_s)$, with respective odd variables
$\eta_x,\eta_y$. Fix symmetric invertible matrices $A_x,A_y$
over $R_0$, and put
\begin{equation}\label{eq:vol2-conditional-potential}
 S(x,y)=\tfrac12x^{\mathrm t}A_xx
             +\tfrac12y^{\mathrm t}A_yy+V(x,y),
 \qquad V\in gR_0[\hbar,x,y][[g]].
\end{equation}
Write $P_x^0$ for the free Wick expectation in $x$. It kills
every term containing an $\eta_x$ and preserves the
$y,\eta_y$ variables. Define
\begin{equation}\label{eq:vol2-conditional-expectation}
 \begin{aligned}
 Z_x(y)&=P_x^0(e^{-V/\hbar}),&
 P_x(f)&=Z_x^{-1}P_x^0(fe^{-V/\hbar}),\\
 S_{\mathrm{eff}}(y)&=\tfrac12y^{\mathrm t}A_yy
                                      -\hbar\log Z_x(y).
 \end{aligned}
\end{equation}
These exponential and logarithmic expressions initially have
Laurent quantum coefficients and are complete in $g$.
The odd-variable rule is the same as for $P_x^0$.

For any potential $W$ in a tuple of even variables $q$, write
\[
 \mathcal B_W=
       \bigl(R_0[\hbar,q]\otimes\Lambda\eta_q\bigr)[[g]]
\]
with differential \eqref{eq:vol2-potential-ward}.
Its binary multiplicative deviation $q_2^W$ is divisible by
$\hbar$; denote the quotient by $\overline q_2^W$.
The operations $(D_W,\overline q_2^W)$ form an
$L_\infty[1]$ structure. Indeed the original square-zero
coderivation has only unary and binary components. Its
differential identity for the binary operation has one
factor of $\hbar$, and its Jacobi identity has two.
Multiplication by $\hbar$ is injective in these polynomial
coefficient modules, even if $R_0$ has zero divisors, so both
identities can be divided.

\begin{theorem}[Integral conditional elimination]
\label{thm:vol2-integral-conditional}
The effective potential belongs to $R_0[\hbar,y][[g]]$.
The normalized map in \eqref{eq:vol2-conditional-expectation}
is a unital cochain quasi-isomorphism
\[
 P_x:(\mathcal B_S,D_S)\longrightarrow
                       (\mathcal B_{S_{\mathrm{eff}}},D_{S_{\mathrm{eff}}})
\]
over $R_0[\hbar][[g]]$. It is the projection in the
filtered perturbation of the free Gaussian contraction.
For every $n\geq1$, its cumulant is divisible by
$\hbar^{n-1}$. The integral maps
\begin{equation}\label{eq:vol2-divided-cumulants}
 \Phi_n^{P_x}=\hbar^{1-n}\kappa_n^{P_x}
\end{equation}
give an $L_\infty[1]$ quasi-isomorphism for the divided
operations $(D_S,\overline q_2^S)$ and
$(D_{S_{\mathrm{eff}}},\overline q_2^{S_{\mathrm{eff}}})$.
They are continuous in $g$, polynomial in $\hbar$ at each
coupling order, and vanish on a unit input when $n>1$.
\end{theorem}
\begin{proof}
First work in Laurent quantum coefficients. The hidden Ward
equations are
\[
 P_x\bigl(S_{x_i}f-\hbar\partial_{x_i}f\bigr)=0.
\]
Differentiating the normalized quotient in a retained variable
gives
\begin{align}
 \partial_{y_j}P_x(f)
 &=P_x(\partial_{y_j}f)
       -\hbar^{-1}\bigl(P_x(fV_{y_j})-P_x(f)P_x(V_{y_j})\bigr),
       \label{eq:vol2-retained-covariance}\\
 (S_{\mathrm{eff}})_{y_j}
 &=(A_yy)_j+P_x(V_{y_j}).\notag
\end{align}
Combining these equations yields
\begin{equation}\label{eq:vol2-retained-ward}
 P_x(S_{y_j}f-\hbar\partial_{y_j}f)
   =(S_{\mathrm{eff}})_{y_j}P_x(f)
                           -\hbar\partial_{y_j}P_x(f).
\end{equation}
The hidden identity deals with the last eliminated odd
generator; \eqref{eq:vol2-retained-ward} deals with every
retained odd generator, with its left-exterior sign.
Terms still containing an eliminated odd generator vanish.
This proves the cochain equation in every exterior degree.

Let $(i_0,P_x^0,H_0)$ be the Gaussian contraction in $x$
tensored with the retained free Ward complex. Apply
Lemma~\ref{lem:vol2-filtered-perturbation} with
$\delta=\sum V_i\partial_{\eta_i}$. Its projection is
$P_x^0(1+\delta H_0)^{-1}$. Each $H_0$ adds one hidden
odd generator. A term surviving the final $P_x^0$ must
start without a hidden odd generator, and each subsequent
$\delta$ must remove the hidden generator just introduced.
Thus this projection acts on retained coefficients by a
normalized functional for the hidden Ward differential.
The uniqueness recursion used in
Theorem~\ref{thm:vol2-formal-state-retract}, now with $y$
as coefficients, identifies it with $P_x$.

In the transferred differential, one more $\delta$ occurs
than $H_0$. A surviving correction therefore removes exactly
one retained odd generator and otherwise only hidden ones.
It has the form $\sum_jF_j(y)\partial_{\eta_{y_j}}$.
Evaluating the cochain equation on $\eta_{y_j}$ identifies
$F_j=\partial_{y_j}(S_{\mathrm{eff}}-\tfrac12y^{\mathrm t}A_yy)$.
The retained second-order term is unchanged. The filtered
lemma now gives the asserted contraction and quasi-isomorphism.
Its projection, inclusion and homotopy are all integral in
$\hbar$.

It remains to prove integrality of the full effective potential
and the stated powers in its cumulants. Expand Wick
expectation with independent marked sources for the observables.
The logarithm selects connected graphs. A graph with $n$
marked vertices, $v$ interaction vertices and $E$ edges has
\[
 E\geq n+v-1.
\]
Every edge contributes a factor $\hbar A_x^{-1}$ and each
interaction vertex contributes at worst $\hbar^{-1}$.
The resulting power is at least $E-v\geq n-1$.
At coupling order $g^k$, the positive coupling orders of
the vertices sum to $k$, so $v\leq k$. For fixed polynomial
inputs there are only finitely many graphs at that order.
Retained variables are coefficients in this count, including
their exterior signs. An eliminated odd input gives zero.

Vacuum components cancel from normalized cumulants. For
$\log Z_x$, a connected vacuum graph satisfies $E\geq v-1$;
the factor $-\hbar$ in the effective potential removes its
one possible negative quantum power. An interaction independent
of $x$ contributes its original potential term, so this argument
also includes zero-valent vertices. It proves both remaining
integrality assertions.

Finally apply Theorem~\ref{thm:vol2-cumulant-comparison} to
$P_x$. In its arity-$n$ morphism equation, a unary term has
the factor $\hbar^{n-1}$. A source binary term has
$\hbar\hbar^{n-2}$, while a target binary term with input
arities $r+s=n$ has
$\hbar\hbar^{r-1}\hbar^{s-1}$. Dividing these common
factors gives the equations for \eqref{eq:vol2-divided-cumulants}.
The division is valid in the integral polynomial modules.
The linear component remains $P_x$; the unit assertion
follows from the corresponding assertion for cumulants.
\end{proof}

The quotients in \eqref{eq:vol2-divided-cumulants} are constructed
before reducing modulo any power $\hbar^M$. This order gives
maps at every finite quantum precision. Division by
$\hbar^{n-1}$ would not be well defined on an arbitrary
class in the truncated ring alone.

\subsection{The determinant force}

For units $a,b\in R_0$, take
\[
 S(x,y)=\tfrac12(a+gy)x^2+\tfrac12by^2.
\]
The conditional Gaussian formulas are
\begin{equation}\label{eq:vol2-determinant-force}
 \begin{aligned}
 Z_x&=(1+gy/a)^{-1/2},\\
 S_{\mathrm{eff}}&=\tfrac12by^2
                         +\tfrac{\hbar}{2}\log(1+gy/a),\\
 P_x(x^2)&=\frac{\hbar}{a+gy}.
 \end{aligned}
\end{equation}
All inverses and the square root are their $g$-adic expansions
with the displayed constant terms. The logarithmic determinant
changes the retained force by $\hbar g/(2(a+gy))$.
It is visible in a single cochain equation: on $\eta_yx^2$,
both $P_xD_S$ and $D_{S_{\mathrm{eff}}}P_x$ have value
\begin{equation}\label{eq:vol2-determinant-force-value}
 \frac{by\hbar}{a+gy}
                  +\frac{3g\hbar^2}{2(a+gy)^2}.
\end{equation}
On the source side use $P_x(x^4)=3\hbar^2/(a+gy)^2$.
On the target side differentiate $\hbar/(a+gy)$ and include
the determinant force. Omitting that force leaves the
nonzero defect $g\hbar^2/(2(a+gy)^2)$.

\subsection{A retained odd generator also changes}

The projection is accompanied by the corrected inclusion
$i'=(1+H_0\delta)^{-1}i_0$ from
\eqref{eq:vol2-perturbation-formulas}. For the formal potential
\[
 S(x,y)=\tfrac12ax^2+\tfrac12by^2+g\,\frac{yx^3}{6},
\]
this correction is already nonzero on $\eta_y$:
\begin{equation}\label{eq:vol2-cubic-antifield-lift}
 \begin{split}
 i'(\eta_y)={}&\eta_y
  -g\eta_x\left(\frac{x^2}{6a}+\frac{\hbar}{3a^2}\right)\\
 &+g^2y\eta_x\left(\frac{x^3}{12a^2}
                          +\frac{5\hbar x}{12a^3}\right)+O(g^3).
 \end{split}
\end{equation}
Indeed $\delta\eta_y=gx^3/6$, and the Gaussian homotopy
sends $x^3$ to $\eta_x(x^2/a+2\hbar/a^2)$.
Applying $H_0\delta$ a second time gives the displayed
quadratic-coupling term.

The same model has
\begin{align}
 Z_x&=1+\frac{5g^2\hbar y^2}{24a^3}
             +\frac{385g^4\hbar^2y^4}{1152a^6}+O(g^6),
             \label{eq:vol2-cubic-partition}\\
 S_{\mathrm{eff}}
 &=\tfrac12by^2-\frac{5g^2\hbar^2y^2}{24a^3}
                       -\frac{5g^4\hbar^3y^4}{16a^6}+O(g^6).
                       \notag
\end{align}
The Gaussian sixth and twelfth moments give $Z_x$.
Subtracting half the square of its order-two term in
$\log Z_x$ gives the fourth-order effective interaction.
Thus the state, the retained force and the lift of an odd
observable change together under elimination.

\section{Sewing the state}
\label{sec:vol2-state-sewing}

Conditional elimination provides maps on observables and an
effective weight for the next integration. Both enter sewing.
The quadratic calculation separates the covariance that crosses
the cut from the determinant that normalizes the result.

\subsection{Covariance and determinant}

Order the coordinates as retained $y$ followed by eliminated
$x$, and let
\[
 A=\begin{pmatrix}B&L\\L^{\mathrm t}&D\end{pmatrix},
 \qquad S=B-LD^{-1}L^{\mathrm t}.
\]
Assume that $D$ and $S$ are invertible. The matrix $S$ is the
\emph{Schur complement} of the eliminated block. Completing
the square gives
\begin{equation}\label{eq:vol2-schur-square}
 \begin{split}
 \tfrac12(y,x)A(y,x)^{\mathrm t}
   =\tfrac12y^{\mathrm t}Sy
    +\tfrac12(x+D^{-1}L^{\mathrm t}y)^{\mathrm t}
                         D(x+D^{-1}L^{\mathrm t}y).
 \end{split}
\end{equation}
Put
\[
 E=\begin{pmatrix}1\\-D^{-1}L^{\mathrm t}\end{pmatrix},
 \qquad
 C_0=\begin{pmatrix}0&0\\0&D^{-1}\end{pmatrix}.
\]
The map $E$ extends a retained configuration to the critical
point in its eliminated fibre.

\begin{proposition}\label{prop:vol2-schur-sewing}
One has
\begin{equation}\label{eq:vol2-schur-covariance}
 A^{-1}=C_0+ES^{-1}E^{\mathrm t},\qquad
 \det A=\det D\det S.
\end{equation}
The normalized Wick functional of covariance $\hbar A^{-1}$
equals successive Wick integration in
$\xi=x+D^{-1}L^{\mathrm t}y$ and then in $y$, with
covariances $\hbar D^{-1}$ and $\hbar S^{-1}$.
Successive Schur complements agree with elimination of the
combined block whenever the intermediate complements are
invertible.
\end{proposition}
\begin{proof}
The triangular matrix
\[
 M=\begin{pmatrix}1&0\\-D^{-1}L^{\mathrm t}&1\end{pmatrix}
\]
has determinant one and satisfies
$M^{\mathrm t}AM=\operatorname{diag}(S,D)$.
Taking inverses gives the first identity in
\eqref{eq:vol2-schur-covariance}; taking determinants gives
the second. The exponential generating function
\eqref{eq:vol2-wick-generating} for two independent Gaussian
blocks then becomes the generating function of covariance
$\hbar(C_0+ES^{-1}E^{\mathrm t})$. Equality of its polynomial
coefficients proves the state identity.

To compare successive block eliminations, solve the linear
critical-point equations in the eliminated variables. Solving
them one block at a time and solving the combined system give
the same solution. Substitution in the quadratic form therefore
gives the same retained form. The determinant identity and
the generating-function identity then compose in the same
order. This proves the assertion for any finite number of
admissible blocks.
\end{proof}

The change of odd generators is also fixed by this calculation.
For the coordinates $(y,\xi)$, write the odd generators as
$\vartheta_y,\vartheta_\xi$. The old ones become
\[
 \eta_{x_a}\longmapsto\vartheta_{\xi_a},\qquad
 \eta_{y_i}\longmapsto\vartheta_{y_i}
       +\sum_a(D^{-1}L^{\mathrm t})_{ai}\vartheta_{\xi_a}.
\]
These are the inverse-Jacobian transformations of coordinate
vector fields. Substitution in the Ward differential gives
the block-diagonal differential for
\eqref{eq:vol2-schur-square}. Thus the conditional comparison
also applies to all odd observables.

At a positive real specialization, the unnormalized generating
integral is
\begin{equation}\label{eq:vol2-unnormalized-gaussian}
 \begin{split}
 Z_A(J)
 &=\int_{\mathbb R^{r+s}}
        e^{(-q^{\mathrm t}Aq/2+J^{\mathrm t}q)/\hbar}\,dq\\
 &=(2\pi\hbar)^{(r+s)/2}(\det A)^{-1/2}
                       e^{J^{\mathrm t}A^{-1}J/(2\hbar)}.
 \end{split}
\end{equation}
The two identities in \eqref{eq:vol2-schur-covariance} preserve
respectively its moments and its normalization. Over a general
coefficient ring, an unnormalized state requires a specified
square-root normalization in addition to the algebraic
determinant identity. The normalized Wick functional itself
needs no such choice.

\subsection{The intermediate weight}

Return to the block potential
\eqref{eq:vol2-conditional-potential}. Let $P_y^0$ denote the
free Wick expectation in the next block. Its effective
conditional expectation is
\[
 P_y^{\mathrm{eff}}(b)
      =\frac{P_y^0(Z_xb)}{P_y^0(Z_x)}.
\]
Thus the factor $Z_x$, which may depend on $y$, is part of
the next integration.

\begin{theorem}\label{thm:vol2-integral-sewing}
Successive normalized elimination equals joint elimination:
\begin{equation}\label{eq:vol2-normalized-fubini}
 P_y^{\mathrm{eff}}P_x(f)
  =\frac{P_y^0P_x^0(fe^{-V/\hbar})}
         {P_y^0P_x^0(e^{-V/\hbar})}.
\end{equation}
For these maps, and for further admissible blocks, the divided
cumulants of \eqref{eq:vol2-divided-cumulants} compose by
\begin{equation}\label{eq:vol2-divided-sewing}
 \Phi_n^{QP}(b_1,\ldots,b_n)
   =\sum_{\pi\in\Pi_n}\epsilon(\pi;b)
      \Phi_{|\pi|}^Q
       \left(\Phi_{|E|}^P(b_j:j\in E):E\in\pi\right).
\end{equation}
These equalities hold over the integral quantum ring, and
iterated sewing is associative in every arity.
\end{theorem}
\begin{proof}
Insert the normalized numerator of $P_x$ into the next
expectation. Its denominator $Z_x$ cancels the intermediate
weight, giving \eqref{eq:vol2-normalized-fubini}.
Free Wick operations on separate blocks commute. At each
coupling order only finitely many polynomial contractions
occur, so this cancellation and Gaussian Fubini are defined
coefficientwise. Theorem~\ref{thm:vol2-integral-conditional}
shows that the resulting normalized maps are integral.

Equation~\eqref{eq:vol2-cumulant-sewing} gives the partition
composition law before division. In the term indexed by
$\pi$, the quantum exponents add to
\[
 (|\pi|-1)+\sum_{E\in\pi}(|E|-1)=n-1.
\]
Dividing that common power gives
\eqref{eq:vol2-divided-sewing}. Composition of the underlying
coalgebra maps is associative; the divided formula preserves
that equality. This argument is also the composition mechanism
for derived BV morphisms in Bandiera
\cite{Bandiera2020}*{\S2.2, Proposition 2.8}.
\end{proof}

For example, let $X=Y+\xi$, where the first expectation
eliminates a centred Gaussian $\xi$ of variance $\sigma$
and leaves $Y$. Its generating function on the square is
\[
 P(e^{tX^2})=(1-2\sigma t)^{-1/2}
                    \exp\!\left(\frac{tY^2}{1-2\sigma t}\right).
\]
The logarithm gives, in every arity,
\begin{equation}\label{eq:vol2-conditional-square-cumulants}
 \kappa_n^P(X^2,\ldots,X^2)
  =2^{n-1}(n-1)!
                    \bigl(\sigma^n+nY^2\sigma^{n-1}\bigr).
\end{equation}
If the next expectation eliminates an independent centred
Gaussian $Y$ of variance $\tau$, the joint variable $X$ has
variance $\sigma+\tau$. Equation~\eqref{eq:vol2-cumulant-sewing}
therefore combines all the conditional terms into
$2^{n-1}(n-1)!(\sigma+\tau)^n$.
Keeping only the conditional mean of $X^2$ would lose these
higher terms.

\section{Coordinate descent of the whole Ward complex}
\label{sec:vol2-ward-coordinates}

A change of coordinates transforms the integration density
and the odd observables at the same time. The logarithm of
the Jacobian enters the potential, while the inverse Jacobian
transforms the odd generators. Their two contributions are
required by the cochain equation.

Let $S(x)=\tfrac12x^{\mathrm t}Ax+V(x)$ be a finite potential
as above. Take a formal change of variables
$x=\phi(q)=q+O(g)$, polynomial in $q$ at each coupling order.
Its Jacobian $J=D_q\phi$ is invertible $g$-adically. Define
\begin{equation}\label{eq:vol2-coordinate-pullback}
 \Phi_\phi(f)=f\circ\phi,\qquad
 \Phi_\phi(\eta_i)=\sum_a(J^{-1})_{ai}\vartheta_a,
 \qquad
 S_\phi=S\circ\phi-\hbar\log\det J.
\end{equation}
Here $\vartheta_a$ are the new degree-minus-one generators.
The logarithm is the unique formal logarithm with zero
constant coupling term, since $\det J=1+O(g)$.

\begin{theorem}\label{thm:vol2-ward-coordinate-descent}
The map $\Phi_\phi$ is a unital graded-algebra and cochain
isomorphism over $R_0[\hbar][[g]]$, and it preserves the
normalized formal state:
\begin{equation}\label{eq:vol2-coordinate-state}
 D_{S_\phi}\Phi_\phi=\Phi_\phi D_S,
 \qquad p_{S_\phi}\Phi_\phi=p_S.
\end{equation}
For a second formal change $q=\psi(r)$,
\[
 (S_\phi)_\psi=S_{\phi\circ\psi},\qquad
 \Phi_{\phi\circ\psi}=\Phi_\psi\Phi_\phi.
\]
\end{theorem}
\begin{proof}
Regard the odd generators as coordinate vector fields and
their products as polyvectors. For the weighted volume form
$\rho=e^{-S/\hbar}dx_1\cdots dx_m$, divergence on polyvectors
has the coordinate formula
\[
 \operatorname{Div}_\rho
    =\sum_i\left(\partial_{x_i}-\hbar^{-1}S_i\right)
                                                    \partial_{\eta_i}.
\]
Thus $D_S=-\hbar\operatorname{Div}_\rho$.
This notation abbreviates the polynomial Ward operator; the
displayed product has no negative quantum power.
Contraction with a volume form identifies this divergence
with the exterior differential, with the exterior signs fixed
by left odd differentiation. It is therefore natural under
coordinate changes. In the present coordinates,
\[
 \phi^*\rho=e^{-S(\phi(q))/\hbar}\det J\,dq
                 =e^{-S_\phi(q)/\hbar}dq,
\]
and vector fields pull back by $J^{-1}$. These two facts
give the first equation in \eqref{eq:vol2-coordinate-state}
in every exterior degree.

For the state, interpolate by
$\phi_u(q)=q+u(\phi(q)-q)$, $0\leq u\leq1$.
The derivative of its pulled-back density times any pulled-back
observable is a total derivative in the integration variables.
At each coupling order, it is a polynomial times the free
Gaussian weight. Formal Gaussian integration by parts
annihilates that derivative. Both the unnormalized numerator
and denominator are consequently constant in $u$.
Their quotient proves the state equation on even observables;
negative cohomological degrees vanish on both sides.
The argument is formal in $g$ and uses no global real
invertibility of $\phi$.

Finally the chain rule gives
\[
 \det D(\phi\circ\psi)
       =((\det D\phi)\circ\psi)\det D\psi.
\]
The factors are commuting units with constant term one,
so their formal logarithms add. The inverse-Jacobian rule
has the same composition law on vector fields. This proves
the two cocycle identities.
\end{proof}

For $S=ax^2/2$ and $\phi(q)=q+gq^2$,
\begin{equation}\label{eq:vol2-nonlinear-jacobian}
 S_\phi=\tfrac a2(q+gq^2)^2-\hbar\log(1+2gq),
 \qquad
 \Phi_\phi(\eta)=\frac{\vartheta}{1+2gq}.
\end{equation}
If the logarithmic term is omitted, the cochain defect on
this odd observable is $2g\hbar/(1+2gq)^2$.
The inverse coordinate begins
$q=x-gx^2+2g^2x^3-5g^3x^4+\cdots$; its fifth-power term
is $14g^4x^5$, which has zero Gaussian expectation.
Consequently
\[
 p_{S_\phi}(q)=-\frac{g\hbar}{a}
                 -\frac{15g^3\hbar^2}{a^2}+O(g^5),
 \qquad p_{S_\phi}(\phi(q))=0.
\]
Without the Jacobian in the density, the last expectation
instead begins with $-2g\hbar/a$. The same missing term
is detected by a cochain and by an ordinary observable.

\subsection{A flat family of state operations}

Let $\phi_s$ depend on finitely many parameters $s_i$ with
commuting derivatives $\partial_i$. Its coefficients may be
formal, or smooth on a specified parameter domain; all
identities remain coefficientwise in $g$. Put
\[
 S_s=S_{\phi_s},\qquad
 w_i=(D_q\phi_s)^{-1}\partial_i\phi_s,
 \qquad \Psi_i=\sum_a w_i^a\vartheta_a.
\]
For a vector field $w$, its Lie derivative on polyvectors
acts by $\mathcal L_w f=w(f)$ on functions and
$\mathcal L_w\vartheta_a=-\sum_b(\partial_{q_a}w^b)\vartheta_b$
on coordinate vector fields. Extend it as a derivation.
Define the connection
\[
 \nabla_i=\partial_i-\mathcal L_{w_i}.
\]

\begin{proposition}\label{prop:vol2-flat-state-family}
The coordinate family satisfies
\begin{equation}\label{eq:vol2-flat-state-family}
 \begin{gathered}
 \partial_iS_s=D_{S_s}\Psi_i
                  =w_i(S_s)-\hbar\operatorname{div}w_i,\\
 [\nabla_i,D_{S_s}]=0,
 \qquad [\nabla_i,\nabla_j]=0,
 \qquad \partial_i p_s(F)=p_s(\nabla_iF).
 \end{gathered}
\end{equation}
Every cumulant is horizontal:
\begin{equation}\label{eq:vol2-horizontal-cumulants}
 \partial_i\kappa_n^{p_s}(F_1,\ldots,F_n)
  =\sum_j\kappa_n^{p_s}
       (F_1,\ldots,\nabla_iF_j,\ldots,F_n).
\end{equation}
The same statement holds for the integral divided cumulants.
\end{proposition}
\begin{proof}
Differentiation of pullback gives
$\partial_i\Phi_s=\mathcal L_{w_i}\Phi_s$.
Differentiating the cochain equation in
Theorem~\ref{thm:vol2-ward-coordinate-descent} proves
$[\nabla_i,D_{S_s}]=0$.
Commuting two parameter derivatives of $\Phi_s$ gives
\[
 \partial_iw_j-\partial_jw_i-[w_i,w_j]=0,
\]
which is the curvature equation for $\nabla$.
Differentiating the pulled-back density gives the first
identity in \eqref{eq:vol2-flat-state-family}.

On a scalar polynomial $F$, the Ward equation gives
\[
 p_s(FD_{S_s}\Psi_i)=\hbar p_s(w_iF),
 \qquad p_s(D_{S_s}\Psi_i)=0.
\]
Indeed $D_{S_s}(F\Psi_i)=F D_{S_s}\Psi_i-\hbar w_iF$,
and its expectation is zero. The derivative of the
normalized state is the covariance with $\partial_iS_s$
divided by $-\hbar$, together with $p_s(\partial_iF)$.
Substituting these Ward identities gives
$p_s(\partial_iF-w_iF)$. The expression is integral;
both sides vanish on negative cohomological degrees.

Apply this state equation to every block of
\eqref{eq:vol2-state-cumulant}. Since $\mathcal L_{w_i}$
is a derivation, each input is differentiated once with
the same partition sign. This proves
\eqref{eq:vol2-horizontal-cumulants}. Dividing the already
integral cumulants by their known quantum powers gives
the final assertion.
\end{proof}

For example,
\[
 \phi_{u,v}(q)=q+ugq^2+vgq^3,
 \qquad
 w_u=\frac{gq^2}{1+2ugq+3vgq^2},\quad
 w_v=\frac{gq^3}{1+2ugq+3vgq^2}
\]
satisfies the displayed curvature equation as a rational
identity, before its $g$-adic expansion. On parameter forms
of degree $p$, the total differential is
$d_\nabla+(-1)^pD_{S_s}$. Flatness and the commutator
equation make its square zero. Restriction to parameter
faces preserves the cochain and cumulant equations.

\subsection{Coordinate changes within an eliminated fibre}

For a conditional reduction, take
$(x,y)=(\phi(q,y),y)$ and write
$J_x=D_q\phi$, $B=D_y\phi$.
The full inverse Jacobian transforms the old odd generators as
\begin{equation}\label{eq:vol2-fibre-antifields}
 \begin{aligned}
 \eta_{x_i}&\longmapsto
                     \sum_a(J_x^{-1})_{ai}\vartheta_{q_a},\\
 \eta_{y_j}&\longmapsto\vartheta_{y_j}
                     -\sum_a(J_x^{-1}B)_{aj}\vartheta_{q_a}.
 \end{aligned}
\end{equation}
The density changes by $\det J_x$.

\begin{proposition}\label{prop:vol2-conditional-coordinate-square}
The fibre substitution preserves the conditional partition
factor $Z_x(y)$ and effective potential. If $P_q$ is the
transformed conditional Ward map, then
\[
 P_q\Phi_\phi=P_x
\]
on the entire complex. The same square commutes on every
cumulant and divided cumulant, and with subsequent eliminations.
\end{proposition}
\begin{proof}
Apply the formal change-of-variables argument of
Theorem~\ref{thm:vol2-ward-coordinate-descent}, with $y$
held as coefficients. It preserves the conditional numerator
and its normalization, hence also $S_{\mathrm{eff}}$.
Projection kills every term containing a $\vartheta_q$.
In the second line of \eqref{eq:vol2-fibre-antifields}, only
$\vartheta_{y_j}$ survives; the same holds term by term in
products of retained odd generators. This proves the cochain
square in every exterior degree. Substitution in the
partition formula proves its cumulant versions, and
Theorem~\ref{thm:vol2-integral-sewing} proves compatibility
with successive blocks.
\end{proof}

\section{Canonical families and their marked homotopies}
\label{sec:vol2-gauge-state-homotopies}

Coordinate descent transports a state by an algebra isomorphism.
A canonical variation of a BV weight gives a more general
comparison: the normalized elimination commutes with parallel
transport up to a specified cochain homotopy. Its higher
components have one additional marked insertion.

\subsection{The integration domain and the master weight}

Let $A,B$ be graded-commutative coefficient algebras over a
characteristic-zero field with central parameter $\hbar$.
Initially allow Laurent coefficients in $\hbar$; integral
restrictions will be specified for the finite Gaussian sector.
Let $\Delta_A,\Delta_B$ have degree one, square zero,
differential-operator order at most two, and value zero on
the unit. In the potential model our convention is
\[
 \Delta_A=-\sum_i\partial_{q_i}\partial_{\eta_i}.
\]
Choose a degree-zero linear integration map $\mathcal P$
on a specified domain of weights in extensions of $A,B$, satisfying
the Stokes equation
\begin{equation}\label{eq:vol2-bv-stokes}
 \Delta_B\mathcal P=\mathcal P\Delta_A.
\end{equation}
Multiplication by the weights and division by the effective
weights below are required to be defined on this domain.

In the finite potential model, the domain consists of Gaussian
weights times polynomial observables, expanded coefficientwise
in the positive coupling filtration. Integration means formal
Wick integration in the eliminated variables. Gaussian
integration by parts proves \eqref{eq:vol2-bv-stokes} there.
For an analytic integral the integration cycle, decay and
absence of boundary terms are hypotheses of the same equation.

A weight $U_A$ of cohomological degree zero satisfying
$\Delta_AU_A=0$ is a \emph{master weight}. Put
$U_B=\mathcal P U_A$ and assume that $U_A$ and $U_B$
are invertible on the chosen domain. Define
\begin{equation}\label{eq:vol2-normalized-bv-map}
 D_A=\hbar U_A^{-1}\Delta_AU_A,
 \qquad D_B=\hbar U_B^{-1}\Delta_BU_B,
 \qquad p(f)=U_B^{-1}\mathcal P(U_Af).
\end{equation}
Conjugation in the two differential formulas denotes
operators: for example $D_Af=\hbar U_A^{-1}\Delta_A(U_Af)$.
For $U_A=e^{-S/\hbar}$ with $S$ independent of odd variables,
this is exactly the Ward differential
\eqref{eq:vol2-potential-ward}.

\begin{proposition}\label{prop:vol2-normalized-bv-map}
The effective weight is a master weight, and $p$ is a unital
cochain map. Successive normalized integrations compose when
each integration uses the effective weight from the preceding
one.
\end{proposition}
\begin{proof}
Stokes gives $\Delta_BU_B=0$. Conjugation therefore gives
$D_A1=D_B1=0$ and $D_A^2=D_B^2=0$. Moreover
\[
 D_Bp
   =\hbar U_B^{-1}\Delta_B\mathcal P U_A
   =\hbar U_B^{-1}\mathcal P\Delta_AU_A=pD_A,
\]
and $p1=1$ by definition. If $U_1=\mathcal P_1U_A$ and
$U_2=\mathcal P_2U_1$, the composite normalized map is
\[
 U_2^{-1}\mathcal P_2
        \bigl(U_1U_1^{-1}\mathcal P_1(U_Af)\bigr)
       =U_2^{-1}\mathcal P_2\mathcal P_1(U_Af).
\]
This is normalization of the composite integration, with
the entire intermediate weight retained.
\end{proof}

Theorem~\ref{thm:vol2-integral-conditional} supplies the
quasi-isomorphism and integral formulas in the finite
potential sector. Proposition~\ref{prop:vol2-normalized-bv-map}
uses only Stokes and the stated weight domain; those
hypotheses alone assert a cochain map.

\subsection{Transporting a canonical variation}

Let $t$ be a parameter and $R_t\in A$ have degree $-1$.
Suppose
\begin{equation}\label{eq:vol2-canonical-master-family}
 \partial_tU_A=\Delta_A(R_tU_A),\qquad \Delta_AU_A=0.
\end{equation}
The integration map and the two BV operators remain fixed.
A varying cycle can be treated in this form only after
its pullback to a fixed domain has been specified.
Write $m_R$ for multiplication by $R$ and put
\[
 L_R=[\Delta_A,m_R]=\Delta_Am_R+m_R\Delta_A,
 \qquad r_t=p_t(R_t).
\]
By Stokes, the effective family obeys
\begin{equation}\label{eq:vol2-effective-master-family}
 \partial_tU_B=\Delta_B(r_tU_B).
\end{equation}

Set $w_A=U_A^{-1}\partial_tU_A$ and
$w_B=U_B^{-1}\partial_tU_B$. Define the degree-zero transports
\begin{equation}\label{eq:vol2-normalized-transports}
 v_A=U_A^{-1}L_RU_A-m_{w_A},\qquad
 v_B=U_B^{-1}[\Delta_B,m_r]U_B-m_{w_B}.
\end{equation}
A commutator of a second-order operator with multiplication
has order at most one. Its value on the unit here is $w$.
Subtracting that value therefore makes each $v$ a derivation.
Since $[\Delta,L_R]=0$, differentiation gives
\begin{equation}\label{eq:vol2-moving-differentials}
 \partial_tD_A=[v_A,D_A],\qquad
 \partial_tD_B=[v_B,D_B].
\end{equation}
Thus the connections $\partial_t-v_A$ and $\partial_t-v_B$
commute with their respective differentials.

\begin{theorem}\label{thm:vol2-canonical-state-homotopy}
The degree-minus-one map
\begin{equation}\label{eq:vol2-marked-linear-homotopy}
 h_t(f)=\frac{p_t(R_tf)-p_t(R_t)p_t(f)}{\hbar}
\end{equation}
satisfies
\begin{equation}\label{eq:vol2-state-gauge-equation}
 \partial_tp_t+p_tv_A-v_Bp_t=D_Bh_t+h_tD_A,
 \qquad h_t(1)=0.
\end{equation}
The formula holds on the declared integration domain.
For polynomial marked inputs in a formal Gaussian
presentation with no negative quantum powers, the quotient
in \eqref{eq:vol2-marked-linear-homotopy} is integral.
\end{theorem}
\begin{proof}
On unnormalized weights define the odd operator
$K=\mathcal Pm_R-m_r\mathcal P$.
Equation~\eqref{eq:vol2-bv-stokes} gives
\begin{equation}\label{eq:vol2-unnormalized-gauge-equation}
 \mathcal PL_R-L_r\mathcal P
                      =\Delta_BK+K\Delta_A.
\end{equation}
Differentiate $p_t=U_B^{-1}\mathcal P U_A$ and add
$p_tv_A-v_Bp_t$. The two logarithmic-derivative
multiplication terms cancel, leaving
\[
 U_B^{-1}(\mathcal PL_R-L_r\mathcal P)U_A.
\]
Conjugating \eqref{eq:vol2-unnormalized-gauge-equation}
by the weights and using
$h_t=\hbar^{-1}U_B^{-1}KU_A$ proves
\eqref{eq:vol2-state-gauge-equation}.
Normalization gives $K(U_A)=0$, hence $h_t1=0$.
In the formal Gaussian sector, the numerator is the
two-input cumulant $\kappa_2^{p_t}(R_t,f)$.
Its connected graphs have at least one quantum power by
the argument of Theorem~\ref{thm:vol2-integral-conditional}.
\end{proof}

When parallel transports $T_A(t,0),T_B(t,0)$ for these
connections exist in the specified parameter domain,
\[
 \widehat p_t=T_B(t,0)^{-1}p_tT_A(t,0)
\]
has derivative homotopic to zero. Its homotopy is
$T_B(t,0)^{-1}h_tT_A(t,0)$.
Integrating this expression along a defined interval gives
an endpoint cochain homotopy. This may be an ordinary
analytic interval or coefficientwise integration in a
positive formal filtration; the evolution and endpoint
evaluation are part of the chosen domain.

\subsection{One marked input in every arity}

Apply the coalgebra construction of
\eqref{eq:vol2-cumulant-coalgebra} to $p_t$.
Let $S^c(p_t;h_t)$ be its unary tangent map: it applies
$h_t$ to exactly one symmetric factor and $p_t$ to every
other factor, with the graded tensor signs. It is a
degree-minus-one coderivation along $S^c(p_t)$.
Define
\[
 \mathcal H_t=\tau_B^{-1}S^c(p_t;h_t)\tau_A.
\]
For a derivation $v$, write $\widetilde v$ for its unary
coderivation. Derivations commute with multiplication
in the sense $\widetilde v\tau=\tau\widetilde v$.

\begin{theorem}\label{thm:vol2-marked-cumulant-homotopy}
The full operation comparison satisfies
\begin{equation}\label{eq:vol2-coalgebra-gauge-equation}
 \partial_t\mathscr K(p_t)
  +\mathscr K(p_t)\widetilde v_A
  -\widetilde v_B\mathscr K(p_t)
     =\mathscr D_B\mathcal H_t+\mathcal H_t\mathscr D_A.
\end{equation}
Its homotopy has Taylor coefficients
\begin{equation}\label{eq:vol2-marked-cumulant-homotopy}
 (\mathcal H_t)_n(b_1,\ldots,b_n)
       =\hbar^{-1}\kappa_{n+1}^{p_t}(R_t,b_1,\ldots,b_n).
\end{equation}
The marked input $R_t$ is placed first; the partition
sign convention in \eqref{eq:vol2-state-cumulant}
determines all signs. In the formal Gaussian sector,
$(\mathcal H_t)_n$ is divisible by $\hbar^{n-1}$.
Dividing gives a homotopy for the divided state operations
without inverting $\hbar$.
\end{theorem}
\begin{proof}
Differentiation of $S^c(p_t)$ differentiates one factor
at a time. Equation~\eqref{eq:vol2-state-gauge-equation}
and the cochain equation for $p_t$ therefore give the
unary-coderivation version of
\eqref{eq:vol2-coalgebra-gauge-equation}.
Conjugate it by $\tau_A,\tau_B$, and use the derivation
identity for $v_A,v_B$. This proves the full equation.

To compute its Taylor coefficient, differentiate the
partition formula for $\kappa_n$ by replacing one factor
$p(\prod_{j\in E}b_j)$ by
\[
 h\left(\prod_{j\in E}b_j\right)
 =\hbar^{-1}\left[
       p\left(R\prod_{j\in E}b_j\right)
           -p(R)p\left(\prod_{j\in E}b_j\right)\right].
\]
The first terms enumerate partitions in which the block
of $R$ is not a singleton. For a partition of the original
inputs into $k$ blocks, the second term occurs $k$ times.
Its coefficient is
$-k(-1)^{k-1}(k-1)!=(-1)^k k!$, precisely the coefficient
of the added singleton block. Moving $R$ to its block
produces the Koszul sign prescribed by placing it first.
This proves \eqref{eq:vol2-marked-cumulant-homotopy}.

A connected graph with $n+1$ marked inputs, $v$ interaction
vertices and $E$ edges has $E\geq n+v$.
Its quantum power after the interaction factors is at
least $n$. Division by the explicit $\hbar$ in
\eqref{eq:vol2-marked-cumulant-homotopy} leaves
$\hbar^{n-1}$. The same arity count as in
Theorem~\ref{thm:vol2-integral-conditional} permits division
of the entire homotopy equation.
\end{proof}

These homotopies also sew. For successive eliminations
$p_t:A\to B$ and $q_t:B\to C$ with the same intermediate
transport, differentiation gives
\begin{equation}\label{eq:vol2-gauge-homotopy-sewing}
 h_t^{\mathrm{tot}}=q_th_t^p+h_t^qp_t,
 \qquad
 \mathcal H^{\mathrm{tot}}
    =\mathscr K(q)\mathcal H^p
                         +\mathcal H^q\mathscr K(p).
\end{equation}
For normalized integration of the same master family,
the intermediate generator is $r_t=p_t(R_t)$.
The two covariance expressions in the first formula then
cancel their intermediate $q_t(r_tp_t(f))$ term, leaving
the marked covariance for $q_tp_t$ itself.
Three eliminations give the same three homotopy terms
in either parenthesization. Thus the path comparison
and the state sewing have a common composition law.

\subsection{The total master equation on parameter faces}

Let $T$ be a finite-dimensional parameter simplex.
Place parameter forms to the left of coefficient
elements, and extend $\Delta_A$ by
\[
 \Delta_A(\alpha\otimes a)
            =(-1)^{|\alpha|}\alpha\otimes\Delta_Aa.
\]
Suppose a total-degree-zero weight has been supplied with
\begin{equation}\label{eq:vol2-total-master-weight}
 \mathcal U_A\in\Omega^\bullet(T)\widehat\otimes A,
 \qquad(d_T+\hbar\Delta_A)\mathcal U_A=0,
\end{equation}
and invertible zero-form weight on the integration domain.
Its pushforward $\mathcal U_B=\mathcal P\mathcal U_A$
satisfies the same total equation by Stokes. The full
weights are invertible wherever their zero-form parts
are, because positive form degree is nilpotent.

\begin{proposition}\label{prop:vol2-master-family-faces}
For a supplied family \eqref{eq:vol2-total-master-weight},
the operators
\[
 \mathbb D_A=\mathcal U_A^{-1}(d_T+\hbar\Delta_A)\mathcal U_A,
 \qquad
 \mathbb D_B=\mathcal U_B^{-1}(d_T+\hbar\Delta_B)\mathcal U_B
\]
and the normalized map
\[
 \mathbb p(F)=\mathcal U_B^{-1}\mathcal P(\mathcal U_AF)
\]
are compatible with restriction to every parameter face.
They give square-zero unital differentials and a cochain
map. Their cumulants supply the higher parameter
homotopies and their boundary equations.
\end{proposition}
\begin{proof}
The form differential and the extended BV operator
anticommute. Hence their sum squares to zero, and
conjugation preserves that identity. The total master
equation makes the conjugated differentials kill the unit.
Stokes gives $\mathbb D_B\mathbb p=\mathbb p\mathbb D_A$
by the calculation of
Proposition~\ref{prop:vol2-normalized-bv-map}.
All these operations commute with pullback to a face.
Apply the finite-length cumulant conjugation to this
cochain equation over the ground coefficient ring.
Its coefficients in parameter form degree are the
successive homotopy equations; integration of their
exterior derivatives gives their face boundaries by
Stokes on the parameter simplex.
\end{proof}

On an interval, \eqref{eq:vol2-canonical-master-family}
gives the actual total weight
\[
 \mathcal U_A=U_A+dt\,R_tU_A/\hbar.
\]
The coefficient of $dt$ in its normalized map is $h_t$,
and its total differential is
$D_A+dt(\partial_t-v_A)$ on zero-form inputs.
For two or more parameters, independently chosen
one-parameter generators need not extend to a total
master weight. If $\mathcal U_1$ is its proposed
one-form part, a two-form part must solve the degree-two
equation of \eqref{eq:vol2-total-master-weight}.
The existence of that part is an additional mathematical
condition. The flat coordinate families of
Proposition~\ref{prop:vol2-flat-state-family} give one
explicit class with all parameter compatibilities.

\subsection{A canonical shear with nonzero homotopy}

Take even variables $x,y$ and odd variables $\eta,\zeta$
of degree $-1$, with
\[
 \Delta=-\partial_x\partial_\eta-\partial_y\partial_\zeta.
\]
For invertible stiffnesses $a,b$, set
\begin{equation}\label{eq:vol2-shear-potential}
 S_t(x,y)=\tfrac12ax^2+\tfrac12b(y-tx)^2,
 \qquad d=a+bt^2,
 \qquad k=bt/d.
\end{equation}
Work formally near $t=0$, or with real $a,b>0$ and real
$t$. The odd generator $R=\zeta x$ satisfies
\[
 [\Delta,m_R]=-x\partial_y+\zeta\partial_\eta.
\]
Since $\partial_tS_t=-bx(y-tx)$, the weight
$U_A=e^{-S_t/\hbar}$ obeys
\eqref{eq:vol2-canonical-master-family}.

Complete the square in $x$, and normalize the eliminated
Gaussian relative to its value at $t=0$. This gives
\begin{equation}\label{eq:vol2-shear-state}
 \begin{aligned}
 p_t(x)&=ky,&
 \operatorname{Var}_{p_t}(x)&=\hbar/d,\\
 S_t^{\mathrm{eff}}(y)
   &=\frac{ab}{2d}y^2+\frac{\hbar}{2}\log(d/a),&
 r_t&=\zeta ky.
 \end{aligned}
\end{equation}
The effective differential and normalized transports are
\begin{equation}\label{eq:vol2-shear-transports}
 \begin{gathered}
 D_B=\left(\frac{ab}{d}y-\hbar\partial_y\right)\partial_\zeta,\\
 v_A=-x\partial_y+\zeta\partial_\eta,
 \qquad v_B=-ky\partial_y+k\zeta\partial_\zeta.
 \end{gathered}
\end{equation}
The determinant term is independent of $y$, so it does
not appear in $D_B$. It is nevertheless required by the
effective master-family equation
\eqref{eq:vol2-effective-master-family}: differentiating
the logarithmic determinant supplies its scalar term.

The map $p_t$ kills $\eta$ and retains $\zeta$.
For an even polynomial $f$, Gaussian integration by parts
gives
\begin{equation}\label{eq:vol2-shear-homotopy}
 h_t(f)=\frac{\zeta}{\hbar}
              \bigl(p_t(xf)-ky\,p_t(f)\bigr)
                  =\frac{\zeta}{d}p_t(\partial_xf).
\end{equation}
The homotopy is zero on terms containing either odd
generator. In particular $h_t(x)=\zeta/d$ is nonzero
and integral in $\hbar$.

On $x$, both sides of
\eqref{eq:vol2-state-gauge-equation} are
\[
 \partial_t(ky)-v_B(ky)=\frac{ab}{d^2}y
                       =D_B(\zeta/d).
\]
On the eliminated odd generator $\eta$, the left side
is $p_t(v_A\eta)=\zeta$, and the right side is
\[
 h_t(D_A\eta)=h_t(dx-bty)=\zeta.
\]
This equation sees a component that cannot be tested on
ordinary even observables alone.

The higher homotopy also acts on composites. If
$m=ky$ and $\sigma=\hbar/d$, then
\[
 \kappa_2(x^2,x^2)=4m^2\sigma+2\sigma^2,
 \qquad
 \kappa_3(x,x^2,x^2)=\sigma\partial_m\kappa_2(x^2,x^2)
                              =8m\sigma^2.
\]
The derivative identity follows by differentiating
the Gaussian generating function with respect to its mean.
Equation~\eqref{eq:vol2-marked-cumulant-homotopy} therefore
gives
\begin{equation}\label{eq:vol2-shear-higher-component}
 (\mathcal H_t)_2(x^2,x^2)
                         =\frac{8\zeta\,ky\,\hbar}{d^2}.
\end{equation}
The linear homotopy and this higher component are parts
of the same normalized comparison.

\section{Comparing normalized Gaussian choices}
\label{sec:vol2-gaussian-choices}

A fixed projection onto retained variables does not determine
Gaussian elimination. For linear observables $x,y$, the
covariance remains in
\[
 R(xy)=p(x)p(y)+\hbar P(x,y).
\]
Two choices with the same projection therefore require a
comparison on products. In a finite paired complex, a
contraction of the discarded coefficients constructs that
comparison and all its relative parameter extensions.

\subsection{A contractible complex of choices}

Let $(V,d)$ and $(H,d_H)$ be bounded complexes of finitely
generated projective $R_0$-modules. Fix a contraction
\[
 H\mathrel{\mathop{\rightleftarrows}^{i}_{p}}V,
 \qquad pi=1,\qquad dh+hd=1-ip,
 \qquad hi=ph=h^2=0.
\]
The maps $i,p$ are cochain maps and $h$ has degree $-1$.
Put $t=ip$ and $\mathsf E=t\otimes t$ on $V\otimes V$.
Every tensor of maps uses the graded rule
$(f\otimes g)(x\otimes y)=(-1)^{|g||x|}f(x)\otimes g(y)$.
For a homogeneous endomorphism $M$ of $V\otimes V$, define
\[
 \delta M=d_{V\otimes V}M-(-1)^{|M|}Md_{V\otimes V}.
\]

The symmetrized tensor contraction is
\begin{equation}\label{eq:vol2-reference-two-coefficients}
 A_0=\frac12\bigl(h\otimes1+t\otimes h
                         +1\otimes h+h\otimes t\bigr).
\end{equation}
The two unsymmetrized sums are exchanged by the graded
interchange of the tensor factors. Each has differential
$1-\mathsf E$, so
\begin{equation}\label{eq:vol2-normalized-two-coefficients}
 \delta A_0=1-\mathsf E,
 \qquad \mathsf EA_0=A_0\mathsf E=0.
\end{equation}
Let $\mathcal B$ be the graded endomorphisms commuting
with that interchange, with differential $\delta$, and set
\[
 \mathcal N=(1-\mathsf E)\mathcal B(1-\mathsf E).
\]
A \emph{normalized two-coefficient homotopy} is an element
$A\in\mathcal N^{-1}$ with $\delta A=1-\mathsf E$.
This condition imposes no equation on $A^2$.

\begin{lemma}\label{lem:vol2-gaussian-choice-contraction}
Left multiplication $S(M)=A_0M$ contracts
$(\mathcal N,\delta)$:
\[
 \delta S+S\delta=1.
\]
In particular, two normalized homotopies $A,B$ have the
specified comparison $L=A_0(B-A)$ with $\delta L=B-A$.
\end{lemma}
\begin{proof}
The graded Leibniz rule gives
\[
 \delta(A_0M)=(1-\mathsf E)M-A_0\delta M
                                      =M-A_0\delta M.
\]
Normalization is preserved: both
$\mathsf EA_0M$ and $A_0M\mathsf E$ vanish.
Composition also preserves commutation with the graded
interchange. Thus $S$ takes values in $\mathcal N$ and
the identity holds there. Apply it to the closed element
$B-A$ to obtain the last assertion.
\end{proof}

\subsection{The covariance homotopy}

Suppose a graded symmetric pairing
$b:V\otimes V\to R_0$ of degree $1$ is closed in the
Hom complex. Write $b_H=b(i\otimes i)$ for its retained
pairing. For any graded symmetric pairing $P$, let
$\partial_P$ denote unordered pair contraction on
$\operatorname{Sym}(V)$, normalized by
$\partial_P(xy)=P(x,y)$. On a product of homogeneous
linear observables, select two factors and move them
to the front with the Koszul sign. Apply $P$ to that pair
and retain the other factors in their order, then sum
over all selected pairs.

Extend $d$ to $\operatorname{Sym}(V)$ as a derivation.
The two paired differentials are
\[
 D_b=d+\hbar\partial_b,
 \qquad D_H=d_H+\hbar\partial_{b_H}.
\]
Closedness of $b$ gives $[d,\partial_b]=0$; constant
coefficient pair contractions supercommute, so
$\partial_b^2=0$. Thus both displayed differentials
square to zero. For a normalized homotopy $A$, put
\begin{equation}\label{eq:vol2-choice-gaussian-map}
 P_A=-bA,
 \qquad R_A=\operatorname{Sym}(p)e^{\hbar\partial_{P_A}}.
\end{equation}
The covariance has degree zero. Its exponential is a
finite sum on every polynomial: each contraction lowers
the number of linear factors by two. These maps are
defined over $R_0[\hbar]$ and extend by scalar extension to
$R_0[[\hbar]]$ without completing polynomial degree.

\begin{proposition}\label{prop:vol2-choice-gaussian-homotopy}
Each $R_A$ is a cochain map from $D_b$ to $D_H$.
For two normalized homotopies $A,B$, define
\[
 L=A_0(B-A),\qquad Q=bL,\qquad
 P_u=(1-u)P_A+uP_B.
\]
Then the degree-minus-one map
\begin{equation}\label{eq:vol2-choice-gaussian-homotopy}
 K_{B,A}=\hbar\operatorname{Sym}(p)\partial_Q
                         \int_0^1e^{\hbar\partial_{P_u}}\,du
\end{equation}
satisfies
\[
 D_HK_{B,A}+K_{B,A}D_b=R_B-R_A.
\]
It vanishes on constants and on polynomials whose
coefficients all lie in $i(H)$.
\end{proposition}
\begin{proof}
The Hom differential obeys the graded composition rule.
Since $|b|=1$ and $\delta b=0$,
\[
 \delta P_A=b(1-\mathsf E)=b-p^*b_H.
\]
Here $p^*b_H=b_H(p\otimes p)$.
The identity $[d,\partial_{P_A}]=\partial_{\delta P_A}$
then gives
\[
 e^{\hbar\partial_{P_A}}D_b
       =(d+\hbar\partial_{p^*b_H})
                                      e^{\hbar\partial_{P_A}}.
\]
There are no further commutator terms, because the
constant coefficient contractions commute with the
degree-zero operator $\partial_{P_A}$.
Application of $\operatorname{Sym}(p)$ proves the
cochain equation for $R_A$.

By Lemma~\ref{lem:vol2-gaussian-choice-contraction},
$\delta L=B-A$. The same graded composition rule gives
\[
 \delta Q=-b(B-A)=P_B-P_A.
\]
Place parameter forms before coefficient maps. The
total-degree-zero covariance $P_u+du\,Q$ therefore
satisfies
\[
 (d_u+\delta)(P_u+du\,Q)=b-p^*b_H.
\]
Its exponential gives a total cochain map by the preceding
calculation. Its coefficient of $du$ is the integrand
in \eqref{eq:vol2-choice-gaussian-homotopy}. The
coefficient cochain equation, integrated from $0$ to $1$,
is the claimed homotopy identity. This integration is
polynomial integration over the rational ground ring.
Finally $L\mathsf E=0$, so $Q$ vanishes on two retained
coefficients; pair contraction also kills constants.
\end{proof}

The factors left after contraction matter. A quadratic
input has homotopy $\hbar Q(x,y)$; a cubic input has terms
$\hbar Q(x,y)p(z)$ and the other two pairings, with their
Koszul signs. A five-dimensional complex makes both
values nonzero. Take a basis
\[
 |u|=-1,\quad |v|=|r|=|e|=0,\quad |s|=1,
 \qquad du=v,\quad dr=s,
\]
with all other differentials zero. Retain $e$, and set
$h(v)=u$, $h(s)=r$, with $h$ zero on the other basis
vectors. Let $b(r,u)=b(u,r)=1$ and let its other basis
values vanish. This is a closed degree-one pairing.

Define the degree-minus-two endomorphism $J$ of
$V\otimes V$ by
\[
 J(s\otimes v)=J(v\otimes s)
                  =\tfrac12(r\otimes u+u\otimes r),
\]
and zero on other tensor basis vectors. It is normalized
and commutes with graded interchange. Thus
$A=A_0$ and $B=A_0+\delta J$ are normalized homotopies.
Their covariance data are
\begin{equation}\label{eq:vol2-choice-five-dimensional}
 P_A(r,v)=-\tfrac12,
 \qquad P_B(r,v)=\tfrac12,
 \qquad Q(s,v)=1.
\end{equation}
Indeed \eqref{eq:vol2-reference-two-coefficients} gives
$A_0(r\otimes v)=r\otimes u/2$, while
$\delta J(r\otimes v)=-(r\otimes u+u\otimes r)/2$.
Also $A_0\delta J(s\otimes v)
=(r\otimes u+u\otimes r)/2$, proving the third value.
Since $D_b(rv)=sv$ and $b_H=0$, the homotopy equation reads
\[
 R_B(rv)-R_A(rv)=\hbar=K_{B,A}(sv).
\]
After multiplication by the retained coefficient,
$K_{B,A}(sve)=\hbar e$.

\subsection{Prescribed parameter faces}

For the cube $I^m$, use polynomial parameter forms
\[
 \Omega_{\mathrm{poly}}^\bullet(I^m)
       =\mathbb Q[t_1,\ldots,t_m,dt_1,\ldots,dt_m],
\]
where the $dt_i$ are exterior generators of degree $1$.
A parameter family of normalized homotopies is a
total-degree-minus-one element $\mathcal A$ with
coefficients in $\mathcal N$ satisfying
\begin{equation}\label{eq:vol2-choice-total-equation}
 (d_{\mathrm{par}}+\delta)\mathcal A=1-\mathsf E.
\end{equation}
At a vertex, it is an ordinary normalized homotopy.
Restriction to a face sets its fixed coordinates to
$0$ or $1$ and their differentials to zero.

\begin{theorem}\label{thm:vol2-choice-relative-fillers}
A family satisfying \eqref{eq:vol2-choice-total-equation}
on a union of codimension-one faces, with agreeing
restrictions on intersections, extends to the cube.
Two extensions admit such a family between them
relative to the prescribed faces. The same assertion
holds after adjoining any finite number of parameters.
\end{theorem}
\begin{proof}
First extend the forms without their differential
equation. For a face $t_i=\varepsilon$, use its coordinate
projection and the scalar factor $1-t_i$ if
$\varepsilon=0$, or $t_i$ if $\varepsilon=1$.
Sum these extensions and subtract repeated extensions
from intersections by inclusion--exclusion. Opposite
faces have disjoint restriction conditions, since the
factor for one vanishes on the other. On each prescribed
face, terms from every other compatible face cancel
against their intersection terms. The resulting
polynomial form $\widetilde{\mathcal A}$ has the given
restrictions and coefficients in $\mathcal N$.

Write $D=d_{\mathrm{par}}+\delta$ and set
\[
 Z=D\widetilde{\mathcal A}-(1-\mathsf E).
\]
This is closed and restricts to zero on the prescribed
faces. Extend the contraction of
Lemma~\ref{lem:vol2-gaussian-choice-contraction} by
\[
 S_{\mathrm{par}}(\alpha\otimes M)
             =(-1)^{|\alpha|}\alpha\otimes A_0M.
\]
The parameter differential anticommutes with this
degree-minus-one map, and the coefficient identity
gives $DS_{\mathrm{par}}+S_{\mathrm{par}}D=1$.
Consequently
\begin{equation}\label{eq:vol2-choice-relative-correction}
 \mathcal A=\widetilde{\mathcal A}-S_{\mathrm{par}}Z
\end{equation}
solves \eqref{eq:vol2-choice-total-equation} and keeps
all the prescribed face values. To compare two
extensions, prescribe them at the ends of one more
interval and keep their common boundary family
constant along that interval. Apply the same correction.
Repeating this procedure gives the higher comparisons.
\end{proof}

Each family gives the total Gaussian map
$\operatorname{Sym}(p)\exp(\hbar\partial_{-b\mathcal A})$.
Composition with $b$ includes its sign past parameter
forms. The covariance calculation in
Proposition~\ref{prop:vol2-choice-gaussian-homotopy}
proves its total cochain equation. Thus the construction
supplies relative homotopies whenever the prescribed
Gaussian comparisons have these normalized coefficient
lifts. An arbitrary prescribed homotopy of observables
need not come with such a lift.

\section{From elimination to supported field operations}
\label{sec:vol2-state-geometric-comparison}

The finite constructions retain three kinds of data:
the relative support class in the normal direction,
the normalized state and all its cumulants, and the
parameter homotopies transporting that state. Conditional
elimination and coordinate changes construct these data
over integral quantum coefficients. A supplied total
master weight extends them over parameter faces, while
the normalized Gaussian choices above provide an explicit
class for which every relative extension exists.

The corresponding field-theoretic comparison requires a
common supported operation complex containing the
regular, meromorphic and contact coefficients of the
previous chapters. Its normal reduction must intertwine
the differential and every allowed collision operation;
its state map must carry the local vacuum and the
prescribed sewing and parameter homotopies. A regulated
master family must make its boundary terms and parameter
face equations meaningful on that same coefficient
domain. The finite endomorphism contraction above does
not yet act on those distributional products. These are
the missing constructions for a state-preserving
geometric comparison.

There is a second source of higher information when
observables obey algebraic constraints. Before applying
elimination to matrix observables, one must retain the
relations among their constraint equations. The resulting
complex is the starting point for Hamiltonian actions on
derived matrix spaces.

\chapter{Hamiltonian actions on derived matrix spaces}
\label{ch:vol2-derived-hamiltonian-actions}

Imposing a matrix equation has two different effects on an action.
A bracket defect can vanish on its zero set, while the complex
that imposes the equation retains a homotopy for that defect.
The homotopies must themselves satisfy identities. For the
commuting-matrix equation, a calculation with cubic Hamiltonians
exhibits the first of these identities. A contraction of a free
associative resolution then solves them in every arity.

The resulting action preserves a cyclic orientation coherently.
Its matrix realization retains both the relation classes and the
gauge stabilizers. Requiring the homotopies to be Hamiltonian adds
a further condition: the scalar charges remember constants that
their vector fields forget. Quantizing the coefficient plane
introduces another distinction. Its finite-dimensional
representations and the quantization of matrix coordinates have
different equations, and consequently different recipients.

\section{The cubic defect and the relation complex}
\label{sec:vol2-matrix-cubic-defect}

Work over $\mathbb C$, with cohomological differentials of degree
one. All commutators are graded:
$[a,b]=ab-(-1)^{|a||b|}ba$ for homogeneous elements.
Put $A=\mathbb C[q,p]$ and give it the Poisson bracket
\[
 \{f,g\}=f_qg_p-f_pg_q.
\]
For a monomial $q^ap^b$, let $\sigma(q^ap^b)$ be the average of
the distinct words with $a$ occurrences of $q$ and $b$ occurrences of $p$.
Thus $\sigma(qp)=(qp+pq)/2$. The map $\sigma$ takes values in
the free associative algebra; it is a linear section of
abelianization.

Fix a positive integer $N$. Let $X,Y$ be $N\times N$ matrices of
independent commuting variables, with
\[
 \{X_{ij},Y_{kl}\}=\delta_{il}\delta_{jk},\qquad
 \{X_{ij},X_{kl}\}=\{Y_{ij},Y_{kl}\}=0.
\]
The scalar charge associated to $f\in A$ is
\begin{equation}\label{eq:vol2-matrix-charge}
 J_N(f)=\operatorname{Tr}\sigma f(X,Y).
\end{equation}
In particular $J_N(1)=N$. Write $\mu=[X,Y]=XY-YX$ for the
commuting-matrix equation. The brackets of cubic charges give
\begin{equation}\label{eq:vol2-cubic-trace-defect}
 \{J_N(q^3),J_N(p^3)\}-J_N(9q^2p^2)
   =-\frac32\operatorname{Tr}(\mu^2).
\end{equation}
Indeed, the first term is $9\operatorname{Tr}(X^2Y^2)$.
Of the six words in $\sigma(q^2p^2)$, four have cyclic trace
$\operatorname{Tr}(X^2Y^2)$ and two have cyclic trace
$\operatorname{Tr}(XYXY)$. Expanding $\operatorname{Tr}\mu^2$
proves the formula. At rank two, take
\[
 X=\begin{pmatrix}1&0\\0&0\end{pmatrix},\qquad
 Y=\begin{pmatrix}0&1\\1&0\end{pmatrix}.
\]
Then $\mu^2=-I_2$, so the defect is $3$.

The defect vanishes after imposing $\mu=0$. To retain how it
vanishes, introduce an odd matrix $\Psi$ of degree $-1$ and the
commutative differential graded algebra
\begin{equation}\label{eq:vol2-derived-matrix-algebra}
 R_N=\mathbb C[X_{ij},Y_{ij}]\otimes\Lambda(\Psi_{ij}),
 \qquad QX=QY=0,\qquad Q\Psi=[X,Y].
\end{equation}
Here and below, matrices with graded coefficients are multiplied
in their displayed order. The entries of $\Psi$ impose the
equations by a differential, including dependencies among them.
For example $\operatorname{Tr}\Psi$ is a cycle. It is not a
boundary: evaluation at $X=Y=0$ is a cochain map to
$\Lambda(\Psi_{ij})$ with zero differential and sends this cycle
to a nonzero element. At $N=1$ the differential vanishes entirely.
Thus the relation complex contains information beyond the
ordinary commuting scheme.

\section{A free resolution of the plane}
\label{sec:vol2-free-plane-resolution}

The homotopies can be constructed before choosing the matrix rank.
Consider the free differential graded associative algebra
\begin{equation}\label{eq:vol2-plane-resolution}
 R=\mathbb C\langle q,p,t\rangle,\qquad
 |q|=|p|=0,\quad |t|=-1,
 \qquad dq=dp=0,\quad dt=qp-pq.
\end{equation}
The projection $\pi:R\to A$ kills $t$ and makes $q,p$ commute.
Its linear section $i$ sends $q^ap^b$ to the normal word $p^bq^a$.
Both are cochain maps, but $i$ is not multiplicative:
$i(q)i(p)-i(qp)=dt$.

Replace the differential temporarily by $d_0t=qp$. Expand every
$t$ in a word into $qp$, marking the occurrence from which it
came. Two occurrences of $qp$ cannot overlap. For a fixed
expanded word, its marked subsets form a tensor product of
two-term complexes. The differential $d_0$ removes a mark, with
sign $(-1)^a$ if $a$ marks precede it.

Define $h_0$ by adding the first mark when that occurrence is
unmarked, and by zero when it is already marked or no occurrence
exists. Let $p_0$ send an unmarked normal word $p^bq^a$ to
$q^ap^b\in A$ and every other word to zero. The two-term
contraction gives
\begin{equation}\label{eq:vol2-first-word-contraction}
 d_0h_0+h_0d_0=1-ip_0,
 \qquad h_0^2=h_0i=p_0h_0=0.
\end{equation}
To see the signs, removing a later mark anticommutes with adding
the first. Removing the first mark contributes the identity.
When no occurrence exists, the word is already $p^bq^a$.

Write $d=d_0+\delta$, where $\delta t=-pq$. Exchanging an
expanded $qp$ for $pq$ decreases the number of pairs in which a
$q$ precedes a $p$. The operator $\delta h_0$ is therefore locally
nilpotent. Define
\begin{equation}\label{eq:vol2-word-contraction}
 h=h_0\sum_{j\geq0}(-\delta h_0)^j,
 \qquad p'=p_0\sum_{j\geq0}(-\delta h_0)^j.
\end{equation}
These are finite sums on every word.

\begin{proposition}\label{prop:vol2-plane-word-resolution}
The projection $p'$ equals $\pi$, and
\[
 \pi i=1,\qquad dh+hd=1-i\pi,
 \qquad h^2=hi=\pi h=0.
\]
Consequently $R\to A$ is a quasi-isomorphism and $h$ contracts
$J=\ker\pi$.
\end{proposition}
\begin{proof}
Substitute $(1+\delta h_0)^{-1}$ in
\eqref{eq:vol2-first-word-contraction}. Multiplying the finite
series by $1+\delta h_0$ cancels consecutive perturbation terms.
Since $\delta i=0$, the inclusion and the zero differential on
normal words remain unchanged. This gives
$dh+hd=1-ip'$ and the side identities.

The surviving reductions replace $qp$ by $pq$. A degree-zero
word therefore projects to its commutative monomial; a word
containing $t$ has negative degree and projects to zero. Thus
$p'=\pi$. On $J$ the homotopy identity becomes $dh+hd=1$.
\end{proof}

For a commutative differential graded algebra $B$, a morphism
$R\to\operatorname{Mat}_N(B)$ amounts to matrices $X,Y,\Psi$
with the differential in \eqref{eq:vol2-derived-matrix-algebra}.
This is the universal property represented by $R_N$.
The resolution of $A$ does not make $R_N$ acyclic in negative
degrees; the trace class already distinguishes the two assertions.

\section{The first action and its binary homotopy}
\label{sec:vol2-first-derived-action}

The Lie algebra of polynomial Hamiltonian vector fields is
$\mathfrak l=A/\mathbb C1$, with the induced Poisson bracket.
For a word $a_1\cdots a_m$ in $q,p$, define its cyclic derivative
by
\[
 \partial_q^{\mathrm{cyc}}(a_1\cdots a_m)
 =\sum_{a_j=q}a_{j+1}\cdots a_ma_1\cdots a_{j-1},
\]
and similarly for $p$. Removing an occurrence makes the remainder
start just after that occurrence. Cyclic rotation therefore does not
change the derivative.

For $f\in\mathfrak l$, define a degree-zero derivation of $R$ by
\begin{equation}\label{eq:vol2-hamiltonian-first-lift}
 F_1(f)q=-\partial_p^{\mathrm{cyc}}\sigma f,
 \qquad F_1(f)p=\partial_q^{\mathrm{cyc}}\sigma f,
 \qquad F_1(f)t=0.
\end{equation}
A constant has zero cyclic derivatives, so this is well defined.
The identity
\begin{equation}\label{eq:vol2-cyclic-derivative-identity}
 [q,\partial_q^{\mathrm{cyc}}W]
   +[p,\partial_p^{\mathrm{cyc}}W]=0
\end{equation}
follows by cancelling consecutive cyclic rotations of $W$.
It implies $[d,F_1(f)]=0$. Abelianization gives
$\pi F_1(f)=X_f\pi$, where $X_f=\{f,-\}$.

Set $\mu_R=qp-pq$. Cyclic differentiation gives
\begin{equation}\label{eq:vol2-cubic-vector-defect}
 \bigl([F_1(q^3),F_1(p^3)]-F_1(9q^2p^2)\bigr)(q,p,t)
   =\bigl(3[\mu_R,q],3[\mu_R,p],0\bigr).
\end{equation}
The derivation
\begin{equation}\label{eq:vol2-cubic-word-primitive}
 F_2(q^3,p^3)(q,p,t)
     =\bigl(3[t,q],3[t,p],3t^2\bigr)
\end{equation}
has degree $-1$ and has this defect as its differential
$[d,F_2]$. On $q,p$ the assertion follows from $dt=\mu_R$.
On $t$, the identity
\[
 d(t^2)=\mu_Rt-t\mu_R
\]
cancels the term $F_2(\mu_R)$. The value on $t$ is therefore
necessary even though the original defect vanishes on $t$.

The failure of strictness cannot be removed by another polynomial
choice of cyclic words. Let
\[
 \mathcal N=
 \mathbb C\langle q,p\rangle/
 \bigl(\mathbb C1+[\mathbb C\langle q,p\rangle,
                         \mathbb C\langle q,p\rangle]\bigr),
\]
where the commutators span a vector subspace. Its elements are
cyclic words modulo constants. The necklace bracket is
\[
 \{[U],[V]\}_{\mathcal N}
   =\bigl[\partial_q^{\mathrm{cyc}}U\,
               \partial_p^{\mathrm{cyc}}V
          -\partial_p^{\mathrm{cyc}}U\,
               \partial_q^{\mathrm{cyc}}V\bigr].
\]
This bracket has a useful derivation description. Send a cyclic
word $U$ to $D_Uq=-\partial_p^{\mathrm{cyc}}U$ and
$D_Up=\partial_q^{\mathrm{cyc}}U$. Differentiating the displayed
bracket gives
$[D_U,D_V]=D_{\{[U],[V]\}_{\mathcal N}}$: terms selecting
two positions in one cyclic word cancel in pairs, while terms
selecting one position in each word give the two derivation
composites. This also proves that the bracket is well defined
under cyclic rotation. The map to derivations is injective
modulo constants, since for a homogeneous word of length $r>0$,
\[
 [q\,\partial_q^{\mathrm{cyc}}U
   +p\,\partial_p^{\mathrm{cyc}}U]=r[U].
\]
Jacobi follows from Jacobi for derivations and this injectivity.
Abelianization is a Lie map $\mathcal N\to\mathfrak l$.

\begin{proposition}\label{prop:vol2-necklace-no-section}
The abelianization $\mathcal N\to\mathfrak l$ has no linear
Lie section, including sections that do not preserve polynomial
degree.
\end{proposition}
\begin{proof}
In degrees at most three there is only one cyclic word for each
commutative monomial. A lift of a cubic therefore differs from
its cyclic cubic by terms of degree at least four. Bracketing
such a correction with a cubic has degree at least five.

Put $C=[q^2p^2]$ and $D=[qpqp]$. They are distinct cyclic words,
and direct differentiation gives
\[
 \{[q^3],[p^3]\}_{\mathcal N}=9C,
 \qquad \{[q^2p],[qp^2]\}_{\mathcal N}=C+2D.
\]
The commutative brackets are respectively $9q^2p^2$ and
$3q^2p^2$. The degree-four part of a Lie section applied to
$q^2p^2$ would have to equal both $C$ and $(C+2D)/3$.
Higher-degree corrections cannot change this contradiction.
\end{proof}

This obstruction concerns a strict polynomial necklace section.
The relation generator in \eqref{eq:vol2-plane-resolution}
provides the larger differential graded carrier on which
\eqref{eq:vol2-cubic-word-primitive} solves the first defect.

\section{Solving the action equations in every arity}
\label{sec:vol2-all-arity-derived-action}

Give $\operatorname{Der}(R)$ its graded commutator and differential
$\partial D=[d,D]$. The derivations with zero projection form
a differential graded Lie algebra
\[
 K=\{D:\pi D=0\}.
\]
A degree-$j$ derivation in $K$ is specified by
$U=Dq,V=Dp\in J^j$ and $W=Dt\in J^{j-1}$.
Put $L(U,V)=[U,p]+[q,V]$. Then
\begin{equation}\label{eq:vol2-derivation-differential}
 (\partial D)(q,p,t)
   =\bigl(dU,dV,dW-(-1)^jL(U,V)\bigr).
\end{equation}
The contraction of words yields a contraction of this entire
derivation complex, including its nonclosed elements.

\begin{lemma}\label{lem:vol2-derivation-kernel-contraction}
The degree-minus-one map
\begin{equation}\label{eq:vol2-derivation-kernel-homotopy}
 (HD)(q,p,t)
   =\bigl(hU,hV,h(W-(-1)^jL(hU,hV))\bigr)
\end{equation}
preserves $K$ and satisfies $\partial H+H\partial=1_K$.
Every value of $H$ is a finite polynomial.
\end{lemma}
\begin{proof}
The ideal $J$ contains all arguments of $h$ in the formula.
Since $\pi h=0$, the resulting derivation again lies in $K$.
On $q,p$ the contraction identity is $dh+hd=1_J$.
Use $dL(U,V)=L(dU,dV)$ to substitute the formula on $t$.
The sum of the third components is
\[
 (dh+hd)W+
 (-1)^j hL\bigl((dh+hd-1)U,(dh+hd-1)V\bigr)=W.
\]
The word contraction is finite, so the formula has the stated
polynomial meaning.
\end{proof}

We specify all higher signs by a coalgebra convention.
For a differential graded Lie algebra $\mathfrak g$, let $s$
be the degree-minus-one shift to $\mathfrak g[1]$.
The cofree conilpotent symmetric coalgebra
$S^c(\mathfrak g[1])=\bigoplus_{n\geq0}S^n(\mathfrak g[1])$
has the unshuffle coproduct. Its Lie codifferential has
corestrictions
\begin{equation}\label{eq:vol2-lie-coalgebra-signs}
 b_1(sa)=-s\partial a,
 \qquad b_2(sa\,sb)=(-1)^{|a|}s[a,b].
\end{equation}
An $L_\infty$ morphism is a coaugmented coalgebra map commuting
with these codifferentials. On degree-zero source inputs, its
arity-$n$ corestriction is $sF_n$, where
$F_n:\Lambda^n\mathfrak l\to\operatorname{Der}(R)$ has degree
$1-n$.

Suppose $F_1,\ldots,F_{n-1}$ are defined. Their corestrictions
extend by set partitions to a coalgebra map $\mathcal F$ with
arity-$n$ corestriction zero. Define $\Omega_n$ by
\[
 \operatorname{pr}_1(b\mathcal F-\mathcal F b)\big|_{S^n}
     =s\Omega_n,
 \qquad F_n=H\Omega_n.
\]
For the first two residuals this says
\begin{align}
 \Omega_2(f,g)
    &=[F_1(f),F_1(g)]-F_1(\{f,g\}),
       \label{eq:vol2-action-second-residual}\\
 \Omega_3(f,g,k)
    &=\sum_{\mathrm{cyc}}
       \bigl([F_1(f),F_2(g,k)]-F_2(\{f,g\},k)\bigr).
       \label{eq:vol2-action-third-residual}
\end{align}
At arity four, the brackets between two binary components contribute
\begin{equation}\label{eq:vol2-action-fourth-binary-part}
 \begin{aligned}
 &-[F_2(f_1,f_2),F_2(f_3,f_4)]
   +[F_2(f_1,f_3),F_2(f_2,f_4)]\\
 &\hspace{3em}-[F_2(f_1,f_4),F_2(f_2,f_3)].
 \end{aligned}
\end{equation}
They accompany the unary--ternary and source-bracket terms.

\begin{theorem}\label{thm:vol2-all-arity-derived-action}
The specified first action and the recursion $F_n=H\Omega_n$
define an $L_\infty$ morphism
\[
 F:\mathfrak l\longrightarrow\operatorname{Der}(R).
\]
Every component has finite polynomial values on polynomial
inputs. For every $N\geq1$, matrix substitution gives a
$GL_N$-equivariant $L_\infty$ action on $R_N$, and hence on
$R_N^{GL_N}$.
\end{theorem}
\begin{proof}
The first component commutes with $d$ by
\eqref{eq:vol2-cyclic-derivative-identity}. Let
$E=b\mathcal F-\mathcal F b$. The identity
$bE+Eb=0$ follows from the square-zero codifferentials.
If $E$ vanishes below arity $n$, its arity-$n$ corestriction
therefore satisfies $\partial\Omega_n=0$.

The first action projects to the strict Hamiltonian action on
$A$, so $\Omega_2$ projects to zero. Moreover $F_1$ normalizes
$K$, and every constructed component of arity at least two
belongs to $K$. The same is consequently true of each subsequent
residual. Lemma~\ref{lem:vol2-derivation-kernel-contraction}
gives $\partial H\Omega_n=\Omega_n$. In the coalgebra equation,
the new term is $-s\partial F_n$, by
\eqref{eq:vol2-lie-coalgebra-signs}; it cancels $s\Omega_n$.
Induction proves all the identities.

At each arity there are finitely many partitions. The word
contraction terminates on every term, proving polynomial
finiteness. A word derivation acts on matrix entries by
substituting $X,Y,\Psi$ for $q,p,t$ in its generator values.
The graded Leibniz rule makes this substitution preserve
differentials and brackets. It also commutes with simultaneous
conjugation, which proves the two matrix assertions.
\end{proof}

The contraction gives explicit higher terms beyond the binary
homotopy. For
$\mathbf f=(q^3,p^3,qp^3)$ and
$\mathbf g=(q^3,p^3,qp^3,q^2p^2)$, one obtains
\begin{equation}\label{eq:vol2-derived-higher-examples}
 \begin{aligned}
 F_3(\mathbf f)q&=-\frac92(tpt+pt^2),\\
 F_4(\mathbf g)q&=\frac{33}{5}t^3p+\frac{39}{5}t^2pt
                   +\frac{69}{5}tpt^2+\frac{99}{5}pt^3.
 \end{aligned}
\end{equation}
The other generator values are
$F_3(\mathbf f)p=F_3(\mathbf f)t=0$ and
$F_4(\mathbf g)p=F_4(\mathbf g)t=0$.
These formulas follow by applying
\eqref{eq:vol2-derivation-kernel-homotopy} to the residuals.
They can also be checked by substituting $dt=qp-pq$ in
$\partial F_n=\Omega_n$ on each generator. The four-input
identity uses \eqref{eq:vol2-action-fourth-binary-part}.
Thus this action does not truncate at its binary term.
The displayed values depend on the contraction; a nonzero
chosen component alone does not establish an invariant
obstruction to changing the action's presentation.

\section{Choices and pointed formal Hamiltonians}
\label{sec:vol2-derived-action-choices}

Let $\operatorname{Der}_\pi(R)$ be the derivations preserving
$J=\ker\pi$. Projection gives an exact sequence of differential
graded Lie algebras
\begin{equation}\label{eq:vol2-derived-action-fibration}
 0\longrightarrow K\longrightarrow\operatorname{Der}_\pi(R)
   \xrightarrow{\pi_*}\operatorname{Der}(A)\longrightarrow0.
\end{equation}
Surjectivity is degreewise. A derivation of $A$, concentrated in
degree zero, lifts by choosing words for its values on $q,p$
and giving value zero on $t$. Such a lift preserves $J$, although
it need not commute with $d$. The contracted kernel shows that
$\pi_*$ is a quasi-isomorphism.

For a differential graded Lie algebra $L$, retain homotopies
between actions by using the arity-completed
convolution differential graded Lie algebra
\[
 \mathcal G(L)=\prod_{n\geq1}
    \operatorname{Hom}(S^n(\mathfrak l[1]),L),
 \qquad
 \mathcal F^r\mathcal G(L)=\prod_{n\geq r}
    \operatorname{Hom}(S^n(\mathfrak l[1]),L).
\]
The differential combines the source Chevalley codifferential
with that of $L$; the bracket uses the reduced coproduct.
For a degree-one element $\alpha$, the Maurer--Cartan equation
$\partial\alpha+[\alpha,\alpha]/2=0$ is the $L_\infty$
morphism equation in the convention
\eqref{eq:vol2-lie-coalgebra-signs}. Polynomial differential
forms on a simplex give parameter families of such equations.
Tensoring arity by arity defines their simplicial space.

The \emph{derived space of lifts} is the homotopy fibre, over
$f\mapsto X_f$, of the map of these spaces induced by $\pi_*$.
This fixes the action on $A$ and the comparison to it; it does
not yet impose a cyclic or Hamiltonian condition on the lifts.

\begin{proposition}\label{prop:vol2-hamiltonian-lift-space}
The derived space of lifts of $X$ to
$\operatorname{Der}_\pi(R)$ is nonempty and contractible.
The word contraction selects one point of this space.
\end{proposition}
\begin{proof}
On the arity-$n$ associated graded, the kernel is
\[
 \operatorname{Hom}(S^n(\mathfrak l[1]),K),
\]
contracted by postcomposition with $H$, with the tensor sign
when parameter forms are present. Polynomial forms on a simplex
restrict degreewise onto compatible forms on a boundary or horn.
One proves this by extending faces successively: coefficients
vanishing on earlier faces lie in the corresponding ideals
$(x_i,dx_i)$, and the barycentric relation $\sum x_i=1$
handles the last face. No cochain primitive in the parameter
direction is required.

Suppose the lower arities and a compatible boundary family are
already lifted. Choose an arity-$n$ component with that boundary
and the prescribed projection. Its curvature lies in the kernel
above, tensored with forms relative to the boundary. The
arity-$n$ Bianchi identity makes the curvature closed. The
contraction gives a relative primitive; subtracting it solves
the equation without changing the boundary or projection.
The same argument fills horns. Induction and arity completeness
give a trivial fibration of simplicial spaces. Its nonempty
fibre is contractible.
\end{proof}

There is also a completion in polynomial weight. Give $q,p$
weight one and $t$ weight two; write $R_w$ for the homogeneous
piece and put $\widehat R=\prod_{w\geq0}R_w$.
The product topology has discrete homogeneous pieces. Assign
the same weights to $X,Y,\Psi$ to define $\widehat R_N$.
A pointed Hamiltonian has no constant or linear term.

\begin{proposition}\label{prop:vol2-pointed-action-completion}
On homogeneous inputs of degrees $d_1,\ldots,d_n$, the
derivation $F_n$ shifts weight by
\[
 \sum_{i=1}^n(d_i-2).
\]
At each arity, the action therefore extends continuously from
pointed polynomial Hamiltonians to
\[
 \widehat{\mathfrak l}_{\mathrm{pt}}
   =\prod_{d\geq2}\operatorname{Sym}^d(\mathbb Cq\oplus\mathbb Cp)
\]
on $\widehat R$, and to the corresponding equivariant action
on $\widehat R_N$.
\end{proposition}
\begin{proof}
Both $d$ and $h$ preserve weight. Cyclic differentiation gives
the claimed shift for $F_1$. Brackets add shifts, and the
derivation contraction preserves them, so induction proves the
formula for every component.

For pointed inputs, each $d_i-2$ is nonnegative. At a fixed
arity and output weight, only finitely many source weights
and input degrees can contribute. Their homogeneous spaces
are finite dimensional. The operation and each of its identities
are thus finite in every output weight, giving the continuous
extension. Matrix substitution preserves these weights.
\end{proof}

\section{Gauge ghosts and Hamiltonian primitives}
\label{sec:vol2-derived-action-bfv}

The matrices also carry simultaneous conjugation. Introduce
an odd matrix $c$ of degree $1$ and the polynomial graded
commutative algebra
\[
 \mathcal B_N=\mathbb C[X_{ij},Y_{ij}]
                  \otimes\Lambda(c_{ij},\Psi_{ij}).
\]
Extend the even Poisson bracket by
\[
 \{c_{ij},\Psi_{kl}\}=\{\Psi_{kl},c_{ij}\}
    =\delta_{il}\delta_{jk},
\]
with all other new generator brackets zero. The symmetry on
these two odd generators is the graded skew-symmetry of an
even Poisson bracket. The degree-one Hamiltonian
\[
 \Theta=\operatorname{Tr}(c[X,Y]-\Psi c^2)
\]
defines $Q=\{\Theta,-\}$, with
\begin{equation}\label{eq:vol2-classical-bfv-differential}
 QX=-[c,X],\quad QY=-[c,Y],\quad Qc=-c^2,
 \quad Q\Psi=[X,Y]-[c,\Psi].
\end{equation}
Here $[c,\Psi]=c\Psi+\Psi c$. Applying the graded Leibniz
rule once more to the four expressions gives $Q^2=0$.
This is the classical BFV, or Hamiltonian BRST, complex.
Equivalently it is the Chevalley--Eilenberg algebra
$C^*_{\mathrm{CE}}(\mathfrak{gl}_N;R_N)$ with the displayed
gauge sign. Its positive ghosts retain the infinitesimal gauge
action in addition to the negative relation generators.

\begin{proposition}\label{prop:vol2-action-bfv-extension}
Every component of the matrix action extends to $\mathcal B_N$
by giving value zero on $c$. These extensions define an
$L_\infty$ action of $\mathfrak l$ on $\mathcal B_N$.
\end{proposition}
\begin{proof}
Let $D$ be a word derivation of degree $j$ with $Dc=0$.
For a homogeneous matrix word $M$, the graded Leibniz rule gives
$D[c,M]=(-1)^j[c,DM]$. Thus the extended derivation commutes
in the graded sense with the gauge part of
\eqref{eq:vol2-classical-bfv-differential}. Extension preserves
both the differential and the bracket of derivations, so it
preserves every action identity.
\end{proof}

The first derivation is Hamiltonian: it equals
$H_{J_N(f)}=\{J_N(f),-\}$. Each $J_N(f)$ is $Q$-closed by
conjugation invariance. The cubic scalar defect has a
Hamiltonian primitive
\begin{equation}\label{eq:vol2-cubic-hamiltonian-primitive}
 K_N=-\frac32\operatorname{Tr}(\Psi\mu),\qquad
 QK_N=-\frac32\operatorname{Tr}(\mu^2).
\end{equation}
Graded cyclicity cancels the gauge terms in $QK_N$;
\eqref{eq:vol2-cubic-trace-defect} identifies its remaining
term with the bracket defect. Its Hamiltonian vector field is
\begin{equation}\label{eq:vol2-cubic-hamiltonian-field}
 \begin{aligned}
 H_{K_N}X&=\tfrac32[\Psi,X],&
 H_{K_N}Y&=\tfrac32[\Psi,Y],\\
 H_{K_N}c&=-\tfrac32\mu,&H_{K_N}\Psi&=0.
 \end{aligned}
\end{equation}
The extended word homotopy, in contrast, has
\begin{equation}\label{eq:vol2-cubic-derived-field}
 DX=3[\Psi,X],\quad DY=3[\Psi,Y],
 \quad Dc=0,\quad D\Psi=3\Psi^2.
\end{equation}
Both satisfy
\[
 [Q,D]=[Q,H_{K_N}]
   =[H_{J_N(q^3)},H_{J_N(p^3)}]-H_{J_N(9q^2p^2)}.
\]
The ghost component in \eqref{eq:vol2-cubic-hamiltonian-field}
is part of its Hamiltonian solution.

For $N\geq2$, $D$ is not Hamiltonian for the displayed bracket.
Indeed, $Dc=0$ would force a polynomial Hamiltonian for $D$ to
be independent of every $\Psi_{ij}$. Such a polynomial has no
component of degree $-1$, while $D$ is a nonzero degree-minus-one
derivation. A coherent derivation action therefore does not by
itself supply Hamiltonian representatives for its homotopies.

There is a prior scalar obstruction to a homotopy moment map
on the quotient $\mathfrak l$. Choose representatives with
constant term zero. Their bracket can acquire a constant:
the scalar cocycle is $\kappa(f,g)=\{f,g\}(0)$.

\begin{proposition}\label{prop:vol2-moment-central-obstruction}
For $N>0$, no $L_\infty$ morphism
$\mathfrak l\to(\mathcal B_N,Q,\{-,-\})$ has first component
$f\mapsto J_N(f)$ on those representatives. The obstruction
persists if the first components are changed by $Q$-boundaries.
\end{proposition}
\begin{proof}
The classes of $q,p$ have zero bracket in $\mathfrak l$, whereas
\[
 \{\operatorname{Tr}X,\operatorname{Tr}Y\}=N.
\]
The binary morphism equation would make $N$ a $Q$-boundary.
Evaluation of $X,Y,c,\Psi$ at zero is a cochain map to
$\mathbb C$, since every value in
\eqref{eq:vol2-classical-bfv-differential} is quadratic.
It sends the alleged boundary to the nonzero scalar $N$.
Exact changes of closed charges leave their cohomology bracket
unchanged, so they cannot remove this obstruction.
\end{proof}

Retaining the central extension $A\to\mathfrak l$, with
$J_N(1)=N$, removes this particular obstruction. The cocycle
also vanishes on the classical pointed subalgebra $(q,p)^2$.
On vector fields constants already act by zero, which explains
why Theorem~\ref{thm:vol2-all-arity-derived-action} is consistent
with the proposition.

The central extension contains a strict finite sector. Write
$e=q^2/2$, $f=p^2/2$ and $h_{\mathrm{quad}}=qp$.
Besides skew versions, its
nonzero brackets are
\[
 \{q,p\}=1,\quad\{e,p\}=q,\quad\{f,q\}=-p,
 \quad\{e,f\}=h_{\mathrm{quad}},
\]
\[
 \{h_{\mathrm{quad}},q\}=-q,\quad
 \{h_{\mathrm{quad}},p\}=p,\quad
 \{h_{\mathrm{quad}},e\}=-2e,\quad
 \{h_{\mathrm{quad}},f\}=2f.
\]
The map $J_N$ sends $1,q,p,e,h_{\mathrm{quad}},f$ to
$N,\operatorname{Tr}X,\operatorname{Tr}Y,
\operatorname{Tr}(X^2)/2,\operatorname{Tr}(XY),
\operatorname{Tr}(Y^2)/2$. Canonical matrix differentiation
gives precisely the same list of brackets. Thus scalar
quadratics act by a strict Hamiltonian Lie map of closed
functions. Higher scalar Hamiltonians require the homotopies
whose first instance is \eqref{eq:vol2-cubic-hamiltonian-primitive}.

\section{A cyclic orientation of the resolution}
\label{sec:vol2-derived-cyclic-orientation}

The symplectic plane has the volume form $dq\wedge dp$.
To preserve it on a resolution, the action must act on a complex
that remembers cyclic order as well as the differential.
Use normalized Hochschild chains
\[
 C_n(R)=R\otimes\overline R^{\otimes n},\qquad
 \overline R=R/\mathbb C1,
\]
written $a_0[a_1|\cdots|a_n]$. A unit in any bar slot represents
zero. A chain of internal degree $m=\sum_i|a_i|$ has total
cohomological degree $m-n$. Let $d_\otimes$ be the tensor
differential with its ordinary Koszul signs. In the raw-chain
convention the Hochschild differential is
\begin{equation}\label{eq:vol2-raw-hochschild-differential}
 d_{\mathrm H}=(-1)^n d_\otimes+b,
\end{equation}
where
\begin{align*}
 b(a_0[a_1|\cdots|a_n])
  ={}&a_0a_1[a_2|\cdots|a_n]\\
    &+\sum_{j=1}^{n-1}(-1)^j
        a_0[a_1|\cdots|a_ja_{j+1}|\cdots|a_n]\\
    &+(-1)^{n+|a_n|\sum_{i<n}|a_i|}
        a_na_0[a_1|\cdots|a_{n-1}].
\end{align*}
For $n=0$ set $b=0$. The normalized Connes operator is
\begin{equation}\label{eq:vol2-raw-connes-operator}
 B(a_0[a_1|\cdots|a_n])
 =\sum_{j=0}^n
 (-1)^{nj+(\sum_{i<j}|a_i|)(\sum_{i\geq j}|a_i|)}
 1[a_j|\cdots|a_n|a_0|\cdots|a_{j-1}].
\end{equation}
Adjacent face cancellations and cyclic rotation give
$d_{\mathrm H}^2=B^2=d_{\mathrm H}B+Bd_{\mathrm H}=0$.
Thus $(C_\bullet(R),d_{\mathrm H},B)$ is a mixed complex,
with the second operator of degree $-1$.

Let $\zeta$ have degree two. The negative-cyclic complex is
\[
 M_R=\bigl(C_\bullet(R)[[\zeta]],d_{\mathrm H}+\zeta B\bigr).
\]
Its degree-minus-two orientation is
\begin{equation}\label{eq:vol2-raw-orientation-cycle}
 \omega_0=1[q|p]-1[p|q]-1[t].
\end{equation}
The first two terms have Hochschild boundary $-1[qp-pq]$.
The internal differential of the last term, with the sign in
\eqref{eq:vol2-raw-hochschild-differential}, is $+1[qp-pq]$.
Its remaining Hochschild boundary is zero. Every Connes term
has a unit in a bar slot. Hence
$d_{\mathrm H}\omega_0=B\omega_0=0$.

The sign of $1[t]$ is tied to this totalization. Multiplication
of a chain of bidegree $(m,n)$ by $(-1)^{mn}$ changes the
differential to $d_\otimes+(-1)^m b$ and the Connes operator
to $(-1)^mB$; in that convention the last sign is positive.
The explicit formulas below use the raw convention throughout.

Let $\Omega_A^a$ denote polynomial differential $a$-forms with
de Rham differential $d_{\mathrm{dR}}$, and place it in
cohomological degree $-a$. The normalized
Hochschild--Kostant--Rosenberg projection is
\begin{equation}\label{eq:vol2-plane-hkr-projection}
 p_{\mathrm H}(a_0[a_1|\cdots|a_n])
  =\frac1{n!}\pi(a_0)\,d_{\mathrm{dR}}\pi(a_1)
                     \wedge\cdots\wedge d_{\mathrm{dR}}\pi(a_n).
\end{equation}
Substitution in the face and cyclic formulas gives
\[
 p_{\mathrm H}d_{\mathrm H}=0,
 \qquad p_{\mathrm H}B=d_{\mathrm{dR}}p_{\mathrm H},
 \qquad p_{\mathrm H}\omega_0=dq\wedge dp.
\]
The factorial is essential to the last normalization.

\begin{proposition}\label{prop:vol2-mixed-plane-contraction}
There is a weight-preserving contraction
$(p_{\mathrm H},i_{\mathrm H},H_{\mathrm H})$ from
$(C_\bullet(R),d_{\mathrm H})$ to $(\Omega_A^\bullet,0)$,
with $i_{\mathrm H}(dq\wedge dp)=\omega_0$ and
\[
 H_{\mathrm H}^2=H_{\mathrm H}i_{\mathrm H}
                =p_{\mathrm H}H_{\mathrm H}=0.
\]
Its negative-cyclic perturbation has unchanged projection and
target differential $\zeta d_{\mathrm{dR}}$. More precisely,
\begin{equation}\label{eq:vol2-mixed-plane-perturbation}
 \begin{aligned}
 i_\zeta&=(1+\zeta H_{\mathrm H}B)^{-1}i_{\mathrm H},\\
 H_\zeta&=H_{\mathrm H}(1+\zeta BH_{\mathrm H})^{-1},
       &p_\zeta&=p_{\mathrm H}.
 \end{aligned}
\end{equation}
The choices can be made by finite rational linear algebra at
each polynomial weight.
\end{proposition}
\begin{proof}
Give $q,p,t$ weights $1,1,2$ and each bar chain the sum of its
generator weights. Normalized bar slots have positive weight, so
the bar-length filtration is finite in each fixed weight.
The quasi-isomorphism $R\to A$ therefore gives a
quasi-isomorphism of Hochschild chains weight by weight:
on each associated-graded piece the comparison is a tensor
product of quasi-isomorphisms over the field.

Write $A^e=A\otimes A^{\mathrm{op}}$, and use subscripts $L,R$
for its two factors. The diagonal of $A$ has the bimodule
Koszul resolution
\[
 A^e\otimes\Lambda(e_q,e_p),\qquad
 de_q=q_L-q_R,\qquad de_p=p_L-p_R.
\]
The two differences form a regular sequence, since they become
independent polynomial variables after a linear change of
coordinates. Tensoring with $A$ over $A^e$ kills this
differential and identifies $e_q,e_p$ with $dq,dp$.
The bar resolution is another free resolution of the diagonal;
its augmented contraction inserts the unit. The comparison sends
$e_q\wedge e_p$ to the antisymmetrized bar
$1[q|p]-1[p|q]$. Consequently
\eqref{eq:vol2-plane-hkr-projection} is a surjective
quasi-isomorphism with the stated normalization.

Choose lifts of a basis of forms to cycles, using $\omega_0$
for the volume. These give $i_{\mathrm H}$. In the acyclic
kernel, choose degreewise splittings
\[
 (\ker p_{\mathrm H})^r=B^r\oplus L^r,
 \qquad d_{\mathrm H}:L^r\xrightarrow{\sim}B^{r+1}.
\]
Take $H_{\mathrm H}$ to be the inverse on $B^{r+1}$ and zero
on $L$ and the chosen cycle lifts. This proves the contraction
and side identities. Every space in a fixed weight is finite
dimensional; lexicographic word bases and row reduction make
the choices rational.

Apply perturbation by $\zeta B$. Its series converge in the
$\zeta$-adic topology. The relation
\[
 p_{\mathrm H}\zeta BH_{\mathrm H}
      =\zeta d_{\mathrm{dR}}p_{\mathrm H}H_{\mathrm H}=0
\]
keeps the projection fixed. It also kills every higher term
in the transferred differential
$p_{\mathrm H}\zeta B(1+\zeta H_{\mathrm H}B)^{-1}i_{\mathrm H}$,
leaving exactly $\zeta d_{\mathrm{dR}}$.
Since $B\omega_0=0$, the perturbed inclusion still sends the
volume to $\omega_0$.
\end{proof}

There is no unbounded filtration exchange in this comparison.
At fixed weight the bar length and the negative internal degree
are bounded. Fixing total degree then bounds the exponent of
$\zeta$. The completed identities can therefore be checked in
finite-dimensional pieces as well as in finite parameter quotients.

\section{Coherent preservation of the volume}
\label{sec:vol2-equivariant-cyclic-volume}

A derivation $D$ of degree $r$ acts on raw chains by
\begin{equation}\label{eq:vol2-raw-chain-lie-derivative}
 \mathcal L_D(a_0[a_1|\cdots|a_n])
  =\sum_{j=0}^n(-1)^{nr+r\sum_{i<j}|a_i|}
       a_0[\cdots|D(a_j)|\cdots|a_n],
\end{equation}
where the $j=0$ term acts on the outer coefficient. The factor
$(-1)^{nr}$ is required by the raw totalization. The tensor
signs give
\[
 [d_{\mathrm H},\mathcal L_D]=\mathcal L_{[d,D]},
 \qquad [B,\mathcal L_D]=0,
 \qquad [\mathcal L_D,\mathcal L_E]=\mathcal L_{[D,E]}.
\]
Thus the action $F$ makes $M_R$ an $L_\infty$ module.

For explicit totalization, write $C=S^c(\mathfrak l[1])$ and
let $b_C$ be its Chevalley codifferential. The twisting cochain
$\alpha:C\to\operatorname{Der}(R)$ has degree one and
\[
 \alpha_n(sf_1\cdots sf_n)=F_n(f_1,\ldots,f_n).
\]
Its equation is
$\partial\alpha+\alpha b_C+[\alpha,\alpha]/2=0$,
where convolution uses the tensor Koszul sign. On the completed
coefficient complex
\[
 \mathcal C_R=\prod_{n\geq0}
       \operatorname{Hom}(S^n(\mathfrak l[1]),M_R),
\]
the differential of a homogeneous map $a$ is
\begin{equation}\label{eq:vol2-equivariant-cyclic-differential}
 D_Fa=d_{M_R}a-(-1)^{|a|}a b_C+\mathcal L_\alpha\star a.
\end{equation}
Here $\star$ denotes action convolution with the unshuffle
coproduct.
The twisting-cochain equation implies $D_F^2=0$.

The projection to forms intertwines the first action and kills
the higher ones:
\begin{equation}\label{eq:vol2-equivariant-hkr}
 p_{\mathrm H}\mathcal L_{F_1(f)}
      =\mathcal L_{X_f}p_{\mathrm H},
 \qquad p_{\mathrm H}\mathcal L_{F_n}=0\quad(n\geq2).
\end{equation}
The first equation is naturality of the differential-form
formula. For the second, a negative-degree value of a
derivation contains $t$ and is killed by $\pi$.

A coherent lift is a $D_F$-closed element of total degree
$-2$ whose projection is the volume, regarded as a constant
arity-zero cochain. Its arity-$n$ component is an alternating
map $\omega_n:\Lambda^n\mathfrak l\to M_R^{-n-2}$.
Polynomial parameter forms retain homotopies over this fixed
projection and define the derived space of such lifts.

\begin{theorem}\label{thm:vol2-coherent-cyclic-orientation}
The orientation $\omega_0$ extends to a $D_F$-closed element
$\omega_{\mathrm{eq}}$ of total degree $-2$.
Over the fixed Hamiltonian action and fixed projection to
$dq\wedge dp$, the derived space of coherent negative-cyclic
lifts is contractible. The construction is continuous for
pointed Hamiltonians in the positive-weight completion.
\end{theorem}
\begin{proof}
Extend the contraction of
Proposition~\ref{prop:vol2-mixed-plane-contraction} to the
convolution complex, with its total signs. Let
$\mathcal H_\zeta$ denote the extended homotopy and let
$\delta_F$ denote the terms in
\eqref{eq:vol2-equivariant-cyclic-differential} beyond the
coefficient differential $d_{M_R}$. They raise arity.
Perturbation in the arity filtration gives
\begin{equation}\label{eq:vol2-equivariant-volume-formula}
 \omega_{\mathrm{eq}}
    =(1+\mathcal H_\zeta\delta_F)^{-1}
          i_\zeta(dq\wedge dp).
\end{equation}
Equation~\eqref{eq:vol2-equivariant-hkr} and
$p_{\mathrm H}\mathcal H_\zeta=0$ imply
$p_{\mathrm H}\delta_F\mathcal H_\zeta=0$.
The transferred differential is therefore precisely the strict
Hamiltonian action on forms together with
$\zeta d_{\mathrm{dR}}$. The volume is closed and its Lie
derivative under $X_f$ is zero, since
$\partial_q(-f_p)+\partial_p(f_q)=0$.
Its perturbed inclusion is consequently closed. Its arity-zero
component is $\omega_0$.

The same perturbation contracts the kernel of the fixed
projection. The homotopy fibre over the specified closed volume
is an affine translate of this acyclic complex. Polynomial
parameter forms and its contraction fill relative families,
so its derived mapping space is contractible.

At fixed arity only finitely many arity-raising factors occur.
At fixed $\zeta$ precision the mixed perturbation is finite.
Finally, the weight estimate in
Proposition~\ref{prop:vol2-pointed-action-completion} makes
every output weight a finite sum for pointed inputs. These
three finiteness statements prove the completion assertion.
\end{proof}

Two coherence levels can be written without a matrix inversion.
Abbreviate $1[a_1|\cdots|a_n]$ by
$\langle a_1|\cdots|a_n\rangle$. The arity-one equation is
\[
 d_{\mathrm H}\omega_1(f)=-\mathcal L_{F_1(f)}\omega_0.
\]
Choose
\begin{equation}\label{eq:vol2-cyclic-cubic-primitives}
 \omega_1(q^3)=-3\langle q|q|q\rangle,
 \qquad \omega_1(p^3)=-3\langle p|p|p\rangle.
\end{equation}
For instance,
$d_{\mathrm H}\langle q|q|q\rangle
=\langle q|q^2\rangle-\langle q^2|q\rangle$.
A compatible quartic prescription is
\begin{align}
 \omega_1(q^2p^2)
 ={}&-\tfrac23\langle tt\rangle
       \label{eq:vol2-cyclic-quartic-primitive}\\
 &+\tfrac23\bigl(
       \langle q|pt\rangle-\langle q|tp\rangle
      -\langle p|qt\rangle+\langle p|tq\rangle\bigr)\notag\\
 &+\tfrac23\bigl(
       \langle qt|p\rangle-\langle pt|q\rangle
      -\langle tq|p\rangle+\langle tp|q\rangle\bigr)\notag\\
 &-2\bigl(\langle q|p|qp\rangle
           +\langle p|qp|q\rangle+\langle qp|q|p\rangle\bigr).
       \notag
\end{align}
All three chains have zero Connes boundary and zero HKR
projection. For $f=q^3,g=p^3$, the second equation in the same
convention is
\begin{align}
 d_{\mathrm H}\omega_2(f,g)
 ={}&\omega_1(\{f,g\})
       -\mathcal L_{F_1(f)}\omega_1(g)
       +\mathcal L_{F_1(g)}\omega_1(f)\notag\\
    &-\mathcal L_{F_2(f,g)}\omega_0.
       \label{eq:vol2-cyclic-second-equation}
\end{align}
Its solution
\begin{align}
 \tfrac19\omega_2(q^3,p^3)
 ={}&\langle t|t\rangle
      +\langle q|p|t\rangle-\langle p|q|t\rangle
       \label{eq:vol2-cyclic-second-primitive}\\
    &-\langle q|t|p\rangle+\langle p|t|q\rangle
      +\langle t|q|p\rangle-\langle t|p|q\rangle\notag\\
    &-\langle q|q|p|p\rangle+\langle q|p|q|p\rangle
      +\langle q|p|p|q\rangle\notag\\
    &+\langle p|q|q|p\rangle-\langle p|q|p|q\rangle
      -\langle p|p|q|q\rangle\notag
\end{align}
also has zero Connes boundary and HKR projection.
Substitution in the raw differential verifies these identities:
internal terms replace $t$ by $qp-pq$, and adjacent mergers
cancel the cyclic faces. In particular the
$\langle t|t\rangle$ term is required in the second equation.
The contracted projection kernel extends these prescribed
lower components to all arities, even if a particular
row-reduction rule chooses different representatives.

\section{The finite matrix tangent pairing}
\label{sec:vol2-matrix-tangent-pairing}

The diagonal Koszul resolution also explains the nondegeneracy
of the orientation. Its top exterior product pairs the
resolution with its dual and gives
\[
 \operatorname{RHom}_{A^e}(A,A^e)\simeq A[-2].
\]
The volume $dq\wedge dp$ fixes this identification. Together
with its negative-cyclic lift, this is the smooth
$2$-Calabi--Yau orientation of the plane: the diagonal is a
perfect bimodule, and its bimodule dual is the displayed shift.
Extension and restriction along the quasi-isomorphism $R\to A$
transport this duality to $R$. Properness would require finite
total cohomology over $\mathbb C$, which $A$ does not have.
The orientation is a bimodule duality, distinct from a
nondegenerate scalar trace on a proper algebra.

Take the quotient by $GL_N$ of the derived affine scheme
represented by $R_N$. At a commuting pair $(X,Y)$ its tangent
complex is
\begin{equation}\label{eq:vol2-derived-matrix-tangent}
 \underbrace{\mathfrak{gl}_N}_{-1}
 \xrightarrow{\ d_0\ }
 \underbrace{\operatorname{Mat}_N\oplus\operatorname{Mat}_N}_{0}
 \xrightarrow{\ d_1\ }
 \underbrace{\mathfrak{gl}_N}_{1},
\end{equation}
with
\[
 d_0(a)=([a,X],[a,Y]),\qquad
 d_1(U,V)=[U,Y]+[X,V].
\]
The first arrow differentiates conjugation; the second
differentiates the relation. Their composite is
$[a,[X,Y]]=0$. Put
\begin{equation}\label{eq:vol2-matrix-tangent-form}
 \omega((U,V),(U',V'))=\operatorname{Tr}(UV'-VU')
\end{equation}
on the middle term. Pair the terms in degrees $-1$ and $1$
by $\operatorname{Tr}(ab)$, symmetrically under reversal as
required by graded skew-symmetry between two odd degrees.
Trace invariance gives
\[
 \omega(d_0a,(U,V))=\operatorname{Tr}(a\,d_1(U,V)).
\]
Thus these pairings form a degree-zero cochain pairing.
They are perfect term by term, including at points with
stabilizers. Discarding either outer term would discard part
of this self-dual tangent complex. This finite tangent
calculation and coherent cyclic invariance do not require
every chosen derivation homotopy to be Hamiltonian for one
fixed chain-level symplectic form.

\section{Deforming the coefficient relation}
\label{sec:vol2-coefficient-weyl-action}

Let $k_\varepsilon=\mathbb C[\varepsilon]$, with
$|\varepsilon|=0$, and define
\begin{equation}\label{eq:vol2-quantum-plane-resolution}
 \begin{gathered}
 W_\varepsilon=k_\varepsilon\langle q,p\rangle/
                              (qp-pq-\varepsilon),\\
 R_\varepsilon=k_\varepsilon\langle q,p,t\rangle,
 \qquad |t|=-1,\qquad dt=qp-pq-\varepsilon,
 \qquad dq=dp=0.
 \end{gathered}
\end{equation}
The projection $\pi_\varepsilon$ kills $t$.
Normal words $p^bq^a$ give a $k_\varepsilon$-basis of
$W_\varepsilon$: the replacement $qp\mapsto pq+\varepsilon$
terminates, and the leading pattern $qp$ has no overlap
ambiguity. They define a linear section $i_\varepsilon$.

Use the same marked-word contraction $h_0$ as before, with
$d_0t=qp$, but now put $\delta t=-pq-\varepsilon$.
Order expanded words first by length and then by the number
of pairs in which $q$ precedes $p$. Exchanging $qp$ and $pq$
decreases the second number; the constant term decreases the
first. Every term of $\delta h_0$ therefore decreases this
well-founded order.

\begin{proposition}\label{prop:vol2-integral-weyl-resolution}
The finite series
\[
 h_\varepsilon=h_0(1+\delta h_0)^{-1},\qquad
 p'=p_0(1+\delta h_0)^{-1}
\]
satisfy
\[
 dh_\varepsilon+h_\varepsilon d=1-i_\varepsilon\pi_\varepsilon,
 \quad h_\varepsilon^2=h_\varepsilon i_\varepsilon
                      =\pi_\varepsilon h_\varepsilon=0,
 \quad p'=\pi_\varepsilon.
\]
They extend continuously to the $\varepsilon$-adic completions,
allowing word support to grow with parameter order.
\end{proposition}
\begin{proof}
The marked-word proof of
Proposition~\ref{prop:vol2-plane-word-resolution} applies with
the decreasing order just described. Each surviving reduction
is precisely the PBW rule $qp\mapsto pq+\varepsilon$, so the
perturbed projection is $\pi_\varepsilon$. No division by
$\varepsilon$ occurs.

An element of the completed source is a series
$\sum_{r\geq0}\varepsilon^r a_r$ with each $a_r$ a finite
word polynomial. Every perturbation has nonnegative parameter
order. At a fixed output order, only finitely many earlier
coefficients contribute, and each has a terminating word
calculation. This defines the continuous extensions. The
identities hold modulo every power of $\varepsilon$, hence in
the separated inverse limit.
\end{proof}

Let $m(f\otimes g)=fg$ be ordinary polynomial multiplication.
On $k_\varepsilon[q,p]$, the Moyal product is
\begin{equation}\label{eq:vol2-coefficient-moyal-product}
 f\star_\varepsilon g
 =m\exp\!\left[\frac\varepsilon2
       (\partial_q\otimes\partial_p-\partial_p\otimes\partial_q)
       \right](f\otimes g).
\end{equation}
It is finite on polynomial inputs. On three tensor factors,
the three constant bidifferential operators commute; both
parenthesizations are the exponential of their sum. This
proves associativity. Symmetrization identifies this product
with multiplication in $W_\varepsilon$. For linear forms $u,v$,
centrality of their Weyl commutator gives
\[
 e^{\sigma u}e^{\sigma v}
   =e^{\varepsilon\{u,v\}/2}e^{\sigma(u+v)}.
\]
Differentiation in a formal scalar parameter proves this
exponential identity. Taking coefficients in $u,v$ gives the
polynomial product identity.

The commutator is divisible by $\varepsilon$ in the polynomial
ring. Write
\[
 B_\varepsilon(f,g)=
          (f\star_\varepsilon g-g\star_\varepsilon f)/\varepsilon,
 \qquad
 \mathfrak l_\varepsilon=k_\varepsilon[q,p]/k_\varepsilon1.
\]
Division here means the polynomial quotient; the parameter is
not inverted. The quotient is a Lie algebra. The cyclic
derivative formula \eqref{eq:vol2-hamiltonian-first-lift}
still defines a closed derivation of $R_\varepsilon$.
Under $\pi_\varepsilon$ it induces the strict Moyal
derivation $B_\varepsilon(f,-)$. Indeed its values on $q,p$
are the Weyl-ordered derivatives $-f_p,f_q$, and those
generators determine a derivation of the Weyl algebra.

\begin{theorem}\label{thm:vol2-coherent-coefficient-quantum-action}
There is an $L_\infty$ morphism over $k_\varepsilon$
\[
 F^\varepsilon:\mathfrak l_\varepsilon
       \longrightarrow\operatorname{Der}_{k_\varepsilon}
                                   (R_\varepsilon)
\]
with the specified first component. Each Taylor component
is polynomial and finite at fixed arity on polynomial inputs.
The action extends continuously to the parameter completions
and gives, by matrix substitution, a $GL_N$-equivariant action
on the derived representation algebra for every $N$.
\end{theorem}
\begin{proof}
Let $K_\varepsilon$ be the derivations whose generator values
lie in $\ker\pi_\varepsilon$. The differential of a
degree-$j$ derivation with values $(U,V,W)$ is still
\eqref{eq:vol2-derivation-differential}: derivations over
$k_\varepsilon$ kill the added scalar.
Replacing $h$ by $h_\varepsilon$ in
\eqref{eq:vol2-derivation-kernel-homotopy} therefore contracts
$K_\varepsilon$ in every degree.

Use the coalgebra signs
\eqref{eq:vol2-lie-coalgebra-signs} and replace the source
Poisson bracket by $B_\varepsilon$. The first action projects
to a strict action on $W_\varepsilon$. Every higher residual
is consequently closed and lies in $K_\varepsilon$.
Applying its contraction supplies the next component, exactly
as in Theorem~\ref{thm:vol2-all-arity-derived-action}.
Finitely many partitions and the terminating word order prove
polynomial finiteness. The maps never decrease parameter
order, giving the continuous completed identities. Matrix
substitution preserves differentials, brackets and
conjugation over the same base ring.
\end{proof}

The first parameter-dependent correction is already visible
on a cubic and a quartic:
\begin{equation}\label{eq:vol2-first-quantum-hamiltonian-contact}
 \begin{aligned}
 B_\varepsilon(q^3,p^4)&=12q^2p^3+6\varepsilon^2p,\\
 F^\varepsilon_2(q^3,p^4)q
     &=6(-pqt+ptq-qtp+tpq-\varepsilon t),\\
 F^\varepsilon_2(q^3,p^4)p&=6(-p^2t+tp^2),\\
 F^\varepsilon_2(q^3,p^4)t&=6(pt^2+tpt).
 \end{aligned}
\end{equation}
The first and third odd Moyal terms give the bracket in the
first line. Applying $\partial=[d,-]$ to the remaining lines
gives the binary defect on all three generators. In particular
$d(-6\varepsilon t)$ contributes $+6\varepsilon^2$ to the
$q$-component, matching the contribution of
$-F_1(6\varepsilon^2p)$ to the defect.
The third component on the ordered triple
$(q^3,p^4,q^2p)$ is
\begin{equation}\label{eq:vol2-quantum-ternary-action}
 F^\varepsilon_3q=-12tqt+12t^2q,
 \qquad F^\varepsilon_3p=-12pt^2,
 \qquad F^\varepsilon_3t=-6t^3.
\end{equation}
Substitution of $dt=qp-pq-\varepsilon$ verifies its ternary
equation with the binary terms in
\eqref{eq:vol2-first-quantum-hamiltonian-contact}.

The classical positive-weight completion and this parameter
completion use different estimates. Higher Moyal terms can
lower polynomial degree, as the term $\varepsilon^2p$ shows.
The parameter-adic theorem does not assert that the classical
pointed polynomial subspace remains closed under the deformed
bracket.

There is a relative negative-cyclic orientation over the
parameter ring. Form normalized chains over $k_\varepsilon$,
so that a parameter multiple of $1$ is zero in a bar slot.
In the raw totalization,
\begin{equation}\label{eq:vol2-relative-quantum-orientation}
 \omega_\varepsilon=1[q|p]-1[p|q]-1[t]
\end{equation}
satisfies
\[
 d_{\mathrm H}\omega_\varepsilon=-\varepsilon1[1]=0,
 \qquad B\omega_\varepsilon=0.
\]
This is a relative mixed-complex assertion. It uses the
parameter in the base ring; an absolute de Rham variation
of that parameter is additional data. The cycle itself also
does not supply an equivariant cyclic comparison to a
matrix-quantized recipient.

\section{Matrix quantization and finite representations}
\label{sec:vol2-two-quantizations}

Matrix quantization starts with the canonical matrix coordinates,
using an independent parameter $\hbar$. Its even Weyl algebra
$W_{N,\hbar}$ over $\mathbb C[\hbar]$ has
\[
 [X_{ij},Y_{kl}]=\hbar\delta_{il}\delta_{jk},\qquad
 [X_{ij},X_{kl}]=[Y_{ij},Y_{kl}]=0.
\]
Weyl ordering is symmetrization of all coordinate factors.
The Moyal formula on this finite-dimensional symplectic vector
space identifies its multiplication with the corresponding
constant-tensor star product. The commutator of two scalar
quadratic polynomials, divided by $\hbar$, is exactly their
Poisson bracket: the third and higher odd Moyal terms
differentiate each quadratic at least three times and vanish.
Thus the strict scalar quadratic sector constructed above
remains strict under matrix quantization.

The ordering of the matrix moment is part of this assertion.
As operator-valued matrices, set
\begin{equation}\label{eq:vol2-ordered-matrix-moment}
 \widehat\mu=XY-YX-N\hbar I_N.
\end{equation}
The Weyl symbols of $XY$ and $YX$ differ from their commuting
symbols by $N\hbar I_N/2$ and $-N\hbar I_N/2$, respectively.
Consequently $\widehat\mu$ has Weyl symbol $[X,Y]$ in
commuting coordinates. In particular
\[
 \operatorname{Tr}(XY-YX)=N^2\hbar,
 \qquad \operatorname{Tr}\widehat\mu=0.
\]
This scalar ordering correction is fixed by the coordinate
commutators. A specified nonzero value of a moment map is a
separate datum.

The full quantum BRST algebra also retains both ghost degrees.
Choose a basis $a_i$ of $\mathfrak{gl}_N$ with
$[a_i,a_j]=\sum_k f_{ij}^ka_k$, and set
$\ell_i=\operatorname{Tr}(a_i\widehat\mu)$.
These are the Weyl-ordered quadratic comoments, with
\[
 [\ell_i,\ell_j]=\hbar\sum_k f_{ij}^k\ell_k.
\]
This follows either by canonical matrix differentiation and
exact quadratic quantization, or by realizing
$X_{ij}=x_{ij}$ and $Y_{ij}=-\hbar\partial_{x_{ji}}$:
$\ell_a=\hbar\tau(a)$, where
$\tau(a)X=[X,a]$ is the infinitesimal pullback action.

Adjoin odd generators $c^i,b_i$ of degrees $1,-1$, commuting
with the even matter generators and satisfying
\[
 c^ic^j+c^jc^i=0,\qquad b_ib_j+b_jb_i=0,\qquad
 b_ic^j+c^jb_i=\hbar\delta_i^j.
\]
Write $\mathcal B_{N,\hbar}$ for the resulting graded tensor
product with $W_{N,\hbar}$. Define
\begin{equation}\label{eq:vol2-quantum-bfv-hamiltonian}
 \Omega=\sum_i c^i\ell_i
          -\frac12\sum_{i,j,k}f_{ij}^kc^ic^jb_k,
 \qquad Q_\hbar=\hbar^{-1}[\Omega,-]_{\mathrm{gr}}.
\end{equation}
Modulo $\hbar$ the algebra is supercommutative, so the
commutator is divisible by $\hbar$ and the formula is integral.
Under the trace identification of the ghost and its dual,
its classical limit is
\eqref{eq:vol2-classical-bfv-differential}.

\begin{proposition}\label{prop:vol2-integral-quantum-brst}
The element $\Omega$ satisfies $\Omega^2=0$.
Consequently $Q_\hbar$ is a degree-one derivation with
$Q_\hbar^2=0$ over $\mathbb C[\hbar]$.
\end{proposition}
\begin{proof}
Normal words, with all $c$ generators before all $b$ generators,
span the ghost algebra. Represent $c^i$ by exterior
multiplication and $b_i$ by $\hbar\partial_{c^i}$, with left
odd derivatives. After inverting $\hbar$, their normal words
form a basis of the endomorphisms of the exterior space.
They are therefore linearly independent over
$\mathbb C[\hbar]$ as well. Together with the Weyl basis this
proves parameter-torsion freeness and faithfulness on the
tensor product of the matter regular module and the ghost
exterior module.

In this representation,
\[
 \Omega=\sum_i c^i\ell_i+\hbar d_{\mathrm{CE}},\qquad
 d_{\mathrm{CE}}c^k=-\frac12\sum_{i,j}f_{ij}^kc^ic^j.
\]
Jacobi gives $d_{\mathrm{CE}}^2=0$. The square of the first
term is
$\frac{\hbar}{2}\sum f_{ij}^kc^ic^j\ell_k$; its mixed term
with $\hbar d_{\mathrm{CE}}$ is the negative of this expression.
Thus $\Omega^2=0$ in a faithful representation and in the
algebra. For odd $\Omega$, the double graded commutator with
$\Omega$ is the commutator with $\Omega^2$.
Hence $\hbar^2Q_\hbar^2=0$, and torsion freeness gives the result.
\end{proof}

The full gauge complex has extra states even at rank one.
Scalar conjugation is trivial, so $\ell=0$, the Lie bracket
vanishes and $Q_\hbar=0$. In its degree-zero ghost algebra,
$1$ and $u=cb$ are independent and obey
\[
 u^2=\hbar u.
\]
Thus the full BRST cohomology is not obtained by simply
replacing it with the ordinary one-pair Weyl reduction.
Removing the ineffective gauge direction changes the input.

In contrast, a finite-dimensional representation of the
\emph{coefficient} Weyl algebra has commuting coordinate
entries and the relation $[X,Y]=\varepsilon I_N$.
Its derived representation algebra is
\begin{equation}\label{eq:vol2-coefficient-derived-representations}
 R_{N,\varepsilon}
   =\mathbb C[X_{ij},Y_{ij},\varepsilon]\otimes\Lambda(\Psi_{ij}),
 \qquad Q\Psi=[X,Y]-\varepsilon I_N,
 \qquad QX=QY=Q\varepsilon=0.
\end{equation}
This is the recipient supplied by matrix substitution in
Theorem~\ref{thm:vol2-coherent-coefficient-quantum-action}.

\begin{proposition}\label{prop:vol2-finite-weyl-representations}
For $N>0$, multiplication by $\varepsilon$ is null-homotopic
on $R_{N,\varepsilon}$. Its cohomology is annihilated by
$\varepsilon$, and its localization at $\varepsilon$ is
contractible.
\end{proposition}
\begin{proof}
Commuting entries give $\operatorname{Tr}[X,Y]=0$.
The element $v=-\operatorname{Tr}\Psi/N$ has degree $-1$
and satisfies $Qv=\varepsilon$. Multiplication by $v$ obeys
\[
 Q(va)+vQa=\varepsilon a
\]
for every homogeneous $a$. After inverting $\varepsilon$,
multiplication by $v/\varepsilon$ contracts the complex.
\end{proof}

This contraction is integral before localization. It does not
contract the special fibre. At rank one the complex is
\[
 \mathbb C[x,y,\varepsilon]\psi
     \xrightarrow{-\varepsilon}
 \mathbb C[x,y,\varepsilon],\qquad |\psi|=-1.
\]
Its cohomology is $\mathbb C[x,y]$ in degree zero, annihilated
by $\varepsilon$. Derived specialization at $\varepsilon=0$
is computed by the displayed free complex and retains the
degree-minus-one class $\psi$. Simply specializing its
degree-zero cohomology would miss that class.

The equations \eqref{eq:vol2-ordered-matrix-moment} and
\eqref{eq:vol2-coefficient-derived-representations} therefore
cannot be interchanged by naming their parameters alike.
One quantizes the matrix coordinate algebra; the other asks
for matrix representations of a quantized coefficient algebra.
A comparison between them must specify a map and its effect
on relation classes.

\section{Framed and infinite-dimensional recipients}
\label{sec:vol2-weyl-action-recipients}

Framing changes the finite-dimensional equation. Add a column
$\mathbf i\in\mathbb C^N$ and a row
$\mathbf j\in(\mathbb C^N)^*$, and use
\begin{equation}\label{eq:vol2-framed-moment-equation}
 [X,Y]+\mathbf i\mathbf j=\varepsilon I_N.
\end{equation}
Simultaneous conjugation extends by
$\mathbf i\mapsto g\mathbf i$ and
$\mathbf j\mapsto\mathbf jg^{-1}$. The extra rank-one
term permits a nonzero trace of the prescribed moment value.

\begin{proposition}\label{prop:vol2-framed-weyl-chart}
On the ordered chart with pairwise distinct $x_i$, the formulas
\begin{equation}\label{eq:vol2-framed-weyl-coordinates}
 \begin{gathered}
 X=\operatorname{diag}(x_i),\qquad
 \mathbf i_i=1,\qquad\mathbf j_i=\varepsilon,\\
 Y_{ii}=y_i,\qquad
 Y_{ij}=-\frac{\varepsilon}{x_i-x_j}\quad(i\ne j)
 \end{gathered}
\end{equation}
give a free framed orbit and a smooth local quotient of
relative dimension $2N$ over the parameter. At fixed
$\varepsilon$, the reduced symplectic form is
$\sum_i dx_i\wedge dy_i$ on this ordered cover.
\end{proposition}
\begin{proof}
The off-diagonal entries of the left side of
\eqref{eq:vol2-framed-moment-equation} are
$-\varepsilon+\varepsilon=0$; its diagonal entries are
$\varepsilon$. A stabilizer commuting with $X$ is diagonal,
and fixing the column with all entries one makes it the
identity. The infinitesimal action is injective.

On the open set with simple $X$ and nonzero components of
$\mathbf i$ in its eigenbasis, diagonalize $X$ and use diagonal
gauge transformations to set those components equal to one.
The moment equation then forces the displayed values of
$\mathbf j$ and the off-diagonal entries of $Y$.
Thus $(x_i,y_i)$ actually parametrize the local quotient.

With the canonical form
$\operatorname{Tr}(dX\wedge dY)
 +\sum_i d\mathbf i_i\wedge d\mathbf j_i$,
the moment differential is the symplectic transpose of the
infinitesimal action, so it is surjective. The quotient is
therefore smooth of relative dimension
$(2N^2+2N)-2N^2=2N$ near this orbit.
On the slice, $d\mathbf i=d\mathbf j=0$ at fixed parameter,
and $dX$ is diagonal. The pullback of the form is exactly
$\sum_i dx_i\wedge dy_i$. Permuting the eigenvalues gives
the remaining finite identification.
\end{proof}

An infinite-dimensional unframed recipient keeps the original
coefficient equation. Complete $W_\varepsilon$ in the
$\varepsilon$-adic topology and write it as
$\widehat W_\varepsilon$. Use complete left modules with
continuous maps and strict exact sequences, resolved by
topologically free modules. The completed polynomial Fock module
\[
 \mathcal P=\mathbb C[p][[\varepsilon]],\qquad
 q\cdot f=\varepsilon\partial_pf,\qquad p\cdot f=pf
\]
satisfies $[q,p]=\varepsilon$. A coefficient at each parameter
order is a polynomial in $p$, with no uniform degree bound.

\begin{proposition}\label{prop:vol2-integral-fock-self-extension}
Continuous derived Hom in this exact category gives
\[
 \operatorname{Ext}^0_{\widehat W_\varepsilon}(\mathcal P,\mathcal P)
     =\mathbb C[[\varepsilon]],
 \qquad
 \operatorname{Ext}^1_{\widehat W_\varepsilon}(\mathcal P,\mathcal P)
     =\mathbb C[p],
\]
with all other degrees zero. Multiplication by $\varepsilon$
annihilates the degree-one group.
\end{proposition}
\begin{proof}
In the completed PBW basis $p^bq^a$, right multiplication by
$q$ increases $a$ by one. It is injective with closed image,
and its cokernel is $\mathcal P$. As a map of left modules
it gives the strict topologically free resolution
\[
 0\longrightarrow\widehat W_\varepsilon
   \xrightarrow{R_q}\widehat W_\varepsilon
   \longrightarrow\mathcal P\longrightarrow0.
\]
Continuous Hom into $\mathcal P$ is the two-term complex
\[
 \mathcal P\xrightarrow{\varepsilon\partial_p}\mathcal P
\]
in degrees zero and one. Coefficientwise polynomial
integration is a continuous right inverse to $\partial_p$.
The kernel is $\mathbb C[[\varepsilon]]$, the image is
$\varepsilon\mathcal P$, and the cokernel is
$\mathcal P/\varepsilon\mathcal P=\mathbb C[p]$.
\end{proof}

The degree-one self-extension disappears when the parameter
is inverted. Its integral presence records information that
generic specialization loses. The framed chart and the Fock
module thus avoid the finite unframed contraction in different
ways: the first changes the equation, and the second changes
the dimension and completion of the module.

\section{Hamiltonian and geometric comparison maps}
\label{sec:vol2-derived-hamiltonian-comparison}

The construction has produced an action on an associative
resolution, an action on derived matrix representations,
and a coherent classical negative-cyclic orientation.
The quotient-by-constants obstruction fixes the source of
the stronger Hamiltonian problem. Retain either the full
classical Poisson algebra $A$ with $J_N(1)=N$, or its pointed
Lie subalgebra, and keep the scalar quadratic map fixed.

\begin{problem}\label{prob:vol2-derived-hamiltonian-moment-map}
Construct a homotopy moment map to the full classical BFV
algebra whose first component is $J_N$, whose cubic binary
component is $-\tfrac32\operatorname{Tr}(\Psi[X,Y])$ up to
specified homotopies, and whose choices respect block-diagonal
sums of matrices. Compare its Hamiltonian derivation action
with Theorem~\ref{thm:vol2-all-arity-derived-action}, retaining
the cyclic orientation.

For matrix quantization, construct a filtered derived
comparison on the full quantum BRST complex that agrees
with the ordinary quantum trace action in degree zero and
preserves the relation classes, stabilizer states and cyclic
operations. Require the prescribed classical limit and cubic
primitive, and extend the comparison across the matrix
discriminant. Determine whether the coefficient-Weyl action
compares to this recipient, to a framed recipient, or to an
infinite-dimensional module by maps preserving those data.
\end{problem}

The positive gauge ghosts in
\eqref{eq:vol2-classical-bfv-differential} already belong to
the action constructed here. Supplying them is not the
missing Hamiltonian comparison. Likewise a contraction of
one finite coefficient representation does not obstruct
every framed, completed or quantum-reduced recipient.

For mixed holomorphic--topological theory there is a further
geometric map to construct. The polynomial cyclic operations
must be realized by supported operations of the independently
defined field theory. Such a comparison must preserve the
boundary polarization, parameter normalization, descendant
operations and sewing. At four inputs it must include the
simultaneous collision support
$\Delta_{13}\cap\Delta_{24}$ from
Chapter~\ref{ch:vol2-supported-compositions}.
The universal word contraction solves the algebraic action
equations; the geometric comparison must carry their
solutions to these specified supports and operations.

\chapter{Relative quantum reduction across the discriminant}
\label{chap:vol2-relative-quantum-reduction}

Diagonalizing commuting matrices suggests that their invariant functions
should be functions of eigenvalue pairs. At repeated eigenvalues,
diagonalization no longer supplies a coordinate chart. The relation complex
also has odd functions that an ordinary quotient cannot see. A comparison
across this locus must therefore be a map of complexes, with its product
and coefficient ring specified.

In rank two the odd functions form a particularly small module. It has
rank one away from a quadratic cone's vertex, but needs two generators at
the vertex. We first compute this module and its restriction map.
A terminating differential operator then deforms the restriction over
$\mathbb C[\hbar]$. The result is a strict quasi-isomorphism, although it
has no strict algebra section. Its higher inverse makes this distinction
explicit. The final comparison with ordinary matrix reduction retains
the quantum parameter and trace rank at every matrix size.

\section{The relation module at a singular pair}
\label{sec:vol2-rank-two-koszul}

Work over $\mathbb C$ with cohomological differentials of degree one.
Identify the adjoint representation of $\mathfrak{sl}_2$ with the
three-dimensional representation of $G=SO_3(\mathbb C)$.
Normalize its invariant form and the equation generators so that the
commuting equation is the cross product. For even vectors $x,y$ and an
odd vector $\psi$, put
\[
 P=\mathbb C[x_1,x_2,x_3,y_1,y_2,y_3],
 \qquad
 K=P\otimes\Lambda(\psi_1,\psi_2,\psi_3),
 \qquad |\psi_i|=-1,
\]
\begin{equation}\label{eq:vol2-rank-two-koszul}
 dx=dy=0,\qquad d\psi=\mu,\qquad \mu=x\times y.
\end{equation}
The orientation is $\epsilon_{123}=1$. Exterior products are written
in the order $\psi_1,\psi_2,\psi_3$, and odd derivatives act on the left.
The relative complex is $K^G$, the invariants of the rational group
action. Taking these invariants imposes no positive gauge ghosts.

Define
\[
 a=x\cdot x,\qquad b=x\cdot y,\qquad c=y\cdot y,
 \qquad S=\mathbb C[a,b,c],\qquad \Delta=ac-b^2,
\]
and
\[
 \xi=x\cdot\psi,\qquad \eta=y\cdot\psi,\qquad z=\mu\cdot\psi.
\]
The elements $\xi,\eta$ are closed because $\mu$ is orthogonal to both
$x$ and $y$. To exhibit their relations, introduce
\[
 \begin{aligned}
 p&=x_1\psi_2\psi_3-x_2\psi_1\psi_3+x_3\psi_1\psi_2,\\
 q&=y_1\psi_2\psi_3-y_2\psi_1\psi_3+y_3\psi_1\psi_2,\\
 w&=\mu_1\psi_2\psi_3-\mu_2\psi_1\psi_3+\mu_3\psi_1\psi_2,
 \qquad r=\psi_1\psi_2\psi_3.
 \end{aligned}
\]
Here $p,q$ denote elements of exterior degree two, independently of the
plane coordinates used in Chapter~\ref{ch:vol2-derived-hamiltonian-actions}.
Their cohomological degrees are $-2$.

\begin{lemma}\label{lem:vol2-rank-two-primitives}
The differential satisfies
\begin{equation}\label{eq:vol2-rank-two-primitives}
 dz=\Delta,\qquad
 dp=b\xi-a\eta,\qquad dq=c\xi-b\eta,\qquad
 dr=w=\xi\eta,\qquad dw=0.
\end{equation}
\end{lemma}

\begin{proof}
Lagrange's identity gives $\mu\cdot\mu=ac-b^2$.
Identify a vector $v$ with the two-form
$v_1\psi_2\psi_3-v_2\psi_1\psi_3+v_3\psi_1\psi_2$.
Its differential is $(v\times\mu)\cdot\psi$.
The vector identities
\[
 x\times(x\times y)=bx-ay,\qquad
 y\times(x\times y)=cx-by
\]
give the middle two formulas. Expanding $d(\psi_1\psi_2\psi_3)$ gives
$w$, and expanding $(x\cdot\psi)(y\cdot\psi)$ gives the same two-form.
The last equation follows from $d^2r=0$.
\end{proof}

The invariant complex has eight coefficient generators. A polynomial
vector covariant is a polynomial map $v(x,y)$ satisfying
$v(gx,gy)=gv(x,y)$ for $g\in G$. The following calculation proves
the completeness of these generators, including on $\Delta=0$.

\begin{lemma}\label{lem:vol2-rank-two-covariants}
The scalar invariant ring is $S$. The polynomial vector covariants are
the free $S$-module generated by $x,y,\mu$. Consequently
\begin{equation}\label{eq:vol2-rank-two-components}
 \begin{aligned}
 (K^G)^0&=S,&
 (K^G)^{-1}&=S\xi\oplus S\eta\oplus Sz,\\
 (K^G)^{-2}&=Sp\oplus Sq\oplus Sw,&
 (K^G)^{-3}&=Sr.
 \end{aligned}
\end{equation}
\end{lemma}

\begin{proof}
Give a polynomial its separate degrees in $x$ and $y$.
The map $(x,y)\mapsto(a,b,c)$ is dominant: every nondegenerate
two-dimensional symmetric form embeds in the standard form on
$\mathbb C^3$. Thus $a,b,c$ are algebraically independent.
The determinant of the three columns $x,y,\mu$ is $\Delta$, a nonzero
polynomial. Hence these columns are linearly independent over the
fraction field and, in particular, over $S$.

It remains to prove that these explicit inclusions exhaust the
invariants. Let $V_\ell$ be the irreducible $SO_3$-representation of
spin $\ell$, of dimension $2\ell+1$. The elementary
$\mathfrak{sl}_2$ character formulas give
\[
 \operatorname{Sym}^m(V_1)
   =\bigoplus_{j=0}^{\lfloor m/2\rfloor}V_{m-2j},
 \qquad
 V_\ell\otimes V_k
   =\bigoplus_{n=|\ell-k|}^{\ell+k}V_n.
\]
For example, the first formula follows by comparing the coefficients
of $u^m$ in
\[
 \frac1{(1-uz)(1-u)(1-uz^{-1})}
 =\frac1{1-u^2}\sum_{\ell\geq0}u^\ell
     (z^\ell+z^{\ell-1}+\cdots+z^{-\ell}).
\]
The second follows by multiplying the two finite weight sums and
subtracting their successive highest weights.

The multiplicity of $V_0$ in $V_\ell\otimes V_k$ is one precisely
when $\ell=k$. The multiplicity of $V_1$ is one precisely when
$|\ell-k|\leq1\leq\ell+k$. Their bivariate generating series in the
polynomial ring are therefore
\[
 \frac1{(1-u^2)(1-uv)(1-v^2)},\qquad
 \frac{u+v+uv}{(1-u^2)(1-uv)(1-v^2)}.
\]
These are exactly the series of the already embedded ring $S$ and
module $Sx\oplus Sy\oplus S\mu$. Each bidegree is finite-dimensional,
so equality of dimensions makes those inclusions surjective.
The invariant form identifies vectors and covectors.
The invariant volume identifies their second exterior powers, and
the top exterior power is trivial. This gives
\eqref{eq:vol2-rank-two-components}.
\end{proof}

Set
\[
 B=S/(\Delta),\qquad
 M=(B\xi\oplus B\eta)/(b\xi-a\eta,\ c\xi-b\eta),
 \qquad |\xi|=|\eta|=-1.
\]
Let the involution on
$\mathbb C[s,t]\otimes\Lambda(\theta)$ change the signs of all three
generators, with $|\theta|=-1$. Restriction to a Cartan axis defines
\begin{equation}\label{eq:vol2-classical-rank-two-restriction}
 \rho:K^G\longrightarrow
     \bigl(\mathbb C[s,t]\otimes\Lambda(\theta)\bigr)^{\mathbb Z/2},
 \qquad
 x=(s,0,0),\quad y=(t,0,0),\quad\psi=(\theta,0,0).
\end{equation}
The target has zero differential.

\begin{theorem}\label{thm:vol2-classical-rank-two-restriction}
The map $\rho$ is a surjective quasi-isomorphism of commutative
differential graded algebras. It induces
\[
 H^0(K^G)=B,\qquad H^{-1}(K^G)=M,\qquad
 H^j(K^G)=0\quad(j\leq-2),
\]
with all products of negative-degree classes zero.
It sends $(a,b,c,\xi,\eta)$ to
$(s^2,st,t^2,s\theta,t\theta)$.
\end{theorem}

This is the traceless rank-two instance of the derived
Harish--Chandra comparison of Berest, Felder, Patotski, Ramadoss and
Willwacher (2015), Theorem~4.1 of~\cite{BFPRW2015}.
Here the coefficient presentation gives a direct proof.

\begin{proof}
The map on $Sz$ is multiplication by $\Delta$, hence injective.
On $Sp\oplus Sq$ the differential has matrix
\[
 \begin{pmatrix}b&c\\-a&-b\end{pmatrix}
\]
with determinant $\Delta$. It too is injective over the polynomial
domain $S$. Finally $Sr\to Sw$ is an isomorphism.
Thus the displayed four-term complex has no cohomology below
degree $-1$. Its degree-zero quotient is $B$, and its odd quotient is
$S^2/\langle(b,-a),(c,-b)\rangle$.
The two column relations imply $\Delta\xi=\Delta\eta=0$, so this
quotient is precisely $M$. The primitive $dr=\xi\eta$ proves
the product assertion.

The target's even part is
$\mathbb C[s^2,st,t^2]\cong B$. Its odd part consists of odd
polynomials in $s,t$, multiplied by $\theta$. It is generated by
$s\theta,t\theta$ over the even ring.
If even polynomials $F,G$ satisfy $Fs+Gt=0$, coprimality gives
$F=tH$ and $G=-sH$ with $H$ odd. Write $H=sP+tQ$ with $P,Q$
even. The relation is then
$P(st,-s^2)+Q(t^2,-st)$, exactly the relation module above.
Restriction is surjective on the displayed generators and induces
these cohomology isomorphisms.
\end{proof}

On $a\ne0$, the odd module is generated by $\xi$, since
$\eta=(b/a)\xi$. The same conclusion holds on $c\ne0$ with the
roles reversed. These two open sets cover the punctured cone.
At its vertex,
\begin{equation}\label{eq:vol2-rank-two-odd-vertex}
 M/(a,b,c)M=\mathbb C\xi\oplus\mathbb C\eta.
\end{equation}
This is the ordinary fibre of a cohomology module. It is not an
assertion that cohomology commutes with derived base change.
Replacing $M$ by a free line would already give the wrong ordinary
fibre. The restriction theorem uses no discriminant localization.

For $\mathfrak{gl}_2$, the scalar matrix coordinates supply a polynomial
factor $\mathbb C[u,v]$. The trace equation is identically zero,
so its relation generator $\zeta$, of degree $-1$, is closed.
Tensoring $\rho$ with
$\mathbb C[u,v]\otimes\Lambda(\zeta)$ gives the full rank-two
invariant-function comparison. Its degree $-2$ includes
$\mathbb C[u,v]\zeta M$. This extra odd direction cannot be discarded.
The statement concerns invariant functions, rather than an equivalence
of quotient stacks or of their module categories.

\section{The quantum differential and its coefficient operator}
\label{sec:vol2-relative-quantum-differential}

Use the matrix quantum parameter $\hbar$, independently of the
coefficient-plane parameter $\varepsilon$ in
Section~\ref{sec:vol2-coefficient-weyl-action}.
On polynomial symbols in the three canonical pairs $(x_i,y_i)$ set
\begin{equation}\label{eq:vol2-relative-moyal}
 F\star G
 =m\exp\left[
   \frac{\hbar}{2}\sum_{i=1}^3
   \left(\partial_{x_i}\otimes\partial_{y_i}
          -\partial_{y_i}\otimes\partial_{x_i}\right)
       \right](F\otimes G).
\end{equation}
Here $m$ is ordinary multiplication. Each product is polynomial in
$\hbar$. The $\psi_i$ still multiply
exteriorly. Define
\[
 \mathcal K_\hbar
   =\bigl(\mathbb C[x,y]\otimes\Lambda(\psi_1,\psi_2,\psi_3)
      \bigr)^G\otimes\mathbb C[\hbar]
\]
with this product and
\begin{equation}\label{eq:vol2-relative-quantum-differential}
 \delta_\hbar
 =\sum_i\mu_i\partial_{\psi_i}
  +\frac{\hbar^2}{4}
    \sum_{i,j,k}\epsilon_{ijk}
       \partial_{x_j}\partial_{y_k}\partial_{\psi_i}.
\end{equation}
The first summand uses ordinary multiplication of symbols.
The second summand is needed even though the moment functions
$\mu_i$ themselves are unchanged.

\begin{proposition}\label{prop:vol2-relative-quantum-dg}
The operator $\delta_\hbar$ is a square-zero derivation of
$\mathcal K_\hbar$ over $\mathbb C[\hbar]$.
\end{proposition}

\begin{proof}
Adjoin ghosts $\gamma^i$ of degree $1$ with
$\psi_i\gamma^j+\gamma^j\psi_i=\hbar\delta_i^j$.
The BRST element for the rotation moment is
\[
 \Theta=\sum_i\gamma^i\mu_i
  -\frac12\sum_{i,j,k}\epsilon_{ijk}\gamma^i\gamma^j\psi_k.
\]
The moment commutators are
$[\mu_i,\mu_j]_\star=\hbar\sum_k\epsilon_{ijk}\mu_k$.
Together with Jacobi they give $\Theta^2=0$, by the
Weyl--Clifford construction of
Section~\ref{sec:vol2-two-quantizations}.
Its polynomial inner derivation $[\Theta,-]/\hbar$ is therefore
square-zero.

Use symmetric symbols for the ghosts. On an expression independent
of $\gamma$, the terms that still contain $\gamma^i$ are its
infinitesimal simultaneous rotations, including the rotation of
$\psi$. They vanish on $G$-invariants. The remaining term is
\[
 \frac12\sum_i\left(
      \mu_i\star\partial_{\psi_i}F
        +\partial_{\psi_i}F\star\mu_i\right).
\]
Since $\mu_i$ is quadratic, the Moyal expansion stops at its second
derivative. The first derivatives cancel in the average, while
the second derivatives give exactly
\eqref{eq:vol2-relative-quantum-differential}.
The cubic ghost term has no higher contraction on a
$\gamma$-independent input. Thus this invariant subalgebra is
preserved by the inner derivation, proving both assertions.
\end{proof}

The presentation \eqref{eq:vol2-rank-two-components} remains a
presentation of polynomial symbols over $\mathbb C[\hbar]$.
It is a statement about coefficient modules, not about ordinary
multiplication being the quantum product. In particular
$\xi\star\eta=w$: the possible first contraction is proportional
to $\sum_i\psi_i^2$, hence vanishes.

Define operators on $S[\hbar]$ by
\begin{equation}\label{eq:vol2-rank-two-coefficient-operator}
 L=\partial_a\partial_c-\frac14\partial_b^2,\qquad
 T=1+\hbar^2L,\qquad
 T^{-1}F=\sum_{j\geq0}(-\hbar^2L)^jF.
\end{equation}
The inverse sum terminates because $L$ lowers polynomial degree
in $a,b,c$ by two. Neither $\hbar$ nor $\Delta$ is inverted.

\begin{lemma}\label{lem:vol2-rank-two-quantum-coefficients}
For every $F\in S[\hbar]$,
\begin{equation}\label{eq:vol2-rank-two-quantum-coefficients}
 \begin{aligned}
 \delta_\hbar(Fz)&=T(\Delta F),\\
 \delta_\hbar(Fp)&=T(bF)\xi-T(aF)\eta,\\
 \delta_\hbar(Fq)&=T(cF)\xi-T(bF)\eta,\\
 \delta_\hbar(Fr)&=T(F)w,\\
 \delta_\hbar(F\xi)&=\delta_\hbar(F\eta)
                  =\delta_\hbar(Fw)=0.
 \end{aligned}
\end{equation}
All products inside the arguments of $T$ are ordinary coefficient
products.
\end{lemma}

\begin{proof}
The chain rule gives
$\partial_{x_i}F=2x_iF_a+y_iF_b$ and
$\partial_{y_i}F=x_iF_b+2y_iF_c$. Therefore
\[
 \sum_{j,k}\epsilon_{ijk}
       \partial_{x_j}\partial_{y_k}F=4\mu_i LF.
\]
Applying the derivatives also to $\mu_i$ gives
\[
 \frac14\sum_{i,j,k}\epsilon_{ijk}
    \partial_{x_j}\partial_{y_k}(F\mu_i)
  =\Delta LF+(a\partial_a+b\partial_b+c\partial_c)F
       +\frac32F
  =L(\Delta F).
\]
This proves the first equation.
On $Fp$, the quantum coefficients of $\xi,\eta$ are
$bLF-F_b/2$ and $-aLF-F_c$.
On $Fq$, they are $cLF+F_a$ and $-bLF+F_b/2$.
These equal $L(bF),-L(aF),L(cF),-L(bF)$, respectively.
The top exterior differential has quantum correction $LF\,w$.
Antisymmetric contraction with $x$ or $y$ kills the corrections
on $F\xi,F\eta$.
Finally $F\mapsto T(F)$ is surjective, so every $Fw$ is the
differential of $T^{-1}(F)r$. Square-zero then gives its stated
differential.
\end{proof}

Let $U$ apply $T^{-1}$ to coefficients in
$S,S\xi,S\eta,Sw$ and act as the identity on
$Sz,Sp,Sq,Sr$. The coefficient formulas give
\begin{equation}\label{eq:vol2-relative-linear-comparison}
 U\delta_\hbar=\delta_0U.
\end{equation}
Thus $U$ is an integral chain isomorphism. Its inverse uses $T$
on the same four summands. This identity already computes the
underlying cohomology modules. To identify their product, the
restriction itself must be multiplicative.

\section{Strict restriction over the quantum ring}
\label{sec:vol2-strict-relative-restriction}

Let $W_\hbar(s,t)$ be the one-pair Weyl algebra in symmetric
symbols, so that $[s,t]_\star=\hbar$. Put
\[
 \mathcal H_\hbar
   =\bigl(W_\hbar(s,t)\otimes\Lambda(\theta)\bigr)^{\mathbb Z/2},
 \qquad |\theta|=-1,
\]
with zero differential and the simultaneous sign involution.
For a scalar symbol write $\rho_0F=F(s^2,st,t^2)$, where
$st$ in this substitution is the commuting symbol.
Define
\begin{equation}\label{eq:vol2-quantum-rank-two-restriction}
 \begin{aligned}
 \rho_\hbar(F)&=\rho_0(T^{-1}F),\\
 \rho_\hbar(F\xi+G\eta+Jz)
   &=\bigl(s\rho_0(T^{-1}F)+t\rho_0(T^{-1}G)\bigr)\theta,\\
 \rho_\hbar(\mathcal K_\hbar^{\leq-2})&=0.
 \end{aligned}
\end{equation}

\begin{theorem}\label{thm:vol2-quantum-rank-two-restriction}
The map $\rho_\hbar:\mathcal K_\hbar\to\mathcal H_\hbar$ is
a surjective strict quasi-isomorphism of differential graded
algebras over $\mathbb C[\hbar]$. Its kernel has a
$\mathbb C[\hbar]$-linear contraction. Consequently every base
change to a commutative $\mathbb C[\hbar]$-algebra preserves the quasi-isomorphism,
including every scalar specialization of $\hbar$.
\end{theorem}

\begin{proof}
As a chain map, $\rho_\hbar=\rho U$. The classical theorem
therefore proves surjectivity and the cohomology comparison.
We prove multiplicativity by the operators of multiplication
by the three invariant quadratics.

In $d$ canonical pairs, left multiplication on invariant scalar
symbols has the following operators:
\begin{align*}
 A_d={}&a+\hbar(a\partial_b+2b\partial_c)
   +\frac{\hbar^2}{4}
      (a\partial_b^2+4b\partial_b\partial_c
                    +4c\partial_c^2+2d\partial_c),\\
 B_d={}&b+\hbar(c\partial_c-a\partial_a)
   -\frac{\hbar^2}{4}
      (2a\partial_a\partial_b+b\partial_b^2
         +4b\partial_a\partial_c+2c\partial_b\partial_c
         +d\partial_b),\\
 C_d={}&c-\hbar(2b\partial_a+c\partial_b)
   +\frac{\hbar^2}{4}
      (4a\partial_a^2+4b\partial_a\partial_b
                       +c\partial_b^2+2d\partial_a).
\end{align*}
They follow by applying the two derivatives in the Moyal
product to $a,b,c$. For instance,
$a\star F=aF+\hbar x\cdot\partial_yF
 +\hbar^2\sum_i\partial_{y_i}^2F/4$,
and the chain rule gives $A_d$.

These differential operators satisfy
\begin{equation}\label{eq:vol2-rank-two-radial-intertwiners}
 A_3T=TA_1,\qquad B_3T=TB_1,\qquad C_3T=TC_1.
\end{equation}
Here is the coefficient verification. The commutators of $L$
with multiplication by $a,b,c$ are
$\partial_c,-\partial_b/2,\partial_a$.
It commutes with the three first-order coefficients of $\hbar$.
Its commutators with the coefficients of $\hbar^2$ in
$A_1,B_1,C_1$ are, respectively,
$L\partial_c,-L\partial_b/2,L\partial_a$.
Thus $[T,A_1]=\hbar^2\partial_cT=(A_3-A_1)T$,
and the other two equations follow with the displayed constants.

Evaluation by $\rho_0$ changes $A_1,B_1,C_1$ into actual
one-pair left multiplication by $s^2,st,t^2$.
Equation~\eqref{eq:vol2-rank-two-radial-intertwiners}
therefore proves that $\rho_\hbar$ intertwines multiplication
by $a,b,c$ on the scalar sector.
These three quadratics generate its invariant Weyl algebra:
a star monomial has its ordinary monomial as leading symbol,
and subtraction lowers polynomial degree. Induction proves
generation over $\mathbb C[\hbar]$.

On the coefficient column $(F,G)$ of $F\xi+G\eta$,
the corresponding left multiplication operators are
\[
 \begin{pmatrix}
 A_d&\hbar+\frac{\hbar^2}{2}\partial_b\\
 0&A_d+\hbar^2\partial_c
 \end{pmatrix},
 \qquad
 \begin{pmatrix}
 B_d-\frac\hbar2-\frac{\hbar^2}{4}\partial_b
    &-\frac{\hbar^2}{2}\partial_a\\
 -\frac{\hbar^2}{2}\partial_c
    &B_d+\frac\hbar2-\frac{\hbar^2}{4}\partial_b
 \end{pmatrix},
\]
\[
 \begin{pmatrix}
 C_d+\hbar^2\partial_a&0\\
 -\hbar+\frac{\hbar^2}{2}\partial_b&C_d
 \end{pmatrix}.
\]
For example, multiplying the coefficient $G$ and the single
$y$ in $\eta$ accounts for the two off-diagonal terms of the
first matrix. The other entries follow from the same product
rule. Their extra entries have constant coefficients, so they
commute with $T$. The scalar intertwiners prove the column
intertwiners. Replacing $\hbar$ by $-\hbar$ gives right
multiplication as well.

Give $x,y$ weight one and $\psi,\hbar$ weight two.
The scalar and $Sz$ sectors have even weight, while
$S\xi,S\eta$ have odd weight. Scalar star multiplication
therefore preserves the killed $Sz$ sector.
Products of total cohomological degree at most $-2$ have
zero image, and two odd target factors have product zero
because $\theta^2=0$. This exhausts the products and proves
strict multiplicativity.

Choose a contraction of the classical kernel in each finite
weight. Such a contraction exists by the classical
quasi-isomorphism: split each vector space into boundaries
and chosen preimages, and send boundaries back to their
preimages. Extend it $\mathbb C[\hbar]$-linearly and conjugate
by $U$. This contracts the quantum kernel.
Together with a chain section of the surjection, it gives a
contraction of the whole comparison. Tensoring these
identities with any commutative coefficient algebra proves the base-change
assertion without a flatness assumption.
\end{proof}

The correction is already visible in three small values:
\begin{equation}\label{eq:vol2-relative-first-quantum-corrections}
 \rho_\hbar(ac)=s^2t^2-\hbar^2,\qquad
 \rho_\hbar(b^2)=s^2t^2+\frac12\hbar^2,\qquad
 \delta_\hbar z=\Delta+\frac32\hbar^2.
\end{equation}
The naive substitution $\rho_0$ would send the last boundary
to $3\hbar^2/2$, so it is not the quantum chain map.
The strict comparison instead identifies the even cohomology
relation
\[
 \frac12(a\star c+c\star a)-b\star b=\frac34\hbar^2.
\]
Its odd relations retain their primitives:
\[
 (b+\tfrac32\hbar)\star\xi-a\star\eta=\delta_\hbar p,
 \qquad
 c\star\xi-(b-\tfrac32\hbar)\star\eta=\delta_\hbar q.
\]
Tensoring with the independent scalar Weyl pair and its closed
central antighost gives the relative $\mathfrak{gl}_2$ comparison.
The absolute complex still has its positive ghosts.

\section{An inverse with higher components}
\label{sec:vol2-relative-higher-inverse}

A strict projection need not have a strict section. To construct
an inverse of $\rho_\hbar$ with its homotopies, choose the classical
section on the Cartan monomial basis by
\[
 i_0(s^mt^n)=
 \begin{cases}
  a^{m/2}c^{n/2},&m,n\text{ even},\\
  a^{(m-1)/2}bc^{(n-1)/2},&m,n\text{ odd},
 \end{cases}
\]
\[
 i_0(s^mt^n\theta)=
 \begin{cases}
  a^{(m-1)/2}c^{n/2}\xi,&m\text{ odd},\ n\text{ even},\\
  a^{m/2}c^{(n-1)/2}\eta,&m\text{ even},\ n\text{ odd}.
 \end{cases}
\]
In each finite weight split the classical complex into these
cohomology representatives, boundaries and chosen preimages of
boundaries. Let $H_0$ send a boundary to its preimage and vanish on
the other two summands. Choose $H_0\Delta=z$ and $H_0w=r$.
Then
\[
 \delta_0H_0+H_0\delta_0=1-i_0\rho,\qquad
 H_0^2=H_0i_0=\rho H_0=0.
\]
Extend these maps $\mathbb C[\hbar]$-linearly and put
\[
 i_\hbar=U^{-1}i_0,\qquad
 H_\hbar=U^{-1}H_0U.
\]
They give the same normalized contraction identities for
$(\rho_\hbar,i_\hbar,H_\hbar)$.

Let $s$ lower cohomological degree by one.
On the tensor coalgebras of the suspended underlying nonunital
algebras use the codifferentials determined by
\[
 b_1(sa)=-s\delta_\hbar a,\qquad
 b_2(sa,sb)=(-1)^{|a|}s(a\star b).
\]
The target algebra has only $b_2$.
These coalgebras have the deconcatenation coproduct; their maps
are determined by their components into tensors of length one.
No scalar augmentation of a Weyl algebra is required.
Put $H_s=-sH_\hbar s^{-1}$ and $I_1=si_\hbar s^{-1}$.
Given the components below arity $n$, set
\begin{equation}\label{eq:vol2-relative-inverse-recursion}
 O_n=(b_{\mathcal K}I_{<n}-I_{<n}b_{\mathcal H})_n,
 \qquad I_n=-H_sO_n.
\end{equation}
Here $I_{<n}$ denotes the coalgebra map with the prescribed
lower components and zero arity-$n$ component.

\begin{theorem}\label{thm:vol2-relative-higher-inverse}
This recursion defines a strictly unital $A_\infty$ morphism
$I:\mathcal H_\hbar\to\mathcal K_\hbar$ whose first component
is a quasi-isomorphism and for which $\rho_\hbar I$ is the
strict identity. Its only possible components are $I_1,I_2,I_4$.
For the commuting target symbol $B=st$,
\begin{equation}\label{eq:vol2-relative-fourth-inverse}
 I_2(sB,sB)=-sz,\qquad
 I_4(sB,sB,sB,sB)=-\hbar sr.
\end{equation}
\end{theorem}

\begin{proof}
The equation $b^2=0$, together with the lower morphism equations,
makes $O_n$ closed for the unary differential.
Strictness of $\rho_\hbar$ and
$\rho_\hbar i_\hbar=1$, $\rho_\hbar H_\hbar=0$
put $O_n$ in the kernel of the suspended projection.
The contraction identity makes $b_1(-H_sO_n)=-O_n$.
Thus \eqref{eq:vol2-relative-inverse-recursion} solves the next
morphism equation. Every map is polynomial on its input.

The special vanishing follows from the weights and the four
cohomological degrees. On inputs with $k$ odd factors, the
unsuspended third residual has degree $-1-k$ and weight parity
$k$. For $k=0,1,2$ it lies in $S[\hbar]z$,
$S[\hbar]p\oplus S[\hbar]q$, and $S[\hbar]r$,
respectively. The differential is injective on each of these
sectors, by the coefficient formulas and the determinant
$\Delta$ of the middle matrix. Closedness forces the residual
to vanish. For larger $k$ its degree is absent. Hence $I_3=0$.

The fourth component has unsuspended degree $-3-k$, so it
can occur only on four even inputs and takes values in
$S[\hbar]r$. All components of arity at least five would
have degree below $-3$. The fifth residual on even inputs
lies in $S[\hbar]r$ and is closed, hence zero by injectivity.
Every later residual has absent degree. This proves all
morphism equations, including those beyond the last nonzero
component. Since the section preserves the unit and the
homotopy annihilates its image, the normalized recursion
vanishes on a unit input in every higher arity.

For the stated example,
\[
 b\star b=b^2-\frac34\hbar^2,\qquad
 i_\hbar(B\star B)=ac+\frac34\hbar^2.
\]
Their difference is $-\delta_\hbar z$, giving $I_2=-sz$.
The moment commutator and exterior signs give
\[
 z\star z=\hbar w,\qquad \delta_\hbar r=w.
\]
Consequently the fourth residual is
$b_2(-sz,-sz)=-\hbar sw$, and the recursion gives
$I_4=-\hbar sr$.
\end{proof}

For even inputs $F,G$, the binary component is particularly
explicit. Set
\[
 D=i_\hbar(F)\star i_\hbar(G)-i_\hbar(F\star G).
\]
Then $T^{-1}D$ is divisible by $\Delta$ in the polynomial
ring. The unsuspended binary value is
$(T^{-1}D/\Delta)z$. This is exact polynomial division,
not localization at the discriminant.
The fourth value is obtained by multiplying two binary
values and applying $-T^{-1}$ to the resulting $w$ coefficient.
The target still has its ordinary binary product and no
higher multiplications.

\section{Changing a section}
\label{sec:vol2-relative-sections}

The nonzero fourth value belongs to a chosen inverse.
There are two separate questions: whether the projection has
a strict section on the whole algebra, and whether a
particular higher value survives a change of section.

\begin{proposition}\label{prop:vol2-relative-no-strict-section}
The map $\rho_\hbar$ has no $\mathbb C[\hbar]$-linear
dg-algebra section.
\end{proposition}

\begin{proof}
A section would specialize at $\hbar=0$ to an algebra
section of $S\to S/(\Delta)$ in degree zero.
Its lifts of $a,b,c$ would have the form
$A=a+\Delta\alpha$, $B=b+\Delta\beta$,
$C=c+\Delta\gamma$. The relation $AC=B^2$ would imply
\[
 0=\Delta\left(
  1+a\gamma+c\alpha-2b\beta
       +\Delta(\alpha\gamma-\beta^2)\right).
\]
Since $S$ is a domain, the parenthesis would be zero.
Its constant term is one, a contradiction.
\end{proof}

\begin{proposition}\label{prop:vol2-relative-polynomial-section}
There is a unital weight-preserving chain section $i'$ of
$\rho_\hbar$ whose restriction to the polynomial star-subalgebra
$\mathbb C[\hbar][B]$, $B=st$, is the algebra map
\[
 i'(B^{\star n})=b^{\star n}.
\]
A normalized contraction with this section has a recursive
$A_\infty$ inverse whose higher components all vanish when
every input belongs to this subalgebra.
\end{proposition}

\begin{proof}
The leading symbols $(st)^n$ show that the star powers
$B^{\star n}$ are linearly independent over $\mathbb C[\hbar]$.
Multiplicativity of $\rho_\hbar$ makes the proposed lift a
section on their polynomial algebra.
An even basis of the entire target is given by
\[
 A^{\star i}\star B^{\star j}\ (i\geq1),\qquad
 B^{\star j},\qquad
 B^{\star j}\star C^{\star k}\ (k\geq1),
 \qquad A=s^2,\quad C=t^2.
\]
Their leading symbols partition all even monomials in $s,t$.
The lower terms are triangular by polynomial degree, so
these star words are a basis over the parameter ring.
Lift them by the same star words in $a,b,c$.
Choose closed lifts of the odd monomial basis using the
surjective restriction on $S\xi\oplus S\eta$.
This gives a unital weight-preserving chain section $i'$.

From the previous normalized contraction put
$H'=H_\hbar(1-i'\rho_\hbar)$.
The chain-section identity gives
\[
 \delta_\hbar H'+H'\delta_\hbar=1-i'\rho_\hbar.
\]
Moreover $H'i'=\rho_\hbar H'=0$.
The identities $\rho_\hbar H_\hbar=H_\hbar^2=0$
give $(H')^2=0$.
The binary defect on the polynomial subalgebra is zero
because $i'$ is multiplicative there. Inductively every
higher residual consists of lower higher components on
that same subalgebra, so the recursion chooses zero.
\end{proof}

The change already appears on the square:
\[
 i_\hbar(B\star B)=ac+\frac34\hbar^2,\qquad
 i'(B\star B)=b^2-\frac34\hbar^2.
\]
Their difference is $-\delta_\hbar z$.
The target product has not changed. In particular the
repeated-$B$ fourth value in
\eqref{eq:vol2-relative-fourth-inverse} can be removed.

Higher data must nevertheless occur somewhere for the
weight-preserving recursive inverses. A classical section
cannot be multiplicative by
Proposition~\ref{prop:vol2-relative-no-strict-section}, so some
even pair $F,G$ has binary defect $\Delta\phi_0$ with
$\phi_0\ne0$. Its quantum binary primitive has the form
$\phi_\hbar z$, where $\phi_\hbar\bmod\hbar=\phi_0$.
Invariant coefficients Poisson-commute with each $\mu_i$.
It follows that
\[
 (\phi_\hbar z)\star(\phi_\hbar z)
   =\bigl(\hbar\phi_0^2+O(\hbar^2)\bigr)w\ne0.
\]
Here $O(\hbar^2)$ denotes a polynomial divisible by $\hbar^2$.
The degree and parity proof above still makes $I_3=0$.
On $(F,G,F,G)$ the fourth residual is this nonzero product,
so injectivity of $\delta_\hbar:S[\hbar]r\to S[\hbar]w$
forces a nonzero fourth value.

These facts are compatible with formality.
A dg algebra is \emph{formal} if a sequence of dg-algebra
quasi-isomorphisms identifies it with its cohomology algebra
with zero differential.
The strict map $\rho_\hbar$ already proves that
$\mathcal K_\hbar$ is formal. Failure of a section in its
reverse direction does not contradict that conclusion.
Neither a chosen fourth value nor its physical interpretation
is an invariant obstruction merely by being nonzero.

\section{The odd coefficient bimodule}
\label{sec:vol2-relative-odd-bimodule}

The factor $\theta$ makes the product of two odd cohomology
classes zero. Removing that factor gives a different
operation worth retaining. Write
\[
 A_\hbar=W_\hbar(s,t)^{\mathrm{even}},\qquad
 M_\hbar=W_\hbar(s,t)^{\mathrm{odd}}.
\]
Parity here refers to the sign involution on the even
coordinates, not to cohomological degree. Multiplication
gives a bimodule map
$M_\hbar\otimes_{A_\hbar}M_\hbar\to A_\hbar$.

\begin{proposition}\label{prop:vol2-relative-odd-bimodule}
Over $\mathbb C[\hbar]$, the image of this multiplication
is the ideal $(s^2,st,t^2,\hbar)$ in $A_\hbar$, with
cokernel $\mathbb C$ at $s=t=\hbar=0$.
After inverting $\hbar$, multiplication is an isomorphism,
with inverse
\[
 a\longmapsto
   \hbar^{-1}(s\otimes t-t\otimes s)a.
\]
\end{proposition}

\begin{proof}
The four basic products $s\star s,s\star t,t\star s,t\star t$
generate $s^2,st,t^2,\hbar$. Every even ordered monomial of
positive degree splits into two odd monomials.
Their products span all the other even monomials modulo
lower terms, by induction on degree. Thus their image is
the stated ideal. Its quotient leaves only the scalar origin.

In the localized ring, multiplication of the displayed
inverse gives $a$ because $s\star t-t\star s=\hbar$.
For the reverse composite on $m\otimes n$, use that
$t\star m$ and $s\star m$ are even. Balancing them across
the tensor product gives numerator
\[
 s\otimes(t\star m)\star n
   -t\otimes(s\star m)\star n
  =(s\star t-t\star s)\star m\otimes n
  =\hbar m\otimes n.
\]
This proves the second identity.
\end{proof}

Only this last invertibility assertion inverts the parameter.
The strict relative comparison was integral.
An interpretation of the coefficient bimodule product as
composition of physical lines requires a realization map;
the algebraic calculation alone does not construct one.

\section{The ordinary reduction at every matrix size}
\label{sec:vol2-integral-matrix-reduction}

The rank-two complex has retained the relation classes.
There is also an ordinary quantum reduction at every matrix
size, defined across the discriminant. Its global trace
action is an input to a derived comparison, rather than
a replacement for that comparison.

Fix $N\geq1$ and set $R=\mathbb C[\hbar]$.
Let $W_{N,\hbar}$ be generated by two $N$ by $N$ matrices
$X,Y$ with
\[
 [X_{ij},X_{kl}]=[Y_{ij},Y_{kl}]=0,\qquad
 [X_{ij},Y_{kl}]=\hbar\delta_{il}\delta_{jk}.
\]
The group $GL_N(\mathbb C)$ acts by simultaneous conjugation.
Use the ordered moment matrix from
Section~\ref{sec:vol2-two-quantizations},
\[
 \widehat\mu=XY-YX-N\hbar I_N,
 \qquad
 I_{N,\hbar}=\sum_{i,j}W_{N,\hbar}\widehat\mu_{ij}.
\]
This is a left ideal. Its invariant intersection is
two-sided in the invariant algebra, since invariant
operators commute with the infinitesimal conjugation
operators. Define the ordinary reduction
\begin{equation}\label{eq:vol2-ordinary-matrix-reduction}
 \mathcal A_{N,\hbar}
 =W_{N,\hbar}^{GL_N}/
       (I_{N,\hbar}\cap W_{N,\hbar}^{GL_N}).
\end{equation}
Replacing the full left ideal by the two-sided ideal it
generates would define a different problem.

Let $\mathfrak g=\mathfrak{gl}_N$ and let $D(\mathfrak g)$
be the polynomial differential operators.
Its order filtration $F_m$ allows arbitrary polynomial
coefficients and at most $m$ derivatives, with $F_{-1}=0$.
The Rees algebra is
\[
 \operatorname{Rees}_F D(\mathfrak g)
  =\bigoplus_{m\geq0}\hbar^mF_mD(\mathfrak g)
       \subset D(\mathfrak g)[\hbar].
\]
The assignments
$X_{ij}\mapsto x_{ij}$ and
$Y_{ij}\mapsto-\hbar\partial_{x_{ji}}$ identify it with
$W_{N,\hbar}$. The Weyl relations span the monomials with
all $X$ before all $Y$. Their differential-operator images
are independent after inverting $\hbar$, so they have no
nonzero $R$-linear relation before inversion either.
They span the Rees algebra together with powers of $\hbar$.
Thus $X$ has Rees order zero, while $Y$ and $\hbar$ have
order one. These are not cohomological degrees.

Let $\tau(a)$ be the differentiated pullback conjugation
action, normalized by $\tau(a)X=[X,a]$ on coordinate
functions. The Weyl relations give
\[
 \operatorname{Tr}(a\widehat\mu)=\hbar\tau(a).
\]
In particular $\operatorname{Tr}\widehat\mu=0$.
Put $J_D=(D(\mathfrak g)\tau(\mathfrak g))^{GL_N}$.
Let $J_{\mathrm{cl}}$ be the ideal generated by the entries
of the commutator of a pair in $\mathfrak g\oplus\mathfrak g$.
The trace pairing identifies the second factor with
cotangent momentum.

Write $\mathfrak t$ for diagonal matrices and $S_N$ for
its Weyl group. On the distinct-eigenvalue locus, the raw
radial operator $\operatorname{rad}(D)$ acts on restrictions
of invariant functions. Its conjugated form is
\[
 \Phi(D)=\delta\,\operatorname{rad}(D)\,\delta^{-1},
 \qquad
 \delta=\prod_{i<j}(\lambda_i-\lambda_j).
\]
The radial definition begins on this open set.
The following theorem is the extension statement that
allows it to be used across the discriminant.

\begin{theorem}[Filtered Harish--Chandra comparison]
\label{thm:vol2-filtered-harish-chandra}
For $\mathfrak g=\mathfrak{gl}_N$, the conjugated radial map
induces an algebra isomorphism
\[
 D(\mathfrak g)^{GL_N}/J_D
   \xrightarrow{\ \sim\ }D(\mathfrak t)^{S_N}.
\]
It is strict for the order filtrations. Its associated
graded map is restriction of commuting pairs to diagonal
pairs, and
\[
 \operatorname{gr}_FJ_D
       =(\sqrt{J_{\mathrm{cl}}})^{GL_N}
       =J_{\mathrm{cl}}^{GL_N}.
\]
\end{theorem}

The algebra isomorphism and associated graded assertions are
the Harish--Chandra results recorded by Etingof and Ginzburg
(2001), Theorems~6.8--6.10 of~\cite{EG2001HC}.
The equality of the invariant ideal with its radical is
Gan--Ginzburg (2005), Theorem~1.2.1 and
Section~2.8 of~\cite{GG2005Almost}.
Strictness follows by lifting a leading symbol and subtracting:
the differential order decreases at each step.
The equality of invariant ideals does not assume that the
whole commuting scheme is reduced.

\begin{lemma}\label{lem:vol2-strict-moment-orders}
Every $D\in J_D\cap F_mD(\mathfrak g)$ is a finite invariant
sum of terms $P_a\tau(a)$ with
$P_a\in F_{m-1}D(\mathfrak g)$.
For $m=0$ the only such $D$ is zero.
\end{lemma}

\begin{proof}
The leading symbol $\sigma_m(D)$ belongs to
$J_{\mathrm{cl}}^{GL_N}$.
Each moment generator has momentum degree one.
Its homogeneous ideal expression can therefore be written
\[
 \sigma_m(D)=\sum_a p_a\,\sigma_1(\tau(a)),
 \qquad \deg_{\mathrm{momentum}}p_a=m-1.
\]
Only finitely many coefficients occur.
The map from coefficient vectors to the ideal is equivariant.
Rational $GL_N$-invariants are exact in characteristic zero:
on a finite-dimensional rational module, average over the
unitary compact form. Applying this to a degree-bounded
module containing the coefficients gives an invariant
coefficient vector with the same image.

Normal quantization, placing coordinates before derivatives,
is equivariant under the linear action. It lifts the $p_a$
to $P_a$ of order at most $m-1$.
The invariant operator $\sum_aP_a\tau(a)$ has the same
leading symbol as $D$. Subtract it and repeat at smaller
order. This process terminates.
The degree-zero momentum part of $J_{\mathrm{cl}}$ is zero,
so the last order-zero remainder vanishes.
\end{proof}

An ideal $J$ is $\hbar$-saturated if $\hbar u\in J$
implies $u\in J$. Let $W^{\mathrm{Cartan}}_{N,\hbar}$
be the Weyl algebra on $N$ pairs $\lambda_i,\pi_i$ with
$[\lambda_i,\pi_j]=\hbar\delta_{ij}$ and all other
generator commutators zero.

\begin{theorem}\label{thm:vol2-integral-harish-chandra}
The conjugated radial map has an $R$-algebra realization
\[
 \Phi_{N,\hbar}:\mathcal A_{N,\hbar}
   \xrightarrow{\ \sim\ }
   (W^{\mathrm{Cartan}}_{N,\hbar})^{S_N}.
\]
Moreover
\[
 I_{N,\hbar}\cap W_{N,\hbar}^{GL_N}
     =\operatorname{Rees}_FJ_D.
\]
This invariant ideal is $\hbar$-saturated, the quotient
is $R$-free, and its special fibre is
\[
 \mathcal A_{N,\hbar}/\hbar\mathcal A_{N,\hbar}
   =\mathbb C[\lambda_1,\pi_1,\ldots,\lambda_N,\pi_N]^{S_N}.
\]
\end{theorem}

\begin{proof}
The conjugation action preserves every filtration step, so
invariants give
$W_{N,\hbar}^{GL_N}
 =\operatorname{Rees}_F(D(\mathfrak g)^{GL_N})$.
In each Rees degree an element of the left moment ideal
has coefficient in $D(\mathfrak g)\tau(\mathfrak g)$.
Its invariant part therefore belongs to
$\operatorname{Rees}_FJ_D$.

Conversely, let $\hbar^mD$ have $D\in F_mJ_D$.
The preceding lemma expresses
$D=\sum_aP_a\tau(a)$ with $P_a\in F_{m-1}D(\mathfrak g)$.
Hence
\[
 \hbar^mD=\sum_a(\hbar^{m-1}P_a)(\hbar\tau(a))
       \in I_{N,\hbar}.
\]
This element is invariant, proving equality of the ideals.
The strict filtered Harish--Chandra isomorphism now
identifies the quotient with the Rees algebra of
$D(\mathfrak t)^{S_N}$.
Its Weyl generators are the displayed $\lambda_i,\pi_i$.

Weyl symmetrization identifies its underlying $R$-module
equivariantly with the ordinary polynomial module.
A complex basis of the $S_N$-invariant polynomial space
therefore gives an $R$-basis of the quotient.
It has no $\hbar$-torsion. Thus if $\hbar u$ maps to zero,
$u$ maps to zero too, which is saturation.
The Rees special fibre is its associated graded algebra,
giving the final formula.
\end{proof}

Saturation here concerns the invariant intersection.
No assertion about saturation of the entire left ideal
was required. At $\hbar=0$, $\pi_i$ remains a momentum
symbol in the intrinsic special fibre. Substituting
$\hbar=0$ into the external representation
$\pi_i=-\hbar\partial_{\lambda_i}$ would instead annihilate
it. Specializing that representation does not compute
the Rees special fibre.

\section{The global ordinary trace action}
\label{sec:vol2-global-ordinary-moyal-action}

For $f,g\in R[q,p]$, let
$B_\hbar(f,g)=(f\star_\hbar g-g\star_\hbar f)/\hbar$
be the scalar Moyal bracket. Its numerator is polynomial
and divisible by $\hbar$.
For a monomial define the matrix-word symmetrized trace
\[
 J_N(q^ap^b)
  =\binom{a+b}{a}^{-1}\sum_{v}\operatorname{Tr}v(X,Y),
\]
where $v$ runs over distinct ordered products with $a$
occurrences of $X$ and $b$ of $Y$.
Extend $J_N$ $R$-linearly. In particular $J_N(1)=N$.
Let $\operatorname{Op}_W(f)(\lambda,\pi)$ denote
one-pair Weyl symmetrization.

\begin{theorem}\label{thm:vol2-global-moyal-trace}
For every matrix size $N$ and every polynomial $f$,
\begin{equation}\label{eq:vol2-global-moyal-trace-image}
 \Phi_{N,\hbar}([J_N(f)])
     =\sum_{i=1}^N\operatorname{Op}_W(f)(\lambda_i,\pi_i).
\end{equation}
Consequently
\[
 \frac{[J_N(f),J_N(g)]}{\hbar}-J_N(B_\hbar(f,g))
       \in I_{N,\hbar}\cap W_{N,\hbar}^{GL_N}.
\]
The induced map
$R[q,p]/R1\to\mathcal A_{N,\hbar}/R1$ is injective
and preserves the normalized Lie brackets.
\end{theorem}

\begin{proof}
Multiplication by $\operatorname{Tr}X^d$ restricts to
$\sum_i\lambda_i^d$.
The radial image of the quadratic momentum trace is
\[
 \operatorname{rad}(\operatorname{Tr}Y^2)
 =\hbar^2\left(
   \sum_i\partial_{\lambda_i}^2
    +2\sum_{i<j}
       \frac{\partial_{\lambda_i}-\partial_{\lambda_j}}
            {\lambda_i-\lambda_j}\right).
\]
To see the mixed derivative coefficient, perturb entries
$(i,j)$ and $(j,i)$ by $u,v$. Their two eigenvalues satisfy
$(z-\lambda_i)(z-\lambda_j)-uv=0$.
Their mixed second derivatives are
$\pm(\lambda_i-\lambda_j)^{-1}$, yielding the displayed
first-order terms.
The Vandermonde polynomial $\delta$ is harmonic:
its Laplacian is alternating of lower degree, hence
divisible by $\delta$ and therefore zero.
Conjugation by $\delta$ removes the first-order terms.
Thus $\Phi_{N,\hbar}([\operatorname{Tr}Y^2])
 =\sum_i\pi_i^2$.

Put $H=-\operatorname{Tr}Y^2/2$.
The polynomial derivation $D=[H,-]/\hbar$ satisfies
$DX=Y$ and $DY=0$.
Its exponential terminates on any polynomial in the
matrix entries and gives
\[
 e^{tD}\operatorname{Tr}X^d
      =\operatorname{Tr}(X+tY)^d.
\]
Multiplicativity of $\Phi_{N,\hbar}$ and torsion-freeness
of its target show that it intertwines $D$ with
$\lambda_i\mapsto\pi_i$, $\pi_i\mapsto0$.
The coefficient of $t^b$, with $a+b=d$, is the sum of
the matrix products defining $J_N(q^ap^b)$.
The same coefficient on the Cartan side is the sum of
the corresponding one-pair Weyl products.
Dividing by the same nonzero binomial coefficient proves
\eqref{eq:vol2-global-moyal-trace-image}.

Different Cartan pairs commute, and one-pair Weyl
quantization intertwines Moyal products. The stated
bracket defect has zero radial image.
The exact kernel in
Theorem~\ref{thm:vol2-integral-harish-chandra} puts it
in the invariant moment ideal.
The division by $\hbar$ is the polynomial quotient of
a divisible numerator throughout.

Finally suppose $[J_N(f)]$ is scalar.
Taking Cartan Weyl symbols gives
$\sum_i f(\lambda_i,\pi_i)=\alpha$ for some $\alpha\in R$.
Set all pairs except the first equal to zero. Then
$f(q,p)+(N-1)f(0,0)=\alpha$, so $f$ is constant.
This proves injectivity on the quotient by constants.
\end{proof}

The theorem constructs a global ordinary trace action and
ideal membership, not a unique expression of every defect
in moment generators. The strict order argument supplies
finite expressions after choices. Compatible primitives
in a resolution must still satisfy their higher equations.
The scalar rank $J_N(1)=N$ also remains present before
passing to the quotient by constants.

The parameter-adic completion has an exact description.
Writing $J=I_{N,\hbar}\cap W_{N,\hbar}^{GL_N}$,
saturation gives
$J\cap\hbar^mW_{N,\hbar}^{GL_N}=\hbar^mJ$.
The $R$-free quotient makes reduction of the defining
short exact sequence modulo $\hbar^m$ exact.
The transition maps $J/\hbar^{m+1}J\to J/\hbar^mJ$ are
surjective, so compatible lifts corrected in these
kernels prove exactness after inverse limit.
Thus $\Phi_{N,\hbar}$ completes to an isomorphism, and
the completed kernel is the closure of $J$.
Every product at fixed $\hbar$-order uses finitely many
polynomial coefficients. Their degrees may grow with
the order.

\clearpage
\section{Relative reduction and the remaining comparison}
\label{sec:vol2-relative-reduction-comparison}

There are now two positive comparisons.
The relative rank-two dg algebra maps strictly to
its Weyl--Cartan cohomology algebra, retaining the odd
relation module and every specialization.
The ordinary reduction at arbitrary rank has its
integral trace action across the discriminant.
The maps have different sources and preserve different
data. Neither identifies the absolute BRST complex
merely by taking rational group invariants.

For rank three, an extension requires a specified
relative quantum relation complex and a map to its
derived Cartan recipient.
Its relation cohomology and possible invariant
Massey or extension classes must be computed.
A value of one inverse component would not decide
whether such a comparison exists.

The full Hamiltonian question also retains the positive
gauge ghosts. It asks for a coherent moment comparison
on the absolute rank-two quantum recipient and its
higher-rank extensions, strict on the scalar
quadratics and with the cubic primitive of
Section~\ref{sec:vol2-derived-action-bfv}.
The comparison must carry the negative relation classes,
stabilizer states and cyclic operations across the
discriminant. The coefficient-plane parameter
$\varepsilon$ is still independent of matrix $\hbar$.
These requirements are additional to the strict relative
restriction and the ordinary trace theorem.
They are the operations to compare after the absolute
complex and its higher model have been constructed.

\chapter{Absolute quantum reduction and higher products}
\label{chap:vol2-absolute-quantum-reduction}

The relative reduction retains relations among the moment
equations. The absolute reduction must also retain the gauge
directions. Even a trivial coefficient module for
$\mathfrak{so}_3$ has a degree-three Lie cohomology class.
Multiplying this class by the relative degree-minus-one relation
class produces a degree-two class. Both disappear if one keeps
only invariant coefficients and discards the positive ghosts.

We compute the full polynomial traceless rank-two complex over
$R=\mathbb C[\hbar]$. Its cohomology is supported in degrees
$-1,0,2,3$. A contraction constructs this answer across the
discriminant and specifies every higher product.
Its projection cannot be a unital graded algebra map to the
cohomology ring. That obstruction has a fixed direction:
whether a zigzag of dg-algebra quasi-isomorphisms exists is a
separate question.

\section{The positive ghosts}
\label{sec:vol2-absolute-ghosts}

Use the traceless coordinates $x_i,y_i$ and antighosts $\psi_i$
of Chapter~\ref{chap:vol2-relative-quantum-reduction}, for
$1\leq i\leq3$. Their degrees are $0,0,-1$.
Adjoin positive ghosts $\gamma^i$ of degree $1$ and impose
\begin{equation}\label{eq:vol2-absolute-ghost-relations}
 [x_i,y_j]=\hbar\delta_{ij},\qquad
 \psi_i\gamma^j+\gamma^j\psi_i=\hbar\delta_i^j,
 \qquad
 \gamma^i\gamma^j=-\gamma^j\gamma^i,\quad
 \psi_i\psi_j=-\psi_j\psi_i.
\end{equation}
The even coordinates commute with the ghosts.
The $x_i$ commute among themselves, as do the $y_i$.
The odd squares are zero. Write
$\mathrm{Cl}_\hbar(\gamma,\psi)$ for this graded Clifford
algebra and put
\[
 \mathcal B_\hbar
 =W_\hbar(x,y)\otimes_R\mathrm{Cl}_\hbar(\gamma,\psi).
\]
Use symmetric Weyl symbols in $x,y$ and place all positive
ghosts to the left of all antighosts. The ordered products
$\gamma^I\psi_J$, for increasing subsets $I,J\subset\{1,2,3\}$,
form an $R$-basis of the ghost algebra.
Indeed the relations span it by these products. After inverting
$\hbar$, their action on $\Lambda(\gamma)$ by multiplication
and $\psi_i=\hbar\partial_{\gamma^i}$ is the full matrix
algebra, whose dimension is $2^6$. Thus they are independent
already over $R$.

For the rotation moment define
\begin{equation}\label{eq:vol2-absolute-brst-charge}
 \mu_i=\sum_{j,k}\epsilon_{ijk}x_jy_k,\qquad
 \Theta=\sum_i\gamma^i\mu_i
       -\frac12\sum_{i,j,k}\epsilon_{ijk}
                    \gamma^i\gamma^j\psi_k,
 \qquad Q=\frac{[\Theta,-]_{\mathrm{super}}}{\hbar}.
\end{equation}
Here $\epsilon_{123}=1$ is alternating, and the bracket is the
graded commutator for the full product.
Division means the polynomial quotient of a divisible numerator.
The algebra modulo $\hbar$ is graded commutative, so every such
commutator is divisible. The Weyl--Clifford basis makes the
quotient unique.

The moment relation is
$[\mu_i,\mu_j]=\hbar\sum_k\epsilon_{ijk}\mu_k$.
Together with Jacobi it gives $\Theta^2=0$, as in
Proposition~\ref{prop:vol2-integral-quantum-brst}.
One can also represent the ghosts on $\Lambda(\gamma)$ after
inverting $\hbar$. The charge becomes the Chevalley operator
for the representation $\mu_i/\hbar$, multiplied by $\hbar$.
Its square vanishes. Faithfulness and polynomial torsion-freeness
give the same identity before inversion.
Consequently $Q^2=0$, and $Q$ is a degree-one $R$-linear
derivation. These statements remain valid at $\hbar=0$ by
specializing the polynomial formulas.

The distinction from invariants already appears in the scalar
Chevalley complex. Its differential is
\begin{equation}\label{eq:vol2-scalar-ghost-differential}
 d\gamma^1=-\gamma^2\gamma^3,\qquad
 d\gamma^2=\gamma^1\gamma^3,\qquad
 d\gamma^3=-\gamma^1\gamma^2.
\end{equation}
It is an isomorphism from exterior degree one to degree two
and is zero in every other degree. Thus its cohomology is
$\Lambda(\chi)$, where
\[
 \chi=\gamma^1\gamma^2\gamma^3,\qquad |\chi|=3.
\]
For the full complex, the positive ghosts also act on the
polynomial coefficients. We next contract that action without
removing the scalar and top-ghost classes.

\section{Normal ordering and the two differentials}
\label{sec:vol2-absolute-normal-order}

Let
$K_{\mathrm{all}}=W_\hbar(x,y)\otimes\Lambda(\psi)$
be the space of coefficients with no positive ghost.
For the basis $e_i$ of $\mathfrak{so}_3$, use
$[e_i,e_j]=\sum_k\epsilon_{ijk}e_k$.
The simultaneous coefficient action $\rho_i$ is
\begin{equation}\label{eq:vol2-absolute-coefficient-action}
 \rho_i x_j=\sum_k\epsilon_{ijk}x_k,
 \qquad \rho_i y_j=\sum_k\epsilon_{ijk}y_k,
 \qquad \rho_i\psi_j=\sum_k\epsilon_{ijk}\psi_k.
\end{equation}
It extends by the derivation rule. On Weyl coefficients it is
$[\mu_i,-]/\hbar$. It does not act on $\gamma$ in this notation.
Each polynomial lies in a finite-dimensional rational rotation
module, so the coefficient representation is locally finite.

For $X\in\mathfrak{so}_3$, let $\mu_X$ and $\psi_X$ denote
the corresponding linear combinations. Define a degree-one
operator on the antighost complex by
\begin{align}\label{eq:vol2-absolute-chain}
 \partial_N(a\psi_{X_1}\cdots\psi_{X_q})
 ={}&\sum_{i=1}^q(-1)^{i-1}(a\star\mu_{X_i})
     \psi_{X_1}\cdots\widehat{\psi_{X_i}}\cdots\psi_{X_q}
       \notag\\
 &+\hbar\sum_{i<j}(-1)^{i+j}a\psi_{[X_i,X_j]}
     \psi_{X_1}\cdots\widehat{\psi_{X_i}}\cdots
                       \widehat{\psi_{X_j}}\cdots\psi_{X_q}.
\end{align}
The bracket factor in the second line stands first.
For example,
\[
 \partial_N(a\psi_i\psi_j)
 =a\star\mu_i\,\psi_j-a\star\mu_j\,\psi_i
          -\hbar a\psi_{[e_i,e_j]}.
\]
Applying $\partial_N$ again gives the moment commutator minus
the same Lie bracket. It is zero. In higher antighost degree,
the terms involving two moments cancel in the same way;
the terms involving three indices cancel by Jacobi.
This proves $\partial_N^2=0$ in every degree.
The moment map and the Lie bracket are equivariant, so
$\partial_N$ commutes with each $\rho_i$.

Let $\iota_i=\partial_{\gamma^i}$ be the left exterior
derivative. Define
\begin{equation}\label{eq:vol2-absolute-chevalley}
 d_{\mathrm{CE}}=\sum_i\gamma^i\rho_i
       -\frac12\sum_{i,j,k}\epsilon_{ijk}
                     \gamma^i\gamma^j\iota_k.
\end{equation}
This is the Chevalley cochain differential with coefficients in
$K_{\mathrm{all}}$. The Lie representation identity and Jacobi
give $d_{\mathrm{CE}}^2=0$.
On positive-ghost degree $p$, put
$\partial^{\mathrm{tot}}=(-1)^p\partial_N$.
Equivariance gives
$d_{\mathrm{CE}}\partial^{\mathrm{tot}}
 +\partial^{\mathrm{tot}}d_{\mathrm{CE}}=0$.

\begin{lemma}\label{lem:vol2-absolute-normal-order}
Normal ordering identifies the underlying complex with
\begin{equation}\label{eq:vol2-absolute-bicomplex}
 (\mathcal B_\hbar,Q)
  =\bigl(\Lambda(\gamma)\otimes K_{\mathrm{all}},
          d_{\mathrm{CE}}+\partial^{\mathrm{tot}}\bigr).
\end{equation}
On invariant positive-ghost-free coefficients, $Q$ is the
relative differential $\delta_\hbar$ of
\eqref{eq:vol2-relative-quantum-differential}.
\end{lemma}

\begin{proof}
Commute each positive ghost produced by $[\Theta,-]$ to the
left. The terms retaining one new ghost give the simultaneous
coefficient action. A contraction with an existing positive
ghost gives the second term of
\eqref{eq:vol2-absolute-chevalley}.
Removing one antighost gives the moment term in
\eqref{eq:vol2-absolute-chain}. Contracting the cubic ghost
term against two antighosts gives its bracket term.
Moving this degree-one operator past $p$ positive ghosts
gives $(-1)^p$.

On generators, for example,
\[
 Q\gamma^1=-\gamma^2\gamma^3,
 \qquad Q\psi_1=\mu_1-\gamma^2\psi_3+\gamma^3\psi_2,
 \qquad Qx_j=\sum_{i,k}\gamma^i\epsilon_{ijk}x_k.
\]
The same formula holds for $Qy_j$.
The ordered computation above applies to an arbitrary product
and hence proves the identity on the whole basis.

For invariant coefficients, $\sum\gamma^i\rho_i$ vanishes.
Invariance combines the ordered moment and bracket terms into
\[
 \frac12\sum_i\left(
     \mu_i\star\partial_{\psi_i}F
       +\partial_{\psi_i}F\star\mu_i\right).
\]
The first Moyal derivatives cancel. The second derivatives of
the bilinear moments give
$\hbar^2\sum\epsilon_{ijk}
 \partial_{x_j}\partial_{y_k}\partial_{\psi_i}/4$,
and all higher derivatives vanish. This is the stated
relative differential.
\end{proof}

Equation~\eqref{eq:vol2-absolute-bicomplex} is an equality of
complexes after choosing an ordered basis.
Multiplication still uses
\eqref{eq:vol2-absolute-ghost-relations}.
Replacing it by the exterior tensor product would remove
terms needed for both the ring comparison and the higher products.

\section{Contracting the coefficient action}
\label{sec:vol2-absolute-casimir}

The scalar Chevalley complex should remain. To contract its
nontrivial coefficient representations, define
\begin{equation}\label{eq:vol2-absolute-casimir}
 \mathsf C=\sum_i\rho_i^2,
 \qquad K=\sum_i\rho_i\iota_i.
\end{equation}
Here $\mathsf C$ acts only on $K_{\mathrm{all}}$.
The diagonal rotation action on both coefficients and positive
ghosts has a different Casimir.

\begin{lemma}\label{lem:vol2-absolute-casimir}
The operators satisfy
\[
 d_{\mathrm{CE}}K+Kd_{\mathrm{CE}}=\mathsf C,
 \qquad
 \partial^{\mathrm{tot}}K+K\partial^{\mathrm{tot}}=0.
\]
Moreover $\mathsf C$ commutes with both differentials.
\end{lemma}

\begin{proof}
Use $\iota_j\gamma^i+\gamma^i\iota_j=\delta_j^i$.
The anticommutator of $K$ with $\sum\gamma^i\rho_i$ is
\[
 \mathsf C+\sum_{i,j,k}\epsilon_{ijk}
                   \gamma^i\rho_k\iota_j.
\]
Its anticommutator with the bracket part of
$d_{\mathrm{CE}}$ is the negative of the second summand.
Indeed commuting $\iota_\ell$ through
$\gamma^i\gamma^j$ gives
$\delta_\ell^i\gamma^j-\delta_\ell^j\gamma^i$;
antisymmetry then gives the displayed cancellation.

The coefficient Casimir is central for the rotation action,
so it commutes with $d_{\mathrm{CE}}$.
Equivariance of $\partial_N$ gives its commutation with
$\rho_i$ and $\mathsf C$.
Finally $K$ lowers positive-ghost degree by one, reversing
the totalization sign of $\partial_N$. This proves the
second anticommutator identity.
\end{proof}

On an irreducible spin-$\ell$ coefficient representation,
$\mathsf C=-\ell(\ell+1)$, for $\ell\geq0$.
These are integer spins because the representation is rational
for $SO(3)$.
To check the normalization, put
$E=\rho_1+i\rho_2$, $F=-\rho_1+i\rho_2$ and
$H_{\mathrm{wt}}=2i\rho_3$. They satisfy
$[H_{\mathrm{wt}},E]=2E$, $[H_{\mathrm{wt}},F]=-2F$,
$[E,F]=H_{\mathrm{wt}}$, and
\[
 \mathsf C=-\frac14H_{\mathrm{wt}}^2-\frac12(EF+FE).
\]
On a highest-weight vector of weight $2\ell$, this is
$-\ell^2-\ell$. Complete reducibility in each finite coefficient
module follows by averaging a Hermitian form over the compact
real rotation group and taking invariant orthogonal complements.
Thus the eigenvalue decomposition used below is an actual direct
sum on every finite input module.
For an input contained in spins at most $L$, define
\begin{equation}\label{eq:vol2-absolute-spectral-polynomials}
 P_L(X)=\prod_{\ell=1}^L
               \left(1+\frac{X}{\ell(\ell+1)}\right),
 \qquad
 G_L(X)=\frac{1-P_L(X)}{X}+P_L'(0)P_L(X).
\end{equation}
Both expressions are polynomials.
The second has value zero at $X=0$ and value $1/X$ at
every nonzero Casimir eigenvalue in this range.
Thus
\[
 P_0=P_L(\mathsf C),\qquad G=G_L(\mathsf C)
\]
are, respectively, projection to coefficient invariants and
the inverse Casimir on their complement, extended by zero.
Enlarging $L$ gives the same operators on the original input.
This constructs them on the union of finite coefficient modules.
They commute with equivariant maps, including $\partial_N$.
No power of $\hbar$ or of the discriminant is inverted.

Set $H_0=GK$. Then
\[
 QH_0+H_0Q=1-P_0,\qquad
 P_0H_0=H_0P_0=0.
\]
To impose the square-zero side condition, replace $H_0$ by
\begin{equation}\label{eq:vol2-absolute-normalized-casimir}
 H_1=H_0d_{\mathrm{CE}}H_0.
\end{equation}
On the complement of $P_0$, write $d=d_{\mathrm{CE}}$ and
$h=H_0$. The identity $dh+hd=1$ gives
$dh^2=h^2d$. Hence
\[
 d(hdh)+(hdh)d=dh+hd=1,
 \qquad (hdh)^2=h(dh^2d)h=0.
\]
The total chain differential anticommutes with $h,d,h$ in
succession, and therefore with $hdh$.
Consequently
\begin{equation}\label{eq:vol2-absolute-first-contraction}
 QH_1+H_1Q=1-P_0,
 \qquad H_1^2=H_1P_0=P_0H_1=0.
\end{equation}

The image of $P_0$ is
$\mathrm{CE}^*(\mathfrak{so}_3,\mathbb C)
 \otimes\mathcal K_\hbar$,
with the tensor differential. The next step contracts its
first factor explicitly, then uses the relative comparison.

\section{The absolute contraction}
\label{sec:vol2-absolute-contraction}

Let $j_\gamma:\Lambda(\chi)\to\Lambda(\gamma)$ include
the scalar and top-ghost monomials, and let $p_\gamma$ project
onto them. The homotopy for
\eqref{eq:vol2-scalar-ghost-differential} is
\begin{equation}\label{eq:vol2-absolute-scalar-homotopy}
 h_\gamma(\gamma^2\gamma^3)=-\gamma^1,
 \quad h_\gamma(\gamma^1\gamma^3)=\gamma^2,
 \quad h_\gamma(\gamma^1\gamma^2)=-\gamma^3,
\end{equation}
and zero on other exterior degrees. Direct substitution gives
$dh_\gamma+h_\gamma d=1-j_\gamma p_\gamma$ and the
normalized side conditions.

Write $(\rho_\hbar,i_{\mathrm{rel}},H_{\mathrm{rel}})$
for the relative contraction in
Section~\ref{sec:vol2-relative-higher-inverse}.
Its target is
\[
 \mathcal H_\hbar
  =\bigl(W_\hbar(s,t)\otimes\Lambda(\theta)\bigr)^{\mathbb Z/2},
 \qquad [s,t]=\hbar,\quad |\theta|=-1.
\]
The involution negates $s,t,\theta$ and fixes $\hbar$.
Choose its displayed section so that
$i_{\mathrm{rel}}(s\theta)=\xi=x\cdot\psi$,
$i_{\mathrm{rel}}(t\theta)=\eta=y\cdot\psi$ and
$H_{\mathrm{rel}}(\xi\star\eta)=r=\psi_1\psi_2\psi_3$.
Its formulas use finite classical weight splittings and the
terminating coefficient operator $T^{-1}$.

Put $Z=\Lambda(\chi)\otimes\mathcal H_\hbar$, initially
as a graded $R$-module with zero differential.
Compose the three contractions: first
\eqref{eq:vol2-absolute-first-contraction}, then
$(j_\gamma,p_\gamma,h_\gamma)$ tensored with the relative
complex, and finally the relative contraction tensored with
$\Lambda(\chi)$.
On the last tensor product, for $e\in\{0,1\}$, the homotopy is
\[
 H'(\chi^e\otimes u)
       =(-1)^{3e}\chi^e
                                      \otimes H_{\mathrm{rel}}u.
\]
For two successive contractions, the composite homotopy is
$H_A+i_AH_Bp_A$.
Expansion of this expression proves the contraction identity;
the mixed terms in its square vanish by the side conditions.
This specifies the final homotopy $H$ on every polynomial input.

\begin{theorem}\label{thm:vol2-absolute-contraction}
There is a normalized contraction over $R$
\[
 (\mathcal B_\hbar,Q)
    \underset{i}{\overset{\Pi}{\rightleftarrows}}(Z,0),
 \qquad QH+HQ=1-i\Pi,
 \qquad \Pi i=1,\quad H^2=Hi=\Pi H=0.
\]
It is polynomial on each input and remains a contraction
after every commutative coefficient base change $R\to R'$.
For $F=\sum_I\gamma^IF_I$ in normal order,
\begin{equation}\label{eq:vol2-absolute-projection}
 \Pi(F)=\rho_\hbar(P_0F_\varnothing)
            +\chi\rho_\hbar(P_0F_{123}),
 \qquad
 i(u+\chi v)=i_{\mathrm{rel}}u+\chi i_{\mathrm{rel}}v.
\end{equation}
\end{theorem}

\begin{proof}
The three contractions just constructed have precisely these
inclusion and projection maps. Composition proves all displayed
identities. The Casimir operators use finite-dimensional
coefficient representations and only nonzero complex integer
denominators. The relative operator $T^{-1}$ terminates on
each polynomial, as do the finite weight splittings.
Every constituent is $R$-linear with polynomial output.
Tensoring the identities with $R'$ therefore preserves them,
without a flatness assumption on $R'$.
The resulting homotopy also proves that the projection kernel
is contractible, rather than merely giving a list of proposed
cohomology representatives.
\end{proof}

Let $E=W_\hbar(s,t)^{\mathrm{even}}$ and
$M=W_\hbar(s,t)^{\mathrm{odd}}$, where parity means total
polynomial degree modulo two. The cohomology groups are
\begin{equation}\label{eq:vol2-absolute-cohomology-degrees}
 H^{-1}=M\theta,
 \qquad H^0=E,
 \qquad H^2=\chi M\theta,
 \qquad H^3=\chi E,
\end{equation}
and all other degrees vanish. The odd module and its singular
fibre from the relative calculation remain present.
The absolute operation adds the two positive degrees.
This is a polynomial statement. Extending the contraction to
a completion requires its own continuity and convergence
conditions.

\section{Products and the failure of a strict projection}
\label{sec:vol2-absolute-products}

Give $Z=\Lambda(\chi)\otimes\mathcal H_\hbar$ the graded
tensor product. Thus
$u\chi=(-1)^{|u|}\chi u$ and $\chi^2=0$.
The product on $\mathcal H_\hbar$ is its Weyl--exterior
product from the relative restriction.

\begin{proposition}\label{prop:vol2-absolute-ring}
Under the contraction, the cohomology algebra is
\[
 H^*(\mathcal B_\hbar,Q)
       \cong\Lambda(\chi)\otimes\mathcal H_\hbar.
\]
This is an isomorphism of cohomology rings. The linear
projection $\Pi$ is not an algebra map on the full complex.
\end{proposition}

\begin{proof}
Every relative class has an invariant closed representative
of antighost degree zero or one. Products of representatives
with no top ghost have the relative cohomology product.
For a product with one top ghost, normal ordering preserves
the term with three positive ghosts, with exactly the graded
tensor sign. All contraction terms have positive-ghost
degree two and are killed by $\Pi$.
The relative projection of the remaining coefficient is
its relative product.

Products with two top ghosts and at most two antighosts
are zero. Six positive ghosts cannot survive among three
exterior generators unless at least three are removed by
contractions. Only two antighosts are available.
This proves all products in the asserted ring.
Since the representatives are closed and $\Pi$ is the
projection of a contraction, these projected products compute
their actual cohomology classes.
\end{proof}

The graded commutation with $\chi$ is a cohomology identity.
Its cochain defect has an explicit primitive:
\begin{equation}\label{eq:vol2-absolute-top-commutators}
 \chi\star\xi+\xi\star\chi
     =\hbar Q\left(\sum_i x_i\gamma^i\right),
 \qquad
 \chi\star\eta+\eta\star\chi
     =\hbar Q\left(\sum_i y_i\gamma^i\right).
\end{equation}
For the first equation, normal ordering gives the left side as
\[
 \hbar\left(x_1\gamma^2\gamma^3
       -x_2\gamma^1\gamma^3+x_3\gamma^1\gamma^2\right).
\]
The coefficient-action part of $Q(\sum x_i\gamma^i)$
gives twice the expression in parentheses, and the scalar
Chevalley part subtracts it once. The antighost differential
does not enter. This proves the first identity, including its
sign; replacing $x$ by $y$ proves the second.

\begin{proposition}\label{prop:vol2-absolute-no-strict-projection}
There is no graded unital $R$-algebra map
$\mathcal B_\hbar\to Z$.
\end{proposition}

\begin{proof}
Equation~\eqref{eq:vol2-absolute-cohomology-degrees} shows
$Z^1=0$. Such a map would send every $\gamma^i$ to zero.
The Clifford relation
$\psi_i\gamma^i+\gamma^i\psi_i=\hbar$ would then send
$\hbar$ to zero. A unital $R$-linear map must preserve it,
and $\hbar\ne0$ in the $R$-free target. This is a contradiction.
\end{proof}

The assertion fixes the direction of the map and the coefficient
ring. At the classical specialization $\hbar=0$, its Clifford
obstruction disappears. It also says nothing about a sequence
of quasi-isomorphisms whose arrows may point in either direction.
A dg algebra is \emph{formal} if such a sequence connects it
to its cohomology ring with zero differential. Determining
formality of this absolute algebra requires an invariant of
that comparison problem.

\section{The higher model}
\label{sec:vol2-absolute-higher-model}

The contraction provides the operation data beyond the ring.
Let $s$ denote suspension of degree $-1$, distinct from the
Cartan coordinate in $W_\hbar(s,t)$, and put
\[
 V=s\mathcal B_\hbar,\qquad W=sZ,
 \qquad d_V=-sQs^{-1},\qquad h_V=-sHs^{-1},
 \qquad i_s=sis^{-1},\quad p_s=s\Pi s^{-1}.
\]
For a graded module $U$, use the tensor coalgebra
$T^c(U)=\bigoplus_{n\geq1}U^{\otimes n}$ with the
deconcatenation coproduct. Tensor products of maps use the
Koszul sign. On $T^c(V)$ let $d_T$ be the tensor differential
and let the multiplication coderivation be determined by
\begin{equation}\label{eq:vol2-absolute-bar-multiplication}
 b_2(sa,sb)=(-1)^{|a|}s(a\star b).
\end{equation}
The derivation property of $Q$ and associativity give
$(d_T+b_2)^2=0$.
This uses the underlying algebra with its unit forgotten;
it assumes no scalar augmentation of a Weyl algebra.

On tensor length $n$ extend the contraction by
\begin{equation}\label{eq:vol2-absolute-tensor-homotopy}
 i_T=i_s^{\otimes n},\quad p_T=p_s^{\otimes n},\qquad
 \mathbf H=
 \sum_{j=1}^n(i_sp_s)^{\otimes(j-1)}
                     \otimes h_V\otimes1^{\otimes(n-j)}.
\end{equation}
The identity $d_Vh_V+h_Vd_V=1-i_sp_s$ telescopes to
\[
 d_T\mathbf H+\mathbf H d_T=1-i_Tp_T,
 \qquad p_Ti_T=1,
 \qquad \mathbf H^2=\mathbf H i_T=p_T\mathbf H=0.
\]
For instance, in a mixed term of $\mathbf H^2$ the earlier
homotopy is followed by $i_sp_s$ or the later one is preceded
by it. The factors $h_Vi_s$ and $p_sh_V$ kill these terms;
the equal-position term contains $h_V^2$.

\begin{theorem}\label{thm:vol2-absolute-higher-model}
The finite-at-each-length formulas
\begin{align}\label{eq:vol2-absolute-transfer}
 b_Z&=p_Tb_2(1+\mathbf H b_2)^{-1}i_T,
       \notag\\
 \mathbf I&=(1+\mathbf H b_2)^{-1}i_T,
 &\mathbf P&=p_T(1+b_2\mathbf H)^{-1}
\end{align}
give a minimal $A_\infty$ structure on $Z$ and
$A_\infty$ comparison maps with $(\mathcal B_\hbar,Q)$.
The binary product is the ring in
Proposition~\ref{prop:vol2-absolute-ring}.
All components are polynomial over $R$ on polynomial inputs
and persist under coefficient base change.
The construction has a strictly unital normalized presentation.
\end{theorem}

\begin{proof}
Each occurrence of $b_2$ reduces tensor length by one.
Thus every inverse in \eqref{eq:vol2-absolute-transfer}
is a finite geometric sum on a tensor of fixed length $n$,
with at most $n-1$ occurrences of $b_2$.
Apply the perturbation identity of
Lemma~\ref{lem:vol2-filtered-perturbation} to the tensor
contraction. Its finite sums give
\[
 b_Z^2=0,
 \quad (d_T+b_2)\mathbf I=\mathbf I b_Z,
 \quad \mathbf P(d_T+b_2)=b_Z\mathbf P,
 \quad \mathbf P\mathbf I=1.
\]
The comparison homotopy is
$\mathbf H'=(1+\mathbf H b_2)^{-1}\mathbf H$.

We also need the coalgebra structures, not just these complex
identities. Cutting a tensor in
\eqref{eq:vol2-absolute-tensor-homotopy} on either side of
its homotopy factor gives
\begin{equation}\label{eq:vol2-absolute-homotopy-coproduct}
 \Delta\mathbf H
   =\bigl(\mathbf H\otimes1+i_Tp_T\otimes\mathbf H\bigr)\Delta.
\end{equation}
The maps $i_T,p_T$ are coalgebra maps, and $b_2$ is a
coderivation. Insert these three identities into the finite
series. Induction on the number of multiplications gives
\[
 \Delta\mathbf I=(\mathbf I\otimes\mathbf I)\Delta,
 \qquad
 \Delta\mathbf P=(\mathbf P\otimes\mathbf P)\Delta,
 \qquad
 \Delta b_Z=(b_Z\otimes1+1\otimes b_Z)\Delta.
\]
To describe the induction explicitly, cut each term at one
of its surviving tensor gaps. The multiplication is on the
left or right of that gap, so the two choices give the two
tensor factors of the preceding identities. A homotopy crossing
the cut splits as the two terms of
\eqref{eq:vol2-absolute-homotopy-coproduct}.
Terms with $\mathbf H i_T$, $p_T\mathbf H$ or
$\mathbf H^2$ vanish. The remaining left and right strings
are exactly the shorter geometric series. Their tensor signs
are inherited from the displayed coproduct identity.
This proves the induction at every length.

There is no unary component of $b_Z$, since the cohomology
differential is zero. The binary component is
$p_sb_2(i_s\otimes i_s)$, which is the computed cohomology
product. Projection to length one gives all higher components.
Each is a finite sum of products and the polynomial contraction,
so it has the asserted coefficient and specialization properties.

Finally $i(1)=1$, $\Pi(1)=1$ and $H(1)=0$.
In a higher tree with a unit input, the first multiplication
meeting that input removes the unit. An adjacent homotopy
then gives $Hi=0$ or $H^2=0$, or the final projection gives
$\Pi H=0$. This is the normalized strict-unit construction:
the binary unit laws remain, while higher operations and higher
comparison components with a unit input vanish.
It removes the tensors degenerate by a unit insertion, without
introducing an algebra augmentation.
\end{proof}

The formulas specify every operation once the relative weight
splittings are chosen. They do not assert that the higher
operations vanish. For the next calculation, keep the relative
section and homotopy fixed in
Section~\ref{sec:vol2-absolute-contraction}.

\subsection{A ternary operation from two trees}

The chosen full homotopy has
\begin{equation}\label{eq:vol2-absolute-tree-homotopies}
 H(\xi\star\eta)=r,
 \qquad
 H(\eta\star\chi)=\hbar\sum_i y_i\gamma^i.
\end{equation}
The first equality is the relative choice.
For the second, write
\[
 \eta\star\chi=-\chi\eta+\hbar\omega_y,
 \qquad
 \omega_y=y_1\gamma^2\gamma^3
          -y_2\gamma^1\gamma^3+y_3\gamma^1\gamma^2.
\]
The first term is a chosen cohomology representative.
On the coefficient spin-one summand,
$G=-1/2$. Direct rotation gives
$K\omega_y=-2\sum_i y_i\gamma^i$ and
$d_{\mathrm{CE}}(\sum_i y_i\gamma^i)=\omega_y$.
Thus the normalized Casimir homotopy sends $\omega_y$ to
$\sum_i y_i\gamma^i$, proving the second equality.

\begin{proposition}\label{prop:vol2-absolute-ternary}
For this contraction, the suspended ternary coefficient is
\begin{equation}\label{eq:vol2-absolute-ternary}
 (b_Z)_3\bigl(s(s\theta),s(t\theta),s\chi\bigr)
     =s\left(\hbar^2st+\frac12\hbar^3\right).
\end{equation}
Inside the parentheses, $s,t$ are the Cartan Weyl symbols.
The outer $s$ is suspension.
\end{proposition}

\begin{proof}
The three suspended inputs have even parity.
The length-three term of \eqref{eq:vol2-absolute-transfer}
is $-p_Tb_2\mathbf H b_2i_T$.
Using \eqref{eq:vol2-absolute-tree-homotopies} and the minus
sign in $h_V=-sHs^{-1}$, its two unsuspended contributions
have a plus sign and are
\[
 \Pi(r\star\chi),\qquad
 \Pi\left(\xi\star\hbar\sum_i y_i\gamma^i\right).
\]
The scalar part of
$\psi_1\psi_2\psi_3\gamma^1\gamma^2\gamma^3$
is $-\hbar^3$. All other terms have one or two positive
ghosts, or have three positive ghosts and the relative factor
$r$, whose projection is zero. Hence
\[
 \Pi(r\star\chi)=-\hbar^3.
\]
In the second tree, only the contractions
$\psi_i\gamma^j\mapsto\hbar\delta_i^j$
contribute to positive-ghost degree zero.
The Weyl product gives
$x_i\star y_i=x_iy_i+\hbar/2$.
The invariant sum $b=\sum_i x_iy_i$ restricts to $st$,
so this tree has projection
\[
 \hbar^2st+\frac32\hbar^3.
\]
Their sum proves the formula.
\end{proof}

Both the complete ghost contraction and the Weyl scalar term
are necessary. The coefficient is a property of this specified
minimal model. To use it as a formality obstruction, one would
have to show that its Hochschild class cannot be removed by
an $A_\infty$ change of coordinates. A single displayed value
does not prove that assertion.
The polynomial-subalgebra section in
Proposition~\ref{prop:vol2-relative-polynomial-section}
already shows how particular relative higher values can change
without changing the underlying quasi-isomorphism type.

\section{The central matrix direction}
\label{sec:vol2-absolute-central-direction}

The traceless calculation uses $\mathfrak{sl}_2$.
For $\mathfrak{gl}_2$, split off the scalar matrix direction
using the trace pairing. Normalize its coordinates $x_0,y_0$
so that $[x_0,y_0]=\hbar$.
Its moment function and adjoint action are zero.
The corresponding ghosts therefore have zero differential,
but their product is
\begin{equation}\label{eq:vol2-absolute-central-clifford}
 \psi_0\gamma^0+\gamma^0\psi_0=\hbar,
 \qquad (\gamma^0)^2=\psi_0^2=0,
 \qquad |\gamma^0|=1,\quad |\psi_0|=-1.
\end{equation}
In this split normalization, the full complex is
\[
 (\mathcal B_\hbar,Q)\otimes_R W_\hbar(x_0,y_0)
       \otimes_R\bigl(\mathrm{Cl}_\hbar(\gamma^0,\psi_0),0\bigr).
\]
Both additional factors are $R$-free with zero differential.
Tensoring the constructed contraction computes its cohomology
ring as the same two factors tensored with $Z$.

At $\hbar=0$ the central ghost product specializes to the
exterior product. Over $R$ it retains
\eqref{eq:vol2-absolute-central-clifford}.
After inverting $\hbar$, the central Clifford algebra acts as
the full endomorphism algebra of the two-dimensional graded
space $\Lambda(\gamma^0)$.
This may be used for a Morita comparison of module categories;
it does not identify the algebra with an exterior algebra by
a dg-algebra quasi-isomorphism.

\section{Absolute formality and the Hamiltonian comparison}
\label{sec:vol2-absolute-remaining-comparison}

The contraction determines the absolute complex, its ring
and one complete higher model over $\mathbb C[\hbar]$.
The formality question asks for a zigzag to the displayed
cohomology ring, or an obstruction invariant under every
allowed change of higher model. The first possible higher
product must be treated as a Hochschild class, including its
indeterminacy under such changes. The strict-map obstruction
and the chosen ternary value have their narrower conclusions.

The Hamiltonian comparison requires additional structure.
Its first component must agree with the global ordinary trace
map, remain strict on the scalar quadratic sector, and retain
the cubic primitive of
Section~\ref{sec:vol2-derived-action-bfv}.
The higher equations must be solved on this full absolute
recipient, with the negative relation classes, positive gauge
classes and cyclic operations present.
The matrix parameter $\hbar$ is independent of the
coefficient-plane parameter $\varepsilon$.
A linear contraction does not supply those equivariances.

Nor does the rank-two calculation give a higher-rank absolute
contraction or a field-theoretic realization.
Those comparisons must specify their own relation complexes,
supported operations and coefficient completions.
The explicit recipient constructed here makes their required
maps precise, including the classes that an ordinary
degree-zero reduction would lose.

\chapter{Cyclic traces and wheel corrections}
\label{chap:vol2-cyclic-traces-wheels}

A matrix trace remembers the size of the matrix:
$\operatorname{Tr}(I_N)=N$.
Invariant quantization also remembers a scalar. For the quadratic
$\mathfrak{sl}_2$ invariant, the Duflo operator adds
$\hbar^2/2$ to the symmetrized Casimir.
These scalars matter when matrices are stabilized or Lie algebras
are completed. Removing them can destroy compatibility with
brackets or continuity.

We first construct trace maps on cyclic chains and specify the
rank transitions they preserve. We then compute the finite wheel
corrections and examine two different infinite-dimensional
coefficient spaces. The cyclic orientation of
Chapter~\ref{ch:vol2-derived-hamiltonian-actions} remains a
chain-level structure throughout. A scalar state requires its own
functional.

\section{Matrix trace on the mixed complex}
\label{sec:vol2-matrix-mixed-trace}

Let $B$ be a unital associative complex algebra and put
$\overline B=B/\mathbb C1$.
Its normalized Hochschild chains are
\[
 C_n(B)=B\otimes\overline B^{\otimes n},
 \qquad a_0[a_1|\cdots|a_n]\in C_n(B).
\]
Place $C_n(B)$ in cohomological degree $-n$.
For $n>0$ the Hochschild boundary is
\begin{align*}
 b(a_0[a_1|\cdots|a_n])
 ={}&a_0a_1[a_2|\cdots|a_n]\\
 &+\sum_{j=1}^{n-1}(-1)^j
      a_0[a_1|\cdots|a_ja_{j+1}|\cdots|a_n]\\
 &+(-1)^na_na_0[a_1|\cdots|a_{n-1}],
\end{align*}
and $b=0$ on $C_0(B)$.
The Connes operator is
\[
 \mathsf B(a_0[a_1|\cdots|a_n])
 =\sum_{j=0}^n(-1)^{nj}
       1[a_j|\cdots|a_n|a_0|\cdots|a_{j-1}].
\]
Units in bar slots are zero. The identities
$b^2=\mathsf B^2=b\mathsf B+\mathsf Bb=0$ give the mixed
complex. With $|\zeta|=2$, its negative-cyclic complex is
\[
 CC^-(B)=\bigl(C_\bullet(B)[[\zeta]],b+\zeta\mathsf B\bigr).
\]
Here the completion is in $\zeta$. In particular, it does not
complete the coefficients of $B$.

For a second unital algebra $D$, define
\begin{equation}\label{eq:vol2-matrix-chain-trace}
 \mathsf T_N\bigl(M_0[M_1|\cdots|M_n]\bigr)
 =
 \sum_{i_0,\ldots,i_n=1}^N
 (M_0)_{i_0i_1}
 [(M_1)_{i_1i_2}|\cdots|(M_n)_{i_ni_0}],
\end{equation}
from matrix chains to chains of $D$.
For $n=0$ the formula is the ordinary matrix trace.
The indices form a closed path, so a change of the first index
does not change the sum.

\begin{proposition}\label{prop:vol2-mixed-matrix-trace}
The map $\mathsf T_N$ is a morphism of mixed complexes
$C_\bullet(\operatorname{Mat}_N D)\to C_\bullet(D)$.
It therefore extends continuously to negative-cyclic complexes.
It is natural in $D$ and additive under block sums of
representations. Its value on the unit in degree zero is $N1$.
\end{proposition}
\begin{proof}
A scalar identity in a bar slot forces the two adjacent
indices to agree. The corresponding output has a scalar unit
in that bar slot and is zero. Thus the formula is defined on
normalized chains.

An internal Hochschild face multiplies two adjacent matrices.
Expanding the product sums precisely over the shared index
in \eqref{eq:vol2-matrix-chain-trace}.
The first face has the same interpretation at the outer
coefficient. For the last face, reindex the closed path
cyclically. These identities prove
$b\mathsf T_N=\mathsf T_Nb$, with the displayed face signs.

In a Connes term the leading identity matrix forces its
adjacent indices to coincide. After deleting that repeated
index, the remaining sum is the cyclic rotation of the same
closed path. The rotation has the same sign $(-1)^{nj}$ on
both sides. Hence
$\mathsf B\mathsf T_N=\mathsf T_N\mathsf B$.
Coefficientwise application now commutes with
$b+\zeta\mathsf B$ and with reduction modulo every power of
$\zeta$.

For block-diagonal matrices, a nonzero closed path lies
entirely in one block. The trace of a block sum is consequently
the sum of the two chain traces. Naturality follows by
applying an algebra map to every matrix entry. Finally the
ordinary trace of the identity is $N$.
\end{proof}

For dg algebras use the raw totalization and graded cyclic
signs of
\eqref{eq:vol2-raw-hochschild-differential} and
\eqref{eq:vol2-raw-connes-operator}.
The same formula is a mixed-complex map: matrix indices
have degree zero, the internal differential acts entrywise,
and every cyclic permutation has the same Koszul sign
before and after the index sum.

This gives a concrete image of the orientation already
constructed on the free resolution.
Let
\[
 Q=\mathbb C\langle q,p,t\rangle,\qquad
 |t|=-1,\qquad dt=qp-pq,
\]
and let $R_N$ be its derived representation algebra:
the graded-commutative algebra on the entries of $X,Y,\Psi$,
with
$dX=dY=0$ and $d\Psi=XY-YX$.
Universal substitution is a dg algebra map
$Q\to\operatorname{Mat}_N(R_N)$.
The trace of its negative-cyclic orientation is
\begin{equation}\label{eq:vol2-traced-orientation}
 \omega_N=
 \sum_{i,j}
 \bigl(1[X_{ij}|Y_{ji}]-1[Y_{ij}|X_{ji}]\bigr)
 -\sum_i1[\Psi_{ii}].
\end{equation}
It is closed for both mixed differentials, by
Proposition~\ref{prop:vol2-mixed-matrix-trace}.
The relation term has total degree $-2$, just as the
two-bar terms do.
The matrix trace also commutes with entrywise derivations,
so it transports the coherent classical action and
orientation of
Theorem~\ref{thm:vol2-coherent-cyclic-orientation}.
These are constructions on derived representation
coefficients. The full quantum BRST and supported
field-theoretic comparisons retain their additional maps.

\section{Rank and scalar reduction}
\label{sec:vol2-reduced-stable-trace}

The degree-zero cyclic space of $B$ is
\[
 \mathfrak c_B=HC_0(B)=B/[B,B],
\]
where $[B,B]$ is the linear span of commutators.
Let $\operatorname{Rep}_NB$ represent unital algebra maps
$B\to\operatorname{Mat}_N(D)$ for commutative complex
algebras $D$, and put
\[
 \mathcal R_N=\mathcal O(\operatorname{Rep}_NB)^{GL_N}.
\]
The ordinary trace is the linear map
\[
 \operatorname{Tr}_N:\mathfrak c_B\longrightarrow\mathcal R_N,
 \qquad
 \operatorname{Tr}_N([b])(\rho)=\operatorname{Tr}(\rho(b)).
\]
The trace of a commutator is zero and conjugation preserves
the trace, proving that this map has the stated source
and target. It is generally not multiplicative.
Even on scalars, its value at $1$ is $N$.

Choose an augmentation $\epsilon:B\to\mathbb C$.
It specifies a rank transition
\[
 s_N:\operatorname{Rep}_NB\longrightarrow\operatorname{Rep}_{N+1}B,
 \qquad \rho\longmapsto\rho\oplus\epsilon.
\]
Define
\[
 \overline{\mathfrak c}_B=B/(\mathbb C1+[B,B]),
 \qquad
 t_N(\overline b)(\rho)=\operatorname{Tr}(\rho(b))-N\epsilon(b).
\]
Let $\mathfrak m_N$ be the ideal in $\mathcal R_N$ of functions
vanishing at $b\mapsto\epsilon(b)I_N$, and set
\[
 \widehat{\mathcal R}_N
   =\varprojlim_{r\geq1}\mathcal R_N/\mathfrak m_N^r.
\]
If $J$ is the ideal generated by
$\overline{\mathfrak c}_B$ in
$\operatorname{Sym}_{\mathbb C}\overline{\mathfrak c}_B$,
write
\[
 \widehat{\operatorname{Sym}}_J\overline{\mathfrak c}_B
 =\varprojlim_{r\geq1}
   \operatorname{Sym}_{\mathbb C}\overline{\mathfrak c}_B/J^r.
\]

\begin{proposition}\label{prop:vol2-reduced-trace-completion}
The maps $t_N$ are well defined, have image in
$\mathfrak m_N$, and satisfy $s_N^*t_{N+1}=t_N$.
They extend to a continuous homomorphism of commutative algebras
\[
 \widehat{\operatorname{Sym}}_J\overline{\mathfrak c}_B
 \longrightarrow \varprojlim_N\widehat{\mathcal R}_N.
\]
All maps are natural for augmentation-preserving algebra maps.
\end{proposition}
\begin{proof}
Both terms defining $t_N$ vanish on commutators, and their
values on a scalar multiple of $1$ cancel.
Their values also cancel at the augmentation representation.
Moreover,
\[
 s_N^*t_{N+1}(\overline b)
 =\operatorname{Tr}_N([b])+\epsilon(b)-(N+1)\epsilon(b)
 =t_N(\overline b).
\]
Since $s_N$ takes the augmentation representation to the
augmentation representation,
$s_N^*(\mathfrak m_{N+1}^r)\subseteq\mathfrak m_N^r$.
The symmetric-algebra extension of $t_N$ sends $J^r$ into
$\mathfrak m_N^r$. It therefore defines a map on each pair of
quotients and hence on their inverse limits.
The stabilization identity makes these completed maps
compatible in $N$. These formulas also prove naturality.
\end{proof}

The subtraction exists on cyclic chains as well.
For a representation $\rho:B\to\operatorname{Mat}_N(D)$,
subtract $N$ times the map induced by
$B\xrightarrow{\epsilon}\mathbb C\to D$ from
$\mathsf T_N\rho_*$.
Both summands are mixed-complex maps.
Their difference vanishes on the scalar subcomplex and is
unchanged by $\rho\mapsto\rho\oplus\epsilon$.
This chain-level reduction has a chosen augmentation as
part of its data.

A Poisson comparison requires more.
Suppose $\mathfrak c_B$ is supplied with a Lie bracket and
$\mathcal R_N$ with Poisson brackets satisfying
\begin{equation}\label{eq:vol2-trace-bracket-datum}
 \{\operatorname{Tr}_N\beta,\operatorname{Tr}_N\gamma\}_N
 =\operatorname{Tr}_N[\beta,\gamma]_{\mathrm{cyc}}.
\end{equation}
Then the ordinary trace is a Lie map.
A Hamiltonian action through it also requires a Lie map
$j:\mathfrak h\to\mathfrak c_B$.
Neither datum is attached to an arbitrary associative
algebra by its definition.
For the free algebra on two generators, the necklace bracket
and trace moment map provide
\eqref{eq:vol2-trace-bracket-datum};
the polynomial-to-necklace comparison still has the cubic
defect discussed in Chapter~\ref{ch:vol2-derived-hamiltonian-actions}.

Scalar reduction need not preserve this bracket.
For $B=\mathbb C[x,y]$ with $\{x,y\}=1$ and augmentation
at the origin, evaluation on $N$ points gives
\[
 \{\operatorname{Tr}_Nx,\operatorname{Tr}_Ny\}=N,
 \qquad t_N(\overline1)=0.
\]
Thus Proposition~\ref{prop:vol2-reduced-trace-completion}
asserts a commutative algebra map.
A reduced Poisson map needs a bracket-compatible reduction,
for example a selected Hamiltonian sector on which the
augmentation of every bracket vanishes.

\section{The stable multisymmetric algebra}
\label{sec:vol2-stable-multisymmetric}

The rank dependence can be retained explicitly.
For $N$ points of the affine symplectic plane put
\[
 S_N=\mathbb C[x_1,y_1,\ldots,x_N,y_N]^{\mathfrak S_N},
 \qquad
 p_{ab}^{(N)}=\sum_{i=1}^Nx_i^ay_i^b.
\]
Give $x_i,y_i$ degree one and, for $a+b>0$, give the
formal generator $p_{ab}$ degree $a+b$.
Evaluation at an extra point $(0,0)$ defines
$S_{N+1}\to S_N$.

\begin{theorem}\label{thm:vol2-stable-multisymmetric-ring}
In each fixed degree $d$, evaluation gives an isomorphism
\[
 \mathbb C[p_{ab}\mid a,b\geq0,\ a+b>0]_d
   \xrightarrow{\ \sim\ }(S_N)_d
 \qquad (N\geq d).
\]
Writing $P$ for this graded polynomial algebra, one has
\[
 \widehat P_{\mathrm{wt}}=\prod_{d\geq0}P_d
 \cong
 \varprojlim_N\prod_{d\geq0}(S_N)_d
\]
as complete commutative algebras.
\end{theorem}
\begin{proof}
A monomial orbit is specified by a finite multiset
$\lambda$ of nonzero pairs $(a,b)$, one pair for each
occupied point. Let $m_\lambda^{(N)}$ be the sum of the
distinct monomials in this orbit.
These sums form a basis when $\lambda$ has at most $N$
members. In total degree $d$, it has at most $d$ members.
The basis is therefore unchanged for $N\geq d$, and
evaluation at the added origin preserves every basis element.

For $\lambda=\{\alpha_1,\ldots,\alpha_r\}$ expand
$p_{\alpha_1}^{(N)}\cdots p_{\alpha_r}^{(N)}$.
Terms using distinct point indices give
\[
 \left(\prod_\alpha m_\alpha!\right)m_\lambda^{(N)},
\]
where $m_\alpha$ is the multiplicity of $\alpha$ in
$\lambda$. Every other term uses fewer than $r$ points,
merging the pairs assigned to the same point.
Ordering orbit sums by the number of occupied points
gives a triangular change of basis with nonzero diagonal.
It proves both generation and independence in degree $d$.

Multiplication adds total degrees. Completing degreewise
therefore gives $\widehat P_{\mathrm{wt}}$.
A compatible family in the right-hand limit determines,
in each degree, the unique stable coefficient just
constructed. Conversely those coefficients determine the
family. These inverse maps preserve multiplication and
the inverse-limit topologies.
\end{proof}

The completion in this theorem is specified by total degree.
It differs from the $J$-adic completion in which every
$p_{ab}$ has filtration one.
For example,
$\sum_{a\geq1}p_{a0}$ belongs to
$\widehat P_{\mathrm{wt}}$, but does not define an element
of the $J$-adic completion: its reduction modulo $J^2$
would be an infinite linear combination in a vector space
whose elements are finite linear combinations.
Nor does the theorem replace the target by the inverse
limit of the uncompleted rings $S_N$.
The series $\sum_{r\geq0}p_{10}^r$ already has a
nonpolynomial image at $N=1$.

These degree-zero trace functions do not record
the relation and gauge cochain degrees of a
derived matrix quotient. A comparison with a
derived centre requires representation and trace
maps on the specified complexes.

For the Poisson bracket, retain a central generator
$r=p_{00}$ and put
\[
 P^{\mathrm{rk}}=\mathbb C[r,p_{ab}\mid a+b>0].
\]
Define on generators
\begin{equation}\label{eq:vol2-rank-poisson}
 \{p_{ab},p_{cd}\}=(ad-bc)p_{a+c-1,b+d-1},
 \qquad \{r,-\}=0.
\end{equation}
A term with a negative index is zero; in that case its
coefficient is already zero.

\begin{proposition}\label{prop:vol2-rank-poisson}
Equation~\eqref{eq:vol2-rank-poisson} defines a Poisson
bracket. The evaluations
$P^{\mathrm{rk}}\to S_N$, with $r\mapsto N$, are Poisson.
Block sum is described by the Poisson coproduct
\[
 \Delta p_{ab}=p_{ab}\otimes1+1\otimes p_{ab},
 \qquad
 \Delta r=r\otimes1+1\otimes r.
\]
Origin stabilization is not Poisson on the full $S_N$.
On the subalgebras generated by $p_{ab}^{(N)}$ with
$a+b\geq2$, it is Poisson, and their stable algebra is
$\mathbb C[p_{ab}\mid a+b\geq2]$ with the same bracket.
\end{proposition}
\begin{proof}
The polynomial Poisson algebra $\mathbb C[x,y]$, regarded
only as a Lie algebra, has basis $x^ay^b$ including $1$.
Its symmetric algebra is $P^{\mathrm{rk}}$, where the
Lie element $1$ becomes the generator $r$, not the
multiplicative unit. The Lie--Poisson construction extends
its bracket by Leibniz and proves Jacobi.

Direct differentiation in $S_N$ gives
\[
 \{p_{ab}^{(N)},p_{cd}^{(N)}\}
 =(ad-bc)p_{a+c-1,b+d-1}^{(N)},\qquad
 p_{00}^{(N)}=N.
\]
This proves the evaluation assertion.
The primitive coproduct preserves a Lie bracket on
generators, so its extension is Poisson.
It also equals the power-sum formula for a block sum.

In rank $N+1$ the bracket of $p_{10}$ and $p_{01}$ is
$N+1$. Its pullback remains $N+1$, whereas the bracket
of the pullbacks is $N$. This proves the failure.
If both input degrees are at least two, the output
degree is at least two. The central rank never occurs,
and setting the last point to zero preserves every
displayed bracket. The triangular basis argument of
Theorem~\ref{thm:vol2-stable-multisymmetric-ring}, restricted
to these generators, proves the final assertion.
\end{proof}

\section{Stable matrix Lie homology}
\label{sec:vol2-stable-lie-homology}

There is a second stabilization, with a different direction.
Let
\[
 \mathfrak{gl}_\infty B=\varinjlim_N\mathfrak{gl}_N B
\]
under block-zero Lie embeddings, and use trivial
coefficients $\mathbb C$ for its Chevalley homology.
This construction requires no augmentation of $B$.
For clarity, ordinary cyclic homology is the homology
of the cyclic complex
\[
 CC(B)=
 \bigl(C_\bullet(B)((\zeta))/\zeta C_\bullet(B)[[\zeta]],
       b+\zeta\mathsf B\bigr),
 \qquad HC_n(B)=H^{-n}CC(B).
\]

The Loday--Quillen--Tsygan theorem identifies its
primitive part:
\[
 \operatorname{Prim}H_n^{\mathrm{CE}}(\mathfrak{gl}_\infty B;\mathbb C)
 \cong HC_{n-1}(B).
\]
Equivalently, if $HC_{n-1}(B)$ is placed in homological
degree $n$, it gives an isomorphism of graded Hopf algebras
\begin{equation}\label{eq:vol2-lqt-hopf}
 H_*^{\mathrm{CE}}(\mathfrak{gl}_\infty B;\mathbb C)
 \cong
 \operatorname{Sym}_{\mathrm{gr}}
       \left(\bigoplus_{n\geq1}HC_{n-1}(B)\right).
\end{equation}
Here $\operatorname{Sym}_{\mathrm{gr}}$ uses the Koszul
sign rule. This is the characteristic-zero theorem for
unital associative algebras of Loday and Quillen
\cite{LodayQuillen1984}*{Theorem~6.2 and Proposition~6.6}.
The product comes from block sum and the coproduct
from the diagonal Lie map.

The matrix-chain trace in
\eqref{eq:vol2-matrix-chain-trace} is also the trace
appearing in this theorem. The invariant-theory
description separates a permutation into its cycles:
one cycle gives a primitive cyclic chain, and disjoint
cycles give its graded symmetric products.
The theorem concerns this stable Lie homology, not the
inverse system of representation-coordinate rings.

For $B=\mathbb C$, normalized Hochschild chains are
concentrated in degree zero. Thus $HC_{2j}(\mathbb C)=\mathbb C$
and $HC_{2j+1}(\mathbb C)=0$, and
\[
 H_*^{\mathrm{CE}}(\mathfrak{gl}_\infty\mathbb C;\mathbb C)
 \cong\Lambda(u_1,u_3,u_5,\ldots),\qquad |u_{2j+1}|=2j+1.
\]
These classes remain although the scalar-reduced cyclic
space in Proposition~\ref{prop:vol2-reduced-trace-completion}
is zero.

\section{Cyclic functionals on the Weyl algebra}
\label{sec:vol2-weyl-cyclic-functional}

Quantization further separates cyclic chains from a scalar
trace. Let $R=\mathbb C[\hbar]$ and let $W_\hbar$ be the
polynomial Weyl algebra generated by $x,y$ with
$[x,y]=\hbar$.
Denote by $[W_\hbar,W_\hbar]_R$ the $R$-linear span
of commutators.

\begin{proposition}\label{prop:vol2-weyl-cyclic-quotient}
One has
\[
 [W_\hbar,W_\hbar]_R=\hbar W_\hbar,\qquad
 HH_0^R(W_\hbar)=W_\hbar/\hbar W_\hbar.
\]
Every $R$-linear cyclic functional $W_\hbar\to R$ is zero.
Every continuous $\mathbb C[[\hbar]]$-linear cyclic
functional on its $\hbar$-adic completion with values
in $\mathbb C[[\hbar]]$ is also zero.
\end{proposition}
\begin{proof}
Reduction modulo $\hbar$ is commutative, so every
commutator belongs to $\hbar W_\hbar$.
For every normal monomial,
\[
 [x,x^ay^{b+1}]=\hbar(b+1)x^ay^b.
\]
Since $b+1$ is invertible in $\mathbb C$, these
commutators span $\hbar W_\hbar$.
The Hochschild boundary in degree one is a commutator,
which gives the quotient formula.

If $\tau$ is cyclic and $R$-linear, then
$\hbar\tau(x^ay^b)=0$ for every $a,b$.
The target has no $\hbar$-torsion, so all these values
vanish. Continuity extends the same argument from
the dense polynomial algebra to its completion.
\end{proof}

The quotient in the proposition is nonzero and has
$\hbar$ acting as zero. It is the universal cyclic
target in degree zero, rather than a torsion-free
scalar ring. The negative-cyclic orientation of the
resolved plane has a different degree and universal
property. Neither it nor a prepared boundary state
is removed by the proposition.

\section{Finite formality and the Duflo map}
\label{sec:vol2-finite-duflo}

Let $\mathfrak g$ be a finite-dimensional complex Lie
algebra. Its linear Poisson tensor on $\mathfrak g^*$ is
$\pi_{\mathfrak g}$.
The polynomial polyvector complex, in unshifted
cochain degrees, is
\[
 \mathcal P_{\mathfrak g}
 =
 \left(
   \bigl(\operatorname{Sym}\mathfrak g
          \otimes\Lambda^\bullet\mathfrak g^*\bigr)[[\hbar]],
   \hbar d_{\mathrm{CE}}
 \right).
\]
The differential is
$[\hbar\pi_{\mathfrak g},-]_{\mathrm{Sch}}$.
Let $A_\hbar$ be the $\hbar$-adic completion of the
Rees enveloping algebra with relations
$xy-yx=\hbar[x,y]_{\mathfrak g}$.
Coefficients remain polynomial at each parameter order.

Fix the affine formality morphism $\mathcal U$ with
first component the normalized
Hochschild--Kostant--Rosenberg map:
\[
 X_1\wedge\cdots\wedge X_p
 \longmapsto
 \frac1{p!}\sum_{\sigma\in\mathfrak S_p}
       \operatorname{sgn}(\sigma)
       X_{\sigma(1)}\otimes\cdots\otimes X_{\sigma(p)}.
\]
We use Kontsevich's affine theorem
\cite{KontsevichDQ}*{Section~6.4}.
Twisting by $\hbar\pi_{\mathfrak g}$ gives a product
$\star$ on $\operatorname{Sym}\mathfrak g[[\hbar]]$.
There is an algebra isomorphism
\[
 T:(\operatorname{Sym}\mathfrak g[[\hbar]],\star)
       \longrightarrow A_\hbar
\]
fixing the linear generators.
Indeed their star commutator is
$\hbar[-,-]_{\mathfrak g}$.
The induced enveloping-algebra map has associated
graded the identity, and its filtered inverse is $T$.
This algebra map includes the change from the star
product to the PBW product.
Raw PBW symmetrization identifies the modules and
does not by itself identify these products.

The first component of the twisted comparison is
\begin{equation}\label{eq:vol2-finite-formality-component}
 \Phi_{\mathfrak g}(v)
 =T_*\sum_{n\geq0}\frac1{n!}
       \mathcal U_{n+1}
       ((\hbar\pi_{\mathfrak g})^{\otimes n},v).
\end{equation}
It maps $\mathcal P_{\mathfrak g}$ to relative
continuous Hochschild cochains of $A_\hbar$ over
$\mathbb C[[\hbar]]$.
The target differential is the Hochschild differential
of the quantized product.
All terms are defined at each parameter order.
The affine theorem proves the twisted
$L_\infty$ identities on the complexes shifted by $[1]$.

The linear component is a quasi-isomorphism.
Modulo $\hbar$ it is the polynomial HKR map, whose
comparison with the full Hochschild complex is given
by the polynomial diagonal Koszul resolution.
Choose a contracting homotopy of its cone over
$\mathbb C$ and extend it coefficientwise.
The difference between the complete cone differential
and this undeformed one raises $\hbar$ order.
The convergent perturbation series of
Lemma~\ref{lem:vol2-filtered-perturbation}
therefore contracts the complete cone.
This argument also explains why the scaled
Chevalley differential cannot be omitted.

Invariant polynomials have a more explicit comparison.
For $x\in\mathfrak g$, define the formal power series
\begin{equation}\label{eq:vol2-duflo-factor}
 j_{\mathfrak g}^{1/2}(x)
 =
 \det_{\mathfrak g}^{1/2}
 \left(
   \frac{\sinh(\operatorname{ad}_x/2)}
        {\operatorname{ad}_x/2}
 \right),
\end{equation}
with constant term one.
Interpret its Taylor coefficients as
constant-coefficient differential operators on
$\operatorname{Sym}\mathfrak g$.
The Duflo theorem gives the algebra isomorphism
\begin{equation}\label{eq:vol2-duflo-map}
 \operatorname{Duf}_\hbar
 =\operatorname{sym}_\hbar\circ
       j_{\mathfrak g}^{1/2}(\hbar\partial):
 (\operatorname{Sym}\mathfrak g)^{\mathfrak g}[[\hbar]]
 \xrightarrow{\ \sim\ } Z(A_\hbar);
\end{equation}
see \cite{KontsevichDQ}*{Section~8.3.4}.
On a polynomial, only finitely many differential
orders act nontrivially.
Its polynomial Rees formula may consequently be
specialized at $\hbar=1$.
This specialization is not defined on arbitrary
formal power series in $\hbar$.

For $\mathfrak g=\mathfrak h_{\leq d}$, these results
apply to the finite Hamiltonian jets of
Chapter~\ref{ch:vol2-hamiltonian-jets}.
The relative bar--Chevalley comparison, the
$L_\infty$ formality map and the invariant algebra
isomorphism in \eqref{eq:vol2-duflo-map} preserve
different specified operations.
Their conclusions must be used with those operations.

The Hochschild brace action on $(C^\bullet(A_\hbar),A_\hbar)$
can be transported when the required operadic comparison
is supplied. In detail, a homotopy-$E_2$ refinement is
a quasi-isomorphism or a zigzag of quasi-isomorphisms
from the selected complete homotopy-$E_2$ structure
on $\mathcal P_{\mathfrak g}$ to the relative
Hochschild complex, refining the cochain comparison.
A mixed refinement also specifies the open operations
and the coherences between the two colors.

\begin{corollary}[Open--closed structure with a refinement]
\label{cor:vol2-open-closed-refinement}
Given these refinements, the pair
$(\mathcal P_{\mathfrak g},A_\hbar)$ has the transported
homotopy Swiss-cheese action.
Its open product is the enveloping product.
Its closed and mixed operations are those carried by
the specified operadic comparison.
\end{corollary}
\begin{proof}
Transport the two-colored operations through the
supplied equivalence of complete operadic algebras.
The comparison includes the homotopies expressing
substitution and the open--closed identities, so
those identities hold on the transported pair.
\end{proof}

\section{The quadratic and quartic corrections}
\label{sec:vol2-quadratic-quartic-wheels}

For $\mathfrak{sl}_2$, use generators $e,f,h$ with
\[
 [h,e]=2e,\qquad [h,f]=-2f,\qquad [e,f]=h.
\]
In the symmetric algebra set $q=2ef+h^2/2$.
In the Rees enveloping algebra write the generators
as $E,F,H$ and put
$\Omega=EF+FE+H^2/2$.
The relations have an additional factor $\hbar$.
They imply that $\Omega$ is central:
for example,
$[E,EF+FE]=\hbar(EH+HE)$ and
$[E,H^2/2]=-\hbar(EH+HE)$.

\begin{proposition}\label{prop:vol2-duflo-quadratic}
With these normalizations,
\[
 \operatorname{Duf}_\hbar(q)=\Omega+\frac{\hbar^2}{2}.
\]
At $\hbar=1$, its eigenvalue on the irreducible
highest-weight-$n$ module is $(n+1)^2/2$.
\end{proposition}
\begin{proof}
For $x=ue+vf+wh$, the adjoint matrices give
\[
 \operatorname{Tr}(\operatorname{ad}_x^2)=8uv+8w^2.
\]
The first nonconstant term of
\eqref{eq:vol2-duflo-factor} is
$\operatorname{Tr}(\operatorname{ad}_x^2)/48$.
Thus its differential operator on $q$ contributes
\[
 \frac{\hbar^2}{6}
       (\partial_e\partial_f+\partial_h^2)
       (2ef+h^2/2)=\frac{\hbar^2}{2}.
\]
Higher differential orders annihilate $q$, and
symmetrization sends $q$ to $\Omega$.
On a highest-weight vector at $\hbar=1$, $EF$ acts
by $n$, $FE$ by zero and $H^2/2$ by $n^2/2$.
Adding $1/2$ gives $(n+1)^2/2$.
Centrality gives this value on the full irreducible module.
\end{proof}

The next invariant distinguishes a multiplicative
quantization from raw symmetrization.
Use complexified Hermitian generators $J_i$ with
\[
 [J_i,J_j]=i\hbar\epsilon_{ijk}J_k,
 \qquad
 p=x_1^2+x_2^2+x_3^2,\qquad C=\sum_iJ_i^2.
\]
Here $x_i$ are commuting symbols.
The change
$E=J_1+iJ_2$, $F=J_1-iJ_2$, $H=2J_3$
gives $q=2p$ and $\Omega=2C$.
Write $\Delta=\sum_i\partial_{x_i}^2$.
The adjoint eigenvalues are $0,\sqrt p,-\sqrt p$,
so the differential factor becomes
\[
 \frac{\sinh(\hbar\sqrt\Delta/2)}
      {\hbar\sqrt\Delta/2}
 =1+\frac{\hbar^2}{24}\Delta
     +\frac{\hbar^4}{1920}\Delta^2+\cdots.
\]
The square root denotes the displayed power series
in $\Delta$.

\begin{proposition}\label{prop:vol2-duflo-quartic}
In the universal Rees algebra,
\begin{align}
 \operatorname{sym}(p^2)&=C^2-\frac{\hbar^2}{3}C,
       \label{eq:vol2-pbw-quartic}\\
 \operatorname{Duf}_\hbar(p)&=C+\frac{\hbar^2}{4},
       \label{eq:vol2-duflo-compact-quadratic}\\
 \operatorname{Duf}_\hbar(p^2)
       &=\left(C+\frac{\hbar^2}{4}\right)^2.
       \label{eq:vol2-duflo-compact-quartic}
\end{align}
\end{proposition}
\begin{proof}
Antisymmetry of $\epsilon_{ijk}$ proves that $C$ is
central. The commutation relations give
\[
 \sum_j[J_j,[J_j,J_i]]=2\hbar^2J_i,
 \qquad
 \sum_jJ_jJ_iJ_j=(C-\hbar^2)J_i.
\]
The second identity follows by expanding the first.
Symmetrizing the two repeated indices in $p^2$ gives
\[
 \operatorname{sym}(p^2)
 =\frac13\sum_{i,j}
 \bigl(J_i^2J_j^2+J_iJ_jJ_iJ_j+J_iJ_j^2J_i\bigr).
\]
The three sums are $C^2$, $C^2-\hbar^2C$ and $C^2$,
respectively. This proves \eqref{eq:vol2-pbw-quartic}.

Direct differentiation gives
$\Delta p=6$, $\Delta p^2=20p$ and
$\Delta^2p^2=120$.
The quadratic correction is therefore $\hbar^2/4$.
For $p^2$, the full expression is
\[
 C^2-\frac{\hbar^2}{3}C
      +\frac{5\hbar^2}{6}C+\frac{\hbar^4}{16}
 =\left(C+\frac{\hbar^2}{4}\right)^2.
\]
\end{proof}

Thus the product of two operators assigned by $p\mapsto C$
differs from raw symmetrization of $p^2$ by
$\hbar^2C/3$.
This is an operator identity before any state is chosen.
In spin $j$, the eigenvalues of $C$ and
$\operatorname{Duf}_\hbar(p)$ are
$\hbar^2j(j+1)$ and $\hbar^2(j+1/2)^2$.
At $\hbar=1$, evolution by the second operator divided
by two differs from evolution by $C/2$ by $e^{-t/8}$.
Both factors obey the semigroup law, so that law alone
does not choose the vacuum energy.

A multiplicative invariant map with $p\mapsto C$ is
instead
\[
 Q(F(p))=
 \operatorname{Duf}_\hbar
       \bigl(F(p-\hbar^2/4)\bigr).
\]
The shift is an automorphism of the invariant polynomial
algebra. This formula makes the ordering choice explicit;
its relation to an independently prepared state requires
the state comparison as well.

\section{Half-densities and matrix wheels}
\label{sec:vol2-half-density-wheels}

The same formal determinant has a finite-dimensional
geometric meaning. Let $G$ be a local complex Lie group
with unimodular Lie algebra $\mathfrak g$.
For a dimensionless exponential coordinate $X$ near zero,
put $L=\operatorname{ad}_X$.
Left trivialization gives
\[
 d\exp_X=\int_0^1e^{-tL}\,dt
       =e^{-L/2}\frac{\sinh(L/2)}{L/2}.
\]
Normalize the tangent volume and left-invariant volume
to agree at the identity.
Since $\operatorname{Tr}L=0$, the exponential Jacobian
and its normalized square root are
\[
 j(X)=\det\frac{\sinh(L/2)}{L/2},
 \qquad
 \mathfrak j(X)=
 \exp\left(\frac12\operatorname{Tr}
                  \log\frac{\sinh(L/2)}{L/2}\right).
\]
Expanding the scalar series yields
\begin{equation}\label{eq:vol2-log-wheel-coefficients}
 \log\mathfrak j(X)
 =\frac{\operatorname{Tr}L^2}{48}
  -\frac{\operatorname{Tr}L^4}{5760}
  +\frac{\operatorname{Tr}L^6}{362880}
  +O(X^8).
\end{equation}
Each trace is the finite contraction of the adjoint
indices around a wheel.

For $\mathfrak g=\mathfrak{gl}_N$, let
$p_a=\operatorname{tr}(X^a)$, including $p_0=N$.

\begin{proposition}\label{prop:vol2-matrix-wheel-traces}
For every $m\geq0$,
\[
 \operatorname{Tr}_{\mathfrak{gl}_N}(\operatorname{ad}_X)^m
 =\sum_{a=0}^m(-1)^a\binom ma p_{m-a}p_a.
\]
In particular,
\[
 \operatorname{Tr}L^2=2Np_2-2p_1^2,\qquad
 \operatorname{Tr}L^4=2Np_4-8p_1p_3+6p_2^2.
\]
For matrices $X,Y$ of sizes $N,M$, the quadratic part of
$\log\mathfrak j(X\oplus Y)-\log\mathfrak j(X)
 -\log\mathfrak j(Y)$ is
\[
 \frac1{24}
 \left(M\operatorname{tr}X^2+N\operatorname{tr}Y^2
             -2\operatorname{tr}X\operatorname{tr}Y\right).
\]
\end{proposition}
\begin{proof}
On $\operatorname{End}(\mathbb C^N)$, left and right
multiplication commute and their difference is
$\operatorname{ad}_X$.
The matrix-unit basis gives
\[
 \operatorname{Tr}(L_X^aR_X^b)
    =\operatorname{tr}(X^a)\operatorname{tr}(X^b).
\]
Expand $(L_X-R_X)^m$ by the binomial theorem.
The cases $m=2,4$ give the stated polynomials.
Their derivation requires no diagonalizability.
Finally the quadratic logarithm is
$(N\operatorname{tr}X^2-(\operatorname{tr}X)^2)/24$;
subtracting its two block values gives the last formula.
\end{proof}

The off-diagonal matrix blocks are responsible for
this failure of additivity. An adjoint wheel density
is therefore different from a primitive matrix trace.

There is also a convergent product for this
finite-dimensional half-density.
On a sufficiently small neighborhood of zero,
\begin{equation}\label{eq:vol2-half-density-mode-product}
 \mathfrak j(X)=\lim_{K\to\infty}
 \prod_{n=1}^K
       \det\left(1+\frac{L^2}{4\pi^2n^2}\right)^{1/2},
\end{equation}
with branches equal to one at zero.
The logarithms are locally bounded by a constant times
$n^{-2}$, so the product converges locally uniformly.
The Euler product for $\sinh(z/2)/(z/2)$, applied to
the eigenvalues or its matrix power series, identifies
the limit with the determinant above.
If $L$ is real diagonalizable with real eigenvalues,
the partial product $\mathfrak j_K$ satisfies
\[
 0\leq\log\mathfrak j-\log\mathfrak j_K
 \leq\frac{\operatorname{Tr}L^2}{8\pi^2K}.
\]
Indeed $0\leq\log(1+u)\leq u$ for $u\geq0$, and
$\sum_{n>K}n^{-2}\leq1/K$.
For imaginary eigenvalues the scalar factors are
sine factors; the normalized local branch is taken
before their first zeros.
This product has a fixed finite-dimensional adjoint
space. It does not sum traces over infinitely many
independent Lie algebra blocks.

\section{A two-step colour algebra}
\label{sec:vol2-two-step-duflo}

A family in which every wheel vanishes distinguishes
the correction from the underlying comparison.
Let $W$ be finite dimensional and let
$C\in\Lambda^3W^*$.
In dual bases define
\[
 \mathfrak n_C=W\oplus W^*,\qquad
 [x_i,x_j]=\sum_kC_{ijk}y^k,\qquad
 [x_i,y^j]=[y^i,y^j]=0.
\]
The image of every bracket is central, so Jacobi holds
and $\operatorname{ad}_a\operatorname{ad}_b=0$.
The pairing $\kappa(x_i,y^j)=\delta_i^j$, zero on
the other two pairs of summands, is invariant because
$C$ is alternating.

Over $R=\mathbb C[\hbar]$, the Rees enveloping algebra
$A_C$ has the symmetric PBW product
\begin{equation}\label{eq:vol2-colour-pbw-product}
 f*_Cg
 =m\exp\left(
       \frac{\hbar}{2}\sum_{i,j,k}
          C_{ijk}Y^k\,\partial_{X_i}\otimes\partial_{X_j}
       \right)(f\otimes g).
\end{equation}
Every polynomial pair gives a finite sum.
The three bidifferential operators acting on pairs
of slots in a triple tensor product commute, since
their coefficients are central $Y$'s and their
derivatives act only on $X$'s. Both bracketings are
therefore the exponential of the same sum.
The generator commutators are
$[X_i,X_j]_{*_C}=\hbar C_{ijk}Y^k$, as required.

The adjoint coefficient Chevalley complex is
\[
 P_C=R[X,Y]\otimes\Lambda(\xi,\eta),
 \qquad |\xi_i|=|\eta^k|=1,
\]
with differential
\[
 Q\xi_i=QY^k=0,\qquad
 Q\eta^k=-\frac12C_{ijk}\xi_i\xi_j,\qquad
 QX_j=C_{ijk}Y^k\xi_i.
\]
Repeated indices are summed.
These are the Chevalley formulas, and Jacobi gives
$Q^2=0$.
The relative Hochschild cohomology of $A_C$ is
computed by
\[
 Z^p=\operatorname{Hom}_R
       (\Lambda_R^p(\mathfrak n_C\otimes R),A_C).
\]
Its differential uses the commutator action and
$\hbar$ times the bracket faces.
The unnormalized alternating Chevalley-to-bar map
of Chapter~\ref{ch:vol2-hamiltonian-jets}
gives this comparison.

Symmetrization is equivariant:
$[T_a,\operatorname{sym}_\hbar f]
  =\hbar\operatorname{sym}_\hbar(a\cdot f)$.
Consequently
\begin{equation}\label{eq:vol2-colour-beta}
 \beta_p:P_C^p\longrightarrow Z^p,\qquad
 \beta_p=\hbar^p\operatorname{sym}_\hbar
\end{equation}
is a cochain map.
The target differential contributes the extra factor
$\hbar$ needed to pass from degree $p$ to degree $p+1$.

\begin{proposition}\label{prop:vol2-colour-duflo}
Symmetrization identifies the invariant polynomial
algebra with the centre of $A_C$ over $R$.
After extending to $R[\hbar^{-1}]$, $\beta$ induces
an isomorphism of graded cup-product algebras
\[
 H^*(P_C)[\hbar^{-1}]
   \xrightarrow{\ \sim\ }HH_R^*(A_C)[\hbar^{-1}].
\]
The second assertion need not be an integral
isomorphism in positive degree.
\end{proposition}
\begin{proof}
If $f$ is invariant, then
$\sum_{i,k}C_{jik}Y^k\partial_{X_i}f=0$ for every $j$.
The bidifferential operator in
\eqref{eq:vol2-colour-pbw-product} therefore kills
$f\otimes g$ for every $g$, as do all its positive
powers. Hence $f*_Cg=fg$.
Conversely the PBW commutator of a linear generator
with $f$ is $\hbar$ times its adjoint derivative.
The PBW module has no $\hbar$-torsion, so centrality
is equivalent to invariance. This proves the
integral degree-zero assertion.

After inverting $\hbar$, identify $A_C$ with
$U(\mathfrak n_C)\otimes R[\hbar^{-1}]$ by
$T_a\mapsto\hbar a$.
Evaluate a degree-$p$ cochain on the rescaled
arguments $\hbar^{-1}T_a$; this contributes
$\hbar^{-p}$ and cancels the factor in $\beta_p$.
The target becomes the ordinary adjoint Chevalley
complex. Under these identifications, $\beta$ is
symmetrization composed with the source automorphism
$f(X,Y)\mapsto f(\hbar X,\hbar Y)$.
This automorphism commutes with $Q$ and preserves
the cup product.
Every adjoint product of length at least two is
zero, so the determinant
\eqref{eq:vol2-duflo-factor} equals one.
The cohomological Duflo theorem of Pevzner and
Torossian \cite{PevznerTorossianDuflo}*{Theorem~2.5}
then proves the cup-product assertion in every
degree. The polynomial argument extends over the
Laurent coefficient ring.

For the last assertion take $C=0$.
Both cochain differentials are zero, and $\beta_p$
is multiplication by $\hbar^p$ on a nonzero free
module whenever $0<p\leq\dim\mathfrak n_C$.
Its cokernel has $\hbar$-torsion and its
specialization at $\hbar=0$ is zero.
\end{proof}

There are algebraic product homotopies after this
localization, although the theorem does not
choose a field-theoretic one.
Here is an explicit way to specify the first.
Let $X=P_C\otimes P_C$ and $Y=Z$, now over
$R[\hbar^{-1}]$, and write
\[
 D=\beta\,m_{P_C}-m_Z(\beta\otimes\beta):X\longrightarrow Y.
\]
The cup product on $Z$ uses the usual Chevalley
shuffle diagonal and the coefficient product.
Choose contractions $(i_X,p_X,h_X)$ and
$(i_Y,p_Y,h_Y)$ onto cohomology.
They exist here: the preceding rescalings identify
the complexes with extensions of complexes of
complex vector spaces, which split into cohomology
and contractible pairs.
Multiplicativity on cohomology says $p_YDi_X=0$.
Then
\begin{equation}\label{eq:vol2-colour-cup-homotopy}
 H=h_YD+i_Yp_YDh_X
 \quad\hbox{satisfies}\quad d_YH+Hd_X=D.
\end{equation}
To verify it, use that $D$ is a cochain map and
expand the two contraction identities.
The first summand gives $(1-i_Yp_Y)D$.
The second gives $i_Yp_YD$, since
$p_YDd_X=0$ and $p_YDi_X=0$.
Their sum is $D$.
This construction depends on the splittings.
Compatibility with braces, geometric collisions
and a prescribed two-point operation is further
structure.

\section{A mode completion that excludes the finite normalization}
\label{sec:vol2-wheel-mode-obstruction}

Finite Duflo maps do not automatically assemble along
projections that discard Lie algebra blocks.
For each $i\geq1$, let $\mathfrak g_i$ be a copy
of $\mathfrak{sl}_2$ with generators $E_i,F_i,H_i$.
Different blocks commute. Put
$\mathfrak g_{\leq N}=\bigoplus_{i=1}^N\mathfrak g_i$
and define
\[
 A_{\mathrm{mode}}=
 \varprojlim_{N,M}
       \operatorname{Sym}(\mathfrak g_{\leq N})
                  \otimes\mathbb C[\hbar]/(\hbar^M),
\]
\[
 U_{\mathrm{mode}}=
 \varprojlim_{N,M}
       U_\hbar(\mathfrak g_{\leq N})/(\hbar^M).
\]
The transition maps kill the omitted blocks.
Each finite quotient has the discrete linear topology,
and the limits have their inverse-limit topologies.

\begin{proposition}\label{prop:vol2-mode-duflo-obstruction}
There is no continuous linear map
$A_{\mathrm{mode}}\to U_{\mathrm{mode}}$ whose
restriction to invariant polynomials in every finite
set of blocks is the Duflo map with the normalization
of \eqref{eq:vol2-duflo-map}.
\end{proposition}
\begin{proof}
The invariant
$c_i=H_i^2/2+2E_iF_i$ tends to zero in
$A_{\mathrm{mode}}$: every fixed block projection
eventually kills it.
Proposition~\ref{prop:vol2-duflo-quadratic} prescribes
its image as
\[
 \operatorname{sym}_\hbar(c_i)+\frac{\hbar^2}{2}.
\]
Project to the first block and reduce modulo
$\hbar^3$. For every $i>1$ the PBW term is zero,
while $\hbar^2/2$ is nonzero.
The proposed images do not tend to zero, contradicting
continuity.
\end{proof}

The obstruction fixes both the topology and the
normalization on every finite block.
It does not exclude a different coefficient domain
or a comparison with a specified renormalization.
For example, using a parameter $q_i$ of increasing
positive weight as the $i$th Rees parameter replaces
the scalar correction by $q_i^2/2$.
Only finitely many such corrections survive at any
fixed weight. The next construction makes that
coefficient domain precise.

\section{Positive weights and mixed coefficients}
\label{sec:vol2-weighted-mixed-formality}

Let $A=\mathbb C[x_1,x_2,\ldots]$, where each polynomial
uses finitely many variables.
Write $T_{\mathrm{fin}}^{p-1}$ for finite sums of
polynomial $p$-vectors and
$D_{\mathrm{fin}}^{p-1}$ for finite sums of operators
\[
 (f_1,\ldots,f_p)\longmapsto
 a(x)\,\partial^{\alpha_1}f_1\cdots\partial^{\alpha_p}f_p.
\]
Every multi-index has finite support.
In degree $-1$, both spaces are $A$.
These are the shifted deformation complexes:
a $p$-vector or $p$-ary cochain has degree $p-1$.
Use the Gerstenhaber insertion bracket on $D_{\mathrm{fin}}$.
If $m$ is ordinary multiplication, its differential is
$b_0=[m,-]_G$.
For the Schouten bracket use
\[
 [X_1\wedge\cdots\wedge X_p,f]_{\mathrm{Sch}}
 =\sum_i(-1)^{i-1}X_i(f)
          X_1\wedge\cdots\widehat X_i\cdots\wedge X_p.
\]

Take a mixed coordinate domain $U=I^a\times D^b$,
with $I$ an open real interval and $D$ a complex disk.
Let $K(U)$ be the algebra of smooth real de Rham
forms and complex $(0,*)$-forms on this product,
with
$d_K=d_{\mathbb R^a}+\bar\partial$
and the total-complex signs.
Its coefficients are arbitrary smooth functions
on the product, not just finite sums of separated
functions. Give it the compact-open $C^\infty$ topology.

Give each finite polynomial operator space the fine
bornology: a bounded set is contained and bounded
in a finite-dimensional subspace.
In $K(U)\widehat\otimes T_{\mathrm{fin}}$ and
$K(U)\widehat\otimes D_{\mathrm{fin}}$, this means
the direct sum of copies of $K(U)$ indexed by a
polynomial operator basis; a bounded set uses only
finitely many basis directions.
Field derivatives act on the $x_i$ only.
No spacetime differential operator or product of
distributions is included in this definition.

Choose degree-zero central parameters $q_i$ with
positive integer weights $w_i$, only finitely many
of which lie below any fixed bound.
Set
\[
 \operatorname{wt}(q^\alpha)=\sum_iw_i\alpha_i,
 \qquad
 \mathscr R_{<M}
   =\mathbb C[q_1,q_2,\ldots]/
       (q^\alpha:\operatorname{wt}(q^\alpha)\geq M),
 \qquad
 \mathscr R=\varprojlim_M\mathscr R_{<M}.
\]
Each $\mathscr R_{<M}$ is finite dimensional.
For the parameter-completed tensor define
$V\widehat\otimes\mathscr R
 =\varprojlim_M(V\otimes\mathscr R_{<M})$ and put
\[
 \widehat T(U)=
   (K(U)\widehat\otimes T_{\mathrm{fin}})
                 \widehat\otimes\mathscr R,\qquad
 \widehat D(U)=
   (K(U)\widehat\otimes D_{\mathrm{fin}})
                 \widehat\otimes\mathscr R.
\]
The filtration $F^M$ is the kernel of projection
to parameter precision $M$.
Every coefficient has finite field support,
although that support can increase with weight.

Finite polynomial differential graphs act on these
spaces. On a bounded input set, only finitely many
coefficient and derivative directions occur.
A graph consequently has finitely many index
assignments. Its spacetime coefficients are products,
and Leibniz's rule bounds their derivatives on each
compact subset by finitely many products of input
seminorms. Thus the graph map is bounded.
A series whose $n$th term raises weight by at least
$n$ has only finitely many terms at precision $M$,
provided that each term has finitely many graphs.
Its compatible finite-precision maps define a
continuous map on the completion.

Extend the affine Taylor maps over the mixed
coefficients by
\begin{equation}\label{eq:vol2-mixed-formality-sign}
 \mathcal U_n^K(k_1\gamma_1,\ldots,k_n\gamma_n)
 =
 (-1)^{\sum_{i<j}|\gamma_i||k_j|
              +(1-n)\sum_i|k_i|}
 k_1\cdots k_n\,\mathcal U_n(\gamma_1,\ldots,\gamma_n).
\end{equation}
The first exponent moves the coefficients to the
left. The second is the tensor sign for a map of
degree $1-n$.
After suspension, the Taylor maps have degree zero.
Their coefficient extensions commute with the
total coderivations because $d_K$ is a derivation.
Undoing suspension proves all the extended
$L_\infty$ identities with the displayed sign.
The finite graph formulas are stable under
adjoining unused field coordinates.

Let $\Pi\in F^1\widehat T(U)$ be a bivector of
mixed-form degree zero satisfying
$d_K\Pi=0$ and $[\Pi,\Pi]_{\mathrm{Sch}}=0$.
Independent finite Poisson blocks $\pi_i$ give the
example $\Pi=\sum_iq_i\pi_i$.

\begin{theorem}\label{thm:vol2-weighted-relative-formality}
The formulas
\begin{align}
 B_\Pi&=\sum_{n\geq1}\frac1{n!}
              \mathcal U_n^K(\Pi^{\otimes n}),
       &m_\star&=m+B_\Pi,
       \label{eq:vol2-weighted-star}\\
 \mathcal F_r(v_1,\ldots,v_r)
       &=\sum_{n\geq0}\frac1{n!}
           \mathcal U_{n+r}^K
                (\Pi^{\otimes n},v_1,\ldots,v_r),
       &&r\geq1
       \label{eq:vol2-weighted-twist}
\end{align}
define an associative product commuting with $d_K$
and a filtered continuous $L_\infty$ quasi-isomorphism
\[
 (\widehat T(U),d_K+[\Pi,-]_{\mathrm{Sch}})
 \longrightarrow
 (\widehat D(U),d_K+[m_\star,-]_G).
\]
The maps commute with restriction in $U$.
The assertion is about the unreduced mixed complexes.
\end{theorem}
\begin{proof}
At precision $M$, the twisting sums have $n<M$.
The input coefficients and the relevant coefficients
of $\Pi$ use a common finite set of field variables.
The calculation is therefore the finite affine
formality calculation with coefficients in $K(U)$
and signs \eqref{eq:vol2-mixed-formality-sign}.
The Maurer--Cartan equation gives
\[
 (d_K+[m,-]_G)B_\Pi+\frac12[B_\Pi,B_\Pi]_G=0.
\]
Since $B_\Pi$ is binary and has mixed-form degree
zero, this equation gives separately
$d_KB_\Pi=0$ and $[m_\star,m_\star]_G=0$.

For the higher maps, conjugate the source symmetric
coalgebra codifferential by multiplication with
$\exp(\Pi)$ and the target one by $\exp(B_\Pi)$.
The coefficient with $r$ free inputs is
\eqref{eq:vol2-weighted-twist}.
All coefficients and identities are finite at
the specified precision. Compatibility in $M$
proves them on the inverse limit.

The untwisted finite-direction HKR map is a
quasi-isomorphism by the filtered colimit of the
finite polynomial HKR comparisons.
Its cone is an acyclic complex of complex vector
spaces. Choose a square-zero contracting
homotopy $H_0$ by splitting it into contractible
pairs. Every linear map between fine bornological
spaces is bounded.
Extend $H_0$ over $K(U)$ with its degree-minus-one
tensor sign. The two terms involving $d_K$ cancel,
so this contracts the unreduced mixed cone.

Write the cone differential of $\mathcal F_1$ as
$d_0+\delta$. Every term of $\delta$ raises
parameter weight.
The homotopy
\[
 H=H_0(1+\delta H_0)^{-1}
   =H_0\sum_{j\geq0}(-\delta H_0)^j
\]
is a finite sum at every precision.
The perturbation identity gives
$(d_0+\delta)H+H(d_0+\delta)=1$.
This proves the quasi-isomorphism and continuity.
Finally the graph maps use only field derivatives
and pointwise products of mixed coefficients,
which commute with restriction in $U$.
\end{proof}

Both sign contributions in
\eqref{eq:vol2-mixed-formality-sign} can occur.
On an interval with coordinate $t$, take
$\pi=\partial_u\wedge\partial_v$ and $\Pi=h\pi$,
where $h$ has weight one.
The constant-bivector normalization gives
$\mathcal U_2(\pi,\pi)(u^2,v^2)=1$.
The order-$h$ term of $\mathcal F_1(t\pi)$ on
these two functions is $ht$.
Its mixed differential is $h\,dt$.
In $\mathcal U_2^K(\pi,dt\,\pi)$, the input-crossing
sign and the degree-minus-one map sign cancel,
giving precisely this value.

The target remains the completed finite-direction
polydifferential complex.
For example,
$\sum_i x_i\partial_{x_i}$ acts on every polynomial
and is bounded for the fine bornology.
It is not a finite-direction derivation.
It is not a Hochschild boundary either, since
the nullary Hochschild differential of a
commutative algebra is zero.
Thus a comparison with the full bounded
Hochschild complex requires a larger source.

\section{A variable Poisson coefficient}
\label{sec:vol2-variable-poisson-test}

A small calculation shows why both the target differential
and the derivative corrections matter.
On $\mathbb C[x,y]$, take
$\pi=x\partial_x\wedge\partial_y$ and use one
parameter $h$.
The first coefficient of the star product is
\[
 B_1(f,g)=\frac{x}{2}(f_xg_y-f_yg_x).
\]
The Schouten and Hochschild formulas give
\[
 [\pi,x]_{\mathrm{Sch}}=x\partial_y,\qquad
 [\pi,\partial_x]_{\mathrm{Sch}}
       =-\partial_x\wedge\partial_y,
\]
\[
 [m_\star,x]_G(y)=hx+O(h^2),\qquad
 [m_\star,\partial_x]_G(x,y)=-h/2+O(h^2).
\]
For the last expression, expand
$(\partial_xx)\star y+x\star(\partial_xy)
 -\partial_x(x\star y)$.
The first-order term is $-h/2$.

There cannot be a cochain map reducing to HKR from
$(T_{\mathrm{fin}}[[h]],h[\pi,-]_{\mathrm{Sch}})$
to the undeformed target
$(D_{\mathrm{fin}}[[h]],b_0)$.
Such a map sends $x$ to a nullary cochain, on which
$b_0$ is identically zero.
Its order-$h$ chain equation would be
$0=h\,\operatorname{HKR}([\pi,x])=hx\partial_y$.
Higher coefficients cannot change that equation.

For the second product term, choose the affine gauge
with
\begin{equation}\label{eq:vol2-variable-poisson-second-term}
 \begin{split}
 B_2(f,g)={}&\frac{x^2}{8}
       (f_{xx}g_{yy}-2f_{xy}g_{xy}+f_{yy}g_{xx})\\
 &+\frac{x}{12}
       (f_xg_{yy}+f_{yy}g_x-f_{xy}g_y-f_yg_{xy}).
 \end{split}
\end{equation}
This is the second-order affine formula
\cite{KontsevichDQ}*{Section~1.4.2}, normalized so
that the first commutator is $h\{-,-\}$.

\begin{proposition}\label{prop:vol2-variable-poisson-associator}
The product $f\star g=fg+hB_1(f,g)+h^2B_2(f,g)$
is associative modulo $h^3$.
In particular,
\[
 (x\star y)\star y=x\star(y\star y)
   =xy^2+hxy+\frac{h^2x}{6}\pmod{h^3}.
\]
Omitting the $x/12$ term in
\eqref{eq:vol2-variable-poisson-second-term}
leaves associator $h^2x/4$ on $(x,y,y)$.
\end{proposition}
\begin{proof}
Leibniz's rule cancels the first-order associator.
The second coefficient on $(f,g,k)$ is
\[
 \begin{split}
 &B_2(f,g)k+B_1(B_1(f,g),k)+B_2(fg,k)\\
 &\hspace{1cm}
   -fB_2(g,k)-B_1(f,B_1(g,k))-B_2(f,gk).
 \end{split}
\]
Using only the $x^2/8$ term and expanding each
derivative by Leibniz gives
$x(f_xk_y-f_yk_x)g_y/4$.
The four terms with coefficient $x/12$ give its
negative. This proves the assertion for arbitrary
polynomials.
For the displayed triple,
$B_2(xy,y)=-x/12$ and $B_2(x,y^2)=x/6$.
The iterated $B_1$ term adds $x/4$ on the left,
giving $x/6$ on each side.
Without the derivative correction, the left side
has $x/4$ and the right side has zero in order two.
\end{proof}

The parameter also affects Lie compatibility.
After inverting $h$, multiplication by $h^{p-1}$
on $p$-vectors identifies the differentials
$[\pi,-]$ and $h[\pi,-]$ and preserves the
Schouten bracket: on a bracket of a $p$-vector and
a $q$-vector both powers are $h^{p+q-2}$.
It sends the function $1$ to $h^{-1}$.
The unit-preserving scaling $h^p$ instead preserves
the bracket if the source bracket is multiplied
by $h$. The two rescalings answer different
comparison questions.

\section{Comparing quantizations and continuous operations}
\label{sec:vol2-continuous-wheel-comparison}

Two quantizations on the same coefficient module
must be compared in their allowed operator space.
Suppose associative products $\mu=m+O(h)$ and
$\mu'=m+O(h)$ agree modulo $h^N$, for $N\geq1$.
Assume their coefficient domains and differentials
have been identified, and put
$\Delta_N=h^{-N}(\mu'-\mu)\bmod h$.
Subtracting associativity in order $h^N$ gives
$b_0\Delta_N=0$.
A permitted degree-zero unary operator $D$ gives
\[
 \mu'=(1+h^ND)^{-1}
        \mu\bigl((1+h^ND)\otimes(1+h^ND)\bigr)
                      \pmod{h^{N+1}}
\]
if and only if
\begin{equation}\label{eq:vol2-unary-gauge-equation}
 \Delta_N=b_0D.
\end{equation}
Indeed expansion of the right side gives
\[
 \mu+h^N\bigl(m(D\otimes1)+m(1\otimes D)-Dm\bigr).
\]
If the mixed differential $d_K$ is fixed, the
same comparison also requires $[d_K,D]=0$.
A primitive in a larger space of operators does
not solve the equation in the prescribed domain.
For a commutative product $b_0D$ is symmetric in
its inputs, so a nonzero antisymmetric first-order
difference cannot be removed this way.
For example,
$b_0(\partial_x^2)(f,g)=-2f_xg_x$.

The Hamiltonian jet system has a different
completion from the finite-direction weighted
model. The two-index construction of
Chapter~\ref{ch:vol2-continuous-cochains}
gives $\mathcal P_{\mathrm{uc}}$ and its cochain
comparison with $C^\bullet_{\mathrm{uc}}(A/R,A)$.
A geometric polyvector realization must first be
compared with that complex in its stated topology.
The compatible-jet criterion of
Theorem~\ref{thm:vol2-compatible-jets} records the
additional transitions and limit identifications
needed for a proposed $\mathcal P_\infty$.

The further question is a continuous
homotopy-$E_2$ quasi-isomorphism into this
two-cutoff Hochschild recipient, with an invariant
Duflo map and a specified renormalization of the
wheels. A negative answer would require a brace
obstruction that persists under every allowed
renormalization. The fixed-normalization
counterexample of
Proposition~\ref{prop:vol2-mode-duflo-obstruction}
does not have that quantifier.

At the binary product level a proposed cochain
map $\Phi$ already needs a homotopy satisfying
\[
 dH(u,v)+H(du,v)+(-1)^{|u|}H(u,dv)
   =\Phi(uv)-\Phi(u)\smile\Phi(v).
\]
The brace substitutions require compatible higher
homotopies on the same completed coefficient
domain. A finite invariant algebra isomorphism
does not specify them.
The mixed-chain traces, cyclic orientation and
finite wheel operators constructed above provide
separate parts of such a comparison.
To identify their action with an independently
defined field theory, one must also preserve
insertion supports, singular coefficients,
collision substitutions and the chosen state.


\chapter{Infinite-mode states and local stabilization}
\label{ch:vol2-infinite-mode-states}

A current commutator does not determine a state. For one oscillator,
\([a,a^*]=1\), the vacuum has expectation zero on \(a^*a\).
A thermal state at inverse temperature \(T\) has expectation
\((e^T-1)^{-1}\) on the same operator. Both give expectation one
on the commutator. The singular current coefficient can therefore
survive a change of preparation while its regular correlations change.

Chapter~\ref{ch:vol2-gaussian-elimination} constructed finite Gaussian
states and their higher operations. Passing to infinitely many modes
requires a domain, an energy estimate and a normalization. We construct
these data on a Heisenberg Fock space. A local cubic interaction then
separates two possible conclusions: every perturbative coefficient
exists and sews, although the numerical cubic heat operator is unbounded.
A local quartic correction gives positive states. Its coefficient
determines whether those states approach the free state continuously.

\section{The common Hilbert space}

Fix a numerical level \(\lambda>0\). Positivity in this chapter
concerns complex Hilbert spaces with this fixed real parameter.
It is separate from an identity over a formal quantum coefficient ring.
Let
\[
 \mathcal H_1=\ell^2(\mathbb N_{>0}),\qquad
 \mathcal F=\Gamma_s(\mathcal H_1)
 =\widehat{\bigoplus}_{r\geq0}\operatorname{Sym}^r\mathcal H_1.
\]
The inner product is conjugate linear in its first argument.
The occupation vectors \(\lvert\mathbf m\rangle\), indexed by
finite-support sequences \(\mathbf m=(m_1,m_2,\ldots)\), form an
orthonormal basis. Define
\[
 a_n\lvert\mathbf m\rangle
 =\sqrt{m_n}\lvert\mathbf m-\mathbf e_n\rangle,\qquad
 a_n^*\lvert\mathbf m\rangle
 =\sqrt{m_n+1}\lvert\mathbf m+\mathbf e_n\rangle.
\]
Thus \([a_m,a_n^*]=\delta_{mn}\). Write
\[
 \alpha_n=\sqrt n\,a_n,\qquad
 \alpha_{-n}=\sqrt n\,a_n^*,\qquad
 J_n=\sqrt\lambda\,\alpha_n\quad(n\ne0),
 \qquad\alpha_0=J_0=0.
\]
The current relations are
\[
 [\alpha_m,\alpha_n]=m\delta_{m+n,0},\qquad
 J_n^*=J_{-n},\qquad
 [J_m,J_n]=\lambda m\delta_{m+n,0}.
\]

The vacuum \(\Omega=\lvert\mathbf0\rangle\) has energy zero.
The energy operator and particle-number operator are
\[
 H_0=\sum_{n\geq1}n a_n^*a_n
     =\sum_{n\geq1}\alpha_{-n}\alpha_n,
 \qquad N=\sum_{n\geq1}a_n^*a_n\leq H_0.
\]
These are diagonal self-adjoint operators on their maximal domains.
For example,
\[
 \operatorname{Dom}(H_0)
 =\left\{\sum_{\mathbf m}v_{\mathbf m}\lvert\mathbf m\rangle:
   \sum_{\mathbf m}\left(\sum_n n m_n\right)^2
              |v_{\mathbf m}|^2<\infty\right\}.
\]
Let \(\mathcal F_E\) be the energy-\(E\) eigenspace and \(P_E\)
its orthogonal projection. Its dimension is the partition number
\(p(E)\). The finite-energy space
\(\mathcal F_{\mathrm{fin}}=\bigoplus_{E\geq0}\mathcal F_E\)
is a common core for the finite-mode operators. Their closed extensions
are used whenever an infinite vector is an input.

For \(f\in\mathcal H_1\), the smeared oscillators are
\[
 a(f)=\sum_n\overline{f_n}a_n,\qquad
 a^*(f)=\sum_n f_n a_n^*.
\]
They act on the algebraic direct sum of the particle-number spaces,
which permits arbitrary square-summable one-particle vectors.
On that domain,
\[
 \|a(f)v\|\leq\|f\|_2\|N^{1/2}v\|,
 \qquad
 \|a^*(f)v\|\leq\|f\|_2\|(N+1)^{1/2}v\|.
\]
Successive application bounds a polynomial of degree \(d\) in finitely
many such oscillators by a constant times \((N+1)^{d/2}\).
This estimate specifies the unbounded observables whose thermal
expectations will be used below.

\section{A prepared trace and its current kernel}

Choose a cylinder length \(T>0\) and put \(q=e^{-T}\).
Its evolution operator is \(K_T=e^{-TH_0}\). This choice supplies
the preparation before any comparison with a bulk theory.

\begin{theorem}\label{thm:vol2-fock-thermal-state}
The operator \(K_T\) is positive and trace class. If
\(E_1 e_n=n e_n\) on \(\mathcal H_1\), then
\begin{equation}\label{eq:vol2-fock-partition}
 Z_0(T)=\operatorname{Tr}_{\mathcal F}K_T
 =\prod_{n\geq1}(1-e^{-Tn})^{-1}
 =\det_{\mathcal H_1}(1-e^{-TE_1})^{-1}.
\end{equation}
The density \(\rho_T=K_T/Z_0(T)\) defines a positive normalized
state on the bounded operators on \(\mathcal F\). Every polynomial
in finitely many square-summably smeared oscillators also has a
thermal expectation
\[
 \varphi_T(A)=Z_0(T)^{-1}\operatorname{Tr}(A K_T).
\]
For the radially ordered current \(J(z)=\sum_nJ_nz^{-n-1}\),
\begin{equation}\label{eq:vol2-thermal-current}
 \varphi_T(J(z)J(w))
 =\frac{\lambda}{zw}\sum_{n\geq1}
 \frac{n\bigl((w/z)^n+q^n(z/w)^n\bigr)}{1-q^n},
 \qquad q<|w/z|<1.
\end{equation}
The series and every fixed holomorphic derivative converge normally
on compact subannuli. Near \(z=w\ne0\), its singular part is
\(\lambda/(z-w)^2\).
\end{theorem}

\begin{proof}
At a mode cutoff \(M\), independent occupation sums give
\(Z_{0,M}(T)=\prod_{n=1}^M(1-q^n)^{-1}\). Moreover,
\begin{equation}\label{eq:vol2-partition-tail}
 0\leq\log\frac{Z_0(T)}{Z_{0,M}(T)}
 \leq\frac{q^{M+1}}{(1-q)(1-q^{M+1})}.
\end{equation}
Indeed \(-\log(1-x)\leq x/(1-x)\), and summing the resulting
geometric bound proves convergence of the logarithmic product.
Monotone summation over all occupation vectors gives its equality
with the trace. The eigenvalues \(q^n\) of \(e^{-TE_1}\) are
summable, so its Fredholm determinant is this same convergent product.

For every \(r\geq0\), the operator \((H_0+1)^r e^{-TH_0}\)
is trace class: the polynomial factor is bounded by a constant
times \(e^{TE/2}\) on energy \(E\), and the remaining trace is
\(Z_0(T/2)\). The oscillator bounds therefore make \(A K_T\)
trace class. They also make \(A K_{T/2}\) Hilbert--Schmidt.
Consequently the expectation of \(A^*A\) is its squared
Hilbert--Schmidt norm divided by \(Z_0(T)\), and is nonnegative.
This extends the positive state to the stated polynomial domain.

The one-mode geometric sums give
\[
 \varphi_T(J_nJ_{-n})=\frac{\lambda n}{1-q^n},\qquad
 \varphi_T(J_{-n}J_n)=\frac{\lambda n q^n}{1-q^n}
 \quad(n>0).
\]
They imply \eqref{eq:vol2-thermal-current}. On a compact subannulus,
both \(|w/z|\) and \(q|z/w|\) are bounded strictly below one.
The denominator is bounded below by \(1-q\), and derivatives add
only polynomial factors in \(n\). Split the first numerator into
its vacuum contribution and its \(q^n\)-weighted remainder.
The vacuum sum \(\sum_{n\geq1}n(w/z)^n\) gives
\(\lambda/(z-w)^2\); the two remaining sums extend across \(z=w\).
\end{proof}

At a cutoff, the normalized-oscillator contractions are bounded by
\((1-q)^{-1}\|f\|_2\|h\|_2\). Wick expansion therefore gives
the same polynomial expectations by a mode limit. The estimates
prove convergence of an actual density and determinant, as well as
convergence of its correlations.

A fixed charge \(b\in\mathbb R\) can be retained by setting
\(J_0=b\) and adding \(b^2/(2\lambda)\) to the energy.
The partition function becomes \(e^{-Tb^2/(2\lambda)}Z_0(T)\),
the mean current is \(b/z\), and the connected two-point function
is unchanged. A sum over charges additionally requires a measure
or lattice and a convergence domain.

\section{Sewing before normalization}

For \(z\in\mathcal H_1\), define the exponential vector
\[
 \mathrm e(z)=\bigoplus_{r\geq0}\frac{z^{\otimes r}}{\sqrt{r!}}.
\]
Its norm squared is \(e^{\|z\|_2^2}\). The cylinder matrix element is
\begin{equation}\label{eq:vol2-coherent-cylinder}
 K_T(z,\overline w)
 =\langle\mathrm e(w),K_T\mathrm e(z)\rangle
 =\exp\!\left(\sum_{n\geq1}e^{-Tn}z_n\overline{w_n}\right).
\end{equation}
Cauchy--Schwarz makes the exponent convergent.

\begin{proposition}\label{prop:vol2-fock-sewing}
For \(T_1,T_2>0\), coherent cylinder kernels sew by
\(K_{T_1}\circ K_{T_2}=K_{T_1+T_2}\). Closing the composite by
the trace gives \(Z_0(T_1+T_2)\).
\end{proposition}

\begin{proof}
For a bounded one-particle contraction \(B\), second quantization is
\(\Gamma(B)=\bigoplus_{r\geq0}\operatorname{Sym}^rB\). It gives
\[
 \Gamma(e^{-T_1E_1})\Gamma(e^{-T_2E_1})
 =\Gamma(e^{-(T_1+T_2)E_1}).
\]
The finite-dimensional coherent Gaussian integral computes this
identity at each mode cutoff. Its limit is Fock-space composition.
The high-mode vacuum projections commute with the heat operators,
and the omitted positive trace tends to zero by
\eqref{eq:vol2-partition-tail}. Thus taking the trace of the
composite gives the stated determinant. Associativity is already
the associativity of operator composition.
\end{proof}

The infinite-mode pairing here is a limit of cylindrical Wick
pairings. No translation-invariant measure on \(\ell^2\) is used.
Normalization is performed after the required trace closure:
\[
 \rho_{T_1}\rho_{T_2}
 =\frac{Z_0(T_1+T_2)}{Z_0(T_1)Z_0(T_2)}\rho_{T_1+T_2}.
\]
The scalar factor is part of sewing data.

The thermal state also satisfies, for finite-mode oscillator
polynomials \(A,B\),
\begin{equation}\label{eq:vol2-fock-kms}
 \varphi_T(AB)
 =\varphi_T\bigl(B e^{-TH_0}A e^{TH_0}\bigr).
\end{equation}
This follows by cyclicity of the trace on the common domain, or
directly from the occupation sums. The conjugation multiplies
\(\alpha_n\) by \(e^{Tn}\). For general infinite smears it
therefore requires a further domain condition.
Equation~\eqref{eq:vol2-fock-kms} is the thermal, or KMS, identity.
Ordinary traciality fails already on \(a_na_n^*\) and \(a_n^*a_n\).
A cyclic chain class from Chapter~\ref{chap:vol2-cyclic-traces-wheels}
and this scalar state are consequently different constructions.

\section{The stress tensor and the cylinder constant}

The vacuum has contraction
\(\langle\Omega,J(z)J(w)\Omega\rangle=\lambda/(z-w)^2\)
in \(|w|<|z|\). Subtract this contraction to define
\[
 T(z)=\frac1{2\lambda}{:}J(z)^2{:}
     =\frac12{:}\alpha(z)^2{:}
     =\sum_{m\in\mathbb Z}L_mz^{-m-2},
 \qquad
 L_m=\frac12\sum_{j\in\mathbb Z}{:}\alpha_j\alpha_{m-j}{:}.
\]
Normal ordering puts negative-index creators to the left.
Terms of index zero vanish. Each sum is finite on a finite-energy
vector, and \(L_0=H_0\).

\begin{proposition}\label{prop:vol2-current-virasoro}
On \(\mathcal F_{\mathrm{fin}}\),
\[
 [L_m,\alpha_n]=-n\alpha_{m+n},\qquad
 [L_m,L_n]=(m-n)L_{m+n}
       +\frac{m^3-m}{12}\delta_{m+n,0}.
\]
If \(z=f(u)\) is a local holomorphic coordinate change with
\(f'\ne0\), the stress tensor defined by subtraction in the
\(u\)-coordinate is
\begin{equation}\label{eq:vol2-stress-schwarzian}
 T_u(u)=f'(u)^2T_z(f(u))+\frac1{12}\{f,u\},
 \qquad
 \{f,u\}=\frac{f'''}{f'}-\frac32\left(\frac{f''}{f'}\right)^2.
\end{equation}
In particular \(f(u)=e^u\) gives cylinder Hamiltonian
\(H_0-1/24\).
\end{proposition}

\begin{proof}
Commuting a current through the two factors of \(L_m\) gives
the first identity. The Jacobi identity then shows that
\([L_m,L_n]-(m-n)L_{m+n}\) commutes with every current.
An operator with this property is scalar: it preserves the common
kernel of the positive modes, which is \(\mathbb C\Omega\), and
then has the same scalar action on all created vectors.
Its energy degree forces that scalar to vanish unless \(m+n=0\).
For \(m>0\), its vacuum expectation is
\[
 \|L_{-m}\Omega\|^2
 =\frac12\sum_{j=1}^{m-1}j(m-j)
 =\frac{m^3-m}{12}.
\]
The two contractions retain the factor \(1/2\), including when
\(j=m-j\). Antisymmetry gives the other signs of \(m\).

The current transforms as a one-form, so its pulled-back vacuum
contraction is \(\lambda f'(u)f'(v)/(f(u)-f(v))^2\).
Taylor expansion at \(v=u\) gives
\[
 \frac{f'(u)f'(v)}{(f(u)-f(v))^2}
 =\frac1{(u-v)^2}+\frac16\{f,u\}+O(v-u).
\]
For completeness, writing \(a=f''/f'\), \(b=f'''/f'\), the
constant term is \(b/2-a^2+3a^2/4-b/3=b/6-a^2/4\).
Subtracting the \(u\)-coordinate pole and dividing by \(2\lambda\)
proves \eqref{eq:vol2-stress-schwarzian}.
Since \(\{e^u,u\}=-1/2\), its constant is \(-1/24\).
\end{proof}

The conformal cylinder convention consequently has partition function
\[
 \operatorname{Tr}e^{-T(H_0-1/24)}=e^{T/24}Z_0(T).
\]
The scalar factor respects sewing and cancels from normalized
expectations. It arises from the specified stress subtraction.
Identifying it with a determinant in another field theory requires
that theory's regulator, finite term, boundary condition and sewing
normalization to be compared.

\section{Formal interactions on two coefficient domains}

Let \(V\) be a finite linear combination of normally ordered
polynomials in \(a(f_j),a^*(f_j)\), for finitely many
\(f_j\in\ell^2\). Define a formal coupling series by
\begin{equation}\label{eq:vol2-state-dyson}
 \begin{split}
 U_T(g)&=\sum_{k\geq0}g^k U_T^{(k)},\qquad U_T^{(0)}=K_T,\\
 U_T^{(k)}&=(-1)^k\int_{0\leq t_1\leq\cdots\leq t_k\leq T}
 e^{-(T-t_k)H_0}V e^{-(t_k-t_{k-1})H_0}\cdots\\
 &\hspace{24mm}\cdots V e^{-t_1H_0}\,dt_1\cdots dt_k.
 \end{split}
\end{equation}
For this coefficient domain, the definition first uses Wick matrix
elements and traces at a mode cutoff, followed by their limits.

\begin{proposition}\label{prop:vol2-smeared-dyson}
Every thermal trace and coherent-vector matrix element of each
coefficient in \eqref{eq:vol2-state-dyson} has a mode-cutoff limit.
The same holds with a fixed polynomial smeared insertion.
In these coefficientwise sewn matrix elements,
\begin{equation}\label{eq:vol2-formal-state-sewing}
 U_{T_1}(g)U_{T_2}(g)=U_{T_1+T_2}(g).
\end{equation}
Dividing observable traces by \(\operatorname{Tr}U_T(g)\) defines
a normalized formal functional.
\end{proposition}

\begin{proof}
A time-ordered thermal contraction of normalized modes contains
one of
\[
 \frac{e^{-n\tau}}{1-e^{-nT}},\qquad
 \frac{e^{-n(T-\tau)}}{1-e^{-nT}},\qquad0\leq\tau\leq T.
\]
Both are bounded by \(C_T=(1-e^{-T})^{-1}\). A smeared contraction
is therefore at most \(C_T\|f\|_2\|h\|_2\), uniformly in the
mode cutoff and time. At fixed order there are finitely many monomial
choices and Wick pairings. Their bounds, multiplied by the simplex
volume \(T^k/k!\), are integrable. The contraction tails converge
by Cauchy--Schwarz, so dominated convergence gives the mode limits.

For coherent matrix elements, each external contraction contains
a one-particle heat contraction \(e^{-sE_1}\), with \(s\geq0\).
The free coherent factor is bounded by \(e^{\|z\|_2\|w\|_2}\).
These bounds prove the same assertion with external vectors.
A fixed insertion only adds finitely many contractions.

At a cutoff, split the time simplex at \(T_2\). Its pieces are
indexed by the number of insertions before the cut and give the
coefficient of \(g^k\) in \eqref{eq:vol2-formal-state-sewing}.
The bounds pass this identity to the limit. The denominator has
constant term \(Z_0(T)>0\), and hence a formal inverse.
\end{proof}

This domain permits infinitely many oscillator modes in a smear,
but only finitely many input directions. It does not permit a
point-current smear: its regularized one-particle norm contains
\[
 \sum_{n\geq1}n e^{-2\varepsilon n}
 =\frac{e^{-2\varepsilon}}{(1-e^{-2\varepsilon})^2}
 \sim\frac1{4\varepsilon^2}.
\]
Local spatial Wick polynomials require a different estimate.

For fixed \(n\in\mathbb Z\), \(d\geq1\) and polynomial \(P\),
define the local mode
\begin{equation}\label{eq:vol2-local-wick-mode}
 W_n=\sum_{n_1+\cdots+n_d=n}
 P(n_1,\ldots,n_d){:}\alpha_{n_1}\cdots\alpha_{n_d}{:}.
\end{equation}
Zero indices give zero terms. The polynomial \(P\) includes factors
from angular derivatives. Finite Fourier combinations of these
operators, together with the identity, form the second domain.

\begin{lemma}\label{lem:vol2-local-mode-bound}
The sum \eqref{eq:vol2-local-wick-mode} is finite on \(\mathcal F_E\),
maps it to \(\mathcal F_{E-n}\), and satisfies
\[
 \|W_nP_E\|\leq C_{n,d,P}(1+E+|n|)^{3d/2-1+\deg P}.
\]
The same estimate holds after discarding any set of its summands.
\end{lemma}

\begin{proof}
The annihilators act first. Their total index is at most \(E\)
whenever the monomial is nonzero. The creators' total absolute index
is that sum minus \(n\), hence at most \(E+|n|\).
Each contributing index is bounded by \(E+|n|\), and the sum
constraint determines the last from the first \(d-1\).
There are at most \(C_d(1+E+|n|)^{d-1}\) tuples.

Intermediate energies first decrease and then increase to \(E-n\).
The norm of each oscillator factor is at most
\((1+E+|n|)^{1/2}\): its square is bounded by the energy before
annihilation or after creation. Multiply the \(d\) bounds and
the polynomial bound for \(P\), then sum over tuples.
The exponent is \((d-1)+d/2+\deg P\). Since this argument bounds
each summand separately, it also bounds every partial sum.
\end{proof}

For example \(W_0=2H_0\) when \(d=2\) and \(P=1\).
Its one-particle compression has eigenvalues \(2,4,6,\ldots\).
A polynomial in finitely many square-summable smears instead has
one-particle compression equal to a scalar plus a bounded finite-rank
operator. The two domains therefore do not coincide.

\begin{theorem}\label{thm:vol2-local-mode-dyson}
Let \(V\) be a finite Fourier combination of local Wick modes.
For every \(T>0\) and \(k\geq0\), the coefficient
\(U_T^{(k)}\) has a unique trace-class extension from finite-energy
vectors. Mode-cutoff coefficients, embedded by the high-mode vacuum,
converge in trace norm. A fixed finite Fourier Wick-polynomial
insertion \(A\) also makes \(A U_T^{(k)}\) trace class.
The formal operators obey \eqref{eq:vol2-formal-state-sewing}, and
their normalized observable traces are well defined.
\end{theorem}

\begin{proof}
There are \(C,M,R\) such that
\[
 V\mathcal F_E\subset\bigoplus_{|r|\leq R}\mathcal F_{E+r},
 \qquad\|VP_E\|\leq C(1+E)^M.
\]
At order \(k\), every intermediate energy lies between
\(E-kR\) and \(E+kR\). For a fixed sequence of energy projections,
the heat factors are bounded by \(e^{-TE+TkR}\), uniformly over
the closed time simplex. The operator factors contribute at most
\(C^k(1+E+kR)^{Mk}\). Absorb the finite number of energy-shift
choices into \(C^k\). A block with source \(\mathcal F_E\)
has rank at most \(p(E)\), so
\begin{equation}\label{eq:vol2-local-dyson-tail}
 \|U_T^{(k)}\|_1
 \leq\frac{T^k}{k!}C^k e^{TkR}
        \sum_{E\geq0}p(E)e^{-TE}(1+E+kR)^{Mk}<\infty.
\end{equation}
The sum is finite because a polynomial energy weight is dominated
by \(e^{TE/2}\). Finite block sums are trace-norm continuous in
the simplex variables; the uniform summable tail gives a trace-norm
Bochner integral, including coincident-time faces.

For fixed \(E,k\), all contributing oscillator indices occur
below \(E+kR\), after enlarging this bound by the fixed finite
shift if necessary. Those blocks stabilize at a mode cutoff.
Lemma~\ref{lem:vol2-local-mode-bound} gives the same tail bound at
every cutoff, proving trace-norm convergence. An insertion only adds
fixed shifts and a polynomial factor to this estimate.

Subdivision of the simplex proves sewing on finite blocks.
The trace-norm convergence passes the identity to the full space,
using \(\|BC\|_1\leq\|B\|_1\|C\|\) and boundedness of each
trace-class coefficient. The free partition function is the
nonzero constant denominator. Uniqueness follows from the density
of finite-energy vectors.
\end{proof}

Both constructions are coefficientwise in \(g\). Neither estimate
asserts convergence as the coupling order tends to infinity.
The local-mode theorem concerns spatially local boundary operators
on the stated Fock carrier. A local interacting half-space theory
also requires its supported counterterms and quantum master equation.

\section{The cubic mode and its spectral obstruction}

Consider the self-adjoint finite-energy matrices of the cubic mode
\begin{equation}\label{eq:vol2-cubic-mode}
 \begin{split}
 V&=\frac16\operatorname{Res}_{z=0}z^2{:}\alpha(z)^3{:}\,dz\\
  &=\frac12\sum_{m,n\geq1}
   \bigl(\alpha_{-m}\alpha_{-n}\alpha_{m+n}
          +\alpha_{-(m+n)}\alpha_m\alpha_n\bigr).
 \end{split}
\end{equation}
It commutes with \(H_0\). The factor \(1/2\) comes from the
three placements of the uniquely signed mode in the residue.
In the polynomial realization \(\alpha_{-n}=p_n\),
\(\alpha_n=n\partial_{p_n}\), it becomes
\begin{equation}\label{eq:vol2-cut-and-join}
 V=\frac12\sum_{m,n\geq1}
 \left((m+n)p_mp_n\partial_{p_{m+n}}
       +mn p_{m+n}\partial_{p_m}\partial_{p_n}\right).
\end{equation}
This is the cut-and-join operator, with the normalization of
Mironov--Morozov--Natanzon \cite{MironovMorozovNatanzon2009}*{equation (51)}.

\begin{lemma}\label{lem:vol2-cubic-norm}
On energy \(E\), the restriction \(V_E\) has norm
\(E(E-1)/2\), and both extreme signs occur.
\end{lemma}

\begin{proof}
For a partition \(\mu=(1^{m_1}2^{m_2}\cdots)\) of \(E\), put
\(p_\mu=\prod_n p_n^{m_n}\) and
\(z_\mu=\prod_n n^{m_n}m_n!\). The Fock pairing is
\(\langle p_\mu,p_\nu\rangle=\delta_{\mu\nu}z_\mu\).
The characteristic map on class functions of \(S_E\),
\[
 \operatorname{ch}(f)=\sum_{\mu\vdash E}\frac{f(\mu)}{z_\mu}p_\mu,
\]
is an isometry for normalized counting measure. Convolution by the
sum of all transpositions preserves class functions. A transposition
joins two cycles or splits one. Counting its marked entries gives
the coefficients \(mn\) and \(m+n\) in
\eqref{eq:vol2-cut-and-join}, with the factor \(1/2\) for the two
orders. Thus this sum corresponds to \(V_E\).

Every transposition acts unitarily and self-adjointly by translation
on functions. The sum has norm at most \(\binom E2\).
The trivial and sign characters attain the positive and negative
values. Their characteristic images are the complete and elementary
symmetric polynomials \(h_E,e_E\), defined by
\[
 \sum_{E\geq0}h_Et^E
 =\exp\!\left(\sum_{n\geq1}\frac{p_nt^n}{n}\right),\qquad
 \sum_{E\geq0}e_Et^E
 =\exp\!\left(\sum_{n\geq1}\frac{(-1)^{n-1}p_nt^n}{n}\right).
\]
They give \(Vh_E=\binom E2h_E\) and
\(Ve_E=-\binom E2e_E\), proving equality in the norm bound.
\end{proof}

For nonzero real \(g\), one of these families has energy
\[
 (H_0+gV)\big|_{\mathrm{extreme},E}
 =E-\frac{|g|}{2}E(E-1)\longrightarrow-\infty.
\]
The self-adjoint direct sum \(H_0+gV\) is therefore unbounded
below, and \(e^{-T(H_0+gV)}\) is unbounded for every \(T>0\).
It cannot prepare a trace-class Gibbs state. The formal coefficients
from Theorem~\ref{thm:vol2-local-mode-dyson} nevertheless exist.
Since \([H_0,V]=0\), they are
\(U_T^{(k)}=(-T)^k e^{-TH_0}V^k/k!\).

The first response already sees the interaction. With
\(\langle A\rangle_q=\varphi_T(A)\), one has
\begin{equation}\label{eq:vol2-cubic-thermal-moment}
 \langle V^2\rangle_q
 =\sum_{m,n\geq1}
 \frac{mn(m+n)q^{m+n}}
      {(1-q^m)(1-q^n)(1-q^{m+n})}.
\end{equation}
To prove the formula, split \(V\) into its joining and splitting
terms. Only a join followed by a split, or the reverse order,
has complete thermal contractions. The two orders and the two
pairings of the \(m,n\) legs cancel the factor \(1/4\).
The same multiplicity remains when \(m=n\). Substitution of the
one-mode expectations proves the formula. Its convergence follows
from the exponential numerator and denominators bounded below by
\(1-q\). Its first coefficients are
\[
 2q^2+16q^3+60q^4+160q^5+360q^6+O(q^7).
\]
Oscillator parity sends \(V\) to \(-V\) and preserves the free
density, so all odd moments vanish. Hence the formal partition
function obeys
\[
 \log\frac{\operatorname{Tr}U_T(g)}{Z_0(T)}
 =\frac{g^2T^2}{2}\langle V^2\rangle_q+O(g^4).
\]

\section{Stabilization by the total energy}

A square of the total energy provides a first positive preparation.
For real \(\varepsilon,g\), define the self-adjoint block operator
\[
 H_{\varepsilon,g}=H_0+\varepsilon H_0^2+gV.
\]
The term \(H_0^2\) is a prescribed nonlocal interaction on the
circle. It changes the Hamiltonian whose state is being constructed.

\begin{theorem}\label{thm:vol2-nonlocal-stability}
For \(T>0\), the operator \(e^{-TH_{\varepsilon,g}}\) is positive
and trace class exactly when
\begin{equation}\label{eq:vol2-nonlocal-threshold}
 \varepsilon\geq |g|/2.
\end{equation}
On this region it gives a positive normalized state on bounded
operators, and its unnormalized heat kernels sew by addition of
cylinder lengths. For fixed \(T>0\), the partition function is
holomorphic in the complex tube
\(\operatorname{Re}\varepsilon>|\operatorname{Re}g|/2\).
\end{theorem}

\begin{proof}
Lemma~\ref{lem:vol2-cubic-norm} gives, on energy \(E\),
\begin{equation}\label{eq:vol2-nonlocal-energy}
 H_{\varepsilon,g}\geq
 (\varepsilon-|g|/2)E^2+(1+|g|/2)E.
\end{equation}
At equality in \eqref{eq:vol2-nonlocal-threshold}, this is still
a positive linear lower bound. Summation against \(p(E)\) proves
trace class by \eqref{eq:vol2-fock-partition}. Below the threshold,
choose the extreme eigenvector with the sign opposite to \(g\).
If \(g=0\), either extreme vector suffices. Its energy equals the
right side of \eqref{eq:vol2-nonlocal-energy} and tends to
\(-\infty\), so the heat operator is unbounded.

Spectral calculus gives positivity and sewing for the allowed
parameters. In a common eigenbasis of \(H_0,V\), each exponential
eigenvalue is entire in \((\varepsilon,g)\). On a compact subset
of the tube, its absolute value is at most
\(e^{-cE^2-TE}\) for some \(c>0\). Every parameter derivative
adds only a polynomial energy factor. Uniform summability proves
holomorphy of the trace.
\end{proof}

This convergence has an explicit tail. Put
\(c=\varepsilon-|g|/2\geq0\). For an integer \(M\geq0\),
\begin{equation}\label{eq:vol2-stabilized-tail}
 \sum_{E>M}\operatorname{Tr}_{\mathcal F_E}e^{-TH_{\varepsilon,g}}
 \leq Z_0(T/2)
 \frac{e^{-T(M+1)/2-Tc(M+1)^2}}
      {1-e^{-T/2-Tc(2M+3)}}.
\end{equation}
Indeed \(p(E)\leq Z_0(T/2)e^{TE/2}\). The remaining dominating
series has first term as displayed and a decreasing ratio, bounded
by its first ratio. The estimate includes the equality case \(c=0\).
At fixed nonzero real \(g\), the stabilizer cannot be removed
within this family while retaining trace class.

\section{The unchanged cubic series and its Borel boundary}

The original cubic formal trace has coefficients
\begin{equation}\label{eq:vol2-cubic-formal-trace}
 Z(T,g)=\sum_{n\geq0}a_n(T)g^n,
 \qquad
 a_n(T)=\frac{(-T)^n}{n!}\operatorname{Tr}(e^{-TH_0}V^n).
\end{equation}
Its odd coefficients vanish and its even coefficients are positive.
Define its ordinary Borel transform by
\[
 \widehat Z_T(\zeta)=\sum_{n\geq0}\frac{a_n(T)}{n!}\zeta^n.
\]
The Borel variable \(\zeta\) is a complex scalar independent of
the current coordinate and of the coupling.

\begin{theorem}\label{thm:vol2-cubic-borel-radius}
For every \(T>0\), \eqref{eq:vol2-cubic-formal-trace} has radius
zero in \(g\). Its Borel transform has exact radius \(T/2\)
and diverges at both real endpoints \(\zeta=\pm T/2\).
\end{theorem}

\begin{proof}
Fix \(0<\nu<T\) and write \(t=T-\nu\). Since
\(p(E)\leq Z_0(\nu)e^{\nu E}\), the norm bound gives
\[
 \frac{|a_n(T)|}{n!}
 \leq Z_0(\nu)\frac{T^n}{2^n(n!)^2}
          \sum_{E\geq0}E^{2n}e^{-tE}
 \leq Z_0(\nu)e^t
       \frac{T^n(2n)!}{2^n(n!)^2t^{2n+1}}.
\]
For the last step, integration over \([E,E+1]\) bounds the
summand by \(e^t\int_E^{E+1}x^{2n}e^{-tx}\,dx\).
The central binomial coefficient has \(n\)th root tending to four.
Consequently
\[
 \limsup_n\left(\frac{|a_n(T)|}{n!}\right)^{1/n}
 \leq\frac{2T}{(T-\nu)^2}.
\]
Letting \(\nu\) decrease to zero gives the upper bound \(2/T\).

For even \(n\), retain one extreme eigenvector at energy
\(E_n=\lfloor2n/T\rfloor\), with \(n\) sufficiently large.
It contributes
\[
 \frac{a_n(T)}{n!}
 \geq\frac{T^n e^{-TE_n}}{(n!)^2}
          \left(\frac{E_n(E_n-1)}2\right)^n.
\]
Stirling's formula makes the \(n\)th root of this expression
tend to \(2/T\). This proves the exact Borel radius.
Multiplication by \(n!\) makes the corresponding coefficient
roots unbounded, proving radius zero for the original series.

For endpoint divergence, retain both extreme vectors at every
energy \(E\geq2\). For \(\zeta>0\), their combined contribution
to the Borel sum is
\[
 e^{-TE}\bigl(I_0(X_E)+J_0(X_E)\bigr),\qquad
 X_E=2\sqrt{T\zeta E(E-1)/2},
\]
where
\[
 I_0(X)=\sum_{n\geq0}\frac{(X^2/4)^n}{(n!)^2},\qquad
 J_0(X)=\sum_{n\geq0}\frac{(-1)^n(X^2/4)^n}{(n!)^2}.
\]
Their sum retains precisely the even powers. At \(\zeta=T/2\),
\(X_E=TE-T/2+O(E^{-1})\). The elementary integral representation
\[
 I_0(X)=\frac1\pi\int_0^\pi e^{X\cos\theta}\,d\theta
\]
gives \(I_0(X)\geq c e^X/\sqrt X\) for large positive \(X\),
by restricting to \(0\leq\theta\leq X^{-1/2}\).
The oscillatory integral for \(J_0\) gives \(|J_0(X)|\leq1\).
Thus the retained contribution is bounded below by a positive
multiple of \(E^{-1/2}\) for large \(E\). Its sum diverges.
All coefficients of the even Borel series are nonnegative, so
this subseries comparison is valid. Evenness gives the other endpoint.
\end{proof}

A fixed vacuum-energy shift multiplies the entire scalar series by
a nonzero constant and changes neither radius. Division by the
free partition function has the same property. The theorem does not
construct continuation outside the disk or choose a Laplace contour.
A positive resummation requires a state functional on an observable
star-algebra, with positivity and sewing for every admitted observable.
Those requirements are stronger than a continuation of the partition
function alone.

\section{Fermionic charges and a local quartic density}

The nonlocal square \(H_0^2\) stabilizes the spectrum, but it does
not express the correction as the integral of a local circle density.
For that purpose, a second realization of the same Fock space exposes
a useful positive charge.

Let \(r\) run over half-integers. A fermionic occupation set differs
from the filled negative sea
\(S_0=\{-1/2,-3/2,\ldots\}\) in finitely many places.
The corresponding wedges form an orthonormal basis. Creation
\(c_r^*\) inserts the mode \(r\), and annihilation \(c_r\)
removes it, with the exterior signs. They satisfy
\(\{c_r,c_s^*\}=\delta_{rs}\).
Normal ordering subtracts the vacuum occupation:
\[
 {:}c_r^*c_s{:}=c_r^*c_s-\delta_{rs}\mathbf1_{r<0}.
\]
The charge is the number of added positive particles minus the
number of holes in the negative sea. Charge-zero occupations are
indexed by partitions \(\mu\), through
\[
 S_\mu=\{\mu_i-i+1/2:i\geq1\}.
\]
Their energy is \(|\mu|\). Write
\(A(\theta)=\sum_{n\ne0}\alpha_n e^{-in\theta}\) for the angular
current, with each Fourier coefficient acting on finite-energy vectors.

\begin{proposition}\label{prop:vol2-boson-fermion-charge}
The charge-zero wedge space is unitarily isomorphic to \(\mathcal F\),
with vacuum sent to the filled sea and
\begin{equation}\label{eq:vol2-fermion-current}
 \alpha_n=\sum_{r\in\mathbb Z+1/2}{:}c_{r-n}^*c_r{:}
 \qquad(n\ne0).
\end{equation}
Define diagonal operators on the partition basis by
\begin{equation}\label{eq:vol2-partition-charges}
 Q_j(\mu)=\sum_{i\geq1}
 \left((\mu_i-i+1/2)^j-(-i+1/2)^j\right).
\end{equation}
Then \(Q_1=H_0\), \(Q_2=2V\), and
\begin{equation}\label{eq:vol2-local-quartic-charge}
 Q_3=\frac14\int_0^{2\pi}\frac{d\theta}{2\pi}
 \left({:}A(\theta)^4{:}
       +{:}(\partial_\theta A(\theta))^2{:}
       -\frac12{:}A(\theta)^2{:}\right).
\end{equation}
All these identities hold on finite-energy vectors.
\end{proposition}

\begin{proof}
For a fixed occupation set, only finitely many summands of
\eqref{eq:vol2-fermion-current} act. The anticommutation relations
give \([\alpha_m,\alpha_n]=m\delta_{m+n,0}\): off-diagonal
matrix units cancel, and for \(m=-n>0\) the vacuum boundary
contributes the \(m\) modes crossing zero. The positive modes
kill the filled sea. The vectors obtained by applying products of
negative modes therefore have the same orthogonal norms
\(z_\mu\) as the bosonic occupation monomials. They are linearly
independent. In energy \(E\), both spaces have dimension \(p(E)\),
so they span. This proves the isometry and its extension to the
Hilbert completions.

The diagonal charges and their vacuum scalar can be calculated
through a generating field. Put \(\varsigma(s)=2\sinh(s/2)\) and
\[
 \mathcal E_n(s)=\sum_r e^{s(r-n/2)}{:}c_{r-n}^*c_r{:}
                   +\frac{\delta_{n0}}{\varsigma(s)}.
\]
Direct multiplication of the fermionic matrix units gives
\begin{equation}\label{eq:vol2-bilocal-commutator}
 [\alpha_m,\mathcal E_n(s)]
       =\varsigma(ms)\mathcal E_{m+n}(s).
\end{equation}
When \(m+n=0\), the boundary sum is
\(\varsigma(ms)/\varsigma(s)\), which accounts for the scalar
in this equation. These are the charge and operator conventions
of Okounkov--Pandharipande \cite{OkounkovPandharipande2002}*{Sections 2.1--2.2}.
Here the vacuum scalar is subtracted from every \(Q_j\).

The normally ordered exponential
\[
 \mathcal B(z,s)=\frac1{\varsigma(s)}
 {:}\exp\!\left(\sum_{n\ne0}
       \frac{\varsigma(ns)}n\alpha_nz^{-n}\right){:}
\]
has commutator
\([\alpha_m,\mathcal B]=\varsigma(ms)z^m\mathcal B\).
The field \(\sum_n\mathcal E_n(s)z^{-n}\) has the same
commutators by \eqref{eq:vol2-bilocal-commutator}.
Both have vacuum matrix element \(1/\varsigma(s)\).
Commuting positive modes from the left and negative modes from the
right reduces any matrix element between created vacuum vectors
to that scalar matrix element. Induction on the total number of
external modes therefore proves equality of the two fields.
This is an identity of formal Laurent series in \(s\), evaluated
between finite-energy vectors.

Let \(\operatorname{CT}_z\) denote the angular constant term.
Taking it in the field identity and subtracting the vacuum gives
\begin{equation}\label{eq:vol2-charge-bilocal}
 \sum_{j\geq0}\frac{s^j}{j!}Q_j
 =\frac1{\varsigma(s)}
 \left[\operatorname{CT}_z
 {:}\exp\!\left(\sum_{n\ne0}
       \frac{\varsigma(ns)}n\alpha_nz^{-n}\right){:}-1\right].
\end{equation}
The left side follows directly from the finite occupation difference
in \eqref{eq:vol2-partition-charges}. In charge zero, \(Q_0=0\).

Write \(A(z)=\sum_{n\ne0}\alpha_nz^{-n}\) and
\(B(z)=\sum_{n\ne0}n^2\alpha_nz^{-n}\). The exponent is
\(sA+s^3B/24+O(s^5)\), while
\(\varsigma(s)^{-1}=s^{-1}-s/24+O(s^3)\).
Since \(\operatorname{CT}A=0\), extraction gives
\begin{align*}
 Q_1&=\tfrac12\operatorname{CT}{:}A^2{:}=H_0,\\
 Q_2&=\tfrac13\operatorname{CT}{:}A^3{:}=2V,\\
 Q_3&=\tfrac14\operatorname{CT}{:}A^4{:}
      +\tfrac14\operatorname{CT}{:}AB{:}
      -\tfrac18\operatorname{CT}{:}A^2{:}.
\end{align*}
Finally,
\[
 \operatorname{CT}{:}AB{:}
 =2\sum_{n>0}n^2\alpha_{-n}\alpha_n
 =\operatorname{CT}{:}(\partial_\theta A)^2{:}.
\]
This proves \eqref{eq:vol2-local-quartic-charge}, including its
derivative correction and quadratic subtraction.
\end{proof}

The isometry in the proof sends the orthonormal Schur basis
\[
 s_\mu(p)=\sum_{\nu\vdash |\mu|}
                  \frac{\chi_\mu(\nu)}{z_\nu}p_\nu
\]
to \(S_\mu\). One can check this directly on finite alternants:
multiplication by \(p_n\) raises an occupied exponent by \(n\),
and reordering supplies the exterior sign. Character orthogonality
gives the stated normalization. Formula~\eqref{eq:vol2-partition-charges}
also gives \(V s_\mu=(\sum_{(i,j)\in\mu}(j-i))s_\mu\),
consistent with the two extreme partitions above.

\section{A positive local preparation}

Let \(d\) be the number of diagonal boxes of \(\mu\), and put
\[
 p_i=\mu_i-i+\tfrac12,\qquad
 h_i=\mu_i'-i+\tfrac12\qquad(1\leq i\leq d).
\]
These positive half-integers are the added particle energies and
the absolute energies of the holes. Separating the two occupation
differences in \eqref{eq:vol2-partition-charges} gives
\begin{equation}\label{eq:vol2-particle-hole-charges}
 E=\sum_i(p_i+h_i),\qquad
 Q_2=\sum_i(p_i^2-h_i^2),\qquad
 Q_3=\sum_i(p_i^3+h_i^3).
\end{equation}
In particular \(Q_3\geq0\), and Cauchy--Schwarz yields
\begin{equation}\label{eq:vol2-charge-inequality}
 Q_2^2\leq E Q_3.
\end{equation}
The two vectors used here have entries
\((\sqrt{p_i},\sqrt{h_i})\) and
\((p_i^{3/2},-h_i^{3/2})\).

\begin{theorem}\label{thm:vol2-local-stable-state}
For real \(g\), define
\begin{equation}\label{eq:vol2-local-stable-hamiltonian}
 H_g=H_0+gV+\frac{g^2}{8}Q_3.
\end{equation}
It is essentially self-adjoint on finite-energy vectors, with its
real diagonal closure in the partition basis, and
\(H_g\geq H_0/2\). For every \(T>0\),
\[
 1\leq Z(T,g)=\operatorname{Tr}e^{-TH_g}\leq Z_0(T/2)<\infty.
\]
The normalized trace is a positive state on bounded operators.
Finite Fourier Wick polynomials also have finite expectations.
The unnormalized heat kernels sew exactly, and their dependence
on real \(g\) is smooth in trace norm.
\end{theorem}

\begin{proof}
All charges are real diagonal in an orthonormal basis. Their sum
has maximal domain
\[
 \left\{(v_\mu):\sum_\mu|H_g(\mu)|^2|v_\mu|^2<\infty\right\}.
\]
Finite-support vectors are a core for this multiplication operator.
By \eqref{eq:vol2-charge-inequality},
\[
 |gV|\leq\frac{|g|}{2}\sqrt{E Q_3}
       \leq\frac E2+\frac{g^2Q_3}{8}.
\]
The last inequality is the nonnegativity of
\((\sqrt{E/2}-|g|\sqrt{Q_3/8})^2\).
It proves the lower bound and hence trace class by summing
\(e^{-TH_g(\mu)}\leq e^{-T|\mu|/2}\).
The vacuum contributes one. Polynomial energy bounds extend the
trace domain to the stated unbounded insertions.

The spectral calculus gives
\(e^{-T_1H_g}e^{-T_2H_g}=e^{-(T_1+T_2)H_g}\).
On compact real coupling intervals, each derivative of order \(k\)
of a diagonal exponential is bounded by a constant times
\((1+E)^{3k}e^{-TE/2}\). Indeed
\(|Q_2|\leq E^2\) and \(Q_3\leq E^3\) follow from
\eqref{eq:vol2-particle-hole-charges}. This dominating sequence is
summable against \(p(E)\). Termwise differentiation is therefore
valid in trace norm, at every order.
\end{proof}

The correction is the integral of the local density
\eqref{eq:vol2-local-quartic-charge}. It first appears at order
\(g^2\). The derivatives of the resulting state agree with the
formal expansion of this modified Hamiltonian.

The real smoothness does not imply a holomorphic trace-class family
near \(g=0\). For \(g=i\gamma\ne0\), the real part of
\(H_g\) is \(H_0-\gamma^2Q_3/8\). On a one-row partition
of size \(E\),
\(Q_3=(E-1/2)^3+1/8\), so the corresponding exponential is
unbounded. Any holomorphic trace-class extension agreeing on a real
interval would have these diagonal matrix entries, by uniqueness
of scalar holomorphic continuation. Such an extension is impossible.

There is also a convergent product expression for the positive trace.
For \(r\in\mathbb Z_{\geq0}+1/2\), set
\[
 \varepsilon_+(r)=r+\frac g2r^2+\frac{g^2}{8}r^3,
 \qquad
 \varepsilon_-(r)=r-\frac g2r^2+\frac{g^2}{8}r^3.
\]
Each is at least \(r/2\). The grand-canonical product
\begin{equation}\label{eq:vol2-local-grand-product}
 \mathcal Z_{T,g}(y)=\prod_{r>0}
 \bigl(1+y e^{-T\varepsilon_+(r)}\bigr)
 \bigl(1+y^{-1}e^{-T\varepsilon_-(r)}\bigr)
\end{equation}
converges absolutely and uniformly on \(|y|=1\). The sum of
the magnitudes of its nonconstant factors is bounded by
\(2\sum_{r>0}e^{-Tr/2}\). Consequently
\begin{equation}\label{eq:vol2-charge-projected-trace}
 Z(T,g)=[y^0]\mathcal Z_{T,g}(y)
       =\frac1{2\pi}\int_0^{2\pi}\mathcal Z_{T,g}(e^{i\theta})\,d\theta.
\end{equation}
Each factor records one particle or one hole. The exponent of \(y\)
records their charge difference, so its zero coefficient retains
exactly the partition basis of \(\mathcal F\). Termwise extraction
is justified by absolute convergence and gives the spectral trace.
At \(g=0\), the expression reduces to \(Z_0(T)\). Omitting the
charge projection would change the Hilbert space.

\section{The sharp coefficient and the critical state}

To determine what the coefficient \(1/8\) achieves, consider
the larger real family
\[
 H_{g,c}=H_0+\frac g2Q_2+c g^2Q_3.
\]
It is defined by the same diagonal closure and local density.

\begin{proposition}\label{prop:vol2-sharp-stabilization}
For \(c>0\), the largest constant \(a_c\) such that
\(H_{g,c}\geq a_cH_0\) for every real \(g\) is
\[
 a_c=1-\frac1{16c}.
\]
Nonnegativity for every coupling holds exactly for \(c\geq1/16\).
The uniform half-energy bound holds exactly for \(c\geq1/8\).
For \(c>1/16\), the heat operators and normalized densities
are continuous in \(g\) in trace norm.
\end{proposition}

\begin{proof}
For a nonempty partition, complete the square:
\[
 H_{g,c}
 =cQ_3\left(g+\frac{Q_2}{4cQ_3}\right)^2
       +E-\frac{Q_2^2}{16cQ_3}
 \geq\left(1-\frac1{16c}\right)E.
\]
The vacuum charges vanish. On the row partition \((N)\), put
\(p=N-1/2\). Then
\[
 E=N,\qquad Q_2=p^2-\tfrac14=N(N-1),\qquad
 Q_3=p^3+\tfrac18.
\]
The ratio \(Q_2^2/(EQ_3)\) tends to one. At
\(g=-Q_2/(4cQ_3)\), the ratio \(H_{g,c}/E\) therefore tends
to \(a_c\), proving optimality. For \(c<0\), the energy of
a nonempty partition becomes negative at large \(|g|\).
For \(c=0\), a row of length at least two has a nonzero cubic
charge and becomes negative for one sign of \(g\).
This proves both coefficient thresholds.

If \(c>1/16\), the diagonal heat entries are bounded uniformly
by \(e^{-Ta_cE}\). Their sum is \(Z_0(Ta_c)\).
For \(g_j\to g\), differences converge entrywise and are
dominated by twice this summable sequence. Their sum of absolute
values is precisely the trace norm. The traces also converge and
are at least one, so division gives continuity of the densities.
\end{proof}

At the critical value, fixed-coupling trace class survives but
this continuity fails.

\begin{theorem}\label{thm:vol2-critical-state-escape}
Let \(c=1/16\), \(K_g(T)=e^{-TH_{g,c}}\), and
\(\rho_g(T)=K_g(T)/\operatorname{Tr}K_g(T)\).
For every real \(g\) and \(T>0\), the heat operator is positive
and trace class. As \(g\to0\), it converges strongly to
\(e^{-TH_0}\). Nevertheless, along
\[
 g_N=-\frac4{N-1/2},\qquad N\geq4,
\]
one has
\begin{align*}
 \operatorname{Tr}K_{g_N}(T)&\longrightarrow\infty,\\
 \|K_{g_N}(T)-K_0(T)\|_1&\longrightarrow\infty,\\
 \|\rho_{g_N}(T)-\rho_0(T)\|_1&\longrightarrow2.
\end{align*}
\end{theorem}

\begin{proof}
The particle and hole contributions are now
\[
 \varepsilon_\pm(r)=r\left(1\pm\frac{gr}{4}\right)^2\geq0.
\]
For fixed \(g\ne0\) and \(r\geq8/|g|\), both are at least
\(r\). Each lower particle or hole mode can occur only once.
Thus \(H_{g,c}\geq E-C_g\), where the finite constant is
\[
 C_g=\sum_{0<r<8/|g|}
 \bigl((r-\varepsilon_+(r))_+
       +(r-\varepsilon_-(r))_+\bigr),\qquad x_+=\max(x,0).
\]
It follows that \(\operatorname{Tr}K_g(T)\leq e^{TC_g}Z_0(T)\).
At \(g=0\), trace class holds directly. Nonnegativity bounds
\(\|K_g(T)\|\) by one. Its entries converge on every fixed
partition vector, so density of their finite span proves strong
convergence on the Hilbert space.

Set \(p=N-1/2\). For the row partitions \((N+k)\), with
\(0\leq k\leq\lfloor\sqrt p\rfloor\), the energy at \(g_N\) is
\begin{equation}\label{eq:vol2-critical-row-energy}
 H_{g_N,1/16}(N+k)
 =\frac{(p+k)k^2}{p^2}
       +\frac12\left(1+\frac1{2p}\right)^2<3.
\end{equation}
The first term is at most \(1+1/\sqrt p<2\), and the second
is less than one for \(p\geq7/2\). These distinct states give
\[
 \operatorname{Tr}K_{g_N}(T)
 \geq(1+\lfloor\sqrt p\rfloor)e^{-3T}\longrightarrow\infty.
\]
The trace difference is a lower bound on the trace-norm difference,
which proves the second limit.

For a fixed finite energy projection \(F\),
\(\operatorname{Tr}(F\rho_{g_N}(T))\to0\), since its numerator
is bounded by \(\operatorname{rank}F\) and the denominator diverges.
Given \(\eta>0\), choose \(F\) with
\(\operatorname{Tr}(F\rho_0(T))>1-\eta\).
Both densities are diagonal and have trace one. Summing absolute
differences on \(F\) and its complement gives
\[
 \|\rho_{g_N}(T)-\rho_0(T)\|_1
 \geq2\left|\operatorname{Tr}(F\rho_{g_N}(T))
              -\operatorname{Tr}(F\rho_0(T))\right|.
\]
The lower limit is at least \(2(1-\eta)\), and the norm is at
most two. Let \(\eta\) decrease to zero.
\end{proof}

The critical example distinguishes strong convergence of heat
operators from convergence of the prepared states. Above the
critical coefficient, uniform energy decay prevents the mass
from escaping to partitions of increasing size. The coefficient
\(1/8\) gives the particularly simple bound \(H_g\geq H_0/2\).

\section{The interacting current and the remaining state map}

Return to the local Hamiltonian \eqref{eq:vol2-local-stable-hamiltonian}.
Its state preserves energy, so a two-current expectation vanishes
unless the two indices sum to zero. The fermionic occupations give
its remaining coefficients explicitly. For \(r>0\), put
\[
 f_r(y)=\frac{y e^{-T\varepsilon_+(r)}}
                   {1+y e^{-T\varepsilon_+(r)}},\qquad
 f_{-r}(y)=\frac1{1+y^{-1}e^{-T\varepsilon_-(r)}}.
\]
These functions record occupancy of a positive mode and occupancy
of a negative mode in the filled sea. Their denominators do not
vanish on \(|y|=1\). Define, for \(n>0\),
\begin{equation}\label{eq:vol2-interacting-current-coefficient}
 A_n(T,g)=
 \frac{[y^0]\left(\mathcal Z_{T,g}(y)
       \sum_{r\in\mathbb Z+1/2}f_{r-n}(y)(1-f_r(y))\right)}{Z(T,g)}.
\end{equation}
For fixed \(n\), the tails decay exponentially and uniformly on
the unit circle. The finitely many modes crossing the sea boundary
are harmless.

\begin{proposition}\label{prop:vol2-interacting-current-pole}
The coefficient \(A_n(T,g)\) is
\(\varphi_{T,g}(\alpha_n\alpha_{-n})\). If
\(c_n=A_n-n\), then \(c_n\geq0\) and
\begin{equation}\label{eq:vol2-interacting-current-bound}
 c_n\leq C_T e^{-Tn/4},\qquad
 C_T=\sum_{E\geq0}E p(E)e^{-TE/4}<\infty.
\end{equation}
The current two-point function is
\begin{equation}\label{eq:vol2-stable-current-kernel}
 \varphi_{T,g}(J(z)J(w))
 =\frac{\lambda}{(z-w)^2}
   +\frac{\lambda}{zw}\sum_{n\geq1}c_n(T,g)
          \bigl((w/z)^n+(z/w)^n\bigr).
\end{equation}
Its radially ordered domain includes \(e^{-T/4}<|w/z|<1\).
The remainder and all fixed holomorphic derivatives extend across
\(z=w\ne0\).
\end{proposition}

\begin{proof}
In the grand-canonical occupation measure, a raising move from
\(r-n\) to \(r\) has weight \(f_{r-n}(1-f_r)\).
By \eqref{eq:vol2-fermion-current}, the squared norm of this move
is the corresponding contribution to \(\alpha_n\alpha_{-n}\).
Different moves from a fixed occupation vector give orthogonal
vectors. Summing their expectations therefore gives the numerator
in \eqref{eq:vol2-interacting-current-coefficient}. The insertion
is charge neutral, so taking the zero-charge coefficient commutes
with it. Absolute convergence justifies the extraction.

The commutator gives
\(c_n=\varphi_{T,g}(\alpha_{-n}\alpha_n)\geq0\).
This positive operator is at most \(H_0\), and it kills energy
spaces below \(n\). The half-energy bound and \(Z(T,g)\geq1\)
give
\[
 c_n\leq\sum_{E\geq n}E p(E)e^{-TE/2}
       \leq e^{-Tn/4}\sum_{E\geq0}E p(E)e^{-TE/4}.
\]
The constant is finite by the free partition estimate.
Inserting \(A_n=n+c_n\) into the current series yields
\eqref{eq:vol2-stable-current-kernel}. The vacuum sum supplies
the pole; the exponential bound makes the two remainder series
normally convergent in a neighborhood of \(|w/z|=1\).
\end{proof}

The modified preparation retains the current pole while changing
the regular coefficients. It need not be Gaussian as a state of
the bosonic current algebra. Its operator sewing and scalar
normalization are fixed by the spectral preparation, including
the charge projection. A conformal shift \(-1/24\) multiplies
its unnormalized kernels by \(e^{T/24}\) and leaves normalized
correlations unchanged.

For the unchanged cubic interaction, the outstanding construction
is a state-valued Borel continuation with a specified Laplace
prescription on a chosen observable star-algebra and sewing
completion. It must give a normalized positive functional on every
allowed observable and preserve the unnormalized cylinder
composition. The scalar Borel boundary does not decide whether
such a prescription exists. The positive states of the modified
local Hamiltonian differ from that interaction at order \(g^2\).

A bulk realization has a further requirement. It must construct
the boundary preparation from independently defined bulk fields,
counterterms and a boundary condition, then intertwine the current,
collision and sewing operations on their declared domains.
An equality of characters or a matching determinant alone does
not supply that state map.

\chapter{Native fields and the boundary action}
\label{ch:vol2-native-bulk-boundary}

A boundary current can be obtained by integrating fields in a normal
direction. The integration also determines how an interior insertion acts
on that current. These two outputs contain different information.
An eliminated field may have zero projection and still contract nontrivially
with a boundary observable. The resulting operation depends on the
propagator, its supports and its quantum normalization.

We construct this comparison for the Gaussian theory on a three-dimensional
half-space. Its Wilson source then gives an independent incoming junction,
including a charged vacuum and a sewn operation. Finite regular jets can be
retained by explicit changes of representatives. The completions used for
these constructions will be specified separately: a smooth smearing, a
polynomial state and a formal series of states have different pole bounds.

\section{The half-space action}

Put \(M=\mathbb R_{\geq0,t}\times\mathbb C_z\), with
\(z=x+iy\) and orientation \(dt\wedge dx\wedge dy\).
The partial forms on \(M\) are generated by \(dt,d\bar z\). Write
\[
 D=dt\,\partial_t+d\bar z\,\partial_{\bar z},\qquad
 \partial=\partial_z,\qquad
 \partial_{\bar z}=\tfrac12(\partial_x+i\partial_y).
\]
The differential \(D\) has degree one, and \(\partial\) has degree zero.
They commute. Fields are smooth up to the boundary. Variations and test
fields have compact support, which may meet that boundary.

For partial forms of compact support use the traces
\begin{equation}\label{eq:vol2-native-traces}
 \tau_M(a)=\frac1{2\pi i}\int_M dz\wedge a,
 \qquad
 \tau_\partial(b)=\frac1{2\pi i}\int_{\mathbb C}dz\wedge b.
\end{equation}
They select partial form degree two and antiholomorphic degree one,
respectively. If \(i:\mathbb C\hookrightarrow M\) is the boundary inclusion,
these conventions give \(\tau_M(Da)=\tau_\partial(i^*a)\).
The sign includes both the outward normal and the position of \(dz\).

Let \(A=\mathbb C[u^a_n:1\leq a\leq r,\ n\geq0]\) be a differential
polynomial algebra, with \(\partial u^a_n=u^a_{n+1}\).
A finite-order matrix differential operator
\[
 H^{ab}(\partial)=\sum_{j=0}^N H_j^{ab}(u,\partial u,\ldots)\partial^j
\]
specifies generator brackets by
\(\{u^a_\lambda u^b\}=H^{ba}(\lambda)\).
Here \(\lambda\) records a holomorphic derivative of a contact.
The extension to all differential polynomials is
\begin{equation}\label{eq:vol2-native-master-formula}
 \{f_\lambda g\}=
 \sum_{a,b,m,n}
 \frac{\partial g}{\partial u^b_n}(\lambda+\partial)^n
 H^{ba}(\lambda+\partial)(-\lambda-\partial)^m
 \frac{\partial f}{\partial u^a_m}.
\end{equation}
Every operator acts on the expression to its right. The sum is finite on
each pair of polynomials.

We call \(H\) Hamiltonian when this bracket is skew and satisfies Jacobi.
Equivalently, \(H^{ab}=-(H^{ba})^*\), where the formal adjoint reverses
composition and sends \(\partial\) to \(-\partial\), and Jacobi holds on
generator triples. The master formula then propagates the identities to
all polynomials \cite[Theorem 1.15 and Proposition 1.16]{BarakatDeSoleKac2009}.
In particular it gives
\[
 \{\partial f_\lambda g\}=-\lambda\{f_\lambda g\},\qquad
 \{f_\lambda\partial g\}=(\lambda+\partial)\{f_\lambda g\}.
\]
The scalar operator \(H=\kappa\partial\) is the current example.

Introduce superfields \(\Phi^a\) and \(\eta_a\) of total degrees zero
and one. In their expansions in \(1,dt,d\bar z,dt\,d\bar z\), their
coordinate degrees are
\[
 \deg\Phi^a=(0,-1,-1,-2),\qquad
 \deg\eta_a=(1,0,0,-1).
\]
The degree-one zero-form in \(\eta_a\) is the gauge ghost. The negative
degrees are antifields. The BV two-form and action are
\begin{equation}\label{eq:vol2-native-action}
 \omega_{\mathrm{BV}}=\tau_M(\delta\Phi^a\delta\eta_a),
 \qquad
 S_H=\tau_M\left(\eta_aD\Phi^a+
              \tfrac12\eta_aH^{ab}(\partial;\Phi)\eta_b\right).
\end{equation}
The two-form has degree \(-1\), and the action has degree zero.
Its inverse pairing defines the degree-one BV bracket on local
functionals. This is the local Poisson vertex sigma model
\cite[\S\S2.2 and 4.1]{KhanZeng2025}.

\begin{proposition}\label{prop:vol2-native-classical}
For a Hamiltonian operator \(H\), the action
\eqref{eq:vol2-native-action} satisfies the classical master equation
on the boundary condition
\[
                         i^*\eta=0.
\]
This is a Lagrangian condition on the whole boundary superfield and is
preserved by the BV differential.
\end{proposition}

\begin{proof}
The kinetic square vanishes by \(D^2=0\). Its bracket with the interaction
is a total partial derivative. Its boundary trace contains \(i^*\eta\),
and vanishes. Pair the interaction bracket with three independent odd
test fields and integrate every holomorphic derivative by parts.
The remaining coefficient is the generator Jacobi expression
\[
 \{u^a_\lambda\{u^b_\mu u^c\}\}
 -\{u^b_\mu\{u^a_\lambda u^c\}\}
 -\{\{u^a_\lambda u^b\}_{\lambda+\mu}u^c\}.
\]
It vanishes by the Hamiltonian-operator identity. This proves
\(\{S_H,S_H\}_{\mathrm{BV}}=0\), and the graded Jacobi identity gives
\(Q^2=0\) for \(Q=\{S_H,-\}_{\mathrm{BV}}\).

Variation has boundary term
\(-\tau_\partial(\eta_a\delta\Phi^a)\).
The condition sets one complementary summand of this nondegenerate
pairing to zero and leaves the other free. It is Lagrangian.
The interaction part of \(Q\eta\) is quadratic in \(\eta\) and its
holomorphic derivatives. Its tangential pullback vanishes on the condition,
and \(\bar\partial\) preserves the same condition.
\end{proof}

For \(H=\kappa\partial\), write the physical fields as a scalar \(\phi\)
and a partial one-form \(a\). Then
\[
 S^{\mathrm{cl}}_\kappa=\tau_M\left(aD\phi+\tfrac\kappa2a\partial a\right),
 \qquad
 D\phi+\kappa\partial a=0,\qquad Da=0.
\]
The gauge variations are \(\delta\phi=\kappa\partial\epsilon\) and
\(\delta a=-D\epsilon\), with \(\epsilon\) zero at the boundary.
For \(\kappa\ne0\), the one-form
\(\mathcal A=a-\kappa^{-1}\phi\,dz\) identifies this action with
\(-\kappa(4\pi i)^{-1}\int\mathcal A\wedge d\mathcal A\).
Indeed expansion and one integration by parts give the displayed scalar
action, with boundary term zero under \(i^*a=0\). Thus the current model
has a pure abelian Chern--Simons realization with this normalization.

The finite Hamiltonian Lie algebra \(\mathfrak p_N\) of
Chapter~\ref{ch:vol2-hamiltonian-jets} supplies another classical example.
If \([e_a,e_b]=f_{ab}^{\ c}e_c\), set
\(\{u^a_\lambda u^b\}=f_{ab}^{\ c}u^c\).
This has differential order zero and satisfies Jacobi by the Lie identity.
Its action is the shifted cotangent action with those coefficients.
The finite-dimensional quantum measure of that chapter concerns its
algebraic model. Quantizing all spacetime modes and identifying a physical
boundary algebra remain additional comparisons.

\section{Regular Gaussian observables}

Keep \(\kappa\) as an indeterminate and put
\(R=\mathbb C[\kappa,\hbar]\). For the scalar Gaussian model use the
tangent complex
\[
 E=\left(\Omega_{\mathrm{HT}}^\bullet e_\Phi
       \oplus\Omega_{\mathrm{HT}}^\bullet e_\eta\right)\otimes\mathbb C[\kappa],
 \qquad |e_\Phi|=0,\quad |e_\eta|=-1.
\]
For a partial form \(\alpha\) of degree \(p\), its differential is
\begin{equation}\label{eq:vol2-native-tangent}
 \begin{split}
 Q(\alpha e_\Phi)&=D\alpha e_\Phi,\\
 Q(\alpha e_\eta)&=D\alpha e_\eta
            +(-1)^{p+1}\kappa\partial\alpha e_\Phi.
 \end{split}
\end{equation}
The two off-diagonal terms in \(Q^2\) have opposite signs. Impose the
relative boundary condition on every tangential \(e_\eta\) component.
The coordinate and tangent degrees are dual conventions.

The degree-minus-one pairing is
\[
 B(\alpha e_\Phi,\beta e_\eta)=\tau_M(\alpha\beta),\qquad
 B(\alpha e_\eta,\beta e_\Phi)
       =-(-1)^{\deg\beta}\tau_M(\alpha\beta),
\]
and it is zero on equal summands. Let \(E_c(U)\) denote the compactly
supported fields on an open \(U\), with the boundary condition when
needed. On \(C(U)=E_c(U)[1]\), write
\[
 d(s\xi)=-sQ\xi,\qquad
 b(s\xi,s\zeta)=(-1)^{|\xi|+1}B(\xi,\zeta).
\]
Thus \(b\) is a graded-symmetric form of degree one.

At fixed field number, a regular observable is a compact smooth joint
kernel in the corresponding graded symmetric tensor power of \(C(U)\).
Complete that tensor at fixed compact support in its smooth seminorms,
then take the inductive limit over compact supports. Denote the result
by \(\operatorname{Sym}_b^m C(U)\). Equivalently one can use the completed
bornological tensor of the compact smooth spaces. The direct sum over
field number remains algebraic. Define
\begin{equation}\label{eq:vol2-native-quantobs}
 \mathcal O^q(U)=
 \left(\bigoplus_{m\geq0}\operatorname{Sym}_b^m C(U)[\hbar],
                       d+\hbar\Delta_b\right).
\end{equation}
Here \(\Delta_b\) contracts an unordered pair with \(b\), with its Koszul
sign. Restriction of a smooth joint kernel to a diagonal and compact
integration define this operation continuously.

Compatibility of the pairing with \(Q\) gives
\([d,\Delta_b]=0\), and odd pair contractions give \(\Delta_b^2=0\).
Consequently \eqref{eq:vol2-native-quantobs} is a cochain complex.
Disjoint opens have products by extension by zero and concatenation of
kernels. There are no cross BV contractions on disjoint supports, so
these products commute with the differential and compose associatively.
The constant observable is their unit.

The descent statement uses the free-field theorem for an elliptic complex
of finite total rank, topological in the normal direction, with a local
Lagrangian elliptic boundary condition
\cite[Definitions 2.1--2.3 and Theorems 4.1--4.2]{GwilliamRabinovichWilliams2021}.
Our tangential differential is the triangular extension of two Dolbeault
complexes, and the normal factor is de Rham. These are precisely the
normally topological Gaussian data. The theorem gives a factorization
algebra of these regular observables and its boundary restriction.
The maps below construct the collar comparison explicitly.

Point evaluation has a delta-distribution kernel and is absent from
\eqref{eq:vol2-native-quantobs}. Its extension to local composite
observables is a renormalization problem. Likewise, the formal completion
\(\mathcal O^q(U)^{\wedge}_{\hbar}\) allows the field-number bound to grow
with the coefficient order. It is larger than completing each fixed
field number and then taking their algebraic sum.

\section{The normal contraction and its quantum correction}

Choose a smooth profile \(\rho\) of integral one supported in
\([0,\varepsilon)\), and put \(F(t)=\int_0^t\rho(s)\,ds\).
For a tangential open \(V\), split its compact fields into
\[
 L=\Omega_c^{0,\bullet}(V)e_\Phi\otimes\mathbb C[\kappa],
 \qquad G=\Omega_c^{0,\bullet}(V)e_\eta\otimes\mathbb C[\kappa].
\]
The tangential differential is
\(d_\partial(\ell,g)=(d_L\ell+Q_{\mathrm{rel}}g,d_Gg)\), where
\[
 Q_{\mathrm{rel}}(\alpha e_\eta)
       =(-1)^{\deg\alpha+1}\kappa\partial\alpha e_\Phi,
 \qquad d_LQ_{\mathrm{rel}}+Q_{\mathrm{rel}}d_G=0.
\]
Write a collar field as \(\gamma(t)+dt\,\beta(t)\), with
\(\gamma_G(0)=0\). Its differential is
\(d_\partial\gamma+dt(\partial_t\gamma-d_\partial\beta)\).

Define maps to and from \((G[-1],-d_G)\) by
\begin{equation}\label{eq:vol2-native-normal-maps}
 P(\gamma+dt\,\beta)=\int_0^\infty\beta_G(s)\,ds,
 \qquad Ig=dt\,\rho g+(F-1)Q_{\mathrm{rel}}g.
\end{equation}
The second term in \(I\) is required. Its normal derivative cancels the
unwanted term \(-dt\,\rho Q_{\mathrm{rel}}g\) in the differential of the
first term. In particular \(QI=-Id_G\) and \(PI=1\).

\begin{lemma}\label{lem:vol2-native-normal}
The degree-minus-one map
\[
 h(\gamma+dt\,\beta)=-\int_t^\infty\beta(s)\,ds
                  +(1-F(t))\int_0^\infty\beta_G(s)\,ds
\]
satisfies
\[
 Qh+hQ=1-IP,
 \qquad hI=Ph=h^2=0.
\]
These identities hold continuously on compact smooth joint kernels and
retain compact tangential support. No inverse of \(\kappa\) is used.
\end{lemma}

\begin{proof}
The correction sets the \(G\)-component of \(h\) to zero at \(t=0\).
All tail integrals vanish above the supports of the input and \(\rho\).
Differentiating the tail gives \(\beta(t)\), while differentiating
\(1-F\) gives \(-\rho\). Substitution in the collar differential proves
\(Qh+hQ=1-IP\). Since \(I\) has normal component \(\rho g\), its two
terms under \(h\) cancel. The output of \(h\) has no normal one-form,
so \(Ph=h^2=0\). Integrals in a compact normal interval are bounded in
each smooth tangential seminorm. The same estimates with other kernel
variables as parameters prove the completed-kernel statement.
\end{proof}

Let \(\mathfrak a(V)=\Omega_c^{0,\bullet}(V)\otimes\mathbb C[\kappa]\).
Fix \(\zeta^2=-1\). The map
\(\alpha\mapsto\zeta(-1)^q\alpha e_\eta\), for Dolbeault degree \(q\),
identifies \((\mathfrak a(V),\bar\partial)\) with \((G[-1],-d_G)\).
After suspension, \eqref{eq:vol2-native-normal-maps} and \(-sh\) give
a contraction \((i,p,h_C)\) from \(C\) to \(\mathfrak a(V)[1]\).

The pairing induced by \(i\) is determined by
\begin{equation}\label{eq:vol2-native-current-cocycle}
 \mu_\kappa(\alpha,\beta)
       =-\kappa\tau_\partial(\alpha\partial\beta),
 \qquad
 b(i(s\alpha),i(s\beta))=(-1)^q\mu_\kappa(\alpha,\beta).
\end{equation}
To check the normalization, only \(q+r=1\) contributes. The two terms
of \(I\) pair to
\(\kappa\tau_\partial(\alpha\partial\beta)\), since
\(2\int\rho(1-F)\,dt=1\) and
\(\tau_M(dt\,\omega)=-\int\tau_\partial(\omega)\,dt\).
The two phase factors give \(\zeta^2(-1)^{q+r}=1\), and suspension
gives \((-1)^{q+1}\). This proves the formula.

Write \(\mu^s=i^*b\) for the suspended pairing and set
\[
 \mathcal A_\kappa(V)=
 \left(\operatorname{Sym}_b(\mathfrak a(V)[1])[\hbar],
                  d+\hbar\Delta_{\mu^s}\right).
\]
Integration by parts makes \(\mu_\kappa\) a graded-skew cocycle.
The inclusion \(i^q=\operatorname{Sym}(i)\) therefore preserves the quantum
differential. The projection needs a correction.

\begin{theorem}\label{thm:vol2-native-quantum-collar}
Regular Gaussian collar observables contract onto
\(\mathcal A_\kappa(V)\) over \(R\). Both chain maps preserve the unit
and products on tangentially disjoint regions. They are continuous and
polynomial in \(\hbar\) at fixed field number, and extend coefficientwise
to the \(\hbar\)-adic completion.
\end{theorem}

\begin{proof}
Put \(e=ip\). Average \(h_C\otimes1+e\otimes h_C\) with its graded
transpose to obtain \(H_2\). Then
\(dH_2+H_2d=1-e\otimes e\). The even pairing \(P_2=-bH_2\) satisfies
\[
 \delta P_2=b-p^*\mu^s,
 \qquad [d,\partial_{P_2}]=\Delta_b-\Delta_{p^*\mu^s},
\]
where \(\partial_{P_2}\) contracts an unordered pair.
Constant contractions graded-commute. Hence for
\(W=\exp(\hbar\partial_{P_2})\) the conjugation has only one commutator:
\[
 W(d+\hbar\Delta_b)W^{-1}=d+\hbar\Delta_{p^*\mu^s}.
\]
Take
\begin{equation}\label{eq:vol2-native-quantum-maps}
 i^q=\operatorname{Sym}(i),\qquad
 p^q=\operatorname{Sym}(p)W.
\end{equation}
The pairing \(P_2\) is zero on \(i\otimes i\), so \(Wi^q=i^q\) and
\(p^qi^q=1\).

Split \(C=\operatorname{im}i\oplus\ker p\). Extend \(h_C\) as a derivation
of the symmetric algebra and divide by the positive number of \(\ker p\)
factors. Set the resulting \(H_0\) to zero on the retained factors alone.
Then \([d,H_0]=1-\operatorname{Sym}(e)\). Only retained factors enter
\(\Delta_{p^*\mu^s}\), so it graded-commutes with \(H_0\). Consequently
\[
 H^q=W^{-1}H_0W,
 \qquad (d+\hbar\Delta_b)H^q+H^q(d+\hbar\Delta_b)=1-i^qp^q.
\]
Each contraction lowers field number by two. All exponentials therefore
terminate on every polynomial. Normal integration and diagonal restriction
are continuous on smooth joint kernels. Their tangential supports remain
those of the uncontracted variables. In particular \(P_2\) has no cross
term on disjoint tangential supports, proving the product assertion.
At a fixed coefficient of \(\hbar\), only finitely many coefficients of a
formal input contribute. This proves the completed assertion as well.
\end{proof}

\section{Currents and supported interior insertions}

The finite-circle-spectrum sector on a centered disk is the algebraic
sum of its rotation eigenspaces. Choose a radial annular cutoff \(\chi\),
equal to one near the inner circle and zero beyond the outer circle.
The compact forms \(\alpha_n=(\bar\partial\chi)z^{-n}\), \(n\geq1\), obey
\[
 \tau_\partial(z^m\alpha_n)=\delta_{m,n-1},\qquad
 -\partial\alpha_n=n\alpha_{n+1}
                    -\bar\partial(z^{-n}\partial\chi).
\]
Stokes' formula proves the first identity. Differentiation proves the second.
They fix the creation modes and translation, including their normalization.

\begin{proposition}\label{prop:vol2-native-heisenberg}
The Gaussian boundary has polynomial vacuum states
\[
 V_\kappa=R[x_1,x_2,\ldots],\qquad
 J_{-n}=x_n,\quad J_n=\hbar\kappa n\partial_{x_n}\ (n>0),\quad J_0=0.
\]
Its current satisfies \(J(z)J(w)\sim\hbar\kappa(z-w)^{-2}\).
Translation is \(T=\sum_{n>0}n x_{n+1}\partial_{x_n}\).
These formulas include the nonsimple zero-level specialization.
\end{proposition}

\begin{proof}
The cocycle \eqref{eq:vol2-native-current-cocycle} on Laurent test functions
is \(-\kappa\operatorname{Res}(f\partial g\,dz)\). Thus the contour modes
obey \([J_m,J_n]=\hbar\kappa m\delta_{m+n,0}\).
The compact representatives supply the polynomial creation states and
\([T,J_m]=-mJ_{m-1}\). Normal-ordered products of derivatives of \(J\)
then give every state field, with the vacuum and translation identities.
Wick contraction proves locality and the stated double pole.

The action is quadratic. If its level term is treated as an insertion,
each kinetic propagator joins \(\Phi\) to \(\eta\), while that insertion
has two \(\eta\) legs. Two insertions cannot join each other.
A boundary pair receives exactly one insertion, and its two attachments
cancel the factor \(1/2\). Hence the current level has no interaction
correction. At \(\kappa=0\), the positive modes vanish and the same
polynomial carrier is a commutative differential vertex algebra.
\end{proof}

The current algebra does not select the Hilbert-space preparation of
Chapter~\ref{ch:vol2-infinite-mode-states}. Even in this Gaussian theory,
normalized polynomial correlations leave determinant and framing factors
undetermined. They are additional data of a physical state.

For an interior open \(U\subset V\times(a,b)\) with
\(0<\varepsilon<a<b\), the collar inclusion \(i^q\) has normal support
disjoint from \(U\). Define
\begin{equation}\label{eq:vol2-native-supported-action}
 \beta_{U;V}(f)(a_0)
       =p^q\bigl(\operatorname{ext}_U(f)i^q(a_0)\bigr).
\end{equation}
Its target is the complex of continuous operations on
\(\mathcal A_\kappa(V)\), with Hom differential
\(T\mapsto dT-(-1)^{|T|}Td\). Both comparison maps are cochain maps.
Disjointness removes the cross BV term in the middle product.
Thus \(\beta_{U;V}\) is a cochain map.

This action detects a contraction that linear restriction misses.
Take a unit-integral normal profile \(\rho_{\mathrm{int}}dt\) above the
boundary profile, and compact tangential tests \(\phi,\beta\) of
Dolbeault degrees zero and one with \(\tau_\partial(\phi\beta)=1\).
Set
\[
 u=\zeta^{-1}\kappa s(\rho_{\mathrm{int}}dt\,\phi e_\Phi),
 \qquad v=s\beta.
\]
Then \(p(u)=0\), but
\(h_Cu=\zeta^{-1}\kappa s(\phi e_\Phi)\) on the boundary profile.
Substitution in \(P_2=-bH_2\) gives
\[
 P_2(u,i(v))=\kappa,
 \qquad \beta(u)(v)=\hbar\kappa,
 \qquad \beta(u)(1)=0.
\]
The cochain \(u\) is not assumed closed. This is an operation on the
supported complexes, rather than a new bulk cohomology class.

For a homogeneous linear input, the complete formula is
\[
 \begin{split}
 \beta(u)(v_1\cdots v_r)
 &=p(u)v_1\cdots v_r\\
 &\quad+\hbar\sum_{j=1}^r
 (-1)^{|u|\sum_{\ell<j}|v_\ell|}P_2(u,i(v_j))
                  v_1\cdots\widehat v_j\cdots v_r.
 \end{split}
\]
Its quantum factor remains over every truncated or completed coefficient
ring. There is no division by \(\hbar\).

If \(f\) is supported above \(g\), the normal homotopy of \(g\) remains
below \(f\). Inserting \(1-i^qp^q=dH^q+H^qd\) gives the ordered
composition homotopy
\begin{equation}\label{eq:vol2-native-action-homotopy}
 K(f,g,a_0)=(-1)^{|f|}p^q\bigl(
       \operatorname{ext}(f)H^q(\operatorname{ext}(g)i^q(a_0))\bigr).
\end{equation}
Its Hom differential is the joint projected product minus the successive
actions. The sign moves the degree-minus-one homotopy past \(f\).
Arbitrary motion of supports requires parameter-dependent operations.
The forest coefficient homotopies of
Chapter~\ref{ch:vol2-supported-compositions} retain their contact terms,
but their decoration by native field operations is a further comparison.

\section{Smooth periods and polynomial states}

A smooth compact smearing need not have finitely many holomorphic moments.
For \(0<r<R_0\), define
\[
 E_r=\left\{\sum_{n\geq1}a_nw^{-n}:
                       \sum_{n\geq1}|a_n|r^{-n}<\infty\right\},
 \qquad \mathscr P(D_{R_0})=\operatorname*{colim}_{r<R_0}E_r.
\]
The norm is the displayed weighted sum. For \(r<r'\), the inclusion has
diagonal coefficients \((r/r')^n\) in normalized bases and is nuclear
and compact. A countable cofinal sequence of radii defines the complete
locally convex inductive limit. Its elements are exterior holomorphic
germs vanishing at infinity. Cauchy estimates identify this description
with the coefficient description above.

Put
\[
 \mathscr F(D_{R_0})=
       \bigoplus_{m\geq0}\operatorname{Sym}_\pi^m\mathscr P(D_{R_0})[\hbar].
\]
Here the projective tensor is completed at each fixed \(m\), and the
field-number sum is algebraic. A finite union of disjoint disks uses the
direct sum of their germ spaces. Common inner radii and compact support
bounds are required locally on parameter families.

The exterior Cauchy map on a compact \((0,1)\)-form is
\begin{equation}\label{eq:vol2-native-cauchy}
 \mathcal C(sa)(w)=\tau_\partial\left(\frac{a(z)}{w-z}\right)
             =\sum_{n\geq0}\tau_\partial(z^na)w^{-n-1}.
\end{equation}
It is zero on shifted functions. The expansion is taken outside the
support disk. A radial cutoff \(\chi\) provides a section
\(j_\chi(f)=s(\bar\partial\chi)f\), on any germ stage lying inside its
transition annulus.

\begin{lemma}\label{lem:vol2-native-dolbeault}
At every fixed field number, the symmetric Cauchy map is a continuous
quasi-isomorphism from the compact Dolbeault current complex to the
corresponding symmetric germ space. Its cohomology is in degree zero,
with its Hausdorff quotient topology, and its differential has closed
range.
\end{lemma}

\begin{proof}
The Cauchy convolution on the whole plane is a smooth function
\(\widetilde{\mathcal C}a\) with
\(\bar\partial\widetilde{\mathcal C}a=-a\).
If its exterior germ is zero, it is a compact primitive. There are no
nonzero compact holomorphic functions.
Choose radii
\(r_0<r<r'<r''<R_0\), with the source supported in \(D_{r_0}\) and
\(\chi=1\) through \(D_{r'}\), supported in \(D_{r''}\). Then
\[
 h_\chi(sa)=s\left(\widetilde{\mathcal C}a
                           -(1-\chi)\mathcal C(a)\right),
 \qquad h_\chi(s\phi)=0
\]
satisfies \(d_1h_\chi+h_\chi d_1=1-j_\chi\mathcal C\), where
\(d_1(s\alpha)=-s\bar\partial\alpha\).
The second term is defined outside \(D_r\) and extended by zero where
\(1-\chi=0\). Stokes' formula gives \(\mathcal Cj_\chi=1\).

Cauchy estimates on the annular gap bound every derivative by finitely
many source seminorms and the \(E_r\) norm. Apply this contraction in each
variable of a joint kernel. The tensor homotopy applies \(h_\chi\) to one
variable, the projection to earlier variables and the identity to later
variables, with Koszul signs. The estimates are uniform in the remaining
smooth parameters. Thus the identity holds on completed projective
tensors. Averaging over the finite permutation group gives symmetric
powers. The ranges are kernels of continuous maps to Hausdorff spaces,
and the stagewise sections make the inverse on cohomology continuous.
\end{proof}

The Gaussian and source terms can also be removed explicitly.
Let \(k_1(z)=1/z\) as a locally integrable distribution and
\(k_2=-\partial_zk_1\). On shifted top forms define
\[
 P_{\mathrm C}(sa,sb)=\kappa\tau_z\tau_w
                      \bigl(k_2(z-w)a(z)b(w)\bigr),
 \qquad \gamma(sa)=\tau_z\left(\frac{a(z)}z\right).
\]
They vanish on other form degrees. Their Hom differentials are
\(\delta P_{\mathrm C}=\mu^s\) and \(\delta\gamma=\epsilon_0\), where
\(\epsilon_0(s\phi)=\phi(0)\) and \(\epsilon_0\) is zero on top forms.
Indeed \(\partial_{\bar z}(1/z)=\pi\delta_0\) and
\(\tau(f\,d\bar z)=-\pi^{-1}\int f\) give
\[
 \tau\left(\frac{\bar\partial\phi}{z}\right)=\phi(0),
 \qquad
 \tau_z\bigl(k_2(z-w)\bar\partial\phi(z)\bigr)=\partial_w\phi(w).
\]

For a scalar Wilson charge \(q\), the boundary differential is
\(d_1+\hbar\Delta_{\mu^s}+\hbar q\epsilon_0\).
The exponential
\[
 N_q=\exp\bigl(\hbar\partial_{P_{\mathrm C}}
                       +\hbar q\partial_\gamma\bigr)
\]
conjugates it to \(d_1\). The proof is the same one-commutator calculation
as in Theorem~\ref{thm:vol2-native-quantum-collar}.
Every contraction removes factors. Products of contracted distributions
use disjoint pairs of variables, and therefore exist as distributional
tensor products. They are paired with smooth joint kernels.

The line source on bulk observables is the constant derivation
\[
 \ell(s\xi)=\zeta^{-1}\int_0^\infty\beta_\eta^{(0)}(t,0)\,dt,
 \qquad \xi=\gamma+dt\,\beta.
\]
Its degree is one. The relative boundary condition gives \(\ell d=0\),
and it graded-commutes with constant contractions. Thus
\(d+\hbar\Delta_b+\hbar q\ell\) squares to zero.
Moreover \(\ell=\epsilon_0p\), and the quantum collar maps still give a
contraction onto the charged boundary complex. The Gaussian identity
\(e^{-q\lambda}(d+\hbar\Delta_b)e^{q\lambda}
=d+\hbar\Delta_b+\hbar q\ell\) expresses the Wilson insertion, where
\(\lambda\) is its distributional linear representative. The identity
defines the operation on regular observables without placing \(\lambda\)
inside the smooth observable algebra.

Consequently
\begin{equation}\label{eq:vol2-native-smooth-period}
             \Pi=\operatorname{Sym}(\mathcal C)N_qp^q
\end{equation}
is a continuous quasi-isomorphism onto \(\mathscr F(D_{R_0})\), with zero
target differential. The statement follows by composing the three
explicit contractions. It is a comparison with the germ completion.
For example \(e^{c/w}-1\), for \(c\ne0\), is the Cauchy transform of
\(s(\bar\partial\chi)(e^{c/z}-1)\) and has infinitely many moments.

On a centered disk, rotation gives \(w^{-n}\) weight \(-n\).
At a fixed total negative weight there are finitely many partitions into
positive mode numbers. Hence the finite-circle-spectrum target is exactly
\(R[x_1,x_2,\ldots]\), with \(x_n=w^{-n}\).
The restricted period map remains a quasi-isomorphism. To see this, choose
a primitive of a finite-spectrum cocycle in its full acyclic cone, then
apply the finite Fourier projector onto the weights of that cocycle.
All rotations lie in a common compact support and germ stage, so the
integral exists and commutes with the differential. This gives a
finite-spectrum primitive.

The same map identifies separated current insertions with
\[
 J(z)=\sum_{n\geq1}x_nz^{n-1}+\frac{\hbar q}{z}
                  +\hbar\kappa\sum_{n\geq1}n z^{-n-1}\partial_{x_n}.
\]
A compact radial current smearing centered at \(z\), disjoint from the
inner support, has Cauchy transform \((w-z)^{-1}\). Its charge contraction
is \(\hbar q/z\), and its covariance with an inner period \(f\) is
\(-\hbar\kappa f'(z)\). These three terms give the formula.
Contour integration extracts each mode. Changing the contour radius within
the separated annulus is a cochain homotopy by Stokes' formula.

\section{An incoming junction from the Gauss constraint}

We now set \(\hbar=\kappa=1\) and retain a numerical charge \(q\).
Consider the incoming Wilson ray in complex abelian Chern--Simons theory
\[
 S=-\frac1{4\pi i}\int\mathcal A\wedge d\mathcal A,
 \qquad W_q=\exp\left(q\int_\gamma\mathcal A\right).
\]
On an oriented boundary disk with coordinate \(w\), orient the incoming
line opposite to the outward normal. Variation of the normal multiplier
then imposes
\begin{equation}\label{eq:vol2-native-gauss}
 d_Da=-2\pi iq\,\delta_0,
 \qquad a=a_q+\alpha,
 \qquad a_q=-q\,\frac{dw}{w},
\end{equation}
where \(\alpha\) is smooth and \(\delta_0\) has integral one.
The singularity class is fixed in this statement. Gauge functions vanish
on the boundary circle. There is no quotient by large gauge
transformations and no added monopole sector.

The fluctuation complex is
\[
 \Omega^0_\partial(D)\xrightarrow{d}\Omega^1(D)
                               \xrightarrow{d}\Omega^2(D)
\]
in degrees \(-1,0,1\). It retracts to
\(C^\infty(S^1)/\mathbb C\) in degree zero. One explicit contraction uses
radial Poincar\'e maps \(h_1,h_2\), with
\(dh_1+h_2d=1\), \(dh_2=1\), \(h_1h_2=0\), and an extension \(E\) of
mean-zero boundary functions. Put
\[
 P\alpha=[\operatorname{tr}h_1\alpha],\qquad Ib=dEb,
 \qquad
 H_0\alpha=h_1\alpha-E(P\alpha)
                    -\operatorname{mean}(\operatorname{tr}h_1\alpha),
 \quad H_1=h_2.
\]
The Poincar\'e identities give \(PI=1\) and \(dH+Hd=1-IP\).
For a gauge function vanishing at the circle, the mean term removes the
constant in \(h_1df=f-f(0)\). Thus the contraction respects the relative
domain, not only the unrestricted disk complex.

The reduced form and currents are
\begin{equation}\label{eq:vol2-native-reduced-form}
 \Omega_{\mathrm{red}}
      =\sum_{n>0}\frac{\delta J_{-n}\wedge\delta J_n}{n},
 \qquad J_n=-\frac1{2\pi i}\oint w^na,
 \qquad J_0=q.
\end{equation}
Indeed the smooth solutions are \(a_q+df\), modulo zero-boundary gauge
functions. Stokes' formula gives
\(\Omega(\delta_1f,\delta_2f)
=(2\pi i)^{-1}\oint\delta_1f\,d\delta_2f\), which gives the nonzero
modes in \eqref{eq:vol2-native-reduced-form}. The source fixes the zero mode.

To retain the constraints before reduction, choose Dirichlet functions
\(\epsilon_j\) orthonormal for the energy pairing and put
\[
 u_j=d\epsilon_j,\qquad v_j=2\pi i\,*d\epsilon_j.
\]
For \(\Omega_D(\alpha,\beta)=(2\pi i)^{-1}\int_D\alpha\wedge\beta\),
these obey \(\Omega_D(u_j,v_k)=\delta_{jk}\) and zero equal-type pairings.
Harmonic gradients are symplectically orthogonal to these modes by
Green's formula. Solving the two Dirichlet Poisson equations gives the
exact, coexact and harmonic splitting.

Write the constraint part of \(\alpha\) as \(\sum_j(s_ju_j+p_jv_j)\).
Its moment is
\[
 -\frac1{2\pi i}\int_D\epsilon_j(d_Da+2\pi iq\delta_0)=p_j.
\]
The constraint \(p_j=0\) translates \(s_j\).
Keep finitely many gauge and boundary modes, quantize in the incoming
polarization \(x_n=J_{-n},s_j\), and take the union of the resulting
polynomial state spaces. With odd ghosts \(\nu_j\) of degree one this
defines
\begin{equation}\label{eq:vol2-native-bfv-junction}
 \mathcal J_q=
 \left(\mathbb C[x_1,x_2,\ldots,s_1,s_2,\ldots]
          \otimes\Lambda(\nu_1,\nu_2,\ldots),
       \quad Q_{\mathrm{BFV}}=\sum_j\nu_j\partial_{s_j}\right).
\end{equation}
All polynomials are finite. The differential is the quantized commuting
moment constraint \(\sum_j\nu_jp_j\). No constant gauge ghost occurs for
the chosen zero-boundary gauge group.

\begin{theorem}\label{thm:vol2-native-polynomial-junction}
Evaluation at zero gauge coordinates and ghosts gives a quasi-isomorphism
\[
 \rho_q:\mathcal J_q\longrightarrow F_q=\mathbb C[x_1,x_2,\ldots],
 \qquad \rho_q(f)=f(x,0,0).
\]
Its inverse up to homotopy is the inclusion \(\iota_q(f)=f(x)\).
These maps and their contraction commute with
\[
 J_{-n}=x_n,\qquad J_n=n\partial_{x_n},\qquad J_0=q,
 \qquad T_q=\sum_{n>0}n x_{n+1}\partial_{x_n}+qx_1.
\]
Thus they preserve every current field, the rooted residues and the
identity action of the local vacuum.
\end{theorem}

\begin{proof}
Let \(N\) count total polynomial and exterior degree in \(s,\nu\),
excluding \(x\). With left odd derivatives set
\[
 H=\frac1N\sum_js_j\partial_{\nu_j}\quad(N>0),\qquad H=0\quad(N=0).
\]
The anticommutator of its numerator with \(Q_{\mathrm{BFV}}\) is the
gauge Euler derivation. Therefore
\(Q_{\mathrm{BFV}}H+HQ_{\mathrm{BFV}}=1-\iota_q\rho_q\).
The sum is finite on each state and passes to the algebraic union of
cutoffs. Every displayed current and \(T_q\) acts only on \(x\), so it
commutes with the contraction. Normal ordering and differentiation give
all current fields. Their Laurent coefficients, including coefficients
tested by multipliers before taking principal parts, give the rooted
residues. The vacuum acts as the identity throughout.
\end{proof}

For the unnormalized disk function \(\epsilon=1-r^2\), its energy is
\(2\pi\). Thus \(u=d\epsilon\), \(v=i*d\epsilon\) form one canonical
pair and give the differential \(\nu\partial_s\).
Removing the Wilson source leaves this contractible gauge complex
unchanged but sets \(J_0=0\). It then has zero-mode defect two on the
unit of the charge-two recipient. Gauge contraction alone does not
recover the missing charge.

\section{The charged operation and its collision lattice}

Place a second unit charge at separation \(z\ne0\).
The source in \eqref{eq:vol2-native-gauss} changes to
\(-dw/w-dw/(w-z)\). Contour integration gives the current Ward identity
\[
 J_m^{\mathrm{out}}\mathcal V(z)-\mathcal V(z)J_m^{\mathrm{in}}
       =z^m\mathcal V(z),\qquad m\in\mathbb Z.
\]
In the incoming polarization, its oscillator solution with highest
coefficient one is
\begin{equation}\label{eq:vol2-native-endpoint-kernel}
 K(z)f(x)=\exp\left(\sum_{n>0}\frac{x_nz^n}{n}\right)
                                      f(x_n-z^{-n}).
\end{equation}
The creation exponential supplies the positive-mode Ward term, and the
shift supplies the negative-mode term. On a polynomial input, the shift
is finite. At a fixed output weight, only finitely many creation modes
contribute. Values lie in the product of oscillator-weight spaces.

Translation fixes the remaining scalar. Direct substitution gives
\[
 T_2K-KT_1=\partial_zK+z^{-1}K.
\]
For a charged operation \(h(z)K(z)\), translation covariance requires
\(h'=h/z\). Normal ordering removes each endpoint's self-contraction,
and normalization at \(z=1\) gives
\begin{equation}\label{eq:vol2-native-charged-operation}
                         \mathcal V(z)=zK(z).
\end{equation}
Insertion of negative current modes gives every descendant of the first
input. Explicitly, if \(I(u,z)\) denotes an endpoint operation, then
\[
 I(J_{-n}u,z)=
 \frac{\partial_z^{n-1}J^-(z)}{(n-1)!}I(u,z)
 +I(u,z)\frac{\partial_z^{n-1}J^+(z)}{(n-1)!}.
\]
Here \(J^-\) contains negative modes, while \(J^+\) contains nonnegative
modes. The indicated left and right operations commute for different
\(n\), so this gives a well-defined extension to every polynomial state.
Its current commutator follows from
\(m\binom{m-1}{n-1}=n\binom mn\). Thus the operation retains descendants,
not just the highest coefficient.

Let \(\mathcal L^-\) be the operation line with normalized frame \(k=K\).
The connection on operations is
\(\nabla^{\mathrm{op}}I=dI-T_2I\,dz+IT_1\,dz\), and hence
\[
 \nabla^-k=-\frac{dz}{z}k,
 \qquad \mathcal L^+=(\mathcal L^-)^\vee,
 \qquad \nabla^+k^\vee=\frac{dz}{z}k^\vee.
\]
The fused state recipient is \(F_2\otimes\mathcal L^+\).
To verify its representing property on finite sums of scalar-charge Fock
modules, apply the zero-mode Ward identity first: the output charge must
be two. Removing the creation exponential and shift leaves a current-linear
map of polynomial Fock modules. Such a map is determined by a scalar on
the highest vector. The descendant recursion supplies every first-input
state. The resulting operation space is
\(\mathcal L^-\otimes\operatorname{Hom}_{\mathrm{cur}}(F_2,W)\), which
is the Hom space from the stated dual-line recipient, with its Hom
connection.

The map \(\iota_2\) therefore gives a quasi-isomorphism of flat families
\[
 F_2\otimes\mathcal L^+
       \longrightarrow \mathcal J_2\otimes\mathcal L^+.
\]
The physical Gaussian operation is carried to \(\mathcal V\) by
\(\rho_2\). Counterclockwise half-transport is \(-1\) on the state line
and on its dual operation line. For general scalar charges \(a,b\) at
level \(\kappa\), the same calculation gives opposite connection forms
\(\pm(ab/\kappa)d\log z\). Their transports are dual, and need not be
equal when \(ab/\kappa\) is nonintegral.

There is a specific extension through collision for the recipient.
Extend the normalized frames over \(\mathbb C[z]\), defining
\[
 \overline{\mathcal E}_F=\mathbb C[z]\otimes F_2\,k^\vee,
 \qquad
 \overline{\mathcal E}_{\mathrm{phys}}
       =\mathbb C[z]\otimes\mathcal J_2\,k^\vee.
\]
The connection is logarithmic:
\(\nabla_{z\partial_z}(a(z)m)=(za'(z)+a(z))m\).
The contraction is independent of \(z\), so
\(1\otimes\iota_2\) is a quasi-isomorphism over \(\mathbb C[z]\),
preserves this connection and has residue \(+1\). Its derived fibre at
zero is the quasi-isomorphism \(F_2\to\mathcal J_2\).
The horizontal frame \(z^{-1}k^\vee\) specifies a different lattice with
residue zero. A meromorphic change of frame cannot be evaluated at zero
to identify their chosen fibres.

The operation has a different limit. If \(f\) has oscillator weight \(w\),
then its output of weight \(d\) has the form
\[
 \operatorname{pr}_d\mathcal V(z)f=z^{1-w+d}v_d(f).
\]
An annihilation removes weight \(a\) with factor \(z^{-a}\), a creation
adds weight \(c\) with factor \(z^c\), and \(d=w-a+c\).
Consequently, separately in every output weight,
\[
 \lim_{z\to0}z^{w-1}\mathcal V(z)f
                     =f(-1,-1,\ldots)|2\rangle.
\]
There is no common finite-power normalization for all inputs:
\begin{equation}\label{eq:vol2-native-unbounded-input-poles}
       \langle2|\mathcal V(z)x_1^m=(-1)^mz^{1-m}.
\end{equation}
The first three values are \(z,-1,z^{-1}\).
Removing the negative shift would remove these poles and violate the
\(J_{-1}\) Ward identity. The framed recipient fibre and the
weight-dependent operation limit are therefore distinct constructions.

\section{Regular Wilson jets and corrected cutoff maps}

Return to a nonzero numerical \(\kappa\), keeping \(\hbar=1\).
For the incoming face use the pairing \(b_-=-b\) and the linear inclusion
\(i_- =\zeta^{-1}i\). These changes preserve the boundary pairing:
\(b_-(i_-\otimes i_-)=\mu^s\).
The ray functional then has no phase factor.
For \(\xi=\gamma+dt\,\beta\), define constant derivations
\begin{equation}\label{eq:vol2-native-ray-jets}
 \begin{aligned}
 \ell_n(s\xi)&=\frac1{n!}\int_0^\infty
                  \partial_z^n\beta_\eta^{(0)}(t,0)\,dt,& |\ell_n|&=1,\\
 r_n(s\xi)&=\frac1{n!}\int_0^\infty
                  \partial_z^n\beta_\Phi^{(0)}(t,0)\,dt,& |r_n|&=0,\\
 e_n(s\xi)&=\frac1{n!}\partial_z^n\gamma_\Phi^{(0)}(0,0),& |e_n|&=1.
 \end{aligned}
\end{equation}
The superscript selects Dolbeault degree zero. Each derivation evaluates
one variable of a smooth kernel. No product of point distributions is
introduced.

\begin{lemma}\label{lem:vol2-native-endpoint}
The ray and endpoint derivations satisfy
\begin{equation}\label{eq:vol2-native-endpoint-transgression}
 [d,\ell_n]=[d,e_n]=0,
 \qquad [d,r_n]=\kappa(n+1)\ell_{n+1}-e_n.
\end{equation}
They graded-commute with each other and with \(\Delta_{b_-}\).
\end{lemma}

\begin{proof}
On a linear test field, suspension and normal integration give
\[
 r_n(d(s\xi))=e_n(s\xi)-\kappa(n+1)\ell_{n+1}(s\xi).
\]
The endpoint term uses
\(\int_0^\infty\partial_t\gamma_\Phi^{(0)}dt
=-\gamma_\Phi^{(0)}(0)\). The other term is the off-diagonal part of
\eqref{eq:vol2-native-tangent}, including the derivative factorial.
Since \(r_n\) has degree zero, its commutator with \(d\) is \(-r_nd\)
on a linear input. The relative condition kills the possible endpoint
terms for \(\ell_n d\) and \(e_nd\). Constant contractions graded-commute.
Continuity extends these calculations to smooth joint kernels.
\end{proof}

For a finite set \(F\subset\mathbb Z_{>0}\), introduce the regular
coefficient algebra
\[
 Y_F=\mathbb C[B_0,B_n:n\in F]\otimes\Lambda(c_{n-1}:n\in F),
 \quad Q_YB_n=\kappa n c_{n-1},\quad Q_YB_0=Q_Yc_j=0.
\]
The \(B\)'s have degree zero and the \(c\)'s degree one. Tensors with
\(Y_F\) are algebraic in the coefficients. Operator expressions use
graded left multiplication, while the retained module structure is the
regular right action.
Set
\[
 D_0=d+\Delta_{b_-}+Q_Y+B_0\ell_0,
 \qquad L_F=\sum_{n\in F}\frac{B_n}{\kappa n}r_{n-1}.
\]
The bare descendant source would add
\(\sum_{n\in F}(B_n\ell_n+c_{n-1}r_{n-1})\) to \(D_0\).
Its square is \(\sum_{n\in F}c_{n-1}e_{n-1}\).
For the first jet it takes the value \(c_0\) on a suspended test function
equal to one at the endpoint. This is a nonzero boundary defect.

Define the corrected differential by
\begin{equation}\label{eq:vol2-native-corrected-source}
 D_F=D_0+\sum_{n\in F}
       \left(B_n\ell_n+c_{n-1}r_{n-1}
                         -\frac{B_n}{\kappa n}e_{n-1}\right).
\end{equation}
The correction is forced by the endpoint transgression.

\begin{proposition}\label{prop:vol2-native-source-gauge}
On regular smooth observables with coefficients in \(Y_F\),
\[
                 D_F=e^{-L_F}D_0e^{L_F},\qquad D_F^2=0.
\]
The conjugation is right \(Y_F\)-linear, fixes the unit and is finite
on every input of finite field number.
\end{proposition}

\begin{proof}
The term \(Q_YB_n=\kappa nc_{n-1}\) cancels the interior part of
\([d,c_{n-1}r_{n-1}]\). Its endpoint part is the bare defect above.
Differentiating \(-B_ne_{n-1}/(\kappa n)\) cancels that defect.
More precisely,
\[
 [D_0,L_F]=D_F-D_0,
 \qquad [L_F,[D_0,L_F]]=0.
\]
The second identity follows from constant contractions and commuting
even coefficients. The conjugation formula follows. Each \(r_n\)
lowers field number by one and kills constants, so its exponential
terminates and fixes \(1\otimes Y_F\).
\end{proof}

Let \((i_0,p_0,H_0)\) be the scalar-source quantum collar contraction in
the incoming phase, extended in the remaining coefficients. Then
\[
 I_F=i_0,
 \qquad P_F=p_0e^{L_F},
 \qquad H_F=e^{-L_F}H_0e^{L_F}
\]
contract \((\mathcal O(U)\otimes Y_F,D_F)\) onto
\[
 \mathcal A_F(V)=
 \left(\operatorname{Sym}_b(\Omega_c^{0,\bullet}(V)[1])\otimes Y_F,
         d_1+\Delta_{\mu^s}+B_0\epsilon_0+Q_Y\right).
\]
Indeed \(L_Fi_0=0\): the normal inclusion has no normal
\(e_\Phi\)-component, and its Wick contractions commute with \(L_F\).
Conjugation now transports every strong contraction identity.
The maps preserve the unit, extension of compact tangential supports
and products on tangentially disjoint regions. In particular all
\(B_n,c_j\) continue to act regularly.

Adding a jet changes the differential of a raw smooth observable.
For \(F\subset G\), coefficient inclusion \(\iota_{F,G}\) is therefore
replaced by
\begin{equation}\label{eq:vol2-native-cutoff-transition}
 j_{F,G}=e^{-(L_G-L_F)}\iota_{F,G}.
\end{equation}
The commuting \(L\)'s give
\(D_Gj_{F,G}=j_{F,G}D_F\) and \(j_{G,H}j_{F,G}=j_{F,H}\).
The maps are injective. With \(Y_\infty=\operatorname{colim}_F Y_F\),
their direct limit has the explicit untwisting
\begin{equation}\label{eq:vol2-native-all-jet-source}
 \Psi:\operatorname{colim}_F(\mathcal C_F,j_{F,G})
       \xrightarrow{\sim}(\mathcal O(U)\otimes Y_\infty,D_0),
 \qquad [f]_F\longmapsto e^{L_F}f.
\end{equation}
For its inverse, choose a finite stage containing an input's coefficient
indices and send \(a\) to \([e^{-L_F}a]_F\). Enlarging the stage gives
the same class. This proves the isomorphism and defines the action of
every coefficient. It is continuous for the final topology of finite
coefficient, field-number and compact smooth stages.

The normal contractions commute with these transitions, since
\(P_Gj_{F,G}=P_F\), \(j_{F,G}I_F=I_G\) and
\(H_Gj_{F,G}=j_{F,G}H_F\) after coefficient extension.
They consequently retract the direct limit onto \(\mathcal A_\infty\).
No inverse-limit cohomology assertion is used.

This construction does not apply the literal infinite source to an
unchanged smooth kernel. On a disk of radius less than one choose
\(\chi=1\) near zero and an interior normal profile of integral one.
The linear observable
\[
 u=s\left(dt\,\rho(t)\frac{\chi(z)}{1-z}e_\Phi\right)
\]
has \(r_n(u)=1\) for every \(n\), and \(\ell_n(u)=e_n(u)=0\).
For \(F=\{1,\ldots,K\}\),
\[
 D_Fu=du+\sum_{j=0}^{K-1}c_j,
 \qquad e^{L_F}u=u+\sum_{n=1}^K\frac{B_n}{\kappa n}.
\]
Neither infinite sum belongs to \(Y_\infty\).
The representative of \([u]_\varnothing\) at stage \(F\) is instead
\(u-\sum_{n\in F}B_n/(\kappa n)\).
Quotient cutoff maps also define an inverse-limit carrier, but that
carrier admits coefficient series absent from the direct limit above.

\section{The incoming cap and its state domain}

Geometric rotations act on the incoming test fields by
\[
 R_\theta(s\alpha e_\eta)=s\rho_\theta^*\alpha e_\eta,
 \qquad
 R_\theta(s\alpha e_\Phi)=e^{i\theta}s\rho_\theta^*\alpha e_\Phi,
 \quad \rho_\theta(t,z)=(t,e^{i\theta}z),
\]
and give \(B_n,c_{n-1}\) weight \(n\), with \(B_0\) of weight zero.
Pullback includes the form factors. The extra factor on \(e_\Phi\)
makes the off-diagonal \(\partial_z\) equivariant. Direct evaluation
shows that every source term, \(L_F\), cutoff map and normal contraction
commutes with this circle action.

Define \(\mathcal J_{\mathrm{rot}}(D_{R_0})\) as the algebraic sum of
the circle eigenspaces in the direct-limit source. This is a geometric
condition before taking periods. Combining
\eqref{eq:vol2-native-smooth-period} with \(\Psi\) gives a continuous
right \(Y_\infty\)-linear quasi-isomorphism
\begin{equation}\label{eq:vol2-native-regular-period}
 \Pi_{\mathrm{rot}}:\mathcal J_{\mathrm{rot}}(D_{R_0})
       \longrightarrow
       (\mathbb C[x_1,x_2,\ldots]\otimes Y_\infty,Q_Y).
\end{equation}
For clarity, the period at stage \(F\) is
\(\operatorname{Sym}(\mathcal C)
\exp(\partial_{P_{\mathrm C}}+B_0\partial_\gamma)p_0e^{L_F}\).
All exponentials are finite on each source input.

To prove the assertion about finite circle spectrum, each coefficient
element uses finitely many monomials of \(Y_\infty\).
Separating their weights reduces the target to the finite-weight germ
calculation above. The finite Fourier projector argument then restricts
the full quasi-isomorphism. Separated current insertions give
\[
 J_{-n}=x_n,\qquad J_n=\kappa n\partial_{x_n},\qquad J_0=B_0.
\]
The same Cauchy, covariance and charge contractions prove this on
cochains, before passing to cohomology.

The state produced by an incoming cap is a further object.
For fixed finite \(F\), write the boundary superfields as
\(\Phi_\Sigma=\phi+d\bar z\,\phi_b\) and
\(\eta_\Sigma=c+d\bar z\,a_b\). The boundary jet is
\(x_n=\partial_z^{n-1}\phi(0)/(n-1)!\).
The incoming boundary charge, with its endpoint, is
\begin{equation}\label{eq:vol2-native-incoming-charge}
 \begin{split}
 S^F_{\mathrm{in}}
 &=-\frac1\pi\int c(\partial_{\bar z}\phi+\kappa\partial_z a_b)\,dx\,dy
       +B_0c(0)\\
 &\quad+\sum_{n\in F}
       \left(\frac{B_n\partial_z^nc(0)}{n!}+c_{n-1}x_n\right).
 \end{split}
\end{equation}
The line coefficients are Casimirs for the boundary pairing. Their
differential \(Q_Y\) is retained externally. The incoming form and
potential have the opposite signs to those of the outward face.
In the polarization with \(\eta_\Sigma\) as momentum directions,
\[
             \widehat a_b=\pi\frac\delta{\delta\phi},
             \qquad \widehat c=\pi\frac\delta{\delta\phi_b}.
\]
Derivatives are placed to the right.

The relative master identity can be checked on smeared fields. For a
test function \(g\), the induced derivation has
\[
 \begin{aligned}
 q_F\phi(g)&=-\kappa c(\partial_zg),&q_Fc(g)&=0,\\
 q_Fa_b(g)&=c(\partial_{\bar z}g)
       -\pi\sum_{n\in F}c_{n-1}
                         \frac{\partial_z^{n-1}g(0)}{(n-1)!},\\
 q_F\phi_b(g)&=\phi(\partial_{\bar z}g)+\kappa a_b(\partial_zg)
       +\pi B_0g(0)+\pi\sum_{n\in F}B_n\frac{\partial_z^ng(0)}{n!}.
 \end{aligned}
\]
Together with \(q_F|_{Y_F}=Q_Y\), its square vanishes.
On \(\phi_b(g)\), the two bulk derivatives cancel,
and the endpoint cancellation is
\[
 -\pi\kappa\sum_n c_{n-1}\frac{\partial_z^ng(0)}{(n-1)!}
 +\pi\sum_n\kappa nc_{n-1}\frac{\partial_z^ng(0)}{n!}=0.
\]
Omitting the line-ghost term \(c_{n-1}x_n\) leaves the second sum.

Let \(\mathcal P_{\mathrm{sm}}\) be polynomial functionals of
\(\phi,\phi_b\) with compact smooth kernels. The state domain is
\begin{equation}\label{eq:vol2-native-cap-domain}
 \mathcal H^F_{\mathrm{in}}
   =\operatorname*{colim}_{N\geq\max(F\cup\{1\})}
       \left(\mathcal P_{\mathrm{sm}}\otimes_{\mathbb C}
                       Y_F[[x_1,\ldots,x_N]]\right).
\end{equation}
The series are completed by jet weight, and each weight coefficient is
polynomial in \(Y_F\). The tensor with smooth functionals is algebraic.
Completion precedes the union over the finite set of jet variables.

\begin{proposition}\label{prop:vol2-native-cap-complex}
On the domain \eqref{eq:vol2-native-cap-domain}, the normal-ordered
quantization \(Q_Y+\Omega^F_{\mathrm{in}}\) of
\eqref{eq:vol2-native-incoming-charge} squares to zero.
On the subspace of junction-jet functionals alone it is
\begin{equation}\label{eq:vol2-native-cap-differential}
                    D_G=Q_Y+\sum_{n\in F}x_nc_{n-1}.
\end{equation}
\end{proposition}

\begin{proof}
At a fixed jet stage, differentiation of a series lowers weight by a
bounded amount. Contraction with a smooth kernel takes finitely many
derivatives at the junction, while its remaining contractions restrict
smooth kernels to diagonals. The operator is therefore defined on the
stated domain.

The charge has no \(\phi_b\)-dependent coefficient. Among its terms,
only the \(a_b\)-momentum can contract with a \(\phi\)-coefficient.
Thus a commutator has at most one such contraction, and equals the
quantization of the relative Poisson bracket. The classical calculation
above gives
\(Q_YS^F_{\mathrm{in}}+\tfrac12
\{S^F_{\mathrm{in}},S^F_{\mathrm{in}}\}_{\mathrm{in}}=0\).
There is no higher-contraction correction to its quantization.
On a junction-jet functional, every momentum term has
\(\delta/\delta\phi_b\) on the right and vanishes. The endpoint
line-ghost multiplication remains, giving
\eqref{eq:vol2-native-cap-differential}.
\end{proof}

The larger completion in all jet variables together with smooth
functionals fails this domain test. Choose \(g\) smooth and compactly
supported, equal to \((1-z)^{-1}\) near zero. The sequence
\(u_M=\phi_b(g)\sum_{n=1}^M x_n\) converges in that larger jet-weight
completion. The scalar part of its quantum image is
\[
 \left.\Omega^F_{\mathrm{in}}u_M\right|_{\phi=\phi_b=0}
      =\pi\kappa\sum_{n=1}^M
                      \frac{\partial_z^ng(0)}{(n-1)!}
      =\pi\kappa\frac{M(M+1)}2.
\]
It does not stabilize even at output weight zero. This is why
\eqref{eq:vol2-native-cap-domain} retains a finite jet stage for each
mixed smooth state.

The incoming zero-momentum Lagrangian has constant flat prequantum section
away from the junction. A normalized exponential cap is \(e^f\), where
\(f\) is a polynomial in finitely many \(x_n\), has coefficients in
\(Y_F^0\), and has zero constant term in the jets. Equation
\eqref{eq:vol2-native-cap-differential} says that it is closed exactly when
\[
 f=-\sum_{n\in F}\frac{B_nx_n}{\kappa n}+h(x),
 \qquad h\in\mathbb C[B_0][x_1,x_2,\ldots].
\]
Indeed the coefficient of \(c_{n-1}\) is
\(\kappa n\partial_{B_n}f+x_n\).
Rotation gives \(x_n\) weight \(-n\), while \(B_n\) has weight \(n\).
Every nonconstant monomial of \(h\) has negative weight.
The unique rotation-invariant cap in this declared exponential class is
therefore
\begin{equation}\label{eq:vol2-native-gaussian-cap}
                G_F=\exp\left(-\sum_{n\in F}\frac{B_nx_n}{\kappa n}\right).
\end{equation}
The constant functional \(1\) is not closed when \(F\ne\varnothing\).

Gluing a linear junction functional \(\sum_ny_nx_n\) by the incoming
BV pairing gives \(\sum_ny_ne_{n-1}\) on regular source observables.
For \eqref{eq:vol2-native-gaussian-cap}, this is exactly the endpoint
correction in \eqref{eq:vol2-native-corrected-source}.
The sign is essential: the opposite exponential leaves twice the bare
endpoint defect. Nonlinear junction functionals leave point observables
after contraction and require a larger observable domain.

On the degreewise oscillator-weight product, the all-jet expression
\[
 G=\exp\left(-\sum_{n>0}\frac{x_nB_n}{\kappa n}\right)
\]
is well-defined: each weight coefficient is a finite polynomial in the
line coefficients. Define its polynomial descendant image
\(\mathsf P_\kappa=G(\mathbb C[x_1,x_2,\ldots]\otimes Y_\infty)\).
Then
\[
 D_G(Gf)=GQ_Yf,
 \qquad
 (\kappa n\partial_{x_n}+B_n)(Gf)=G\kappa n\partial_{x_n}f.
\]
Multiplication by \(G\) is a chain isomorphism onto this image, with
inverse \(G^{-1}\). Composing it with
\eqref{eq:vol2-native-regular-period} gives the regular incoming cap map,
sending \(1\) to \(G\) and retaining every line coefficient and
separated current. It identifies the cap image. It does not identify the
whole incoming state completion with a polynomial Fock module.

\section{Formal states and their Laurent fields}

The preceding completions control different variables. Completing the
coefficient ring introduces a further distinction, even for the constant
Heisenberg theory. Fix a numerical \(\kappa\ne0\), specialize the quantum
parameter to one, and introduce an independent formal parameter
\(\varepsilon\). In this section put
\(R_\varepsilon=\mathbb C[[\varepsilon]]\) and
\(R_{\varepsilon,r}=R_\varepsilon/(\varepsilon^r)\).
The parameter \(\varepsilon\) is neither a spacetime coordinate nor the
quantum parameter used to define the current commutator.

Let \(H=\mathbb C[x_1,x_2,\ldots]\) be the polynomial vacuum space,
with \(J=x_1\) and translation \(Tx_n=nx_{n+1}\).
Its current has singular product
\(Y(J,z)J=\kappa z^{-2}+\text{regular terms}\).
Every finite-precision state space
\(H\otimes R_{\varepsilon,r}\) has the usual Laurent fields.
The compatible inverse limit of the state spaces is \(H[[\varepsilon]]\).

\begin{proposition}\label{prop:vol2-native-completed-state-pole}
There is a state in \(H[[\varepsilon]]\) whose field on \(J\) has
no pole bound independent of coefficient precision. Consequently the
finite-precision fields do not take all completed state pairs into
\(H[[\varepsilon]]((z))\).
\end{proposition}

\begin{proof}
The state
\[
 a=\sum_{n\geq0}\frac{(-\varepsilon)^n}{(n+1)!}T^nJ
\]
is well-defined coefficientwise. Translation covariance gives
\(Y(T^nJ,z)=\partial_z^nY(J,z)\). Therefore its singular value is
\begin{equation}\label{eq:vol2-native-completed-pole}
             Y(a,z)J\big|_{\mathrm{sing}}
                =\kappa\sum_{n\geq0}\varepsilon^nz^{-n-2}.
\end{equation}
Modulo \(\varepsilon^r\), the pole has order \(r+1\).
All its coefficients are forced by the compatible reductions. An element
of the displayed uniform Laurent space has only finitely many negative
powers of \(z\), so it cannot have these reductions.
\end{proof}

This state is excluded by a fixed finite-weight restriction. Thus a
finite-weight state core and the total formal state completion answer
different questions. A natural field recipient for the latter is the
restricted Laurent module
\begin{equation}\label{eq:vol2-native-restricted-laurent}
 \mathcal L_\varepsilon(M)
 =\varprojlim_r(M/\varepsilon^rM)((z))
 =\left\{\sum_{j\in\mathbb Z}m_jz^j:
   \forall r\ \exists N_r\quad
   m_j\in\varepsilon^rM\text{ for }j<-N_r\right\},
\end{equation}
where \(M\) is separated and complete. The equality follows by lifting
each compatible coefficient and applying the finite pole bound at each
precision. Formula \eqref{eq:vol2-native-completed-pole} belongs to this
module. Coefficientwise Wick operations extend here because, at every
precision, they are the original finite-precision operations. This is a
formal recipient; it asserts no numerical convergence of the series.

The same distinction occurs in a separated Wilson junction. At
\(\kappa=1\), take rank-one coefficient modules with zero differential
and
\[
 B'_0=0,\quad B'_1=\varepsilon,\quad B'_n=0\ (n>1),
 \qquad B''_0=1,\quad B''_n=0\ (n>0),\qquad c'_j=c''_j=0.
\]
They satisfy the line relation \([d,B_n]=nc_{n-1}\).
The separated scalar twist is \(s(z)=\varepsilon/z\); the two
connections are
\[
                 \nabla_0=d,\qquad
                 \nabla_1=d-\varepsilon z^{-2}dz.
\]
At precision \(r\), the horizontal comparison and its inverse are
\begin{equation}\label{eq:vol2-native-formal-gauge}
 G_r(z)=\sum_{n=0}^{r-1}\frac{(-\varepsilon)^n}{n!z^n},
 \qquad G_r(z)^{-1}=\sum_{n=0}^{r-1}\frac{\varepsilon^n}{n!z^n}.
\end{equation}
Their product is one by the binomial identity, and
\(G'_r-\varepsilon z^{-2}G_r=0\).

Over \(R_\varepsilon((z))\) there is no nonzero horizontal comparison.
If \(g=\sum_{j\geq j_0}g_jz^j\ne0\), with \(g_{j_0}\ne0\), then
the coefficient of \(z^{j_0-2}\) in
\(g'-\varepsilon z^{-2}g\) is
\(-\varepsilon g_{j_0}\ne0\).
A current-linear endomorphism of the polynomial Fock factor is scalar:
its value on the vacuum is killed by every positive mode, hence is a
constant, and negative modes generate all other states. It follows that
\[
 \operatorname{Hom}^{\mathrm{flat,current}}_{R_\varepsilon((z))}=0,
 \qquad
 \varprojlim_r
 \operatorname{Hom}^{\mathrm{flat,current}}_{R_{\varepsilon,r}((z))}
       =R_\varepsilon.
\]
The comparison \(e^{-\varepsilon/z}\) does exist in
\eqref{eq:vol2-native-restricted-laurent}. These are comparisons of
separated flat current families. Extending them through coincidence
requires the supported collision operation as additional data.

There is a useful sufficient bound. If all twist coefficients lie in
an ideal whose \(m\)-th power is zero, and their largest inverse-separation
degree is \(k\), the exponential has pole order at most \(k(m-1)\).
The exponent must be uniform. In \eqref{eq:vol2-native-formal-gauge},
the nilpotence exponent grows with \(r\), and so does the pole order.

\section{Normal support after coefficient completion}

The analogous support issue is visible on one normal coordinate \(u\).
For this section abbreviate \(R=R_\varepsilon\), \(R_r=R_{\varepsilon,r}\),
and use right modules over the Weyl algebra
\[
 D_R=R\langle u,\partial_u\rangle/
                  (\partial_u u-u\partial_u-1).
\]
The parameter \(\varepsilon\) commutes with both operators.
The principal-part module and its density action are
\begin{equation}\label{eq:vol2-native-normal-principal-parts}
 P=R[u,u^{-1}]/R[u]=\bigoplus_{m\geq0}R\beta_m,
 \qquad
 \beta_0u=0,\quad \beta_mu=\beta_{m-1}\ (m>0),\quad
 \beta_m\partial_u=(m+1)\beta_{m+1}.
\end{equation}
Right differentiation is minus ordinary differentiation of the density.
In particular \(P=D_R/uD_R\), with right-module resolution
\[
             0\longrightarrow D_R\xrightarrow{L_u}D_R
                     \longrightarrow P\longrightarrow0.
\]
Here \(L_u\) is left multiplication by \(u\), and is right linear.

The coefficient completion is the restricted mode product
\begin{equation}\label{eq:vol2-native-normal-completion}
 \widehat P=
 \left\{\sum_{m\geq0}a_m\beta_m:
    a_m\in R,\quad
    \forall r\text{ all but finitely many }a_m\in\varepsilon^rR\right\}.
\end{equation}
It is derived complete because the reductions have surjective transitions
and \(P\) is scalar free. Its ordinary normal support has changed.

\begin{proposition}\label{prop:vol2-native-normal-open}
For \(Q=\widehat P/P\), both \(u\) and \(\varepsilon\) act invertibly,
and
\[
 Q\ne0,\qquad \widehat P[u^{-1}]=Q,
 \qquad Q\otimes_R^{\mathbb L}R_r=0.
\]
Thus \(P_r\) has no lift that is both derived complete with those
compatible reductions and ordinarily supported at \(u=0\).
\end{proposition}

\begin{proof}
Lowering the modes is surjective on \(\widehat P\). An element killed
by \(u^N\) has no modes above \(\beta_{N-1}\), so the union of these
kernels is exactly \(P\). Localization is therefore the quotient \(Q\),
on which \(u\) is invertible. The element
\(p=\sum_{m\geq0}\varepsilon^m\beta_m\) has nonzero image there and
satisfies \(pu=\varepsilon p\).

Subtracting a finite representative modulo \(\varepsilon\) and dividing
its remaining coefficients by \(\varepsilon\) proves surjectivity of
\(\varepsilon\) on \(Q\). If \(\varepsilon x\) has finite mode support,
so does \(x\), since \(R\) has no \(\varepsilon\)-torsion. This proves
injectivity. Derived reduction is the cone of multiplication by
\(\varepsilon^r\), and hence vanishes on \(Q\).
Finally, a derived complete lift of the compatible tower is its derived
inverse limit \(\widehat P\), whose localization is nonzero.
\end{proof}

The scalar vacuum \(R\) still exists. The proposition concerns its
normal coefficient, and the perfect test \(D_R\) detects the nonzero
open quotient by \(\operatorname{Hom}_{D_R}(D_R,Q)=Q\).

For a complex \(M\), write
\[
 M_r=M\otimes_R^{\mathbb L}R_r,
 \qquad \Lambda M=R\!\varprojlim_r M_r.
\]
Call \(M\) \emph{derived adically supported} at \(u=0\) when
\begin{equation}\label{eq:vol2-native-derived-support}
                         M_r[u^{-1}]\simeq0
                         \quad\text{for every }r.
\end{equation}
Open restriction is flat localization of the underlying
\(R[u]\)-module with its induced differential action. It does not treat
\(u\) as central in \(D_R\).
Every cohomology module of every derived reduction is tested in
\eqref{eq:vol2-native-derived-support}. For a scalar-flat module this is
equivalent to
\[
 \forall v\ \forall r\ \exists N:
                         vu^N\in\varepsilon^rM.
\]
For a general module the displayed condition controls only ordinary
quotients, and misses the negative cohomology of derived reduction.

Here is an explicit instance of that difference. Let \(M\) consist of
families in
\(\prod_{n\geq1,m\geq0}(R/\varepsilon^n)\beta_m^{(n)}\)
whose coefficients are nonzero modulo \(\varepsilon^r\) at only
finitely many pairs \((n,m)\), for each \(r\). Give each summand the
normal action \eqref{eq:vol2-native-normal-principal-parts}.
Division of coefficients shows that
\[
 M/\varepsilon^rM
       =\bigoplus_{n\geq1}P/\varepsilon^{\min(n,r)}P.
\]
These quotients are supported, and their inverse limit is \(M\).
The negative cohomology of \(M_r\) is
\(T_r=M[\varepsilon^r]\), with transition multiplication by
\(\varepsilon\). Its inverse limit is zero by separatedness. Moreover,
for \(y_r\in T_r\), the convergent series
\(x_r=\sum_{j\geq0}\varepsilon^jy_{r+j}\) belongs to \(T_r\) and
satisfies \(x_r-\varepsilon x_{r+1}=y_r\). Thus
\(\varprojlim^1T_r=0\), proving derived completeness of \(M\).

Nevertheless, in \(N=M[u^{-1}]\) the elements
\begin{equation}\label{eq:vol2-native-torsion-witness}
 w_r=\sum_{n\geq r}\varepsilon^{n-r}\beta_{n-1}^{(n)}
\end{equation}
give nonzero compatible classes in \(N[\varepsilon^r]\).
Indeed \(\varepsilon^rw_r=0\), and
\(w_r-\varepsilon w_{r+1}=\beta_{r-1}^{(r)}\) vanishes after
localization. No power of \(u\) kills \(w_r\): choose an index
\(n\geq r\) larger than that power. Hence the compatible class in
\(\varprojlim_rN[\varepsilon^r]\) is nonzero, and
\(H^{-1}(\Lambda N)\ne0\). In fact the Milnor exact sequence in
degree \(-1\) identifies these two groups, since the reductions have
no degree \(-2\) cohomology. This proves that the complete module
\(M\) fails \eqref{eq:vol2-native-derived-support} despite its
supported ordinary quotients.

\section{Reconstruction on the normal Weyl carrier}

The derived support condition has a positive reconstruction theorem.
It applies to normal coefficients; identifying them with all native
fields or their physical fusion requires additional comparisons.
Let \(\mathcal G_\varepsilon\) be the category of derived complete
scalar complexes, and let \(\mathcal C_\varepsilon\) be the category
of derived complete right \(D_R\)-complexes satisfying
\eqref{eq:vol2-native-derived-support}.

First, completion preserves every derived reduction:
\begin{equation}\label{eq:vol2-native-completion-reduction}
                         (\Lambda X)_r\simeq X_r.
\end{equation}
To see this without a boundedness assumption, tensor the inverse limit
with the perfect complex \([R\xrightarrow{\varepsilon^r}R]\).
For \(n\geq r\), the resulting two-term model of
\(X_r\otimes_R^{\mathbb L}R_n\) has zero differential between its
two copies of \(X_r\). The transitions are multiplication by
\(\varepsilon\) on the negative copy and identity on the other copy.
The first tower has zero homotopy limit: its shift on the product has
\(r\)-th power zero, so \(1-\mathrm{shift}\) has a finite geometric-series
inverse. The second tower is constant. This proves
\eqref{eq:vol2-native-completion-reduction} with its natural map.
The fibre of completion is an inverse limit of multiplication by
\(\varepsilon\), hence has invertible \(\varepsilon\). Maps from such
an object to a complete object vanish, by computing them into its
reductions. Consequently \(\Lambda\) is an exact reflection onto
complete complexes.

Localization commutes with scalar reduction. Therefore
\eqref{eq:vol2-native-derived-support} can also be tested by
\(M_1[u^{-1}]\simeq0\), or by \(\Lambda(M[u^{-1}])\simeq0\).
Indeed the first condition says that \(\varepsilon\) is invertible on
the open restriction, which is precisely the kernel of completion.

\begin{theorem}\label{thm:vol2-native-normal-reconstruction}
The functors
\[
 I(E)=\Lambda(P\otimes_R^{\mathbb L}E),\qquad
 K(M)=R\operatorname{Hom}_{D_R}(P,M)
\]
give inverse equivalences
\(\mathcal G_\varepsilon\simeq\mathcal C_\varepsilon\).
No scalar flatness or cohomological boundedness is required.
\end{theorem}

\begin{proof}
Over a commutative complex algebra \(B\), let \(V\) be a right Weyl
module with locally nilpotent \(u\). Define the finite Taylor projector
\[
 \pi(v)=\sum_{j\geq0}\frac{(-1)^j}{j!}(vu^j)\partial_u^j.
\]
The identity \(\partial_u^ju=u\partial_u^j+j\partial_u^{j-1}\)
gives \(\pi(v)u=0\) by consecutive cancellation.
The binomial identity gives the inverse formula
\[
                    v=\sum_{m\geq0}
                         \frac{\pi(vu^m)\partial_u^m}{m!}.
\]
Thus \(P_B\otimes_B\ker u\to V\), sending
\(\beta_m\otimes e\) to \(e\partial_u^m/m!\), is surjective.
Applying the largest power of \(u\) in a finite relation extracts
its highest coefficient, so the map is injective. It also proves
surjectivity of \(u\) on \(V\).

The resolution of \(P_B\) computes
\[
             K_B(V)=[V\xrightarrow{u}V]_{[0,1]}.
\]
This is a complex of scalar modules; its differential is not claimed
to be right Weyl linear. For a complex with supported cohomology,
the long exact sequence identifies
\(H^qK_B(V)=\ker(u:H^qV\to H^qV)\), since \(u\) is surjective
on that cohomology. Scalar freeness of \(P_B\) now proves that
the evaluation counit is an isomorphism on all cohomology groups.
The unit is an equivalence because lowering on \(P_B\) is split
surjective with kernel \(B\beta_0\).

The finite free resolution makes \(K\) commute with derived scalar
reduction and completion. Tensor--Hom adjunction and the reflection
property give \(I\dashv K\) on complete objects. At precision \(r\),
the unit and counit are the equivalences just proved over \(B=R_r\).
Their limits give the assertion.
\end{proof}

For example, this equivalence supplies an intrinsic supported product.
It need not be the ambient product. The latter fails already for
\(W=\prod_{i\geq0}\widehat P\). It is scalar flat and complete, and
\(W_1=\prod_iP_1\). The element
\(e=(\beta_0,\beta_1,\ldots)\) has a nonzero \(i=N\) coordinate
after multiplication by \(u^N\), for every \(N\). Thus \(W_1\) is
not supported. The intrinsic product is instead
\(I(\prod_iR)\): at each coefficient precision it imposes a mode bound
uniform over \(i\). The example concerns all indices; every finite
product is supported.

There is also an actual geometric operation compatible with the
equivalence. Let \(a:\mathbb A_x^1\times\mathbb A_y^1\to\mathbb A_t^1\)
be addition, \(t=x+y\), and define
\[
                   M*_+N=\Lambda a_+
                             (M\boxtimes_R^{\mathbb L}N),
\]
where \(a_+\) is right differential-module direct image, computed by
the relative de Rham complex. For complete scalar complexes put
\(E\widehat\otimes_R F=\Lambda(E\otimes_R^{\mathbb L}F)\).

\begin{proposition}\label{prop:vol2-native-normal-convolution}
There is a natural equivalence
\[
                    I(E)*_+I(F)\simeq
                    I(E\widehat\otimes_R F).
\]
It is compatible with associativity, symmetry and the convolution unit
\(\widehat P\). Relative convolution over an algebra \(I(B)\)
corresponds to \(I(\Lambda(E\otimes_B^{\mathbb L}F))\), with its
specified right and left module actions.
\end{proposition}

\begin{proof}
At precision \(r\), change variables to \(t=x+y\), \(s=y\).
Since \(\partial_x=\partial_t\) and
\(\partial_y=\partial_t+\partial_s\), the density modes transform by
\[
 \beta_i(x)\beta_j(y)\longmapsto
       \sum_{k=0}^j\binom{i+k}{k}
                 \beta_{i+k}(t)\beta_{j-k}(s).
\]
The relative de Rham differential on \(P_s\), in degrees \(-1,0\),
is \(\partial_s\); it raises the mode by one with coefficient \(m+1\).
It is split injective over \(R_r\), with cokernel \(R_r\beta_0\).
Projection to that cokernel sends the displayed expression to
\(\binom{i+j}{i}\beta_{i+j}(t)\). This remains a contraction after
tensoring with any derived scalar ground.
The three-input coefficients agree by the multinomial identity.
Addition and iterated relative de Rham direct image provide the
coherent comparisons, symmetry and unit.

Scalar reduction commutes with this construction. Equation
\eqref{eq:vol2-native-completion-reduction} then gives the completed
comparison. Finally its two-sided bar has the actual action maps and
terms \(E\widehat\otimes B^{\widehat\otimes n}\widehat\otimes F\).
Its realization is the completed relative tensor in the statement,
because completion is a monoidal localization: its kernel of
\(\varepsilon\)-invertible objects is a tensor ideal.
\end{proof}

This additive normal convolution retains the coordinate \(t\).
The balanced Weyl tensor consumes the normal coordinate instead:
\[
 R\operatorname{Hom}_{D_R}(P,P)=R,
 \qquad P\otimes_{D_R}^{\mathbb L}P^t=R[1].
\]
Here transpose sends \(u\) to \(u\) and \(\partial_u\) to
\(-\partial_u\). The same lowering operator computes both formulas,
in degrees \(0,1\) for Hom and \(-1,0\) for tensor.
Neither operation alone is the physical boundary fusion constructed
from separated insertions and their collision maps.

Nor is scalar flatness preserved by every relative tensor. For
\(B=\widehat{R[t]}^{\,(\varepsilon)}\), take the two rank-one modules
with \(t=0\) and \(t=\varepsilon\). The free resolution for the first
gives their relative tensor as
\([R\xrightarrow{\varepsilon}R]_{[-1,0]}\simeq R/(\varepsilon)\).
At precision \(r\), both its kernel and cokernel are nonzero.
Derived normal support retains both; ordinary scalar reduction loses
one of them.

\section{The interacting comparison}

Return to a nonconstant Hamiltonian operator \(H\).
The classical master identity constructs its local bulk action.
The Gaussian calculation constructs a quantum action on supported
boundary observables, a charged junction and an incoming cap on their
declared domains. Extending all of them to the interacting theory is
a coupled problem.

To formulate its perturbative obstruction, suppose a regulator and
a local subtraction prescription give a BV complex of bulk, boundary
and defect functionals, with the boundary terms included in its
differential. Let \(S_0\) be its classical solution and let
\(s=Q+\{S_0,-\}\). Write a proposed quantum action as
\(S=S_0+\sum_{n\geq1}\hbar^nS_n\).
In a regularized BV model its equation is
\begin{equation}\label{eq:vol2-native-joint-qme}
                 QS+\tfrac12\{S,S\}+\hbar\Delta S=0.
\end{equation}
At order \(n\), the equation for the degree-zero correction is
\[
 sS_n=-\mathcal O_n,
 \qquad
 \mathcal O_n=\Delta S_{n-1}
       +\tfrac12\sum_{i=1}^{n-1}\{S_i,S_{n-i}\}.
\]
The BV identities and the lower-order equations give
\(s\mathcal O_n=0\), by taking the order-\(n\) term of the identity
that the differential twisted by \(S\) annihilates its curvature.
Thus the obstruction is a class in degree one of this joint local
complex. A class in the bulk-only complex does not test endpoint
counterterms or the boundary Ward identities.

The required local prescription must also respect restrictions to
nested regions. Its annulus amplitudes must sew with the chosen cap,
and collision extensions must obey the same current and defect Ward
identities. The forest contractions of
Chapter~\ref{ch:vol2-supported-compositions} provide explicit Gaussian
homotopies. They do not solve the interacting master equation on all
these supports.

An algebraic Hochschild action poses a different comparison. To
identify it with the physical action, one needs a map from the
renormalized bulk observables to the coherent operations on the
boundary category, preserving the unit, translation, collision maps
and their homotopies. If a centre represents those operations, the map
must factor through that representing object. Equivalence further
requires that the chosen boundary objects detect and reconstruct the
bulk. Neither a shared current bracket nor a solution of the master
equation supplies these maps by itself.

The Gaussian normal contraction shows why the operation is the relevant
datum: a field removed by linear projection can still act through its
contraction with a boundary insertion. Formal support and Laurent
completion then specify where that action lives. The general interacting
bulk-to-boundary evaluator, and its identification with a universal
holomorphic centre, remain open beyond the constructions proved here.

\chapter{Marked reconstruction, descent and sewing}
\label{ch:vol2-marked-reconstruction}

A boundary junction can detect some boundary objects and miss others.
Its endomorphism algebra determines the objects generated by the
junction. Recovering a larger category requires a generation theorem;
recovering local fields also requires mapping objects that respect
quasi-isomorphisms and geometric composition. These requirements have
different tests.

The simplest test uses the complex
\[
 P_0=\bigl(k[x]\theta\xrightarrow{\ x\ }k[x]\bigr),
 \qquad |\theta|=-1.
\]
It reconstructs the sector supported at \(x=0\), but it cannot detect
the nonzero module \(k[x,x^{-1}]\). For the Ramond junction, the
failure persists without localization: nonzero polynomial boundary
modules with nonnegative energy have zero derived Hom from the junction.
A local residue gives a second obstruction. Its cone belongs to the
stable category generated by the junction but admits no compatible
action of the original differential graded quartet.

We first establish the positive reconstruction statement and retain
its marking. We then prove these obstructions on their full
coefficient categories. The final sections specify the comparisons
needed for collision descent, holomorphic globalization and sewing
with a prepared state. Throughout, complexes are cohomological,
\(M[1]^n=M^{n+1}\), and \(k\) is a field of characteristic zero.

\section{The sector detected by one object}
\label{sec:vol2-generated-sector}

Put \(A=k[x]\) and regard \(P_0\) as a left \(A\)-module complex.
Its augmentation to \(A/(x)\) is a quasi-isomorphism. The derived
category \(D(A)\) is obtained from complexes of left \(A\)-modules
by inverting quasi-isomorphisms. Define its \(x\)-torsion sector by
\begin{equation}\label{eq:vol2-torsion-sector}
 D_{x\text{-tor}}(A)
 =\{M\in D(A):A[x^{-1}]\otimes_A^{\mathbb L}M\simeq0\}.
\end{equation}
This condition is imposed on the derived localization. It does not
impose \(x\)-adic completion.
An object is compact when its derived Hom functor preserves
coproducts. It generates a stable sector when the smallest
full stable subcategory containing it and closed under
coproducts is that sector.

\begin{proposition}\label{prop:vol2-torsion-generator}
The object \(P_0\) is a compact generator of
\(D_{x\text{-tor}}(A)\). It does not generate \(D(A)\).
\end{proposition}
\begin{proof}
Let \(\operatorname{Loc}(P_0)\) be the smallest full stable
subcategory of \(D(A)\) containing \(P_0\) and closed under
coproducts. Localization preserves triangles and coproducts and
kills \(P_0\), so
\(\operatorname{Loc}(P_0)\subset D_{x\text{-tor}}(A)\).
The localization fibre
\[
 \Gamma_xA=[A\longrightarrow A[x^{-1}]]
\]
has its terms in degrees \(0,1\). It is quasi-isomorphic to
\((A[x^{-1}]/A)[-1]\), and
\[
 A[x^{-1}]/A
 =\varinjlim_n(A/x^n,\ \times x).
\]
Each \(A/x^n\) is a finite extension of copies of \(A/x\).
After tensoring with \(M\), each successive factor is a complex
of \(k\)-vector spaces with the quotient \(A\)-action. Such a
complex belongs to \(\operatorname{Loc}(A/x)\): over a field
it is a direct sum of shifts of its cohomology. Consequently
\(\Gamma_xA\otimes_A^{\mathbb L}M\) belongs to
\(\operatorname{Loc}(P_0)\). The localization triangle identifies
this object with \(M\) whenever \(M\) satisfies
\eqref{eq:vol2-torsion-sector}. This proves equality of the two
subcategories. The two-term finite free model makes \(P_0\)
compact.

On the other hand,
\[
 \operatorname{RHom}_A(P_0,A[x^{-1}])
 =[A[x^{-1}]\xrightarrow{\ -x\ }A[x^{-1}]]
\]
is contracted by multiplication by \(-x^{-1}\). Its nonzero
target is therefore invisible to \(P_0\), which excludes
generation of \(D(A)\).
\end{proof}

The general statement uses the same evaluation map. Let \(A\)
now be a dg algebra over a commutative dg ring \(R\), and let
\(P\in D(A)\). Choose a cofibrant representative and write
\[
 E=\operatorname{REnd}_A(P).
\]
We use right \(E\)-modules. Precomposition makes
\(\operatorname{RHom}_A(P,M)\) a right \(E\)-module; the
evaluation action makes \(P\) a left \(E\)-module commuting,
with the graded sign, with its left \(A\)-action. The adjunction is
\begin{equation}\label{eq:vol2-morita-adjunction}
 G(N)=N\otimes_E^{\mathbb L}P,\qquad
 F(M)=\operatorname{RHom}_A(P,M),\qquad G\dashv F.
\end{equation}
Let \(\operatorname{thick}(P)\) denote closure under finite
extensions, shifts and retracts. Let \(\operatorname{Perf}(E)\)
denote the corresponding closure of the free right module \(E\).

\begin{theorem}[Reconstruction of the generated sector]
\label{thm:vol2-marked-morita}
The unit and counit of \eqref{eq:vol2-morita-adjunction} give
inverse equivalences
\[
 \operatorname{thick}(P)\simeq\operatorname{Perf}(E).
\]
If \(P\) is compact in \(D(A)\), they extend to
\[
 \operatorname{Loc}(P)\simeq D(E^{\mathrm{op}}).
\]
Here \(D(E^{\mathrm{op}})\) means the category of right
\(E\)-modules. The image of \(P\) is the free module \(E\);
all specified arrows between objects of the sector are carried
with their full derived mapping spaces.
\end{theorem}
\begin{proof}
There are canonical identifications \(F(P)=E\) and \(G(E)=P\).
The adjunction unit on \(E\) and counit on \(P\) are identities
under these identifications. Both functors are exact and preserve
retracts. The subcategories on which their unit and counit are
equivalences are therefore closed under finite extensions, shifts
and retracts. They contain the respective generators, proving
the first assertion and full faithfulness.

If \(P\) is compact, \(F\) also preserves coproducts; \(G\)
does so as a left adjoint. The same argument now applies to
localizing subcategories. The free right module \(E\) generates
its whole derived category: every module admits a semifree
resolution formed from coproducts of shifts of \(E\), with
successive free attachments. Thus the unit is an equivalence
on every right \(E\)-module. The image of \(G\) is
\(\operatorname{Loc}(P)\), where the counit is an equivalence.
The equivalence is enriched, since the adjunction was formed
from derived tensor and Hom; it therefore carries arrows and
all their homotopies.
\end{proof}

This is the module form of derived Morita theory
\cite{KellerDG2006}. Its domain is part of the theorem. If
\(N\ne0\) and \(F(N)=0\), the counit at \(N\) is \(0\to N\).
Such an \(N\) is a \emph{right orthogonal} to \(P\).
It excludes reconstruction on any category containing \(N\),
without affecting reconstruction on \(\operatorname{thick}(P)\).

\subsection{The unit and the coefficient action}

An equivalence of module categories identifies their centres
when it preserves the specified coefficient action. For ordinary
dg algebras over \(R\), the centre is the Hochschild cochain
\(E_2\)-algebra
\[
 Z_{\mathrm{Mor},R}(A)=\operatorname{RHom}_{A^e}(A,A),
 \qquad A^e=A\otimes_R^{\mathbb L}A^{\mathrm{op}}.
\]
Its \(E_2\) operations arise from composition and the coherent
interchange of endomorphisms of the identity functor.

\begin{proposition}\label{prop:vol2-morita-centre-transport}
An \(R\)-linear derived Morita equivalence identifies these
centres and their actions on modules. If it is represented by
a \(B\)--\(A\) bimodule \(U\) with inverse \(A\)--\(B\)
bimodule \(V\), let \(K\) be an \(R\)-complex. The image of
a bimodule map \(z:K\otimes_R^{\mathbb L}A\to A\) is
\[
 \begin{aligned}
 K\otimes_R^{\mathbb L} B
 &\simeq U\otimes_A^{\mathbb L}
       (K\otimes_R^{\mathbb L} A)\otimes_A^{\mathbb L}V\\
 &\xrightarrow{\,1\otimes z\otimes1\,}
 U\otimes_A^{\mathbb L}V\simeq B.
 \end{aligned}
\]
The comparison uses both inverse bimodule equivalences and
their triangle homotopies.
\end{proposition}
\begin{proof}
A colimit-preserving \(R\)-linear functor on module categories
is determined by its value on the free module. Applying it to
the free bar resolution identifies it with derived tensor by
that value, including the right algebra action and all natural
transformations. In particular, identity-functor endomorphisms
are diagonal-bimodule endomorphisms. Conjugation by the
equivalence and its inverse preserves composition, the
identity functor and the interchange of transformations.
It therefore preserves the centre's \(E_2\) structure.
The displayed formula is this conjugation on evaluation;
tensoring it with a module proves compatibility with its action.
\end{proof}

For example, take \(B=\operatorname{Mat}_2(k)\), \(U=k^2\)
and \(V=(k^2)^*\). The outer product and evaluation give
\[
 U\otimes_k V\xrightarrow{\sim}B,\qquad
 V\otimes_B U\xrightarrow{\sim}k.
\]
The first map sends the four standard tensors to the matrix
units. In the second tensor product, balancing kills the
off-diagonal tensors and identifies the two diagonal tensors.
The inverse images of the units are
\(\sum_i e_i\otimes e_i^*\) and \([e_1^*\otimes e_1]\).
The triangle composites are
\(\sum_i e_i e_i^*(u)=u\) and
\(\sum_i v(e_i)e_i^*=v\).
There is nevertheless no unital algebra map \(B\to k\):
both square-zero off-diagonal units would map to zero,
and so would their products and their sum \(1_B\).
Morita units do not supply scalar augmentations.

\section{Reconstruction with a marked arrow}
\label{sec:vol2-arrow-reconstruction}

An action on an object and an action preserving a chosen
preparation are different data. Fix a closed arrow
\(\nu:Q\to P\) of \(A\)-bimodules, with cofibrant endpoints.
The marking is the entire arrow, including its domain.
Write \(\operatorname{Hom}\) for the bimodule mapping complex
in this subsection. A concrete model for its derived end algebra is
\begin{equation}\label{eq:vol2-arrow-end-algebra}
 E_\nu=\operatorname{End}(Q)\oplus\operatorname{End}(P)
                     \oplus\operatorname{Hom}(Q,P)[-1].
\end{equation}
For a homogeneous triple of degree \(n\), its third component
has degree \(n-1\). Set
\begin{align}
 d(a,b,h)&=(da,db,\nu a-b\nu-dh),
       \label{eq:vol2-arrow-end-d}\\
 (a,b,h)(a',b',h')&=(aa',bb',ha'+(-1)^n bh'),
       \label{eq:vol2-arrow-end-product}\\
 1&=(1_Q,1_P,0).\nonumber
\end{align}
A closed degree-zero triple records two endomorphisms and the
homotopy making the marked square commute.

\begin{proposition}\label{prop:vol2-arrow-end-model}
Equations \eqref{eq:vol2-arrow-end-d} and
\eqref{eq:vol2-arrow-end-product} define a unital dg algebra
model of the endomorphisms of \(\nu\) in the derived arrow
category.
\end{proposition}
\begin{proof}
Put \(F(a,b)=\nu a-b\nu\). The differential is the homotopy-fibre
differential of \(F\), so it squares to zero. The identity
\[
 F(aa',bb')=F(a,b)a'+bF(a',b')
\]
and the Hom differential's Leibniz rule verify the dg Leibniz
identity with the sign in
\eqref{eq:vol2-arrow-end-product}. The third component of
either triple product, for degrees \(n,m\), is
\[
 ha'a''+(-1)^n bh'a''+(-1)^{n+m}bb'h'',
\]
which proves associativity.

To identify the multiplication in the derived category, replace
\(\nu\) by the cofibrant arrow \(i:Q\to T\), where
\[
 \begin{gathered}
 T=P\oplus Q\oplus Q[1],\qquad i(q)=(0,q,0),\\
 d_T(p,q,s)=(d_Pp-\nu s,d_Qq+s,-d_Qs).
 \end{gathered}
\]
The source \(Q\) and the cokernel \(\operatorname{Cone}(-\nu)\)
are cofibrant, which verifies the assertion about the arrow.
The map \(\pi:T\to P\), \(\pi(p,q,s)=p+\nu q\), is a
quasi-isomorphism with contractible kernel and \(\pi i=\nu\).
For a degree-\(n\) triple, define
\[
 \Phi(a,b,h)=(a,B),\qquad
 B(p,q,s)=(bp+(-1)^n hs,\ aq,\ (-1)^n as).
\]
Substitution into the differentials and composition proves that
\(\Phi\) is a strict dg algebra map to \(\operatorname{End}(i)\).
Postcomposition with \((1,\pi)\) identifies its source complex
with \(\operatorname{Hom}(i,\nu)\): the three free coefficients
are \(b,\nu a,(-1)^n h\). Since \(i\) is cofibrant,
\(\operatorname{End}(i)\to\operatorname{Hom}(i,\nu)\)
is a quasi-isomorphism. Thus \(\Phi\) is a quasi-isomorphism,
as required.
\end{proof}

To prescribe the scalar action on the domain of the marking,
form the homotopy pullback of unital \(E_1\)-algebras
\begin{equation}\label{eq:vol2-fixed-mark-classifier}
 Z_\nu=E_\nu\times^h_{\operatorname{REnd}_{A^e}(Q)}R.
\end{equation}
The map from \(R\) is scalar multiplication; the projection
\(Z_\nu\to R\) is an augmentation.

\begin{theorem}\label{thm:vol2-fixed-mark-classifier}
Let \(B\) be an independently specified augmented unital
\(E_1\)-algebra over \(R\), with augmentation \(\varepsilon\).
Then
\[
 \operatorname{Map}_{\operatorname{Alg}_R/R}(B,Z_\nu)
 \simeq \operatorname{Act}^{E_1}_{\varepsilon}(B;\nu).
\]
The right side is the full derived space of actions on the
arrow \(\nu\), together with an identification of the action
on \(Q\) with \(\varepsilon\,1_Q\). It retains the marked
square's homotopy and every higher identification.
\end{theorem}
\begin{proof}
Maps from a cofibrant resolution of \(B\) to the derived end
algebra of an object classify homotopy coherent associative
actions on that object. Applying this construction in the
arrow category gives \(\operatorname{Map}(B,E_\nu)\).
Mapping into the homotopy pullback
\eqref{eq:vol2-fixed-mark-classifier}, and taking the fibre over
the prescribed augmentation, gives
\[
 \operatorname{Map}(B,E_\nu)
 \times^h_{\operatorname{Map}(B,\operatorname{REnd}(Q))}
 \{\varepsilon\,1_Q\}.
\]
This is exactly the space in the statement. The construction
uses derived mapping spaces, so it also identifies their
homotopies in every degree.
\end{proof}

\subsection{Two actions on the same marked object}

The Hochschild centre
\(\mathcal H_A=\operatorname{REnd}_{A^e}(A)\) has two
natural unital \(E_1\) maps to
\(\operatorname{REnd}_{A^e}(P)\):
\begin{equation}\label{eq:vol2-two-tensor-actions}
 \lambda_P(z)=z\otimes_A^{\mathbb L}1_P,\qquad
 \rho_P(z)=1_P\otimes_A^{\mathbb L}z.
\end{equation}
They use the two unit identifications with \(P\).
Derived tensor is an enriched functor, so it preserves
composition and all homotopies. It acts naturally on a supplied
arrow \(\nu\), giving lifts to \(E_\nu\). A lift through
\(Z_\nu\) additionally requires the prescribed scalar action
on \(Q\). No \(E_2\) structure on the end algebra of an arbitrary
object is implied by these \(E_1\) maps.

\begin{example}\label{ex:vol2-split-mark-two-actions}
Let \(A=k[x]\), \(P=A\oplus A^e\), and let \(\nu:A\to P\)
include the diagonal summand. The class \(x\in HH^0(A)\)
acts identically by \(\lambda_P\) and \(\rho_P\) on that
summand. On \(A^e=k[x_L,x_R]\), their difference sends \(1\)
to \(x_L-x_R\ne0\). Since \(P\) is concentrated in degree
zero, this is a nonzero class in
\(H^0\operatorname{REnd}_{A^e}(P)\). Hence the two algebra
maps are not homotopic, even though the marking is split.
\end{example}

The example also explains why the comparison paths in
Chapter~\ref{ch:vol2-currents} cannot be discarded.
An equivalence preserving its polynomial-root object must
preserve the full Massey coset \(-1+pB\). An equivalence over
the level ring must preserve the arrow to the operator complex,
its nonsplit cohomology extension and its derived special
fibre. Retaining only an ordinary scalar centre loses these
data. A fixed-mode gauge equivalence must also retain its
residual automorphisms; equality of orbit sets is weaker.

\section{The Ramond junction and its boundary categories}
\label{sec:vol2-ramond-reconstruction}

We now test generation on a junction defined by its fields.
Let \(R=\mathbb C\llbracket q\rrbracket\) and
\(S_\ell=R/(q^\ell)\), \(\ell\ge1\).
All states over \(S_\ell\) are algebraic.
Completion means \(V\llbracket q\rrbracket\), with a finite
algebraic state at each coefficient; no product over oscillator
energy or internal shifts is taken.

Fix a Ramond coordinate and the normalized coefficient lines.
The boundary algebras have zero differential and regular
self-products:
\[
 A_0=S_\ell[\partial^j\gamma:j\ge0],\qquad
 A_1=\Lambda_{S_\ell}(\partial^j C:j\ge0).
\]
The physical odd field is \(c=qC\). The normalized action
of \(C\) is specified independently, without inverting \(q\).
The junction state complex is
\begin{equation}\label{eq:vol2-ramond-junction}
 P_\ell=S_\ell[x_i,y_n]\otimes\Lambda(\theta_i,\eta_n),
 \qquad
 D=\sum_{i\ge0}x_i\partial_{\theta_i}
             -\sum_{n\ge1}\eta_n\partial_{y_n}.
\end{equation}
The degrees of \(x,y,\theta,\eta\) are \(0,0,-1,1\).
Exterior derivatives are left derivatives. Its two internal
gradings \(W,E\) are
\[
 (W,E)(x_i)=(W,E)(\theta_i)=(1,2i+1),\qquad
 (W,E)(y_n)=(W,E)(\eta_n)=(-1,2n-1).
\]
Here \(E\) is the specified oscillator grading.

The independent boundary actions are
\[
 X_i=x_i,\quad G_n=-\partial_{y_n},\qquad
 K_i=\partial_{\theta_i},\quad H_n=\eta_n.
\]
The even modes commute, and the odd modes supercommute.
Their generating fields are
\begin{align}
 \Gamma(z)&=\sum_{i\ge0}X_i z^{i-1/2}
                  +\sum_{n\ge1}G_nz^{-n-1/2},\label{eq:vol2-ramond-even}\\
 C(z)&=\sum_{n\ge1}H_nz^{n-1/2}
                  +\sum_{i\ge0}K_i z^{-i-1/2}.\label{eq:vol2-ramond-odd}
\end{align}
The even action is originally a right action. Its commutativity
allows left-module notation here; no odd opposite-algebra
identification is made.

Define \(\mathcal C_{0,\ell}\) to consist of internally
bigraded dg modules with this unital even field action such that,
for every vector \(v\), there are \(N,h\) with
\begin{equation}\label{eq:vol2-positive-mode-support}
 (G_n:n>N)v=0,\qquad (G_1,\ldots,G_N)^h v=0.
\end{equation}
Equivalently, \(J_{N,h}v=0\), where
\(J_{N,h}=(G_n:n>N)+(G_1,\ldots,G_N)^h\).
Define \(\mathcal C_{1,\ell}\) using the odd field action and
the condition \(K_iv=0\) for sufficiently large \(i\), for each
\(v\). A finite exterior family is already jointly nilpotent.
All derivative, product and regular substitution identities
are part of these field actions.

The restricted enriched mapping complex is the direct sum
over internal shifts,
\[
 \underline{\operatorname{Hom}}^{\mathrm{res}}(M,N)
 =\bigoplus_{u,v\in\mathbb Z}
   \operatorname{Hom}_{\mathrm{fields,bigr}}
                    (M,N\langle u,v\rangle).
\]
The shift means that a map raises \(W,E\) by \(u,v\).
A degree-\(r\) odd-boundary map obeys
\(f(am)=(-1)^{r|a|}af(m)\); its differential is
\(d_Nf-(-1)^rfd_M\). Derive these mapping complexes at each
\(S_\ell\), then take \(\mathbf R\varprojlim_\ell\).
Weak equivalences are underlying quasi-isomorphisms.

Translation is additional structure: an even operator \(T_M\)
of internal degree \((0,2)\), commuting with \(d_M\), and
\([T_M,Y_M(a,z)]=\partial_zY_M(a,z)\).
Use a superscript \(T\) when it is retained. On \(P_\ell\),
\begin{align}
 T_P={}&\sum_{i\ge0}(i+\tfrac12)
       (x_{i+1}\partial_{x_i}+\theta_{i+1}\partial_{\theta_i})
       \nonumber\\
 &+\sum_{n\ge1}(n+\tfrac12)
       (y_{n+1}\partial_{y_n}+\eta_{n+1}\partial_{\eta_n})
       -\tfrac12y_1x_0+\tfrac12\eta_1\theta_0.
       \label{eq:vol2-ramond-translation}
\end{align}
All sums are finite on an input. The two final terms are
multiplication operators. In particular,
\begin{equation}\label{eq:vol2-preparation-translation}
 T_P1=D(-y_1\theta_0/2)\ne0.
\end{equation}
Both boundary unit fields still act by the identity.
Write \(\mathcal D^T_{a,\ell}\) for the corresponding
category localized at the underlying quasi-isomorphisms.

\subsection{A positive reverse mode comparison}

There is an exact reconstruction of the reverse mode algebra
at the junction. Its recipient is the derived category of
all bigraded modules over
\(\mathsf B_\ell=\Lambda_{S_\ell}(\kappa_i,\eta_n)\),
where \(\kappa_i=K_i\), with the restricted sum of internal
shifts defined above. Write \(P=\varprojlim_\ell P_\ell\)
and \(\mathsf B=\varprojlim_\ell\mathsf B_\ell\) for their
coefficientwise completions.
This paragraph does not impose the additional odd finite-tail
condition on every object.

Put \(a_i=2i+1\), \(b_n=2n-1\). Define an independent
algebra of coefficientwise continuous operators by
\begin{equation}\label{eq:vol2-restricted-reverse-actor}
 (\mathsf A_{\mathrm{res}})^{m,e}
 =
 \left(
 \prod_{|\beta|\le m}
 \ \bigoplus_{\substack{|\alpha|=m-|\beta|\\
                 a\cdot\alpha-b\cdot\beta=e}}
 \mathbb C\,x^\alpha g^\beta
 \right)\llbracket q\rrbracket,\qquad m\ge0.
\end{equation}
Take the coefficientwise direct sum over \(m,e\); all
multiindices have finite support. For fixed \(\beta\),
the sum over \(\alpha\) is finite because the positive
weights \(a_i\) tend to infinity. On an input, only finitely
many \(\beta\) act through \(g_n=-\partial_{y_n}\).
Thus \eqref{eq:vol2-restricted-reverse-actor} is an actual
operator algebra, with composition adding internal degrees.

\begin{proposition}\label{prop:vol2-reverse-mode-equivalence}
The independent right action induces an equivalence
\[
 \mathsf A_{\mathrm{res}}
 \xrightarrow{\sim}
 \underline{\operatorname{REnd}}_{\mathsf B}^{\mathrm{res}}
                                      (P)^{\mathrm{op}}.
\]
It is a statement in this mode category, with the specified
coefficient completion.
\end{proposition}
\begin{proof}
At coefficient level \(S_\ell\), take the semifree resolution
\[
 Q_{\mathsf B}=\mathsf B_\ell[t_i,s_n],\qquad
 d_Q=\sum_i\kappa_i\partial_{t_i}
                         +\sum_n\eta_n\partial_{s_n}.
\]
The new generators have degree zero and internal degrees
\((-1,-a_i)\), \((-1,b_n)\). Use divided monomials.
Its augmentation to \(S_\ell\) is a quasi-isomorphism by
the positive-degree Euler contraction. Filtration by total
\(t,s\)-degree is semifree.

The restricted dual
\(\underline{\operatorname{Hom}}_{\mathsf B}^{\mathrm{res}}
(Q_{\mathsf B},S_\ell)\) is concentrated in degree zero
with zero differential. A map of shift \((m,e)\) can be
nonzero on \(t^{[\alpha]}s^{[\beta]}\) exactly when
\[
 |\alpha|+|\beta|=m,\qquad a\cdot\alpha-b\cdot\beta=e.
\]
Its independent coefficient values give precisely
\eqref{eq:vol2-restricted-reverse-actor} modulo \(q^\ell\).
The primitive coproduct on divided powers identifies
convolution with multiplication of the dual symbols.

To identify the actual action, define the
\(\mathsf B_\ell\)-linear quasi-isomorphism
\[
 f(t^{[\alpha]}s^{[\beta]})
       =\delta_{\alpha,0}(-y)^{[\beta]}
       \quad:Q_{\mathsf B}\longrightarrow P_\ell.
\]
It preserves the constant cohomology. Define a degree-minus-one
map by
\[
 H_i(t^{[\alpha]}s^{[\beta]})
       =\theta_i\delta_{\alpha,0}(-y)^{[\beta]},
 \qquad H_i(bv)=(-1)^{|b|}bH_i(v).
\]
The Clifford identity gives
\[
 d_{\operatorname{Hom}}H_i=x_if-f\partial_{t_i},
 \qquad g_nf=f\partial_{s_n}.
\]
For \(F=x_{i_1}\cdots x_{i_r}g^\beta\), telescope:
\[
 H_F=\sum_{j=1}^r x_{i_1}\cdots x_{i_{j-1}}
 H_{i_j}\partial_{t_{i_{j+1}}}\cdots
             \partial_{t_{i_r}}\partial_s^\beta,
 \qquad dH_F=F_Pf-fF_Q.
\]
The finite-input bounds of
\eqref{eq:vol2-restricted-reverse-actor} make this homotopy
valid for every admitted coefficient sum. Composing with
the augmentation of \(Q_{\mathsf B}\) identifies the actual
action with the corresponding dual coefficient. It induces
an isomorphism on cohomology and was already an associative
derived-algebra map, proving the equivalence.
The coefficient reductions of these Hom models are
surjective, so the ordinary inverse limit computes the
prescribed derived limit. All maps commute with reduction.
\end{proof}

Even this equivalence does not reconstruct all local fields.
The bidegree \((2,0)\) consists of operators
\(\sum_i\lambda_i x_i g_{i+1}\).
Modes of finite differential-polynomial fields in \(\gamma\),
and finite products of such modes, have coefficient sequences
in
\[
 \mathbb C^{(\mathbb N_0)}+\mathbb C[i(i+1)]
\]
at each coefficient of \(q\). Indeed a quadratic derivative
field contributes
\[
 (-1)^{a+b}
 \bigl[(\tfrac12-i)_a(i+\tfrac32)_b
       +(i+\tfrac32)_a(\tfrac12-i)_b\bigr],
\]
where \((z)_a\) is the rising factorial. These polynomials
are invariant under \(i\mapsto-i-1\); the cases \((a,b)=(2r,0)\)
have degree \(2r\) and leading coefficient \(2\), so they
span \(\mathbb C[i(i+1)]\). Products of two linear modes
supply exactly the finitely supported terms. Higher
polynomial degree cannot contribute at weight two.
The sequence \(\lambda_i=2^i\) is absent: every sufficiently
high finite difference kills a polynomial but preserves
\(2^i\). It therefore gives a nonzero action-cone class,
not killed by any power of \(q\). The mode completion and
the local field algebra retain different information.

\section{Polynomial right orthogonals}
\label{sec:vol2-ramond-orthogonals}

Consider the zero-differential single-boundary modules
\[
 N_{0,\ell}=S_\ell[x_i,y_n],\qquad
 N_{1,\ell}=\Lambda_{S_\ell}(\theta_i,\eta_n).
\]
Their modes and gradings are inherited from \(P_\ell\).
Their translations are the even and odd parts, respectively,
of \eqref{eq:vol2-ramond-translation}, including the
multiplication terms \(-y_1x_0/2\) and \(+\eta_1\theta_0/2\).
They satisfy \eqref{eq:vol2-positive-mode-support} or the
odd finite-tail condition. Each has nonnegative energy and
finite free energy components, because all its state generators
have positive energy. Their constant classes are nonzero.

\begin{theorem}\label{thm:vol2-ramond-nongeneration}
In both boundary categories, with translation either retained
or forgotten,
\[
 \underline{\operatorname{RHom}}^{\mathrm{res}}
          (P_\ell,N_{0,\ell})=0,\qquad
 \underline{\operatorname{RHom}}^{\mathrm{res}}
          (P_\ell,N_{1,\ell})=0.
\]
The same holds after the prescribed derived coefficient
completion. Thus the junction does not generate either full
category, even on the objects with nonnegative energy and
finite free energy components.
\end{theorem}

The proof needs the infinite set of modes. A finite cutoff has
a surviving top Koszul class, so its vanishing cannot serve as
a substitute.

\subsection{Cofree reduction}

Put \(B_0=S_\ell[x_i]\) and \(B_1=\Lambda_{S_\ell}(\eta_n)\).
Let \(C_+\) be the coalgebra with finite-support divided-power
basis \(e_\alpha\) and
\[
 \Delta e_\alpha=\sum_{\beta+\gamma=\alpha}
                         e_\beta\otimes e_\gamma.
\]
Let \(C_\theta=\Lambda_{S_\ell}(\theta_i)\) be the exterior
coalgebra with primitive generators. There are exact
presentations
\begin{equation}\label{eq:vol2-supported-comodule-categories}
 \begin{aligned}
 \mathcal C_{0,\ell}&\simeq
 C_+\text{-}\operatorname{Comod}(B_0\text{-}\operatorname{Mod}_{\mathrm{bigr}}),\\
 \mathcal C_{1,\ell}&\simeq
 C_\theta\text{-}\operatorname{Comod}(B_1\text{-}\operatorname{Mod}_{\mathrm{bigr}}).
 \end{aligned}
\end{equation}
Indeed, the even coaction
\(\rho(v)=\sum_\alpha e_\alpha\otimes G^\alpha v\) is finite
precisely by \eqref{eq:vol2-positive-mode-support}. The odd
coaction uses finite subsets and the exterior signs.
Conversely these modes define the fields of every differential
polynomial. Move annihilators to the right: their indices and
powers are bounded on an input, and a fixed Laurent coefficient
then bounds the sum of the remaining creation indices.
Every coefficient is finite. Zero self-contractions give the
product and substitution identities. Intertwining all modes
is equivalent to intertwining these fields, proving
\eqref{eq:vol2-supported-comodule-categories}.

The forgetful functor \(U\) is exact. Its cofree right adjoint
\(F(V)=C\otimes_{S_\ell}V\) is exact because either coalgebra
is a direct sum of shifts of free coefficient modules.
For a K-injective base resolution \(I\) of \(V\) and an
acyclic comodule \(Z\),
\[
 \operatorname{Hom}_{\mathcal C}(Z,FI)
 =\operatorname{Hom}_{B}(UZ,I)
\]
is acyclic. Hence \(FI\) computes the derived cofree object,
including for unbounded complexes. In every internal shift,
\begin{equation}\label{eq:vol2-cofree-derived-adjunction}
 \operatorname{RHom}_{\mathcal C}(M,FV)
 \simeq\operatorname{RHom}_{B}(UM,V).
\end{equation}
Taking the exact direct sum gives the restricted enrichment.

The targets are \(F(B_0)\) and \(F(B_1)\); on the even side
use \(e_\alpha=(-y)^\alpha/\alpha!\).
The appropriate Koszul or polynomial de Rham factor of
\(UP_\ell\) contracts by the Euler homotopy, leaving \(S_\ell\).
For instance, over \(B_0\) set \(y,\eta\) to zero and retain
\(B_0\otimes\Lambda(\theta_i)\), \(d\theta_i=x_i\);
over \(B_1\) set \(x,\theta\) to zero and retain
\(B_1[y_n]\), \(dy_n=-\eta_n\).
Thus it remains to prove
\begin{equation}\label{eq:vol2-infinite-derived-duals}
 \underline{\operatorname{RHom}}_{B_0}^{\mathrm{res}}(S_\ell,B_0)=0,
 \qquad
 \underline{\operatorname{RHom}}_{B_1}^{\mathrm{res}}(S_\ell,B_1)=0.
\end{equation}

\subsection{A finite head in every cohomological degree}

For \(B_0\), use the semifree resolution
\(B_0\otimes\Lambda(\theta_i)\) with \(d\theta_i=x_i\).
Its positive-degree Euler contraction proves exactness.
Semifreeness gives K-projectivity: against an acyclic target,
a primitive of a closed Hom map is chosen successively on
free generators in each internal shift.

A degree-\(p\) Hom cochain is an alternating family of finite
polynomials indexed by \(p\)-element subsets; the family itself
need not have finite support. With dual exterior symbols
\(\epsilon_i\), its raw differential is
\((-1)^{p+1}\sum_i x_i\epsilon_i\wedge-\).
Conjugation in degree \(p\) by \((-1)^{p(p+1)/2}\) changes it to
\[
 \delta=\sum_i x_i\epsilon_i\wedge-.
\]
Each output coefficient uses only the finitely many indices
in its indexing subset. Define
\[
 Q_jf=\frac{f-f|_{x_j=0}}{x_j},\qquad
 h_j=\iota_{\epsilon_j}Q_j.
\]
Let \(\pi_j\) retain terms containing \(\epsilon_j\) and
set \(x_j=0\) in their coefficients. Polynomial division gives
\begin{equation}\label{eq:vol2-head-contraction}
 \delta h_j+h_j\delta=1-\pi_j.
\end{equation}
On a term without \(\epsilon_j\), use \(Q_j(x_jf)=f\).
On a term with it, the two sides subtract the constant
coefficient in \(x_j\). Cross-index terms anticommute.

For a head of length \(m\), put
\[
 H_m=h_0+\pi_0h_1+\cdots+\pi_0\cdots\pi_{m-2}h_{m-1},
 \qquad \Pi_m=\pi_0\cdots\pi_{m-1}.
\]
The projections commute with \(\delta\), so telescoping gives
\begin{equation}\label{eq:vol2-head-telescope}
 \delta H_m+H_m\delta=1-\Pi_m.
\end{equation}
On a degree-\(p\) cocycle choose \(m=p+1\).
The image of \(\Pi_m\) requires \(m\) exterior symbols and
therefore vanishes. This constructs a primitive of every
cocycle. Each primitive still has finite polynomial output
on each input and preserves the two internal shifts.

For \(B_1\), use
\[
 Q=B_1\otimes S_\ell[s_1,s_2,\ldots],\qquad
 d_Q=\sum_n\eta_n\partial_{s_n},\qquad |s_n|=0,
\]
with divided monomials and
\((W,E)(s_n)=(-1,2n-1)\). The positive-degree Euler contraction
and filtration by total \(s\)-degree prove that this is a
semifree resolution. A degree-\(p\) Hom cochain is a family
\(F_\alpha\in B_1^p\). Its raw differential is
\(-\sum_{\alpha_n>0}\eta_nF_{\alpha-e_n}\):
the Hom sign and graded left-linearity cancel their
\((-1)^p\) factors. Conjugating by \((-1)^p\) gives
\[
 (\delta F)_\alpha
       =\sum_{\alpha_n>0}\eta_nF_{\alpha-e_n}.
\]
Now set
\[
 (h_jF)_\alpha=\iota_{\eta_j}F_{\alpha+e_j},\qquad
 (\pi_jF)_\alpha=
 \begin{cases}
 \eta_j\wedge\iota_{\eta_j}F_\alpha,&\alpha_j=0,\\
 0,&\alpha_j>0.
 \end{cases}
\]
The exterior contraction identity proves
\eqref{eq:vol2-head-contraction}, separately for
\(\alpha_j=0\) and \(\alpha_j>0\).
The same telescoping proof applies to the head \(1,\ldots,m\).
Its top projection forces \(m\) distinct exterior output
generators, hence vanishes for \(m=p+1\). These homotopies
query finitely many input coefficients and preserve both
internal shifts. This proves
\eqref{eq:vol2-infinite-derived-duals}.

\subsection{Translation and completion}

Let \(\delta_T\) be the derivation of boundary modes induced
by translation. Adjoin \(T\) by \(Ta-aT=\delta_T(a)\).
Induction from the category forgetting translation has
underlying module \(S_\ell[T]\otimes M\), and
\[
 a(T^r\otimes v)=
 \sum_{j=0}^r(-1)^j\binom rj
 T^{r-j}\otimes\delta_T^j(a)v.
\]
It remains inside the admitted support class. If
\(J_{N,h}\) kills \(v\), then \(J_{N+r,h+r}\) kills
\(T^r\otimes v\): at most \(r\) annihilators can be
differentiated, leaving at least \(h\) original ones.
Their indices decrease by at most \(r\), and the even creation
modes at the seam commute with them. The odd tail bound
follows from the same finite index change.

For a translated module, the relative sequence
\[
 \begin{gathered}
 0\longrightarrow\operatorname{Ind}(UM)^{(E+2)}
 \xrightarrow{\ \partial\ }\operatorname{Ind}(UM)
 \longrightarrow M\longrightarrow0,\\
 \partial(a\otimes m)=aT\otimes m-a\otimes T_Mm
 \end{gathered}
\]
is exact by division by the monic operator \(T-T_M\).
The first term's generators have energy increased by two.
The covariance identity makes the map balanced; commutation
with the differential makes it a chain map.
Derived adjunction therefore computes translated Hom by a
two-term totalization of underlying derived Hom objects.
Both terms vanish by the preceding calculation. Finally
\(\mathbf R\varprojlim_\ell\) of the zero mapping objects
is zero, while \(N_a\) retains its constant class.
This completes the proof of
Theorem~\ref{thm:vol2-ramond-nongeneration}.

The proof excludes generation in the stated categories.
It does not calculate a dual for an unspecified native tensor.
Regular ambient mode modules and supported boundary modules
also have different admissibility conditions. An assertion
about one of their derived duals cannot be transferred to the
other without a comparison.

\section{A local residue whose cone has no quartet lift}
\label{sec:vol2-quartet-cone}

The polynomial right orthogonals have zero differential.
They cannot carry the original dg quartet, since its relations
would require nonzero boundary actions to be nullhomotopic.
Imposing quartet liftability on each object might therefore
appear to select a reconstructible sector. The next calculation
shows that this condition is not stable under the boundary
cones already generated by \(P_\ell\).

On \(P_\ell\), the original conjugate fields are
\[
 \begin{aligned}
 \widehat\beta(z)&=\sum_{n\ge1}y_nz^{n-1/2}
                     +\sum_{i\ge0}\partial_{x_i}z^{-i-1/2},\\
 b(z)&=\sum_{i\ge0}\theta_i z^{i-1/2}
                     +\sum_{n\ge1}\partial_{\eta_n}z^{-n-1/2}.
 \end{aligned}
\]
A compatible quartet extension on another boundary object means
fields with the same normalized Weyl and Clifford relations,
all original zero mixed supercommutators, and
\[
 [d,b(z)]=\Gamma(z),\qquad [d,\widehat\beta(z)]=-C(z).
\]
In particular its mode \(b_0\) is even-boundary-linear and
satisfies \([d,b_0]=X_0\). Thus the \(X_0\)-action vanishes
in the derived even-boundary category after forgetting
translation. This necessary condition is invariant under
boundary quasi-isomorphism.

Define the degree-two operator
\begin{equation}\label{eq:vol2-ramond-residue}
 \Omega=\sum_{i\ge0}(2i+1)\eta_{i+1}\partial_{\theta_i}.
\end{equation}
It is finite on each input, has internal degree \((-2,0)\),
and commutes with \(D\) and the even boundary modes.

\begin{lemma}\label{lem:vol2-residue-translation}
The operator \(\Omega\) is a translation-invariant local
residue:
\[
 \Omega=-\operatorname{Res}_{z=0}C(z)\partial_z C(z),
 \qquad [T_P,\Omega]=0.
\]
Within the family \(\sum_i a_i\eta_{i+1}\partial_{\theta_i}\),
translation invariance is equivalent to
\(a_i=(2i+1)a_0\).
\end{lemma}
\begin{proof}
A residue contains one creation mode and one annihilation
mode. The pair with index \(i\) contributes
\[
 (i+\tfrac12)\partial_{\theta_i}\eta_{i+1}
 -(i+\tfrac12)\eta_{i+1}\partial_{\theta_i}
 =-(2i+1)\eta_{i+1}\partial_{\theta_i}.
\]
The field \(C\) has zero self-contraction, so this is a mode
of its differential-polynomial local field.
Translation gives
\[
 \begin{gathered}
 [T_P,\partial_{\theta_0}]=\eta_1/2,\qquad
 [T_P,\partial_{\theta_i}]
       =(\tfrac12-i)\partial_{\theta_{i-1}}\quad(i\ge1),\\
 [T_P,\eta_n]=(n+\tfrac12)\eta_{n+1}.
 \end{gathered}
\]
The coefficient of
\(\eta_{i+2}\partial_{\theta_i}\) in the commutator is
\(a_i(i+\tfrac32)-a_{i+1}(i+\tfrac12)\).
These operators are independent, as evaluation on \(\theta_i\)
shows. Their vanishing is the stated recurrence. The remaining
seam is \(a_0\eta_1^2/2=0\).
\end{proof}

Let \(a\in S_\ell\), and adjoin one exterior variable
\(\epsilon\) with degree \(-1\) and internal degree \((2,0)\).
Set
\begin{equation}\label{eq:vol2-residue-cone}
 \mathcal E_{\ell,a}=P_\ell\otimes\Lambda(\epsilon),
 \qquad d_a=D+a\epsilon\Omega,\qquad
 T_{\mathcal E}=T_P\otimes1.
\end{equation}
The even boundary modes are unchanged. Equivalently,
\[
 d_a\theta_i=x_i+a(2i+1)\epsilon\eta_{i+1},\qquad
 d_a y_n=-\eta_n,\qquad
 d_a x_i=d_a\eta_n=d_a\epsilon=0.
\]
The differential squares to zero, commutes with translation,
and preserves the support and internal gradings. The short
exact sequence
\[
 0\longrightarrow\epsilon P_\ell
 \longrightarrow\mathcal E_{\ell,a}
 \longrightarrow P_\ell\longrightarrow0
\]
has connecting morphism \(a\Omega\). Its first term is a
cohomological shift of \(P_\ell\), with weight increased by two.
Thus \(\mathcal E_{\ell,a}\) belongs to the stable closure of
the admitted shifts of \(P_\ell\).
The new variable has energy zero and is exterior, so each
energy component is the direct sum of two finite free
components with different internal shifts. Energy remains
nonnegative.

Its underlying coefficient cohomology is
\(\Lambda_{S_\ell}(\epsilon)\). Indeed, the homogeneous
polynomial change
\[
 x_i'=x_i+a(2i+1)\epsilon\eta_{i+1}
\]
leaves independent pairs \(d\theta_i=x_i'\) and
\(dy_n=-\eta_n\). The Euler homotopy contracts all positive
pair-degree, with \(\epsilon\) a spectator.
This change is not linear for the original boundary action.
To test that action, one must retain a boundary coordinate.

\begin{proposition}\label{prop:vol2-residue-derived-detector}
Forget translation and restrict to the two even modes
\(x=X_0\) and \(g=G_1\). The derived functor
\[
 \mathcal L_g(M)
   =M\otimes^{\mathbb L}_{S_\ell[g]}S_\ell,\qquad g\mapsto0,
\]
retains the \(x\)-action and satisfies
\[
 H^*(\mathcal L_g\mathcal E_{\ell,a})
       \cong e\,\Lambda_{S_\ell}(\epsilon,\eta_1),
 \qquad |e|=-1.
\]
In the ordered basis
\((e,\epsilon e,\eta_1e,\epsilon\eta_1e)\), \(x\) has matrix
\begin{equation}\label{eq:vol2-residue-coordinate-matrix}
 X_a=
 \begin{pmatrix}
 0&0&0&0\\
 0&0&0&0\\
 0&0&0&0\\
 -a&0&0&0
 \end{pmatrix}.
\end{equation}
It is nonzero whenever \(a\ne0\).
\end{proposition}
\begin{proof}
The bounded free resolution
\([S_\ell[g]e\xrightarrow{g}S_\ell[g]]\), with
\((W,E)(e)=(1,-1)\), computes the functor even on
unbounded complexes. On \(\mathcal E_{\ell,a}\),
\(g=-\partial_{y_1}\) is degreewise surjective and commutes
with \(d_a\). Its two-term Koszul tensor is therefore
quasi-isomorphic to its kernel in the indicated shift.
This follows from the short exact sequence of complexes;
polynomial integration need not be a chain map.

The kernel contains precisely the states independent of
\(y_1\). The coordinates \(x_i'\) contract every
\((\theta_i,x_i')\) pair, and the remaining
\((y_n,-\eta_n)\) pairs for \(n\ge2\). What survives is
\(\Lambda(\epsilon,\eta_1)\). Since
\(x_0=x_0'-a\epsilon\eta_1\), its action is exactly
\eqref{eq:vol2-residue-coordinate-matrix}.

For an explicit last contraction, put
\[
 B=\Lambda_{S_\ell}(\epsilon,\eta),\qquad
 K_a=B[x]\otimes\Lambda(\theta),\qquad
 d\theta=x+a\epsilon\eta.
\]
The multiplier \(x+a\epsilon\eta\) is monic, hence injective
on \(B[x]\), including over the nonreduced ring \(S_\ell\).
Polynomial division gives
\[
 Q_a(u)=\frac{u-u|_{x=-a\epsilon\eta}}{x+a\epsilon\eta},
 \qquad h(u)=\theta Q_a(u),\qquad h(\theta u)=0.
\]
With \(\rho\) setting \(\theta=0,x=-a\epsilon\eta\) and
\(\iota\) including exterior representatives,
\(dh+hd=1-\iota\rho\). Thus the four surviving classes are
free over \(S_\ell\), and a nonzero entry \(-a\) cannot
vanish by coefficient torsion.
\end{proof}

\begin{theorem}\label{thm:vol2-no-stable-quartet-sector}
No full replete stable subcategory of
\(\mathcal D^T_{0,\ell}\), containing \(P_\ell\) and
closed under its admitted internal shifts, can consist
entirely of objects admitting the original dg quartet
extension, even up to boundary quasi-isomorphism.
The assertion survives coefficient completion and any further
support restrictions on the extra quartet operators.
\end{theorem}
\begin{proof}
Fullness admits the translated local morphism \(\Omega\)
between shifts of \(P_\ell\); stability admits
\(\mathcal E_{\ell,1}\). Any quartet lift would make its
\(x_0\)-action zero by the boundary-linear homotopy \(b_0\).
Restriction to two modes and derived tensor send a zero
morphism to zero. Proposition~\ref{prop:vol2-residue-derived-detector}
contradicts this. No support condition on \(b_0\) was used.
The coefficient \(-1\) survives already modulo \(q\), so
reduction proves the completed assertion.
\end{proof}

More precisely, \(\mathcal E_{\ell,a}\) is quartet-liftable
if and only if \(a=0\). At \(a=0\) the original operators act
on the first tensor factor. Every nonzero \(a\) is excluded
by the same matrix. The underlying vector complex still has
the primitive
\[
 d_a(\theta_0-a\epsilon y_1)=x_0,
 \qquad
 [g_1,\theta_0-a\epsilon y_1]=a\epsilon.
\]
The second equation is the failure of boundary-linearity.

The integral physical residue is
\[
 -\operatorname{Res}c(z)\partial_zc(z)=q^2\Omega.
\]
Its cone is obstructed for \(\ell\ge3\); for \(\ell=1,2\)
it splits. As a map of complex vector spaces,
\(\operatorname{rank}X_{q^2}=\max(\ell-2,0)\).
Thus the complete lattice is obstructed already modulo
\(q^3\), without cancelling or inverting \(q\).

\subsection{The extension class and its local contact}

The two-mode class has a direct Ext interpretation. Over
\(A=S_\ell[x,g]\), let
\[
 P^{(2)}=S_\ell[x,y]\otimes\Lambda(\theta,\eta),
 \qquad D\theta=x,\quad Dy=-\eta,\quad g=-\partial_y.
\]
The finite free resolution of its constant cohomology is
\[
 Q=A\otimes\Lambda(u,v),\qquad
 |u|=|v|=-1,\quad du=x,\quad dv=g.
\]
The \(A\)-linear map \(f:Q\to P^{(2)}\) given by
\(f(1)=1,f(u)=\theta,f(v)=f(uv)=0\) is a quasi-isomorphism.
For \(\omega=\eta\partial_\theta\), define the degree-one
map \(H\) by \(H(u)=-y\) and zero on the other generators.
Then
\[
 (d_{\operatorname{Hom}}H)(u)=\eta,\qquad
 (d_{\operatorname{Hom}}H)(uv)=-1.
\]
Consequently \(\omega f-d_{\operatorname{Hom}}H\) sends
\(uv\) to \(1\) and the other basis elements to zero.
It is the image under \(f\) of the degree-two endomorphism
of \(Q\) with that same formula. Compose this endomorphism
with the augmentation \(Q\to S_\ell\). The resulting
Hom complex has zero differential, and the value \(1\)
on \(uv\) is its top generator. Thus the class generates
\(\operatorname{Ext}^2_A(S_\ell,S_\ell)=S_\ell\).
The discarded all-mode residue tail has the explicit
two-mode-linear homotopy
\[
 \left[D,-\sum_{i\ge1}(2i+1)y_{i+1}\partial_{\theta_i}\right]
   =\sum_{i\ge1}(2i+1)\eta_{i+1}\partial_{\theta_i}.
\]
Every sum is finite on an input. Thus the Ext class is the
restriction of the actual all-mode residue.

There is a local correction of the field differential, but
it changes a mixed contact. Put
\[
 b_a(z)=b(z)-2a\epsilon\partial_z\widehat\beta(z).
\]
Oscillator commutators give
\[
 [\Omega,b(z)]=2\partial_zC(z),\qquad
 [\Omega,\widehat\beta(z)]=0.
\]
The tensor sign gives
\([d_a,\epsilon\partial_z\widehat\beta]=\epsilon\partial_zC\).
Hence
\[
 [d_a,b_a]=\Gamma,\qquad
 [d_a,\widehat\beta]=-C,\qquad [T_{\mathcal E},b_a]=\partial_zb_a.
\]
However, if
\(\Delta(z,u)=\sum_{r\in\mathbb Z}z^{r-1/2}u^{-r-1/2}\),
then
\begin{equation}\label{eq:vol2-quartet-derivative-contact}
 [\Gamma(u),b_a(z)]=2a\epsilon\,\partial_z\Delta(z,u).
\end{equation}
For example its modes are
\[
 [g_n,b_{a,-i}]=a(2i+1)\epsilon\delta_{n,i+1},\qquad
 [x_i,b_{a,n}]=-a(2n-1)\epsilon\delta_{i,n-1}.
\]
This is a local derivative contact. It is not the original
zero mixed bracket, and
Theorem~\ref{thm:vol2-no-stable-quartet-sector} excludes
every correction restoring all original relations on this
derived boundary object.

\section{Local fields must respect derived equivalence}
\label{sec:vol2-polynomial-fields-derived}

The mapping object is as consequential as the object class.
Local polynomial endomorphism fields at a fixed junction do
not automatically extend to a derived bifunctor on all pairs.
The obstruction is visible on the acyclic localization
\[
 Q_\ell=P_\ell[x_0^{-1}].
\]
The variables retain their gradings, with
\((W,E)(x_0^{-1})=(-1,-1)\).
Laurent differentiation extends all quartet operators and
translation. The boundary support conditions still hold:
inverting \(x_0\) changes neither finite \(y\)-degree nor
finite \(\theta\)-support. The multiplication operator
\(h=m_{\theta_0/x_0}\) satisfies
\[
 Dh+hD=1
\]
as an operator anticommutator. Thus \(Q_\ell\to0\) is a
quasi-isomorphism preserving all specified structures.
The contracting homotopy need not preserve them.

Fix the following polynomial field prescription. A normal
symbol uses the original creation and annihilation variables,
with polynomial creation coefficients and a uniform bound on
total symbol degree at each coefficient of \(q\). It can
contain infinitely many mode monomials, provided its value
on each input is a field with finite lower Laurent order and
honest state coefficients. Inverse creation variables are
not among its generators. Impose quartet locality,
translation covariance and exact linearity for the selected
boundary. For an endomorphism, locality means the first
equation below for each quartet field \(f\) and some finite
order \(K\); covariance means the second:
\[
 (u-z)^K[f(u),F(z)]=0,\qquad [T,F(z)]=\partial_zF(z),
\]
together with supercommutation with every boundary mode.
For \(F:P_\ell\to Q_\ell((z^{1/2}))\), use the corresponding
rectangular equations and differential
\(DF-(-1)^{|F|}FD\).
Call these complexes \(\mathcal E_{a,\mathrm{poly}}^T(M)\)
and \(\mathcal H_{a,\mathrm{poly}}^T(P_\ell,Q_\ell)\).

\begin{proposition}\label{prop:vol2-localization-polynomial-fields}
Extension and restriction identify
\[
 \mathcal E_{a,\mathrm{poly}}^T(Q_\ell)
 \cong\mathcal E_{a,\mathrm{poly}}^T(P_\ell)
 \cong\mathcal H_{a,\mathrm{poly}}^T(P_\ell,Q_\ell)
 \qquad(a=0,1).
\]
The identity of \(Q_\ell\) and inclusion
\(\iota:P_\ell\to Q_\ell\) represent nonzero degree-zero
classes in their respective complexes.
\end{proposition}
\begin{proof}
A normal expression has bounded order in
\(\partial_{x_0}\). On a fixed Laurent input, only finitely
many other annihilator rows can act, by the input's finite
support. Each such row has a lower Laurent bound in the field
variable and finite polynomial coefficients. The expression
therefore extends to \(Q_\ell\).

Successive evaluation on polynomial monomials recovers each
normal row: subtract the proper subrows already recovered,
then divide by the monomial factorial. This proves
faithfulness over \(S_\ell\). Every polynomial expression on
\(Q_\ell\) preserves \(P_\ell\); every rectangular polynomial
expression lands in \(P_\ell\). The field equations are
identities of these bounded-degree polynomial operators,
coefficient by coefficient, so they extend and restrict
faithfully.

Let \(\varepsilon:P_\ell\to S_\ell\) extract the scalar
coefficient. Then \(D1=0\) and \(\varepsilon D=0\).
If a degree-minus-one polynomial field had \(dF=\iota\),
its value on \(1\) would lie in \(P_\ell((z^{1/2}))\).
Applying \(\varepsilon\) to \(D(F(z)1)=1\) would give
\(0=1\). The diagonal identity is excluded in the same
way by restriction.
\end{proof}

\begin{theorem}\label{thm:vol2-polynomial-fields-not-derived}
No reduced bifunctor invariant under underlying
quasi-isomorphisms can be quasi-isomorphic to the unchanged
polynomial field complex on \((P_\ell,Q_\ell)\), or on
every diagonal.
Here reduced means that its value is zero if either input
is zero. The obstruction holds with translation and after
coefficient completion.
\end{theorem}
\begin{proof}
Derived invariance and \(Q_\ell\simeq0\) force both
cross-object and diagonal values involving \(Q_\ell\) to
vanish. Proposition~\ref{prop:vol2-localization-polynomial-fields}
contradicts the proposed comparisons. A completed primitive
would reduce to an impossible primitive modulo \(q\).
\end{proof}

This does not exclude a derived replacement agreeing with
the polynomial actor only at \(P\). It also does not apply
to a recipient imposing nonnegative energy on every object:
\(Q_\ell\) contains \(x_0^{-m}\), and its energy-zero part
contains the independent states
\(x_i x_0^{-(2i+1)}\). The polynomial right orthogonals of
Theorem~\ref{thm:vol2-ramond-nongeneration} require neither
of these features.

\subsection{An inverse field and a remaining odd obstruction}

An enlarged local field class supplies an even-boundary
primitive explicitly. Write
\[
 a(z)=\sum_{i\ge0}x_i z^i,\qquad
 p(z)=\sum_{n\ge1}\partial_{y_n}z^{-n},\qquad
 \Gamma(z)=z^{-1/2}(a(z)-p(z)).
\]
Since \(x_0\) is invertible, set
\begin{equation}\label{eq:vol2-local-inverse-field}
 I(z)=z^{1/2}\sum_{r\ge0}a(z)^{-r-1}p(z)^r,\qquad
 H(z)=b(z)I(z).
\end{equation}

\begin{proposition}\label{prop:vol2-local-inverse-homotopy}
These are fields with finite Laurent-polynomial state
coefficients on every input. They satisfy
\[
 \Gamma I=I\Gamma=1,\qquad [D,H]=1.
\]
The field \(H\) is translation-covariant, local with the
quartet and exactly even-boundary-linear.
There is no degree-minus-one strictly odd-boundary-linear
finite-output field \(P_\ell\to Q_\ell((z^{1/2}))\) with
differential \(\iota\), nor such an endomorphism field of
\(Q_\ell\) with differential \(1\).
\end{proposition}
\begin{proof}
On a vector of \(y\)-degree \(d\), only \(r\le d\) contributes.
At a fixed coefficient of \(z\),
\[
 a(z)^{-s}=x_0^{-s}
       \left(1+\sum_{i\ge1}(x_i/x_0)z^i\right)^{-s}
\]
has finitely many partition terms. The geometric-series
telescoping is therefore finite on every input and proves
the inverse identities. The factors \(b,I\) have zero
mutual contractions; applying inverse-field coefficients
does not increase \(y\)-degree or \(\eta\)-support.
Their product is a field, with \([D,I]=0\), and
\([D,H]=[D,b]I=\Gamma I=1\).

Put \(J=(a-p)^{-1}\). Modewise differentiation gives
\[
 [\partial_{x_i},J]=-z^iJ^2,\qquad
 [y_n,J]=-z^{-n}J^2.
\]
With
\(\Delta(u,z)=\sum_{r\in\mathbb Z}u^{r-1/2}z^{-r-1/2}\),
the nonzero commutators are
\[
 \begin{aligned}
 [\widehat\beta(u),H(z)]&=-\Delta(u,z)b(z)I(z)^2,\\
 [C(u),H(z)]&=\Delta(u,z)I(z).
 \end{aligned}
\]
Each mode coefficient is finite on an input, and
\((u-z)\Delta(u,z)=0\) proves locality.
Exact commutation with \(\Gamma\) proves even-boundary
linearity. Differentiating \(\Gamma I=1\), and commuting it
with \(T\), gives
\([T,I]=-I(\partial_z\Gamma)I=\partial_zI\), then the
covariance of \(H\).

For the odd obstruction, every \(K_i\) kills \(1\).
Strict graded linearity of a degree-minus-one field gives
\(K_iF(z)1=0\) for every \(i\). The common kernel of the
\(K_i=\partial_{\theta_i}\) is
\[
 S_\ell[x_i,x_0^{-1},y_n]\otimes\Lambda(\eta_n),
\]
which has no negative cohomological degree. Thus \(F(z)1=0\),
and its differential cannot send \(1\) to \(1\).
No locality or translation assumption is needed for this
last argument.
\end{proof}

The enlargement in \eqref{eq:vol2-local-inverse-field} is
substantial. On the constant state,
\[
 \begin{gathered}
 H(z)1=\frac{\sum_i\theta_i z^i}{\sum_i x_i z^i},\\
 H_0(1)=\theta_0/x_0,\quad
 H_1(1)=\theta_1/x_0-\theta_0x_1/x_0^2.
 \end{gathered}
\]
Already the \(\theta_0,x_0,x_1\) terms have pole order
\(m+1\) at coefficient \(z^m\). They are honest finite
states coefficientwise, with no uniform pole or symbol-degree
bound. Enlarging symbols can repair the even primitive;
it cannot repair strict odd-boundary linearity.

\section{What full quartet linearity does reconstruct}
\label{sec:vol2-full-quartet-category}

There is a positive reconstruction with a different mapping
category. Let \(\mathscr W_\ell\) be the algebra of finite
polynomial oscillator operators
\[
 x_i,y_n,\theta_i,\eta_n,\quad
 \xi_i=\partial_{x_i},\quad\zeta_n=\partial_{y_n},
 \quad\kappa_i=\partial_{\theta_i},\quad\pi_n=\partial_{\eta_n},
\]
with the normalized Weyl and Clifford relations. All other
supercommutators vanish. Its differential \(\delta=[D,-]\)
satisfies
\[
 \delta\theta_i=x_i,\qquad \delta y_n=-\eta_n,\qquad
 \delta\xi_i=-\kappa_i,\qquad\delta\pi_n=-\zeta_n,
\]
and kills the other generators.

\begin{proposition}\label{prop:vol2-full-quartet-morita}
The unit inclusion \(S_\ell\to\mathscr W_\ell\) and
evaluation \(\mathscr W_\ell\to P_\ell\), \(w\mapsto w1\),
are quasi-isomorphisms. For full quartet-linear morphisms,
\[
 D(\mathscr W_\ell)\simeq D(S_\ell),\qquad
 \operatorname{REnd}_{\mathscr W_\ell}(P_\ell)\simeq S_\ell.
\]
In the bigraded category, the internal shifts of \(P_\ell\)
form the compact generating family.
\end{proposition}
\begin{proof}
Normal ordering gives a free superpolynomial symbol basis.
The differential preserves this order and has exactly the
four displayed pair differentials. In positive total symbol
degree \(N\), the homotopy is
\[
 h=N^{-1}\left(
 \sum_i\theta_i\partial_{x_i}-\sum_n y_n\partial_{\eta_n}
 -\sum_i\xi_i\partial_{\kappa_i}-\sum_n\pi_n\partial_{\zeta_n}
 \right).
\]
Its anticommutator with \(\delta\) is the identity in positive
degree. It is finite on each symbol and preserves internal
degrees. Only the identity class survives; evaluation on \(1\)
identifies it with the constant cohomology of \(P_\ell\).

Tensoring the coefficient contraction with any complex \(V\)
shows that \(V\to\mathscr W_\ell\otimes_{S_\ell}V\) is a
quasi-isomorphism. For a dg \(\mathscr W_\ell\)-module \(M\),
the multiplication map
\(\mathscr W_\ell\otimes_{S_\ell}M\to M\) is a module map
whose composite with the coefficient unit is the identity.
It too is a quasi-isomorphism. These are the derived unit
and counit for extension and restriction along
\(S_\ell\to\mathscr W_\ell\).
The assertions about the end and compact generators follow.
\end{proof}

The scalar symbol projection used in the proof is not an
algebra augmentation. Indeed,
\(\xi_i x_i=x_i\xi_i+1\) has scalar part \(1\), while
\(\xi_i\) and \(x_i\) each have scalar part zero.
The multiplicative quasi-isomorphism is the unit inclusion.

One may also impose polynomial joint support: on every
vector all but finitely many annihilators vanish, and the
remaining even annihilators are jointly locally nilpotent.
For one even pair \([a,z]=1\), the finite-on-input operator
\[
 \Pi_{z,a}=\sum_{r\ge0}\frac{(-1)^r}{r!}z^ra^r
\]
projects to \(\ker a\); for an odd pair \(az+za=1\),
the projector is \(1-za\). The projectors of distinct
pairs commute. Their product \(\Pi_{\mathrm{vac}}\) is
finite on an input and projects to the common annihilator
kernel \(\operatorname{Vac}(M)\).

Taylor expansion gives
\[
 \begin{gathered}
 m=\sum_{\alpha,I}\frac{z^\alpha z_I}{\alpha!}
         \Pi_{\mathrm{vac}}(a_Ia^\alpha m),\\
 P_{\ell,\mathrm{gr}}\otimes_{S_\ell}\operatorname{Vac}(M)
       \xrightarrow{\sim}M.
 \end{gathered}
\]
Odd annihilators are ordered opposite to the creation monomial.
Applying projected derivatives proves uniqueness. Projecting
a lift of a vacuum vector proves exactness of
\(\operatorname{Vac}\). The differential preserves its kernel
and becomes \(D_P\otimes1+1\otimes d_{\mathrm{vac}}\).
This gives the same derived coefficient category on the
jointly supported modules. Fixing oscillator translation
makes the comparison translation-compatible. An arbitrary
extra covariant translation is instead an extra degree-zero
chain operator on the vacuum complex and must be retained.

This positive theorem uses full quartet-linear maps.
Forgetting to the boundary category is not full:
\([\Omega,\theta_0]=\eta_1\ne0\), and the scalar end has
no degree-two class mapping to the boundary residue.
Nor does primitive diagonal tensoring give a normalized
quartet tensor:
\[
 [\xi\otimes1+1\otimes\xi,\ x\otimes1+1\otimes x]
       =2\,1.
\]
The construction therefore does not identify a native
relative sewing product.

\section{Endomorphisms under coefficient and collision maps}
\label{sec:vol2-marked-base-change}

The endomorphisms of a marked object need not commute with
restriction of its coefficients. The comparison is induced
by evaluation and exists before it is known to be invertible.
Let \(R\to S\) be a map of commutative dg rings, \(B\) a
dg \(R\)-algebra, and \(P,Q\) left \(B\)-modules. Write
\(B_S=S\otimes_R^{\mathbb L}B\) and similarly for modules.
Extending evaluation and taking its adjoint gives
\begin{equation}\label{eq:vol2-marked-base-change}
 S\otimes_R^{\mathbb L}\operatorname{RHom}_B(P,Q)
 \longrightarrow
 \operatorname{RHom}_{B_S}(P_S,Q_S).
\end{equation}

\begin{proposition}\label{prop:vol2-perfect-marked-base-change}
If \(P\) is perfect over \(B\), the comparison
\eqref{eq:vol2-marked-base-change} is an equivalence for
every \(Q\). For \(Q=P\), it is an equivalence of unital
\(E_1\)-algebras. For two perfect modules over
\(R\)-algebras \(B_1,B_2\), the external-product comparison
\[
 \operatorname{REnd}_{B_1}(P_1)
       \otimes_R^{\mathbb L}\operatorname{REnd}_{B_2}(P_2)
 \longrightarrow
 \operatorname{REnd}_{B_1\otimes_R^{\mathbb L}B_2}
                (P_1\otimes_R^{\mathbb L}P_2)
\]
is also an equivalence. Both assertions retain a finite
diagram of marking arrows when all its source objects
are perfect.
\end{proposition}
\begin{proof}
For \(P=B\), \eqref{eq:vol2-marked-base-change} is the
identity on \(S\otimes_R^{\mathbb L}Q\). Both sides, as
functors of \(P\), take finite extensions to finite limits
and preserve retracts. Scalar extension preserves finite
limits in a stable category. Induction on finite cells
proves the assertion for every perfect \(P\).
For the external product, replace the target copies of
\(P_1,P_2\) by arbitrary \(Q_1,Q_2\). The comparison is the
identity for the two free source modules; apply the same
finite-cell argument in each source variable.
The maps are adjoint to evaluation, so they preserve
composition, units and successive base changes. Cross Hom
has the same proof, and the arrow end model uses a finite
homotopy fibre. This proves the marked assertion.
\end{proof}

Perfectness cannot simply be omitted. Let \(R=k[t]\),
\(S=R[t^{-1}]\) and \(P=\bigoplus_{n\ge0}Re_n\).
The endomorphism \(e_n\mapsto t^{-n}e_n\) of \(P_S\)
has no common denominator, so it is absent from
\(S\otimes_R\operatorname{End}_R(P)\).
There is a separate tensor defect. On
\(V=\bigoplus_{i\ge0}ke_i\), the endomorphism
\[
 T(e_i\otimes e_j)=\delta_{ij}e_i\otimes e_j
\]
does not belong to \(\operatorname{End}(V)\otimes
\operatorname{End}(V)\). A sum of \(r\) elementary tensors
would make every finite diagonal coefficient matrix have
rank at most \(r\), contradicting its \((r+1)\)-square
identity corner.

\subsection{The collision comparisons}

Let \(X\) be a smooth curve and \(C\) a connected
coaugmented factorizing coalgebra on its ordered collision
carrier. Write \(\mathscr R_I\) for the carrier with a finite
nonempty set of roots \(I\). Suppose
a marked bicomodule category \(\mathcal M_I\) and its
regular object \(\mathbb V_I\) have been specified, with
a closed action of the normalized coefficient category.
On a smooth chart \(Y\) of dimension \(d\), and for a
collision map \(f:Y'\to Y\), use
\[
 L_Y(M)=\omega_Y\otimes_{\mathcal O_Y}M[d],\qquad
 f^\dagger=L_{Y'}^{-1}f^!L_Y.
\]
This fixes the left-to-right normalization; in particular
\(f^\dagger\mathcal O_Y=\mathcal O_{Y'}\).
Define the internal end
\[
 E_I=\underline{\operatorname{RHom}}_{\mathcal M_I}
                   (\mathbb V_I,\mathbb V_I).
\]
Its enrichment, rather than a global scalar Hom, determines
the coefficient object.

If collision restriction takes the regular object to the
regular object and is linear over normalized pullback,
evaluation induces
\begin{equation}\label{eq:vol2-root-collision-end-map}
 c_\alpha:\delta_\alpha^\dagger E_I\longrightarrow E_{I'}
 \qquad(\alpha:I\twoheadrightarrow I').
\end{equation}
For a partition \(\mathcal P=(I_1|\cdots|I_r)\), write
\(j_{\mathcal P}:U_{\mathcal P}\to\mathscr R_I\) for the
separated open. Let \(E_{I,\mathcal P}^{\mathrm{loc}}\)
be the end of the regular object in the restricted marked
category. Restriction and external-product evaluation give
the span
\begin{equation}\label{eq:vol2-separated-end-span}
 j_{\mathcal P}^\dagger E_I
 \xrightarrow{\ b_{\mathcal P}\ }E_{I,\mathcal P}^{\mathrm{loc}}
 \xleftarrow{\ e_{\mathcal P}\ }
 \left.\boxtimes_a E_{I_a}\right|_{U_{\mathcal P}}.
\end{equation}
Neither arrow is inverted by its construction.

\begin{proposition}\label{prop:vol2-perfect-collision-criterion}
Suppose the entire marked category on each root chart is
presented as \(\operatorname{Mod}_{B_I}\), carrying
\(\mathbb V_I\) to a perfect \(P_I\). Suppose each collision
and separated-open restriction is the corresponding derived
extension of scalars, and the separated categories and regular
objects are the external tensor products of their blocks.
If these identifications preserve coefficient actions,
chart transitions, promotions of labels and every ordered
collision diagram, then \(c_\alpha,b_{\mathcal P},e_{\mathcal P}\)
are coherent equivalences preserving composition and evaluation.
\end{proposition}
\begin{proof}
Apply Proposition~\ref{prop:vol2-perfect-marked-base-change}
on each chart. The maps obtained there are the maps adjoint
to the same evaluations as
\eqref{eq:vol2-root-collision-end-map} and
\eqref{eq:vol2-separated-end-span}. The supplied coherent
presentations therefore identify them on overlaps and on
composite restrictions. Equivalences of underlying
coefficient complexes detect equivalences of the normalized
differential modules. They descend with their algebra and
evaluation structures.
\end{proof}

Perfectness at each fixed number of ordinary labels is
weaker than the premise. It neither identifies the restriction
functor of the entire marked category with scalar extension
nor commutes an infinite homotopy end with pullback.

Let \(\mathcal C^f\) be the category of factorizing coalgebras
on this carrier, with transverse tensor \(\otimes^!\)
and local unit \(J\). A transverse monoid is a unital
associative algebra object for this tensor.
The space \(\operatorname{Act}(F;C)\) consists of its
unital associative actions on \(C\), with all their
comparison homotopies.

\begin{theorem}[Conditional factorizing recognition]
\label{thm:vol2-factorizing-recognition}
Suppose the comparisons
\eqref{eq:vol2-root-collision-end-map} and
\eqref{eq:vol2-separated-end-span} are equivalences with all
mixed collision and factorization coherences. Suppose their
evaluations descend, preserve the local unit, and commute
with promotion of ordinary labels to roots. Then
\[
 Z_I=L_{\mathscr R_I}(E_I),\qquad Z_\varnothing=k,
\]
defines a factorizing transverse monoid with an action on \(C\).
It represents the full transverse action functor if, in
addition, evaluation induces equivalences
\begin{equation}\label{eq:vol2-native-action-recognition}
 \operatorname{Map}_{\operatorname{Alg}(\mathcal C^f)}(F,Z)
       \xrightarrow{\sim}\operatorname{Act}(F;C)
\end{equation}
for every actor \(F\) in the specified category
\(\mathcal C^f\).
\end{theorem}
\begin{proof}
The maps \(b_{\mathcal P}^{-1}e_{\mathcal P}\) give the
separated-product equivalences, and the \(c_\alpha\) give
collision equivalences. Their stipulated coherences are the
descent diagrams. Composition and the unit descend because
the comparisons preserve them; the evaluation hypothesis
supplies the action. Composing a map \(F\to Z\) with this
action constructs \eqref{eq:vol2-native-action-recognition}
on all simplices, naturally in \(F\). Its additional
equivalence assertion is exactly representability of the
full action functor.
\end{proof}

The additional assertion is substantive. The full pointed
and ordinary multiplicative action categories considered in
Chapter~\ref{chap:vol2-recipients} have explicit
nonfinality obstructions. A proposed restriction must explain
which actions or maps it excludes and why that restriction is
geometrically appropriate. The marked associative classifier
of Theorem~\ref{thm:vol2-fixed-mark-classifier} remains valid
in its arrow category. Identifying that category with the
supported factorizing action category is a separate map.
The two-object question of a pair of actions with no common
recipient also retains its own quantifier; an unequalizable
parallel pair does not by itself supply that pair of objects.

\section{Complete filtrations and formal comparison}
\label{sec:vol2-complete-comparison}

A comparison of finite coefficients becomes a comparison of
formal objects only when the limit and its operations are
controlled. Let \(C\) be a complex with descending
subcomplex filtration
\(C=F^0C\supset F^1C\supset\cdots\).
Call it complete and separated here when the natural map
\[
 C\longrightarrow\varprojlim_n C/F^nC
\]
is an isomorphism degree by degree. The differential is
continuous for this filtration. Set
\(\operatorname{gr}^nC=F^nC/F^{n+1}C\).

\begin{lemma}\label{lem:vol2-complete-filtered-comparison}
Let \(f:C\to D\) preserve complete separated filtrations
of this form. If every associated-graded map
\(\operatorname{gr}^nf\) is a quasi-isomorphism, then \(f\)
is a quasi-isomorphism. The conclusion is an equivalence
of structured algebras if \(f\) was already a continuous
map of those structured algebras.
\end{lemma}
\begin{proof}
Give the cone of \(f\) the componentwise filtration. Its
associated graded is acyclic. For a cone cocycle \(z\),
acyclicity of \(\operatorname{gr}^0\) gives \(b_0\) with
\(z-db_0\in F^1\). This error is closed. Inductively,
acyclicity of \(\operatorname{gr}^n\) gives
\(b_n\in F^n\) such that
\[
 z-d(b_0+\cdots+b_n)\in F^{n+1}.
\]
Completeness defines \(b=\sum_{n\ge0}b_n\), and continuity
with separatedness gives \(db=z\). Thus the cone is
acyclic. No new operation is supplied by this argument;
it proves equivalence for the continuous structured map
assumed in the statement.
\end{proof}

The existence of a compatible map is a different problem.
For example, let \(\mathfrak g\) be a complete dg Lie
algebra with
\([F^i\mathfrak g,F^j\mathfrak g]\subset F^{i+j}\mathfrak g\).
A deformation is a Maurer--Cartan element
\(\alpha\in F^1\mathfrak g^1\):
\[
 d\alpha+\tfrac12[\alpha,\alpha]=0.
\]
If this equation holds modulo \(F^n\), \(n\ge1\), its error
\(r_n\in F^n\mathfrak g^2\) is closed modulo \(F^{n+1}\).
Indeed the Bianchi identity
\(dr_n+[\alpha,r_n]=0\) makes the second term one
filtration step higher. A correction \(u_n\in F^n\mathfrak g^1\)
changes the error by \(du_n\) modulo \(F^{n+1}\).
Thus the next obstruction lies in
\(H^2(\operatorname{gr}^n\mathfrak g)\). When it vanishes,
choices differ by degree-one classes; their gauge
identifications involve degree-zero classes.
Completeness sums actual compatible corrections. It does
not force these obstruction or automorphism groups to vanish.

This distinction preserves the Hamiltonian comparisons of
the preceding chapters. A complete cochain map can be a
quasi-isomorphism while its proposed brace refinement remains
unconstructed. A selected higher coefficient can change under
transfer. A full Massey coset, a nonzero action-cone class or
a residual gauge stabilizer requires its own invariant
calculation.

\section{Holomorphic descent with a specified comparison}
\label{sec:vol2-marked-holomorphic-descent}

The Hamiltonian Lie-cochain model and the symplectic
polyvector model have different carriers. To formulate the
second, let \((S,\omega)\) be a holomorphic symplectic
surface and \(R=\mathbb C\llbracket\hbarq\rrbracket\).
Write \(\pi=\omega^{-1}\). A chosen formal Poisson tensor
\(\pi_\hbar=\hbarq\pi+O(\hbarq^2)\), with
\([\pi_\hbar,\pi_\hbar]=0\), defines
\[
 \mathcal P_S=
 \bigl(\operatorname{PV}_S\llbracket\hbarq\rrbracket,
                  [\pi_\hbar,-]\bigr).
\]
Polyvector degree is its cohomological degree; the Schouten
bracket has degree \(-1\), so this differential has degree one.
Here \(\operatorname{PV}_S^p=\bigwedge^pT_S\).
The choice \(\pi_\hbar=\hbarq\pi\) is a particular formal
class, not the only quantization datum.

A Weyl algebroid \(\mathcal W\) is specified by local
formal Weyl algebras, equivalences on overlaps and their
coherent cocycle data. A Fedosov presentation uses formal
fibre coordinates and a flat resolution differential, locally
written
\[
 D_F=\nabla-\delta+\hbarq^{-1}[\gamma,-].
\]
Here \(\nabla\) is the chosen local symplectic connection,
\(\gamma\) is a Weyl-valued one-form, and \(\delta\) is the
fibre Koszul differential. The
commutator is divisible by \(\hbarq\), so the displayed
operator need not invert the coefficient parameter.
Flatness is an equation on \(\gamma\) and the prescribed
curvature data. In the holomorphic setting these resolution
and descent data must be supplied, for example through
compatible Dolbeault models; they are not obtained by
assuming a global holomorphic symplectic connection exists.

Dolgushev's smooth formality construction uses Fedosov
resolutions and gives an \(L_\infty\) comparison of
polyvectors with polydifferential operators
\cite{DolgushevCov}*{Theorems 2--4}. That result alone
does not give a holomorphic homotopy-\(E_2\) comparison
with every marking and cyclic structure. We state precisely
the data used below.

Fix a cover, all its multiple intersections, and continuous
homotopy-\(E_2\) resolutions of \(\mathcal P_S\) and
\(\mathcal C_{\mathcal W}\), the relative Hochschild
cochain sheaf of \(\mathcal W\). Suppose a homotopy coherent
map of their descent diagrams
\begin{equation}\label{eq:vol2-fedosov-comparison-data}
 \Phi_\bullet:\mathcal P_\bullet\longrightarrow
                            \mathcal C_{\mathcal W,\bullet}
\end{equation}
has been constructed. Its leading maps are required to be
quasi-isomorphisms for complete separated positive
filtrations as in
Lemma~\ref{lem:vol2-complete-filtered-comparison}.
The coefficient completion, chosen formality refinement
and formal Poisson class are fixed throughout.

\begin{theorem}[Descent of a supplied formality comparison]
\label{thm:vol2-fedosov-marked-descent}
Under these hypotheses, \(\Phi_\bullet\) descends to an
equivalence of the complete homotopy-\(E_2\) sheaves and
induces
\[
 \mathbf R\Gamma(S,\mathcal P_S)
       \simeq \mathbf R\Gamma(S,\mathcal C_{\mathcal W}).
\]
If the chosen continuous \(R\)-linear Morita theory
identifies the latter with the centre of
\(\operatorname{Perf}(\mathcal W)\), this is the corresponding
global centre comparison. Coherent markings, cyclic maps
and state maps descend whenever
\(\Phi_\bullet\) is supplied in their diagram category.
\end{theorem}
\begin{proof}
The filtered lemma makes the map an equivalence on every
member of the descent diagram. A homotopy limit sends an
objectwise equivalence of diagrams to an equivalence.
Taking the descent totalization therefore gives the sheaf
equivalence, and derived global sections give the displayed
map. Limits of algebra objects carry their operations
objectwise, so the homotopy-\(E_2\) structure is preserved.
The Morita assertion uses its stated centre model.
The same limit argument in a diagram category retains
markings and all their comparison homotopies.
\end{proof}

A change of Fedosov data gives the same conclusion when
a compatible gauge equivalence of the chosen resolutions
and its descent homotopies are supplied. An arbitrary
change of the Fedosov one-form is not such a hypothesis:
it can change the formal curvature class.
Likewise, equivariance requires a coherent equivariant
choice of \eqref{eq:vol2-fedosov-comparison-data}.
These qualifications apply on a K3 surface as well.
Its holomorphic symplectic form supplies \(\pi\), not all
of the displayed comparison data.

There is also a formal algebraic realization in two
topological directions. For each holomorphic open, the
specified \(E_2\)-algebra has its factorization realization
on framed two-disks. The framed structure may be promoted
to an oriented one when the required rotation-equivariant
structure is supplied. If these assignments satisfy the
required cosheaf descent in the holomorphic variable, they
define a mixed factorization theory on
\(\mathbb R^2_{\mathrm{top}}\times S\).
Theorem~\ref{thm:vol2-fedosov-marked-descent} then realizes
the supplied centre equivalence in those topological
directions. Neither this realization nor Morita invariance
alone identifies a universal physical boundary category.

For the Hamiltonian BF action
\[
 S_{\mathrm{BF}}(A,B)=
 \int_{\mathbb R^2\times S}
 \left\langle B,(d+\bar\partial)A+\tfrac12[A,A]\right\rangle,
\]
the first comparison is a map from its Hamiltonian Chevalley
observables to the chosen symplectic polyvector complex.
It must preserve the operations and coefficient completion.
Gauge fixing and a coupled quantum master equation are
additional requirements for a quantum factorization theory.
When \(S\) is a complex surface, this spacetime has six real
dimensions. The three-dimensional half-space theory of
Chapter~\ref{ch:vol2-native-bulk-boundary} has a different
geometric input.

\section{Sewing with a prepared state}
\label{sec:vol2-marked-state-sewing}

A reconstruction can preserve an algebra and lose its state.
Let \(B\) be a dg \(R\)-algebra, \(M\) a \(B\)-module,
\(\nu:R\to M\) a chosen closed preparation in \(R\)-complexes,
and \(s:M\to R\) a closed \(R\)-linear state map with
\(s\nu=1\). An equivalence of such data consists of the
algebra and module comparisons, a homotopy carrying the
preparation, a homotopy carrying \(s\), and compatibility
of these homotopies with \(s\nu=1\).
Neither the unit of \(B\) nor an augmentation of \(B\)
specifies \(s\).

To express sewing, fix boundary algebras \(B_0,B_1,B_2\),
a \(B_0\)--\(B_1\) bimodule \(M_{01}\), and a
\(B_1\)--\(B_2\) bimodule \(M_{12}\). The sewn object
\(M_{02}\) is a \(B_0\)--\(B_2\) bimodule.
A sewing map has domain
\[
 \mu_{012}:
 M_{01}\widehat\otimes^{\mathbb L}_{B_1}M_{12}
 \longrightarrow M_{02},
\]
where the completed balanced tensor has been independently
defined on the allowed objects. Preparations must obey
\(\mu_{012}(\nu_{01}\otimes\nu_{12})\simeq\nu_{02}\).
If \(\sigma_{012}\) is the supplied contraction functional
on that balanced tensor, state compatibility is the
specified homotopy
\begin{equation}\label{eq:vol2-marked-sewing-state}
 s_{02}\mu_{012}\simeq\sigma_{012}.
\end{equation}
For three regions these homotopies must agree with the
associativity comparison; higher regions require their
higher coherence. A tensor of scalar functionals descends
to a balanced tensor only when it satisfies the actual
balancing relation. This must be checked in the state model.

\begin{proposition}\label{prop:vol2-sewing-transport}
Suppose two such systems are related by equivalences of
their full diagrams of boundary algebras, bimodules,
markings, state maps and sewing maps, including every
coherence. Then all finite sewn states agree under those
equivalences. If the systems are diagrams of derived
complete objects and their comparisons commute with the
coefficient quotients, the same is true after the prescribed
derived completion.
\end{proposition}
\begin{proof}
For two regions the assertion is
\eqref{eq:vol2-marked-sewing-state} transported through
the given square. Inductively compose this square with
the next sewing map. The associativity homotopies identify
the different parenthesizations, and their higher
coherences identify the resulting paths of identifications.
Thus the full finite diagram, including its final state
map, is carried by the equivalence. At the complete level,
apply the derived inverse limit to the compatible
coefficient diagrams. It preserves their equivalences
and homotopies.
\end{proof}

This proposition requires states on the completed domain in
question. Finite Gaussian integration constructs the
normalization and balanced contractions in
Chapter~\ref{ch:vol2-gaussian-elimination}; the determinant
cannot be removed from that state comparison.
The infinite-mode and stabilized states have the specific
analytic domains of
Chapter~\ref{ch:vol2-infinite-mode-states}.
A formal coefficient equivalence by itself supplies
neither positivity nor a Laplace prescription, and
coefficientwise existence does not establish continuity
on a larger state completion.

\begin{problem}[Local mapping and physical reconstruction]
\label{prob:vol2-marked-physical-reconstruction}
Construct a stable boundary category with homotopy-invariant
cross-object local mapping objects, independently defined
relative sewing and evaluation against both boundary
algebras. The construction must retain the coefficient
lines, local units, preparations and regular, localized,
supported and vacuum collision terms. Determine its
comparison with the marked associative action classifier
and with independently defined bulk observables.

For the Ramond junction, this requires deciding how the
local residue \(q^2\Omega\) and its cone occur in that
category, including the derivative contact
\eqref{eq:vol2-quartet-derivative-contact}.
A category containing the acyclic localization \(Q_\ell\)
cannot have unchanged polynomial local actors on every
diagonal and still be derived invariant. Within the
specified even-boundary derived category, no full stable
sector containing the junction and its shifts can require
every object to lift to the original quartet.
An admissible alternative must construct its mapping and
sewing operations and prove its comparison, rather than
infer them from the end algebra of one object.
\end{problem}

\chapter{Quantum symmetries and geometric comparisons}
\label{ch:vol2-quantum-interfaces}

An algebra of boundary operators determines how those operators compose.
It does not determine how two representations combine.
The latter operation uses a coproduct, or a more general tensor
construction on the category of representations.
Nor does an exchange operator determine the coefficient lattice on
which multiplication acts. These distinctions persist in explicit
polynomial examples.

We begin with one algebra carrying two different Hopf structures.
We then compare ordinary bars, polynomial shuffle lattices and
matrix-coefficient reconstruction. The final examples concern
geometric and gravitational comparisons.
A tangent complex does not determine its nonlinear equations.
The spectrum of a connection does not identify its holonomy with
its coefficients.

Throughout, $k$ is a field of characteristic zero.
All tensors are algebraic until a different completion or localization
is specified.

\section{One operator algebra and two representation tensors}

Put $A=k[x,y,z]$, with augmentation
$\varepsilon(x)=\varepsilon(y)=\varepsilon(z)=0$.
For $\eta=0,1$, define algebra maps
$\Delta_\eta:A\to A\otimes_k A$ by
\begin{equation}\label{eq:vol2-two-hopf}
\begin{aligned}
 \Delta_\eta x&=x\otimes1+1\otimes x,\\
 \Delta_\eta y&=y\otimes1+1\otimes y,\\
 \Delta_\eta z&=z\otimes1+1\otimes z+\eta x\otimes y.
\end{aligned}
\end{equation}
A Hopf structure requires coassociativity, the two counit identities,
and an antipode $S$ with
$m(S\otimes1)\Delta=m(1\otimes S)\Delta=\mathbf1\varepsilon$.
Here $m$ is multiplication and $\mathbf1:k\to A$ is the unit.

\begin{proposition}\label{prop:vol2-two-hopf}
Both coproducts in \eqref{eq:vol2-two-hopf} define Hopf algebras.
Their antipodes are
\[
 S_\eta x=-x,\qquad S_\eta y=-y,\qquad
 S_\eta z=-z+\eta xy.
\]
They have the same augmented algebra and the same ordinary bar
coalgebra, but they are not isomorphic as Hopf algebras.
\end{proposition}
\begin{proof}
It suffices to check the generators.
The only nonprimitive generator is $z$.
Both iterated coproducts give its three single-factor terms and
\[
 \eta\bigl(x\otimes y\otimes1
       +x\otimes1\otimes y+1\otimes x\otimes y\bigr).
\]
This proves coassociativity.

The augmentation kills each extra tensor factor, giving the counits.
On $z$, either antipode composite is
$(-z+\eta xy)+z-\eta xy=0$.
The checks on $x,y$ are immediate, and the identities extend
multiplicatively.

The coproduct $\Delta_0$ is cocommutative, whereas
\[
 \Delta_1z-\Delta_1^{\mathrm{op}}z=x\otimes y-y\otimes x\ne0.
\]
Cocommutativity is preserved by a Hopf isomorphism, so none exists.
The ordinary augmented bar uses multiplication, augmentation and
unit, with word deconcatenation as its coalgebra structure.
Those data are identical for the two algebras.
\end{proof}

The difference is visible on one-dimensional modules.
Let $M$ have character $(x,y,z)=(1,0,0)$ and let $N$ have character
$(0,1,0)$.
On $M\otimes_k N$, the action induced by $\Delta_0$ has $z=0$.
The action induced by $\Delta_1$ has $z=1$.
Thus even the ordinary tensor of representations is extra information
relative to the augmented algebra and its bar.

Assigning weights $1,1,2$ to $x,y,z$ makes both coproducts graded.
Their common Hilbert series
\[
 \frac{1}{(1-t)^2(1-t^2)}
\]
does not distinguish them either.

\section{Ordinary bars and geometric coefficients}

Let $V$ be finite-dimensional and concentrated in degree zero.
Give its generators weight one.
The tensor algebra and the square-zero algebra are
\[
 T(V)=\bigoplus_{r\ge0}V^{\otimes r},
 \qquad
 A_{\mathrm{sq}}=k\oplus V,\qquad V^2=0.
\]
Both are ordinary augmented algebras over $k$.

We use a cohomological bar convention: a word of bar length $n$
has degree $-n$, and $sV=V[1]$ has degree $-1$.
For an augmentation ideal $\bar A$, the bar has terms
$(s\bar A)^{\otimes n}$ and differential given by adjacent products.
Its coproduct cuts a word at every position, including both empty cuts.

\begin{proposition}\label{prop:vol2-free-bar}
There is a quasi-isomorphism of coaugmented coalgebras
\[
 k\oplus sV\longrightarrow B(T(V)),
\]
where $sV$ is primitive.
Consequently the derived augmentation end algebra is
\[
 \operatorname{RHom}_{T(V)}(k,k)
 \simeq k\oplus V^*[-1]
\]
with square-zero reduced multiplication.
For the square-zero algebra,
\[
 B(A_{\mathrm{sq}})=T^c(sV),
\]
with zero differential. Its continuous weight dual is
$\widehat T(V^*[-1])=\prod_{n\ge0}(V^*[-1])^{\otimes n}$,
with concatenation product.
\end{proposition}
\begin{proof}
Choose a basis of $V$.
A bar element of $T(V)$ is a nonempty word in basis elements,
together with a set of cuts separating its nonempty algebra factors.
Fix an underlying word of length $r$.
Its possible cuts form subsets
$C\subset\{1,\ldots,r-1\}$, and its bar length is $|C|+1$.

Choose the bar sign for which
\[
 dC=\sum_{j=1}^{|C|}(-1)^{j-1}
       \bigl(C\setminus\{c_j\}\bigr),
 \qquad c_1<\cdots<c_{|C|}.
\]
For $r\ge2$, define $hC=\{1\}\cup C$ when $1\notin C$, and
$hC=0$ otherwise.

If $1\notin C$, the term deleting the new first cut is $C$.
The other terms cancel those of $hdC$.
If $1\in C$, only the term of $dC$ deleting that cut contributes
to $hdC$, again giving $C$.
Hence $dh+hd=1$ on this word.
Only words of length one and the empty word survive.
Their inclusion respects deconcatenation and proves the first assertion.

The ordinary two-sided free bar resolution computes the derived Hom.
Dualizing the displayed coalgebra quasi-isomorphism gives the
square-zero multiplication, since the surviving reduced classes are
primitive. Dualization is exact over $k$.
For $A_{\mathrm{sq}}$, every product of two augmentation-ideal
elements is zero, so the multiplication differential vanishes.
Each weight is finite-dimensional. Dualizing the algebraic weight sum
therefore gives the indicated product and the dual concatenation law.
\end{proof}

These formulas require an ordinary scalar augmentation.
An empty word, a local vacuum and a prepared state have different
definitions. None can replace that augmentation without a map.
Likewise, the coefficient dual
$(\bigoplus_n W_n)^*=\prod_n W_n^*$ changes a sum to a product.
It does not by itself identify a field dual, an internal line dual
or the boundary-junction construction of
Chapter~\ref{ch:vol2-marked-reconstruction}.

\subsection{Componentwise tensor algebras and factorization}

For a factorization coefficient $F$, restriction to two distinct
points identifies $F_{\{1,2\}}$ with
$F_{\{1\}}\otimes F_{\{2\}}$.
Suppose these three coefficients are $k$.
The componentwise sum $F\oplus F$ has coefficient $k^2$ at
both one and two points, whereas factorization would require $k^4$
at two points.

The same difficulty applies to the naive componentwise tensor algebra.
For the coefficient $F$ just described, it has $k[t]$ at every
nonempty configuration. The inherited separated map is
\[
 k[t]\longrightarrow k[t_1,t_2],
 \qquad t^n\longmapsto t_1^nt_2^n.
\]
It misses $t_1$ and is not a factorization equivalence.
Weight completion changes the map to
$k[[t]]\to k[[t_1,t_2]]$ with the same missing coefficient.

The free and square-zero calculation remains valid in its ordinary
algebraic category. Placing it in a factorization category requires
the actual free construction there, or a comparison functor that
preserves the required structure.
Pointwise tensor closure alone does not give closure under the
coproducts used to build a tensor algebra.

\begin{proposition}\label{prop:vol2-bar-functor}
Let $F:\mathcal C\to\mathcal D$ be a strong monoidal functor
between monoidal categories in which the augmented simplicial bars
and their geometric realizations exist.
Suppose $F$ preserves those realizations.
Then for every augmented associative algebra $A$ there is a natural
equivalence
\[
 F(B_{\mathcal C}A)\simeq B_{\mathcal D}(F A).
\]
It preserves the specified unit, augmentation and all simplicial maps.
\end{proposition}
\begin{proof}
In simplicial degree $n$, the strong monoidal structure identifies
$F(A^{\otimes n})$ with $(F A)^{\otimes n}$.
It carries adjacent-multiplication faces to faces, augmentation
faces to augmentation faces, and unit insertions to degeneracies.
The coherence axioms make these identifications a map of simplicial
objects. Apply preservation of geometric realization.
\end{proof}

An ordered collision atlas and a category of consecutive word cuts
have different diagrams.
Applying the proposition to them first requires a functor on those
diagrams and a comparison of their coefficient tensors.
It also requires preservation of any support condition and completion
used in the realization.
An ordinary bar--cobar equivalence on either carrier does not supply
this functor.
The Lie-to-commutative-coalgebra chiral duality has yet another operadic
source and target. It cannot replace the associative comparison.

\subsection{A finite Chevalley model is not a finite full bar}

For $A=U(\mathfrak{sl}_2)$, PBW gives the basis
$e^ah^bf^c$, excluding the constant when forming the augmentation
ideal $I$.
Thus $I^{\otimes n}$ is countably infinite-dimensional for every
$n\ge1$.
Its finite associated-graded pieces satisfy
\[
 \sum_{d\ge0}
 \dim\operatorname{gr}^{\mathrm{PBW}}_d(I^{\otimes n})\,t^d
 =\bigl((1-t)^{-3}-1\bigr)^n.
\]
The numbers $3^n$ occur only in the least possible PBW degree $d=n$.

The Chevalley complex instead has dimensions $1,3,3,1$.
With $d(x\wedge y)=[x,y]$,
\[
 d(e\wedge h)=-2e,\qquad
 d(e\wedge f)=h,\qquad
 d(h\wedge f)=-2f,
\]
while
\[
 d(e\wedge h\wedge f)
 =-2e\wedge f-h\wedge h-2f\wedge e=0.
\]
Its homology has dimensions $1,0,0,1$.
The bar--Chevalley comparison identifies homology, not these
different chain spaces.

\section{Zeros, poles and polynomial shuffle lattices}

Let $K(u)$ be a rational kernel.
For symmetric functions in $m$ and $n$ variables, sum over
the $\binom{m+n}{m}$ ordered two-block partitions.
Insert one factor $K$ for each cross-block pair.
This defines their shuffle product.
There is no factorial divisor.

An iterated product in three blocks has one summand for each ordered
three-block partition and the same cross factors in either bracketing.
This proves associativity wherever the stated coefficient space is
closed under the formula.

Consider first the meromorphic kernel
\[
 K_{\mathrm{mer}}(u)
 =\frac{(u+\epsilon_1)(u+\epsilon_2)(u+\epsilon_3)}{u^3},
 \qquad \epsilon_1+\epsilon_2+\epsilon_3=0.
\]
Put $\sigma_2=\sum_{i<j}\epsilon_i\epsilon_j$ and
$\sigma_3=\epsilon_1\epsilon_2\epsilon_3$.
Then
\begin{equation}\label{eq:vol2-shuffle-divisor}
 K_{\mathrm{mer}}(u)=1+\frac{\sigma_2}{u^2}
                         +\frac{\sigma_3}{u^3},
 \qquad
 K_{\mathrm{mer}}(u)-K_{\mathrm{mer}}(-u)
 =\frac{2\sigma_3}{u^3}.
\end{equation}
If the parameters are nonzero and pairwise distinct,
$u=-\epsilon_a$ is a zero, not a pole.
The reciprocal has residue
\begin{equation}\label{eq:vol2-shuffle-reciprocal}
 \operatorname{Res}_{u=-\epsilon_a}
    K_{\mathrm{mer}}(u)^{-1}\,du
 =\frac{(-\epsilon_a)^3}
        {\prod_{b\ne a}(\epsilon_b-\epsilon_a)}.
\end{equation}
Both statements follow by factoring the numerator at that point.

Degenerate parameter fibers must be calculated before cancelling
coincident zeros and poles.

For the convention using $K_{\mathrm{mer}}(z_1-z_2)$, put
$u=z_1-z_2$ and $w=z_2$.
The commutator of the degree-one inputs $z^r,z^s$ is
\[
 \frac{2\sigma_3}{u^3}
   \bigl((w+u)^rw^s-(w+u)^sw^r\bigr).
\]
For $(r,s)=(1,0)$ it is $2\sigma_3/u^2$.
For $(r,s)=(2,0)$ it is
$4\sigma_3w/u^2+2\sigma_3/u$.
The latter has residue $2\sigma_3$ at the true diagonal.

Neither expression has a pole at a generic nonzero shifted zero
$u=-\epsilon_a$.
To assign shifted residues to a geometric operation one must first
specify the meromorphic coefficient of that operation.

\subsection{The coefficient lattice retained by multiplication}

Now use the distinct polynomial Hall carriers
\begin{equation}\label{eq:vol2-hall-ring}
\begin{gathered}
 R=\mathbb Q[h_1,h_2],\qquad h_3=-h_1-h_2,\\
 a=-(h_1^2+h_1h_2+h_2^2),\qquad
 b=-h_1h_2(h_1+h_2),\\
 N(v)=v^3+av+b,\qquad
 K_0(v)=N(v)/v,\qquad K_1(v)=vN(v).
\end{gathered}
\end{equation}
For $j=0,1$, independently set
\[
 \mathcal A^{(j)}=\bigoplus_{n\ge0}A_n^{(j)},\qquad
 A_n^{(j)}=R[x_1,\ldots,x_n]^{\mathfrak S_n}.
\]
Every cochain degree and parity is zero.
Charge is the separate grading $n$. Each element has finite charge
support and polynomial coefficients.
No parameter or root is inverted.

The product is
\begin{equation}\label{eq:vol2-hall-shuffle}
 (F*_jG)(X)
 =\sum_{\substack{I\sqcup J=[d+e]\\ |I|=d}}
     F(X_I)G(X_J)
       \prod_{\substack{i\in I\\ k\in J}}K_j(x_k-x_i).
\end{equation}
The unit is $1\in A_0^{(j)}=R$.
The constant polynomial $e_0=1\in A_1^{(j)}$ has charge one and
is not that unit.

\begin{proposition}\label{prop:vol2-polynomial-shuffles}
Formula \eqref{eq:vol2-hall-shuffle} gives a polynomial,
unital associative algebra.
Both kernels have exchange ratio
\[
 r(v)=\frac{K_j(-v)}{K_j(v)}
     =\frac{v^3+av-b}{v^3+av+b}.
\]
\end{proposition}
\begin{proof}
For $K_1$ polynomiality is immediate.
For $K_0$, each possible denominator is a single root difference
in a summand.
Pair two summands by exchanging the labels of that pair between
the two blocks.
On their diagonal all other factors agree and the two residues
are opposite.
Every possible denominator therefore cancels in the polynomial
unique factorization domain over $R$.

Symmetry follows by permuting the summands.
The ordered-partition argument proves associativity in every charge.
An empty block proves the unit identities.
Finally $K_1=v^2K_0$, and this even factor cancels in the ratio.
\end{proof}

The ordinary zero-potential Hall product has precisely the
numerator Euler factors from moving arrows and denominator factors
from the flag tangent space \cite[Section~2.4, Theorem~2]{KSCoHA2011}.
Three moving arrow weights $v+h_i$ give $K_0$.
Adding two zero-weight loops multiplies this kernel by $v^2$.
This geometric interpretation does not identify these ordinary
polynomial algebras with a critical Hall algebra or a native field
theory.

In charge two put
\[
 c=(x+y)/2,\qquad v=y-x,\qquad t=v^2,\qquad B=R[c,t].
\]
The entire input tensor is $R[x,y]=B\oplus vB$.
For $p=A+vC$, its linearized product is
\begin{equation}\label{eq:vol2-hall-binary}
 \mu_0(p)=2\bigl((t+a)A+bC\bigr),\qquad
 \mu_1(p)=2t\bigl((t+a)A+bC\bigr).
\end{equation}
This follows by adding $K_j(v)p(x,y)$ and
$K_j(-v)p(y,x)$.
Hence the complete binary images are
$J=(t+a,b)B$ and $tJ$.

\begin{theorem}\label{thm:vol2-hall-lattices}
There is no charge-preserving $R$-linear algebra isomorphism
$\mathcal A^{(0)}\to\mathcal A^{(1)}$.
The conclusion still holds for maps semilinear under one automorphism
of $R$.
Nevertheless the maps
\begin{equation}\label{eq:vol2-discriminant-map}
 \Phi_n(F)=D_nF,\qquad
 D_n=\prod_{i<k}(x_k-x_i)^2,\qquad D_0=D_1=1,
\end{equation}
give a unital injective algebra map.
They identify the two subalgebras generated by charge one.
They also give an isomorphism after inverting $D_n$ in each
charge, and commute with simultaneous root translation.
\end{theorem}
\begin{proof}
An isomorphism would identify the cokernels of the two full
binary products.
But
\[
 B/J=R[c]/(b),\qquad
 B/tJ\twoheadrightarrow R[c],\quad F\longmapsto F(c,0).
\]
The first has parameter annihilator $(b)$ and the second has
annihilator zero.
These cannot be isomorphic as $R$-modules.
Under a coefficient automorphism, the nonzero annihilator becomes
$(\sigma(b))$, so the semilinear case has the same contradiction.

For the positive map, the discriminant identity
\[
 D_{d+e}(X)
 =D_d(X_I)D_e(X_J)
    \prod_{\substack{i\in I\\ k\in J}}(x_k-x_i)^2
\]
holds for every ordered partition.
Replace $K_1$ by $v^2K_0$ in each summand of
$\Phi_d(F)*_1\Phi_e(G)$.
The discriminants combine to $D_{d+e}$, independent of the
partition. This proves multiplicativity.

The polynomial rings are domains, so each $\Phi_n$ is injective.
It is the identity in charges zero and one.
Its image on the charge-one-generated subalgebra is exactly the
corresponding generated subalgebra. Hence the restricted map is
an isomorphism.
Its inverse divides by $D_n$ only on this specified image.

After inverting $D_n$, the same division is defined on the whole
charge-$n$ carrier.
Products remain in this localization: every internal and cross
root difference divides a power of the total discriminant.
Finally $\sum_i\partial_{x_i}D_n=0$, which proves translation
compatibility.
\end{proof}

This is a distinction between the full polynomial realization and
the smaller algebra generated by its currents.
Equality of exchange ratios decides neither.

The quotient
$Q_n=A_n^{(1)}/D_nA_n^{(1)}$ retains the discarded collision
coefficients.
It is an $\mathcal A^{(0)}$-bimodule through
\[
 a\cdot[z]=[\Phi(a)*_1z],\qquad
 [z]\cdot a=[z*_1\Phi(a)].
\]
These are well-defined because $\Phi$ is multiplicative.

The quotient is not an algebra quotient: the image of $\Phi$
contains the unit and is a proper subalgebra.
In charge two its cone is
$[B\xrightarrow{t}B]$, with cohomology $B/(t)$ in degree zero.
On the ordered cover this is the length-two normal quotient
$R[c,v]/(v^2)$.

\section{Cooperation depends on the tensor}

The product on a polynomial shuffle carrier does not automatically
admit a coproduct into the ordinary tensor product.
The obstruction uses only charge one and charge two.
Write $e_m=x^m\in A_1^{(j)}$.
Formula \eqref{eq:vol2-hall-binary} gives
\begin{equation}\label{eq:vol2-cooperation-relation}
 [e_0,e_3]_{*_j}-3[e_1,e_2]_{*_j}
       +a[e_0,e_1]_{*_j}
 =b(e_0*_je_0).
\end{equation}
Indeed, the three commutators on the left are
$2bt^j(x^2+xy+y^2)$, $2bt^jxy$ and $2bt^j$, while
$e_0*_je_0=2t^j(t+a)$.

\begin{theorem}\label{thm:vol2-shuffle-no-cooperation}
Neither $\mathcal A^{(j)}$ admits a charge-preserving unital
$R$-algebra map
\[
 \Delta:\mathcal A^{(j)}
       \longrightarrow\mathcal A^{(j)}\otimes_R\mathcal A^{(j)}
\]
with the ordinary tensor multiplication and with two-sided counit
equal to projection onto charge zero.
The same obstruction applies to the current-generated subalgebras,
including a derived algebra map with the ordinary Koszul tensor.
Coassociativity is not needed for this conclusion.
\end{theorem}
\begin{proof}
Charge preservation leaves only the splits $(1,0)$ and $(0,1)$
for each $e_m$.
The two counits force
$\Delta e_m=e_m\otimes1+1\otimes e_m$.
Commutators of even primitive elements are primitive, since their
mixed terms cancel.
Thus the reduced coproduct of the left side of
\eqref{eq:vol2-cooperation-relation} is zero.
Multiplicativity makes the reduced coproduct of its right side
\[
 2b\,e_0\otimes e_0\ne0
       \quad\text{in }A_1\otimes_R A_1=R[x,y].
\]
This is a contradiction over the original coefficient ring.

For the ambient algebras all charge components are $R$-free.
For a current-generated subalgebra, taking degree-zero cohomology
of a proposed derived map still gives the same identity in the
free $(1,1)$ component.
Hence the derived formulation has the same contradiction.
\end{proof}

The detector can vanish on a special parameter fiber.
The theorem asserts no uniform fibrewise obstruction when $b=0$.
At $(h_1,h_2,h_3)=(1,1,-2)$, however, it becomes
$-4e_0\otimes e_0$, so the obstruction is already visible over
$\mathbb Q$.
Neither this result nor its proof concerns every possible
cut-convolution coalgebra or every shifted coefficient tensor.

\subsection{A rational cooperation with a different interchange}

There is an explicit positive construction after localizing the
spectral coefficients and changing the middle interchange.
For ordered root lists $X,Y$, put
\[
 Q_{d,e}^{(j)}(X,Y)
   =\prod_{x\in X,\ y\in Y}K_j(y-x),\qquad
 \Delta_{d,e}^{\mathrm{rat}}F
   =\frac{F(X,Y)}{Q_{d,e}^{(j)}(X,Y)}.
\]
The recipient consists of block-symmetric polynomials localized
at all pair differences and their translates by $\pm h_i$.
Parameters themselves are not inverted.

Let $F$ have bidegree $(d,e)$ and $G$ have bidegree $(d',e')$.
Define their product by summing over
$X=X_1\sqcup X_2$, $Y=Y_1\sqcup Y_2$ of these sizes, with summand
\begin{equation}\label{eq:vol2-rational-interchange}
\begin{aligned}
 &F(X_1,Y_1)G(X_2,Y_2)
   \prod_{\substack{x\in X_1\\x'\in X_2}}K_j(x'-x)\\
 &\qquad\cdot
   \prod_{\substack{y\in Y_1\\y'\in Y_2}}K_j(y'-y)
   \prod_{\substack{y\in Y_1\\x'\in X_2}}r(y-x').
\end{aligned}
\end{equation}
The last product is the middle interchange.
Its factors would be absent from the ordinary tensor multiplication.

\begin{proposition}\label{prop:vol2-rational-cooperation}
The multiplication \eqref{eq:vol2-rational-interchange} is
associative and unital.
The sum of the maps $\Delta_{d,e}^{\mathrm{rat}}$ is multiplicative
into this recipient and is rationally coassociative.
It has both charge-zero counits.
\end{proposition}
\begin{proof}
For multiplicativity, consider a term with four ordered groups
$X_1,Y_1,X_2,Y_2$.
Divide its source shuffle product by the cross kernel between
$X_1\sqcup X_2$ and $Y_1\sqcup Y_2$.
The factors from $X_1$ to $Y_2$ cancel.
The remaining middle ratio is
\[
 \frac{K_j(x'-y)}{K_j(y-x')}=r(y-x'),
 \qquad y\in Y_1,\ x'\in X_2.
\]
The other factors are exactly those in
\eqref{eq:vol2-rational-interchange}, together with
$\Delta F$ and $\Delta G$.
Summing over all four groups proves the identity.

For three inputs, either bracketing includes the same internal
kernel factors and the same factors from an earlier right block
to a later left block.
This proves associativity.
For three output lists, either iterated cooperation divides by
$Q_{12}Q_{13}Q_{23}$ in the common localization.
This proves coassociativity.
The empty lists give the unit and both counits.
Every input has finite charge support, so these sums are finite.
\end{proof}

The four-current relation is now compatible with
\[
 (1-r(v))v(v^2+a)=b(1+r(v)),
\]
which follows directly from
$r(v)=(v^3+av-b)/(v^3+av+b)$.
The mixed contribution excluded by the ordinary tensor is present.
The construction still leaves the polynomial carrier:
\[
 \Delta_{1,1}^{\mathrm{rat}}(1\in A_2)
  =\frac{v}{N(v)}\quad(j=0),\qquad
 \Delta_{1,1}^{\mathrm{rat}}(1\in A_2)
  =\frac{1}{vN(v)}\quad(j=1).
\]
It also removes the closed loci at its inverted denominators.
An equivalence on this open coefficient domain is not yet a
supported comparison across those loci.

\subsection{Shifted collisions retain their excess}

Keep the same coefficient ring $R$.
On three position variables consider
\[
 f_{12}=z_1-z_2+h_1,\qquad
 f_{23}=z_2-z_3+h_1.
\]
Their intersection is $R[z_3]$, with
$z_2=z_3-h_1$ and $z_1=z_3-2h_1$.
The composed endpoint shift is $-2h_1$.
Thus the finite list of shifts $0,\pm h_1,\pm h_2,\pm h_3$
is not closed under composition of these shifted collisions.

For an ordered list of equations, use the Koszul complex with
degree-minus-one generators $\xi_{ij}$ and
$d\xi_{ij}=f_{ij}$.
Adding the induced endpoint equation gives
\[
 f_{13}=z_1-z_3+2h_1=f_{12}+f_{23}.
\]
The change of exterior generator
$\eta=\xi_{13}-\xi_{12}-\xi_{23}$ yields
\begin{equation}\label{eq:vol2-shifted-excess}
 K(f_{12},f_{23},f_{13})
 \simeq R[z_3]\otimes\Lambda(\eta),
 \qquad |\eta|=-1,\qquad d\eta=0.
\end{equation}
Indeed, $f_{12},f_{23},z_3$ are polynomial coordinates over $R$,
so their first two Koszul pairs contract. $\eta$ remains.
Discarding the redundant equation would discard its nonzero
negative-degree class.

The true endpoint diagonal instead restricts to $-2h_1$.
Its derived restriction is
\[
 [R[z_3]\xrightarrow{-2h_1}R[z_3]]
\]
in degrees $-1,0$.
Before specialization its only cohomology is $R[z_3]/(h_1)$
in degree zero.
Derived specialization to $h_1=0$ makes the differential zero and
retains both degrees.

There is a separate obstruction to viewing the scalar quotient as
the unchanged supported differential module.
For $P=R[z_1,z_2]$ and $f=z_1-z_2-\lambda$, with
$\lambda\in R$, a relative differential-module structure on
$P/(f)$ would satisfy $[\partial_{z_1},f]=1$.
But $f$ acts by zero on that module, so its commutator with every
operator is zero.
The quotient is nonzero, giving a contradiction.

The same argument applies to a differential-module complex whose
underlying $P$-complex is quasi-isomorphic to $P/(f)$.
Its degree-zero cohomology would have that impossible action.
A supported differential-module realization must retain its normal
derivatives and the maps attaching them.
The scalar Koszul calculation does not supply those maps.

\section{Matrix coefficients reconstruct a specified tensor category}

The preceding examples distinguish a product, its cooperation and
its coefficient domain.
Tensor reconstruction begins with all of them encoded in a category.
Let $\mathcal C$ be an essentially small $k$-linear monoidal category,
and let
\[
 \omega:\mathcal C\longrightarrow\operatorname{Vect}_{\mathrm{fd}}(k)
\]
be a $k$-linear strong monoidal functor.
Write $V_M=\omega(M)$ and use its monoidal identifications
$V_M\otimes V_N\simeq V_{M\otimes N}$ and $V_{\mathbf1}\simeq k$.
Choose a small skeleton when forming the direct sum below.

The matrix-coefficient coend is the vector space
\begin{equation}\label{eq:vol2-coend}
 H_\omega=
 \left(\bigoplus_{M\in\mathcal C}V_M^*\otimes V_M\right)
 \big/
 \left\langle
 c_M(\lambda\omega(f),v)-c_N(\lambda,\omega(f)v)
 \right\rangle.
\end{equation}
Here $f:M\to N$, $v\in V_M$, $\lambda\in V_N^*$, and
$c_M(\lambda,v)$ denotes the class of the indicated tensor.
The relations remember every morphism of $\mathcal C$.

\begin{proposition}\label{prop:vol2-coend-operations}
The formulas
\begin{equation}\label{eq:vol2-coend-coalgebra}
\begin{aligned}
 \Delta c_M(\lambda,v)
   &=\sum_i c_M(\lambda,e_i)\otimes c_M(e^i,v),\\
 \varepsilon c_M(\lambda,v)&=\lambda(v)
\end{aligned}
\end{equation}
are independent of a basis $(e_i)$ and its dual $(e^i)$.
They define a counital coassociative coalgebra.
Its universal coaction is
\[
 \delta_M(v)=\sum_i e_i\otimes c_M(e^i,v).
\]
The monoidal structure supplies the multiplication
\begin{equation}\label{eq:vol2-coend-product}
 c_M(\lambda,v)c_N(\mu,w)
   =c_{M\otimes N}(\lambda\otimes\mu,v\otimes w),
\end{equation}
with unit $c_{\mathbf1}(1,1)$.
Together these operations make $H_\omega$ a bialgebra.
If $\mathcal C$ is rigid, it has an invertible antipode.
\end{proposition}
\begin{proof}
The tensor $\sum_i e_i\otimes e^i$ is the identity of $V_M$,
so its insertion in \eqref{eq:vol2-coend-coalgebra} is
basis-independent.
For $f:M\to N$, inserting this identity before or after $\omega(f)$
gives the same linear map.
The coend relation therefore makes $\Delta$ well-defined.
The same relation respects evaluation, so it also defines
$\varepsilon$.

Inserting the identity twice gives
\[
 \sum_{i,j}
 c_M(\lambda,e_i)\otimes
 c_M(e^i,e_j)\otimes
 c_M(e^j,v)
\]
for either iterated coproduct.
Evaluating either outer pair gives $c_M(\lambda,v)$, proving
both counits.
These same calculations prove the coaction identities.
The coend relation says exactly that each $\omega(f)$ is
colinear.

Naturality of the monoidal constraints makes
\eqref{eq:vol2-coend-product} compatible with the relations.
Their associativity and unit coherences give the algebra laws.
Using the basis $e_i\otimes f_j$ of $V_M\otimes V_N$ shows
that $\Delta$ preserves multiplication. Evaluation shows the
same for the counit.

For a left dual $M^\vee$, identify $V_{M^\vee}$ with $V_M^*$.
Let $\widetilde v\in V_M^{**}$ be evaluation at $v$.
Then
\[
 S\,c_M(\lambda,v)=c_{M^\vee}(\widetilde v,\lambda).
\]
The evaluation morphism $M^\vee\otimes M\to\mathbf1$ and the
coevaluation morphism $\mathbf1\to M\otimes M^\vee$ are
morphisms of $\mathcal C$.
Applying their coend relations gives the two antipode identities:
the two sums contract to $\lambda(v)c_{\mathbf1}(1,1)$.
Right duals give the inverse antipode by the corresponding
right-dual formula.
\end{proof}

Coevaluation gives the coproduct, whereas the tensor product of
objects gives multiplication.
Composition has a third role.
The full linear dual $H_\omega^*$ is the algebra of linear natural
endomorphisms of $\omega$.
An endomorphism $T$ gives the functional
$c_M(\lambda,v)\mapsto\lambda(T_Mv)$, and
\[
 (\phi_T*\phi_U)c_M(\lambda,v)
   =\sum_i\lambda(T_Me_i)e^i(U_Mv)
   =\lambda(T_MU_Mv).
\]
Thus convolution is composition.

Tensor-compatible natural automorphisms must additionally preserve
the monoidal constraints and the unit.
They are not all elements of this dual algebra.

\begin{theorem}[Tensor reconstruction]\label{thm:vol2-tensor-reconstruction}
Suppose $\mathcal C$ is $k$-linear abelian, with
finite-dimensional morphism spaces and finite-length objects.
Suppose it is rigid, its tensor product is exact in both variables,
and $\omega$ is exact and faithful.
Then the comparison given by $\delta_M$ is a monoidal equivalence
\[
 \mathcal C\simeq
       \operatorname{Comod}_{\mathrm{fd}}(H_\omega).
\]
\end{theorem}
\begin{proof}
The coalgebra reconstruction theorem for an exact faithful functor
to finite-dimensional vector spaces identifies $\mathcal C$ with
the finite-dimensional right comodules of the coend
\eqref{eq:vol2-coend}
\cite[Theorem~1.10.1]{EGNO2015}.
Its comparison is exactly the coaction
\eqref{eq:vol2-coend-coalgebra}.
Equation \eqref{eq:vol2-coend-product} identifies the coaction
on $V_{M\otimes N}$ with the tensor product of the two coactions.
The given monoidal constraints therefore make the equivalence
strong monoidal, including its unit.
Proposition~\ref{prop:vol2-coend-operations} supplies the antipode,
as in the Hopf reconstruction theorem
\cite[Theorem~5.4.1]{EGNO2015}.
\end{proof}

For a braided $\mathcal C$, write
$\beta_{M,N}:V_M\otimes V_N\to V_N\otimes V_M$ for the induced
linear exchange.
Its coefficient functional is
\begin{equation}\label{eq:vol2-coend-braiding}
 \mathfrak r\bigl(c_M(\lambda,v),c_N(\mu,w)\bigr)
   =(\mu\otimes\lambda)\beta_{M,N}(v\otimes w).
\end{equation}
Naturality makes this well-defined on both coends.
The inverse braiding gives its convolution inverse.

Taking matrix coefficients of the two hexagon diagrams gives the
two coquasitriangular identities. Taking coefficients of the
colinearity square gives compatibility with multiplication.
Thus \eqref{eq:vol2-coend-braiding} is the coquasitriangular form.

Under the reconstruction equivalence it recovers the braiding by
\[
 v\otimes w\longmapsto
    \sum w_{(0)}\otimes v_{(0)}
      \mathfrak r(v_{(1)},w_{(1)}),
 \qquad
 \delta(v)=\sum v_{(0)}\otimes v_{(1)}.
\]

When $H_\omega$ is infinite-dimensional, the coefficient functional
lives in $(H_\omega\otimes H_\omega)^*$.
Regard $H_\omega^*$ with the inverse-limit topology from its finite
coefficient coalgebras, and define
$H_\omega^*\widehat\otimes H_\omega^*
=(H_\omega\otimes H_\omega)^*$ with that topology.
The functional is then the corresponding completed $R$-operator.
This does not place it in the ordinary algebraic tensor
$H_\omega^*\otimes H_\omega^*$.
A meromorphic version additionally needs specified expansion maps,
domains and continuity for its spectral parameter.

\subsection{An exchange relation is not a presentation}

For a selected object, the naturality of its exchange gives the
matrix relation usually written
$RT_1T_2=T_2T_1R$.
These relations define an algebra with a comparison map to the coend.
They form a presentation precisely when that map is an isomorphism.
The category's morphisms can impose further relations.

For example, compare finite-dimensional vector spaces with
finite-dimensional $\mathbb Z$-graded vector spaces, using the
ordinary flip and the forgetful fiber functor in both cases.
For vector spaces, every map $k\to V$ is allowed.
The coend relations therefore give
\[
 c_V(\lambda,v)=\lambda(v)c_k(1,1),\qquad H_\omega=k.
\]

For graded vector spaces, decompose into homogeneous lines.
There are no morphisms between distinct degrees.
There is one coefficient $u_n$ in each degree, and the formulas give
\[
 \Delta u_n=u_n\otimes u_n,\qquad
 u_nu_m=u_{n+m},\qquad u_0=1,\qquad S(u_n)=u_{-n}.
\]
Hence $H_\omega=k[t,t^{-1}]$ with $u_n=t^n$.

The exchange on a chosen one-dimensional fiber is the scalar $1$
in both examples.
The one-by-one RTT relation imposes no further condition on its
coefficient $t$.
In the first category the morphisms additionally impose $t=1$.
The same scalar exchange therefore leads to different reconstructed
Hopf algebras.

The geometric action algebra of a collection of kernels records
the operators they generate on a state space.
The coend instead records natural symmetries respecting every
chosen kernel morphism.
Identifying these two constructions requires a comparison of their
actions and tensor structures.
A list of root multiplicities, or a common exchange ratio, provides
neither that map nor its inverse.

\section{Geometric comparison preserves equations, not only counts}

Two positively graded Lie algebras can have the same graded PBW
character and different brackets.
Give $x,y$ weight one and $z$ weight two.
The abelian Lie algebra on these generators and the Heisenberg
Lie algebra with $[x,y]=z$ have the same PBW Hilbert series
$((1-t)^2(1-t^2))^{-1}$.
One enveloping algebra is commutative and the other is not.
The equality of their characters therefore supplies no algebra
isomorphism.

Similarly, the zero super vector space and a space with one even
and one odd basis vector both have superdimension zero.
An index cannot recover the forgotten parity dimensions.

An equivalence of underlying tangent complexes is also weaker
than a comparison of deformation equations.
Let $L$ have basis $u$ in degree one and $v$ in degree two, with
zero differential.
Consider the two dg Lie structures
\[
 L_{\mathrm{ab}}:\ [\, ,\,]=0,\qquad
 L_{\mathrm{q}}:\ [u,u]=2v,\quad [u,v]=[v,v]=0.
\]
The graded bracket on the odd element $u$ is symmetric.
Every iterated bracket is zero because $v$ is central, so both
structures satisfy Jacobi.
Their tangent complexes are identical.

\begin{proposition}\label{prop:vol2-nonlinear-comparison}
The first-order deformation $\epsilon u$ over
$k[\epsilon]/(\epsilon^2)$ extends to
$k[\epsilon]/(\epsilon^3)$ for $L_{\mathrm{ab}}$ and does not
extend for $L_{\mathrm{q}}$.
Thus the identity of tangent complexes cannot extend to an
$L_\infty$ equivalence between these deformation problems.
\end{proposition}
\begin{proof}
The Maurer--Cartan equation is
$d\alpha+\tfrac12[\alpha,\alpha]=0$.
It vanishes identically in the abelian case.
In the second algebra every possible lift has the form
$\alpha=(\epsilon+c\epsilon^2)u$.
Its curvature modulo $\epsilon^3$ is $\epsilon^2v\ne0$.
There is no degree-zero element giving a gauge correction.

Equivalently, an $L_\infty$ map with identity first component
would have to preserve the binary bracket up to the differential
of its second component.
Both differentials are zero, so it would identify the zero
bracket with $[u,u]=2v$, which is impossible.
\end{proof}

A geometric field comparison must consequently carry the nonlinear
operations as well as its tangent map.
For a cotangent BV theory it must also carry the pairing, action
and boundary term.
The classical master equation of a supplied coupled algebra proves
the consistency of that algebra.
It does not construct the action of an independently specified
gravitational or matter field theory on it.
At the quantum level, the chosen regularization and the local
quantum master equation remain additional comparisons, as in
Chapter~\ref{ch:vol2-native-bulk-boundary}.

\subsection{A projector must respect the actual transition maps}

A finite decomposition model has a simple idempotent:
on a sum indexed by pairs $(d_1,d_2)$, retain precisely the
summands with $d_1=d_2$.
Writing its inclusion and retraction as $\iota,\pi$ gives
\[
 \pi\iota=1,\qquad \iota\pi=e_{\mathrm{al}}.
\]
It is an equivalence on the full object exactly when the
off-aligned summands vanish.
These equations determine the invertibility of this canonical
comparison. They do not rule out every possible noninvertible
distributive law.

To descend the projector, equality of indices must be compatible
with each transition correspondence.
Preserving equality alone is insufficient.
The map $\{0,1\}\to\{*\}$ preserves equal pairs but sends the
off-diagonal pair $(0,1)$ to the diagonal pair $(*,*)$.
Extension along this map does not commute with the two projectors.

In a six-functor construction, assume the actual squares with the
diagonal index locus are Cartesian.
Assume also that their base-change maps are isomorphisms.
Under those hypotheses, base change identifies the pull--push
idempotents across every transition. The triple-diagonal identity
then descends as well.
Finite combinatorics alone does not verify those geometric squares.

\section{Gravitational spectra, holonomy and boundary data}

Three-dimensional Cartan gravity begins with independent fields
$e,\omega$.
For negative cosmological constant, the two connections are
$A^\pm=\omega\pm\ell^{-1}e$.
On the invertible-coframe locus, their flatness equations give
torsion zero and constant curvature.
This is a field construction with an action comparison.
Specifying two boundary central charges or a list of current spins
does not supply these fields or their action on the boundary.

\subsection{The principal embedding and its spin spectrum}

Let $V_m=\operatorname{Sym}^m(k^2)$ be the irreducible
$\mathfrak{sl}_2$ module of highest weight $m$.
The principal embedding in $\mathfrak{sl}_N$ acts on its defining
representation as $V_{N-1}$.
Its adjoint therefore decomposes as
\begin{equation}\label{eq:vol2-principal-spectrum}
 \mathfrak{sl}_N
 \cong\bigoplus_{s=2}^{N}V_{2s-2},\qquad
 \sum_{s=2}^{N}(2s-1)=N^2-1.
\end{equation}
Indeed, the tensor product
$V_{N-1}\otimes V_{N-1}^*$ contains once each highest weight
$0,2,\ldots,2N-2$.

This follows by subtracting successively the characters of those
highest weights from its weight character. Finite-dimensional
$\mathfrak{sl}_2$ modules are semisimple.
Removing the scalar endomorphisms removes $V_0$, proving
\eqref{eq:vol2-principal-spectrum}.

For a specified principal Chern--Simons theory linearized about
its gravitational $\mathfrak{sl}_2$ background, these summands
are the frame-like fields of spins $2,\ldots,N$.
With principal Drinfeld--Sokolov boundary conditions, the reduced
classical currents have the same spins
\cite{CampoleoniHigherSpin,HenneauxRey}.
This matches the representations in that construction.
Arbitrary supplied $W_N$ boundary data need a comparison from this
boundary reduction.
It must preserve the Poisson operations and the bulk coupling.
The dimension identity alone does not construct it.

\subsection{Connection coefficients and holonomy are different functions}

Work over $\mathbb C$ in this subsection, and let $L>0$.
Consider a constant traceless companion matrix $A(q)$ of size $N$:
\[
 A_{i,i+1}=1\quad(1\le i<N),\qquad
 A_{N,j}=q_{N+1-j}\quad(1\le j<N),\qquad A_{N,N}=0,
\]
with all remaining entries zero.
These are the constant coefficients of the corresponding
companion-gauge oper.
Expansion of the characteristic determinant gives
\begin{equation}\label{eq:vol2-companion-polynomial}
 \det(\lambda I-A(q))
   =\lambda^N-\sum_{s=2}^{N}q_s\lambda^{N-s}.
\end{equation}
One can also recover the same polynomial from the eigenvector
recursion $v_{i+1}=\lambda v_i$ and the final row.

For a circle of length $L$ and the parallel-transport convention
$\partial_\theta\psi=A(q)\psi$, the holonomy is
$H(q)=\exp(LA(q))$.
It is an exponential of the matrix, not its characteristic
polynomial.

\begin{proposition}\label{prop:vol2-holonomy-logarithm}
Holonomy conjugacy does not in general determine the coefficients
$q_s$.
Suppose $A(q)$ is regular semisimple.
Choose a domain containing its eigenvalues on which
$\lambda\mapsto e^{L\lambda}$ is injective.
The corresponding logarithm branch then recovers its coefficients
from holonomy conjugacy.
\end{proposition}
\begin{proof}
For $N=2$,
\[
 A(q)=\begin{pmatrix}0&1\\q&0\end{pmatrix},\qquad A(q)^2=qI,
\]
so the exponential series gives
\begin{equation}\label{eq:vol2-sl2-holonomy}
 H(q)=\cosh(L\sqrt q)\,I
       +\frac{\sinh(L\sqrt q)}{\sqrt q}\,A(q),
 \qquad
 \operatorname{Tr}H(q)=2\cosh(L\sqrt q).
\end{equation}
Both coefficients are entire functions of $q$. The second has
value $L$ at $q=0$.
For $L=2\pi$ and $q=-n^2$, every positive integer $n$ gives
$H(q)=I$.
Distinct connection coefficients therefore have the same holonomy.
This is already a counterexample among constant companion
connections.

Under the additional hypothesis, the distinct eigenvalues of
$H(q)$ give the unordered eigenvalues
$\lambda_i=L^{-1}\log\lambda_i(H)$ of $A(q)$.
They determine its characteristic polynomial.
Equation \eqref{eq:vol2-companion-polynomial} then recovers every
$q_s$.
\end{proof}

For nonconstant boundary fields the holonomy is a path-ordered
exponential. Its traces are global functions of those fields.
The local Drinfeld--Sokolov current brackets are obtained from
Hamiltonian reduction. They are not identified with the algebra
of these holonomy traces by a matrix characteristic polynomial.
Boundary conditions, a chosen logarithm or spectral family, and
the actual Hamiltonian comparison are therefore part of any
charge-identification theorem.
The constant example does not decide which of its connections
satisfies a separately imposed physical regularity condition.

\subsection{Algebraic removal of a central curvature}

Let $\mathfrak g$ be a dg Lie algebra and let
$\kappa\in\mathfrak g^2$ be closed and central.
Adjoin a central generator $u$ of degree one with $du=\kappa$.
The differential remains square-zero and a derivation:
the only new derivation equation is $[\kappa,x]=0$.
The class of $\kappa$ becomes exact in the larger algebra.
If $\kappa$ is the curvature term of the corresponding curved
Maurer--Cartan equation, twisting by $-u$ changes that term to
\[
 \kappa+d(-u)+\tfrac12[-u,-u]=0.
\]
This is an explicit algebraic construction.

It can change the deformation problem: for
$\mathfrak g=k\kappa$ with zero differential, the original
degree-two cohomology is nonzero, whereas the extension
$[ku\xrightarrow{d}k\kappa]$ is acyclic.

A geometric anomaly cancellation requires more.
Suppose an independently constructed invertible field theory has
the opposite anomaly.
A comparison must identify both classes in the same complex and
normalization.
Under these hypotheses, tensoring the theories cancels their sum.
Adjoining $u$ does not construct that invertible theory, its
descent data or the comparison of anomaly complexes.
The algebraic curvature, a determinant-line curvature and a
local BV obstruction retain their separate meanings until
such maps are supplied.

\section{The remaining geometric comparisons}

The constructions above preserve definite operations.
Ordinary bars retain augmented multiplication.
The discriminant map retains the shuffle product and translation
on its stated image.
Rational cooperation retains a changed middle interchange on an
open coefficient domain.
Tensor reconstruction retains a specified category, fiber
functor and all its morphisms.
Cartan gravity and companion-matrix holonomy begin with actual
connection fields.

To compare any of these with native chiral quantum symmetry, the
map must be defined before its equivalence is asserted.
For the shuffle models it must retain the integral binary images,
the shifted collisions and the excess class
\eqref{eq:vol2-shifted-excess}, with a specified tensor and counit.
For boundary reconstruction it must retain the derived local
mapping objects and relative sewing of
Problem~\ref{prob:vol2-marked-physical-reconstruction}.
For a gravitational realization it must carry the nonlinear
field equations, pairing, action and boundary reduction,
followed by the chosen quantum and state comparisons.
These requirements leave room for restricted coefficient domains,
larger mapping categories and interfaces between different
realizations. Each possibility has its own comparison map and
its own test of invertibility.

\chapter*{Hamiltonian jets and boundary quantization}
\addcontentsline{toc}{chapter}{Hamiltonian jets and boundary quantization}

Quantization of a symplectic boundary changes the product of observables while retaining the classical Poisson bracket to first order. The formal symplectic plane gives a concrete setting in which to compare Hamiltonian operations with operations on a quantized brane. Pointed Hamiltonians vanish to order at least two; translations belong to a separately adjoined affine sector. Finite Hamiltonian jets and their shifted cotangent fields provide the local algebraic models.

The universal boundary brane leads to a Rees enveloping algebra. Its relative Hochschild complex is compared with a Chevalley complex whose differential carries the same Rees parameter. Duflo quantization gives the separate invariant comparison, and the quadratic Moyal calculation determines the oscillator Casimir quotient. These finite constructions distinguish the underlying cochain comparison from higher operations on it.

The all-arity open--closed problem asks for compatible closed and mixed operations on the boundary algebra. Its stated refinement supplies the operadic maps and coherence homotopies. Passing through the Hamiltonian jet limit also requires actual cochain transition systems and their continuous limit identifications. A geometric brane realization must supply the boundary module and compare its centre action with the chosen observable or trace maps; global realization further requires descent and quantization data.

\cleardoublepage
\phantomsection
\addcontentsline{toc}{part}{Mixed Holomorphic--Topological Deligne Theory}
\chapter*{Book entrance: Mixed Holomorphic--Topological Deligne Theory}
\addcontentsline{toc}{chapter}{Book entrance: Mixed Holomorphic--Topological Deligne Theory}
\begin{summarybox}
Hamiltonian jets, shifted cotangent fields, universal branes, Duflo quantization, and the all-arity open--closed map.
\end{summarybox}
\part{The finite-jet laboratory}

\chapter{Moyal quantization in coordinates}
\section{The bidifferential formula}
For polynomials on $\C^2$, the Moyal product is
\begin{equation}\label{eq:vol2-moyal}
 f\star g
 =\sum_{n\ge0}\frac1{n!}\left(\frac{\hbarq}{2}\right)^n
 \sum_{r=0}^n(-1)^r\binom nr
 (\partial_1^{n-r}\partial_2^rf)
 (\partial_1^r\partial_2^{n-r}g).
\end{equation}
Every polynomial pair gives a finite sum.

\begin{theorem}[Quadratic exactness]
For polynomials $f,g$ of degree at most two,
\[
 [f,g]_\star=\hbarq\{f,g\}.
\]
The quadratic generators therefore define a Lie representation of the Rees-scaled $\mathfrak{sl}_2$.  Their associative algebra has an additional Casimir relation.
\end{theorem}
\begin{proof}
The commutator retains the odd powers of the Poisson bidifferential.  Derivatives of order three vanish on quadratics, so the first-order term is the entire commutator.
\end{proof}

\begin{proposition}[The oscillator quotient]\label{prop:vol2-oscillator-quotient}
Let $R_N=\C[\hbarq]/(\hbarq^{N+1})$ and let $W_N=R_N[q,p]$
have the Moyal product for $\{q,p\}=1$.
Let $U_N$ have generators $E,F,H$ with relations
\[
 [H,E]=2\hbarq E,\qquad [H,F]=-2\hbarq F,\qquad [E,F]=\hbarq H.
\]
Set $\Omega=EF+FE+H^2/2$.
The assignments
\[
 E\longmapsto e=q^2/2,\qquad
 F\longmapsto f=-p^2/2,\qquad
 H\longmapsto h=-qp
\]
induce an algebra isomorphism
\[
 U_N/(\Omega+3\hbarq^2/8)
 \xrightarrow{\sim} W_N^{(q,p)\mapsto(-q,-p)}.
\]
These compatible isomorphisms also hold after $\hbarq$-adic completion.
If linear Hamiltonians are retained, the generated associative algebra is the full Weyl algebra.
\end{proposition}
\begin{proof}
Quadratic exactness proves the three commutator relations.
The associative products are
\[
\begin{split}
 e\star f&=-q^2p^2/4-\hbarq qp/2-\hbarq^2/8,\\
 f\star e&=-q^2p^2/4+\hbarq qp/2-\hbarq^2/8,\\
 h\star h&=q^2p^2-\hbarq^2/4.
\end{split}
\]
Thus $\Omega$ maps to $-3\hbarq^2/8$.
Its commutator with each of $E,F,H$ vanishes by the defining relations.
Give these generators PBW degree one.
Give an even monomial $q^ap^b$ degree $(a+b)/2$.
The leading symbol of $\Omega+3\hbarq^2/8$ is
$C=2EF+H^2/2$.
It is a non-zero-divisor in $R_N[E,F,H]$, since $2C$ is monic in $H$.
Centrality and this property give
\[
 \operatorname{gr}\bigl(U_N/(\Omega+3\hbarq^2/8)\bigr)
 =R_N[E,F,H]/(C).
\]
Indeed, every element of the two-sided ideal is a multiple of the central generator.
Its leading symbol is the product of the two leading symbols.
The quotient has basis $E^aH^bF^c$ with $a,c\ge0$ and $b=0,1$.
Their leading symbols map to unit multiples of the distinct monomials
$q^{2a+b}p^{2c+b}$, precisely the even monomials.
Hence the associated-graded map is an isomorphism.
Lifting highest-degree terms proves surjectivity and comparing them proves injectivity.
Every element at fixed $N$ has finite polynomial degree.
The formulas commute with reduction in $N$, so their inverse limit gives the formal isomorphism.
Finally $q,p$ generate $W_N$ and topologically generate its formal limit.
\end{proof}

The primitive coproduct on $U_N$ does not descend to this quotient.
Modulo $\hbarq$, it sends the classical relation to
\[
 \Delta(C)=C\tensor1+1\tensor C
   +2E\tensor F+2F\tensor E+H\tensor H.
\]
The cross-term is nonzero in the tensor square of $\C[E,F,H]/(C)$.
For example, evaluation at $(E,F,H)=(1,0,0)$ in the first factor
and $(0,1,0)$ in the second gives $2$.
Thus the oscillator quotient is an associative algebra, with no coproduct inherited from these primitive generators.

\begin{calculation}[Cubic correction]
For $f=z_1^3$ and $g=z_2^3$,
\[
 f\star g-g\star f
 =9\hbarq z_1^2z_2^2+\frac32\hbarq^3.
\]
The first term quantizes the quartic Hamiltonian channel.  The constant is the first wheel-free ordering correction visible beyond the quadratic Lie algebra.
\end{calculation}
\begin{proof}
In \eqref{eq:vol2-moyal}, the odd orders are $n=1$ and $n=3$.  The first contributes $9\hbarq z_1^2z_2^2$.  The third contributes $2(\hbarq/2)^3(3!)=3\hbarq^3/2$.
\end{proof}

\section{Weyl generators}
With the convention $\{q,p\}=1$, the symbols satisfy
$[q,p]_\star=\hbarq$.
Their Weyl-ordered differential operators are
\[
 \widehat q=q,\qquad \widehat p=-\hbarq\partial_q,
 \qquad [\widehat q,\widehat p]=\hbarq.
\]
The abstract Weyl algebra is defined by this commutator relation.
Weyl ordering transports its multiplication to the Moyal product.
The differential-operator representation realizes these relations over
$\C[\hbarq]$ and its formal completion.
The finite quotient argument above uses the abstract $R_N$-algebra,
so it does not require a faithful differential-operator representation at nilpotent $\hbarq$.

\chapter{The quadratic invariant}
\section{The \texorpdfstring{$\mathfrak{sl}_2$}{sl2} invariant}
For the quadratic jet, identify the coordinate functions with $e,f,h$.  The invariant
\[
 q=2ef+\frac12h^2
\]
is $d_{\CE}$-closed.  Its formal Duflo image is $\Omega+\hbarq^2/2$, where $\Omega=ef+fe+h^2/2$ in the Rees enveloping algebra.  This is the degree-zero invariant comparison.  Under the oscillator quotient of Proposition~\ref{prop:vol2-oscillator-quotient}, its image is $\hbarq^2/8$.

\chapter{All-arity open--closed operations}
\section{Hochschild braces}
For Hochschild cochains $F\in C^p(A,A)$ and $G_i\in C^{q_i}(A,A)$, the brace
\[
 F\{G_1,\ldots,G_r\}
\]
is the signed sum over ordered insertions of the $G_i$ into the inputs of $F$.  The brace identities state that a nested insertion is the sum over all ways of distributing the inner insertions among the outer slots.

\begin{theorem}[Brace realization with a refinement]
For $A=A_d$, supply the homotopy-$\Etwo$ and mixed refinements of
Corollary~\ref{cor:vol2-open-closed-refinement}.
They identify the selected homotopy Gerstenhaber operations on
$\mathcal P_d$ with the relative Hochschild centre operations.
The chosen mixed refinement also identifies their action on the boundary algebra.
\end{theorem}
\begin{proof}
The relative Hochschild complex has its brace operations.
The supplied operadic equivalence preserves their coherent homotopies.
An equivalence of complexes alone would not imply this conclusion.
\end{proof}

\section{The first mixed identities}
A relative Hochschild $1$-cocycle $D$ of an algebra concentrated in degree zero
is an $R$-linear derivation:
\[
 D(ab)=D(a)b+aD(b).
\]
For two such cocycles $D,E$, their Gerstenhaber bracket acts by
\[
 [D,E](a)=D(E(a))-E(D(a)).
\]
For a graded algebra and homogeneous derivations, these formulas become
\[
\begin{split}
 D(ab)&=D(a)b+(-1)^{|D||a|}aD(b),\\
 [D,E](a)&=D(E(a))-(-1)^{|D||E|}E(D(a)).
\end{split}
\]
Higher Hochschild cochains act through multilinear evaluations and braces.
They need not act by derivations or carry a coproduct defining a diagonal action.
A compatible mixed refinement transports these operations, including their homotopies.

\begin{corollary}[Continuous open--closed structure with compatible jets]
Under the jet comparison and mixed refinement hypotheses of
Theorem~\ref{thm:vol2-compatible-jets}, the pair
$(\mathcal P_\infty,\widehat U_\hbar(\h))$
has the corresponding complete homotopy Swiss-cheese action.
Its closed color is equivalent to the continuous relative derived centre.
\end{corollary}
\begin{proof}
The supplied finite operations commute with the transition maps
and converge in the specified completion.
Their limit acts on the two limit objects.
The comparison theorem identifies the closed complex with continuous relative Hochschild cochains.
\end{proof}

\chapter{Finite Chan--Paton ranks}
\section{Rank one}
For $N=1$, the commuting equation is automatic.  The boundary coordinate algebra is $\C[x,y]$, and
\[
 J_1(f)=f(x,y).
\]
The necklace bracket reduces to the ordinary Poisson bracket.  Restriction to $f\in\mathfrak m^2$ gives the defining pointed-Hamiltonian representation.  The two linear Hamiltonians act by translations and recover the affine center-of-mass extension.

\section{Rank two and Cayley--Hamilton reduction}
For a $2\times2$ matrix $X$,
\[
 X^2-(\Tr X)X+(\det X)I=0.
\]
Hence every trace containing two consecutive powers of $X$ reduces to traces with $X$-degree at most one after adjoining $\Tr X$ and $\det X$.  The same statement holds for $Y$.  On the commuting locus, simultaneous diagonalization identifies the invariant algebra with the symmetric polynomials in two ordered pairs.

\begin{proposition}[First finite-rank trace relation]
For commuting $2\times2$ matrices,
\[
 p_{20}=p_{10}^2-2\det X,
 \qquad
 p_{02}=p_{01}^2-2\det Y,
\]
and
\[
 p_{11}^2-p_{20}p_{02}
 =-\bigl(x_1y_2-x_2y_1\bigr)^2.
\]
The stable generators remain independent until a relation forced by the rank is reached.
\end{proposition}
\begin{proof}
Write the joint eigenvalues as $(x_1,y_1)$ and $(x_2,y_2)$ and expand both sides.
\end{proof}

\section{Stable passage}
Adding the origin defines the rank transition.
Theorem~\ref{thm:vol2-stable-multisymmetric-ring} identifies the
stable graded algebra and its total-degree completion.
Proposition~\ref{prop:vol2-rank-poisson} retains the central rank
and identifies the stable Poisson sector of degree at least two.
A quotient map from a derived centre requires a representation
and trace comparison on the specified complexes; the stable
ring calculation alone does not construct that map.


\part{Global holomorphic symplectic descent}

\chapter{Fedosov descent}
\section{The Weyl algebroid}
All Hochschild cochains and Morita centres in this chapter are continuous
and relative to $R=\C\llbracket\hbarq\rrbracket$.
Let $(S,\omega)$ be a holomorphic symplectic surface.
Choose a formal Poisson tensor
$\pi_\hbar=\hbarq\omega^{-1}+O(\hbarq^2)$ and a Weyl algebroid
$\mathcal W_{S,\hbar}$, with specified holomorphic or Dolbeault
Fedosov resolutions and coherent overlap maps.
The formal Poisson class is part of these data.
It is not determined by the finite Hamiltonian Lie-cochain comparison.

\begin{theorem}[Global comparison with formality and descent data]
Suppose the resolutions carry the coherent complete
homotopy-$\Etwo$ comparison of
Theorem~\ref{thm:vol2-fedosov-marked-descent}, with
quasi-isomorphic associated graded maps. Then there is an equivalence
\[
 \Phi_S:
 \left(\PV_S\llbracket\hbarq\rrbracket,[\pi_\hbar,-]\right)
 \xrightarrow{\sim}
 C^\bullet_{\HH}(\mathcal W_{S,\hbar},\mathcal W_{S,\hbar}),
\]
and the corresponding equivalence on derived global sections.
A compatible gauge equivalence of the supplied resolutions
transports this comparison. Equivariance requires equivariant
comparison and descent data.
\end{theorem}
\begin{proof}
The complete filtered comparison is an equivalence on every
multiple intersection by Lemma~\ref{lem:vol2-complete-filtered-comparison}.
The homotopy limit of its descent diagram preserves that equivalence.
Apply derived global sections and the supplied gauge or equivariance
homotopies.
\end{proof}

\begin{corollary}[Global Morita centre with a continuous model]
Assume in addition that the chosen continuous Morita theory computes
the centre of $\Perf(\mathcal W_{S,\hbar})$ by derived global
sections of these relative Hochschild cochains. Then
\[
 Z_{\mathrm{Mor}}\bigl(\Perf(\mathcal W_{S,\hbar})\bigr)
 \simeq
 \R\Gamma\!
 \left(S,\PV_S\llbracket\hbarq\rrbracket,[\pi_\hbar,-]\right).
\]
This applies to a K3 surface when the same comparison data are supplied.
A Hamiltonian BF identification additionally requires a map from its
Chevalley observables to this symplectic polyvector complex.
\end{corollary}
\begin{proof}
Use the stated Morita model and the preceding derived global equivalence.
\end{proof}

\chapter{Factorization and descent}
\section{The algebraic mixed-HT realization}
Retain the complete homotopy-$\Etwo$ sheaf
\[
 \mathcal P_S=
 \left(\PV_S\llbracket\hbarq\rrbracket,[\pi_\hbar,-]\right).
\]
For each holomorphic open, its factorization realization uses two
framed topological directions. Suppose these realizations have the
cosheaf descent needed in the holomorphic variable.
They then define an algebraic mixed factorization theory
\[
 \Acal^{\mathrm{alg}}_{S,\mathrm{bulk}}
\quad\text{on}\quad
 \R^2_{\mathrm{top}}\times S.
\]
The supplied formality comparison realizes the corresponding
equivalence with the relative Hochschild centre.
An oriented realization requires rotation-equivariant structure.
Identifying $\Perf(\mathcal W_{S,\hbar})$ with a universal physical
boundary category requires the independent module, action and
generation comparisons of Chapter~\ref{ch:vol2-marked-reconstruction}.
Neither the factorization realization nor the scalar centre
calculation supplies them.

\section{BV realization}
The cyclic Lie sheaf of formal Hamiltonians has the shifted-cotangent action
\[
 S_{\mathrm{BF}}(A,B)
 =\int_{\R^2\times S}
 \left\langle B,(\dd+\dbar)A+\frac12[A,A]\right\rangle.
\]
Here $S$ is a complex symplectic surface, so $\mathbb R^2\times S$ has six real dimensions. A comparison of the classical observables with $\Acal^{\mathrm{alg}}_{S,\mathrm{bulk}}$ must specify a map from Hamiltonian Chevalley cochains to the symplectic polyvector complex. The finite Lie algebra identification does not supply it. Gauge fixing, extensions of graph weights and a solution of the quantum master equation would construct a perturbative factorization theory. Identifying its bulk--boundary action with the algebraic map additionally requires a comparison preserving support, units and collision homotopies. The finite jet zero modes carry the exact quantum measure constructed above. The three-dimensional half-space construction and its interacting comparison problem are treated separately in Chapter~\ref{ch:vol2-native-bulk-boundary}.

\part{Matrix branes and trace geometry}

\chapter{The derived commuting stack}
Let $X,Y\in\mathfrak{gl}_N$ and let $\psi\in\mathfrak{gl}_N[1]$.  The Koszul differential
\[
 QX=QY=0,
 \qquad
 Q\psi=[X,Y]
\]
resolves the commuting equation.  The simultaneous conjugation action gives the boundary cdga
\[
 A_N=C^*_{\CE}\!\left(
 \mathfrak{gl}_N,
 \Sym(\mathfrak{gl}_N\oplus\mathfrak{gl}_N\oplus\mathfrak{gl}_N[1])
 \right).
\]

Cyclic differentiation gives moment-preserving matrix derivations for
cyclic Hamiltonians. A polynomial Hamiltonian requires a cyclic lift,
and these lifts do not give a strict Lie map on all polynomials:
the cubic defect is \eqref{eq:vol2-cubic-trace-defect}.
Theorem~\ref{thm:vol2-all-arity-derived-action} constructs the
coherent action on the free resolution, and
Proposition~\ref{prop:vol2-action-bfv-extension} extends its matrix
realization through the gauge complex. Their higher relation
homotopies are part of the action. A strict continuous representation
of the completed enveloping algebra would require an additional
comparison in that same recipient.

\chapter{Necklace Hamiltonians}
Let
\[
 F=\C\angles{x,y},
 \qquad
 F_{\mathrm{cyc}}=F/[F,F].
\]
For a cyclic word $f$, define cyclic derivatives by cutting the word at every occurrence of the selected letter.  Put
\[
 \delta_f(x)=-\partial_y^{\mathrm{cyc}}f,
 \qquad
 \delta_f(y)=\partial_x^{\mathrm{cyc}}f.
\]

\begin{theorem}[Necklace Lie algebra]
The bracket
\[
 \{f,g\}_{\mathrm{neck}}
 =\left[
 \partial_x^{\mathrm{cyc}}f\,
 \partial_y^{\mathrm{cyc}}g-
 \partial_y^{\mathrm{cyc}}f\,
 \partial_x^{\mathrm{cyc}}g
 \right]_{\mathrm{cyc}}
\]
makes $F_{\mathrm{cyc}}/\C$ a Lie algebra.  The assignment $f\mapsto\delta_f$ is a Lie morphism.  It preserves the moment map:
\[
 \delta_f([x,y])=0.
\]
\end{theorem}
\begin{proof}
The cyclic Euler identity is
\[
 [x,\partial_x^{\mathrm{cyc}}f]+[y,\partial_y^{\mathrm{cyc}}f]=0.
\]
It gives the moment-map statement.  The canonical double Poisson bracket on the free algebra descends to the displayed necklace bracket; its Jacobi identity gives the first assertion.  Evaluation on generators gives $[\delta_f,\delta_g]=\delta_{\{f,g\}_{\mathrm{neck}}}$.
\end{proof}

\begin{theorem}[Trace moment map]
For
\[
 J_N(f)=\Tr f(X,Y),
\]
one has
\[
 \{J_N(f),J_N(g)\}_{\mathcal M_N}
 =J_N\!\left(\{f,g\}_{\mathrm{neck}}\right),
\]
where $\mathcal M_N=\mathfrak{gl}_N^2$ carries $\Tr(\dd X\wedge\dd Y)$.
\end{theorem}
\begin{proof}
The differential of a trace is the matrix evaluation of the cyclic derivative.  Contracting the two differentials with the inverse trace symplectic form gives the necklace formula.
\end{proof}

\chapter{The two-brane cone}
Separate center-of-mass and relative coordinates.  Put
\[
 u=x_1-x_2,
 \qquad
 v=y_1-y_2,
 \qquad
 E=\frac12u^2,
 \quad
 F=-\frac12v^2,
 \quad
 H=-uv.
\]
Then
\[
 H^2+4EF=0.
\]

\begin{proposition}[The $A_1$ Kirillov--Kostant cone]
With $\{u,v\}=2$,
\[
 \{H,E\}=4E,
 \qquad
 \{H,F\}=-4F,
 \qquad
 \{E,F\}=2H.
\]
After the evident rescaling this is the $\mathfrak{sl}_2$ Lie--Poisson bracket on the nilpotent cone.
\end{proposition}
\begin{proof}
The three identities follow by direct differentiation.  The quadratic equation is the determinant-zero equation for the traceless matrix
\[
 \begin{pmatrix}H&2E\\2F&-H\end{pmatrix}.
\]
\end{proof}

\part{Wheel corrections and exact quantum structure}

\chapter{The local centre with comparison data}
Assume a compatible jet comparison and a complete homotopy-$\Etwo$
refinement as in Theorem~\ref{thm:vol2-compatible-jets}.
Put
\[
 \mathcal A=\widehat U_\hbar(\h),\qquad
 \Acal_{\mathrm{bulk}}^{\mathrm{alg}}
 =\operatorname{Fact}_{\Etwo}(\mathcal P_\infty).
\]
The factorization envelope uses two framed topological directions.
An oriented refinement requires the corresponding rotation-equivariant structure.
The formal holomorphic coefficient remains at the Darboux point.
Let $\Ccal_\partial$ be the category of perfect complete $\mathcal A$-modules
in the chosen continuous $R$-linear Morita theory.
Assume its centre is computed by
$C^\bullet_{\HH,\mathrm{cont}}(\mathcal A/R,\mathcal A)$.

\begin{theorem}[Local derived-centre comparison]
These data give an equivalence of complete homotopy-$\Etwo$ factorization algebras
\[
 \Acal_{\mathrm{bulk}}^{\mathrm{alg}}
 \xrightarrow{\sim}
 Z^{\mathrm{der}}_{\mathrm{Mor},R}(\Ccal_\partial).
\]
At the formal insertion point the map is the selected refinement of $\Phi_\infty$.
Every continuous boundary module has the induced derived centre action.
\end{theorem}
\begin{proof}
The compatible jet theorem gives the complete homotopy-$\Etwo$ comparison.
The factorization envelope preserves this equivalence.
The assumed continuous Morita model identifies its target with the centre.
A centre element is a derived natural endomorphism of the identity functor,
so evaluation gives the action on each boundary module.
\end{proof}

\begin{corollary}[Global K3 comparison with descent data]
Let $S$ be a K3 surface with holomorphic symplectic form.
Suppose the local comparisons extend to a sheafwise equivalence with the
relative Hochschild centre of $\Perf(\mathcal W_{S,\hbar})$.
Suppose further that the selected cyclic and mixed refinements extend
and satisfy descent in the same completed category.
Then derived global sections give the corresponding global centre,
cyclic, and open--closed comparisons.
\end{corollary}
\begin{proof}
Derived global sections preserve equivalences of sheaves.
Apply them to the supplied sheafwise maps and their coherence homotopies.
\end{proof}

\chapter{Synthesis}
At each fixed Hamiltonian jet, the Rees enveloping algebra has relative
Hochschild complex equivalent to the Chevalley complex with differential
$\hbarq d_{\CE}$.
PBW antisymmetrization proves this over every $R_N$ and in the formal coefficient limit.
The Duflo theorem gives the separate multiplicative comparison on invariant functions.
Quadratic Moyal Hamiltonians generate the oscillator quotient of the enveloping algebra.
Their associative Casimir relation prevents the primitive coproduct from descending.

The Hamiltonian jet limit requires the compatible cochain systems of
Theorem~\ref{thm:vol2-compatible-jets}.
With its operadic refinement and continuous Morita model, the centre acts
on every supplied continuous boundary module $M$ through
\[
 \mathcal P_\infty
 \xrightarrow{\ \Phi_\infty\ }
 C^\bullet_{\HH,\mathrm{cont}}(\mathcal A/R,\mathcal A)
 \longrightarrow\RHom_{\mathcal A}(M,M).
\]
The last arrow evaluates a derived natural endomorphism at $M$.
A matrix-brane realization must construct such a module and compare its
centre action with the selected trace functions.
A commuting square with scalar representation invariants requires these
additional comparison maps.
Likewise, Fedosov descent concerns symplectic polyvectors and a Weyl algebroid.
Its relation to Hamiltonian Lie cochains requires the specified map between their carriers.

\part{Technical appendices to the preceding book}
\chapter{Bibliographic note}
The finite affine formality map and the separate symplectic globalization use Kontsevich formality and its Fedosov resolution \cite{KontsevichDQ,DolgushevCov,CalaqueDolgushevHalbout}.  The relative Hochschild model is the Cartan--Eilenberg comparison, with both Chevalley terms scaled by $\hbarq$.  Its finite-base resolution proof is Theorem~\ref{thm:vol2-relative-rees}.  Jet limits and operadic refinements require the additional data stated above.  The all-arity centre action is the Deligne--Tamarkin brace action \cite{McClureSmith,KontsevichSoibelmanDeligne}.  The global deformation-quantization category is the holomorphic Weyl algebroid.  The matrix trace and necklace formulas are the noncommutative symplectic moment-map construction.

\cleardoublepage
\phantomsection
\addcontentsline{toc}{part}{Ordered and Associative Chiral Geometry}
\chapter*{Book entrance: Ordered and Associative Chiral Geometry}
\addcontentsline{toc}{chapter}{Book entrance: Ordered and Associative Chiral Geometry}
\begin{summarybox}
Collision on a complex curve, order on a transverse line, exchange in a plane, and the three corresponding bar constructions.
\end{summarybox}
\part{The invariant object}

\chapter{Two directions of locality}
\section{The local experiment}
Let $X$ be a smooth complex algebraic curve and let $I\subset\R$ be an oriented interval. Place two protected observables at $(z_1,t_1)$ and $(z_2,t_2)$ in $X\times I$. The limits $z_1\to z_2$ and $t_1\to t_2$ carry independent information. The first limit produces a singular distribution along the diagonal of $X^2$. The second composes operators in the order selected by the orientation of $I$.

\begin{principle}[Two-direction principle]
Holomorphic collision belongs to the factorization coefficient on $X$. Ordered fusion belongs to an internal $\Eone$ multiplication. Exchange belongs to an additional two-dimensional transport structure.
\end{principle}

This separation gives five geometric carriers:
\[
\begin{array}{c|c}
\text{carrier}&\text{operation}\\ \hline
\text{diagonal in }X^2&\text{singular chiral operation}\\
\text{oriented interval}&\text{ordered fusion and bar construction}\\
\text{oriented plane}&\text{braid transport}\\
\text{circle}&\text{Hochschild trace}\\
\text{stable curve}&\text{modular sewing}.
\end{array}
\]
Each later structure is introduced by the datum carried by its row.

\section{Carrier and operation}
For the diagonal embedding $\Delta:X\hookrightarrow X^2$ and its complement $j:X^2\setminus\Delta(X)\hookrightarrow X^2$, a Beilinson--Drinfeld chiral Lie operation is a morphism
\[
 \mu:j_*j^*(\g\boxtimes\g)\longrightarrow \Delta_!\g
\]
in the chiral pseudo-tensor category of right $D$-modules \cite{BD04}. The $D$-module is the carrier. The morphism $\mu$ is the operation. A factorization object packages carriers on all powers $X^I$ together with their disjoint-locus structure maps. The complete operator product is the pair consisting of that system and its multiplication maps.

In a formal coordinate $z$ near the diagonal, the pole filtration of $\Ocal_{X^2}(*\Delta)$ has
\[
 \gr_m^F\Ocal_{X^2}(*\Delta)
 \cong \Ocal_{\Delta}(m\Delta)
 \cong N_{\Delta/X^2}^{\tensor m},\qquad m\ge 1.
\]
A principal part of order $m$ therefore has an intrinsic normal symbol. Coordinate changes act triangularly on the filtered object. Residues recover the first symbol; normal jets recover the higher symbols.

\chapter{The Ran coefficient category}
\section{The fixed six-functor setting}
Fix a presentable stable category $\Dcal_X$ of right $D$-modules on $\Ran(X)$ with the following operations:
\begin{enumerate}
\item the pointwise tensor product $\tensor^{!}$;
\item chiral convolution $\tensor^{\ch}$ along disjoint union;
\item base change and projection formulae for the disjoint-union correspondences;
\item geometric realizations and totalizations preserved by the tensor products used below.
\end{enumerate}
The standard renormalized $D$-module setting of chiral Koszul duality has these operations \cite{FG12}. Every construction below is formed in this fixed setting.

Let
\[
  \Ran(X)\times\Ran(X) \xleftarrow{\ j\ }
  (\Ran(X)\times\Ran(X))_{\mathrm{disj}}
  \xrightarrow{\ u\ }\Ran(X)
\]
be the disjoint-union correspondence. Define
\[
 F\tensor^{\ch}G=u_!j^!(F\boxtimes G).
\]
The pointwise product is induced by the diagonal of $\Ran(X)$.

\begin{definition}[Aligned-decomposition projector]
For a finite set $I$, let $\Dec(I)$ be the finite groupoid of ordered decompositions $I=A\sqcup B$ and let
\[
 a_I:\Dec(I)\longrightarrow\Dec(I)\times\Dec(I),\qquad d\longmapsto(d,d)
\]
be the diagonal.  In the linear species model define
\[
 e_I=a_{I,!}a_I^!.
\]
Thus $e_I$ projects a pair of decompositions onto the summands in which the two decompositions agree.  In the six-functor Ran model the same formula denotes extraordinary pullback to the diagonal correspondence followed by the available renormalized pushforward.
\end{definition}

\begin{theorem}[Aligned sector in the finite decomposition model]
Let $\mathcal S_I$ be a stable category indexed by ordered decompositions of a finite set $I$, and assume the direct-sum decomposition
\[
 \mathcal M_I\simeq\bigoplus_{d_1,d_2\in\Dec(I)}\mathcal M_{d_1,d_2}.
\]
The projection
\[
 e_I:\mathcal M_I\longrightarrow\mathcal M_I,
 \qquad
 e_I\big|_{\mathcal M_{d_1,d_2}}=
 \begin{cases}\id,&d_1=d_2,\\0,&d_1\ne d_2,
 \end{cases}
\]
is idempotent.  Its image is the aligned sector
\[
 \mathcal M_I^{\mathrm{al}}
 =\bigoplus_{d\in\Dec(I)}\mathcal M_{d,d}.
\]
The family $e_I$ is natural under bijections and compatible with iterated disjoint union.  A mixed distributive law in this model is a natural transformation defined on $\mathcal M_I^{\mathrm{al}}$ whose two restrictions to the triple diagonal $d_1=d_2=d_3$ agree.
\end{theorem}
\begin{proof}
Idempotence is checked on each direct summand.  Naturality follows because a bijection carries an ordered decomposition to an ordered decomposition and preserves equality of the two indices.  For three products, both iterated alignments project
\[
 \bigoplus_{d_1,d_2,d_3}\mathcal M_{d_1,d_2,d_3}
 \quad\text{onto}\quad
 \bigoplus_d\mathcal M_{d,d,d}.
\]
The two projectors are therefore equal.  This equality is precisely the triple-alignment coherence.
\end{proof}

\begin{theorem}[Descent of an aligned projector]
Suppose the common-decomposition inclusions are clopen on the chosen charts.
Suppose every transition correspondence gives Cartesian squares with these
inclusions and the corresponding Beck--Chevalley maps are isomorphisms.
Then the functors
\[
 e_I=a_{I,!}a_I^!
\]
assemble to a natural idempotent $e_{\mathrm{al}}$ on the mixed correspondence object. Its image is the common-decomposition correspondence. The two iterated projectors on three decomposition indices coincide, and the assumed base-change maps transport this equality through the transitions.
\end{theorem}
\begin{proof}
Restriction and extension by zero along a clopen inclusion form a projector.
The Cartesian-square hypothesis and base change identify these composites
across transitions. For three indices, both iterated projectors select the
same locus $d_1=d_2=d_3$, so the identification respects triple alignment.
\end{proof}

The Cartesian-square hypothesis must be checked on the actual correspondences.
Preserving equal indices alone does not suffice: a map can also send distinct
indices to equal ones, as the map $\{0,1\}\to\{*\}$ does.
On a common disjoint-union correspondence, the projection formula gives the
aligned comparison. Its inclusion and retraction are inverse on the full
object exactly when every off-aligned summand vanishes.
This criterion concerns that canonical comparison. It does not exclude every
noninvertible distributive law.

\section{Factorization coalgebras}
A chiral commutative coalgebra is a commutative coalgebra object for $\tensor^{\ch}$. It is factorizing when every coproduct becomes an equivalence on the relevant disjoint locus. Let
\[
 \Fact_D(X)\subset \operatorname{CoCAlg}_{\ch}(\Dcal_X)
\]
be the full subcategory of factorizing objects. Chiral Koszul duality identifies the corresponding Lie-theoretic and factorization descriptions on the completed Ran category \cite{FG12,Heuts}.

\begin{proposition}[Pointwise closure]
The componentwise tensor product of two factorization objects is a factorization object. Thus $\Fact_D(X)$ is symmetric monoidal under $\tensor^!$ whenever the fixed Ran enhancement is closed under that tensor product.
\end{proposition}
\begin{proof}
For a decomposition $I=I_1\sqcup I_2$, restrict to its disjointness locus. The factorization equivalences of $F$ and $G$ give
\[
\begin{aligned}
 j^!(F_I\tensor^!G_I)
 &\simeq j^!F_I\tensor^!j^!G_I\\
 &\simeq (F_{I_1}\boxtimes F_{I_2})\tensor^!
          (G_{I_1}\boxtimes G_{I_2})\\
 &\simeq (F_{I_1}\tensor^!G_{I_1})\boxtimes
          (F_{I_2}\tensor^!G_{I_2}).
\end{aligned}
\]
Compatibility under refinement and permutation follows componentwise from the corresponding compatibilities for $F$ and $G$.
\end{proof}


\chapter{The two noncommutative chiral objects}
\section{Transverse order}
Two operators at the same chiral support can compose in an independent
oriented real direction. Their multiplication is internal to the pointwise
coefficient category. This is the structure carried by a boundary or line
defect in a holomorphic--topological theory on $X\times\R$.

\begin{definition}[Transversely ordered chiral algebra]
A transversely ordered chiral algebra is an augmented object
\[
 A^\perp\in\Alg_{\Eone}^{\mathrm{aug}}(\Fact_D(X),\tensor^!).
\]
The superscript $\perp$ records the transverse oriented direction.  The shorter phrase \emph{ordered chiral algebra} will always mean this object.
\end{definition}

\section{Associativity inside chiral convolution}
Chiral convolution combines disjoint finite subsets before allowing them to collide.  Associativity for this product is intrinsic to the curve and has a different source and a different bar construction.

\begin{definition}[Associative chiral algebra]
An associative chiral algebra is an augmented algebra
\[
 A^{\ch}_{\mathsf{Ass}}
 \in\Alg_{\Eone}^{\mathrm{aug}}(\Dcal_X,\starCh).
\]
A chiral Lie algebra is a Lie algebra object
\[
 \g^{\ch}\in\Alg_{\mathsf{Lie}}(\Dcal_X,\starCh).
\]
The commutator of an associative chiral algebra is a chiral Lie algebra.
\end{definition}

The operation $m^\perp:A\tensor^!A\to A$ composes transverse intervals.  The operation
$m^{\ch}:A\starCh A\to A$ composes chiral configurations.  They have different domains, different units, and different deformation complexes.

\begin{theorem}[Independence of the two associative structures]
Let $M$ be an object for which the split square-zero extensions below exist.  A pointwise associative product $\mu_\perp:M\tensor^!M\to M$ defines a transversely ordered algebra on $\mathbbm1_\perp\oplus M$, while the reduced chiral-convolution product may independently be zero.  A chiral-convolution product $\mu_{\ch}:M\starCh M\to M$ defines an associative chiral algebra on $\mathbbm1_{\ch}\oplus M$, while the reduced pointwise product may independently be zero.  Hence neither associative structure determines the other.

A common carrier equipped with both products is a mixed object after a mixed distributive datum has been supplied and its unit, binary, and triple coherences have been verified.
\end{theorem}
\begin{proof}
Adjoining the unit to an associative product on the augmentation ideal gives a unital split extension.  The zero product is associative in either monoidal category.  These two constructions vary independently.  On a common carrier, compatibility is additional structure because the mixed composites have different a priori factorizations through the two monoidal products.
\end{proof}

\begin{definition}[Mixed associative chiral datum]
A mixed associative chiral datum is a factorization object $A$ equipped with
\begin{enumerate}
\item a pointwise multiplication $m^\perp:A\tensor^!A\to A$ and a unit map $\eta_\perp:\mathbbm1_\perp\to A$;
\item a chiral-convolution multiplication $m^{\ch}:A\starCh A\to A$ and a unit map $\eta_{\ch}:\mathbbm1_{\ch}\to A$;
\item an aligned mixed distributive law identifying the two composites on $\operatorname{im}(e)$.
\end{enumerate}
The coherence diagrams are the two color-$3$ faces of the mixed forest complex.  A bare symbol $\mathbbm1$ below denotes the unit of the monoidal category in which the displayed construction is formed.
\end{definition}

\section{Formal-disc realization}
Fix a formal coordinate on a disc.  Translation-equivariance turns a chiral-convolution product into a state--field operation $Y(a,z)b$ with coefficients in Laurent series.  The following theorem packages the exact relationship with the vertex-algebra language.

\begin{theorem}[Formal-disc realizations]
Let $V$ be a translation-equivariant state--field system with vacuum and injective state--field map.
\begin{enumerate}
\item If its fields are weakly associative in the sense of Li, then $V$ is a nonlocal vertex algebra.
\item If, in addition, $V$ is topologically free over $\kfield[[\hbar]]$, carries a unitary rational quantum Yang--Baxter operator $S(z)$, and satisfies $S$-locality and the hexagon compatibility, then $V$ is a quantum vertex algebra.
\item A translation-equivariant associative chiral-convolution algebra produces such a state--field system whenever its formal-disc restriction satisfies the reconstruction, lower-truncation, and weak-associativity hypotheses.  Conversely, a nonlocal or quantum vertex algebra produces the corresponding formal-disc factorization system by its products of fields.
\end{enumerate}
\end{theorem}
\begin{proof}
The first assertion is Li's reconstruction theorem for quasi-compatible fields \cite{LiNonlocal}.  The second is the quantum-vertex reconstruction theorem: unitarity gives inverse exchange, the quantum Yang--Baxter equation compares the two three-field exchange paths, and the hexagon relates exchange to the vertex operation \cite{LiQVA,DSGKQVA}.  For the third assertion, the formal-disc chiral product supplies the iterated field products; the stated analytic-formal hypotheses are exactly those used by the reconstruction theorems.  The converse assigns to a configuration of formal points the iterated field product and uses weak associativity, or braided associativity, for factorization under refinement.
\end{proof}

\section{Three separation laws}
The complete principal-part filtration, transverse order, and planar exchange form a strict hierarchy of supplied data.  The $m$th normal symbol of a pole is recovered by the $m$th normal-jet map.  An oriented line supplies contractible order chambers and hence $\Eone$ fusion.  A plane supplies braid paths and hence exchange.  These positive constructions locate each operation in the first geometry that carries it.

\chapter{Transversely ordered chiral algebras}
\begin{definition}[Ordered chiral algebra]
An ordered chiral algebra on $X$ is an augmented $\Eone$-algebra object
\[
 A\in \Alg_{\Eone}^{\mathrm{aug}}\bigl(\Fact_D(X),\tensor^!\bigr).
\]
Equivalently, $A$ is a factorization object equipped with multiplication maps
\[
 m_I:A_I\tensor^!A_I\longrightarrow A_I
\]
and a unit, compatible with all diagonal transition maps and all disjoint-locus factorization equivalences.
\end{definition}

For every decomposition $I=I_1\sqcup I_2$, locality is the componentwise identity
\[
 \Delta_{I_1,I_2}\circ m_I
 = (m_{I_1}\boxtimes m_{I_2})\circ\sigma\circ
   (\Delta_{I_1,I_2}\tensor^!\Delta_{I_1,I_2}),
\]
where $\sigma$ reorders the middle two factors. This equation says that ordered fusion commutes with separation of chiral supports.

\begin{theorem}[Locality recognition]
To give an ordered chiral algebra is equivalent to giving a factorization family $A_I$, an associative unital multiplication $m_I$ on every component, and the displayed locality identities, compatibly with refinements, permutations, and diagonal transitions.
\end{theorem}
\begin{proof}
An $\Eone$-algebra in the symmetric monoidal category $\Fact_D(X)$ is, by definition, an associative multiplication and unit whose structure maps are morphisms in $\Fact_D(X)$. A morphism in $\Fact_D(X)$ is exactly a componentwise morphism compatible with the transition and factorization maps. Expanding this compatibility for a binary decomposition gives the displayed identity; the higher decompositions follow by iteration.
\end{proof}

\section{Product field theories}
Let $\Fact_{\mathrm{hol,lc}}(X^{\mathrm{an}}\times\R)$ denote factorization algebras that are holomorphic along $X$ and locally constant along $\R$. Product additivity gives
\[
 \Fact_{\mathrm{hol,lc}}(X^{\mathrm{an}}\times\R)
 \simeq
 \Alg_{\Eone}\bigl(\HolFact(X^{\mathrm{an}})\bigr)
\]
for product-generated objects. Algebraization by regular holonomic Riemann--Hilbert then gives ordered chiral algebras in $\Fact_D(X)$ whenever the coefficient system is regular holonomic and the comparison preserves the required colimits.

\begin{theorem}[Holomorphic--topological recognition]
Let $\F$ be a product-generated factorization algebra on $X^{\mathrm{an}}\times\R$ that is holomorphic on $X$ and locally constant on $\R$. Then the restriction of $\F$ to a transverse point is an $\Eone$-algebra in holomorphic factorization algebras on $X$. A regular-holonomic algebraization produces an ordered chiral algebra on $X$.
\end{theorem}
\begin{proof}
An ordered family of disjoint intervals in $\R$ acts on the restriction by factorization. Contractibility of each ordered configuration chamber gives the coherent $\Eone$ structure. Factorization in the $X$ direction is preserved because product embeddings commute. Riemann--Hilbert transports the structure in the regular-holonomic subcategory.
\end{proof}

\chapter{Block-ordered operations}
The preceding definition is intrinsic. A geometric model arises after one chooses an infinity-operad $\mathsf{Ch}_X$ representing the selected chiral factorization theory. The Boardman--Vogt tensor product
\[
 \mathsf{OrdCh}_X=\Eone\tensor\mathsf{Ch}_X
\]
then represents compatible ordered and chiral operations.

\begin{theorem}[Operadic exponential law]
Let $P$ and $Q$ be infinity-operads whose Boardman--Vogt tensor product exists, and let $\Vcal$ be a compatible presentable symmetric monoidal infinity-category.  Then
\[
 \Alg_{P\tensor Q}(\Vcal)
 \simeq
 \Alg_P\bigl(\Alg_Q(\Vcal)\bigr).
\]
In particular, after a chiral operad $\mathsf{Ch}_X$ has been constructed, the tensor product $\Eone\tensor\mathsf{Ch}_X$ classifies $\Eone$-algebra objects in $\mathsf{Ch}_X$-algebras.
\end{theorem}
\begin{proof}
The tensor product of infinity-operads is characterized by the displayed universal property \cite{LurieHA}.  The second statement is the specialization $P=\Eone$ and $Q=\mathsf{Ch}_X$.  This formal equivalence applies to each geometric model of $\mathsf{Ch}_X$ after that model and its operadic tensor product have been constructed.
\end{proof}

A profile is an ordered list $(I_1,\ldots,I_r)$ of finite chiral label sets. Its geometric carrier is $X^{I_1}\times\cdots\times X^{I_r}$. Chiral morphisms refine and collide labels inside one block. Transverse morphisms merge adjacent blocks. The tensor-product operad imposes their mixed squares.

\section{Mixed forests}
Fix a cofibrant dg-operad $Q_X\to\mathsf{Ch}_X$. Apply the functorial $W$-construction to $\Eone\tensor Q_X$. A cell is represented by a levelled forest with transverse and chiral internal edges. Put
\[
 \operatorname{or}(F)=\det(E_\perp(F))\tensor\det(E_{\ch}(F)).
\]
The differential is the sum of the screen differential and one-edge contractions.

\begin{theorem}[Two-color forest differential]
Let $C_*$ be the chain complex freely generated by oriented two-color forests, with one determinant-line generator for every transverse edge and one for every chiral edge.  Suppose contraction of an edge is functorial on the coefficient system and every codimension-two contraction square commutes before orientations are inserted.  With the product orientation, the transverse and chiral contraction operators satisfy
\[
 d_\perp^2=0,\qquad d_{\ch}^2=0,\qquad
 d_\perp d_{\ch}+d_{\ch}d_\perp=0.
\]
Thus $d=d_{\mathrm{int}}+d_\perp+d_{\ch}$ squares to zero whenever the internal differential anticommutes with both contraction operators.
\end{theorem}
\begin{proof}
Every term of $d_\perp^2$ contracts the same pair of transverse edges in two orders.  The underlying coefficient maps agree by functoriality, while exchanging the two transverse determinant generators reverses the sign.  The terms cancel in pairs.  The same argument proves $d_{\ch}^2=0$.  For one edge of each color the coefficient maps again form a commuting square and exchanging the two determinant generators contributes the Koszul sign $-1$, proving the mixed anticommutator identity.  The internal terms vanish by the chain-map assumption.
\end{proof}

This construction supplies a geometric model whenever a concrete $Q_X$ and its screen coefficients are available. The intrinsic definition of ordered chiral algebra remains independent of that choice.

\part{Bar, cobar, and the boundary state}


\chapter{Three bar constructions}
The phrase \emph{the chiral bar} denotes three functors once its monoidal product and algebraic operad have been named.

\begin{definition}[Transverse, associative-chiral, and Lie-chiral bars]
For augmented objects of the indicated types define
\[
 \BarPerp(A^\perp)
   :=\mathbbm1\tensor_{A^\perp}^{\mathbb L,!}\mathbbm1,
\]
\[
 \BarAssCh(A^{\ch}_{\mathsf{Ass}})
   :=\mathbbm1\starCh_{A^{\ch}_{\mathsf{Ass}}}^{\mathbb L}\mathbbm1,
\]
and
\[
 \BarLieCh(\g^{\ch})
   :=C_*^{\ch}(\g^{\ch})
   =\Sym^{\ch}(\g^{\ch}[1])
\]
with its chiral Chevalley differential.  The first bar is formed with $\tensor^!$, the second with $\starCh$, and the third is the Lie-to-commutative-coalgebra Koszul functor.
\end{definition}

\begin{theorem}[Typed bar theorem]
The three bars have the following chain models.
\begin{enumerate}
\item $\BarPerp(A^\perp)$ is the tensor coalgebra $T^c_{\tensor^!}(s\bar A)$ with coderivation induced by $m^\perp$.
\item $\BarAssCh(A^{\ch}_{\mathsf{Ass}})$ is the tensor coalgebra $T^c_{\starCh}(s\bar A)$ with coderivation induced by $m^{\ch}$.
\item $\BarLieCh(\g^{\ch})$ is the chiral symmetric coalgebra with differential induced by the chiral Lie bracket.
\end{enumerate}
Each differential squares to zero by the defining relation of its own operad.  The outputs live in different coalgebra categories and are compared by a displayed monoidal or PBW morphism.
\end{theorem}
\begin{proof}
Normalize the two simplicial two-sided bars in their respective monoidal categories.  Degeneracies remove units, leaving the two tensor coalgebras, and faces contract adjacent factors using the relevant multiplication.  The Lie formula is the ordinary operadic bar of the Lie operad internal to chiral convolution.  Associativity and Jacobi are exactly the quadratic identities that make the corresponding coderivations square to zero.
\end{proof}

\begin{definition}[PBW-admissible chiral Lie algebra]
A chiral Lie algebra $\g^{\ch}$ is PBW-admissible when its chiral enveloping algebra $U^{\ch}(\g)$ has a complete separated PBW filtration, the canonical map
\[
 \Sym^{\ch}(\g)\longrightarrow\gr_{\mathrm{PBW}}U^{\ch}(\g)
\]
is an equivalence, and the bar filtration is strongly convergent in the chosen nilcomplete Ran category.
\end{definition}

\begin{theorem}[Filtered chiral PBW comparison]
Let $\g$ be a chiral Lie algebra in a complete filtered chiral category.  Assume:
\begin{enumerate}
\item its chiral enveloping algebra $U^{\ch}(\g)$ exists;
\item the PBW filtration is exhaustive, complete, and separated;
\item the chiral PBW morphism
\[
 \Sym^{\ch}(\g)\longrightarrow \gr U^{\ch}(\g)
\]
is an isomorphism;
\item the induced filtrations on the two bar complexes are strongly convergent.
\end{enumerate}
Then antisymmetrization induces a filtered quasi-isomorphism
\[
 \BarAssCh\bigl(U^{\ch}(\g)\bigr)
 \simeq
 \BarLieCh(\g).
\]
\end{theorem}
\begin{proof}
The associated graded of the associative chiral bar is the Koszul bar of $\Sym^{\ch}(\g)$.  The usual Koszul antisymmetrization identifies this coalgebra with the chiral Chevalley coalgebra.  Hypotheses (2)--(4) give a complete, separated, strongly convergent comparison spectral sequence.  Since its associated-graded map is a quasi-isomorphism, the filtered map is a quasi-isomorphism \cite{BD04,FG12}.
\end{proof}

\begin{definition}[Bar-admissible mixed datum]
A mixed associative chiral datum is bar-admissible when its aligned distributive law acts on the full bisimplicial two-sided bar and both geometric realizations commute with the tensor products occurring in that bisimplicial object.
\end{definition}

\begin{theorem}[Mixed Fubini theorem]
Let $A$ be a mixed algebra in a presentable stable category.  Assume the mixed distributive datum extends to the bisimplicial two-sided bar, and that both monoidal products preserve its geometric realizations separately.  Then the two iterated realizations and the diagonal realization are canonically equivalent:
\[
 \left|\,p\mapsto\left|q\mapsto B_{p,q}(A)\right|_{\ch}\right|_\perp
 \simeq
 \left|\operatorname{diag}B_{\bullet,\bullet}(A)\right|
 \simeq
 \left|\,q\mapsto\left|p\mapsto B_{p,q}(A)\right|_\perp\right|_{\ch}.
\]
\end{theorem}
\begin{proof}
The distributive datum makes the two simplicial directions commute, so $B_{\bullet,\bullet}(A)$ is a genuine bisimplicial object.  Geometric realization is a colimit.  Since the two tensor products preserve the relevant colimits, Fubini for colimits identifies both iterated realizations with the realization of the diagonal simplicial object.
\end{proof}

The notation used in the remainder of this part is
\[
 \Barc_X^{\ord}=\BarPerp.
\]
It always refers to transverse order.

\section{Realization and factorization}

Factorization compares the coefficient at a disjoint configuration with
an external tensor product of its coefficients at the separate configurations.
This condition is nonlinear. For example, suppose those coefficients are
$k$ at one point and at two distinct points, where $k$ is a field.
The componentwise sum of two copies has coefficient $k^2$ at both configurations.
Factorization would instead require $k^2\otimes_k k^2\cong k^4$ at two points.
A componentwise shift also fails: its two-point coefficient is $k[1]$,
whereas the tensor product of its singleton coefficients is $k[2]$.
Thus an ambient stable coefficient category does not make its
factorization objects closed under componentwise sums or shifts.

Geometric realization has a different colimit property.
A small indexing $\infty$-category $J$ is \emph{sifted} when it is nonempty
and its diagonal $J\to J\times J$ is cofinal.
Here cofinal means that, for every diagram $D$ on $J\times J$,
the canonical map $\operatorname*{colim}_{J}(D\circ\Delta)
\to\operatorname*{colim}_{J\times J}D$ is an equivalence whenever these colimits exist.
Iteration gives the same property for $J\to J^m$, for every finite $m\geq1$.
The simplex category with arrows reversed is sifted, so simplicial
geometric realization is a sifted colimit.

\begin{proposition}
Fix coefficient categories on configuration charts, with specified
restriction and collision functors and their coherent identifications.
Suppose the coefficient categories admit $J$-indexed colimits
and these functors preserve them.
Suppose external tensor products preserve colimits separately in each variable.
Let $(F_j)_{j\in J}$ be a diagram of factorization objects and
morphisms preserving their factorization equivalences.
If $J$ is sifted, the componentwise colimit is a factorization object.
\end{proposition}
\begin{proof}
Write $F_I=\operatorname*{colim}_{j\in J}(F_j)_I$ on a chart indexed by $I$.
The collision maps and their coherences pass to the colimit by hypothesis.
For a decomposition $I=I_1\sqcup\cdots\sqcup I_m$, let $u$ denote
restriction to its disjointness locus.
Restriction, factorization, cofinality, and separate preservation give
\[
\begin{aligned}
 uF_I
 &\simeq\operatorname*{colim}_{j\in J}u(F_j)_I\\
 &\simeq\operatorname*{colim}_{j\in J}
        \mathop{\boxtimes}_{a=1}^m(F_j)_{I_a}\\
 &\simeq\operatorname*{colim}_{(j_1,\ldots,j_m)\in J^m}
        \mathop{\boxtimes}_{a=1}^m(F_{j_a})_{I_a}
 \simeq\mathop{\boxtimes}_{a=1}^m F_{I_a}.
\end{aligned}
\]
Each comparison is natural in the diagram and the decomposition.
Refinements and permutations therefore preserve the given coherence diagrams.
If an empty-configuration unit is specified, its constant diagram also
has that unit as colimit, since a sifted category has contractible groupoid completion.
\end{proof}

For a simplicial two-sided bar, this proposition applies when every
simplicial term and every face and degeneracy are already in the
specified factorization category.
It also requires the actual restriction and collision functors to preserve realization.
Normalization may then be computed in an ambient stable category.
Its individual shifted summands need not be factorization objects.
A coproduct on the normalized complex requires its own compatibility
with the factorization equivalences.

\chapter{The internal ordered bar}
Let $A$ be an augmented ordered chiral algebra.
Assume the restriction and collision functors preserve geometric realizations,
and the external and pointwise tensor products preserve them separately.
The two-sided simplicial bar has
\[
 B_r(\mathbbm1,A,\mathbbm1)=
 \mathbbm1\tensor^! A^{\tensor^!r}\tensor^!\mathbbm1.
\]
Every face and degeneracy is a factorization morphism by the algebra and augmentation axioms.
The sifted-colimit proposition therefore constructs its realization
$\Barc_X^{\ord}(A)$ as a factorization object.

\begin{theorem}[Normalized bar]
Choose a symmetric monoidal dg model for the ambient coefficients in which the displayed
tensors compute derived tensors and direct-sum normalization computes realization.
Write $\bar A$ for the augmentation ideal in that model.
The underlying bar complex is
\[
 \Barc_X^{\ord}(A)\simeq
 \bigl(T^c(s\bar A),d_{\mathrm{int}}+b_m\bigr).
\]
Here $s$ decreases cohomological degree by one.
The direct sum and its individual shifted summands are formed in the ambient coefficient category.
Deconcatenation makes this complex a conilpotent coassociative coalgebra there:
\[
 \Delta_{\mathrm{dec}}[a_1|\cdots|a_n]
 =\sum_{i=0}^n[a_1|\cdots|a_i]\tensor[a_{i+1}|\cdots|a_n].
\]
\end{theorem}
\begin{proof}
Normalization removes words containing a unit and leaves the shifted augmentation ideal in every slot.
The faces give the internal differential and adjacent multiplication.
The simplicial identities give a square-zero differential.
Two successive cuts enumerate the same three subwords in either order.
An internal differential or multiplication lies on one side of any surviving cut,
with the tensor sign of its prefix. Thus the differential is a coderivation.
Every element has finite word support, so sufficiently many reduced cuts give zero.
\end{proof}

\section{The bar comparison on separated configurations}
The normalized word decomposition is not itself a decomposition in $\Fact_D(X)$.
Its compatibility with factorization is expressed by a map between whole bar complexes.
For augmented dg algebras $P,Q$ in the chosen model, define
\[
 \operatorname{sh}:BP\tensor BQ\longrightarrow B(P\tensor Q)
\]
by shuffling the letters $p_i\tensor1$ and $1\tensor q_j$,
keeping their respective orders.
The sign is the Koszul sign for the permutation of the suspended letters.
For degree-zero letters,
\[
 \operatorname{sh}([p]\tensor[q])
 =[p\tensor1|1\tensor q]-[1\tensor q|p\tensor1].
\]

\begin{proposition}[Shuffle and factorization]
The shuffle is a natural quasi-isomorphism and a coalgebra map,
where its source has the tensor-product coalgebra structure.
For the geometric coefficient models above, suppose these maps commute with the specified
restriction and collision functors.
Then they give factorization equivalences for the whole bar and make
deconcatenation a morphism of factorization objects in the derived category.
\end{proposition}
\begin{proof}
In the bar differential of a shuffle, a merge of letters from different factors
occurs twice with opposite suspended signs.
The equality $(p\tensor1)(1\tensor q)=(-1)^{|p||q|}(1\tensor q)(p\tensor1)$ cancels these terms.
Merges within one factor give the differential of that factor.
Internal differentials commute with the same signed permutation.
Thus the shuffle is a cochain map.

A cut of a shuffled word determines one cut in each input word,
a shuffle of their prefixes, and a shuffle of their suffixes.
Conversely these four choices determine exactly that cut shuffle.
The signs on both sides come from the same permutation of suspended letters.
This proves the coalgebra identity with the tensor-product coproduct.

To verify the quasi-isomorphism, form the bisimplicial object
$B_p(k,P,k)\tensor B_q(k,Q,k)$, with $k$ denoting the tensor unit.
Its iterated realization is $BP\tensor BQ$.
Its diagonal is the simplicial bar of $P\tensor Q$.
Cofinality of the simplicial diagonal identifies these realizations.
On normalized chains the map from the iterated realization to the diagonal
is the shuffle just defined: the monotone paths from $(0,0)$ to $(p,q)$
triangulate the product of its two simplices.
Each path inserts precisely the degeneracies of its shuffle.
Adjacent paths have opposite orientations on a shared face,
so their signed sum has exactly the product boundary.
This is the realization comparison, hence an equivalence.

For $I=I_1\sqcup I_2$ on the disjointness locus, the given algebra factorization
identifies $A_I$ with $A_{I_1}\boxtimes A_{I_2}$.
The shuffle therefore gives the comparison
$BA_{I_1}\boxtimes BA_{I_2}\to BA_I$.
Its inverse in the derived category is the bar factorization map.
The coalgebra identity just proved gives compatibility with deconcatenation.
For several blocks, iterated shuffles enumerate the same ordered multishuffles,
with the same Koszul signs. This proves refinement coherence.
Naturality and the stated functor compatibility give the collision coherence.
\end{proof}

The restriction and collision compatibility in this proposition is a hypothesis on the
chosen geometric model. It does not follow from naming an ambient stable category.
Conilpotence here refers to finite normalized words in that model.
No word-length completion is implicit.


\section{Three coproducts}
The theorem system uses three coproducts with distinct geometric meanings:
\[
\begin{aligned}
 \Delta_{\mathrm{fact}}&:\text{separation of chiral support},\\
 \Delta_{\mathrm{dec}}&:\text{cutting an ordered bar word},\\
 \Delta_{\mathrm{Hopf}}&:\text{fusion of representations}.
\end{aligned}
\]
The first two are present on the transverse bar.  The third is additional bialgebra data on an acting algebra.

\begin{proposition}[Coproduct separation]
Factorization separation is carried by every ordered coefficient.  Deconcatenation is generated functorially by the bar construction.  A Hopf coproduct is supplied by an independent bialgebra structure.  A coalgebra comparison morphism relates the Hopf coproduct to either geometric coproduct whenever a theory identifies their carriers.
\end{proposition}
\begin{proof}
Factorization separation is part of the coefficient object, and deconcatenation is the canonical cut map on tensor words.  A Hopf coproduct is a map from the algebra to a completed tensor square satisfying algebraic coassociativity.  Its source, target, and defining universal property differ from the first two, so a relation among them is exactly the data of a comparison morphism.
\end{proof}

\section{Interval factorization homology}
For a right $A$-module $R$ and a left $A$-module $L$, factorization homology of an interval with endpoint labels gives
\[
 \int_{([0,1],\{0,1\})}(A;R,L)
 \simeq R\tensor_A^{\mathbb L}L.
\]
This statement holds internally in every presentable monoidal category supporting factorization homology \cite{AFT,LurieHA}. Taking $R=L=\mathbbm1$ identifies the interval amplitude with the ordered bar.

\chapter{Cobar and twisting cochains}
Let $C$ be a coaugmented conilpotent differential graded coalgebra in the ambient
coefficient model. The ordinary, direct-sum tensor cobar algebra is
\[
 \Cobarc_X^{\ord}(C)=
 \bigl(T(s^{-1}\bar C),d_{\mathrm{int}}+d_\Delta\bigr).
\]
A twisting cochain $\tau:C\to A$ has degree $1$, vanishes on the coaugmentation,
and has zero composite with the augmentation of $A$. It satisfies
\[
 d\tau+\tau d+m(\tau\tensor\tau)\bar\Delta=0.
\]
The universal twisting cochains are
\[
 \tau_A:\Barc_X^{\ord}(A)\to A,
 \qquad
 \tau_C:C\to\Cobarc_X^{\ord}(C).
\]

\begin{theorem}[Bar--cobar adjunction]
For ordinary augmented dg algebras and coaugmented conilpotent dg coalgebras
in the specified coefficient model, tensor cobar is left adjoint to tensor bar.
Their mapping set is the set of twisting cochains.
An internal adjunction on factorization objects additionally requires both constructions
to produce factorization objects and their twisting maps to preserve the geometric data.
\end{theorem}
\begin{proof}
An algebra map from the tensor algebra $T(s^{-1}\bar C)$ is determined by its restriction to $s^{-1}\bar C$. Compatibility with the cobar differential is the Maurer--Cartan equation above. The same equation describes coalgebra maps from $C$ to the cofree conilpotent coalgebra $T^c(s\bar A)$.
\end{proof}

\chapter{Complete bar--cobar reconstruction}
Ordinary tensor cobar uses finite words. Completion changes that carrier.
Let $k$ be a commutative ring and
$A_{\mathrm{alg}}=\bigoplus_{w\ge0}A_w$ an augmented algebra in degree zero,
with $A_0=k$, augmentation given by projection onto $A_0$, and multiplication adding weights.
Put $D=\Omega B(A_{\mathrm{alg}})$ using ordinary direct-sum tensors,
and let $D_w$ be its weight-$w$ subcomplex.
The complete algebra and its complete resolution are
\[
 \widehat A=\prod_{w\ge0}A_w,\qquad
 \widehat D=\prod_{w\ge0}D_w.
\]
Products of complexes are taken degreewise, with the weight topology.

\begin{theorem}[Complete reconstruction]
The componentwise counit $\widehat D\to\widehat A$ is a quasi-isomorphism
of complete differential graded algebras.
\end{theorem}
\begin{proof}
A weight-$w$ block word has at most $w$ positive-weight entries.
Filter by their number. Multiplication lowers this number and splitting preserves it.
At a fixed list of $N\ge2$ entries, the splitting complex is the tensor product
of the coefficient module with $N-1$ copies of $[k\xrightarrow{1}k]$ in degrees $-1,0$.
Changing the first outer cut to an inner gap contracts it.
The other differential terms cancel by their prefix signs.
This identity survives tensoring with any coefficient module.
Subtracting a primitive at the largest entry number terminates, so negative cohomology vanishes.
In degree zero, the boundaries impose $c(a)c(b)=c(ab)$,
leaving exactly $A_w$.
Taking the degreewise product preserves kernels and images of these differentials,
and therefore preserves the computed cohomology.
\end{proof}

For $k[x]$ with weight one on $x$, the ordinary counit recovers $k[x]$,
whereas this weight completion recovers $k[[x]]$.
A completed geometric bar--cobar equivalence needs a specified completed tensor product,
coalgebra category, and comparison maps compatible with collision and inverse limits.
An operadic nilcompletion or a boundary double centralizer is not identified
with the ordinary tensor counit by this weight calculation.
The coalgebra-unit and geometric duality statements require their own convergence and comparison arguments.

\section{Quadratic recognition}
Let $A=T(V)/(R)$ with $R\subset V\tensor V$. Its quadratic dual coalgebra $A^{\mathrm{q!}}$ is cogenerated by $sV$ with corelations $s^2R$. The canonical map $A^{\mathrm{q!}}\to\Barc(A)$ gives the Koszul complex.

\begin{theorem}[Quadratic Koszul criterion]
Assume positive weight and degreewise finite duals. The following conditions are equivalent:
\begin{enumerate}
\item $A^{\mathrm{q!}}\to\Barc(A)$ is a quasi-isomorphism;
\item $\Cobarc(A^{\mathrm{q!}})\to A$ is a quasi-isomorphism;
\item the left and right Koszul complexes have homology concentrated in weight zero;
\item $\Tor_i^A(\mathbbm1,\mathbbm1)$ is concentrated in internal weight $i$.
\end{enumerate}
\end{theorem}
\begin{proof}
The universal twisting morphism $A^{\mathrm{q!}}\to A$ is Koszul exactly when the associated twisted tensor products resolve the unit. The fundamental theorem of twisting morphisms identifies this condition with the bar and cobar quasi-isomorphisms. Weight purity is the homological form of the same statement \cite{Priddy,LV}.
\end{proof}

\chapter{Modules, comodules, and Morita theory}
The ordinary tensor counit reconstructs an algebra. Reconstruction of its
modules requires the unit and counit of the twisted module functors as well.
The following weight conditions make these maps finite in every component.

Let $k$ be a field and let $A_{\mathrm{alg}}=\bigoplus_{w\ge0}A_w$
be an associative algebra in cohomological degree zero.
Assume $A_0=k$, multiplication adds weights, and augmentation projects onto $A_0$.
Put $C_{\mathrm{alg}}=B(A_{\mathrm{alg}})$, with its bar weight grading, and set
\[
 A=\prod_{w\ge0}A_w,\qquad C=\prod_{w\ge0}(C_{\mathrm{alg}})_w.
\]
Products are degreewise in complexes.
The universal twisting map $\tau:C\to A$ has degree one,
sends $[a]$ to $a$, and vanishes on the empty word and all longer words.
The bar differential has the earlier convention $b_2[ a|b]=[ab]$.

A complete weight-graded left $A$-module here means a product
$M=\prod_{w\ge w_0}M_w$ of complexes, for some integer $w_0$,
with differential preserving weight and actions $A_u\otimes M_v\to M_{u+v}$.
A complete weight-graded left $C$-comodule means such a product $N$
with a counital coassociative coaction
$\rho:N\to C\widehat\otimes N$ preserving weight and differential.
The completed tensor product is defined componentwise by
\[
 (C\widehat\otimes N)_w
 =\bigoplus_{u+v=w}(C_{\mathrm{alg}})_u\otimes N_v.
\]
This sum is finite since $u\ge0$ and $v\ge w_0$.
Each reduced coaction lowers the weight of its final $N$-component.
It therefore vanishes after more than $w-w_0$ iterations on $N_w$.
These comodules need no chosen coaugmentation.

Write $D_{\mathrm{wt}}(A)$ and $D_{\mathrm{wt}}(C)$ for the categories
obtained by inverting maps that are quasi-isomorphisms in every weight,
respectively for these modules and comodules.
This specifies the weak equivalences; no coderived localization is implicit.
Define
\[
 F(M)=C\widehat\otimes_\tau M,\qquad
 G(N)=A\widehat\otimes_\tau N.
\]
The first has the coaction induced by $\Delta_C$ and the second the left $A$-action.
Writing $\Delta(c)=\sum c_{(1)}\otimes c_{(2)}$ and
$\rho(n)=\sum n_{(-1)}\otimes n_{(0)}$, their differentials are
\[
\begin{aligned}
 d_F(c\otimes m)={}&d_Cc\otimes m+(-1)^{|c|}c\otimes d_Mm
  +\sum(-1)^{|c_{(1)}|}c_{(1)}\otimes\tau(c_{(2)})m,\\
 d_G(a\otimes n)={}&a\otimes d_Nn
  -\sum a\tau(n_{(-1)})\otimes n_{(0)}.
\end{aligned}
\]
All sums are finite in each weight.
The twisting equation, associativity, and coassociativity make both differentials square to zero.

\begin{theorem}[Koszul module equivalence in positive weights]
These functors induce inverse equivalences
\[
 D_{\mathrm{wt}}(A)\simeq D_{\mathrm{wt}}(C).
\]
Their counit and unit are
\[
 \epsilon_M(a\otimes c\otimes m)=a\epsilon_C(c)m,\qquad
 \eta_N(n)=\sum n_{(-1)}\otimes1\otimes n_{(0)}.
\]
\end{theorem}
\begin{proof}
The differential formulas show that these are cochain maps,
respectively $A$-linear and $C$-colinear.
The twisting terms cancel in the unit equation; the counit equation uses the module action.
Coassociativity and the counit identities give the two triangle identities.

The complex $GF(M)$ is the two-sided normalized bar resolution
$A\widehat\otimes B(A_{\mathrm{alg}})\widehat\otimes M$.
For its contraction, use the simplicial boundary convention
$\partial=\sum_i(-1)^i d_i$.
Multiplying a bar word of length $r$ by $(-1)^r$ translates this convention
to the displayed twisting differentials.
Writing $\bar a=a-\epsilon_A(a)1$, the extra degeneracy is
\[
 h\bigl(a_0[a_1|\cdots|a_r]m\bigr)
 =1[\bar a_0|a_1|\cdots|a_r]m.
\]
The first face returns the original term minus its scalar first factor.
The other faces cancel against applying $h$ after the boundary,
including the module-action face.
The internal differential of $M$ cancels by the suspension sign.
Thus $\partial h+h\partial=1-i\epsilon_M$, with $i(m)=1[\,]m$.
This is a linear contraction; $i$ need not be $A$-linear.
It preserves total weight, so the counit is a quasi-isomorphism in every weight.

At fixed total weight, filter $FG(N)$ by the weight of its final $N$-factor.
This filtration is finite.
The reduced-coaction term lowers that weight and disappears in the associated graded.
The associated graded unit is therefore the inclusion into
$B(k,A_{\mathrm{alg}},A_{\mathrm{alg}})\otimes N_v$.
The rightmost extra degeneracy contracts this augmented bar onto $k$:
\[
 h\bigl([a_1|\cdots|a_r]a\bigr)
 =(-1)^{r+1}[a_1|\cdots|a_r|\bar a]1
\]
in the same simplicial convention.
The last face gives $1-i\epsilon$; all other faces cancel in pairs.
Consequently the unit is a quasi-isomorphism on each associated graded component,
and the finite filtration proves the same for each weight of $N$.

Both functors preserve weightwise quasi-isomorphisms.
Filter by the weight of the final module or comodule component.
The twisting terms strictly change that weight in one direction.
On associated graded components, the functors tensor complexes of $k$-vector spaces,
which preserves quasi-isomorphisms.
The filtration is finite at fixed total weight.
Hence $F,G$ descend to the stated localizations, where their unit and counit are isomorphisms.
\end{proof}

This theorem uses an explicit weight decomposition and a lower weight bound on every module and comodule.
Derived $I$-adic completion, with $I=\ker(A\to k)$, instead uses
$M\to\varprojlim_p(A/I^p\otimes_A^{\mathbb L}M)$.
The two completions require a comparison of their filtrations.
A coderived category also has a separately specified class of coacyclic complexes.
Neither identification follows from the algebra counit alone.
For right and left modules, the two-sided bar computes their relative derived tensor product
in the same finite weight components.
An identification with a geometric derived cotensor product additionally requires
the two-sided twisting comparison and its compatibility with the chosen completion and collision maps.

Let $\Ccal$ be a presentable factorization category and let $b$ be a compact factorization generator. The internal endomorphism object
\[
 A_b=\RHom_{\Ccal}(b,b)
\]
has ordered multiplication by composition.

\begin{theorem}[Factorization Morita theorem]
Assume $b$ generates each local category and the local Morita equivalences commute with factorization descent. Then
\[
 \Ccal\simeq \Mod_{A_b}(\Fact_D(X)).
\]
\end{theorem}
\begin{proof}
The compact-generator Morita theorem applies on every factorization chart. The descent hypothesis identifies the equivalences on intersections and hence on the Ran colimit.
\end{proof}

\chapter{Derived centre and bulk action}
The ordered Hochschild cochains are
\[
 Z_{\ord}(A)=\RHom_{A\tensor A^{\op}}(A,A).
\]

\begin{theorem}[Higher Deligne centre]
Let the ambient coefficient category be presentable, stable, symmetric monoidal, closed, and let tensor product preserve geometric realizations separately. Then $Z_{\ord}(A)$ carries a canonical $\Etwo$-algebra structure. It is the universal $\Etwo$ algebra acting centrally on $A$.
\end{theorem}
\begin{proof}
Hochschild cochains carry the brace operations, and the brace operad is equivalent to chains on little two-disks. This is the Deligne theorem \cite{McClureSmith,KontsevichSoibelmanDeligne}. Universality is the identification of Hochschild cochains with endomorphisms of the identity $A$-bimodule.
\end{proof}

A boundary condition in a field theory gives a morphism
\[
 \beta:\Obs_{\mathrm{bulk}}|_{\partial}\longrightarrow Z_{\ord}(A_{\partial}).
\]
Locality makes the action central. A boundary condition is open--closed complete on a protected sector when $\beta$ is an equivalence on that sector.

For a Calabi--Yau category, Hochschild cochains carry the framed $\Etwo$ refinement of the cyclic Deligne theorem \cite{BravRozenblyum}.

\chapter{Cyclic closure and modular sewing}
A cyclic trace becomes a surface theory after dualizability and nondegeneracy are supplied.

\begin{definition}[Perfect Calabi--Yau ordered algebra]
A perfect Calabi--Yau ordered algebra is a dualizable ordered algebra $A$ together with a trace
\[
 \operatorname{tr}:\HH_0(A)\longrightarrow\mathbbm1[-d]
\]
whose induced pairing $A\tensor A\to\mathbbm1[-d]$ is nondegenerate.
\end{definition}

\begin{theorem}[Ribbon-graph realization of a Calabi--Yau algebra]
Let $A$ be a homologically finite, smooth and proper cyclic $\Ainf$ algebra with a nondegenerate pairing of degree $d$.  Then the cyclic bar construction of $A$ defines an open topological conformal field theory.  Its operations are indexed by chains on moduli spaces of bordered surfaces, equivalently by the corresponding ribbon-graph complexes.  The value on the circle is Hochschild homology $\HH_*(A)$.
\end{theorem}
\begin{proof}
The nondegenerate cyclic pairing identifies outgoing and incoming boundary labels and contracts the two flags of every internal ribbon edge.  The $\Ainf$ identities are exactly the boundary identities for metric ribbon graphs.  Smoothness and properness give the dualizability and finite pairing required for all contractions.  Costello's construction then produces the open TCFT and identifies the closed state complex with Hochschild chains \cite{CostelloCY}.
\end{proof}

\begin{remark}
The trace functional supplies the annular operation.  Smoothness, properness, and the nondegenerate Calabi--Yau pairing extend it to the complete ribbon-surface theory of the theorem.
\end{remark}

Stable-curve sewing is supplied by a modular operad action. Its contractions form the modular structure carried by the chiral coefficient. The two sewing geometries coexist when their contraction maps supercommute.

\part{Exact noncommutative models}

\chapter{Free and square-zero pairs}
Work in a presentable stable linear closed monoidal coefficient category $\mathcal C$ whose
tensor preserves coproducts separately and is exact.
For $M\in\mathcal C$, put
\[
 T_X(M)=\mathbbm1\oplus M\oplus M^{\tensor2}\oplus\cdots
\]
This is the free augmented algebra in $\mathcal C$.
All sums and shifts are formed in $\mathcal C$.
Even when $M$ is a factorization object, this componentwise construction
need not be one, as the separated-locus calculation in
Chapter~\ref{ch:vol2-quantum-interfaces} shows.

\begin{theorem}[Free--square-zero duality]
Assume $M$ is dualizable and use the continuous dual for word weight.
Let $\Barc$ in this statement denote the ordinary augmented bar in
$\mathcal C$. Then
\[
 \Barc(T_X(M))\simeq \mathbbm1\oplus sM,
 \qquad
 \RHom_{T_X(M)}(\mathbbm1,\mathbbm1)
 \simeq \mathbbm1\oplus M^\vee[-1]
\]
with square-zero multiplication on the reduced summand.
Conversely, for $A=\mathbbm1\oplus M$ with $M^2=0$,
\[
 \Barc(A)=T^c(sM),\qquad A^!\simeq \widehat T(M^\vee[-1]).
\]
\end{theorem}
\begin{proof}
In weight $r$, bar terms are copies of $M^{\tensor r}$ indexed by cuts
of a word of length $r$. For $r\ge2$, insertion of the first possible
cut contracts the cut complex. The suspension signs give the total
complex homotopy. In weight one the primitive class $sM$ survives.

The explicit cut identity is proved in
Proposition~\ref{prop:vol2-free-bar}.
Dualizing this primitive coalgebra gives square-zero multiplication.
For the square-zero algebra all multiplication faces vanish, leaving
the cofree tensor coalgebra. Its continuous weight dual is the
completed tensor algebra.
\end{proof}

\chapter{The quantum plane}
Let
\[
 A_q=\kfield\angles{x,y}/(yx-qxy),\qquad q\in\kfield^\times.
\]
Give $x$ and $y$ weight one.

\begin{theorem}[PBW basis of the quantum plane]
The monomials $x^ay^b$, $a,b\ge0$, form a basis of $A_q$.
\end{theorem}
\begin{proof}
Use the rewriting rule $yx\mapsto qxy$. Every reduction decreases the inversion number of a word. The leading word $yx$ has distinct initial and terminal letters, so every overlap decomposes into independent reductions. Bergman's diamond argument therefore gives the ordered monomials as normal forms.
\end{proof}

The quadratic dual is
\[
 A_q^!=\kfield\angles{\xi,\eta}/(\xi^2,\eta^2,\xi\eta+q\eta\xi).
\]

\begin{theorem}[Koszulness]
$A_q$ is Koszul, and its Koszul complex is
\[
 0\to A_q\tensor\kfield(\xi\eta)
 \xrightarrow{d_2}
 A_q\tensor(\kfield\xi\oplus\kfield\eta)
 \xrightarrow{d_1}A_q\to\kfield\to0,
\]
with
\[
 d_1(a\tensor\xi+b\tensor\eta)=ax+by,
\quad
 d_2(a\tensor\xi\eta)=ay\tensor\xi-qax\tensor\eta.
\]
\end{theorem}
\begin{proof}
Direct substitution gives $d_1d_2=ayx-qaxy=0$. The image of $d_1$ is the augmentation ideal $(x,y)$, so the complex is exact at $A_q$ and has cokernel $\kfield$. We now compute $\ker d_1$.

Write
\[
 u=\sum_{i,j\ge0}u_{ij}x^iy^j,
 \qquad
 v=\sum_{i,j\ge0}v_{ij}x^iy^j.
\]
The relation $y^jx=q^jxy^j$ gives
\[
 ux+vy=\sum_{a,b\ge0}
 \bigl(q^b u_{a-1,b}+v_{a,b-1}\bigr)x^ay^b,
\]
where coefficients with a negative index are zero. Thus $ux+vy=0$ implies $u_{i,0}=0$, $v_{0,j}=0$, and
\[
 v_{a,b}=-q^{b+1}u_{a-1,b+1}.
\]
Define $c=\sum_{i,j\ge0}u_{i,j+1}x^iy^j$. Then $u=cy$ and, using $x^iy^jx=q^jx^{i+1}y^j$,
\[
 v=-qcx.
\]
Hence $(u,v)=d_2(c)$. Finally, $d_2(c)=0$ implies $cy=0$; the PBW basis and $q\ne0$ show that $A_q$ is a domain, so $c=0$. The complex is exact, and the quadratic Koszul criterion proves that $A_q$ is Koszul.
\end{proof}

\begin{theorem}[Yoneda algebra of the quantum plane]
Use the convention that a Yoneda product $\alpha\beta$ means first $\beta$ and then $\alpha$. The augmentation module has
\[
 \Ext^*_{A_q}(\kfield,\kfield)
 \cong
 \kfield\angles{\xi,\eta}/(\xi^2,\eta^2,\xi\eta+q\eta\xi)
 =A_q^!
\]
as a bigraded algebra, where $|\xi|=|\eta|=1$ in cohomological degree and each generator retains its internal weight one.
\end{theorem}
\begin{proof}
Apply $\Hom_{A_q}(-,\kfield)$ to the minimal Koszul resolution in the preceding theorem. Every matrix entry of $d_1$ and $d_2$ lies in the augmentation ideal, so the dual differential vanishes. The cohomology has basis
\[
 1,\quad \xi,\eta,\quad \xi\eta
\]
in cohomological degrees $0,1,1,2$. Splicing the two degree-one extensions gives the degree-two class. The syzygy $y\otimes\xi-qx\otimes\eta$ identifies the two splicings by
$\xi\eta=-q\eta\xi$, while the absence of degree-two $x^2$ and $y^2$ syzygies gives $\xi^2=\eta^2=0$. These relations account for the complete four-dimensional cohomology, and therefore present the Yoneda algebra.
\end{proof}

The bar, quadratic dual and derived augmentation boundary of $A_q$ are
ordinary algebraic constructions.
Tensoring with a coefficient object gives an ordered algebra in the
ambient pointwise coefficient category when that tensor is defined.
A factorization realization additionally requires its separated-locus
equivalences and collision maps. These do not follow from componentwise
tensoring with $A_q$.

\chapter{Quiver path algebras}
Let $Q$ be a finite quiver over a field $\kfield$ and
$S=\bigoplus_{i\in Q_0}\kfield e_i$.
Assign a vector space $M_a$ to each arrow $a:i\to j$ and set
\[
 M=\bigoplus_{a\in Q_1}e_jM_ae_i,
 \qquad
 T_S(M)=S\oplus M\oplus M\tensor_S M\oplus\cdots.
\]
Words are composable paths. These are ordinary $S$-bimodules.
Placing the construction in geometric coefficients requires a functor
preserving the relative tensors, bar realization and intended
factorization structures, as in Proposition~\ref{prop:vol2-bar-functor}.

For the quiver $1\xrightarrow{a}2$, the path algebra has basis $e_1,e_2,a$ and relation $a=e_2ae_1$. The simple modules $S_1,S_2$ satisfy
\[
 \Ext^1(S_1,S_2)\cong\kfield,
 \qquad
 \Ext^1(S_i,S_i)=0,
 \qquad
 \Ext^1(S_2,S_1)=0.
\]
The bar differential records concatenation of paths, and the quadratic dual reverses the arrow with cohomological shift.

\begin{theorem}[Relative path-algebra bar]
For a tensor path algebra $T_S(M)$ over semisimple $S$,
\[
 S\tensor_{T_S(M)}^{\mathbb L}S\simeq S\oplus sM
\]
as an $S$-bicomodule. A quadratic quotient $T_S(M)/(R)$ is Koszul precisely when its relative quadratic complex is exact away from weight zero.
\end{theorem}
\begin{proof}
The last-arrow contraction is $S$-bilinear and gives the same extra degeneracy as in the free case. The quotient statement follows from the relative twisting-morphism criterion.
\end{proof}

\chapter{Weyl and Clifford deformations}
Let $V$ be finite-dimensional and let $\omega:V\tensor V\to\kfield$ be alternating. The Weyl algebra is
\[
 W_\hbar(V,\omega)=T(V)[[\hbar]]/
 \angles{uv-vu-\hbar\omega(u,v)}.
\]
For a purely odd space and symmetric $\omega$, the same formula with the supercommutator $uv+vu$ gives a Clifford algebra.

\begin{theorem}[Filtered PBW theorem]
For alternating $\omega$ in the even Weyl case, or symmetric $\omega$ in the purely odd Clifford case, the associated graded algebra is $\Sym(V)$ in the Weyl case and $\bigwedge(V)$ in the Clifford case.
\end{theorem}
\begin{proof}
Filter by word length. The constant term lowers filtration by two, so the associated graded relation is the commutative or skew-commutative quadratic relation. The cubic ambiguities reduce to the cocycle identity for a central scalar form, which holds because scalars commute. The ordered monomials give the PBW basis.
\end{proof}

The quadratic-linear-constant presentation has a curved Koszul dual. The constant term becomes the curvature functional
\[
 m_0:s^2R\longrightarrow\kfield[2].
\]
This construction is the algebraic origin of the curvature; geometric anomaly curvature enters through a separate comparison morphism.


\chapter{Associative chiral algebras on the formal disc}
Let $V$ be a complete topological vector space with vacuum $\mathbf1$ and translation $T$.  A state--field map
\[
 Y(-,z):V\longrightarrow\End(V)((z))
\]
is a coordinate presentation of an associative chiral-convolution product when its iterated products are expansions of one factorization kernel in the nested regions $|z_1|>|z_2|$ and $|z_2|>|z_1|$.

\begin{proposition}[Associative convolution from weak associativity]
Suppose $Y(\mathbf1,z)=\id$, $Y(a,z)\mathbf1\in a+zV[[z]]$, $[T,Y(a,z)]=\partial_zY(a,z)$, and for every $a,b,c$ there is $N$ such that
\[
 (z_1-z_2)^N Y(a,z_1)Y(b,z_2)c
 =(z_1-z_2)^N Y(Y(a,z_1-z_2)b,z_2)c
\]
in the common formal expansion ring.  Then the fields define an associative chiral algebra on the formal disc.  The converse holds for a vacuum-generated translation-equivariant associative chiral algebra admitting state--field reconstruction.
\end{proposition}
\begin{proof}
The displayed identity is the equality of the two nested restrictions of the three-point chiral product.  Vacuum and translation give the unit and coordinate covariance.  Factorization on higher configuration spaces is obtained by iterating the equality; coherence follows from associativity of the chiral-convolution product.
\end{proof}

\begin{example}[Braided current fields]
Let $R(z)$ be an invertible solution of the spectral Yang--Baxter equation, normalized by $R^{21}(-z)R(z)=1$.  Generating fields $a_i(z)$ with exchange
\[
 a_i(z_1)a_j(z_2)
 =\sum_{k,l}R_{ij}^{kl}(z_1-z_2)a_l(z_2)a_k(z_1)
\]
form a quantum vertex algebra whenever the vacuum, translation, and weak associativity relations hold.  The $R$-matrix supplies the comparison between the two formal expansion regions and records planar exchange in addition to transverse order.
\end{example}

\chapter{Zamolodchikov--Faddeev exchange}
Let $V$ be finite-dimensional and let $R(u)\in\End(V\tensor V)\tensor\kfield(u)$ be invertible. Introduce generators $Z_i(u)$ with exchange relation
\[
 Z_1(u)Z_2(v)=R_{12}(u-v)Z_2(v)Z_1(u).
\]

\begin{theorem}[Yang--Baxter confluence]
The exchange algebra has a consistent normal ordering for three generators precisely when
\[
 R_{12}(u-v)R_{13}(u-w)R_{23}(v-w)
 =R_{23}(v-w)R_{13}(u-w)R_{12}(u-v).
\]
\end{theorem}
\begin{proof}
Reverse the word $Z_1(u)Z_2(v)Z_3(w)$. Exchanging positions in the order $(12),(13),(23)$ gives the left side. Exchanging in the order $(23),(13),(12)$ gives the right side. These are the two reductions of the unique length-three overlap. Equality is therefore the diamond condition.
\end{proof}

For the rational permutation solution
\[
 R(u)=1-\frac{\hbar P}{u},
\]
the Yang--Baxter equation follows from $P_{ij}^2=1$ and
$P_{12}P_{13}P_{23}=P_{23}P_{13}P_{12}$.

\chapter{The RTT Yangian of \texorpdfstring{$\mathfrak{gl}_2$}{gl2}}
Let
\[
 T(u)=1+\sum_{r\ge1}T^{(r)}u^{-r},
 \qquad R(u)=1-\frac{\hbar P}{u}.
\]
The RTT Yangian $Y_\hbar(\mathfrak{gl}_2)$ is generated by the coefficients $t_{ij}^{(r)}$ subject to
\[
 R(u-v)T_1(u)T_2(v)=T_2(v)T_1(u)R(u-v).
\]
Equivalently,
\[
 (u-v)[t_{ij}(u),t_{kl}(v)]
 =\hbar\bigl(t_{kj}(u)t_{il}(v)-t_{kj}(v)t_{il}(u)\bigr).
\]

\begin{proposition}[Coefficient relations]
Put $t_{ij}^{(0)}=\delta_{ij}$. The RTT identity is equivalent to
\[
 [t_{ij}^{(r+1)},t_{kl}^{(s)}]-[t_{ij}^{(r)},t_{kl}^{(s+1)}]
 =\hbar\bigl(t_{kj}^{(r)}t_{il}^{(s)}-t_{kj}^{(s)}t_{il}^{(r)}\bigr),
 \qquad r,s\ge0.
\]
In particular,
\[
 [t_{ij}^{(1)},t_{kl}^{(1)}]
 =\hbar(\delta_{kj}t_{il}^{(1)}-\delta_{il}t_{kj}^{(1)}),
\]
so the first modes form the enveloping algebra of the copy $\hbar\mathfrak{gl}_2$ inside the Yangian.
\end{proposition}
\begin{proof}
Multiply the generating-series relation by $u-v$ and compare the coefficient of $u^{-r-1}v^{-s-1}$. The two terms contributed by multiplication by $u-v$ produce the two commutators on the left, and the two ordered products on the right are the corresponding coefficients of the RTT numerator. Setting $r=0$ gives the displayed first-mode relation.
\end{proof}

\begin{theorem}[Hopf structure]
The formulas
\[
 \Delta(T(u))=T(u)\dot\tensor T(u),\qquad
 \varepsilon(T(u))=1,\qquad
 S(T(u))=T(u)^{-1}
\]
define a Hopf algebra structure in the $u^{-1}$-adic completion. Coefficientwise,
\[
 \Delta(t_{ij}^{(r)})
 =\sum_{a=1}^{2}\sum_{p=0}^{r}
 t_{ia}^{(p)}\tensor t_{aj}^{(r-p)}.
\]
\end{theorem}
\begin{proof}
Apply the RTT relation in each tensor factor. Matrix multiplication in the auxiliary space gives the RTT relation for $\Delta(T)$. Coassociativity is associativity of matrix multiplication. Expanding the matrix product gives the coefficient formula. The inverse series exists uniquely because $T(u)=1+O(u^{-1})$; its defining identities give the antipode axioms.
\end{proof}

The evaluation map
\[
 \operatorname{ev}_a:T(u)\longmapsto 1+\frac{\hbar E}{u-a},
 \qquad E=(E_{ij}),
\]
lands in $U(\mathfrak{gl}_2)[[u^{-1}]]$ and gives the standard evaluation modules. The PBW theorem identifies the associated graded with $U(\mathfrak{gl}_2[t])$ \cite{Molev,DrinfeldYangian}.

This example contains three structures: ordered multiplication in the Yangian, coproduct on tensor products of modules, and exchange by the rational $R$-matrix.

\chapter{The Hall algebra of the \texorpdfstring{$A_2$}{A2} quiver}
Let $Q:1\to2$ over $\mathbb F_q$. Write $S_1,S_2$ for the simples and $M$ for the unique non-split extension
\[
 0\to S_2\to M\to S_1\to0.
\]
In the twisted Hall algebra, multiplication counts subobjects with Euler-form normalization. Direct counting gives
\[
 [S_1]*[S_2]=v^{-1}[S_1\oplus S_2]+v^{-1}[M],
 \qquad
 [S_2]*[S_1]=[S_1\oplus S_2],
 \qquad v^2=q.
\]
Hence the root vector
\[
 E_{12}=E_1E_2-v^{-1}E_2E_1
\]
is represented by the indecomposable extension.

\begin{theorem}[Quantum Serre relation]
The Hall generators satisfy
\[
 E_1^2E_2-(v+v^{-1})E_1E_2E_1+E_2E_1^2=0,
\]
and the analogous relation with $1$ and $2$ exchanged.
\end{theorem}
\begin{proof}
Every representation of dimension vector $(2,1)$ is a linear map $f:\mathbb F_q^2\to\mathbb F_q$ of rank zero or one. Read a word from the top quotient to the bottom quotient. The complete-flag counts are
\[
\begin{array}{c|ccc}
\operatorname{rk}f&112&121&211\\ \hline
0&q+1&q+1&q+1\\
1&q+1&1&0.
\end{array}
\]
For rank zero every one-dimensional subspace of $\mathbb F_q^2$ is admissible, giving $q+1$ flags in each word. For rank one, the word $121$ forces the unique kernel line, the word $211$ is incompatible with the final subrepresentation, and the word $112$ allows every line at the remaining stage.

The total Euler-form exponents of the three words are respectively $-1,0,1$. Their twisted coefficients are therefore $v^{-1},1,v$. On the rank-zero stratum the Serre combination has coefficient
\[
 (q+1)\bigl(v^{-1}-(v+v^{-1})+v\bigr)=0.
\]
On the rank-one stratum it has coefficient
\[
 v^{-1}(q+1)-(v+v^{-1})
 =v^{-1}(v^2+1)-(v+v^{-1})=0.
\]
The two isomorphism classes exhaust the dimension vector $(2,1)$, proving the relation. This is Ringel's realization of $U_v^+(\mathfrak{sl}_3)$ \cite{Ringel,Green}.
\end{proof}

\chapter{Matrix traces and the necklace bracket}
Let $F=\kfield\angles{x,y}$ and $F_{\mathrm{cyc}}=F/[F,F]$. Cyclic derivatives are defined by
\[
 \partial_x^{\mathrm{cyc}}[a_1\cdots a_n]
 =\sum_{a_i=x}a_{i+1}\cdots a_na_1\cdots a_{i-1}
\]
and similarly for $y$. Define
\[
 \{f,g\}_{\mathrm{neck}}
 =\bracks{
 \partial_x^{\mathrm{cyc}}f\,\partial_y^{\mathrm{cyc}}g
 -\partial_y^{\mathrm{cyc}}f\,\partial_x^{\mathrm{cyc}}g}_{\mathrm{cyc}}.
\]

\begin{theorem}[Trace realization]
For matrices $X,Y\in\mathfrak{gl}_N$ and the canonical symplectic bracket on matrix entries,
\[
 \{\Tr f(X,Y),\Tr g(X,Y)\}
 =\Tr\{f,g\}_{\mathrm{neck}}(X,Y).
\]
\end{theorem}
\begin{proof}
Differentiate the trace by moving the varied matrix entry cyclically to the first position. This gives the cyclic derivative. Contract the $X$ derivative of one trace with the $Y$ derivative of the other and subtract the reversed contraction. The matrix indices close into the trace of the cyclic product.
\end{proof}

In the stable commuting-matrix quotient, let
\[
 p_{ab}=\Tr(X^aY^b).
\]
Then
\[
 \{p_{ab},p_{cd}\}=(ad-bc)p_{a+c-1,b+d-1}.
\]
For two eigenvalues, the relative invariants
\[
 A=(x_1-x_2)^2,
 \quad B=(x_1-x_2)(y_1-y_2),
 \quad C=(y_1-y_2)^2
\]
satisfy $B^2=AC$ and carry the standard $\mathfrak{sl}_2$ Poisson bracket after linear rescaling.


\part{The nonlocal ladder}

\chapter{Spectral exchange as an associative chiral operation}
\section{The coefficient field and the exchange operator}
Let $V$ be a finite-dimensional vector space and let $K$ be a field of meromorphic functions in spectral differences.  Fix an invertible operator
\[
 R(u)\in\End(V\tensor V)\tensor K
\]
that satisfies the spectral Yang--Baxter equation
\begin{equation}\label{eq:vol1-spectral-ybe}
 R_{12}(u-v)R_{13}(u-w)R_{23}(v-w)
 =R_{23}(v-w)R_{13}(u-w)R_{12}(u-v).
\end{equation}
The equation states that two routes through the chamber complex of three labelled points give the same transport.

\begin{definition}[Zamolodchikov--Faddeev algebra]
The exchange algebra $\mathcal Z_R$ is generated by the components of a current $A(u)\in V\tensor\mathcal Z_R$ with relation
\begin{equation}\label{eq:vol1-zf}
 A_1(u)A_2(v)=R_{12}(u-v)A_2(v)A_1(u).
\end{equation}
A normal word has spectral labels in a fixed chamber $u_1\succ\cdots\succ u_n$.
\end{definition}

\begin{theorem}[Confluence is Yang--Baxter]
The quadratic rewriting system defined by \eqref{eq:vol1-zf} is confluent on words with pairwise distinct spectral labels precisely when \eqref{eq:vol1-spectral-ybe} holds.  In that case the normal ordered words form a $K$-basis of $\mathcal Z_R$.
\end{theorem}
\begin{proof}
Every two-letter inversion is removed by one application of \eqref{eq:vol1-zf}.  The minimal overlap occurs in a word of length three.  Moving the leftmost label through the other two gives the coefficient
$R_{12}R_{13}R_{23}$, while moving the rightmost label in the opposite sequence gives $R_{23}R_{13}R_{12}$.  Equality of the two coefficients is \eqref{eq:vol1-spectral-ybe}.  Bergman's diamond lemma then gives the normal-word basis.
\end{proof}

\section{A diagonal exchange model}
Choose a basis $e_1,\ldots,e_N$ of $V$ and functions $q_{ij}(u)$ satisfying
\[
 q_{ij}(u)q_{ji}(-u)=1,
 \qquad q_{ii}(u)=1.
\]
The diagonal operator
\[
 R(u)(e_i\tensor e_j)=q_{ij}(u)e_i\tensor e_j
\]
satisfies the Yang--Baxter equation because all three factors are scalars.  The generators $a_i(u)$ obey
\[
 a_i(u)a_j(v)=q_{ij}(u-v)a_j(v)a_i(u).
\]
Every ordered monomial is a basis vector.  The quantum plane is the constant two-generator specialization $q_{12}=q$, $q_{21}=q^{-1}$.

\begin{calculation}[Three-particle coefficient]
For $u\succ v\succ w$,
\[
 a_i(w)a_j(v)a_k(u)
 =q_{ij}(w-v)q_{ik}(w-u)q_{jk}(v-u)
  a_k(u)a_j(v)a_i(w).
\]
The coefficient is independent of the sequence of adjacent exchanges.  This scalar identity is the diagonal Yang--Baxter equation.
\end{calculation}
\begin{proof}
Each inversion of two adjacent letters contributes its scalar exchange factor.  Both reduced words for the longest permutation invert the same three pairs.  Scalar commutativity makes the products identical.
\end{proof}

\chapter{RTT and reflection algebras}
\section{The RTT bialgebra}
Let $T(u)=(t_{ij}(u))$ be a matrix of formal currents.  The RTT algebra $\mathcal T_R$ is defined by
\begin{equation}\label{eq:vol1-rtt-general}
 R_{12}(u-v)T_1(u)T_2(v)
 =T_2(v)T_1(u)R_{12}(u-v).
\end{equation}

\begin{theorem}[RTT coproduct]
The formula
\[
 \Delta(T(u))=T_{13}(u)T_{12}(u),
 \qquad \epsilon(T(u))=I,
\]
defines a coassociative bialgebra structure on $\mathcal T_R$.
\end{theorem}
\begin{proof}
Apply \eqref{eq:vol1-rtt-general} in each tensor factor.  The Yang--Baxter equation moves the three $R$-operators through the three matrix factors.  The resulting identity is the RTT relation for $\Delta(T)$.  Matrix multiplication is associative, hence the two iterated coproducts both equal $T_{14}T_{13}T_{12}$.
\end{proof}

For the Yang matrix $R(u)=1-\hbar P/u$, expansion at $u=\infty$ gives the RTT Yangian.  The coefficient of $u^{-1}v^{-1}$ is
\[
 [t_{ij}^{(1)},t_{kl}^{(1)}]
 =\hbar\bigl(\delta_{kj}t_{il}^{(1)}-\delta_{il}t_{kj}^{(1)}\bigr),
\]
so the first current layer is the Rees enveloping algebra of $\mathfrak{gl}_N$.

\section{Boundary exchange}
A reflecting boundary reverses the spectral path.  For a constant invertible $R\in\End(V^{\tensor2})$ satisfying Yang--Baxter, define the reflection algebra $\mathcal K_R$ by
\begin{equation}\label{eq:vol1-reflection}
 R_{21}K_1R_{12}K_2
 =K_2R_{21}K_1R_{12}.
\end{equation}

\begin{theorem}[Boundary confluence]
The relation \eqref{eq:vol1-reflection} gives a coherent exchange law for three boundary insertions.  The two reductions of every length-three overlap agree by the Yang--Baxter equation and its inverse.
\end{theorem}
\begin{proof}
A boundary exchange consists of an incoming crossing, reflection, and outgoing crossing.  The two reductions of three insertions sweep the same braid in the doubled half-plane.  Algebraically, one inserts \eqref{eq:vol1-reflection} three times and uses Yang--Baxter to pass the middle $R$-matrix through the outer pair.  The remaining factors cancel in inverse pairs.
\end{proof}

The RTT algebra governs bulk line exchange.  The reflection algebra governs boundary line exchange.  The ZF algebra governs charged fields in a fixed representation.  Their common coefficient is the same $R$-matrix, while their products live in three different endomorphism objects.

\chapter{Quantum vertex realization}
\section{Two expansion regions}
Let $V$ be a complete topological vector space with vacuum $\mathbf1$, translation $T$, state-field map $Y$, and a unitary rational operator
\[
 S(z):V\tensor V\longrightarrow V\tensor V((z)).
\]
The two formal regions $|z_1|>|z_2|$ and $|z_2|>|z_1|$ are related by $S(z_1-z_2)$.

\begin{definition}[Rational quantum vertex datum]
A rational quantum vertex datum satisfies vacuum and translation covariance together with
\begin{equation}\label{eq:vol1-s-locality}
 (z_1-z_2)^N Y(a,z_1)Y(b,z_2)
 =(z_1-z_2)^N
 Y^{21}\!\left(S(z_1-z_2)(a\tensor b);z_2,z_1\right)
\end{equation}
for a sufficiently large $N$, and weak associativity in every triple of states.
\end{definition}

\begin{theorem}[Exchange coherence]
Suppose $S$ is unitary,
\[
 S^{21}(-z)S(z)=1,
\]
and satisfies the spectral Yang--Baxter equation.  Then the pairwise continuation maps obtained from \eqref{eq:vol1-s-locality} define a braid-group action on every $n$-point field product.  Weak associativity makes the action compatible with iterated operator products.
\end{theorem}
\begin{proof}
Unitarity identifies a crossing followed by its reverse with the identity.  Yang--Baxter identifies the two continuations around every codimension-two triple-collision stratum.  These are the defining relations of the braid group.  Weak associativity identifies the collision limit in each chamber with the corresponding iterated field.
\end{proof}

When $S$ is the graded flip, the datum is an ordinary local vertex algebra.  A nontrivial $S$ retains the ordered chiral product and its exchange monodromy.  This is the formal-disc realization of an associative chiral algebra with braiding data in the sense of quantum vertex algebra theory \cite{LiNonlocal,LiQVA}.

\section{The Yangian field}
Let $T(u)$ act on a family of evaluation modules.  The RTT relation becomes $S$-locality for the matrix coefficients of $T(u)$.  The normalized Yang operator
\[
 S(u)=f(u)\left(1-\frac{\hbar P}{u}\right),
 \qquad
 f(u)f(-u)=\left(1-\frac{\hbar^2}{u^2}\right)^{-1},
\]
is unitary.  Its expansion at infinity gives the quantum vertex exchange of Yangian currents.  The scalar factor changes normalization and preserves the Yang--Baxter equation.

\chapter{Shuffle realization and the positive current algebra}
\section{Pairwise kernel}
Let $K(z)$ be a rational function and put
\[
 \operatorname{Sh}_n=\C(z_1,\ldots,z_n)^{\mathfrak S_n}.
\]
Write $\operatorname{Sym}$ for the sum over ordered two-block
partitions, with no factorial divisor.
For $f\in\operatorname{Sh}_m$ and $g\in\operatorname{Sh}_n$ define
\begin{equation}\label{eq:vol1-shuffle}
 (f*g)(z_1,\ldots,z_{m+n})
 =\operatorname{Sym}
 \left(
 f(z_1,\ldots,z_m)g(z_{m+1},\ldots,z_{m+n})
 \prod_{i\le m<j}K(z_i-z_j)
 \right).
\end{equation}

\begin{theorem}[Associativity of the shuffle product]
The product \eqref{eq:vol1-shuffle} is associative on every subspace closed under the displayed symmetrization and pole conditions.
\end{theorem}
\begin{proof}
In a triple product, every ordered pair of variables belonging to different blocks contributes exactly one factor $K(z_i-z_j)$.  The product of all pairwise factors is independent of the parenthesization.  Final symmetrization gives the same function.
\end{proof}

For the tripled Jordan quiver, one uses
\[
 K(z)=\prod_{a=1}^3\frac{z+\epsilon_a}{z},
 \qquad \epsilon_1+\epsilon_2+\epsilon_3=0.
\]
A critical Hall or Yangian realization additionally specifies its
coefficient lattice, wheel conditions and comparison map
\cite{SVYangians,Tsymbaliuk}.
The displayed meromorphic kernel alone does not identify that realization.
The ordinary polynomial Hall kernels $N(v)/v$ and $vN(v)$ of
Chapter~\ref{ch:vol2-quantum-interfaces} have their own distinct lattices.

\begin{calculation}[Degree one against degree one]
For $f(z_1)=z_1^r$ and $g(z_1)=z_1^s$,
\[
 f*g
 =z_1^rz_2^sK(z_1-z_2)+z_2^rz_1^sK(z_2-z_1).
\]
Put $u=z_1-z_2$ and $\sigma_3=\epsilon_1\epsilon_2\epsilon_3$.
The antisymmetric part is
\[
 \frac{2\sigma_3}{u^3}
 (z_1^rz_2^s-z_1^sz_2^r).
\]
For generic nonzero parameters, $u=-\epsilon_a$ is a zero of
$K(u)$ and contributes no residue of this commutator.
For $(r,s)=(2,0)$, the residue at the true diagonal $u=0$,
with $z_2$ fixed, is $2\sigma_3$.
\end{calculation}
\begin{proof}
The two ordered partitions give the displayed product.
Subtract the reversed product and use
$K(u)-K(-u)=2\sigma_3/u^3$.
For $(r,s)=(2,0)$, substitute $z_1=z_2+u$.
The result is $4\sigma_3z_2/u^2+2\sigma_3/u$.
Shifted residues of the reciprocal kernel are instead given by
\eqref{eq:vol2-shuffle-reciprocal}. A geometric operation must specify
which meromorphic coefficient it uses.
\end{proof}

\section{Coproduct and double}
A chamber expansion fixes a coefficient completion, but it does not
prove compatibility of a proposed coproduct with multiplication.
Theorem~\ref{thm:vol2-shuffle-no-cooperation} gives a charge-two
obstruction for the ordinary tensor on the polynomial Hall carriers.
Proposition~\ref{prop:vol2-rational-cooperation} constructs a positive
rational cooperation with a specified changed middle interchange.
A Drinfeld double additionally requires compatible positive and negative
Hopf data, a continuous pairing and a completion in which its canonical
tensor exists. Under those hypotheses the double construction gives
the quasitriangular structure.

\chapter{The ladder in one diagram}
The nonlocal structures form the chain
\[
\text{ordered coefficient}
\longrightarrow \mathcal Z_R
\longrightarrow \mathcal T_R
\longrightarrow (V,Y,S),
\]
\[
(V,Y,S)
\longrightarrow \operatorname{Sh}^{+}
\longrightarrow D\!\left(\operatorname{Sh}^{+}\right).
\]
The arrows retain different data:
\begin{center}
\small
\begin{tabularx}{0.96\textwidth}{@{}>{\raggedright\arraybackslash}p{0.18\textwidth}>{\raggedright\arraybackslash}X>{\raggedright\arraybackslash}p{0.29\textwidth}@{}}
\toprule
\textbf{Object}&\textbf{Defining operation}&\textbf{Coherence}\\
\midrule
$\mathcal Z_R$&exchange of charged fields&Yang--Baxter equation\\
$\mathcal T_R$&matrix-current product and coproduct&RTT and coassociativity\\
$\mathcal K_R$&boundary exchange&reflection equation\\
$(V,Y,S)$&formal field product&$S$-locality and weak associativity\\
$\operatorname{Sh}^{+}$&extension shuffle&pairwise-kernel associativity\\
$D(\operatorname{Sh}^{+})$&positive--negative pairing&quasitriangularity\\
\bottomrule
\end{tabularx}
\end{center}
This ladder gives a concrete generator-and-relation theory of doubly noncommutative chiral observables.  The first noncommutativity is ordered composition.  The second is exchange between distinct order chambers.  Chiral singularity remains the diagonal coefficient throughout.


\part{Field theory and master equations}

\chapter{A finite cotangent BF model}
Let $L$ be a finite-dimensional dg Lie algebra. Its cotangent BV space is
\[
 \mathcal E_L=L[1]\oplus L^*[-2],
\]
with the canonical degree $-1$ symplectic pairing. Write $A\in L[1]$ and $B\in L^*[-2]$. The cotangent action is
\[
 S_L(A,B)=\angles{B,d_LA+\frac12[A,A]}.
\]

\begin{theorem}[Cotangent master equation]
The Hamiltonian vector field of $S_L$ is square-zero, and consequently $\{S_L,S_L\}=0$.
\end{theorem}
\begin{proof}
The $A$ component is the Maurer--Cartan vector field $d_LA+\frac12[A,A]$; the $B$ component is its cotangent lift. The square of the first component is the sum of $d_L^2$, the derivation identity for $d_L$, and the Jacobi identity. The cotangent lift preserves commutators of vector fields, so the full Hamiltonian vector field is square-zero.
\end{proof}

\begin{proposition}[Coordinate form of the modular term]
Let $L$ be an ordinary finite-dimensional Lie algebra with basis $e_a$ and structure constants $[e_a,e_b]=f^c_{ab}e_c$. Let $c^a$ be the odd coordinates on $L[1]$ and $b_a$ their BV-conjugate coordinates. With
\[
 S=\frac12 b_c f^c_{ab}c^ac^b,
 \qquad
 \Delta=-\sum_a\frac{\partial}{\partial c^a}
                 \frac{\partial}{\partial b_a},
\]
one has
\[
 \Delta S=\sum_b\Tr(\ad_{e_b})c^b.
\]
Thus the one-loop class of the finite cotangent theory is exactly the modular character.
\end{proposition}
\begin{proof}
Differentiation first in $b_a$ gives $\frac12f^a_{ij}c^ic^j$. Odd differentiation in $c^a$ gives $f^a_{aj}c^j$, and the chosen sign in $\Delta$ converts this to $f^a_{ja}c^j=\Tr(\ad_{e_j})c^j$. This is the displayed formula.
\end{proof}

A finite-dimensional dg Lie algebra is unimodular when the supertrace of every adjoint endomorphism vanishes. Every finite-dimensional nilpotent Lie algebra in degree zero is unimodular because each $\ad_x$ is nilpotent.

\begin{theorem}[Finite-dimensional quantum master equation]
Let $L$ be unimodular and let $\Delta$ be the BV Laplacian determined by the translation-invariant Berezinian on $\mathcal E_L$. Then
\[
 \Delta S_L=0,
 \qquad
 \frac12\{S_L,S_L\}+\hbar\Delta S_L=0,
\]
and equivalently
\[
 \hbar\Delta\bigl(e^{S_L/\hbar}\bigr)=0.
\]
The quantum observable differential
\[
 \delta_{S_L}=\{S_L,-\}+\hbar\Delta
\]
satisfies $\delta_{S_L}^2=0$.
\end{theorem}
\begin{proof}
The classical master equation gives $\{S_L,S_L\}=0$. The scalar $\Delta S_L$ is the divergence of the Chevalley vector field on $L[1]$. Its linear coefficient is the modular character $x\mapsto\operatorname{str}(\ad_x)$, which vanishes by unimodularity. The BV exponential identity
\[
 \hbar^2e^{-S_L/\hbar}\Delta(e^{S_L/\hbar})
 =\frac12\{S_L,S_L\}+\hbar\Delta S_L
\]
proves the exponential form. The standard BV commutator calculation then gives
\[
 \delta_{S_L}^2=\left\{\frac12\{S_L,S_L\}+\hbar\Delta S_L,-\right\}=0.
\]
\end{proof}

Now take the local dg Lie algebra
\[
 \Lcal_X=\Omega^{0,*}(X)\widehattensor\Omega^*(\R)\tensor\flie
\]
for a finite-dimensional nilpotent $\flie$, and form its shifted cotangent theory with compactly supported density dual. The preceding finite theorem governs every pairing-preserving finite-dimensional deformation retract. The local classical observables are Chevalley cochains and are holomorphic on $X$ and locally constant on $\R$.

Choose the standard interval polarization in which the $L[1]$ variables are fixed at one end and the cotangent variables at the other. Canonical quantization of the boundary linear Poisson algebra gives the associative zero-mode algebra $U(\flie)$. Gluing the two augmentation boundary states across an interval gives the derived tensor product
\[
 \mathbbm1\tensor_{U(\flie)}^{\mathbb L}\mathbbm1
 \simeq\Barc(U(\flie))
 \simeq C_*(\flie),
\]
the Chevalley--Eilenberg chains. A solution of the renormalized quantum master equation in the local BV deformation complex produces the corresponding infinite-dimensional quantum factorization algebra.

\chapter{BRST and ordered reduction}
Let $(A,Q)$ be a dg ordered algebra, with $Q$ a derivation. Extend $Q$ to the bar factorwise. Write $b$ for the multiplication bar differential.

\begin{proposition}[BRST--bar bicomplex]
With the suspension sign convention,
\[
 Q_{\Barc}b+bQ_{\Barc}=0.
\]
Hence $D=Q_{\Barc}+b$ satisfies $D^2=0$.
\end{proposition}
\begin{proof}
Apply $Q$ to a bar face that multiplies adjacent entries. The derivation formula splits the result into the two terms obtained by applying $Q$ before that face. The totalization sign makes their sum the anticommutator.
\end{proof}

The filtration by bar length gives
\[
 E_1\cong\Barc(H^*(A,Q))
 \Longrightarrow H^*(\Barc(A),D)
\]
under boundedness and completeness. Quantum Drinfeld--Sokolov reduction is a current-algebra realization of this bicomplex.

\chapter{Stokes and algebraic identities}
Let $C_*(K)$ be the oriented cellular chains of a regular manifold with corners. Suppose each codimension-one face $F$ is assigned a degree-one operator $\rho_F$ on a coefficient complex $E$, and the assignments are compatible with restrictions to corners.

\begin{theorem}[Stokes-to-master]
Assume that each $\rho_F$ is a chain map of total degree one, that its restriction to a corner depends on the induced contraction, and that disjoint contractions act with the Koszul sign prescribed by the oriented normal bundle. Then the component of
\[
 D^2,\qquad D=d_E+d_K+\sum_F\rho_F,
\]
supported on a codimension-two corner $S$ is the oriented sum
\[
 \sum_{F\supset S}\epsilon(F,S)\,\rho_{F/S}\rho_F.
\]
Thus $D$ is a differential precisely when these corner sums vanish.
\end{theorem}
\begin{proof}
Expand $D^2$ and group terms by support strata. The chain-map condition cancels the mixed internal terms. A term supported on $S$ is obtained by crossing two incident faces. The two orders carry the boundary orientations of the two sides of the oriented normal square. Their signed sum is the displayed expression.
\end{proof}

Associahedra give Stasheff identities. Fulton--MacPherson collision strata give chiral and graph identities after a coefficient map is supplied. Stable-graph boundaries give BV and modular master equations.

\chapter{Two edge colors and two loop operators}
Let $\mathsf G^{\perp,\ch}$ be the completed vector space spanned by stable graphs with transverse and chiral edges. Contracting an edge that joins distinct vertices gives brackets $[-,-]_\perp$ and $[-,-]_{\ch}$. Contracting a loop gives operators $\Delta_\perp$ and $\Delta_{\ch}$. Product orientation gives
\[
 \Delta_\perp\Delta_{\ch}+\Delta_{\ch}\Delta_\perp=0.
\]

\begin{theorem}[Bicolored master equation]
Let $\mathsf P^{\perp,\ch}$ be a two-colored modular operad whose edge contractions carry the product determinant orientation, and let $\mathfrak g_{\mathsf P}$ be the completed convolution Lie algebra of its Feynman transform with an endomorphism modular operad. Algebra structures on the transform are Maurer--Cartan elements $\Theta\in\mathfrak g_{\mathsf P}^1$ satisfying
\[
 d\Theta+\frac12[\Theta,\Theta]_\perp
 +\frac12[\Theta,\Theta]_{\ch}
 +\hbar_\perp\Delta_\perp\Theta
 +\hbar_{\ch}\Delta_{\ch}\Theta=0.
\]
\end{theorem}
\begin{proof}
The Feynman-transform differential contracts one edge. Separating edges give the two brackets and loop edges give the two BV operators. The product orientation makes mixed contractions anticommute. The condition that the structure map commute with the transform differential is exactly the displayed Maurer--Cartan equation \cite{GetzlerKapranov}.
\end{proof}

\chapter{Typed anomalies}
An anomaly is a cohomology class in the deformation complex that governs the attempted construction.
\begin{center}
\begin{tabularx}{\textwidth}{@{}>{\raggedright\arraybackslash}p{0.22\textwidth}>{\raggedright\arraybackslash}X>{\raggedright\arraybackslash}p{0.25\textwidth}@{}}
\toprule
construction&complex&class\\ \midrule
central extension&local Lie algebra cochains&cocycle\\
BRST quantization&BRST deformation complex&nilpotence obstruction\\
BV quantization&local BV cochains&quantum master anomaly\\
conformal-block family&Atiyah algebra of its determinant line&projective curvature\\
curved Koszul duality&curved convolution Lie algebra&curvature $m_0$\\
\bottomrule
\end{tabularx}
\end{center}
A comparison theorem consists of a morphism between two complexes and an equality of the transported classes. This formulation keeps the geometric origin of every scalar visible.

\part{Calabi--Yau and automorphic interfaces}

\chapter{Calabi--Yau categories}
Let $\Ccal$ be a smooth linear category with a $d$-Calabi--Yau structure. Two universal consequences are available.

\begin{theorem}[Cyclic and shifted outputs]
The Hochschild cochains $CC^*(\Ccal)$ form a framed $\Etwo$ algebra. The derived moduli stack of objects carries a $(2-d)$-shifted symplectic structure whenever the representability hypotheses hold.
\end{theorem}
\begin{proof}
The first assertion is the cyclic Deligne theorem \cite{BravRozenblyum}. The second is the PTVV--Brav--Dyckerhoff construction: the Calabi--Yau trace pairs the tangent complex $\RHom(E,E)[1]$ with itself in degree $2-d$ \cite{PTVV,BravDyckerhoff}.
\end{proof}

A geometric BV theory or an open factorization category supplies an ordered chiral algebra through the recognition theorem of Chapter 3. This is the bridge from Calabi--Yau geometry to the present volume.

\chapter{Hall and BPS algebras}
A Hall multiplication starts from the extension correspondence
\[
 \M_{\alpha}\times\M_{\beta}
 \xleftarrow{\ p\ }\M_{\alpha,\beta}^{\mathrm{ext}}
 \xrightarrow{\ q\ }\M_{\alpha+\beta}.
\]
With a coefficient theory for which $p^*$ and $q_*$ are defined,
\[
 a*b=q_*p^*(a\boxtimes b).
\]
Associativity follows from the stack of two-step filtrations.

For 2-Calabi--Yau categories with good moduli spaces, the BPS algebra is the positive enveloping algebra of a generalized Kac--Moody Lie algebra \cite{DHS}. For quivers and surface resolutions, preprojective and surface CoHAs recover positive Yangian halves in established cases \cite{SVYangians,DPPSSV}.

The ordered chiral structure enters when Hall multiplication is promoted to a factorization coefficient over a curve or defect parameter. A coproduct, pairing, and completion then produce the corresponding double.

\chapter{The Igusa denominator as a bar character}
Let $\g_\Delta$ be the Gritsenko--Nikulin Borcherds--Kac--Moody superalgebra and let $\nilp$ be its positive nilpotent part. Root spaces are finite-dimensional and the positive cone is locally finite. Fix a proper integral height $h$ that is positive on the effective root cone, and set
\[
 F^{>H}\nilp=\bigoplus_{h(\alpha)>H}\g_{\Delta,\alpha},
 \qquad
 \nilp^{(H)}=\nilp/F^{>H}\nilp.
\]
Height additivity makes $F^{>H}\nilp$ an ideal; properness makes $\nilp^{(H)}$ finite-dimensional and nilpotent. Every bar, Chevalley, enveloping, and character formula in this chapter is interpreted coefficientwise in the separated height completion.

\begin{theorem}[Bar--Chevalley identification]
For each height $H$ there is a natural quasi-isomorphism of conilpotent coalgebras
\[
 \Barc\bigl(U(\nilp^{(H)})\bigr)\simeq C_*(\nilp^{(H)}).
\]
The quotient maps are compatible with these quasi-isomorphisms, and therefore induce a quasi-isomorphism in the separated height-complete category
\[
 \Barc_h^\wedge\bigl(U(\nilp)\bigr):=\varprojlim_H\Barc\bigl(U(\nilp^{(H)})\bigr)
 \simeq
 C_{*,h}^\wedge(\nilp):=\varprojlim_H C_*(\nilp^{(H)}).
\]
The coefficientwise Euler supercharacter of the completed Chevalley complex is
\[
 \chi\,C_{*,h}^\wedge(\nilp)
 =\prod_{\alpha>0}(1-e^{-\alpha})^{\sdim\g_{\Delta,\alpha}}.
\]
\end{theorem}
\begin{proof}
For every $H$, the PBW filtration on $U(\nilp^{(H)})$ identifies the associated graded with $\Sym(\nilp^{(H)})$. The bar complex of the enveloping algebra and the Chevalley complex have the same associated-graded Koszul resolution. The standard antisymmetrization map is a filtered quasi-isomorphism. Naturality for the quotient maps gives the inverse-limit comparison; proper height makes every homogeneous coefficient stabilize. For an even root vector the exterior Chevalley factor contributes $1-e^{-\alpha}$; for an odd root vector it contributes its inverse. Multiplication over the root spaces gives the product.
\end{proof}

With the positive-root convention $e^{-\rho}=q^{1/2}r^{1/2}s^{1/2}$, the denominator theorem reads
\[
 e^{-\rho}\,\chi\,C_{*,h}^\wedge(\nilp)=D_5(2\Omega)=2^{-6}\Delta_5(2\Omega)
\]
\cite{GN97,GN98b}. This identity is the exact meeting point of ordered bar theory and automorphic products.

\chapter{Finite-height currents and an exact physical model}
Retain the finite-height quotients $\nilp^{(H)}$ fixed in the preceding chapter.

The current Lie conformal superalgebra is
\[
 \operatorname{Cur}(\nilp^{(H)})=\kfield[\partial]\tensor\nilp^{(H)},
 \qquad [a_\lambda b]=[a,b].
\]
Its universal enveloping vertex superalgebra is $V_H$. The current generators carry the cocommutative vertex-bialgebra coproduct
\[
 \Delta\bigl(a(-1)\mathbbm1\bigr)
 =a(-1)\mathbbm1\tensor\mathbbm1
 +\mathbbm1\tensor a(-1)\mathbbm1,
 \qquad a\in\nilp^{(H)},
\]
extended uniquely as a vertex-algebra morphism.

\begin{theorem}[Height-complete current realization]
Let $\mathfrak n=\bigoplus_{\alpha>0}\mathfrak n_\alpha$ be a root-graded Lie superalgebra with finite-dimensional root spaces and a proper additive height function.  Put
\[
 \mathfrak n^{(H)}=\mathfrak n/\!\bigoplus_{h(\alpha)>H}\mathfrak n_\alpha.
\]
Then every $\mathfrak n^{(H)}$ is finite-dimensional and nilpotent, the quotient maps form an inverse system, and
\[
 \widehat U_h(\mathfrak n)=\varprojlim_HU(\mathfrak n^{(H)}),\quad
 \widehat C_{*,h}(\mathfrak n)=\varprojlim_HC_*(\mathfrak n^{(H)}),\quad
 \widehat V_h(\mathfrak n)=\varprojlim_HU^{\ch}(\Cur_X\mathfrak n^{(H)})
\]
are separated complete algebras, coalgebras, and current algebras.  Their associated graded objects are the corresponding completed symmetric objects, and the bar--Chevalley comparison is compatible with the inverse system.
\end{theorem}
\begin{proof}
Properness gives finitely many roots below a fixed height.  Additivity implies that the roots of height $>H$ form a Lie ideal and that the quotient is nilpotent.  PBW and Chevalley antisymmetrization are natural under quotient maps.  Taking inverse limits in each finite total height therefore gives the displayed complete objects and preserves the coefficientwise quasi-isomorphism.
\end{proof}

Apply the finite BF model of Chapter 19 to every $\nilp^{(H)}$. The inverse system gives a height-complete holomorphic--topological classical field theory with completed boundary algebra and interval complex
\[
 U_h^\wedge(\nilp):=\varprojlim_H U(\nilp^{(H)}),
 \qquad
 C_{*,h}^\wedge(\nilp):=\varprojlim_H C_*(\nilp^{(H)}).
\]
This gives an exact algebraic realization, a local classical holomorphic--topological realization, and an exact quantum realization of the finite-dimensional zero-mode sectors of the positive enveloping algebra and its denominator character. The compact $K3\times E$ identification is the geometric comparison studied in the companion volumes.

\chapter{Character, bracket, parity, and orientation}
A character is an Euler invariant.  Faithful reconstruction combines it with the bracket, parity decomposition, and orientation line carried by the intended object.

\begin{theorem}[Three reconstruction fibres]
The following fibres contain distinct structured objects with the same scalar image.
\begin{enumerate}
\item The PBW-character fibre of a graded vector space with basis $x,y,z$ in degrees $a,b,a+b$ contains both the abelian Lie algebra and the Heisenberg Lie algebra $[x,y]=z$.
\item The superdimension-zero fibre contains both the zero super vector space and $\kfield^{1|1}$.
\item The squaring fibre of a line-valued section contains $s$ and $s\otimes\chi$ for every sign local system $\chi$.
\end{enumerate}
Consequently a faithful reconstruction datum consists respectively of a bracket, a parity-refined dimension, and an oriented square-root line.
\end{theorem}
\begin{proof}
PBW identifies the graded vector-space character of an enveloping algebra with the symmetric-superalgebra character of its underlying graded Lie space, so the two displayed brackets have the same PBW character.  The two super vector spaces in the second item both have superdimension zero and have parity dimensions $(0,0)$ and $(1,1)$.  Finally $\chi^{\otimes2}\cong\mathbbm1$, hence
\[
 (s\otimes\chi)^{\otimes2}\cong s^{\otimes2}.
\]
The extra structures listed in the conclusion distinguish each pair.
\end{proof}

\begin{recognition}[Minimal PBW recognition]
Let $\phi:\mathfrak g\twoheadrightarrow\mathfrak h$ be a surjective morphism of positively graded Lie superalgebras with finite-dimensional homogeneous pieces.  If the parity-refined Hilbert series of $U(\mathfrak g)$ and $U(\mathfrak h)$ agree, then $\phi$ is an isomorphism.  Equality of signed PBW characters suffices when the parity of every primitive degree is prescribed in advance.
\end{recognition}
\begin{proof}
If the kernel is nonzero, choose its first positive degree.  In that degree all decomposable PBW contributions agree by minimality, while the primitive even or odd dimension in $U(\mathfrak h)$ is strictly smaller.  This contradicts equality of parity-refined Hilbert series.  With prescribed parity, the formal logarithm of the PBW product recovers the signed primitive dimension inductively in height, and prescribed parity recovers the ordinary dimension.
\end{proof}

\chapter{Exact dependency theorem}
\begin{theorem}[Ordered-chiral closure]
Fix:
\begin{enumerate}
\item a six-functor Ran category and its factorization subcategory;
\item a transversely ordered algebra, an associative chiral algebra, or a chiral Lie algebra, as appropriate;
\item every mixed distributive, completion, duality, trace, or analytic datum named by the intended construction.
\end{enumerate}
Then the three bars, their complete cobars, their module/comodule Koszul transforms, the Hochschild centre, the formal-disc realization, the ribbon theory, and the finite-height current realization are defined functorially and satisfy the identities proved in this volume.  Every passage between the pointwise and chiral-convolution monoidal products is carried by the displayed mixed distributive datum.
\end{theorem}
\begin{proof}
Each output is one of the functors constructed in the preceding chapters.  The source and target monoidal categories are displayed in every definition.  The comparison theorems invoke the explicitly supplied distributive, PBW, completion, duality, or trace datum.  Composition of these functors gives the asserted functorial system.
\end{proof}

\chapter{Synthesis}
The invariant language has three algebraic axes. Transverse order is an
$\Eone$ structure for $\tensor^!$, associative chirality an $\Eone$
structure for $\starCh$, and chiral Lie theory a Lie structure for
$\starCh$. Their bars are $\BarPerp$, $\BarAssCh$ and $\BarLieCh$.
The aligned projector describes the canonical common-decomposition
comparison under its stated descent hypotheses.
A mixed distributive law and comparisons among the three bars retain
their own maps and coherence equations.

The constructed operations form the two functorial rows
\[
 \text{HT boundary theory}
 \xrightarrow{\ \mathrm{Obs}_{\partial}\ }
 \Alg_{\Eone}(\Fact_D(X))
 \xrightarrow{\ \Barc^{\ord}\ }
 \Coalg_{\mathrm{conil,cpl}}(\Fact_D(X)),
\]
\[
 \Alg_{\Eone}(\Fact_D(X))
 \xrightarrow{\ Z_{\ord}\ }
 \Alg_{\Etwo}(\Fact_D(X)).
\]
A primitive Lie algebra enters when a separate Chevalley or Hall comparison identifies the bar coalgebra with the Chevalley chains of that Lie algebra. Hall geometry supplies this comparison through extension correspondences. Quantum groups require an additional coproduct and exchange structure. Automorphic forms arise from graded Euler characters after a grading and completion have been fixed. Each arrow therefore preserves the geometric operation named at its source.

\part{Technical appendices to the preceding book}
\chapter{Sign calculation for the length-three bar}
For homogeneous $a,b,c\in\bar A$, write $[a|b|c]=sa\tensor sb\tensor sc$. With the standard suspension convention,
\[
 b_m[a|b|c]
 =(-1)^{|a|}[ab|c]
 +(-1)^{|a|+|b|+1}[a|bc].
\]
Applying $b_m$ again gives
\[
 b_m^2[a|b|c]
 =(-1)^{|b|}\bigl[(ab)c-a(bc)\bigr],
\]
which vanishes by associativity.

\chapter{Rational \texorpdfstring{$R$}{R}-matrix verification}
Let $P_{ij}$ permute tensor factors $i$ and $j$. Expanding
\[
 R_{12}(u)R_{13}(u+v)R_{23}(v)
\]
and its reversed product shows equality coefficientwise. The cubic coefficient uses
\[
 P_{12}P_{13}P_{23}=P_{23}P_{13}P_{12},
\]
the quadratic coefficients use the symmetric-group relations, and the linear coefficients agree termwise.

\chapter{Carriers and comparison maps}
The five-carrier geometry, transverse bar construction, mixed-forest signs, and noncommutative examples use four kinds of categorical data:
\begin{enumerate}
\item a factorization carrier is paired explicitly with its operation;
\item the Ran category is fixed before its tensor products are used;
\item the block-ordered operad is introduced after a representing chiral operad is chosen;
\item ribbon and stable-curve sewing are derived from nondegenerate Calabi--Yau and modular data.
\end{enumerate}
A comparison between these constructions therefore specifies its source and target categories and the operations preserved by its map.

\chapter{Enveloping bars and planar path sequences}
\label{ch:enveloping-audit}

The dimension of a bar chain space and the dimension of its homology are different invariants. For an enveloping algebra, the Chevalley comparison gives a Lie-theoretic model for the homology. This makes it possible to compare Motzkin and Riordan path counts with the ordered bar on a specified carrier.

Bar length and Chevalley degree are homological degrees: $H_n$ denotes degree-$n$ homology. The cohomology of the dual Chevalley complex is written $H^n$. For finite-dimensional weight pieces their dimensions agree.

\section{Chevalley comparison}

\begin{theorem}[Bar--Chevalley comparison]
\label{thm:bar-chevalley-spine}
Let $k$ be a field of characteristic zero and let $\g$ be an ordinary Lie algebra over $k$.
Give $U(\g)$ its augmentation killing $\g$.
There is a natural quasi-isomorphism of ordinary conilpotent differential graded coalgebras
\[
 \kappa:C_*(\g)\longrightarrow\BarPerp\bigl(U(\g)\bigr),
\]
where both constructions use direct sums and ordinary tensors over $k$.
It preserves the PBW filtration.
If $\g=\bigoplus_{w\ge1}\g_w$ and $[\g_u,\g_v]\subseteq\g_{u+v}$,
the map is a quasi-isomorphism in each total weight.
It extends to the degreewise product over weights, using completed tensor products by total weight.
The completed coalgebras are conilpotent in each weight; no global conilpotence is asserted for infinite sums.
\end{theorem}
\begin{proof}
Put $C_n(\g)=\Lambda^n\g$ in homological degree $n$, with differential
\[
 d(x_1\wedge\cdots\wedge x_n)
 =\sum_{i<j}(-1)^{i+j+1}[x_i,x_j]\wedge
 x_1\wedge\cdots\widehat{x_i}\cdots\widehat{x_j}\cdots\wedge x_n.
\]
Its coproduct splits the exterior factors between two ordered sublists,
with the sign of the corresponding permutation.
The ordinary normalized bar differential joins positions $i,i+1$ with sign $(-1)^{i-1}$.
Define
\[
 \kappa(x_1\wedge\cdots\wedge x_n)
 =\sum_{\sigma\in S_n}\operatorname{sgn}(\sigma)
 [x_{\sigma(1)}|\cdots|x_{\sigma(n)}].
\]
At degree zero this map preserves the unit.
At degree two its differential is $[x_1x_2-x_2x_1]=[[x_1,x_2]]$.
In every degree, pair permutations whose merged adjacent entries occur in opposite orders.
Their difference is the bracket of those entries, with the displayed exterior sign.
This proves the chain-map equation.
Cutting an antisymmetrized word amounts to choosing an ordered subset and
antisymmetrizing its two subwords.
The two procedures give the same permutations and signs, proving the coalgebra identity.

Filter $U(\g)$ by PBW degree, the bar by total PBW degree, and $\Lambda^n\g$ by degree $n$.
PBW identifies the associated graded algebra with $S(\g)$.
The associated graded comparison is the antisymmetrization
$\Lambda(\g)\to B(S(\g))$, with zero differential on its source.
To verify that it is a quasi-isomorphism, consider the free Koszul resolution
$S(V)\otimes\Lambda(V)$ for $V=\g$.
For a basis $(x_i)$ of $V$, its differential is
$\delta=\sum_i x_i\iota_i$, where $\iota_i$ contracts the exterior factor $x_i$.
On polynomial degree $p$ and exterior degree $q$, with $p+q>0$, the map
\[
 h=\frac{1}{p+q}\sum_i (x_i\wedge -)\,\frac{\partial}{\partial x_i}
\]
satisfies $\delta h+h\delta=1$.
Only finitely many summands act on each element, even when $V$ is infinite-dimensional.
The identity follows by cancellation of cross terms and the Euler identity in the two degrees.
Thus this free resolution has augmentation homology $k$ only.
The free normalized bar resolution has the same homology by its extra-degeneracy contraction.
Antisymmetrization defines an augmentation-preserving comparison between these free resolutions.
The simplicial bar sign convention is related to the displayed bar convention by multiplication by $(-1)^n$ in degree $n$.
Free bases permit a reverse comparison and homotopies to be lifted inductively,
since every positive-degree cycle in a resolution is a boundary.
Tensoring these homotopies with $k$ proves the asserted quasi-isomorphism on associated graded complexes.

Every cycle in the cone of $\kappa$ has finite PBW filtration degree.
A primitive for its highest associated graded part lowers that degree;
iteration terminates because the filtration is bounded below.
The cone is acyclic, proving the ordinary assertion.
For a positive weight grading, PBW degree is at most total weight.
The same argument is therefore finite within each weight.
Finally, kernels and images of a degreewise product differential are products of kernels and images,
so the componentwise quasi-isomorphisms extend to the stated completion.
\end{proof}

The ordinary comparison does not itself construct a morphism in $\FactD(X)$.
For geometric coefficient systems, one must give compatible enveloping and Chevalley constructions on every chart
and coherent factorization equivalences for the whole Chevalley complex.
On each disjointness locus, with restriction denoted by $j^!$, the required equation is
\[
 f_{B,I}\,j^!\kappa_I
 =(\kappa_{I_1}\boxtimes\kappa_{I_2})\,f_{C,I},
 \qquad I=I_1\sqcup I_2.
\]
Here $f_{C,I}$ and $f_{B,I}$ are the specified Chevalley and bar factorization maps.
The bar map may use the signed shuffle comparison under its geometric functor hypotheses.
The Chevalley factorization maps, collision compatibility, refinement coherence,
and continuity of the completed maps remain additional geometric data to construct or verify.
Internal tensor formulas and PBW alone do not provide them.


\begin{calculation}[Finite-dimensional simple case]
\label{calc:simple-spine}
For a simple Lie algebra of rank $r$ with exponents $m_1,\dots,m_r$, the
Chevalley cohomology is the exterior algebra on primitive classes of degrees
$2m_i+1$. The corresponding homology Poincar\'e polynomials are
\[
 \BarPerp U(\mathfrak{sl}_2):\ 1+t^3,
 \qquad
 \BarPerp U(\mathfrak{sl}_3):\ 1+t^3+t^5+t^8,
 \qquad
 \BarPerp U(\mathfrak g_2):\ 1+t^3+t^{11}+t^{14}.
\]
The corresponding Betti sequences for $\mathfrak{sl}_2$ and $\mathfrak{sl}_3$ are
$1,0,0,1$ and $1,0,0,1,0,1,0,0,1$.
For a finite-dimensional Lie algebra, the Chevalley chain dimensions are
$\binom{\dim\g}{n}$.

For an ordinary augmented algebra $A$ over $\mathbb C$, let
$I=\ker(A\to\mathbb C)$.  Its normalized two-sided bar with trivial
coefficients has homological chain spaces
\[
 B_0(A)=\mathbb C,\qquad B_n(A)=I^{\otimes n}\quad(n\ge1).
\]
Here $n$ is bar length, independently of any internal weight or PBW degree.
If $I$ has finite dimension $d$, then $\dim B_n(A)=d^n$; in particular,
an augmentation ideal of dimension three gives the chain counts $3^n$.
In a cohomological convention that places $B_n$ in degree $-n$, its homology
is denoted $H^{-n}$ instead of $H_n$.

For $A=U(\mathfrak{sl}_2)$, PBW gives the basis
$e^a h^b f^c$, with $a,b,c\ge0$, where
$[h,e]=2e$, $[h,f]=-2f$, and $[e,f]=h$.
The augmentation ideal has all these basis elements except $1$.
Consequently $B_n(A)$ is countably infinite-dimensional for every $n\ge1$.
Its total PBW filtration has finite-dimensional associated graded pieces,
with generating function
\[
 \sum_{d\ge0}\dim\operatorname{gr}^{\mathrm{PBW}}_d B_n(A)\,z^d
 =\bigl((1-z)^{-3}-1\bigr)^n.
\]
In particular, $3^n$ is the dimension in the least possible total PBW degree
$d=n$. These least-degree dimensions are $3,9,27,81,243$ for
$n=1,2,3,4,5$; the full bar chain spaces are infinite-dimensional. The quotient of
indecomposables is zero: the perfectness of $\mathfrak{sl}_2$ implies $I/I^2=0$.

The Chevalley chain model is different: its degree-$n$ space is
$\Lambda^n\mathfrak{sl}_2$, of dimensions $1,3,3,1$ in degrees $0,1,2,3$.
With the convention $d(x\wedge y)=[x,y]$, the map
$d:\Lambda^2\mathfrak{sl}_2\to\mathfrak{sl}_2$ sends
\[
 e\wedge h\longmapsto-2e,\qquad
 e\wedge f\longmapsto h,\qquad
 h\wedge f\longmapsto-2f,
\]
and is an isomorphism.  The next differential vanishes, since
\[
 d(e\wedge h\wedge f)
 =-2e\wedge f-h\wedge h-2f\wedge e=0.
\]
Together with $d:\mathfrak{sl}_2\to\mathbb C$ equal to zero, this gives
Chevalley homology dimensions $1,0,0,1$.  The ordinary bar--Chevalley
comparison identifies these homology groups with those of $U(\mathfrak{sl}_2)$;
it does not identify the two chain spaces.

\end{calculation}

\section{The Virasoro row}

The transverse presentation of the Virasoro family is the enveloping algebra of
the positive Witt algebra
\[
 L_1=\bigoplus_{n\ge1}\mathbb C\,e_n,
 \qquad
 [e_i,e_j]=(j-i)e_{i+j},
\]
to which Theorem~\ref{thm:bar-chevalley-spine} applies: $L_1$ is positively
graded with one-dimensional weight spaces, so every weight has finitely many
decompositions.

\begin{theorem}[Pentagonal rigidity]
\label{thm:pentagonal-spine}
The bar homology of $U(L_1)$ satisfies
\[
 \dim H_n\bigl(\BarPerp U(L_1)\bigr)=2
 \qquad\text{for every }n\ge1,
\]
one dimension in each of the weights $(3n^2-n)/2$ and $(3n^2+n)/2$. Its graded
Euler character is
\begin{equation}
\label{eq:witt-euler-spine}
 \chi_q\bigl(\BarPerp U(L_1)\bigr)=\prod_{n\ge1}\bigl(1-q^n\bigr),
\end{equation}
and \eqref{eq:witt-euler-spine} is accounted for degree by degree with no
residual freedom.
\end{theorem}
\begin{proof}
By Theorem~\ref{thm:bar-chevalley-spine} the bar of $U(L_1)$ is quasi-isomorphic to $C_*(L_1)$.
Each weight subcomplex is finite-dimensional, so its homology has the same
dimensions as the weightwise dual cohomology. Goncharova's calculation gives
two dimensions in each degree $n\ge1$, in the weights $(3n^2\mp n)/2$. Identity \eqref{eq:witt-euler-spine} is the
denominator product applied to $L_1$, where every root space is one dimensional
and even, so $\smult(L_1)_n=1$. Euler's pentagonal number theorem expands the
product as $\sum_{m\in\Z}(-1)^mq^{m(3m-1)/2}$, whose support is exactly
$\{(3n^2\pm n)/2\}$ with coefficient $(-1)^n$ at each of $(3n^2-n)/2$ and
$(3n^2+n)/2$. Comparing with $\sum_n(-1)^n\dim H_{n,w}$, each weight of the Euler
character is saturated by a single class in the degree Goncharova predicts, so
no further classes and no cancellation are available.
\end{proof}

For $n=1,2,3,4$, the nonzero homology groups occupy the weight pairs
$(1,2)$, $(5,7)$, $(12,15)$, $(22,26)$, one dimension in each.
The corresponding nonzero coefficients of $\prod_{n\ge1}(1-q^n)$ occur at
$0,1,2,5,7,12,15,22,26$ with signs $+,-,-,+,+,-,-,+,+$.

\section{Planar path sequences}

A Motzkin path is a lattice path from $(0,0)$ to $(n,0)$, staying at
nonnegative height, with steps $(1,1)$, $(1,0)$, and $(1,-1)$.
Let $M(n)$ count all such paths of length $n$, and let $R(n)$ count those with
no horizontal step at height zero.  Both counts include the empty path:
$M(0)=R(0)=1$, while $M(1)=1$ and $R(1)=0$.
For $n\ge1$ consider the two sequences
\[
 a_n=M(n+1)-M(n)=1,2,5,12,30,76,\ldots,
 \qquad
 b_n=R(n+3)=3,6,15,36,91,\ldots.
\]

\begin{proposition}[Motzkin and Riordan identities]
\label{prop:one-sequence-spine}
For every $n\ge0$,
\[
 M(n)=R(n)+R(n+1),\qquad
 M(n+1)-M(n)=R(n+2)-R(n).
\]
The ordinary generating functions of the two path counts, and hence those
of $a_n$ and $b_n$, belong to
$\mathbb Q(z,\sqrt{1-2z-3z^2})$.
\end{proposition}
\begin{proof}
Work in the formal power-series ring $\mathbb Q[[z]]$, and write
$M(z)=\sum_{n\ge0}M(n)z^n$ and $R(z)=\sum_{n\ge0}R(n)z^n$.
A nonempty Motzkin path either starts with a horizontal step followed by
an arbitrary Motzkin path, or has a unique first-return decomposition
$U P D Q$, where $U$ and $D$ are the up and down steps,
$P$ and $Q$ are arbitrary Motzkin paths, and $P$ is drawn
one unit above the horizontal axis.  Thus
\[
 M=1+zM+z^2M^2.
\]
A nonempty path counted by $R$ has no initial horizontal step.  Its first
return similarly has the form $U P D Q$, with $P$ arbitrary and $Q$ counted
by $R$.  Consequently
\[
 R=1+z^2MR,\qquad R=(1-z^2M)^{-1}.
\]
The inverses here exist because their denominators have constant term one.
Using $z^2M^2=M-1-zM$, compute
\[
 (1+zM)(1-z^2M)
 =1+zM-z^2M-z^3M^2=1+z.
\]
It follows that $(1+z)R=1+zM$.  Comparing coefficients of $z^{n+1}$ gives
$R(n+1)+R(n)=M(n)$ for all $n\ge0$.  Subtracting consecutive identities
gives the difference formula.

Solving the quadratic equation for the formal series with constant term
one gives
\[
 M(z)=\frac{1-z-\sqrt{1-2z-3z^2}}{2z^2},\qquad
 R(z)=\frac{1+zM(z)}{1+z},
\]
where the square root has constant term one.  The apparent pole in the
formula for $M$ is removable as a formal series.  Finite index shifts and
differences preserve the indicated quadratic function field, proving the
last assertion.  These identities are consequences solely of the path
decompositions.
\end{proof}

\begin{remark}[Low-degree homology]
Theorem~\ref{thm:pentagonal-spine} gives homology dimensions $2,2,2$ in degrees $1,2,3$, whereas $a_1,a_2,a_3$ are $1,2,5$. These triples agree only in degree two. The degree-one discrepancy already excludes equality of the full homology dimension sequences.
\end{remark}

\begin{theorem}[Enveloping bars and planar path counts]
\label{thm:carrier-spine}
Let $B_\bullet(A)$ be the ordinary normalized bar over $\mathbb C$, with
trivial coefficients and homological degree given by bar length.  With
$a_n=M(n+1)-M(n)$ and $b_n=R(n+3)$ indexed by $n\ge1$ as above,
\[
 \dim H_1 B_\bullet(U(\mathfrak{sl}_2))=0\ne b_1=3,
 \qquad
 \dim H_1 B_\bullet(U(L_1))=2\ne a_1=1,
\]
where $L_1=\bigoplus_{m\ge1}\mathbb C e_m$ and
$[e_i,e_j]=(j-i)e_{i+j}$.
Hence neither sequence is the bar homology dimension sequence of the
corresponding enveloping algebra.  Neither is its full bar chain dimension
sequence: in both cases $B_n$ is countably infinite-dimensional for every
$n\ge1$.
\end{theorem}
\begin{proof}
For any augmented algebra $A$ with augmentation ideal $I$, the bar
differential out of $B_1=I$ is zero and the image from $B_2=I\otimes I$ is
$I^2$.  Therefore
\[
 H_1B_\bullet(A)=I/I^2.
\]
For an ordinary Lie algebra $\mathfrak g$, the inclusion of generators
induces an isomorphism
\[
 \mathfrak g/[\mathfrak g,\mathfrak g]\ \xrightarrow{\ \sim\ }
 I/I^2,\qquad I=\ker(U(\mathfrak g)\to\mathbb C).
\]
Indeed the generators span $I/I^2$, and
$[x,y]=xy-yx$ vanishes there.  For injectivity, equip
$\mathbb C\oplus\mathfrak g/[\mathfrak g,\mathfrak g]$ with the square-zero
product on its second summand.  The quotient map from $\mathfrak g$ to that
summand is a Lie map into the commutator Lie algebra of this associative
algebra.  The universal property of $U(\mathfrak g)$ extends it to an
augmented algebra map, which kills $I^2$ and provides an inverse on the
displayed quotient.

The brackets $[h,e]=2e$, $[h,f]=-2f$, and $[e,f]=h$ show that
$\mathfrak{sl}_2$ equals its commutator algebra.  For $L_1$, every nonzero
bracket has weight at least three.  Conversely, for $m\ge3$,
\[
 e_m=\frac{1}{m-2}[e_1,e_{m-1}].
\]
Thus $[L_1,L_1]=\bigoplus_{m\ge3}\mathbb C e_m$, and its abelianization
has basis the classes of $e_1,e_2$, in weights one and two.  This proves
both degree-one assertions without using any higher-degree calculation.

PBW gives a countably infinite basis of the augmentation ideal of either
enveloping algebra.  In the $L_1$ case the basis consists of nonempty
ordered monomials in the countable set $e_1,e_2,\ldots$, each monomial
having finite length.  Finite tensor powers of either augmentation ideal
are therefore countably infinite-dimensional.  The path counts are finite
in every degree, so they cannot count these full chain spaces.
\end{proof}

In particular, a complex whose homology dimensions are $b_n$ cannot be
quasi-isomorphic to $B_\bullet(U(\mathfrak{sl}_2))$, and one whose homology
dimensions are $a_n$ cannot be quasi-isomorphic to $B_\bullet(U(L_1))$,
with the stated degree convention.  A discrepancy of chain dimensions
alone would not imply this conclusion; the degree-one homology discrepancy
does.

The weight grading gives a different, finite comparison.  An ordered PBW
monomial of weight $w$ specifies a partition of $w$, so for $n\ge1$,
\[
 \sum_{w\ge0}\dim B_n(U(L_1))_w q^w
 =\left(\prod_{m\ge1}(1-q^m)^{-1}-1\right)^n.
\]
Each coefficient is finite.  Thus a weightwise dimension, a dimension
after a specified weight cutoff, and an ungraded dimension are distinct
invariants.  In a weight completion, dimensions must likewise be interpreted
weightwise; the countable algebraic dimensions above refer to the ordinary
direct-sum enveloping algebras.

\begin{remark}[Chiral coefficients and transverse multiplication]
Independently of these scalar calculations, a Virasoro chiral coefficient does not determine the transverse bar: a fixed coefficient may be combined with free, square-zero, enveloping, or quantum-plane transverse algebras whose bars differ, while the coefficient and the central charge are unchanged. A chain-level comparison with an ordered bar therefore also requires a specified transverse multiplication and a quasi-isomorphism from the proposed complex to that bar.
\end{remark}

\chapter{The obstruction constitution}
\label{ch:obstructions}

Each statement below names what is \emph{not} determined, and therefore what
must be supplied. Without them the preceding theorems admit false corollaries.

\begin{proposition}[Invertibility of the aligned comparison]
\label{prop:lambda-sector-spine}
In the finite decomposition model, the canonical aligned inclusion and
retraction are inverse on the full object exactly when its off-aligned
summands vanish. Descent to geometric coefficients additionally requires
compatibility with the actual transition correspondences.
\end{proposition}
\begin{proof}
The canonical maps satisfy $\pi\iota=\id$ and
$\iota\pi=e_{\mathrm{al}}$. The latter is the identity exactly when all
off-aligned summands vanish. On geometric charts the required
base-change squares must identify these projectors across transitions.
The conditional descent theorem states those hypotheses.
\end{proof}

Failure of this canonical comparison to be invertible does not exclude
every noninvertible distributive law. Such a law must be given as a
natural transformation satisfying its own unit and composition diagrams.
For Zamolodchikov--Faddeev algebras, the spectral variable and transverse
word supply the two inputs to this separate construction problem.

\begin{theorem}[No quantum group from dimension loss]
\label{nogo:dimension-loss-spine}
The dimensions of the kernels of $V^{\otimes n}\to\Sym^nV$ do not
determine a multiplication or a Hopf structure. Such a determination
requires operations and their compatibility maps in addition to dimensions.
\end{theorem}
\begin{proof}
The abelian and Heisenberg Lie algebras with generator weights $1,1,2$
have enveloping algebras with the same PBW character
$((1-t)^2(1-t^2))^{-1}$.
The first algebra is commutative and the second is not.
Their graded vector spaces, and hence the stated kernel dimensions,
are nevertheless isomorphic.
Even fixing the multiplication and augmentation is insufficient:
Proposition~\ref{prop:vol2-two-hopf} gives two nonisomorphic Hopf
structures on the same polynomial algebra.
\end{proof}

\begin{theorem}[No modular object from topology alone]
The identity $\partial^2=0$ on $\overline{\mathcal M}_{g,n}$ assigns no
coefficient operations to boundary strata. A modular master element exists only
after cyclic contractions and clutching-compatible coefficient maps have been
supplied.
\end{theorem}

\begin{theorem}[No bulk from a boundary centre]
The derived centre is terminal among algebraic central actions on the boundary
algebra. It does not detect a decoupled bulk sector, so a physical bulk is not
determined by the boundary chiral algebra unless a bulk-to-centre comparison is
supplied and proved to be an equivalence.
\end{theorem}
\begin{proof}
Terminality is the universal property of the centre. It is a statement about
the boundary category and is insensitive to any sector acting trivially on it,
so a bulk with trivial boundary action is invisible.
\end{proof}

\begin{theorem}[No classification by central charge]
Central charge determines the Virasoro anomaly coefficient. It does not
determine the chiral operation, the representation category, the braid
monodromy, or the higher-genus conformal blocks. Distinct holomorphic theories
at equal central charge exist.
\end{theorem}

\begin{theorem}[No identification of anomalies of different types]
\label{nogo:anomalies-spine}
A current-algebra level, a Virasoro central charge, a determinant-line
curvature, a BRST ghost-number anomaly, and a Borcherds-product weight are
invariants in different theories, living in different deformation complexes.
They may be related by a theorem in a particular family. They are not one
universal scalar, and no identity among them follows from numerical agreement.
\end{theorem}
\begin{proof}
Each is a class in a distinct complex: an Atiyah-algebra extension class, a
Virasoro cocycle, the curvature of a determinant line, a class in $H^1$ or
$H^2$ of a BRST convolution complex, and a weight in the character lattice of
an automorphic form. A map between two of these complexes is additional data;
in its absence, equality of numbers is a coincidence of values, not an
identification of classes.
\end{proof}

\begin{theorem}[No unnamed comparison of bars]
An asserted equality $\BarPerp(A)=\BarAssCh(A)$ has no meaning without a
functor carrying the pointwise multiplication problem to the chiral-convolution
multiplication problem. The two have different underlying gradings, counting
interval subdivisions and collision trees respectively. Agreement of Euler
characteristics is not a comparison map.
\end{theorem}

\begin{theorem}[A scalar sequence is not homology data]
\label{nogo:scalar-spine}
A sequence of integers becomes bar homology only after a chain complex, a
differential, a grading convention, and a rank computation or quasi-isomorphism
are specified. Matching initial values, discriminants, or growth rates
constructs no such map.
\end{theorem}
\begin{proof}
Chain-space dimensions and homology dimensions differ by the ranks of the
differentials, which are not functions of the chain dimensions.
Calculation~\ref{calc:simple-spine} exhibits an algebra whose positive
bar chain spaces are countably infinite-dimensional and whose bar
homology dimensions are $1,0,0,1$.
The numbers $3^n$ count only the least-PBW-weight piece in bar length $n$.
Proposition~\ref{prop:one-sequence-spine} exhibits two sequences with the same
discriminant and growth rate that are transforms of one another and of neither
homology.
\end{proof}

Theorem~\ref{nogo:anomalies-spine} together with the central-charge statement
retracts any programme treating a tuple of scalars drawn from these sources as
a single classifying invariant, and any trace identity asserted among them
without a named comparison morphism.

\cleardoublepage
\phantomsection
\addcontentsline{toc}{part}{The Igusa Square Root}
\chapter*{Book entrance: The Igusa Square Root}
\addcontentsline{toc}{chapter}{Book entrance: The Igusa Square Root}
\begin{summarybox}
The Jacobi seed, Borcherds denominator, root superalgebra, local--wrapped currents, quantum group, Hall model, and Fredholm realization.
\end{summarybox}
\part{The automorphic determinant}

\chapter{The protected Jacobi input}
\section{The K3 elliptic genus}
Let $S$ be a K3 surface. The right-moving supersymmetric index of a unitary $\mathcal N=(4,4)$ sigma model is the elliptic genus
\[
 Z_{K3}(\tau,z)=2\phiweak(\tau,z).
\]
The normalized weak Jacobi form has weight zero, index one, and Fourier expansion
\[
 \phiweak(\tau,z)=\sum_{n\ge0}\sum_{\ell\in\Z}f(n,\ell)q^nr^\ell.
\]
In Eichler--Zagier normalization,
\[
 \phiweak(\tau,z)=4\sum_{a=2}^{4}
 \frac{\vartheta_a(\tau,z)^2}{\vartheta_a(\tau,0)^2}
 \in J^{\mathrm w}_{0,1}
\]
\cite{EZ}. The coefficient depends on the discriminant $4n-\ell^2$ and on $\ell$ modulo $2$.

\begin{theorem}[Exact low Fourier coefficients]
The first four $q$-rows are
\[
\begin{aligned}
\phiweak={}&r^{-1}+10+r\\
&+q\bigl(10r^{-2}-64r^{-1}+108-64r+10r^2\bigr)\\
&+q^2\bigl(r^{-3}+108r^{-2}-513r^{-1}+808-513r+108r^2+r^3\bigr)\\
&+q^3\bigl(-64r^{-3}+808r^{-2}-2752r^{-1}+4016\\
&\hspace{7.5em}{}-2752r+808r^2-64r^3\bigr)+O(q^4).
\end{aligned}
\]
In particular,
\[
 f(0,0)=10,
 \qquad \frac{f(0,0)}2=5.
\]
\end{theorem}
\begin{proof}
Insert the theta-series expansions into the preceding quotient formula. Work with $t=q^{1/8}$ and $x=r^{1/2}$; each quotient is an exact formal power series in $t$ with Laurent-polynomial coefficients in $x$. The fractional powers cancel in the sum. Division through $t^{24}$ gives the displayed rows. The arithmetic calculation is reproduced in Appendix~\ref{app:theta-computation}.
\end{proof}

\section{The virtual half-index}
For every $(n,\ell)$ choose a finite-dimensional super vector space $U_{n,\ell}$ satisfying
\[
 \sdim U_{n,\ell}=f(n,\ell).
\]
The integer \(f(n,\ell)\) is its superdimension. In the ordinary
Grothendieck group of finite-dimensional super vector spaces, the
even and odd dimensions remain separate. The signed integer is their
image under superdimension. The determinant depends only on that image.

\begin{lemma}[Fock character and Berezinian]
Let $V=\bigoplus_{\gamma}V_\gamma$ be a locally finite graded super vector space. Then
\[
 \sch\bigl(\Sym(V_{\bar0})\tensor\bigwedge(V_{\bar1})\bigr)
 =\prod_\gamma(1-x^\gamma)^{-\sdim V_\gamma},
\]
and
\[
 \sdet_V(1-x)=\prod_\gamma(1-x^\gamma)^{\sdim V_\gamma}.
\]
\end{lemma}
\begin{proof}
An even one-dimensional summand contributes the geometric series $(1-x^\gamma)^{-1}$. An odd one-dimensional summand contributes $1-x^\gamma$. Multiplication over a homogeneous basis proves both identities.
\end{proof}

\chapter{The Borcherds product}
\section{The effective chamber}
Define
\[
\begin{aligned}
\Gammaeff={}&\set{(n,\ell,m)\in\Z^3:m>0,\ n\ge0}\\
&\cup\set{(n,\ell,0):n>0}
 \cup\set{(0,\ell,0):\ell<0}.
\end{aligned}
\]
The ordering is lexicographic in $(m,n,-\ell)$ and is precisely the ordering selected by the cusp chamber.
For $\gamma=(n,\ell,m)$ put
\[
 V_\gamma=U_{nm,\ell},
 \qquad \sdim V_\gamma=f(nm,\ell).
\]

\begin{definition}[Monic virtual determinant]
The monic determinant of the half-index is
\[
 \Dmon(Z)=q^{1/2}r^{1/2}s^{1/2}
 \sdet_V(1-x)
 =q^{1/2}r^{1/2}s^{1/2}
 \prod_{\Gammaeff}(1-q^nr^\ell s^m)^{f(nm,\ell)}.
\]
\end{definition}

\begin{theorem}[Borcherds--Gritsenko--Nikulin product]
The preceding product converges normally on the chosen cusp chamber and extends to a holomorphic genus-two modular form of weight five with a character of order two. It is the monic Igusa form $\Dmon$. The theta-normalized form is
\[
 \Dtheta=2^6\Dmon.
\]
\end{theorem}
\begin{proof}
The singular theta lift of the vector-valued form corresponding to $\phiweak$ produces the product, divisor, weight, and automorphic character. The weight is $f(0,0)/2=5$. The genus-two identification with the Igusa form and the product formula are the Gritsenko--Nikulin theorems \cite{Borcherds95,Borcherds98,GN97,GN98b}. The product of the ten even theta constants has coefficient $2^6$ at $q^{1/2}r^{1/2}s^{1/2}$, fixing the scalar.
\end{proof}

\section{The theta constant normalization}
Let $\mathfrak E$ be the ten even genus-two characteristics. Define
\[
 \Dtheta(Z)=\prod_{\mathfrak m\in\mathfrak E}\vartheta[\mathfrak m](Z,0).
\]
The cusp leading terms of the ten factors contribute six factors of $2$ and total exponent $(1/2,1/2,1/2)$. Thus
\[
 [q^{1/2}r^{1/2}s^{1/2}]\Dtheta=64.
\]
This coefficient selects a canonical scalar normalization of the automorphic section.

\begin{proposition}[Ten-factor leading-term calculation]
Group the ten even characteristics by the first half-characteristic $a\in(\Z/2)^2$. Their leading contributions in the chosen cusp chamber are
\[
\begin{array}{c|c|c|c}
a&\#&\text{scalar}&\text{monomial}\\ \hline
(0,0)&4&1&1\\
(1,0)&2&2&q^{1/8}\\
(0,1)&2&2&s^{1/8}\\
(1,1)&2&2&q^{1/8}r^{1/4}s^{1/8}.
\end{array}
\]
Multiplication gives
\[
 1^4\,2^2\,2^2\,2^2\,
 q^{1/2}r^{1/2}s^{1/2}
 =64q^{1/2}r^{1/2}s^{1/2}.
\]
\end{proposition}
\begin{proof}
Insert each characteristic into the genus-two theta series and retain the unique terms of least order in the lexicographic cusp valuation $(m,n,-\ell)$. For $a=(0,0)$ the zero lattice vector contributes one. For each nonzero $a$, the two shortest lattice representatives are exchanged by sign and contribute the scalar two. Their quadratic exponents are the entries in the final column. Summing the exponents and multiplying the six factors of two gives the formula.
\end{proof}

\begin{proposition}[Boundary factor of the Borcherds product]
The $m=0$ part of the monic product is
\[
 \prod_{n>0,\ell\in\Z}(1-q^nr^\ell)^{f(0,\ell)}
 \prod_{\ell<0}(1-r^\ell)^{f(0,\ell)}
 =(1-r^{-1})\prod_{n\ge1}
 (1-q^nr^{-1})(1-q^n)^{10}(1-q^nr).
\]
Consequently the first Fourier--Jacobi coefficient is
\[
 q^{1/2}r^{1/2}(1-r^{-1})
 \prod_{n\ge1}(1-q^nr^{-1})(1-q^n)^{10}(1-q^nr)
 =\eta(\tau)^9\vartheta_{11}(\tau,z),
\]
where the odd theta function is fixed by the exact convention
\[
 \vartheta_{11}(\tau,z)
 =-q^{1/8}r^{-1/2}
 \prod_{n\ge1}(1-q^{n-1}r)(1-q^nr^{-1})(1-q^n).
\]
\end{proposition}
\begin{proof}
The $q^0$ row of $\phiweak$ has exactly three nonzero coefficients:
$f(0,-1)=1$, $f(0,0)=10$, and $f(0,1)=1$. Substitution into the $m=0$ factors gives the first formula. The Jacobi triple product and $\eta(\tau)=q^{1/24}\prod_{n\ge1}(1-q^n)$ give the second, including its leading sign and monomial.
\end{proof}

\begin{corollary}[The scalar square]
The monic weight-ten form is
\[
 \ChiTen=\Dmon^2=2^{-12}\Dtheta^2.
\]
Its automorphic character is trivial.
\end{corollary}
\begin{proof}
The identity $\Dtheta=2^6\Dmon$ gives the scalar factor.  The character of $\Dmon$ has order two, so its square is trivial.
\end{proof}

\chapter{Divisor, character, and the meaning of the square root}
The diagonal Humbert divisor is
\[
 H_1=\set{Z\in\HHtwo:z=0}/\operatorname{Sp}_4(\Z).
\]

\begin{theorem}[Diagonal divisor]
The divisor of $\Dmon$ is the $\operatorname{Sp}_4(\Z)$-orbit of $H_1$, with multiplicity one. The form transforms by the Maass character $\nu_5$ of order two, and $\ChiTen=\Dmon^2$ is scalar-valued.
\end{theorem}
\begin{proof}
This is the defining diagonal-divisor property of the Igusa square root and the first row of the Cl\'ery--Gritsenko classification \cite{Igusa62,Igusa64,GC}.
\end{proof}

Real roots of the denominator algebra define Weyl reflections. Imaginary roots contribute root spaces and product exponents. This division of labor is intrinsic to a Borcherds algebra: the Weyl group is generated by real simple roots, while the imaginary simple roots enter the correction term of the denominator identity.

\section{Automorphic section, orientation line, and Pfaffian operator}
An automorphic square root is a section of a line bundle with character. An orientation is a square root of a determinant line. A Pfaffian operator is a skew family whose regularized Pfaffian is a section of that orientation line. These form a sequence of successive refinements:
\[
\begin{gathered}
 \text{automorphic section}
 \longrightarrow \text{determinant-line identification}\\
 \longrightarrow \text{oriented square-root line}
 \longrightarrow \text{Fredholm family and Pfaffian section}.
\end{gathered}
\]

\begin{proposition}[Logical separation of the four levels]
The product theorem constructs the first object in the sequence. A determinant-line comparison constructs the second. A coherent square root of that line constructs the third. A first-order operator requires a symbol, domain, ellipticity or Fredholm condition, and regularization, and thereby produces the fourth refinement.
\end{proposition}
\begin{proof}
Each stage adds a specified geometric structure. A determinant realization adds an isomorphism from a geometric determinant line to the automorphic line. An oriented realization adds a coherent square-root line and its parity refinement. An operator realization adds a principal symbol, a Fredholm domain, a spectrum, and a regularization whose Pfaffian is the resulting section.
\end{proof}

The title names the monic automorphic section $\Dmon$, characterized by the exact identity $\Dmon^2=\ChiTen$. Part~IV formulates an analytic Pfaffian realization as a further geometric refinement.

\part{The denominator algebra}

\chapter{The Lorentzian root lattice}
Let
\[
 L=\Z f_2\oplus\Z f_3\oplus\Z f_{-2}
\]
with bilinear form
\[
 (f_2,f_{-2})=-1,
 \qquad (f_3,f_3)=2,
\]
and all remaining basis pairings zero. Its signature is $(2,1)$. Define the type-II root lattice
\[
 L_{II}=2\Z f_2\oplus\Z f_3\oplus2\Z f_{-2}
 \cong U(4)\oplus\angles{2}.
\]
The Gram-index map is
\[
 \alpha:\Gammaind=\Z^3\longrightarrow L_{II},
 \qquad
 \alpha(n,\ell,m)=2nf_2-\ell f_3+2mf_{-2}.
\]
Then
\[
 (\alpha(n,\ell,m),\alpha(n',\ell',m'))
 =-2\bigl(2nm'+2n'm-\ell\ell'\bigr),
\]
and
\[
 (\alpha(n,\ell,m),\alpha(n,\ell,m))
 =-2(4nm-\ell^2).
\]

\section{Real simple roots and Weyl vector}
Put
\[
 \delta_1=2f_2-f_3,
 \qquad
 \delta_2=2f_{-2}-f_3,
 \qquad
 \delta_3=f_3.
\]

\begin{proposition}[Type-II Cartan matrix]
The Gram matrix of the real simple roots is
\[
 \bigl((\delta_i,\delta_j)\bigr)_{i,j=1}^3
 =\begin{pmatrix}
 2&-2&-2\\
 -2&2&-2\\
 -2&-2&2
 \end{pmatrix}.
\]
The reflections
\[
 s_i(x)=x-(x,\delta_i)\delta_i
\]
are integral involutions preserving $L_{II}$ and its form.
\end{proposition}
\begin{proof}
Substitution into the basis pairing gives the matrix. Since $(\delta_i,\delta_i)=2$, the reflection formula has integral coefficients. Direct expansion gives $s_i^2=1$ and $(s_ix,s_iy)=(x,y)$.
\end{proof}

The Weyl vector is
\[
 \rho=f_2-\frac12f_3+f_{-2},
 \qquad (\rho,\delta_i)=-1.
\]
In the Gram-index variables, $e^{-\rho}=q^{1/2}r^{1/2}s^{1/2}$.

\chapter[The denominator algebra]{The Gritsenko--Nikulin Lie superalgebra}
The Gritsenko--Nikulin construction specifies a generalized
Kac--Moody Lie superalgebra \(\gIg\), with real simple roots
\(\delta_1,\delta_2,\delta_3\) and the imaginary correction
attached to \(\phiweak\). Its generators and relations supply the
bracket whose denominator is compared with the Borcherds product.

\begin{theorem}[Root supermultiplicities and denominator]
On the active positive support,
\[
 \smult\gIg_{\alpha(n,\ell,m)}=f(nm,\ell).
\]
In the Lorentzian coordinate convention of Gritsenko--Nikulin,
\[
 \den(\gIg)=\Dmon(2Z)=2^{-6}\Dtheta(2Z).
\]
\end{theorem}
\begin{proof}
The real roots and Weyl chamber are the type-II case of the Gritsenko--Nikulin construction. The imaginary simple-root correction is chosen so that its Weyl--Kac--Borcherds product has exponent $f(nm,\ell)$. The resulting product is the doubled-variable monic Igusa form \cite{GN97,GN98a,GN98b}.
\end{proof}

The theorem gives a Lie superalgebra with a bracket, Cartan, Weyl group, root parity, and denominator. The scalar product alone records the signed dimensions of the root spaces; the Gritsenko--Nikulin construction supplies the algebraic operations.

\chapter{The ordered bar as the denominator complex}
Let $\nIg$ be the positive nilpotent subalgebra of $\gIg$. Fix a proper integral height $h$ on the positive root cone and define
\[
 F^{>H}\nIg=\bigoplus_{h(\alpha)>H}\gIg_\alpha,
 \qquad
 \nIg^{(H)}=\nIg/F^{>H}\nIg.
\]
Height additivity makes $F^{>H}\nIg$ an ideal. Properness and local finiteness make $\nIg^{(H)}$ finite-dimensional and nilpotent.

\begin{theorem}[Height-complete Bar--Chevalley comparison]
For every $H$, antisymmetrization gives a natural quasi-isomorphism of conilpotent dg coalgebras
\[
 \Barc\bigl(U(\nIg^{(H)})\bigr)\simeq C_*(\nIg^{(H)}).
\]
The quotient maps preserve these comparisons and induce a quasi-isomorphism in the separated height-complete category
\[
 \Barc_h^\wedge\bigl(U(\nIg)\bigr)
 :=\varprojlim_H\Barc\bigl(U(\nIg^{(H)})\bigr)
 \simeq
 C_{*,h}^\wedge(\nIg)
 :=\varprojlim_H C_*(\nIg^{(H)}).
\]
Its root-graded Euler supercharacter is
\[
 \chi_{\mathrm{root}}\,C_{*,h}^\wedge(\nIg)
 =\prod_{\alpha>0}
 (1-e^{-\alpha})^{\sdim\gIg_\alpha}.
\]
\end{theorem}
\begin{proof}
For each $H$, the antisymmetrization map sends $\Sym(\nIg^{(H)}[1])$ to the normalized bar of $U(\nIg^{(H)})$. The PBW filtrations have the same associated-graded Koszul resolution of $\Sym(\nIg^{(H)})$, so antisymmetrization is a filtered quasi-isomorphism. Naturality gives a morphism of inverse systems. Every fixed root degree is visible in all sufficiently large height quotients and is unchanged thereafter; hence the inverse-limit map is a quasi-isomorphism coefficientwise. An even root vector becomes odd after the Chevalley shift and contributes $1-e^{-\alpha}$; an odd root vector becomes even and contributes its reciprocal. Multiplication over the finite-dimensional root spaces gives the product.
\end{proof}

Let $\mathbb T_{L_{II}}$ be the formal root torus and write $Z_{\mathrm{root}}$ for its logarithmic coordinate, normalized by
\[
 e^{-\alpha(n,\ell,m)}=Q^nR^\ell S^m,
 \qquad
 e^{-\rho}=Q^{1/2}R^{1/2}S^{1/2}.
\]

\begin{corollary}[Bar denominator identity]
On the formal root torus,
\[
 e^{-\rho}\chi_{\mathrm{root}}
 C_{*,h}^\wedge(\nIg)
 =\Dmon(Z_{\mathrm{root}}).
\]
The Gritsenko--Nikulin embedding of the root torus into the Siegel period domain sends $Z_{\mathrm{root}}$ to $2Z$ \cite{GN97,GN98b}. Hence
\[
 e^{-\rho}\chi_{\mathrm{root}}
 C_{*,h}^\wedge(\nIg)
 =\den(\gIg)=\Dmon(2Z)
\]
in Siegel period coordinates.
\end{corollary}
\begin{proof}
Multiply the Euler product in the preceding theorem by the Weyl monomial $e^{-\rho}$.  The denominator theorem identifies the result on the root torus, and the stated embedding gives the Siegel-coordinate formula.
\end{proof}

The denominator is therefore the Weyl-shifted Euler supercharacter of the height-complete interval bar of the positive enveloping algebra.

\chapter{Finite-height currents}
Retain the finite-height quotients $\nIg^{(H)}$ of the preceding chapter.  Define the current Lie conformal superalgebra
\[
 \Cur(\nIg^{(H)})=\kfield[\partial]\tensor\nIg^{(H)},
 \qquad [a_\lambda b]=[a,b],
\]
and its universal chiral enveloping algebra
\[
 V_H=U^{\ch}\!\left(\Cur_E(\nIg^{(H)})\right).
\]

\begin{theorem}[Height-complete current envelope]
The quotient maps induce a coefficientwise inverse limit
\[
 \widehat V_{\Delta_5}=\varprojlim_HV_H.
\]
Every $V_H$ is PBW-admissible as an associative chiral algebra, the current generators are primitive for the standard current coproduct, and
\[
 \sch\widehat V_{\Delta_5}
 =\prod_{j\ge1}\prod_{\alpha>0}
 (1-t^je^{-\alpha})^{-\sdim\gIg_\alpha}.
\]
Moreover,
\[
 \BarAssCh(V_H)
 \simeq
 \BarLieCh\!\left(\Cur_E(\nIg^{(H)})\right),
\]
and these equivalences pass to the separated height completion.
\end{theorem}
\begin{proof}
PBW orders the negative modes $a(-j)$, $j\ge1$.  Even modes contribute symmetric powers and odd modes contribute exterior powers, giving the product.  Proper height makes every coefficient stationary.  The current coproduct sends $a(z)$ to $a(z)\tensor1+1\tensor a(z)$.  Chiral antisymmetrization identifies the associative-chiral bar of the chiral enveloping algebra with the chiral Chevalley coalgebra; the complete separated PBW filtration is coefficientwise finite.
\end{proof}

\section{Classical local theory and exact quantum zero modes}
For every $H$, the finite-dimensional nilpotent Lie superalgebra $\nIg^{(H)}$ defines a cotangent BF theory on $E\times\R$.  Its local dg Lie algebra is
\[
 \Lcal_H=\Omega^{0,*}(E)\widehattensor\Omega^*(\R)\tensor\nIg^{(H)}
\]
and the classical action is
\[
 S_H(A,B)=\int_{E\times\R}\Omega_E\wedge
 \left\langle B,(\dbar+\dd_{\R})A+\frac12[A,A]\right\rangle .
\]

\begin{theorem}[Finite-height BF realization]
The action $S_H$ satisfies the classical master equation and defines a classical holomorphic--topological factorization algebra.  Harmonic reduction along a compact interval produces the finite-dimensional cotangent BV theory of $\nIg^{(H)}$.  The positive-root height grading makes every even adjoint endomorphism strictly height-increasing and every odd adjoint endomorphism odd; hence the modular supercharacter vanishes and the reduced theory satisfies the quantum master equation.

Choose the standard complementary Lagrangian boundary polarizations, fixing the $\nIg^{(H)}[1]$ coordinate at one endpoint and its cotangent coordinate at the other.  The algebraic boundary quantization, the interval amplitude between the two augmentation states, and the holomorphic boundary current envelope are then
\[
 U(\nIg^{(H)}),\qquad
 \mathbbm1\tensor^{\mathbb L}_{U(\nIg^{(H)})}\mathbbm1
 \simeq C_*(\nIg^{(H)}),\qquad V_H.
\]
The height quotient maps preserve the polarizations and yield the inverse systems
\[
 \widehat U_h(\nIg),\qquad
 C_{*,h}^{\wedge}(\nIg),\qquad
 \widehat V_{\Delta_5}.
\]
\end{theorem}
\begin{proof}
The classical master equation is the Jacobi identity paired with invariance of the cotangent pairing.  Elliptic descent gives the classical factorization algebra.  Harmonic reduction along the interval gives a finite-dimensional BV space.  If $x$ is even of positive root height, $\operatorname{ad}_x$ raises height and is nilpotent on the finite quotient, so its restrictions to the even and odd summands both have trace zero; if $x$ is odd, $\operatorname{ad}_x$ is odd and has supertrace zero.  Thus the modular character vanishes and the finite BV quantum master equation follows.

The chosen boundary polarization quantizes the linear Kirillov--Kostant boundary bracket to $U(\nIg^{(H)})$.  Gluing the two augmentation states gives the two-sided derived bar, and the PBW antisymmetrization identifies it with $C_*(\nIg^{(H)})$.  Holomorphic descendants generate $U^{\mathrm{ch}}(\operatorname{Cur}_E\nIg^{(H)})=V_H$.  All maps preserve root height, so they pass coefficientwise to the separated inverse limits.
\end{proof}

\part{The compact scalar character}

\chapter{Reduced curve counting on \texorpdfstring{$K3\times E$}{K3 x E}}
Let $S$ be a nonsingular projective K3 surface and $E$ an elliptic curve. For a primitive class $\beta_h\in H_2(S,\Z)$ with $\beta_h^2=2h-2$, let $N_{g,h,d}$ be the reduced disconnected Gromov--Witten invariant of class $(\beta_h,d)$ with the standard incidence insertion. Define
\[
 \mathcal Z(u,q,s)=
 \sum_{g,h,d}N_{g,h,d}u^{2g-2}q^{h-1}s^{d-1}.
\]

\begin{theorem}[Primitive Igusa cusp-form theorem]
Under $r=e^{iu}$, the partition function is the Laurent expansion
\[
 \mathcal Z(u,q,s)=-\frac1{\ChiTen(r,q,s)}.
\]
Equivalently, in theta normalization,
\[
 \mathcal Z=-2^{12}\Dtheta^{-2}.
\]
\end{theorem}
\begin{proof}
Oberdieck and Pixton identify the primitive reduced partition function with the Laurent expansion of $-\ChiTen^{-1}$ \cite{OP}; their normalization is the monic Borcherds product $\ChiTen$. The second formula is the identity $\ChiTen=2^{-12}\Dtheta^2$.
\end{proof}

The theorem identifies the numerical generating series in its stated
primitive sector. The Borcherds root algebra supplies a separate
algebraic character to compare with that series.

\begin{corollary}[Exact character comparison]
Let $\jmath:\HHtwo\to\HHtwo$ be the half-coordinate map $\jmath(Z)=Z/2$, and let
\[
 \Fock(\nIg)=
 \Sym(\nIg_{\bar0})\tensor\bigwedge(\nIg_{\bar1}).
\]
Then
\[
 \mathcal Z^{\red}_{S\times E}
 =-\jmath^*\!\left(
 e^{2\rho}\sch\bigl(\Fock(\nIg)^{\tensor2}\bigr)
 \right)
 =-\ChiTen^{-1}.
\]
\end{corollary}
\begin{proof}
Put
\[
 P=\prod_{\alpha>0}(1-e^{-\alpha})^{\sdim\gIg_\alpha}.
\]
The denominator theorem gives $\den(\gIg)=e^{-\rho}P=\Dmon(2Z)$, while PBW gives $\sch\Fock(\nIg)=P^{-1}$.  Hence
\[
 e^{2\rho}\sch\bigl(\Fock(\nIg)^{\tensor2}\bigr)
 =\den(\gIg)^{-2}.
\]
Pullback by $\jmath$ sends $\den(\gIg)$ to $\Dmon(Z)$, whose square is $\ChiTen$.  The reduced incidence convention supplies the overall minus sign in the Oberdieck--Pixton theorem.
\end{proof}

The equality compares the completed root-Fock supercharacter with the
numerical reduced generating series. Its proof uses the denominator
identity and the half-coordinate pullback. A map from compact geometric
states to this Fock space, preserving operators and their trace,
is an additional comparison.

\chapter{Extensions over the elliptic classifying stack}
Work on locally noetherian complex schemes with the fppf topology.
Let \(BE\) be the stack of right \(E\)-torsors. For a left
\(E\)-stack \(M\), an object of \([M/E](T)\) is a right
torsor \(P\to T\) and an equivariant map \(P\to M\).
Forgetting the map gives \([M/E]\to BE\). All products of
stacks below are two-fibre products.

\begin{proposition}[The common translation torsor]
\label{igusa:platonic-common-torsor}
There is a natural equivalence
\[
 [(M\times N)/E_{\mathrm{diag}}]
 \simeq[M/E]\times_{BE}[N/E].
\]
The equivalences for any finite number of factors commute with
regrouping and insertion of the terminal \(E\)-stack.
\end{proposition}
\begin{proof}
An object on the right is \((P,m),(Q,n)\) together with a
torsor isomorphism \(\alpha:P\to Q\). Send it to
\((P,m,n\alpha)\). Conversely, use a common torsor twice
and its identity isomorphism. A pair of torsor maps \(a,b\)
on the right satisfies \(b\alpha=\alpha'a\), so \(b\)
is determined by \(a\). The same construction identifies the
maps and their descent isomorphisms. It commutes with base change.
For several factors, every regrouping gives one torsor with the
same list of equivariant maps, proving all stated coherences.
\end{proof}

For \(M=N=E\) with translation, the diagonal quotient is
\(E\), with coordinate \(y-x\). Both separate quotients
are points. Their fibre product over \(BE\) retains the
relative translation, while their ordinary product loses it.

Let \(\mathfrak M\) be the stack of coherent sheaves on
\(X=S\times E\), with families flat over the parameter
scheme, and put \(\mathfrak Q=[\mathfrak M/E]\).
An object of \(\mathfrak Q(T)\) is a torsor \(P\to T\)
and a flat sheaf on \(S\times P\). Fix a polarization and
write \(\mathfrak Q_\alpha\) for a Hilbert-polynomial component.
Translation preserves these polynomials. Let
\(\mathfrak H_{\alpha,\beta}\) classify short exact sequences
\(0\to A\to F\to B\to0\) on a common \(S\times P\),
with all terms flat over \(T\) and polynomials
\(\alpha,\alpha+\beta,\beta\). Its span over \(BE\) is
\[
 \mathfrak Q_\alpha\times_{BE}\mathfrak Q_\beta
 \xleftarrow{p}\mathfrak H_{\alpha,\beta}
 \xrightarrow{q}\mathfrak Q_{\alpha+\beta}.
\]
The category of spans of fppf stacks over \(BE\) has composition
by fibre product and tensor product \(\times_{BE}\). Its unit
object is the identity map of \(BE\).

\begin{theorem}[Geometric Hall correspondence]
\label{igusa:platonic-common-torsor-hall}
For fixed \(\alpha,\beta\), \(q\) is representable by
algebraic spaces and proper. The spans form a coherently associative
unital Hall object in the category of spans over \(BE\).
Its unit is the zero-sheaf span
\(BE\leftarrow BE\to\mathfrak Q\).
\end{theorem}
\begin{proof}
Over a fixed \((P,F)\), the fibre of \(q\) classifies
quotients \(F\twoheadrightarrow B\) of polynomial \(\beta\).
Its kernel is flat, because \(F\) and \(B\) are flat.
An automorphism fixing the quotient and \(F\) is the identity.
Trivializing \(P\) fppf-locally identifies this fibre with the
fixed-polynomial Quot scheme, projective over the parameter scheme
by Grothendieck's theorem \cite[Theorem~3.2]{IgusaGrothendieckQuot1961}.
The quotient functor descends to an algebraic space, and properness
descends with it.

Both composites of three Hall spans classify
\(0\subset A\subset D\subset F\) on the same torsor.
One bracketing concatenates inclusions. The other starts with
\(A\subset F\) and a subobject \(B\subset F/A\), and sets
\(D=F\times_{F/A}B\). These constructions are inverse on
objects and isomorphisms. They preserve flatness and commute with
base change. For longer flags all regroupings identify the same
subobjects and quotients. Their comparison maps agree, proving the
pentagon and higher coherences. Finally \(\mathfrak Q_0=BE\).
The sequences with zero subobject or zero quotient give the unit
and its coherences.
\end{proof}

Properness here concerns the full coherent-sheaf stack on a fixed
pair of Hilbert-polynomial components. A stability restriction can
replace its fibre by an open part of a Quot scheme, which need not
be proper. A bounded charge window also needs closure under the
extensions and flags used in its operations. A reduced critical-coefficient realization must supply pullback,
pushforward, and orientation comparisons on these same flags. Its unit has
stabilizer \(E\). A coefficient theory for this quotient must
include that nonaffine stabilizer, or work relative to \(BE\),
on the unit and every flag.


\chapter{Additive physical charge and quadratic automorphic degree}
Let
\[
 \Lam=\widetilde H(S,\Z)
\]
be the Mukai lattice of signature $(4,20)$. The electric--magnetic charge lattice is
\[
 \Gammphys=\Lam\oplus\Lam,
 \qquad c=(Q,P).
\]
Its addition is the physical superselection law. Define the Gram map
\[
 \PiGram(Q,P)=
 \left(\frac{Q^2}{2},Q\cdot P,\frac{P^2}{2}\right)
 \in\Gammaind.
\]

\begin{proposition}[Polarization identity]
For $c_i=(Q_i,P_i)$,
\[
 \PiGram(c_1+c_2)=\PiGram(c_1)+\PiGram(c_2)+B(c_1,c_2),
\]
where
\[
 B(c_1,c_2)=
 \bigl(Q_1\cdot Q_2,
 Q_1\cdot P_2+Q_2\cdot P_1,
 P_1\cdot P_2\bigr).
\]
The cross-term is a symmetric vector-valued two-cocycle:
\[
 B(c_1,c_2)+B(c_1+c_2,c_3)
 =B(c_2,c_3)+B(c_1,c_2+c_3).
\]
\end{proposition}
\begin{proof}
Expand the three quadratic expressions. The cocycle identity is bilinearity of $B$.
\end{proof}

\section{The polarized central extension}
Define the abelian central extension
\[
 0\longrightarrow\Gammaind\longrightarrow
 \widetilde\Gamma_\Pi
 \longrightarrow\Gammphys\longrightarrow0,
 \qquad
 \widetilde\Gamma_\Pi=\Gammphys\times\Gammaind,
\]
with multiplication
\[
 (c,\lambda)\cdot(d,\mu)
 =\bigl(c+d,\lambda+\mu+B(c,d)\bigr).
\]
The cocycle identity makes this multiplication associative. The polarization identity makes the graph map
\[
 j_\Pi:\Gammphys\longrightarrow\widetilde\Gamma_\Pi,
 \qquad c\longmapsto(c,\PiGram(c))
\]
a homomorphism.

Let $A=\bigoplus_{c\in\Gammphys}A_c$ be a charge-graded associative algebra, and let $\kfield[\Gammaind]$ have monomials $u^\lambda$. Define the polarized central-extension algebra
\[
 \mathcal R_\Pi(A)=A\tensor\kfield[\Gammaind]
\]
with product
\[
 (a_cu^\lambda)\star(b_du^\mu)
 =a_cb_d\,u^{\lambda+\mu+B(c,d)}.
\]
Give $a_cu^\lambda$ degree $(c,\lambda)\in\widetilde\Gamma_\Pi$.

\begin{theorem}[Graph-multiplicative polarization]
The product $\star$ is associative and $\widetilde\Gamma_\Pi$-graded. The map
\[
 \iota_\Pi:A\longrightarrow\mathcal R_\Pi(A),
 \qquad a_c\longmapsto a_cu^{\PiGram(c)}
\]
is an algebra homomorphism whose homogeneous degree is the graph element $j_\Pi(c)$.
\end{theorem}
\begin{proof}
The two parenthesizations differ by
\[
 B(c,d)+B(c+d,e)-B(d,e)-B(c,d+e),
\]
which vanishes by the cocycle identity. The polarization formula gives
\[
 \iota_\Pi(a_c)\star\iota_\Pi(b_d)
 =a_cb_d\,u^{\PiGram(c+d)}.
\]
Thus the graph of the quadratic Gram map is multiplicative inside the central extension.
\end{proof}

The same construction applies to Lie superalgebras and their enveloping algebras. It retains additive physical charge, Gram degree, and the interaction cross-term in one grading group.

\chapter{Point-local and wrapped sectors}
Let $p_E:S\times E\to E$ be projection.

\begin{lemma}[Positive elliptic degree is wrapped]
Let $C\subset S\times E$ be a proper one-dimensional subscheme whose pushforward $(p_E)_*[C]$ has positive degree. Then $p_E(C)=E$.
\end{lemma}
\begin{proof}
Properness makes $p_E(C)$ closed. Positive pushforward degree implies that at least one irreducible component of $C$ dominates a one-dimensional closed subset of $E$. Every one-dimensional closed subset of the irreducible curve $E$ is $E$ itself, so the image of that component, and hence the image of $C$, is all of $E$.
\end{proof}

Projection-finite objects define point-supported sectors on $\Ran(E)$. Positive elliptic degree defines wrapped sectors supported over the entire elliptic curve.

\begin{definition}[The local--wrapped factorization operad]
Let $\mathsf{Ch}_E$ denote the chosen chiral/factorization operad on $E$ and let $\mathsf{Ass}$ denote the ordered associative operad. Define the colored operad $\mathsf{Hyb}_E$ with colors
\[
 \mathsf l,\qquad \mathsf w_m\quad(m\ge1)
\]
by the following operation objects.
\begin{enumerate}
\item Operations with local output are
\[
 \mathsf{Hyb}_E(\mathsf l^{\otimes k};\mathsf l)=\mathsf{Ch}_E(k).
\]
\item For $r\ge1$ and $m=m_1+\cdots+m_r$, operations with wrapped output are
\[
 \mathsf{Hyb}_E
 (\mathsf l^{\otimes k},\mathsf w_{m_1},\ldots,\mathsf w_{m_r};\mathsf w_m)
 =\mathsf{Ch}_E(k)\tensor\mathsf{Ass}(r).
\]
All other operation objects are initial.
\end{enumerate}
Operadic composition uses chiral insertion and disjoint union in the first factor and ordered substitution in the second. The equality
\[
 (\gamma_1\circ\gamma_2)\tensor(a_1\circ a_2)
 = (\gamma_1\tensor a_1)\circ(\gamma_2\tensor a_2)
\]
is the mixed interchange law. The arities $(k,r)=(k,0),(0,2),(k,1)$ recover, respectively, local factorization, wrapped multiplication, and the action of local insertions on a wrapped sector.
\end{definition}

\begin{definition}[Hybrid local--wrapped algebra]
A hybrid local--wrapped algebra over $E$ is an algebra over $\mathsf{Hyb}_E$. It consists of a factorization algebra $\F_{\loc}$, a height- or charge-complete family $W=\prod_{m>0}W_m$, ordered products
\[
 \mu_{m_1,\ldots,m_r}:W_{m_1}\tensor\cdots\tensor W_{m_r}
 \longrightarrow W_{m_1+\cdots+m_r},
\]
and factorization-compatible actions of the local operation objects $\mathsf{Ch}_E(k)$ on every wrapped product. The operad axioms state associativity, local factorization, and compatibility of the two operations.
\end{definition}

A compact Hall theory of $S\times E$ produces a correspondence-valued $\mathsf{Hyb}_E$-algebra when disjoint-support direct sums, wrapped extensions, and the mixed extension correspondences satisfy the same base-change identities.


\part{The Borcherds algebra and its quantum group}

\chapter{The local--wrapped root algebra}
The positive cone is graded by $\gamma=(n,\ell,m)$, with elliptic degree $m$.  Set
\[
 \nIg_{\loc}=\bigoplus_{\substack{\alpha>0\\m(\alpha)=0}}\gIg_\alpha,
 \qquad
 \nIg_{\wrap}=\bigoplus_{\substack{\alpha>0\\m(\alpha)>0}}\gIg_\alpha.
\]

\begin{proposition}[Semidirect local--wrapped decomposition]
The space $\nIg_{\loc}$ is a Lie subsuperalgebra, $\nIg_{\wrap}$ is a Lie ideal, and
\[
 \nIg=\nIg_{\loc}\ltimes\nIg_{\wrap}.
\]
The proper height filtration preserves both summands.
\end{proposition}
\begin{proof}
The bracket adds root degrees.  The sum of two roots of elliptic degree zero again has degree zero, and addition of a nonnegative degree to a positive degree remains positive.  Height is additive on every nonzero bracket.
\end{proof}

Define the complete enveloping Hopf algebra and current envelope
\[
 \mathbb U_{\Delta_5}=\widehat U_h(\nIg),
 \qquad
 \mathbb V_{\Delta_5,E}:=\widehat V_{\Delta_5}=\varprojlim_HU^{\ch}\!\left(\Cur_E(\nIg^{(H)})\right).
\]
The adjoint action of $\nIg_{\loc}$ on $\nIg_{\wrap}$ extends by derivations to $\widehat U_h(\nIg_{\wrap})$.  Let $\# $ denote the corresponding completed smash product.

\begin{definition}[Current-generated local--wrapped operad]
The colored operad $\mathsf{Hyb}^{\cur}_{E,\Delta}$ is generated by the chiral current operations of $\Cur_E(\nIg_{\loc})$, the ordered multiplications of $\widehat U_h(\nIg_{\wrap})$, and the adjoint derivation action of the former on the latter, modulo the chiral Jacobi, associative, derivation, and semidirect-product relations.
\end{definition}

\begin{theorem}[Local--wrapped Borcherds decomposition]
Let $\nIg_{\loc}$ be the sum of positive root spaces of elliptic degree $m=0$ and let $\nIg_{\wrap}$ be the sum of those with $m>0$.  Then
\[
 \nIg=\nIg_{\loc}\ltimes\nIg_{\wrap},
\]
where $\nIg_{\loc}$ is a Lie subsuperalgebra and $\nIg_{\wrap}$ is a Lie ideal.  For the proper-height completion,
\[
 \widehat U_h(\nIg)
 \cong
 \widehat U_h(\nIg_{\wrap})\#\widehat U_h(\nIg_{\loc}),
\]
and
\[
 \widehat C_{*,h}(\nIg)
 \simeq
 \Tot C_*(\nIg_{\loc};\widehat C_{*,h}(\nIg_{\wrap})).
\]
The primitive Lie superalgebra of the completed enveloping algebra is the height completion $\widehat{\nIg}_h=\varprojlim_H\nIg^{(H)}$, and the Weyl-shifted Euler supercharacter of the completed Chevalley complex is $\Dmon(2Z)$.
\end{theorem}
\begin{proof}
Root addition adds elliptic degree.  Thus the $m=0$ summand is closed under the bracket, and every bracket with an $m>0$ root again has $m>0$.  PBW orders wrapped generators before local generators and gives the completed smash product.  The Chevalley complex of a semidirect product is the Hochschild--Serre total complex.  Primitive elements of a completed enveloping algebra are recovered height by height.  The denominator theorem supplies the last character identity.
\end{proof}

\begin{remark}
The theorem constructs a Borcherds root algebra with local and wrapped colors.  The word Hall is reserved for an algebra obtained from geometric extension correspondences.
\end{remark}

\chapter{The formal Igusa quantum group}
Let $\mathfrak b_+$ be the positive Borcherds Borel half.  Form the restricted graded dual $\mathfrak b_-:=\mathfrak b_+^{\vee,\mathrm{gr}}$ with bracket and cobracket transported by the invariant pairing, identify the two Cartan copies along their common radical quotient, and complete by a proper root height.  The completed Drinfeld double of this paired Borel datum is denoted $\gIg^{\mathrm{pair}}$.  Its positive and negative halves, diagonal Cartan, invariant form, and canonical tensor form a height-admissible Manin triple and hence a quasitriangular Lie superbialgebra.

\begin{definition}[Height-admissible Lie superbialgebra]
A root-graded Lie superbialgebra is height-admissible when every root space is finite-dimensional, every bounded-height region contains finitely many roots, and the bracket, cobracket, invariant tensor, and all iterated operations are coefficientwise finite in the height completion.
\end{definition}

\begin{proposition}[Height admissibility]
The paired Gritsenko--Nikulin Borcherds Lie superbialgebra $\gIg^{\mathrm{pair}}$ is height-admissible.
\end{proposition}
\begin{proof}
The Borcherds root spaces are finite-dimensional and the chosen height is proper.  The bracket and cobracket preserve total root degree.  In a bounded total height, finitely many decompositions of a root occur.  The canonical tensor is therefore a coefficientwise finite sum in every height quotient.
\end{proof}

\begin{theorem}[Quantized Borcherds superalgebra]
Choose the symmetrizable Borcherds--Cartan superdatum and multiplicity index set that present $\gIg$.  Hong's construction gives a quantized enveloping Hopf superalgebra
\[
 U_q(\gIg)
\]
over $\Q(q)$, generated by its Cartan elements and the real and imaginary simple-root generators, with triangular decomposition
\[
 U_q(\gIg)\cong U_q^-(\gIg)\tensor U_q^0(\gIg)\tensor U_q^+(\gIg).
\]
Its PBW root spaces have the same graded ranks as the classical Borcherds superalgebra.  The canonical pairing of the two Borel halves yields a universal $R$-matrix in the root-height completion of
\[
 U_q^+(\gIg)\widehattensor U_q^-(\gIg).
\]
Consequently the quantization preserves the root-multiplicity data that enter $\Dmon(2Z)$.
\end{theorem}
\begin{proof}
The Gritsenko--Nikulin algebra is a symmetrizable Borcherds--Kac--Moody Lie superalgebra.  Apply the generators-and-relations construction of the quantized Borcherds superalgebra.  The nondegenerate bilinear form, triangular decomposition, centre, and universal $R$-matrix are established by Hong \cite{Hong98}.  The PBW theorem identifies the associated graded root spaces with their classical counterparts.  Since the denominator exponents are the classical signed root multiplicities, they are unchanged by the flat root-graded deformation.
\end{proof}

\begin{remark}
The quantum group is attached to the Borcherds generalized Cartan superdatum.  A geometric Hall quantum group is identified with it through a Hall coproduct, pairing, completion, and primitive comparison.
\end{remark}

\chapter{A holomorphic--topological BF realization}
For every height bound $H$, let $\nIg^{(H)}$ be the finite-dimensional nilpotent quotient.  On $E\times\R$ form the shifted cotangent theory of
\[
 \Lcal_H=\Omega^{0,*}(E)\widehattensor\Omega^*(\R)\tensor\nIg^{(H)}
\]
with action
\[
 S_H(A,B)=\int_{E\times\R}\Omega_E\wedge
 \left\langle B,(\dbar+\dd_{\R})A+\frac12[A,A]\right\rangle .
\]

\begin{theorem}[Classical local and exact zero-mode realization]
For every $H$:
\begin{enumerate}
\item $S_H$ satisfies the classical master equation and defines a classical holomorphic--topological factorization algebra;
\item harmonic reduction along a compact interval gives the finite-dimensional cotangent BV theory of $\nIg^{(H)}$, whose positive-root grading makes its modular supercharacter zero and hence gives the quantum master equation;
\item after choosing the complementary Lagrangian boundary polarizations, the boundary algebra is $U(\nIg^{(H)})$, the interval amplitude is $\mathbbm1\tensor^{\mathbb L}_{U(\nIg^{(H)})}\mathbbm1\simeq C_*(\nIg^{(H)})$, and the holomorphic boundary current envelope is $V_H$.
\end{enumerate}
The transition maps yield a pro-object with boundary algebra $\widehat U_h(\nIg)$, interval complex $C_{*,h}^{\wedge}(\nIg)$, and current envelope $\widehat V_{\Delta_5}$.
\end{theorem}
\begin{proof}
The classical master equation is the Jacobi identity and invariance of the cotangent pairing.  The classical observables satisfy factorization descent.  Harmonic reduction is finite-dimensional.  An even adjoint operator raises positive-root height and is nilpotent on the finite quotient; an odd adjoint operator is odd.  Both have supertrace zero, so the modular term in the finite BV Laplacian vanishes.  The chosen polarization quantizes the boundary Lie--Poisson bracket to the enveloping algebra; the two augmentation states give the two-sided bar, identified with Chevalley chains by PBW.  Holomorphic descendants give the current algebra.  All structure maps preserve height.
\end{proof}

The theorem gives the precise field-theoretic realization used in this volume: a classical local theory together with an exact quantum protected zero-mode sector.

\chapter{The determinant and Fredholm square-root operators}
Let
\[
 \Vcal_H=\bigoplus_{\substack{\alpha>0\\h(\alpha)\le H}}\gIg_\alpha
\]
and define
\[
 \mathscr D_H|_{\gIg_\alpha}=(1-e^{-\alpha})\id.
\]

\begin{theorem}[Universal determinant section]
The finite-height Berezinian is
\[
 \Ber(\mathscr D_H)
 =\prod_{\substack{\alpha>0\\h(\alpha)\le H}}
 (1-e^{-\alpha})^{\sdim\gIg_\alpha}.
\]
The coefficientwise limit satisfies
\[
 e^{-\rho}\Ber(\mathscr D_\infty)=\Dmon(2Z).
\]
After pullback by $\jmath(Z)=Z/2$,
\[
 \jmath^*\!\left(e^{-\rho}\Ber(\mathscr D_\infty)\right)=\Dmon(Z),
 \qquad \Dmon(Z)^2=\ChiTen(Z).
\]
\end{theorem}
\begin{proof}
The Berezinian of a diagonal endomorphism is the product of its eigenvalues with exponent equal to the superdimension of the eigenspace.  Proper height makes every coefficient stationary.  The denominator theorem supplies the Weyl line and the root-to-Siegel coordinate map.
\end{proof}

\section{Trace-class realization on the root Hilbert superspace}
Choose homogeneous orthonormal bases of the finite-dimensional root spaces and set
\[
 \Hcal_{\Delta}=\widehat\bigoplus_{\alpha>0}\gIg_\alpha
\]
with its induced $\Z/2$ grading.  Let $\mathfrak C_{\mathrm{abs}}$ be the subchamber on which
\[
 \sum_{\alpha>0}\dim\gIg_\alpha\,
       \abs{e^{-\alpha(Z)}}<\infty.
\]
For $Z\in\mathfrak C_{\mathrm{abs}}$, define the trace-class diagonal operator
\[
 K_Z|_{\gIg_\alpha}=e^{-\alpha(Z)}\id,
 \qquad D_Z=1-K_Z.
\]
On each parity summand form the skew doubling
\[
 \mathbb A_{Z,\bar\epsilon}=
 \begin{pmatrix}0&D_{Z,\bar\epsilon}\\-D_{Z,\bar\epsilon}^{\mathsf t}&0\end{pmatrix}.
\]
Orient every finite-height truncation by the fixed root order and normalize its Pfaffian to $1$ at $K_Z=0$.  The oriented Fredholm super-Pfaffian is the stabilized ratio
\[
 \operatorname{sPf}_{F}(\mathbb A_Z)
 =\frac{\Pf_F(\mathbb A_{Z,\bar0})}
        {\Pf_F(\mathbb A_{Z,\bar1})}.
\]

\begin{theorem}[Polarized Fredholm realization]
Choose the protected parity refinement and the root-space Hilbert norms above.  For every $Z$ in the absolute-convergence chamber, the diagonal operator $K_Z|_{\gIg_\alpha}=e^{-\alpha(Z)}\id$ is trace class on the even and odd root Hilbert spaces.  For $D_Z=1-K_Z$,
\[
 \sdet_F(D_Z)
 =\frac{\det_F(D_Z|_{\bar0})}{\det_F(D_Z|_{\bar1})}
 =\prod_{\alpha>0}(1-e^{-\alpha(Z)})^{\sdim\gIg_\alpha}.
\]
After choosing the standard polarization of $\Hcal_\Delta\oplus\Hcal_\Delta^*$ and its orientation, the paritywise skew block operators $\mathbb A_{Z,\bar0}$ and $\mathbb A_{Z,\bar1}$ have an oriented Fredholm super-Pfaffian
\[
 \operatorname{sPf}_F(\mathbb A_Z)=\frac{\Pf_F(\mathbb A_{Z,\bar0})}{\Pf_F(\mathbb A_{Z,\bar1})}
 =\sdet_F(D_Z).
\]  Hence
\[
 e^{-\rho}\sdet_F(D_Z)=\Dmon(2Z).
\]
\end{theorem}
\begin{proof}
The defining summability condition for \(\mathfrak C_{\mathrm{abs}}\)
is precisely the trace norm of the diagonal operator, summed in both
parities.  The Fredholm determinant of $1-K_Z$ is the convergent product over eigenvalues; taking the even/odd ratio gives the superdeterminant.  On each finite root-height truncation the Pfaffian of the oriented skew doubling equals the determinant of $D_Z$.  Trace-norm convergence of the truncations passes this identity to the Fredholm limit.  The final equality is the Gritsenko--Nikulin denominator formula.
\end{proof}

\begin{remark}
The polarization and orientation in this theorem belong to the root Hilbert space.  A compact moduli-space orientation is a separate line bundle and is compared with this one by a specified intertwiner.
\end{remark}

This constructs the square-root operator on the root Hilbert superspace
where its stated summability condition holds. The nonemptiness estimate
below concerns the minimal virtual superspace, whose dimensions are
\(\lvert\sdim\gIg_\alpha\rvert\). Applying that estimate to the
full root superspace requires a bound on both parity dimensions.
A compact geometric realization requires an operator intertwiner and
a compatible map of determinant lines.

\chapter{What the automorphic character determines}
\begin{theorem}[Denominator data and reconstruction data]
The Borcherds product determines:
\begin{enumerate}
\item the signed root multiplicity $\sdim\gIg_\alpha$ in every active root degree;
\item the Weyl vector and denominator character;
\item the divisor and automorphic multiplier of the section.
\end{enumerate}
A Lie bracket is supplied by the Borcherds--Kac--Moody presentation, parity dimensions by a parity-refined root datum, Hall multiplication by extension correspondences, and an oriented square root by a square-root line with monodromy.  These reconstruction data refine the scalar product independently.
\end{theorem}
\begin{proof}
The first three items are read from the exponent, prefactor, and automorphy law.  Independence is witnessed by explicit pairs: two brackets on the same graded vector space have the same PBW character; an even--odd pair has signed dimension zero; distinct Hall correspondences may have the same Euler product; and twisting a square-root line by a sign local system preserves its square.  Hence each refined object is reconstructed from the additional datum named in the statement.
\end{proof}

\begin{recognition}[Complete primitive recognition]
Let $\mathfrak p$ be a positive graded Lie superalgebra and let $\Psi:\nIg\twoheadrightarrow\mathfrak p$ be a graded Lie map.  The map is an isomorphism if either:
\begin{enumerate}
\item the parity-refined Hilbert series of $U(\nIg)$ and $U(\mathfrak p)$ agree; or
\item each primitive degree has prescribed pure parity and the signed PBW characters agree.
\end{enumerate}
The induced height completions, Chevalley complexes, and current envelopes are then isomorphic.
\end{recognition}
\begin{proof}
Choose the first height containing a homogeneous kernel vector.  All decomposable PBW terms in that height are determined by lower heights and agree.  The primitive parity dimension therefore drops, contradicting the first hypothesis.  In the second case the logarithm of the PBW product recovers signed primitive dimensions inductively, and prescribed parity recovers ordinary dimensions.  Naturality of completion, Chevalley chains, and current envelopes gives the final assertion.
\end{proof}

\part{The universal Borcherds Hall and three-Calabi--Yau models}

\chapter{Finite simple-root windows}
Let $I_\Delta$ be the multiset of simple roots of the Gritsenko--Nikulin Borcherds superalgebra.  A repeated imaginary simple root is written once for each copy.  The parity map is
\[
 p:I_\Delta\longrightarrow\Z/2.
\]
Choose a proper height function
\[
 h:I_\Delta\longrightarrow\Z_{>0}
\]
whose sublevel sets are finite.  Put
\[
 I_H=\set{i\in I_\Delta:h(i)\le H}.
\]
The restriction of the invariant form gives a finite symmetric Borcherds--Cartan matrix
\[
\begin{aligned}
 A_H&=(a_{ij})_{i,j\in I_H},\\
 a_{ii}&=2 &&\text{for real }i,\\
 a_{ii}&\leq0 &&\text{for imaginary }i,\\
 a_{ij}&\leq0 &&\text{for }i\neq j.
\end{aligned}
\]
The evenness of the root lattice makes every imaginary diagonal entry even.

\begin{construction}[The Igusa Borcherds quiver]
Order $I_H$.  For $i<j$ place $-a_{ij}$ arrows from $i$ to $j$.  At an imaginary vertex $i$ place
\[
 g_i=1-\frac{a_{ii}}2
\]
loops.  A real vertex has loop number zero.  Denote the resulting quiver by $Q_{\Delta,H}$.
\end{construction}

\begin{proposition}[Euler form]
The symmetrized Euler form of $Q_{\Delta,H}$ is $A_H$:
\[
 \angles{e_i,e_j}+\angles{e_j,e_i}=a_{ij}.
\]
\end{proposition}
\begin{proof}
For $i\ne j$, the total number of arrows between the two vertices is $-a_{ij}$, so the symmetrized Euler entry is $a_{ij}$.  At $i$ the value is $2-2g_i=a_{ii}$.
\end{proof}

The inclusions $I_H\subset I_{H+1}$ identify $Q_{\Delta,H}$ with a full subquiver of $Q_{\Delta,H+1}$.  Representations supported on the old vertices form an exact extension-closed subcategory.  This gives compatible Hall embeddings.


\chapter{The real-root seed window}
The three real simple roots are
\[
 \delta_1=2f_2-f_3,
 \qquad
 \delta_2=2f_{-2}-f_3,
 \qquad
 \delta_3=f_3,
\]
with Gram matrix
\[
 A_{\mathrm{re}}=
 \begin{pmatrix}
 2&-2&-2\\
 -2&2&-2\\
 -2&-2&2
 \end{pmatrix}.
\]
Choose the order $\delta_1<\delta_2<\delta_3$.  The real-root quiver $Q_{\mathrm{re}}$ has two arrows from vertex $i$ to vertex $j$ for each $i<j$.

\begin{proposition}[Real-root Hall seed]
The symmetrized Euler form of $Q_{\mathrm{re}}$ is $A_{\mathrm{re}}$.  The Hall composition algebra generated by its three simple representations is the positive quantum Kac--Moody algebra with Cartan matrix $A_{\mathrm{re}}$.
\end{proposition}
\begin{proof}
The diagonal Euler entries are $2$.  The total number of arrows between distinct vertices is two, hence the symmetrized off-diagonal entries are $-2$.  The Ringel--Green Hall theorem identifies the composition algebra with the positive quantum algebra defined by that symmetric Cartan matrix.
\end{proof}

The three vertices already contain the Weyl group and the reflective divisor.  Imaginary vertices add the automorphic correction.  The finite-height tower therefore begins with a completely explicit hyperbolic Hall algebra and enlarges it by one finite collection of imaginary simple roots at each stage.

\section{The first Ginzburg differential}
Let $a_{ij}^{(1)},a_{ij}^{(2)}:i\to j$ be the two arrows for $i<j$.  Add reverse arrows $(a_{ij}^{(r)})^*:j\to i$ of degree $-1$ and loops $t_i$ of degree $-2$.  The canonical three-Calabi--Yau completion has
\[
 d a_{ij}^{(r)}=0,
 \qquad
 d(a_{ij}^{(r)})^*=0,
\]
and
\begin{equation}\label{eq:vol3-ginzburg-differential}
 d t_i=
 \sum_{a:s(a)=i}aa^*
 -\sum_{a:t(a)=i}a^*a.
\end{equation}
The identity $d^2=0$ follows immediately.  Formula \eqref{eq:vol3-ginzburg-differential} is the derived moment-map equation for the doubled real-root quiver.

\begin{calculation}[Vertex one]
At vertex $1$,
\[
 d t_1=
 \sum_{r=1}^2
 \left(
 a_{12}^{(r)}(a_{12}^{(r)})^*
 +a_{13}^{(r)}(a_{13}^{(r)})^*
 \right).
\]
At vertex $3$,
\[
 d t_3=-\sum_{r=1}^2\left(
 (a_{13}^{(r)})^*a_{13}^{(r)}+
 (a_{23}^{(r)})^*a_{23}^{(r)}\right).
\]
The sum over vertices is $\sum_a(aa^*-a^*a)$.
It vanishes in the cyclic quotient and after matrix trace,
but generally not in the path algebra.
\end{calculation}
\begin{proof}
The outgoing arrows at vertex $1$ have targets $2$ and $3$.
The incoming arrows at vertex $3$ have sources $1$ and $2$.
The displayed formulas follow from these incidence lists.
Each arrow contributes one commutator to the sum over vertices.
The trace of a commutator vanishes, giving the asserted trace identity.
\end{proof}


\chapter{Hall realization of the positive algebra}
Let $\mathcal A_H=\Rep_{\mathbb F_q}(Q_{\Delta,H})$.  Its twisted Hall product is
\[
 [M]*[N]
 =v^{\angles{\dim M,\dim N}}
 \sum_{[L]}
 \abs{\set{N'\subset L:N'\cong N,\ L/N'\cong M}}
 [L],
 \qquad v^2=q.
\]

The class $[0]$ is the unit: the required subobject is unique on either
side. Associativity follows by counting flags $N_1\subset N_2\subset L$
in two orders. The twist factors agree because the Euler form is
bilinear:
\[
 \langle\alpha,\beta\rangle+\langle\alpha+\beta,\gamma\rangle
 =\langle\beta,\gamma\rangle+\langle\alpha,\beta+\gamma\rangle.
\]
Thus this normalization defines the unital twisted Hall algebra on
isomorphism classes.

\begin{theorem}[Finite-window composition algebra]
At each vertex $i$ of $Q_{\Delta,H}$, choose the one-dimensional
representation $S_i$ supported at $i$, with every loop acting by zero.
Use the generic composition algebra generated by these classes across
finite fields, with the twisted subobject-count multiplication above.
After extending coefficients from $\mathbb Z[v,v^{-1}]$ to
$\mathbb C(v)$, this is the positive quantum generalized Kac--Moody
algebra for the even Borcherds--Cartan matrix $A_H$.
The parameter in Kang--Schiffmann is $v_{\mathrm{KS}}=v^{-1}$.
Their specified semisimple perverse-sheaf construction on nilpotent
representation varieties gives the corresponding canonical basis.
\end{theorem}
\begin{proof}
This is the selected-simple composition construction of
Kang--Schiffmann, Section~2.3 and Theorem~2.2, with one selected
zero-loop simple at each vertex. Their geometric realization uses
its split graded Grothendieck group over
$\mathbb Z[v_{\mathrm{KS}},v_{\mathrm{KS}}^{-1}]$
(Theorem~4.1), with the stated inversion of parameter
\cite{KangSchiffmann}. The theorem does not use every simple object
of the full representation category. A vertex with $c$ loops has
$q^c$ one-dimensional simple representations over $\mathbb F_q$,
whereas only the chosen zero-loop representation is the generator here.
\end{proof}

A semi-derived Hall realization of a full quantum group additionally
specifies two-periodic complexes, the invertible acyclic classes and
Cartan reduction \cite{LuBorcherdsHall,AxtellLee}.
Its coproduct, pairing and completion are separate from the positive
composition algebra used here.

\section{Parity and the quantized superdatum}
The Gritsenko--Nikulin denominator algebra carries a parity map on its simple-root multiset.  It is additive on the free lattice of labelled simple generators.
Projection to the Lorentzian degree can identify contributions of
different parity. Both parity dimensions must therefore be retained
in each projected root space.  Hong's construction applied to this Borcherds superdatum gives the positive Borel
\[
 U_q^+(\gIg^{(H)})
\]
and the full quantized Borcherds superalgebra $U_q(\gIg^{(H)})$.

\begin{theorem}[Ordinary and super presentations]
The ordinary quiver Hall presentation and Hong's quantum
superalgebra use the same specified Cartan matrix and labelled
simple generators. Their relation ideals and parity conventions
are different inputs. Equality of their Cartan gradings alone does
not identify their PBW characters.
\end{theorem}
\begin{proof}
The two presentations are those of the preceding Hall construction
and Hong's superalgebra construction. The difference already occurs
for a single isotropic generator. The enveloping algebra of a
one-dimensional even abelian Lie algebra is \(\mathbb C[x]\).
Declaring its generator odd only as a graded vector space leaves
a nonzero vector \(x^2\) in degree two. The enveloping algebra
of a one-dimensional odd abelian Lie superalgebra is
\(\bigwedge(\xi)\), since its defining relation is
\(2\xi^2=[\xi,\xi]=0\). Its degree-two component is zero.
Thus parity assignment to an ordinary PBW basis does not produce
the super relation ideal or its character. A comparison must
respect these relations degree by degree.
\end{proof}

\chapter{The three-Calabi--Yau completion}
Let $kQ_{\Delta,H}$ be the path algebra. Its algebraic
three-Calabi--Yau completion is
\[
 \Gamma_{\Delta,H}=\Pi_3(kQ_{\Delta,H}).
\]
Write $J$ for the ideal generated by all arrows in its graded quiver.
Its completed version is the separate topological dg algebra
$\widehat\Gamma_{\Delta,H}=\varprojlim_n\Gamma_{\Delta,H}/J^n$.
The algebraic Calabi--Yau theorem is used for $\Gamma_{\Delta,H}$;
a topological Calabi--Yau assertion for $\widehat\Gamma_{\Delta,H}$
requires its completed bimodule category.
This is the Ginzburg model with zero potential: a degree-minus-one
reverse arrow for every arrow and a degree-minus-two loop at each
vertex. The loop differential is the signed sum of incident
arrow--reverse-arrow products.

\begin{theorem}[Borcherds--Ginzburg projections]
For each finite quiver \(Q_{\Delta,H}\), its undeformed
three-Calabi--Yau completion \(\Gamma_{\Delta,H}\) is the
Ginzburg dg algebra with zero potential. If \(H\le H'\) and
\(Q_{\Delta,H}\) is the full subquiver on a subset of vertices,
there is a surjective dg algebra map
\[
 \pi_{H'H}:\Gamma_{\Delta,H'}\longrightarrow\Gamma_{\Delta,H}.
\]
It kills the added vertices, every arrow incident to them, the
corresponding reverse arrows, and the added degree-minus-two loops.
These maps compose strictly and extend continuously to the completed algebras.
\end{theorem}
\begin{proof}
The finite Calabi--Yau assertion is the Ginzburg model of Keller's
completion theorem \cite{KellerCY,GinzburgCY}. We concatenate
paths from left to right, so \(aa^*\) is based at the source
of \(a\). Original and reverse arrows have zero differential.
For the degree-minus-two loop \(t_i\), the differential is
\[
 d t_i=\sum_{s(a)=i}aa^*-\sum_{t(a)=i}a^*a.
\]
The projection deletes exactly the terms involving an added vertex,
leaving the same sum for the full subquiver. It therefore commutes
with the differential. It preserves path length and extends to the
completed path algebra. Successive projections kill the union of
their omitted vertices and arrows, proving strict composition.
\end{proof}

The evident inclusion of graded quivers does not give the reverse
dg map. A single old vertex has loop differential zero. Adding
an arrow \(a:i\to j\) changes that differential to \(aa^*\)
at \(i\). The map fixing the old loop therefore fails to commute
with the differential.

Write \(\Gamma_\Delta=(\Gamma_{\Delta,H},\pi_{H'H})_H\)
for this pro-system of dg algebras. It is separate from the
ordinary Hall embeddings obtained by extension by zero on
representations of a full subquiver. A direct system of
Calabi--Yau dg categories compatible with those Hall embeddings
requires additional dg functors. It does not follow from the
inclusion of the displayed Ginzburg generators.

The quantum group retains its separately specified root-height
completion. Its transition maps must preserve the Cartan generators,
coproduct, and pairing. Quotienting the positive nilpotent Lie
algebra by height and truncating a full quantum group are different
operations.

\chapter{Hall states and protected superstates}
The representation varieties of $Q_{\Delta,H}$ carry the
perverse-sheaf complexes used in the composition-algebra construction.
The full equivariant derived category is a larger carrier.
Already for one real vertex and dimension one, its quotient is
$B\mathbb G_m$, with $H^*(B\mathbb G_m;\mathbb Q)=\mathbb Q[c_1]$,
where $|c_1|=2$. Its total equivariant cohomology is infinite-dimensional.
Finite root multiplicities therefore refer to the specified composition
algebra, not to total equivariant cohomology of this stack.

\begin{definition}[Igusa Hall category]
For every dimension vector $\mathbf d$, let
\[
 \mathfrak R_{H,\mathbf d}=
 [\operatorname{Rep}_{\mathbf d}(Q_{\Delta,H})/G_{\mathbf d}].
\]
The finite-height Hall category is the direct sum of the equivariant constructible derived categories on the $\mathfrak R_{H,\mathbf d}$, equipped with Hall convolution along the stack of short exact sequences.  Denote it by $\mathcal B_{\Delta,H}$ and set
\[
 \mathcal B_\Delta=\varinjlim_H\mathcal B_{\Delta,H}.
\]
\end{definition}

Independently, the protected one-particle superspace is the root-graded superspace
\[
 \mathcal H_{\mathrm{1p},H}
 =\bigoplus_{\alpha>0}\gIg^{(H)}_\alpha,
\]
with the parity of the Borcherds superalgebra.  Its multiparticle completion is the PBW/Fock superspace of $U(\nIg^{(H)})$.

\begin{theorem}[Hall and protected-state comparison]
The composition subalgebra in the finite-window Hall theorem
is the ordinary positive quantum generalized Kac--Moody algebra.
This statement does not identify the Grothendieck group of the full
category $\mathcal B_{\Delta,H}$ with that composition algebra.  The protected one-particle superspace has signed root multiplicities
\[
 \sdim\mathcal H_{\mathrm{1p},H,\alpha}
 =\sdim\gIg^{(H)}_\alpha.
\]
Its height-complete PBW character is
\[
 \sch\Fock(\nIg)
 =\prod_{\alpha>0}(1-e^{-\alpha})^{-\sdim\gIg_\alpha},
\]
and its Weyl-shifted inverse is $\Dmon(2Z)$.  The composition canonical basis and the protected root presentation use the specified simple-root indexing and Cartan grading; the parity refinement belongs to the Borcherds superdatum.
\end{theorem}
\begin{proof}
The composition-algebra assertion is the finite-window theorem. The protected assertion is PBW for the Borcherds Lie superalgebra.  Proper height passes both constructions to the limit.  The Gritsenko--Nikulin denominator theorem gives the final product identity.
\end{proof}

\chapter{Extension determinants and critical orientations}
Let \(\mathcal C\) be a three-Calabi--Yau dg category over
\(\mathbb C\), with its chosen trace. For families of objects
\(A,B\) over a base \(T\), let \(C(A,B)\) denote the relative
derived Hom complex on \(T\). Suppose these complexes are perfect,
and that the trace gives base-change-compatible duality
\[
 C(B,A)\simeq C(A,B)^\vee[-3].
\]
Put \(L(A,B)=\det C(A,B)\) and \(D(A)=L(A,A)\).
The determinant convention is
\(\det C=\bigotimes_i(\det C^i)^{(-1)^i}\), with symmetry sign
specified by \(\chi(C)\bmod2\).

\begin{theorem}[Extension determinant]
For an exact triangle \(A\to F\to B\to A[1]\), there is a
base-change-compatible isomorphism
\[
 D(F)\simeq D(A)\otimes D(B)\otimes L(A,B)^{\otimes2}.
\]
The determinant comparisons obtained by refining a flag agree, with
the stated determinant symmetry signs.
\end{theorem}
\begin{proof}
Apply derived Hom to the triangle in each variable. Determinant
additivity gives
\[
 D(F)\simeq L(B,A)\otimes D(A)\otimes D(B)\otimes L(A,B).
\]
Dualization and the odd shift each invert the determinant line.
Thus Calabi--Yau duality identifies \(L(B,A)\) with \(L(A,B)\).
Reorder by determinant symmetry to obtain the formula. For a full
flag, all iterated decompositions identify with the same ordered
tensor product of \(C(A_i,A_j)\), for its successive quotients
\(A_i\). Determinant additivity for refined filtrations proves
coherence. Each operation commutes with base change.
\end{proof}

A square-root orientation consists of lines \(o_A\) and
isomorphisms \(o_A^{\otimes2}\simeq D(A)\). Extension
compatibility requires maps
\[
 o_A\otimes o_B\otimes L(A,B)\longrightarrow o_F
\]
whose squares are the determinant comparisons. Choosing coordinates
on a smooth critical chart trivializes its ambient canonical line.
Hall convolution requires transition maps compatible with the cross
line and the zero object. On three-step flags, the two orientation
maps can differ by an element of \(\mu_2\), even when their
squares agree. That sign must be trivialized coherently.

\begin{example}[A nontrivial cross line]
Let \(X\) be a smooth projective Calabi--Yau threefold, with a
chosen volume form, and let \(x\in X\). On
\(X\times T\), where \(T=\mathbb P^1\), take
\[
 A=\mathcal O_X\boxtimes\mathcal O_T(1),\qquad
 B=\mathcal O_x\boxtimes\mathcal O_T.
\]
Then \(C(A,B)=\mathcal O_T(-1)[0]\), and relative Serre
duality gives \(C(B,A)=\mathcal O_T(1)[-3]\).
Thus \(L(A,B)=L(B,A)=\mathcal O_T(-1)\).
The self-Hom complexes of \(A\) and \(B\) have constant
determinant lines, since the factors \(\mathcal O_T(1)\)
cancel in the first self-Hom. Nevertheless
\[
 D(A\oplus B)=D(A)\otimes D(B)\otimes\mathcal O_T(-2).
\]
The orientation cross factor is \(\mathcal O_T(-1)\).
A sign change cannot replace this line of nonzero degree.
\end{example}

\begin{corollary}[Oriented critical coefficients]
Suppose the chosen critical charts have square-root orientation
maps on overlaps and extension flags, with unit and flag coherences.
Suppose their vanishing-cycle comparison maps preserve these data.
Then the local vanishing-cycle complexes descend to the specified
oriented coefficient object. Their Hall convolution is associative
whenever its pullback and pushforward maps are defined and satisfy
base change on the flag diagrams.
\end{corollary}
\begin{proof}
The overlap cocycle identifies the local coefficient complexes by
descent. On a three-step flag, orientation coherence identifies
the two Thom--Sebastiani coefficient maps. Proper or compact-support
base change, on the stipulated domain, identifies the two composites
of convolution. The zero flag gives the unit equation.
\end{proof}

For the quiver charts, the remaining construction is an equivariant
system of square roots and coefficient maps on the actual
representation and extension stacks, compatible with changes of
height. For compact \(S\times E\), the reduced complexes and
the common elliptic torsor must also be retained. A coordinate
volume on an affine chart specifies neither descent system by itself.

\chapter{The colored elliptic Ran space}
The elliptic degree is a global charge carried by the whole configuration.  Define
\[
 \Ran^{\mathrm{wr}}(E)=\coprod_{m\ge0}\Ran(E)_m
\]
with monoidal operation
\[
 (I,m)\sqcup(J,n)=(I\sqcup J,m+n)
\]
on disjoint supports.  The component $m=0$ is point-local.  The components $m>0$ are wrapped.

Let
\[
 \nIg=\bigoplus_{m\ge0}\nIg[m]
\]
be the elliptic-degree grading.  The root bracket obeys
\[
 [\nIg[m],\nIg[n]]\subset\nIg[m+n].
\]

\begin{construction}[Hybrid factorization algebra]
Form the chiral Chevalley factorization coalgebra
\[
 \mathcal F^{\mathrm{hyb}}_{\Delta,E}
 =C_*^{\mathrm{ch}}\!\left(\Cur_E(\nIg)\right).
\]
Retain the elliptic-degree grading on every Chevalley word and place its degree-$m$ summand on $\Ran(E)_m$.  Disjoint union adds the degree labels.
\end{construction}

\begin{theorem}[Local--wrapped factorization]
$\mathcal F^{\mathrm{hyb}}_{\Delta,E}$ is a graded factorization coalgebra on $\Ran^{\mathrm{wr}}(E)$.  Its degree-zero color is the current theory of $\nIg[0]$.  Every degree-$m$ color is a module for the degree-zero color.  The factorization product of degrees $m$ and $n$ lands in degree $m+n$.  Its chiral primitives are $\Cur_E(\nIg)$, and constant modes recover $\nIg$.
\end{theorem}
\begin{proof}
The chiral Chevalley construction is a factorization coalgebra for every chiral Lie algebra.  The bracket preserves elliptic-degree addition, so its differential and coproduct preserve the graded Ran placement.  The primitive theorem for the Chevalley coalgebra recovers the current Lie algebra.
\end{proof}

This construction retains the nonnegative root-degree label on a
point-supported current theory. It does not replace a sheaf projecting
onto the whole elliptic curve by a point-supported sheaf. The compact
wrapped comparison must map its common-torsor extension correspondences
to these graded current operations, preserving support and charge.
The bracket on the present carrier is the Borcherds bracket.

\chapter{Physical charge and Gram degree}
Let $\Gammphys$ be an additive physical charge lattice and let
\[
 \PiGram(c)=\left(\frac{Q^2}{2},Q\cdot P,\frac{P^2}{2}\right)
\]
be the Gram map.  Its polarization is
\[
 B(c,d)=\PiGram(c+d)-\PiGram(c)-\PiGram(d).
\]
Define the central extension
\[
 \widetilde\Gamma_\Pi=\Gammphys\times\Gammaind
\]
with multiplication
\[
 (c,\lambda)(d,\mu)
 =(c+d,\lambda+\mu+B(c,d)).
\]

\begin{theorem}[Graph homomorphism]
The map
\[
 j_\Pi:\Gammphys\longrightarrow\widetilde\Gamma_\Pi,
 \qquad
 c\longmapsto(c,\PiGram(c))
\]
is a group homomorphism.  For every $\Gammphys$-graded algebra $A$, the Rees lift
\[
 a_c\longmapsto a_cu^{\PiGram(c)}
\]
turns multiplication into an additive $\widetilde\Gamma_\Pi$-grading with cocycle $u^{-B(c,d)}$.
\end{theorem}
\begin{proof}
The defining multiplication gives
\[
 j_\Pi(c)j_\Pi(d)
 =(c+d,\PiGram(c)+\PiGram(d)+B(c,d))
 =j_\Pi(c+d).
\]
The algebra identity follows immediately.
\end{proof}

The universal Igusa category is graded directly by the Lorentzian root lattice.  A geometric $K3\times E$ theory enters through a charge functor to $\widetilde\Gamma_\Pi$.

\chapter{The universal protected field theory}
For each $H$ take cotangent BF theory with gauge Lie superalgebra $\nIg^{(H)}$ on $\R^2_{\mathrm{top}}\times E$.  Its fields are
\[
 \Omega^\bullet(\R^2)\widehattensor
 \Omega^{0,\bullet}(E)
 \widehattensor
 \left(\nIg^{(H)}[1]\oplus(\nIg^{(H)})^\vee[2]\right).
\]
The action is
\[
 S_H=\int_{\R^2\times E}
 \angles{B,(\dd+\dbar)A+\frac12[A,A]}.
\]

\begin{theorem}[Classical local theory and exact quantum harmonic sector]
For every finite height $H$ the shifted-cotangent action defines a classical mixed holomorphic--topological factorization algebra on $\R^2_{\mathrm{top}}\times E$.  Harmonic reduction on a compact product cell gives the finite-dimensional cotangent BV theory of $\nIg^{(H)}$, and this reduced theory satisfies the quantum master equation.  Complementary Lagrangian boundary polarizations give the boundary enveloping algebra $U(\nIg^{(H)})$, the two-sided bar complex $C_*(\nIg^{(H)})$, and the current envelope on $E$.  The proper-height inverse system is coefficientwise finite and defines a complete classical local theory together with an exact quantum harmonic sector.
\end{theorem}
\begin{proof}
The invariant cotangent pairing and the Jacobi identity give the classical master equation and factorization descent.  Harmonic reduction leaves a finite positive nilpotent Lie superalgebra.  Every even adjoint endomorphism raises proper height and is strictly upper triangular; every odd adjoint endomorphism has vanishing supertrace.  The reduced BV divergence therefore vanishes.  Canonical quantization in the selected polarization gives the enveloping algebra, and the augmentation boundary states give the two-sided bar, identified with Chevalley chains by PBW.  Each map preserves height, so the inverse system is defined coefficientwise.
\end{proof}


\chapter{Absolute convergence of the virtual root operator}
For each positive root $\alpha$, let $V_\alpha$ be the minimal parity realization of the signed multiplicity: it is purely even of dimension $\sdim\gIg_\alpha$ when this integer is nonnegative and purely odd of dimension $-\sdim\gIg_\alpha$ when it is negative.  Set
\[
 \Vcal_\Delta=\widehat\bigoplus_{\alpha>0}V_\alpha.
\]
This virtual root superspace is the canonical linear carrier of the denominator character; the Lie bracket remains carried by $\gIg$.

Let $h(\alpha)$ be the proper height.  Jacobi-form coefficient estimates give constants $C_0,C_1>0$ such that
\begin{equation}\label{eq:vol3-root-growth}
 \dim V_\alpha=\abs{\sdim\gIg_\alpha}\le C_0\exp(C_1h(\alpha))
\end{equation}
for every positive root.  The number of Gram triples of height $H$ is polynomial in $H$.

\begin{theorem}[Nonempty trace-class chamber]
Let $Z_0$ lie in the interior of the positive Weyl chamber.  For sufficiently large $t>0$,
\[
 \sum_{\alpha>0}\dim V_\alpha\,
 \abs{e^{-\alpha(tZ_0)}}<\infty.
\]
Hence the absolute-convergence chamber contains a conical tail in every interior direction.
\end{theorem}
\begin{proof}
Interior positivity gives a constant $c(Z_0)>0$ with
$\Re\alpha(Z_0)\ge c(Z_0)h(\alpha)$.
Combine this with \eqref{eq:vol3-root-growth}.  The contribution of roots of height $H$ is bounded by a polynomial in $H$ times
\[
 \exp\bigl((C_1-tc(Z_0))H\bigr).
\]
The series converges when $tc(Z_0)>C_1$.
\end{proof}

\begin{corollary}[Trace-class root Hamiltonian]
On the conical tail, the diagonal operator
\[
 K_Z\big|_{V_\alpha}=e^{-\alpha(Z)}\id
\]
is trace class on both parity summands.  The Fredholm product is holomorphic and equals the absolutely convergent Borcherds product.
\end{corollary}
\begin{proof}
The trace norm is the summable series of the preceding theorem, separated by parity.  Standard trace-class determinant theory gives holomorphy, and diagonalization gives the product over root eigenvalues.
\end{proof}


\chapter{The virtual root Fredholm operator}
Let
\[
 \mathcal H^{\mathrm{vir}}_\Delta
 =\widehat\bigoplus_{\alpha>0}V_\alpha
\]
with an orthonormal basis in each minimal parity realization.
This is the virtual root Hilbert space used in the preceding
absolute-convergence theorem.  In the absolute-convergence chamber define
\[
 K_Z\big|_{V_\alpha}=e^{-\alpha(Z)}\id,
 \qquad
 D_Z=1-K_Z.
\]

\begin{theorem}[Oriented Fredholm square root]
The operator $K_Z$ is trace class on the even and odd summands of the virtual root Hilbert space.  Its Fredholm Berezinian is
\[
 \sdet_F(D_Z)
 =\prod_{\alpha>0}(1-e^{-\alpha(Z)})^{\sdim\gIg_\alpha}.
\]
Choose an order of the positive roots, first by height and then lexicographically.  This choice orients every finite-height root space.  The skew doubling
\[
 \mathbb A_Z=
 \begin{pmatrix}0&D_Z\\-D_Z^{\mathsf t}&0\end{pmatrix}
\]
has an oriented Fredholm super-Pfaffian satisfying
\[
 \operatorname{sPf}_F(\mathbb A_Z)=\sdet_F(D_Z).
\]
Consequently
\[
 e^{-\rho}\operatorname{sPf}_F(\mathbb A_Z)=\Dmon(2Z).
\]
\end{theorem}
\begin{proof}
Absolute convergence of the Borcherds product is the trace-norm summability condition.  The determinant formula follows from the spectrum.  At finite height the Pfaffian of the oriented skew doubling is the determinant.  Trace-norm convergence passes the identity to the limit.  The denominator theorem supplies the Weyl shift.
\end{proof}

\section{Automorphic orientation}
The product defines a section of the automorphic line $\mathcal L^5\tensor\nu_5$.  On the absolute-convergence chamber, the oriented Fredholm Pfaffian is a trivialization of this line.  Analytic continuation uses the automorphy factor of $\Dmon$ and therefore glues the chamberwise Fredholm models with transition character $\nu_5$.

\begin{corollary}[Orientation character]
The automorphic orientation line of the root Dirac family is the sign local system $\nu_5$.  On every real simple reflection it has value $-1$.
\end{corollary}
\begin{proof}
The Fredholm product equals $D_5$ on the convergence chamber, and $D_5$ transforms with automorphy factor $\nu_5$.  This transformation law is the gluing rule for the automorphic orientation line.  A real simple root occurs with multiplicity one, so reflection reverses its oriented normal line.
\end{proof}

\chapter{The universal Igusa protected effective theory}
\begin{definition}[Universal Igusa completion]
The protected completion selected by the K3 half-index is the quintuple
\[
 \mathfrak K_{\Delta}=
 \left(
 \Gamma_\Delta,
 \mathcal B_\Delta,
 \mathcal H_{\mathrm{1p}},
 \mathcal F^{\mathrm{hyb}}_{\Delta,E},
 \mathbb A_Z
 \right).
\]
Its components are the pro-system of finite three-Calabi--Yau completions, the ordinary quiver Hall category, the Borcherds one-particle Lie superalgebra, the local--wrapped current factorization algebra, and the virtual root Dirac family.
\end{definition}

\begin{theorem}[Igusa protected completion]
The quintuple $\mathfrak K_\Delta$ satisfies:
\begin{enumerate}
\item its one-particle Lie superalgebra is $\gIg$;
\item its ordinary Hall presentation and Hong's quantum supergroup retain
the stated Cartan grading; comparison of their relation ideals and
parity-refined PBW spaces requires an additional map;
\item its chiral primitive algebra on $E$ is the local--wrapped current algebra of $\gIg$;
\item its raw multiparticle supercharacter is $P^{-1}$, where $e^{-\rho}P=\DenIg$; equivalently $e^{\rho}\sch\Fock(\nIg)=\DenIg^{-1}$;
\item the half-coordinate pullback of the doubled Weyl-normalized character is $-\ChiTen^{-1}$ in the Oberdieck--Pixton normalization;
\item its oriented first-order determinant is $\DenIg$ and its orientation monodromy is $\nu_5$.
\end{enumerate}
\end{theorem}
\begin{proof}
The first three statements are the Borcherds root, quiver Hall, Hong quantum-supergroup, Ginzburg, and colored-Ran theorems.  PBW gives the fourth.  The identity $\ChiTen=\Dmon^2$ and the established reduced scalar theorem give the fifth.  The Fredholm theorem gives the sixth.
\end{proof}

The tuple retains algebraic constructions associated with the Igusa
presentation and their scalar characters. A geometric type-IIA or
topological-string realization must construct compatible functors
from its oriented extension category to these targets. Primitive
states and extension amplitudes are inputs to those functors.
Their action on morphisms, convolution, orientations, and chiral
operations must also be specified and proved compatible.

\chapter{The primitive comparison and its extensions}
Let \(\mathfrak p\) be a positive-root-graded Lie superalgebra
with finite-dimensional homogeneous pieces and a proper positive
height. A primitive realization of the Igusa presentation is a
parity-preserving homogeneous map from its labelled simple-root
spaces to \(\mathfrak p\), satisfying the defining relations.

\begin{theorem}[Universal presentation map]
\label{igusa:primitive-presentation-map}
Let \(\mathfrak p\) be a positive-root-graded Lie superalgebra
over \(\mathbb C\), with finite-dimensional homogeneous pieces.
Suppose a homogeneous parity-preserving map from every simple-root
generating space of \(\nIg\) to \(\mathfrak p\) satisfies
the defining Borcherds--Serre relations. It extends uniquely to
a graded Lie-superalgebra map
\[
 f:\nIg\longrightarrow\mathfrak p.
\]
It induces continuous maps on height completions, enveloping
algebras, Chevalley chains, and level-zero currents. The map \(f\)
is an isomorphism if and only if it is surjective and the root-graded
PBW series agree in each parity.
\end{theorem}
\begin{proof}
Extend the chosen maps to the free Lie superalgebra on the
simple-root spaces. The relation hypothesis makes the defining
ideal lie in its kernel, giving \(f\). The generators prove
uniqueness. On enveloping algebras the map sends
\(x_1\cdots x_j\) to \(f(x_1)\cdots f(x_j)\).
On Chevalley chains it is the graded symmetric map on the shifted
generators. It commutes with the Chevalley differential because
\(f\) preserves brackets. The current map is
\(\partial^j x\mapsto\partial^j f(x)\), preserving
\([x_\lambda y]=[x,y]\). Root degree preservation gives
compatible maps at each height and hence on the stated limits.

If \(f\) is surjective, \(U(f)\) is surjective. Proper
positive height makes each root-graded PBW component
finite-dimensional. Equality of its even and odd dimensions makes
\(U(f)\) bijective in every component. PBW embeds each Lie
superalgebra into its enveloping algebra, so \(f\) is injective.
The converse follows by applying the constructions to its inverse.
\end{proof}


A geometric comparison also requires functors on the actual Hall
and critical categories, a compatible coproduct and Hopf pairing,
and maps of determinant lines and operator domains. These categories
and lines are not the domain of the primitive Lie map. The
additional maps must be constructed and checked against that map.
They are not consequences of its universal property.

\chapter{The eight diagonal-divisor forms}
Cl\'ery and Gritsenko classify eight genus-two modular forms whose divisor is the diagonal with order one \cite{GC}.  Gritsenko--Nikulin and the CHL denominator constructions give the sourced reflective presentations used for the corresponding rows \cite{GN97,GN98a,GN98b,GovindarajanKrishna2009,Govindarajan2011,GovindarajanSamanta2019}.  Every row has a minimal virtual product superspace obtained from its Fourier exponents, together with a Fock determinant and a diagonal Fredholm operator.  A sourced denominator presentation refines this product model by its Borcherds--Cartan bracket, Weyl group, and quantum supergroup.  The half-integral rows are formed on their multiplier covers.
\[
\begin{array}{c|c|c|c}
\text{form}&\text{group}&\text{weight}&\text{character order}\\ \hline
\Delta_5&\operatorname{Sp}_4(\Z)&5&2\\
\Delta_2&\Gamma_2&2&4\\
\Delta_1&\Gamma_3&1&6\\
\Delta_{1/2}&\Gamma_4&1/2&8\\
\nabla_3&\Gamma^{(2)}_0(2)&3&2\\
\nabla_2&\Gamma^{(2)}_0(3)&2&2\\
\nabla_{3/2}&\Gamma^{(2)}_0(4)&3/2&4\\
Q_1&\Gamma_2(2)&1&4
\end{array}
\]

\begin{theorem}[Eight-row algebraic realization]
Each diagonal-divisor row has a proper-height virtual root superspace, a super-Fock determinant, a colored abelian current theory, and a trace-class Fredholm denominator operator whose character is the corresponding automorphic product.  For every row equipped with its reflective denominator presentation, the finite-window Borcherds quiver, ordinary Hall category, three-Calabi--Yau completion, and parity-refined quantum supergroup refine this product model and recover the same denominator character.
\end{theorem}
\begin{proof}
Local finiteness of the product exponents gives the graded superspace and its proper-height completion.  The Fock and Fredholm constructions depend only on these signed multiplicities and reproduce the product coefficientwise.  A reflective denominator presentation supplies a symmetrizable Borcherds superdatum; the quiver, Hall, Ginzburg, current, and Hong constructions then apply to that datum and preserve its PBW denominator.
\end{proof}

\chapter{Algebraic and physical closure}
\begin{theorem}[Universal Igusa realization theorem]
The constructions have the following inputs and outputs. An arrow
labelled by a construction records its dependence on that input:
\[
\begin{gathered}
 \phiweak\xmapsto{\mathrm{Borch}}\Dmon,
 \qquad
 \phiweak\xmapsto{\mathrm{root\ datum}}\gIg,
 \qquad
 A_\Delta\xmapsto{\mathrm{quiver}}Q_{\Delta,H}
 \xmapsto{\mathrm{Hall}}\mathcal B_\Delta,\\
 kQ_{\Delta,H}\xmapsto{\Pi_3}\Gamma_{\Delta,H},
 \qquad
 \gIg\xmapsto{\Cur_E}\mathcal F^{\mathrm{hyb}}_{\Delta,E},
 \qquad
 \gIg\xmapsto{\mathrm{Hong}}\widehat U_q(\gIg),\\
 \gIg\xmapsto{\mathrm{root\ spectrum}}\mathbb A_Z,
 \qquad
 e^{-\rho}\operatorname{sPf}_F(\mathbb A_Z)=\DenIg(Z)=\Dmon(2Z),\\
 \mathcal Z^{\red}_{K3\times E}=-(\jmath^*\DenIg)^{-2}=-\ChiTen^{-1}.
\end{gathered}
\]
The system contains parity-refined one-particle spaces, an ordinary Hall extension product, the Borcherds root bracket, a local--wrapped collision law, a quantum supergroup, an automorphic orientation line, a first-order Fredholm determinant, and the reduced scalar square.
\end{theorem}
\begin{proof}
The Borcherds lift, Gritsenko--Nikulin presentation, and primitive
reduced curve-counting theorem give the scalar identities. The
enveloping and Chevalley maps are induced by the presentation.
The quiver Hall algebra and finite Calabi--Yau completion have their
respective quiver and path-algebra inputs. Their common grading
specifies the degrees a comparison must preserve. A map between
those constructions requires the additional functors, coefficient
maps, and pairings described above.
\end{proof}


\part{Technical appendices to the preceding book}
\chapter{Exact theta-series calculation}\label{app:theta-computation}
Set $t=q^{1/8}$ and $x=r^{1/2}$. Then
\[
\begin{aligned}
\vartheta_2(\tau,z)&=\sum_{n\in\Z}t^{4n^2+4n+1}x^{2n+1},\\
\vartheta_3(\tau,z)&=\sum_{n\in\Z}t^{4n^2}x^{2n},\\
\vartheta_4(\tau,z)&=\sum_{n\in\Z}(-1)^nt^{4n^2}x^{2n}.
\end{aligned}
\]
Squaring numerator and denominator, dividing in $\Q[x^{\pm1}][[t]]$, and summing the three quotients gives the coefficient table in Chapter~1. The calculation uses integer convolution and formal-series inversion; fractional powers cancel before conversion from $(t,x)$ to $(q,r)$.

\chapter{The root reflection matrices}
In the basis $(f_2,f_3,f_{-2})$, the form is
\[
 G=\begin{pmatrix}0&0&-1\\0&2&0\\-1&0&0\end{pmatrix}.
\]
The three reflection matrices are
\[
 S_1=\begin{pmatrix}1&4&4\\0&-1&-2\\0&0&1\end{pmatrix},\quad
 S_2=\begin{pmatrix}1&0&0\\-2&-1&0\\4&4&1\end{pmatrix},\quad
 S_3=\begin{pmatrix}1&0&0\\0&-1&0\\0&0&1\end{pmatrix}.
\]
Direct multiplication gives
\[
 S_i^2=I,
 \qquad S_i^{\mathsf T}GS_i=G.
\]

\chapter{Automorphic, homological, and Hall comparisons}
The half-index determinant, theta normalization, Borcherds denominator, reduced scalar theorem, charge Gram map, and geometric Hall comparison datum belong to the following constructions:
\begin{enumerate}
\item the automorphic square root is completed entirely within the Borcherds product theorem;
\item the BKM bracket comes from the Gritsenko--Nikulin algebra;
\item the bar--Chevalley comparison gives the denominator a homological source;
\item the current and BF constructions give an exact first-order algebraic model;
\item orientation data refine determinant lines, while an operator Pfaffian requires an actual Fredholm family;
\item Weyl reflections are attached to real roots, and imaginary roots enter through multiplicities;
\item the compact Hall comparison uses additive Mukai charge, the polarized central extension, and hybrid local--wrapped sectors.
\end{enumerate}

\cleardoublepage
\phantomsection
\addcontentsline{toc}{part}{Calabi--Yau Quantum Groups}
\chapter*{Book entrance: Calabi--Yau Quantum Groups}
\addcontentsline{toc}{chapter}{Book entrance: Calabi--Yau Quantum Groups}
\begin{summarybox}
An extension of geometric objects defines a correspondence. Its action on cohomology depends on the coefficient theory and the allowed pushforward. Tensor products and exchange require a coproduct and compatible intertwiners. A trace on a specified representation then produces an amplitude. The reconstruction problem is to retain these operations through the comparisons between Calabi--Yau categories, boundary fields, Hall algebras, and quantum groups.
\end{summarybox}
\part{Categorical and field-theoretic foundations}

\chapter{Calabi--Yau categories}
\section{Smoothness, properness, and the Calabi--Yau class}
Let $\Ccal$ be a small stable $\kfield$-linear category. Smoothness means that the diagonal bimodule is perfect. Properness means that the morphism complexes are perfect over $\kfield$. A $d$-Calabi--Yau structure is a nondegenerate negative-cyclic class
\[
 \theta_{\CY}\in HC^-_d(\Ccal)
\]
whose underlying Hochschild trace identifies the diagonal bimodule with its dual shifted by $d$.

For $E\in\Ccal$, the tangent complex of the derived moduli stack of objects is
\[
 T_E\M_{\Ccal}\simeq\RHom_{\Ccal}(E,E)[1].
\]
The Calabi--Yau trace pairs two tangent vectors by composition and trace.

\begin{theorem}[The two universal categorical outputs]
Let $\Ccal$ be a smooth proper $d$-Calabi--Yau category.
\begin{enumerate}
\item The Hochschild cochains $CC^*(\Ccal)$ carry a framed $\Etwo$-algebra structure.
\item Whenever the moduli stack of objects is representable in the PTVV setting, $\M_{\Ccal}$ carries a $(2-d)$-shifted symplectic form.
\end{enumerate}
\end{theorem}
\begin{proof}
The cyclic Deligne theorem upgrades the brace $\Etwo$ structure on Hochschild cochains by the circle action supplied by the cyclic trace \cite{BravRozenblyum}. The shifted form at $E$ is
\[
 \RHom(E,E)[1]^{\tensor2}
 \xrightarrow{\ \circ\ }
 \RHom(E,E)[2]
 \xrightarrow{\ \Tr\ }
 \kfield[2-d].
\]
Nondegeneracy is the Calabi--Yau duality isomorphism, and closedness is the negative-cyclic lift \cite{PTVV,BravDyckerhoff}.
\end{proof}

The integer $d$ controls the shift of the moduli form. The Hochschild operation remains framed $\Etwo$. These structures meet through deformation theory while retaining their separate operadic degrees.

\chapter{Local BV theories and factorization observables}
\section{Cyclic elliptic local Lie algebras}
Let $X$ be a complex manifold. A classical local field theory is specified by a sheaf $\Lcal$ of elliptic $L_\infty$ algebras, a local invariant pairing, and an action functional $S$ satisfying the classical master equation. For an open set $U\subset X$, the classical observables are
\[
 \Obs^{\mathrm{cl}}(U)=C^*_{\CE,\mathrm{cont}}(\Lcal(U)).
\]
Here $\Lcal(U)$ consists of unrestricted smooth fields with its
nuclear smooth-section topology. Continuous cochains use the completed
symmetric powers of its continuous dual, with the completed tensor
products of the elliptic-observable construction \cite{CG1,CG2}.
The dual consists of compactly supported distributions.
For $U\subset V$, restriction of fields $\Lcal(V)\to\Lcal(U)$
induces the covariant map on observables.
For pairwise disjoint $U_i\subset V$, restrict to
$\bigoplus_i\Lcal(U_i)$ and multiply the pulled-back cochains.
These are the prefactorization maps.

\begin{theorem}[Factorization observables]
An elliptic classical BV theory produces a factorization algebra $\Obs^{\mathrm{cl}}$. A renormalized quantization satisfying the quantum master equation produces a factorization algebra $\Obs^{\mathrm{q}}$ over $\kfield[[\hbar]]$.
\end{theorem}
\begin{proof}
Restriction is a morphism of local $L_\infty$ algebras, so its
contravariant continuous Chevalley map commutes with the differential.
For disjoint opens, the completed cochains of the finite direct sum
are the completed tensor product of the cochains. This constructs
the displayed maps and their composition identities.
The elliptic descent theorem in this completed coefficient category
supplies Weiss descent \cite{CG1,CG2}. In the quantum theory, renormalized time-ordered products define the structure maps, and the quantum master equation makes their differential square-zero and compatible with nested ultraviolet contractions \cite{CG1,CG2}.
\end{proof}

\section{Boundary and defect restriction}
Let $Y=C\times\R_{\ge0}$, with $C$ a complex curve, and let a boundary condition be given at $C\times\{0\}$. Boundary observables are holomorphic along $C$ and ordered along the inward normal line. A compact boundary object $b$ in the boundary factorization category gives
\[
 A_b=\RHom(b,b).
\]

\begin{theorem}[Ordered boundary algebra]
When boundary observables are holomorphic on $C$ and locally constant along the oriented normal direction, $A_b$ is an ordered chiral algebra
\[
 A_b\in\Alg_{\Eone}\bigl(\Fact_D(C),\tensor^!\bigr).
\]
If $b$ is a local compact generator and the Morita equivalences obey factorization descent, the boundary category is equivalent to $A_b$-modules in factorization coefficients.
\end{theorem}
\begin{proof}
Composition of boundary endomorphisms gives the internal $\Eone$ multiplication. Holomorphic collision gives the factorization coefficient on $C$. The componentwise locality identity makes multiplication compatible with every disjoint-support factorization map. Local compact-generator Morita theory descends over the Ran diagrams. These are the recognition and Morita theorems of Volume~I.
\end{proof}


\chapter{Three curve-level algebraic outputs}
Let $C$ be a smooth curve and let $\Dcal_C$ be the fixed Ran category.  The same geometric source can produce three algebraic outputs, and their distinction is functorial.

\begin{definition}[Curve-level faces]
A transverse face is
\[
 A^\perp\in\Alg_{\Eone}(\Fact_D(C),\tensor^!).
\]
An associative chiral face is
\[
 A^{\ch}_{\mathsf{Ass}}\in\Alg_{\Eone}(\Dcal_C,\starCh).
\]
A chiral Lie face is
\[
 \g^{\ch}\in\Alg_{\mathsf{Lie}}(\Dcal_C,\starCh).
\]
Their bars are $\BarPerp$, $\BarAssCh$, and $\BarLieCh$.
\end{definition}

\begin{theorem}[Source separation]
A boundary normal direction determines the transverse face.  A chiral-convolution multiplication determines the associative chiral face.  A Lie bracket for chiral convolution determines the Lie face.  The commutator functor sends the associative chiral face to the Lie face.  A chiral enveloping algebra $U^{\ch}(\g)$ satisfies
\[
 \BarAssCh(U^{\ch}(\g))\simeq\BarLieCh(\g)
\]
on the nilcomplete PBW locus.  A mixed source joins the transverse and associative chiral faces through an aligned distributive law.
\end{theorem}
\begin{proof}
Each statement is the corresponding operadic restriction or adjunction in the named monoidal category.  The PBW equivalence is the bar comparison for chiral enveloping algebras.  The mixed statement is the distributive-law theorem of Volume~I.
\end{proof}

The source geometry selects the available directions.  The $\En$ level records the operations and coherences constructed on them; dimension alone does not specify those structures.

\chapter{Holomorphic Chern--Simons theory}
Let $X$ be a complex threefold with holomorphic volume form $\Omega_X$, and let $\mathfrak a$ be a quadratic Lie algebra. The BV fields are
\[
 \mathcal E_X=\Omega^{0,*}(X,\mathfrak a)[1],
\]
with degree $-1$ pairing
\[
 \omega(A,B)=\int_X\Omega_X\wedge\angles{A,B}.
\]
The action is
\[
 S_{\hCS}(A)=\int_X\Omega_X\wedge
 \left(\frac12\angles{A,\dbar A}+\frac16\angles{A,[A,A]}\right).
\]

\begin{theorem}[Classical hCS factorization algebra]
The action $S_{\hCS}$ satisfies the classical master equation. Its classical observables form a holomorphic factorization algebra on $X$.
\end{theorem}
\begin{proof}
The Hamiltonian vector field is the Maurer--Cartan differential $\dbar A+\frac12[A,A]$. Its square is the Dolbeault square, the derivation identity, and the Jacobi identity. The elliptic Dolbeault complex then gives factorization observables by the preceding theorem.
\end{proof}

A perturbative quantization is controlled by the local BV deformation complex. In complex dimension three, its one-loop gauge anomaly occupies the degree-four invariant-polynomial summand; mixed and gravitational components enter when the background fields carry the corresponding couplings. A quantization is a trivialization of the complete local anomaly class, compatible with restriction and factorization \cite{CostelloLiAnomaly}.

\chapter{Pushforward to a curve}
Let $p:X\to C$ be a holomorphic map to a complex curve. For a factorization algebra $\F$ on $X$, define
\[
 (p_*\F)(U)=\F(p^{-1}U).
\]

\begin{theorem}[Factorization pushforward]
The assignment $p_*\F$ is a factorization algebra on $C$. Suppose in addition that $p$ is proper and locally trivial in a topology compatible with the factorization theory, that $\F$ is locally constant in the fiber directions, and that product-chart excision and base change hold. Then on every trivializing open $U\subset C$ with fiber $F_p$ there is a natural equivalence
\[
 (p_*\F)(U)\simeq
 \int_{F_p}\F|_{F_p\times U},
\]
compatible on overlaps. If the resulting object is holomorphic and regular holonomic along $C$, Riemann--Hilbert algebraization gives a factorization $D$-module on $C$.
\end{theorem}
\begin{proof}
Disjoint opens in $C$ have disjoint inverse images, so the factorization products push forward. Weiss descent is inherited from inverse images of Weiss covers. On a product chart, local constancy in the fiber and factorization excision identify the value with fiberwise factorization homology. The base-change hypothesis glues these local equivalences. Regular-holonomic Riemann--Hilbert transports the holomorphic factorization object to the algebraic $D$-module setting.
\end{proof}

Thus a curve shadow is attached to an actual map and an actual field theory. Different fibrations and boundary conditions produce different ordered chiral algebras from the same categorical source.

\chapter{The derived centre and the physical bulk}
For an ordered boundary algebra $A$, define
\[
 Z_{\ord}(A)=\RHom_{A\tensor A^{\op}}(A,A).
\]
It is a canonical $\Etwo$ algebra. A bulk theory acting on the boundary gives
\[
 \beta:\Obs_{\mathrm{bulk}}|_{\partial}\longrightarrow Z_{\ord}(A).
\]

\begin{definition}[Open--closed completeness]
A boundary condition is open--closed complete on a protected sector when $\beta$ is an equivalence on that sector.
\end{definition}

The definition turns the phrase ``bulk equals centre'' into a comparison theorem. The centre exists algebraically for every ordered boundary algebra; the physical bulk is identified with it by the open--closed map.

\part{Hall and BPS algebras}

\chapter{Extension correspondences}
Let $\M=\coprod_{\gamma\in\Gamma}\M_\gamma$ be a moduli stack in an exact or derived category. Let $\M^{\mathrm{ext}}_{\alpha,\beta}$ classify extensions
\[
 0\longrightarrow E_\beta\longrightarrow E\longrightarrow E_\alpha\longrightarrow0.
\]
The correspondence is
\[
 \M_\alpha\times\M_\beta
 \xleftarrow{\ p\ }
 \M^{\mathrm{ext}}_{\alpha,\beta}
 \xrightarrow{\ q\ }
 \M_{\alpha+\beta}.
\]

\begin{theorem}[Hall associativity]
Assume a coefficient theory in which $p^*$ and $q_*$ are defined and satisfy base change for the two-step filtration stack. Then
\[
 a*b=q_*p^*(a\boxtimes b)
\]
defines an associative product.
\end{theorem}
\begin{proof}
The two parenthesizations are the two projections from the stack of filtrations $0\subset E_1\subset E_2\subset E_3$. Base change identifies both composites with pullback to this stack followed by pushforward to $E_3$.
\end{proof}

For critical CoHAs, the coefficient theory includes vanishing cycles and a coherent orientation. For symplectic CoHAs, it includes the selected equivariant Borel--Moore theory. Completion is chosen so that each output degree receives a convergent sum.

\chapter{Orientation data and the threefold Hall product}
For compact Calabi--Yau threefolds, natural orientation data on moduli of coherent sheaves and perfect complexes is constructed by Joyce and Upmeier \cite{JoyceUpmeier}.  A general cohomological Hall algebra for a $3$-Calabi--Yau category uses strong orientation data, which carries the coherent multiplicativity required by the attractor and extension correspondences \cite{KPS}.

\begin{theorem}[Threefold Hall algebra from strong orientation]
Let $\Ccal$ be a $3$-Calabi--Yau category satisfying the geometric hypotheses of Kinjo--Park--Safronov, and assume strong orientation data.  Then its cohomological Donaldson--Thomas complexes carry an associative cohomological Hall product.  The product is the parabolic-induction correspondence supplied by the attractor Lagrangian correspondence.
\end{theorem}
\begin{proof}
Kinjo--Park--Safronov prove the functoriality of the Donaldson--Thomas perverse sheaf under the attractor Lagrangian correspondence and use strong orientation to define the parabolic-induction morphism.  Their base-change and associativity theorem gives the CoHA \cite{KPS}.
\end{proof}

\begin{definition}[Reduced $K3\times E$ Hall datum]
A reduced $K3\times E$ Hall datum is a strong-oriented threefold Hall datum together with descent along the elliptic translation action, a reduced deformation theory, and a separated completion in physical charge.  A full quantum-group datum additionally includes a coproduct, Cartan extension, negative half, and nondegenerate continuous Hopf pairing.
\end{definition}

\chapter{The geometric \texorpdfstring{$E$}{E}-equivariant Hall source}
Let $S$ be a projective K3 surface, let $E$ be an elliptic curve, and put
\[
 X=S\times E.
\]
Write $\Coh_{\leq1}(X)$ for coherent sheaves whose support has dimension at most one.  This is an extension-closed abelian category.  Its numerical charge group is denoted $\Gamma_X$, and the degree of the projected one-cycle on $E$ is the additive function
\[
 m:\Gamma_X\longrightarrow\Z_{\geq0}.
\]
For $\gamma\in\Gamma_X$, let $\M_\gamma$ be the moduli stack of objects of charge $\gamma$.  Translation on the elliptic factor gives an action of $E$ on every $\M_\gamma$ and on every extension stack.

Oberdieck and Shen construct the $E$-equivariant Grothendieck ring of stacks and its Hall convolution.  In the present notation the completed algebra is
\[
 \Hcal_X^{E,\mathrm{mot}}
 =\widehat\bigoplus_{\gamma\in\Gamma_X}
 K_0^E(\operatorname{St}/\M_\gamma),
\]
where the completion is taken in effective curve class and Euler characteristic.

\begin{theorem}[Geometric reduced Hall algebra]
The extension correspondence in $\Coh_{\leq1}(X)$ defines an associative $E$-equivariant motivic Hall product on $\Hcal_X^{E,\mathrm{mot}}$.  The Oberdieck--Shen reduced integration morphism
\[
 I_X^E:\Hcal_X^{E,\mathrm{mot}}
 \longrightarrow
 \widehat{\Q[\Gamma_X]}[\epsilon]/(\epsilon^2)
\]
is multiplicative on the regular subalgebra on which their integration theorem is defined.  The coefficient of $\epsilon$ is the reduced contribution of the positive-dimensional $E$-orbits.
\end{theorem}
\begin{proof}
The stack of short exact sequences is $E$-equivariant because translation preserves exactness and numerical charge.  The stack of two-step filtrations identifies the two composites of three Hall products, so equivariant base change proves associativity.  The integration statement is the $E$-equivariant motivic Hall theorem of Oberdieck--Shen \cite{OberdieckShen}.  Their stabilizer stratification separates fixed contributions from moving elliptic orbits and packages the latter in the square-zero parameter $\epsilon$.
\end{proof}

\chapter{Local and wrapped geometry inside the Hall algebra}
Define the completed degree-zero and positive-degree pieces
\[
 \Hcal_X^{\mathrm{loc}}
 =\widehat\bigoplus_{m(\gamma)=0}
 K_0^E(\operatorname{St}/\M_\gamma),
 \qquad
 \Hcal_X^{\mathrm{wr}}
 =\widehat\bigoplus_{m(\gamma)>0}
 K_0^E(\operatorname{St}/\M_\gamma).
\]
A pure one-dimensional support of degree zero projects to a finite subset of $E$.  A support of positive degree projects onto $E$.  The first sector is point-local on the elliptic curve; the second is wrapped.

\begin{theorem}[Geometric local--wrapped decomposition]
Hall multiplication respects elliptic degree:
\[
 \Hcal_{m_1}*\Hcal_{m_2}\subseteq\Hcal_{m_1+m_2}.
\]
Consequently $\Hcal_X^{\mathrm{loc}}$ is a subalgebra, $\Hcal_X^{\mathrm{wr}}$ is a closed two-sided ideal, and
\[
 \Hcal_X^{E,\mathrm{mot}}
 \cong
 \Hcal_X^{\mathrm{loc}}\ltimes\Hcal_X^{\mathrm{wr}}
\]
as a complete vector space with multiplication determined by the Hall correspondences.  The reduced integration morphism preserves this decomposition.
\end{theorem}
\begin{proof}
The cycle of the middle term of a short exact sequence is the sum of the cycles of its subobject and quotient.  Projection to $E$ therefore adds degrees.  Degree zero is closed under addition, and positive degree absorbs addition by a nonnegative degree.  The same grading is present on the target group algebra of $I_X^E$.
\end{proof}

This theorem supplies the geometric version of the local--wrapped splitting constructed from the Borcherds root algebra in the companion volume.  The two constructions have the same additive elliptic degree and become comparable after a primitive root framing is chosen.

\chapter{Cohomological refinement and the Gram character}
A cohomological refinement uses the Donaldson--Thomas sheaf on the derived moduli stack.  The datum required for the convolution is the strong orientation coherence of Kinjo--Park--Safronov.  A strong orientation enhancement of the natural Joyce--Upmeier orientation therefore produces
\[
 \Hcal_X^{E,\mathrm{coh}}
 =\widehat\bigoplus_{\gamma\in\Gamma_X}
 H_*^{\mathrm{BM},E}(\M_\gamma,\Phi_\gamma),
\]
with its cohomological Hall product.

\begin{theorem}[Strong-oriented cohomological lift]
For every strong orientation enhancement on the $E$-equivariant derived moduli problem, the attractor correspondence constructs an associative cohomological Hall algebra $\Hcal_X^{E,\mathrm{coh}}$.  Elliptic degree again gives a subalgebra in degree zero and a closed ideal in positive degree.  Forgetting the cohomological coefficient system recovers the corresponding motivic Hall correspondences.
\end{theorem}
\begin{proof}
The parabolic-induction theorem of Kinjo--Park--Safronov gives the cohomological product from the strong-oriented attractor Lagrangian correspondence \cite{KPS}.  Numerical charge and elliptic degree are constant on connected components of every correspondence, so the proof of the preceding decomposition applies verbatim.
\end{proof}

For the dyonic charge lattice
\[
 \Gamma_{\mathrm{phys}}
 =\widetilde H(S,\Z)\oplus\widetilde H(S,\Z),
\]
write
\[
 \Pi(Q,P)=\left(\frac{Q^2}{2},Q\cdot P,\frac{P^2}{2}\right).
\]
To compare sheaf charges with this lattice, choose an additive map
\[
 c:\Gamma_X\longrightarrow\Gamma_{\mathrm{phys}}.
\]
Put $B_\Pi(a,b)=\Pi(a+b)-\Pi(a)-\Pi(b)$ and give
$\widetilde\Gamma_{\mathrm{phys}}=\Gamma_{\mathrm{phys}}\times\mathbb Z^3$
the group law
\[
 (a,u)(b,v)=(a+b,u+v+B_\Pi(a,b)).
\]
Then $j_\Pi(a)=(a,\Pi(a))$ is a homomorphism.
Assume $j_\Pi c$ preserves the selected charge completions.
On the regular subalgebra $\mathcal H_{\mathrm{reg}}$ of the reduced
integration theorem, define
\[
 I_{X,\Pi,c}^E=(j_\Pi c)_*I_X^E:
 \mathcal H_{\mathrm{reg}}\longrightarrow
 \widehat{\mathbb Q[\widetilde\Gamma_{\mathrm{phys}}]}
 [\epsilon]/(\epsilon^2).
\]
Write $I_{X,\Pi}^E$ for this map when $c$ has been fixed.

\begin{proposition}[Polarized reduced character]\label{cyiv:polarized-character}
The map $I_{X,\Pi,c}^E$ is multiplicative.
Its coefficients are the reduced invariants supplied by $I_X^E$,
regraded by $j_\Pi c$.
If $g\in1+F^1H$ is group-like in a complete filtered Hall bialgebra
and the logarithm converges, then $\log g$ is primitive.
\end{proposition}
\begin{proof}
Expanding $\Pi(a+b+d)$ gives
\[
 B_\Pi(a,b)+B_\Pi(a+b,d)=B_\Pi(b,d)+B_\Pi(a,b+d).
\]
This proves associativity of the group law.
The definition gives $j_\Pi(a+b)=j_\Pi(a)j_\Pi(b)$.
Composing its group-algebra map with reduced integration proves
multiplicativity on $\mathcal H_{\mathrm{reg}}$.
For the last assertion, the two factors in
$\Delta g=(g\widehat\otimes1)(1\widehat\otimes g)$ commute.
Taking the convergent logarithm gives
$\Delta\log g=\log g\widehat\otimes1+1\widehat\otimes\log g$.
\end{proof}

The primitive reduced Gromov--Witten theorem supplies the separate
character $-\chi_{10}^{-1}$ in classes primitive in the K3 factor
\cite{OP}.
Comparing this character with the Hall character requires the same
charge map, sector, and normalization.
A primitive root framing must also preserve the extension bracket
and parity of the chosen BPS states.

\chapter{BPS primitives}
A connected graded bialgebra has primitive Lie algebra
\[
 \Prim(H)=\set{x:\Delta x=x\tensor1+1\tensor x}.
\]
In cohomological Donaldson--Thomas theory, the BPS sheaf isolates the primitive part of the Hall algebra.

\begin{theorem}[Two-Calabi--Yau BPS algebra]
For the geometric two-Calabi--Yau categories covered by the BPS-algebra theorem, the BPS algebra is the positive enveloping algebra of a generalized Kac--Moody Lie algebra. This scope includes semistable sheaves on K3 and other symplectic surfaces under the stated good-moduli hypotheses.
\end{theorem}
\begin{proof}
The theorem is the Davison--Hennecart--Schlegel Mejia construction \cite{DHS}. The simple BPS sheaves give the simple-root spaces, extension correspondences give the bracket, and the integrality/PBW theorem identifies the Hall algebra with the positive enveloping algebra.
\end{proof}

This result supplies a global K3 generalized Kac--Moody algebra at the two-Calabi--Yau Hall level. Yangian structures arise from additional equivariant correspondences and coproducts in the surface and quiver models described below.

\chapter{The integral Hall--Borcherds amalgam}
Let $\mathfrak p_X^+$ be a complete charge-graded BPS Lie superalgebra supplied by a geometric Hall source, and let $\nDelta$ be the positive Igusa Borcherds Lie superalgebra.  The polarized charge extension places both gradings over one abelian Cartan algebra $\mathfrak h_\Pi$.  Form the two complete Borels
\[
 \mathfrak b_X=\mathfrak h_\Pi\ltimes\mathfrak p_X^+,
 \qquad
 \mathfrak b_\Delta=\mathfrak h_\Pi\ltimes\nDelta.
\]
The Cartan acts through the physical charge on the first Borel and through Gram degree on the second.

\begin{definition}[Integral protected Lie algebra]
The integral Hall--Borcherds algebra is the completed amalgamated coproduct
\[
 \mathfrak b_{X\boxtimes\Delta}
 =\mathfrak b_X
  \widehat\amalg_{\mathfrak h_\Pi}
  \mathfrak b_\Delta
\]
in complete charge-filtered dg Lie superalgebras.
\end{definition}

\begin{theorem}[Hall--Borcherds amalgamation]
The completed amalgam exists and represents pairs of continuous Lie morphisms from $\mathfrak b_X$ and $\mathfrak b_\Delta$ that agree on $\mathfrak h_\Pi$.  The Cartan projections of the two Borels induce canonical retractions
\[
 r_X:\mathfrak b_{X\boxtimes\Delta}\longrightarrow\mathfrak b_X,
 \qquad
 r_\Delta:\mathfrak b_{X\boxtimes\Delta}\longrightarrow\mathfrak b_\Delta.
\]
Its completed enveloping algebra is the amalgamated free product
\[
 \widehat U(\mathfrak b_{X\boxtimes\Delta})
 \cong
 \widehat U(\mathfrak b_X)
 \widehat *_{\widehat U(\mathfrak h_\Pi)}
 \widehat U(\mathfrak b_\Delta).
\]
The geometric Hall operations survive under $r_X$, while the Igusa bracket, PBW character, and automorphic determinant survive under $r_\Delta$.
\end{theorem}
\begin{proof}
Construct the ordinary Lie coproduct from the free product of the two Borels and divide by the ideal identifying their two Cartan copies.  Complete by total positive charge and proper root height.  Every fixed filtered degree uses finitely many Lie words, so the completion is separated and has the stated mapping property.  The projection $\mathfrak b_X\to\mathfrak h_\Pi\to\mathfrak b_\Delta$ together with the identity on $\mathfrak b_\Delta$ gives $r_\Delta$ by universality; the second retraction is symmetric.  The enveloping functor preserves coproducts of Lie algebras, and PBW identifies the completed algebraic filtration with the displayed completed free product.
\end{proof}

\begin{definition}[Comparison ideal]
A primitive comparison $\varphi:\nDelta\to\mathfrak p_X^+$ preserving Cartan degree defines the closed ideal
\[
 \mathcal J_\varphi
 =\overline{\angles{
   i_\Delta(x)-i_X(\varphi(x))\mid x\in\nDelta
 }}
 \subset\mathfrak b_{X\boxtimes\Delta}.
\]
The diagonal realization is
\[
 \mathfrak b_{X\simeq\Delta,\varphi}
 =\mathfrak b_{X\boxtimes\Delta}/\mathcal J_\varphi.
\]
\end{definition}

\begin{proposition}[Graph quotient of a primitive comparison]\label{cyiv:graph-quotient}
Let $\varphi:\mathfrak b_\Delta\to\mathfrak b_X$ be a continuous Lie
morphism restricting to the identity on $\mathfrak h_\Pi$.
Use a completion in which the amalgam and closed quotient have
their stated universal properties.
Then
\[
 \mathfrak b_{X\boxtimes\Delta}/\mathcal J_\varphi
 \cong\mathfrak b_X.
\]
The two factor maps become $\id$ and $\varphi$.
If $\varphi$ is an isomorphism, the quotient is also isomorphic
to $\mathfrak b_\Delta$.
\end{proposition}
\begin{proof}
The maps $\id$ and $\varphi$ agree on $\mathfrak h_\Pi$.
They give $q_\varphi:\mathfrak b_{X\boxtimes\Delta}\to\mathfrak b_X$,
which kills the generators of $\mathcal J_\varphi$.
The inclusion of $\mathfrak b_X$ gives a map $s$ in the other direction.
The composite $q_\varphi s$ is the identity.
The other composite is the identity on the geometric factor.
On the Borcherds factor, the defining relation identifies $x$ with
$s\varphi(x)$, so it is also the identity there.
The two factor images topologically generate the quotient.
Continuity proves that the maps are inverse on the completion.
\end{proof}

The graph quotient implements a specified primitive comparison.
Constructing that comparison requires geometric simple states whose
extension brackets satisfy the Borcherds relations.
The quotient formula alone produces neither those states nor their
orientation data.

\chapter{Coproducts, pairings, and doubles}
A positive Hall algebra becomes a quantum group after additional structure is supplied. Let $H^+$ be a complete graded bialgebra, let $H^0$ be a Cartan algebra, and let $H^-$ be the opposite or dual half. Let
\[
 \angles{-,-}:H^+\widehattensor H^-\longrightarrow\kfield
\]
be a nondegenerate Hopf pairing.

\begin{definition}[Completed Hall double]
The completed Drinfeld double is the topological algebra
\[
 \widehat D(H^+)=H^-\widehattensor H^0\widehattensor H^+
\]
with internal products in the three factors and cross-relations determined by the Hopf pairing. Its representation category is monoidal by the coproduct and braided by the universal $R$-matrix whenever the canonical tensor converges.
\end{definition}

\begin{theorem}[Double construction]
Let $H^+$ and $H^-$ be paired complete graded bialgebras, let $H^0$ be a Cartan bialgebra acting by the stated weight grading, and let
\[
 \angles{-,-}:H^+\widehattensor H^-\longrightarrow\kfield
\]
be a continuous nondegenerate Hopf pairing. Assume the triangular multiplication map from $H^-\widehattensor H^0\widehattensor H^+$ is separated and complete. Then the pairing cross-relations determine a unique completed double with this triangular decomposition. On every continuous module category on which the canonical tensors converge, the completed canonical element gives the universal $R$-matrix.
\end{theorem}
\begin{proof}
The Hopf pairing defines the cross-relations between the paired halves and makes the three canonical embeddings bialgebra morphisms. The Hopf identities identify the two reductions of every triple product, proving associativity. Separated triangular decomposition gives uniqueness. Finite-degree canonical tensors form a compatible inverse system; convergence on the selected module category yields the quasitriangular element \cite{DrinfeldDouble,EtingofSchiffmann}.
\end{proof}

\begin{theorem}[Quantum Hall--Borcherds amalgam]
Let $H_X$ and $H_\Delta$ be complete Hopf algebras containing the same complete Cartan Hopf algebra $H^0$, with continuous Hopf retractions onto $H^0$.  Their completed amalgamated free product
\[
 H_{X\boxtimes\Delta}
 =H_X\widehat *_{H^0}H_\Delta
\]
has a unique complete Hopf structure extending the two coproducts, counits, and antipodes.  It represents pairs of continuous Hopf morphisms agreeing on $H^0$ and has canonical Hopf retractions to both factors.  A mixed universal $R$-matrix is equivalent to a continuous cross-pairing satisfying the two hexagon identities; the intrinsic $R$-matrices of the two retracts remain defined independently.
\end{theorem}
\begin{proof}
The algebraic amalgamated free product exists because complete filtered algebras are obtained degreewise from finite quotients.  The two coproducts agree on $H^0$, so the universal property extends them to a coproduct on the free product.  Coassociativity is checked on the two generating Hopf subalgebras.  The same argument extends the counit and antipode.  The Hopf retractions are induced by the factor retractions onto $H^0$.  A quasitriangular structure extending both factors requires the additional commutation data between the two generating halves; writing this data as a cross-pairing turns the quasitriangular and Yang--Baxter equations into the two hexagon identities.
\end{proof}

The deconcatenation coproduct on a bar complex and the representation-theoretic coproduct on a quantum group occupy different structures. A geometric comparison between them is a theorem about a selected model.

\chapter{Ordered chiral bialgebras}
A tensor product of representations requires compatible actions on
threefold products and the unit.
Retain the factorization coefficient category
$(\Fact_D(C),\tensor^!)$ on the smooth complex curve $C$.
Fix a category $\widehat{\mathcal V}_C$ of complete filtered objects
with specified presentations $F\simeq\varprojlim_r F_r$,
where $F_r\in\Fact_D(C)$.
Use derived limits in the fixed Ran enhancement and define
\[
 F\widehat{\tensor}^{!}G
 =\varprojlim_{r,s}(F_r\tensor^!G_s).
\]
Assume these limits remain factorizing and give
$\widehat{\mathcal V}_C$ a symmetric monoidal structure with unit $\mathbf1$.
Its associativity, symmetry, and unit equivalences must preserve
the presentations and the Ran structure.
Morphisms are continuous maps compatible with these presentations.
These are hypotheses on the chosen completion.

\begin{definition}[Ordered chiral bialgebra]\label{def:vol4-ordered-chiral-bialgebra}
An ordered chiral bialgebra is a counital $\Eone$-coalgebra object
in complete ordered chiral algebras:
\[
 A\in\operatorname{CoAlg}_{\Eone}
 \left(\Alg_{\Eone}
 (\widehat{\mathcal V}_C,\widehat{\tensor}^{!})\right).
\]
It has continuous multiplication $m$ and unit $\eta$, and continuous
morphisms of unital ordered chiral algebras
\[
 \Delta:A\longrightarrow A\widehat{\tensor}^{!}A,
 \qquad \epsilon:A\longrightarrow\mathbf1.
\]
The counit is the chosen augmentation.
Write $\tau$ for the coefficient symmetry, including Koszul signs.
Suppressing ambient associativity and unit equivalences, compatibility reads
\[
\begin{aligned}
 \Delta m&=(m\widehat{\tensor}^{!}m)
  (\id\widehat{\tensor}^{!}\tau\widehat{\tensor}^{!}\id)
  (\Delta\widehat{\tensor}^{!}\Delta),\\
 \Delta\eta&=\eta\widehat{\tensor}^{!}\eta,
 \qquad \epsilon m=\epsilon\widehat{\tensor}^{!}\epsilon,
 \qquad \epsilon\eta=\id_{\mathbf1}.
\end{aligned}
\]
Coassociativity and the counit identities read
\[
\begin{aligned}
 (\Delta\widehat{\tensor}^{!}\id)\Delta
   &=(\id\widehat{\tensor}^{!}\Delta)\Delta,\\
 (\epsilon\widehat{\tensor}^{!}\id)\Delta
   &=\id_A=(\id\widehat{\tensor}^{!}\epsilon)\Delta.
\end{aligned}
\]
For a homotopy-coherent presentation, these identities include the
higher coherences of the displayed coalgebra object.
Every operation and coherence preserves diagonal transitions and
disjoint-locus factorization maps.
An antipode is a continuous convolution inverse $S:A\to A$ of $\id_A$:
\[
 m(S\widehat{\tensor}^{!}\id)\Delta
 =\eta\epsilon
 =m(\id\widehat{\tensor}^{!}S)\Delta.
\]
A bialgebra with an antipode is a Hopf object in this coefficient category.
\end{definition}

A continuous left $A$-module is an object $M\in\widehat{\mathcal V}_C$
with a unital $\Eone$ action
$\rho_M:A\widehat{\tensor}^{!}M\to M$.
Its action and coherences are morphisms in these factorization coefficients.
Write $\operatorname{Mod}_A(\widehat{\mathcal V}_C)$ for their category.

\begin{theorem}[Representation fusion]\label{thm:vol4-representation-fusion}
An ordered chiral bialgebra makes
$\operatorname{Mod}_A(\widehat{\mathcal V}_C)$ monoidal.
Its product is $M\widehat{\tensor}^{!}N$ with action
\[
\begin{aligned}
 A\widehat{\tensor}^{!}(M\widehat{\tensor}^{!}N)
 &\xrightarrow{\Delta\widehat{\tensor}^{!}\id}
 (A\widehat{\tensor}^{!}A)\widehat{\tensor}^{!}
 (M\widehat{\tensor}^{!}N)\\
 &\xrightarrow{\id\widehat{\tensor}^{!}\tau_{A,M}
                       \widehat{\tensor}^{!}\id}
 (A\widehat{\tensor}^{!}M)\widehat{\tensor}^{!}
 (A\widehat{\tensor}^{!}N)\\
 &\xrightarrow{\rho_M\widehat{\tensor}^{!}\rho_N}
 M\widehat{\tensor}^{!}N.
\end{aligned}
\]
The unit is $\mathbf1$ with action induced by $\epsilon$.
The ambient associator and unit equivalences are module equivalences.

Suppose a continuous universal $R$-matrix is also given.
It is an invertible morphism
$R:\mathbf1\to A\widehat{\tensor}^{!}A$ for the tensor-product algebra,
satisfying
\[
\begin{gathered}
 R\Delta(a)=\Delta^{\op}(a)R,\qquad
 \Delta^{\op}=\tau_{A,A}\Delta,\\
 (\Delta\widehat{\tensor}^{!}\id)(R)=R_{13}R_{23},
 \qquad
 (\id\widehat{\tensor}^{!}\Delta)(R)=R_{13}R_{12},\\
 (\epsilon\widehat{\tensor}^{!}\id)(R)=\eta
 =(\id\widehat{\tensor}^{!}\epsilon)(R).
\end{gathered}
\]
The first equation is an identity of morphisms from $A$ to its completed tensor square.
Leg notation uses the unit and symmetry of $\widehat{\mathcal V}_C$.
In a homotopy-coherent presentation, require compatible lifts of these identities.
Then $c_{M,N}=\tau_{M,N}\circ R_{M,N}$ is a braiding,
where $R_{M,N}$ acts through the external tensor-product module action.
\end{theorem}
\begin{proof}
The symmetry combines the independent module actions into an
$A\widehat{\tensor}^{!}A$-action.
Restriction along the unital algebra map $\Delta$ gives the displayed action.
For three modules, the two actions use
$(\Delta\widehat{\tensor}^{!}\id)\Delta$ and
$(\id\widehat{\tensor}^{!}\Delta)\Delta$.
Coassociativity makes the ambient associator $A$-linear.
The counit identities make the ambient unit equivalences $A$-linear.
The coalgebra coherences give the pentagon and triangle coherences.
The completion hypotheses keep these maps continuous and factorizing.
Quasitriangularity makes $c_{M,N}$ an intertwiner and invertibility makes it an equivalence.
The two coproduct identities for $R$ give the hexagons.
Its counit identities give compatibility with the unit.
\end{proof}

The tensor in this construction is the completed transverse tensor
$\tensor^!$ on factorization coefficients.
A chiral-convolution or Hall coproduct requires its own construction
and a comparison in the specified category.
A spectral version additionally requires a difference coordinate,
a Laurent-series completion, and compatible domains for iterated expansions.
Its identities must hold in those domains before defining module actions.

\begin{example}[A counital algebra map without coassociativity]
In ordinary complex vector spaces, let $A=\C[x]$ and put
\[
 \Delta(x)=x\tensor1+1\tensor x+x\tensor x^2,
 \qquad \epsilon(x)=0.
\]
Polynomial substitution extends these assignments to unital algebra maps.
Writing $F(u,v)=u+v+uv^2$, the equations
$F(0,v)=v$ and $F(u,0)=u$ prove both counit identities.
But
\[
\begin{split}
 F(F(u,v),w)-F(u,F(v,w))
 &=-uv^2w^4-uv^2w^2-2uvw^3-2uvw\\
 &=-uvw(1+w^2)(vw+2)\ne0.
\end{split}
\]
On three regular modules, the two actions differ already on
$1\tensor1\tensor1$.
Thus the canonical tensor associator is not $A$-linear.
This ordinary example shows why multiplicativity and the counit identities
must be supplemented by coassociativity.
It makes no assertion about a geometric realization in factorization coefficients.
\end{example}

\part{Exact model cases}

\chapter{The \texorpdfstring{$A_2$}{A2} Hall algebra}
For the quiver $1\to2$ over $\mathbb F_q$, the twisted Hall algebra has generators $E_1,E_2$ and $v^2=q$. Counting subrepresentations gives
\[
 E_1^2E_2-(v+v^{-1})E_1E_2E_1+E_2E_1^2=0.
\]
For a representation of dimension vector $(2,1)$, the arrow has rank zero or one. The relevant flag counts are
\[
\begin{array}{c|ccc}
\rank&112&121&211\\ \hline
0&q+1&q+1&q+1\\
1&q+1&1&0.
\end{array}
\]
After the Hall twists $v^{-1},1,v$, both rows give zero. This is the elementary geometric origin of the quantum Serre relation \cite{Ringel,Green}.

\chapter{The tripled Jordan quiver and \texorpdfstring{$\mathbb C^3$}{C3}}
Let $Q$ have one vertex and loops $x,y,z$, with potential
\[
 W=\Tr x[y,z].
\]
The Jacobi algebra is $\C[x,y,z]$. The critical CoHA uses the vanishing cycles of $\Tr W$ on the representation spaces. Its positive currents are generated by the one-vertex dimension sectors.

\begin{theorem}[Affine-Yangian positive half]
After the standard equivariant localization and completion, the critical CoHA of the tripled Jordan quiver is identified with the positive half of the affine Yangian of $\mathfrak{gl}_1$. The independently defined affine Yangian carries its standard Cartan and negative current halves; the positive-half comparison embeds the CoHA into that triangular presentation.
\end{theorem}
\begin{proof}
The CoHA product becomes the shuffle product under localization. Dimensional reduction identifies the tripled-quiver critical theory with the equivariant surface CoHA, and Davison's affine-BPS theorem identifies its completed positive algebra with the strictly positive affine Yangian. The shuffle kernel satisfies the affine-Yangian current relations with equivariant parameters $\epsilon_1,\epsilon_2,\epsilon_3$ and Calabi--Yau relation $\epsilon_1+\epsilon_2+\epsilon_3=0$. PBW and current comparison identify the positive algebra. The independently presented full affine Yangian is the triangular algebra generated by its negative currents, Cartan currents, and this positive half \cite{DavisonAffineBPS,DavisonMeinhardt,SVYangians,Tsymbaliuk}.
\end{proof}

The vacuum representation produces a $W_{1+\infty}$-type vertex algebra through an evaluation homomorphism. The CoHA, its double, and the vacuum algebra occupy successive levels: positive algebra, paired quantum group, and distinguished representation.

\chapter{Surface CoHAs and ADE Yangians}
Let $S_\Gamma\to\C^2/\Gamma$ be the minimal resolution for a finite subgroup $\Gamma\subset\operatorname{SL}_2(\C)$. The exceptional curves have ADE intersection matrix. Moduli of sheaves supported on these curves carry convolution correspondences.

\begin{theorem}[Curve-supported surface CoHA]
The completed positive CoHA of one-dimensional sheaves supported on the exceptional curves is identified with a nonstandard positive half of the affine Yangian of the corresponding ADE Lie algebra.
\end{theorem}
\begin{proof}
The theorem is the curve-supported surface construction of Diaconescu--Porta--Sala--Schiffmann--Vasserot \cite{DPPSSV}. The simple curve sheaves give the simple generators, Hecke correspondences give the current operators, and the completed PBW comparison gives the affine Yangian presentation.
\end{proof}

This theorem is a local K3 model whenever an ADE configuration occurs on a K3 surface. The global K3 surface adds compact curve classes and monodromy beyond the local root lattice.

\chapter{Coulomb branches and shifted Yangians}
For a three-dimensional $\mathcal N=4$ gauge theory, the BFN Coulomb branch is defined by equivariant Borel--Moore homology of an affine-Grassmannian correspondence. Its convolution algebra quantizes the Coulomb branch \cite{BFN}. In quiver gauge theories, shifted Yangians and their truncations provide presentations of these quantizations. Boundary conditions select the shift and truncation data.

The BFN construction supplies a second geometric source of Yangians. Its grading and coproduct arise from loop rotation and convolution, while the surface CoHA source arises from sheaf extensions. A comparison between the two is a model-specific equivalence.

\part{K3 symmetries}

\chapter{The Mukai lattice}
For a K3 surface $S$, the Mukai lattice is
\[
 \Mukai(S,\Z)=H^0(S,\Z)\oplus H^2(S,\Z)\oplus H^4(S,\Z)
\]
with pairing
\[
 ((r,c,s),(r',c',s'))=c\cdot c'-rs'-r's.
\]
It is even unimodular of signature $(4,20)$ and rank $24$.

\chapter{The Nakajima Heisenberg algebra}
Let
\[
 \mathbb H_S=\bigoplus_{n\ge0}H^*(S^{[n]}).
\]
For $m\in\Z\setminus\{0\}$ and $\alpha\in H^*(S)$, the Nakajima correspondence gives an operator $p_m(\alpha)$.

\begin{theorem}[K3 Heisenberg action]
The operators satisfy
\[
 [p_m(\alpha),p_n(\beta)]
 =m\,\delta_{m+n,0}\int_S\alpha\cup\beta\;\id.
\]
They generate the rank-twenty-four Heisenberg algebra acting irreducibly on $\mathbb H_S$ with vacuum $1\in H^*(S^{[0]})$.
\end{theorem}
\begin{proof}
The commutator is the excess intersection of the two incidence correspondences. Opposite creation and annihilation correspondences meet along the diagonal with multiplicity $m$ and intersection pairing $\int_S\alpha\beta$. Correspondences of the same sign commute. The vacuum construction gives the Fock representation \cite{Nakajima97}.
\end{proof}

This action is global, compact, and abelian at the Lie level. It is an exact K3 symmetry and supplies the rank-twenty-four Fock sector of any broader K3 algebra.

\chapter{The orthogonal LLV algebra}
For a compact hyperkähler manifold, every Lefschetz class gives an $\mathfrak{sl}_2$ triple on cohomology. The Lie algebra generated by all such triples is the Looijenga--Lunts--Verbitsky algebra.

\begin{theorem}[Orthogonal LLV symmetries]
For a K3 surface $S$, the LLV algebra acting on $H^*(S)$ is
\[
 \mathfrak g_{\mathrm{LLV}}(S)
 \cong\mathfrak{so}\bigl(\widetilde H(S),(-,-)_{\mathrm{Muk}}\bigr)
 \cong\mathfrak{so}(4,20).
\]
For an irreducible hyperk\"ahler manifold $M$ of $K3^{[n]}$ type with $n\ge2$, one has $b_2(M)=23$ and
\[
 \mathfrak g_{\mathrm{LLV}}(M)\cong\mathfrak{so}(4,21).
\]
Moduli spaces of stable objects on $S$ carry their own LLV actions; Mukai and monodromy correspondences relate these actions to the K3 lattice.
\end{theorem}
\begin{proof}
The Looijenga--Lunts--Verbitsky theorem identifies the Lie algebra generated by all Lefschetz triples with the orthogonal algebra of the Mukai extension $H^2(M)\oplus U$. For $S$, the extension has signature $(4,20)$. For $K3^{[n]}$ type with $n\ge2$, the Beauville--Bogomolov lattice has signature $(3,20)$, so adjoining the hyperbolic plane gives signature $(4,21)$ \cite{LooijengaLunts,Verbitsky96}.
\end{proof}

The LLV algebra is a classical cohomological symmetry. A Yangian deformation of the full orthogonal action consists of two further structures: a Drinfeld current presentation and a compatible spectral coproduct.

\chapter{Local ADE Yangians}
An ADE configuration of $(-2)$-curves on $S$ gives a root sublattice of $H^2(S,\Z)$. Sheaves supported on the configuration produce the surface CoHA of the preceding part and hence a completed positive affine ADE Yangian. Nakajima quiver varieties supply compatible Yangian actions on their equivariant cohomology.

Thus K3 geometry contains local nonabelian Yangian sectors indexed by ADE enhancements of the Picard lattice. Picard monodromy transports the ADE root sublattices over the corresponding Noether--Lefschetz loci. A family of sheaf moduli and Hecke correspondences compatible with that monodromy transports the Yangian actions themselves.

\chapter{The K3 symmetry trichotomy}
The three exact structures are
\begin{center}
\small
\begin{tabularx}{0.96\textwidth}{@{}>{\raggedright\arraybackslash}p{0.23\textwidth}>{\raggedright\arraybackslash}X>{\raggedright\arraybackslash}p{0.28\textwidth}@{}}
\toprule
\textbf{Structure}&\textbf{Carrier}&\textbf{Operation}\\
\midrule
Mukai Heisenberg&$\bigoplus_nH^*(S^{[n]})$&creation and annihilation\\
LLV $\mathfrak{so}(4,20)$&$H^*(S)$; separately $\mathfrak{so}(4,21)$ on $K3^{[n]}$ type&Lefschetz Lie action\\
Local ADE Yangian&curve-supported equivariant moduli&Hecke and current correspondences\\
\bottomrule
\end{tabularx}
\end{center}
A global K3 quantum algebra is a system that contains these three actions with explicit comparison homomorphisms. The two-Calabi--Yau BPS theorem supplies a global generalized Kac--Moody positive algebra. A full Yangian or double is specified by a compatible coproduct, spectral parameter, negative half, Cartan part, pairing, and completion.

\chapter{K3 correspondences on fixed coefficients}
Fix rational cohomology and a collection $\mathscr M_S$ of smooth
projective moduli spaces on a projective K3 surface $S$.
Include the Hilbert schemes $S^{[n]}$ and close the collection under
finite products.
For $M,N\in\mathscr M_S$, define the graded kernel space
\[
 K(M,N)=H^*(M\times N,\mathbb Q)[2\dim_{\mathbb C}M].
\]
A kernel of degree $r$ acts from $H^j(M)$ to $H^{j+r}(N)$ by
\[
 T_\alpha(v)=(p_N)_*(p_M^*v\cup\alpha).
\]
The shift in $K(M,N)$ accounts for integration over $M$.
All pushforwards in this construction are proper.

\begin{proposition}[Kernel composition]\label{cyiv:k3-kernels}
The product
\[
 \beta\circ\alpha=(p_{13})_*
 (p_{12}^*\alpha\cup p_{23}^*\beta)
\]
defines an associative graded category with diagonal identities.
The kernel action identifies $K(M,N)$ with
$\operatorname{Hom}^*(H^*(M),H^*(N))$.
\end{proposition}
\begin{proof}
The K\"unneth isomorphism and Poincar\'e duality identify
$H^*(M\times N)[2\dim M]$ with $H^*(M)^\vee\otimes H^*(N)$.
Under this identification, $T_\alpha$ is the associated linear map.
Integration over the middle factor contracts its dual pair, so
$T_{\beta\circ\alpha}=T_\beta T_\alpha$.
Faithfulness then proves associativity and the diagonal identity.
External products give tensor products with their Koszul signs.
\end{proof}

A geometric symmetry is specified by selected kernels in this category.
Nakajima incidence kernels connect different Hilbert schemes.
Lefschetz multiplication and its adjoint act on each individual
cohomology space.
Using every cohomology kernel would instead give all linear maps;
the selected geometric kernels carry the representation-theoretic content.

\chapter{The algebra generated by geometric actions}
Put $V_S=\bigoplus_{M\in\mathscr M_S}H^*(M,\mathbb Q)$.
Choose a set of homogeneous operators $T_\lambda$ on $V_S$, each sending
a vector with finite support to a vector with finite support.
The Nakajima operators on the Hilbert-scheme summand satisfy this
condition: each changes the point number by a fixed integer.
Individual Lefschetz operators, extended by zero on other summands,
also satisfy it.
Let $F$ be the free unital graded algebra on symbols $t_\lambda$.

\begin{proposition}[Geometric action algebra]\label{cyiv:k3-action}
The map
\[
 \rho:F\longrightarrow\operatorname{End}_{\mathbb Q}(V_S),
 \qquad t_\lambda\longmapsto T_\lambda
\]
defines a faithful action of $A_S=F/\ker\rho$.
For every algebra $B$, maps $A_S\to B$ are exactly maps $F\to B$
that kill $\ker\rho$.
Products and commutators of the selected kernels are their convolution
products and commutators on the relevant moduli spaces.
\end{proposition}
\begin{proof}
Finite words in the selected operators preserve finite support, so
$\rho$ is defined by composition.
The quotient acts faithfully by definition of its kernel.
The factorization property is the universal property of an algebra
quotient.
The convolution assertion follows from Proposition~\ref{cyiv:k3-kernels}.
\end{proof}

An equivariant ADE construction uses a separate coefficient field,
for example $\operatorname{Frac}H_T^*(\mathrm{pt},\mathbb Q)$,
and its prescribed localization and proper pushforwards.
Joining it to $A_S$ requires maps between the actual state spaces that
intertwine those operators and their coefficients.
A completion of $A_S$ requires a specified topology and continuity of
every generator.
The indefinite Mukai norm alone supplies neither condition.
The desired common K3 quantum algebra must construct these local-to-global
actions and their mixed convolution relations.

\chapter{Tensor reconstruction and geometric kernels}
Let $\mathscr C$ be an essentially small rigid linear monoidal category
and let $\omega:\mathscr C\to\operatorname{Vect}_{\mathrm{fd}}(\mathbb Q)$
be a strong monoidal functor.
A choice of geometric kernels specifies the morphisms of $\mathscr C$;
it is part of the input.
Define
\[
 H=\left(\bigoplus_{M\in\mathscr C}\omega(M)^*\otimes\omega(M)\right)
 \big/\left\langle
 [\lambda\circ\omega(f)\otimes v]_M
 -[\lambda\otimes\omega(f)v]_N\right\rangle,
\]
where $f:M\to N$, $v\in\omega(M)$, and $\lambda\in\omega(N)^*$.

\begin{proposition}[Matrix-coefficient coalgebra]\label{cyiv:k3-coend}
For a basis $(e_i)$ with dual basis $(e^i)$, the formulas
\[
 \Delta[\lambda\otimes v]_M
 =\sum_i[\lambda\otimes e_i]_M\otimes[e^i\otimes v]_M,
 \qquad \epsilon[\lambda\otimes v]_M=\lambda(v)
\]
are basis-independent and make $H$ a coalgebra.
Tensor product gives its multiplication, and duals give its antipode.
Each $\omega(M)$ is a comodule by
\[
 v\longmapsto\sum_i e_i\otimes[e^i\otimes v]_M.
\]
Every selected kernel is a comodule morphism.
\end{proposition}
\begin{proof}
The identity tensor $\sum_i e_i\otimes e^i$ is basis-independent.
Inserting it twice proves coassociativity; evaluation proves both
counit identities.
The defining relations express naturality under $f$ and are preserved
by these formulas.
The strong monoidal constraints identify the coefficient of $M\otimes N$
with the product of coefficients of $M$ and $N$.
Their pentagon and unit identities give the bialgebra axioms.
The evaluation and coevaluation identities for duals give the antipode
identities.
The same relations prove that $\omega(f)$ intertwines the coactions.
\end{proof}

The linear dual of $H$ represents linear natural endomorphisms of
$\omega$.
Tensor-compatible natural automorphisms are the corresponding
multiplicative, unit-preserving families, not all elements of this dual.
An equivalence with a comodule category requires the additional exactness
and reconstruction hypotheses stated in the Tannaka theorem below.
The action algebra $A_S$ acts by selected kernels; the coend records
symmetries commuting with those kernels.

\part{The exact \texorpdfstring{$K3\times E$}{K3 x E} protected completion}

\chapter{The protected input}
The K3 elliptic genus is $2\phi_{0,1}$.
The Borcherds half determines the weight-five denominator $D_5$ and its
Borcherds Lie superalgebra $\mathfrak g_\Delta$.
Fix its positive root spaces with their parity and proper height.
The positive algebra and its root Fock space are defined from this Lie
superalgebra, independently of a compact Hall realization.

For each finite zero-potential quiver window $Q_{\Delta,H}$, form the
Ginzburg algebra $\Gamma_{\Delta,H}$.
Restriction to full subquivers gives the projections of
Proposition~\ref{cyiv:ginzburg-projection}.
Write $\Gamma_\Delta$ for this specified inverse system.
A nondegenerate continuous Calabi--Yau structure on a completed limit
is additional data.
An ordinary Hall category $\mathcal B_\Delta$ must carry its chosen
quiver coefficients and extension correspondences.
Its identification with a root presentation requires the corresponding
Hall theorem on that precise datum.

\begin{definition}[Root-datum realizations]
A protected root-datum realization records
\[
 \mathfrak P(K3\times E)=
 (\Gamma_\Delta,\mathcal B_\Delta,\mathcal H_{\mathrm{1p}},
 \mathcal F^{\mathrm{hyb}}_{\Delta,E},\widehat U_q(\mathfrak g_\Delta),
 \mathbb A_Z).
\]
Here $\mathcal H_{\mathrm{1p}}=\bigoplus_{\alpha>0}(\mathfrak g_\Delta)_\alpha$
with its parity.
The current-factorization, paired quantum-group, and Fredholm entries
retain their own coefficient categories, completions, and defining
hypotheses.
Maps among these entries are specified realization maps, not consequences
of using the same root labels.
\end{definition}

The enveloping and Chevalley constructions are functorial under the
primitive maps of Theorem~\ref{cyiv:primitive-comparison}.
The root Fock trace is computed from the actual graded root spaces.
A Fredholm realization adds a family of operators on a specified domain
and the convergence and orientation data identifying its Pfaffian with
the denominator.
These are the analytic inputs to the square-root comparison below.

\begin{definition}[Integral protected realization]
Let $\mathfrak p_X^+$ be the BPS primitive Lie algebra of a selected strong-oriented cohomological lift of the geometric $E$-equivariant Hall source.  The integral protected realization is
\[
 \mathfrak P_{\mathrm{int}}(X,\Delta)
 =\left(
   \Hcal_X^{E,\mathrm{mot}},
   \mathfrak P(K3\times E),
   \mathfrak b_{X\boxtimes\Delta},
   r_X,r_\Delta,
   I_{X,\Pi}^E
 \right).
\]
At the motivic level the first, second, and character components are defined.  A strong orientation enhancement supplies $\mathfrak p_X^+$ and hence the amalgamated primitive component.
\end{definition}

\begin{proposition}[Primitive graph comparison]
Suppose the two primitive Borels, their common Cartan action, and their
completed amalgam are supplied as above.
Their Cartan retractions induce retractions from the amalgam to each
factor.
For each specified primitive map $\varphi$ preserving Cartan action,
its graph quotient is canonically the geometric Borel.
\end{proposition}
\begin{proof}
The identity on one Borel and the Cartan projection of the other agree
on the common Cartan, so the amalgam gives the retraction.
The quotient statement is Proposition~\ref{cyiv:graph-quotient}.
\end{proof}

The motivic integration map and primitive comparison are different maps.
Strong orientation gives the coefficient coherence of a cohomological
Hall construction under its stated geometric hypotheses.
Extracting a primitive Lie algebra additionally requires the appropriate
BPS or coproduct construction.

\chapter{Cyclic open-string operations}
Let $A$ be a finite-dimensional, uncurved cyclic $A_\infty$ algebra
over $\mathbb C$, with a nondegenerate pairing of degree $-3$.
Use the cohomological suspension $s:A\to A[1]$ of degree $-1$.
The suspended operations $b_n:(A[1])^{\otimes n}\to A[1]$ have degree one.
Our sign convention is the coderivation convention $b^2=0$ on the
tensor coalgebra, with the Koszul rule for moving homogeneous inputs.
Write $\omega$ for the resulting degree-$-1$ symplectic form on $A[1]$.
Fix the Hamiltonian convention $\iota_{X_F}\omega=dF$ and
$[X_F,X_G]=X_{\{F,G\}}$.

For the universal degree-zero coordinate $x$ on $A[1]$, put
\[
 S(x)=\sum_{n\ge1}\frac1{n+1}\omega(x,b_n(x^{\otimes n})).
\]
Cyclic invariance is read with the suspended Koszul signs.
All functions are formal power series at the origin.
An equivalent unsuspended convention writes this as
$\frac12\langle a,m_1a\rangle+
\sum_{n\ge2}\frac1{n+1}\langle a,m_n(a^{\otimes n})\rangle$.

\begin{proposition}[Cyclic master equation]\label{cyiv:cyclic-master}
The action satisfies $\{S,S\}=0$.
\end{proposition}
\begin{proof}
Cyclicity makes the $n+1$ derivatives of the arity-$n+1$ term equal,
with the suspended Koszul signs.
Thus $X_S(x)=\sum_{n\ge1}b_n(x^{\otimes n})$.
The coderivation identity implies $X_S^2=0$.
The Hamiltonian identity gives $X_{\{S,S\}}=2X_S^2=0$.
Nondegeneracy of $\omega$ makes $\{S,S\}$ constant.
Uncurvedness removes $b_0$, so $S$ starts in order two and its bracket
has no constant term.
The constant is therefore zero.
\end{proof}

A finite-height Ginzburg realization requires such a cyclic model on
the selected proper objects before this scalar action is applied.
A Hall realization additionally requires a moduli problem with
critical coefficient complexes.
For an extension $A_0\to F\to B_0$, put
$L(A_0,B_0)=\det R\operatorname{Hom}(A_0,B_0)$.
Three-Calabi--Yau duality gives
\[
 R\operatorname{Hom}(B_0,A_0)
 \simeq R\operatorname{Hom}(A_0,B_0)^\vee[-3].
\]
Applying the determinant functor to the two derived Hom triangles gives
\[
 \det R\operatorname{Hom}(F,F)
 \cong \det R\operatorname{Hom}(A_0,A_0)
 \otimes\det R\operatorname{Hom}(B_0,B_0)
 \otimes L(A_0,B_0)^{\otimes2}.
\]
Indeed, the four subquotients are the two diagonal Hom complexes and
the two cross complexes.
Dualization and the odd shift cancel on their determinant lines,
so the two cross determinants coincide.
Orientation lines require square-root isomorphisms along this formula,
compatible on triple flags and with any quotient descent.
The cyclic master equation does not select these isomorphisms.

\chapter{An algebraic open--closed comparison}
A finite-dimensional ordinary Lie algebra $\mathfrak l$ gives a
precise coefficient model for an enveloping-algebra centre.
Let $U=U(\mathfrak l)$ over $\mathbb C$.
The adjoint module $U^{\mathrm{ad}}$ has action
$x\cdot u=xu-ux$.
Its Chevalley cochains are
\[
 C^n(\mathfrak l,U^{\mathrm{ad}})
 =\operatorname{Hom}_{\mathbb C}(\Lambda^n\mathfrak l,U),
\]
with differential
\[
\begin{split}
 (df)(x_1,\ldots,x_{n+1})
 &=\sum_i(-1)^{i+1}[x_i,f(x_1,\widehat x_i,\ldots,x_{n+1})]\\
 &\quad+\sum_{i<j}(-1)^{i+j}
 f([x_i,x_j],x_1,\widehat x_i,\widehat x_j,\ldots,x_{n+1}).
\end{split}
\]
The second line places the bracket first and retains the order of the
remaining arguments.

\begin{theorem}[Hochschild--Chevalley comparison]\label{cyiv:hh-ce}
Restriction followed by antisymmetrization defines a quasi-isomorphism
\[
 C^\bullet_{\mathrm{HH}}(U,U)
 \longrightarrow C^\bullet(\mathfrak l,U^{\mathrm{ad}}),
 \qquad
 \phi\longmapsto\left[(x_1,\ldots,x_n)\mapsto
 \sum_{\sigma\in S_n}\operatorname{sgn}(\sigma)
 \phi(x_{\sigma(1)},\ldots,x_{\sigma(n)})\right].
\]
This is a comparison of ordinary cochain complexes.
\end{theorem}
\begin{proof}
The Chevalley bimodule resolution has terms
$K_n=U\otimes\Lambda^n\mathfrak l\otimes U$.
Its differential has the left-minus-right multiplication terms and
the bracket terms with the signs dual to the displayed differential.
Filter $K_n$ by total PBW degree, counting each exterior generator
as degree one.
The multiplication terms preserve this degree and bracket terms lower it.
The associated graded is the Koszul resolution of the diagonal in
$\operatorname{Sym}(\mathfrak l)\otimes\operatorname{Sym}(\mathfrak l)$,
with regular sequence $x_i\otimes1-1\otimes x_i$.
That Koszul complex is exact away from degree zero.
The filtration is bounded below and exhaustive, and each element has
finite filtration degree, so lifting a leading term proves exactness
of $K_\bullet$ by induction.
Its augmentation has image $U$.
Thus $K_\bullet$ is a free bimodule resolution.

Antisymmetrization gives a map from $K_\bullet$ to the ordinary bar
resolution, lifting the identity of $U$.
Adjacent products in its bar differential pair to commutators;
$xy-yx=[x,y]$ gives exactly the bracket terms of $K_\bullet$.
The outer products give its left-minus-right terms.
A comparison between projective resolutions lifting the identity is
a homotopy equivalence.
Applying $\operatorname{Hom}_{U\otimes U^{\mathrm{op}}}(-,U)$ gives
the stated formula and quasi-isomorphism.
\end{proof}

This is the classical Cartan--Eilenberg comparison.
For a finite Lie superalgebra, the same construction uses the
super-exterior powers and permutation Koszul signs.
Odd generators produce symmetric exterior directions, so the resolution
can be unbounded; its algebraic direct sums and cochain products remain
part of the convention.
The PBW and super-Koszul proof applies with those conventions.

A physical mixed holomorphic--topological realization requires a map
from its independently defined bulk observables to this cochain object.
The theorem specifies the adjoint coefficients and the actual algebraic
comparison.
An $E_2$ refinement, a current-factorization realization, and compatibility
with root-height transitions require maps preserving their additional
operations.
A quotient of Lie algebras alone does not provide a covariant map of
Hochschild centres.

\chapter{The exact scalar and first-order observables}
Set
\[
 \DenIg(Z)=D_5(2Z),
 \qquad \jmath(Z)=Z/2,
 \qquad \jmath^*\DenIg=D_5.
\]
The one-particle root spaces give
\[
 \sch\Fock(\gDelta_+)
 =\prod_{\alpha>0}(1-e^{-\alpha})^{-\sdim\gDelta_\alpha}.
\]
The oriented root operator gives
\[
 e^{-\rho}\operatorname{sPf}_F(\mathbb A_Z)=D_5(2Z).
\]
The primitive reduced curve-counting theorem gives
\[
 \mathcal Z^{\red}_{K3\times E}=-(\jmath^*\DenIg)^{-2}=-\chi_{10}^{-1}.
\]

\begin{theorem}[Protected square-root theorem]
Assume the stated root-space Fredholm comparison with its orientation,
Weyl shift, and convergence domain.
Within that realization, the denominator section $\DenIg$ is simultaneously:
\begin{enumerate}
\item the Weyl denominator of $\gDelta$;
\item the Euler character of the bar--Chevalley complex;
\item the oriented Fredholm Pfaffian of $\mathbb A_Z$;
\item the PBW determinant of the parity-refined root superspace $\mathcal H_{\mathrm{1p}}$;
\item the automorphic denominator whose half-coordinate pullback squares to the reduced scalar denominator.
\end{enumerate}
All five identifications preserve root degree and Weyl transport.
\end{theorem}
\begin{proof}
The first two are the denominator and bar--Chevalley theorems.  The third is the Fredholm theorem.  PBW for the Borcherds Lie superalgebra identifies the root superspace determinant, giving the fourth.  The half-coordinate identity $\jmath^*\DenIg=D_5$, the theta normalization, and the reduced curve-counting theorem give the fifth.
\end{proof}

\chapter{Physical meaning}
The root-datum realizations retain several distinct constructions.
The finite Ginzburg algebras and their projections form an inverse system.
A completed three-Calabi--Yau sector requires a continuous nondegenerate
Calabi--Yau structure and compatible cyclic models on the selected proper objects.
The quiver Hall category retains its chosen coefficients and extension
correspondences. A Hall theorem on that datum must identify its root presentation.
The Borcherds root superspace defines the root-graded Fock space and its
supercharacter. Its identification with physical one-particle states
requires a parity-preserving graded comparison intertwining the primitive actions.

A realization on the elliptic curve must construct the current object
with its local and wrapped supports, collision operations, and comparison maps.
The quantum Borcherds algebra, with its specified pairing and completion,
acts as a line-operator symmetry only through a representation preserving
fusion and exchange. The root operator gives the Fredholm product on its
specified convergence domain. A physical determinant interpretation requires
an operator comparison and an orientation-preserving identification of the
determinant lines and sections. The Fock character gives the scalar square.
An operator trace additionally requires its state space and trace domain.

Simple-root states satisfying the defining relations give a primitive Lie map.
The explicit enveloping and Chevalley extensions preserve products and brackets.
A compact string realization must additionally construct its brane functor,
Hall orientation maps, paired quantum coproduct, and trace-line comparison.

\chapter{Primitive comparisons and their algebraic extensions}
Let $\mathfrak n_\Delta$ be the positive Borcherds Lie superalgebra.
Choose homogeneous parity-preserving images of its simple generators
in a charge-graded Lie superalgebra $\mathfrak p$.
Assume these images satisfy every defining relation.

\begin{theorem}[Extension of a primitive comparison]\label{cyiv:primitive-comparison}
The chosen images determine a unique Lie-superalgebra map
$f:\mathfrak n_\Delta\to\mathfrak p$.
It induces
\[
 U(f)(x_1\cdots x_r)=f(x_1)\cdots f(x_r),
 \qquad C_*(f)=\operatorname{Sym}^c(f[1]).
\]
These are respectively an enveloping-algebra map and a Chevalley
chain-coalgebra map.
Filtration-preserving maps extend to the specified inverse-limit
completions.
If $f$ is surjective and both algebras have finite charge components
and proper positive height, equality of their parity-refined PBW
characters implies that $f$ and $U(f)$ are isomorphisms.
\end{theorem}
\begin{proof}
The defining presentation gives $f$.
Its tensor-algebra extension carries
$xy-(-1)^{|x||y|}yx-[x,y]$ to the corresponding relation.
It therefore descends to $U(f)$.
The shifted symmetric-coalgebra map commutes with the Chevalley
differential because $f$ preserves brackets.
Compatible maps on filtration quotients give maps on their inverse limits.
Surjectivity of $f$ gives surjectivity of $U(f)$.
Character equality makes each finite charge component an isomorphism,
separately in both parities.
The PBW embeddings then imply injectivity of $f$.
\end{proof}

At level zero the current Lie conformal algebra is
$\mathbb C[\partial]\otimes\mathfrak n_\Delta$ with
$[x_\lambda y]=[x,y]$ on constant generators.
The map $\partial^j x\mapsto\partial^j f(x)$ preserves this bracket
and sesquilinearity.
A nonzero central level additionally requires preservation of the
specified invariant pairing.
Hall convolution, Calabi--Yau structures, quantum coproducts, and
determinant lines require their own comparison data.

\chapter{Finite quivers and Calabi--Yau completions}
For a finite quiver $Q$ with zero potential, let $\Gamma(Q,0)$ be its
path-length-completed Ginzburg algebra.
Path products are read from left to right: if $a:i\to j$ and
$b:j\to k$, then $ab:i\to k$.
For each arrow $a:i\to j$, add $a^*:j\to i$ of degree $-1$.
At each vertex add a loop $t_i$ of degree $-2$.
Original arrows have degree zero, and
\[
 da=da^*=0,\qquad
 dt_i=e_i\left(\sum_a[a,a^*]\right)e_i.
\]

\begin{proposition}[Restriction to a full subquiver]\label{cyiv:ginzburg-projection}
If $Q_0$ is a full subquiver, retaining its generators and killing all
other vertices, arrows, reverse arrows, and loops defines a continuous
dg-algebra projection
\[
 \Gamma(Q,0)\longrightarrow\Gamma(Q_0,0).
\]
The inclusion of retained graded generators need not commute with
the differential.
\end{proposition}
\begin{proof}
The projection respects source and target idempotents and extends
multiplicatively.
At a retained vertex, every term involving a discarded arrow vanishes.
The other terms are precisely $dt_i$ in the smaller algebra.
All arrow differentials vanish, so the projection is a chain map.
It preserves the path-length filtration and is continuous.
For the final assertion, start with one vertex $i$ and no arrows,
where $dt_i=0$.
Add a vertex $j$ and an arrow $a:i\to j$.
In the larger algebra $dt_i=aa^*$, a nonzero length-two path.
Thus the graded inclusion is not a dg map.
\end{proof}

Ginzburg completion supplies a three-Calabi--Yau algebra for each finite
quiver in its stated setting.
The displayed projections specify an inverse system for compatible
zero-potential windows.
A Calabi--Yau structure on a completed limit requires its own
nondegeneracy and continuity argument.
Comparisons with Hall categories and quantized Borcherds algebras must
also preserve orientations, root relations, and the paired Hopf data.

\part{Quantum groups, exchange, and representation theory}

\chapter{From positive half to full quantum group}
Suppose the compact or local positive algebra $H^+$ is paired continuously with a negative half $H^-$, carries a completed coproduct, and admits Cartan data and a separated triangular completion. The completed double
\[
 \qgroup=\widehat D(H^+)
\]
has triangular decomposition $H^-\widehattensor H^0\widehattensor H^+$. Its universal $R$-matrix acts on continuous tensor products of modules. The classical limit is a Lie bialgebra whose cobracket is the first-order term of the coproduct asymmetry.

For $\C^3$, $H^+$ is the positive affine Yangian. For local ADE surfaces, it is the completed positive affine ADE Yangian. For the Igusa algebra, the paired Borel construction of the preceding part supplies a height-complete formal quantum group; a compact Hall realization maps to it through primitive, coproduct, and pairing comparisons.

\chapter{Tannaka reconstruction}
Let $\mathcal L$ be an essentially small rigid $\kfield$-linear abelian monoidal category with finite-dimensional morphism spaces, finite-length objects, and exact tensor product. Let
\[
 \omega:\mathcal L\to\Vect_{\mathrm{fd}}
\]
be a faithful exact strong monoidal functor. Assume the coend
\[
 H=\int^{M\in\mathcal L}\omega(M)^*\tensor\omega(M)
\]
is formed by the matrix-coefficient relations. Coevaluation gives its
coproduct, tensor product gives multiplication, and rigidity supplies
the antipode. A braiding gives a coquasitriangular form.
A meromorphic version additionally specifies its coefficient domains,
expansions and continuous completed tensors.

\begin{theorem}[Rigid Tannaka reconstruction]
Under the preceding hypotheses, the comparison functor is a monoidal equivalence
\[
 \mathcal L\simeq\operatorname{Comod}_{\mathrm{fd}}(H).
\]
Braiding transports to the universal coquasitriangular form on $H$.
\end{theorem}
\begin{proof}
The exact faithful coalgebra reconstruction theorem identifies
$\mathcal L$ with finite-dimensional comodules of this coend.
The explicit coaction, multiplication and antipode are those of
Proposition~\ref{prop:vol2-coend-operations}.
The monoidal comparison is proved in
Theorem~\ref{thm:vol2-tensor-reconstruction}.
Taking matrix coefficients of the hexagon identities gives the
coquasitriangular identities.
\end{proof}

This reconstruction explains why a line-operator category, a fiber functor, and exchange together produce a quantum group. The ordered boundary algebra provides ordered composition.  A compatible bialgebra coproduct supplies tensor products of its modules, and exchange supplies additional braiding data.

\chapter{KZ and rational exchange}
For a simple Lie algebra $\mathfrak g$, the Knizhnik--Zamolodchikov connection is
\[
 \nabla=d-\frac{\hbar}{2\pi i}
 \sum_{i<j}\Omega_{ij}\,d\log(z_i-z_j).
\]
Flatness follows from the infinitesimal braid relations. Drinfeld--Kohno identifies its monodromy tensor category with the corresponding Drinfeld--Jimbo category after a choice of associator.

Rational difference exchange instead gives the Yangian $R$-matrix
\[
 R(u)=1-\frac{\hbar P}{u}
\]
in the fundamental representation. The Yang--Baxter equation is the confluence identity for three exchanges. The KZ and Yangian constructions share the geometry of configuration spaces while using different flat connections.

\chapter{The centre and the double}
The Hochschild centre $Z_{\ord}(A)$ controls central operations on the
ordered algebra in its specified coefficient category.
To form a Drinfeld centre of modules, first specify a monoidal module category.
For an ordered chiral bialgebra satisfying
Definition~\ref{def:vol4-ordered-chiral-bialgebra}, use
\[
 \mathcal M_A=\operatorname{Mod}_A(\widehat{\mathcal V}_C)
\]
with the monoidal product of Theorem~\ref{thm:vol4-representation-fusion}.
An object of $Z(\mathcal M_A)$ is a module $M$ together with natural equivalences
\[
 c_{M,N}:M\widehat{\tensor}^{!}N
 \xrightarrow{\sim}N\widehat{\tensor}^{!}M
 \qquad(N\in\mathcal M_A)
\]
satisfying the half-braiding hexagon and unit identities, with their coherences.
This categorical centre is braided.
A comparison with the Hochschild centre must specify a realization functor,
the source and target of the comparison, and the operations it preserves.
Dualizability alone does not supply such a comparison.
Derived loop-space descriptions apply in the geometric categories and
under the hypotheses of their respective centre theorems
\cite{BenZviFrancisNadler}.

The centre and the double answer complementary questions.
The centre classifies central actions in the chosen category.
The double reconstructs a quantum group from paired positive and negative halves.
A field-theory comparison requires explicit realizations of the relevant
line and bulk operations and a map compatible with both constructions.

\part{A compact six-dimensional protected source}

\chapter{Holomorphic cotangent BF theory on \texorpdfstring{$K3\times E$}{K3 x E}}
Let $S$ be a K3 surface with holomorphic symplectic form $\sigma_S$, let $E$ be an elliptic curve with holomorphic form $\dd z$, and put
\[
 X=S\times E,
 \qquad
 \Omega_X=\sigma_S\wedge\dd z.
\]
For a finite proper-height quotient $\nIg^{(H)}$ of the positive Igusa Borcherds algebra, take fields
\[
 A\in\Omega^{0,\bullet}(X,\nIg^{(H)})[1],
 \qquad
 B\in\Omega^{0,\bullet}(X,(\nIg^{(H)})^\vee)[1].
\]
The invariant evaluation pairing and $\Omega_X$ define the BV action
\begin{equation}\label{eq:vol4-compact-hbf}
 S_H(A,B)=
 \int_X\Omega_X\wedge
 \angles{B,\dbar A+\frac12[A,A]}.
\end{equation}
This is the holomorphic cotangent theory of the dg Lie algebra of $\nIg^{(H)}$-valued Dolbeault forms.

\begin{theorem}[Classical compact-background source]
The action \eqref{eq:vol4-compact-hbf} satisfies the classical master equation and defines a holomorphic factorization algebra on the compact Calabi--Yau threefold $X$.
\end{theorem}
\begin{proof}
The shifted cotangent pairing is invariant.  The classical master equation is the Jacobi identity of $\nIg^{(H)}$ together with $\dbar^2=0$ and the derivation property of $\dbar$.  Local observables of an elliptic cyclic $L_\infty$ algebra form a factorization algebra by the standard BV construction.
\end{proof}

\chapter{Tree exactness and the quantum master equation}
Every interaction vertex in \eqref{eq:vol4-compact-hbf} has one $B$-leg and two $A$-legs.  The propagator pairs an $A$-leg with a $B$-leg.

\begin{lemma}[Loop bound]
A connected Feynman graph of the cotangent theory has loop number at most one.
\end{lemma}
\begin{proof}
A graph with $V$ interaction vertices has exactly $V$ internal or external $B$-half-edges.  Every internal edge consumes one of them, hence $E\le V$.  The loop number is $E-V+1\le1$.
\end{proof}

The positive algebra is graded by strictly positive root height.  For homogeneous $x\in\nIg^{(H)}$, the operator $\ad_x$ raises height.  Its coadjoint operator lowers dual height.  A product of adjoint operators labelled by positive roots has total nonzero height shift and therefore has zero trace.

\begin{proposition}[Vanishing of one-loop Lie weights]\label{cyiv:loop-weights}
For every finite positive-height quotient, the Lie weight of each
one-loop graph in the cotangent theory vanishes.
\end{proposition}
\begin{proof}
A one-loop graph has $E=V$.
Every $B$-leg is consumed by an internal edge, so its external legs are
$A$-legs.
Its Lie weight is a cyclic trace of adjoint operators, or the dual
trace with the conventionally equivalent signs.
Each external label has positive root height.
The product therefore strictly raises height on the finite-dimensional
adjoint representation and has zero trace.
The same argument applies to the dual representation.
\end{proof}

This vanishing removes the one-loop Lie coefficients in a perturbative
construction with the stated propagator and vertex types.
A renormalized quantum factorization algebra still requires admissible
analytic kernels, local extensions, the quantum master equation, and
compatible height transitions on the chosen observable spaces.
These are the remaining quantum realization data.

\chapter{Classical reduction along K3}
Choose a K\"ahler metric on $S$ and normalize a closed harmonic
$(0,2)$-form $\eta$ by $\int_S\sigma_S\wedge\eta=1$.
Let $K=\mathbb C[\eta]/(\eta^2)$ with $|\eta|=2$ and zero differential.
There is a unital dg-algebra map
\[
 i:K\longrightarrow\Omega^{0,*}(S),\qquad \eta\longmapsto\eta.
\]
Its image is closed under multiplication because $(0,4)$-forms vanish.
The Dolbeault cohomology of a K3 surface is one-dimensional in degrees
zero and two and zero in degree one, so $i$ is a quasi-isomorphism.
It preserves the trace defined by integration against $\sigma_S$.

Let $\mathfrak l$ be a finite-dimensional Lie algebra, and write
$\mathfrak d=\mathfrak l\ltimes\mathfrak l^*$ for its cotangent Lie
algebra, with invariant evaluation pairing.
For an open $U\subset E$, tensor with $\Omega^{0,*}(U)$, using the
completed smooth tensor product in the $S$ direction.
The field comparison is
\[
 K\otimes\Omega^{0,*}(U)\otimes\mathfrak d
 \xrightarrow{i\otimes\id}
 \Omega^{0,*}(S\times U)\otimes\mathfrak d.
\]
The bracket is multiplication of forms followed by the Lie bracket,
with Koszul signs.

\begin{proposition}[Classical field comparison]\label{cyiv:k3-classical-reduction}
The displayed map is a trace-preserving dg Lie quasi-isomorphism,
compatible with restriction in $U$.
Its induced minimal operations in the K3 factor have no higher products.
\end{proposition}
\begin{proof}
Multiplicativity of $i$ proves bracket compatibility.
Hodge projection $p$ and $h=\bar\partial^*G$ on the compact surface satisfy
$pi=\id$ and $\bar\partial h+h\bar\partial=\id-ip$.
Tensoring these continuous operators with forms on $U$ gives the same
contraction, with the tensor-product signs.
Thus $i\otimes\id$ is a quasi-isomorphism.
Restriction acts only on the $U$ factor and commutes with it.
The normalized trace on $K$ is the restriction of the Dolbeault trace.
Choose the contraction with $hi=0$.
In every higher transferred tree, the first product of inputs from
$i(K)$ remains in $i(K)$ and the next internal homotopy kills it.
Thus every higher transferred product vanishes.
\end{proof}

The coefficient algebra $K$ retains both the degree-zero and degree-two
K3 modes.
The proposition compares classical field complexes and their brackets.
Constructing an equivalence of quantum observable factorization algebras
requires compatible quantization, support conventions, and the induced
maps on observable completions.
The ordinary Hochschild comparison of Theorem~\ref{cyiv:hh-ce} is a
separate algebraic calculation, not that quantum field map.

\chapter{Protected states and the Igusa trace}
Let
\[
 \mathcal H_{\mathrm{1p}}
 =\widehat\bigoplus_{\alpha>0}\gIg_\alpha
\]
with its Borcherds parity.  The multiparticle state space is the completed super-Fock space
\[
 \mathcal H_{\mathrm{Fock}}=\widehat\Sym(\mathcal H_{\mathrm{1p},\bar0})
 \tensor\widehat\Lambda(\mathcal H_{\mathrm{1p},\bar1}).
\]
For $Z$ in the absolute-convergence chamber, let $e^{-H_Z}$ act by $e^{-\alpha(Z)}$ on the root space of degree $\alpha$.

\begin{theorem}[Protected Hamiltonian trace]
The supertrace is
\[
 \str_{\mathcal H_{\mathrm{Fock}}}(e^{-H_Z})
 =\prod_{\alpha>0}
 (1-e^{-\alpha(Z)})^{-\sdim\gIg_\alpha}.
\]
After multiplication by the inverse Weyl line, its reciprocal is $\DenIg(Z)=D_5(2Z)$.  After the half-coordinate pullback and doubling, the disconnected protected scalar is $-(\jmath^*\DenIg)^{-2}=-\chi_{10}^{-1}$ in the Oberdieck--Pixton normalization.
\end{theorem}
\begin{proof}
An even oscillator contributes $(1-e^{-\alpha})^{-1}$ and an odd oscillator contributes $1-e^{-\alpha}$.  Multiplication over the root basis gives the product.  The denominator identity and the reduced curve-counting theorem supply the final two equalities.
\end{proof}

\chapter{The compact-source comparison}
The compact threefold provides the classical cotangent field theory
and the dg Lie comparison of Proposition~\ref{cyiv:k3-classical-reduction}.
The root datum separately gives the Borcherds algebra, its enveloping
algebra, and the root-graded Fock space.
The latter has the trace computed above.
Proposition~\ref{cyiv:loop-weights} supplies the finite one-loop Lie
cancellation, while Theorem~\ref{cyiv:hh-ce} computes the ordinary
algebraic centre model with adjoint coefficients.

\begin{problem}[Geometric compact comparison]\label{cyiv:compact-comparison}
Construct objects $E_i$ representing the simple branes and maps on their
derived morphism complexes that preserve units, composition, and higher
composition homotopies.
Realize these maps on extension stacks with the reduced critical
coefficients and orientation isomorphisms on triple flags.
Construct compatible tensor and unit constraints on the resulting
boundary functor.
Compare the chosen quantum bulk action with the independently defined
boundary centre, including support, collisions, and completion.
Finally construct a trace and a map of its value line to the automorphic
line in the same charge normalization.
\end{problem}

A primitive map satisfying the root relations gives the enveloping and
Chevalley maps of Theorem~\ref{cyiv:primitive-comparison}.
A strong monoidal equivalence transports categorical centres.
Neither construction determines the brane functor or its cyclic trace.
These are the additional maps required for compact geometric
reconstruction.

\part[Physics]{Physics of the comparison maps}

\chapter{Holomorphic--topological boundary theories}
A three-dimensional holomorphic--topological theory on $C\times\R$ has local observables that are holomorphic along $C$ and locally constant along $\R$. A boundary produces an ordered chiral algebra. A second topological direction permits exchange and gives an $\Etwo$ enhancement or a meromorphic braided line category. The bulk acts through the derived centre.

The finite cotangent BF model gives the exact local calculation:
\[
 \Barc(U(\mathfrak l))\simeq C_*(\mathfrak l).
\]
For $\mathfrak l=\nDelta^{(H)}$, this calculation realizes every finite-height Igusa denominator sector.

\chapter{Six-dimensional hCS on \texorpdfstring{$K3\times E$}{K3 x E}}
Choose a holomorphic volume form $\Omega_S\wedge dz_E$. The hCS theory has classical fields
\[
 \Omega^{0,*}(S\times E,\mathfrak a)[1]
\]
and action $S_{\hCS}$. Its factorization observables are native to the threefold. Pushforward along $p_E$ gives a factorization algebra on $E$ whose coefficient complex contains the Dolbeault modes of $S$.

A protected reduction to a chiral algebra consists of a deformation retract of the K3 Dolbeault complex, compatibility of the transferred $L_\infty$ operations with the trace, a boundary or defect selecting ordered fusion, and a quantization trivializing the local anomaly. These data produce the ordered chiral algebra to which the Hall/BPS comparison may be applied.

The Hodge equality $\chi(\Ocal_{S\times E})=0$ controls one compact supertrace. The gauge and gravitational anomalies are classes in the local BV deformation complex and are evaluated by their own characteristic forms.

\chapter{Omega backgrounds}
Let a torus $T$ act holomorphically on the local geometry and preserve the selected supercharge. For a generator $V$ of the action, the equivariant differential is
\[
 Q_\Omega=Q+\iota_V,
 \qquad Q_\Omega^2=\mathcal L_V.
\]
After localization, fixed loci produce deformation parameters and often a spectral $R$-matrix. Toric threefolds and ADE surface resolutions carry such actions. A deformation family equipped with monodromy-compatible moduli correspondences and wall-crossing equivalences transports the localized algebra from algebraic or equivariant K3 fibers to the selected generic fiber.

\chapter{Anomalies as typed classes}
The principal anomaly classes are
\[
\begin{array}{c|c}
\text{construction}&\text{deformation complex}\\ \hline
\text{hCS quantization}&\text{local BV cochains}\\
\text{critical CoHA orientation}&\text{Picard/gerbe theory of the critical stack}\\
\text{Hall coproduct}&\text{correspondence and localization complex}\\
\text{projective conformal blocks}&\text{Atiyah algebra of the determinant line}\\
\text{automorphic multiplier}&\text{character group of the modular cover}.
\end{array}
\]
A comparison theorem is a morphism between the corresponding complexes together with equality of the transported classes. This formulation turns numerical coincidences into geometric identities precisely when an intertwining map exists.

\part[Automorphic forms]{Automorphic characters and scalar shadows}

\chapter{Borcherds products as characters}
Let $\mathfrak n_+$ be a locally finite positive Lie superalgebra. The Chevalley Euler character is
\[
 \chi C_*(\mathfrak n_+)
 =\prod_{\alpha>0}(1-e^{-\alpha})^{\sdim\mathfrak g_\alpha}.
\]
A Borcherds product becomes the Weyl-shifted character of $C_*(\mathfrak n_+)$. This construction produces a character of an existing algebra. Hall geometry supplies a second algebra, and a primitive comparison theorem identifies the two.

\chapter{The universal Borcherds weight formula}
For a weight-zero weak Jacobi form
\[
 \varphi(\tau,z)=\sum c(n,\ell)q^nr^\ell,
\]
the Borcherds lift has weight
\[
 \wt(\operatorname{Borch}(\varphi))=\frac{c(0,0)}2.
\]
For the K3 half-index, $c(0,0)=10$ and the weight is $5$. This formula lives entirely on the automorphic layer. Hodge Euler characteristics, Heisenberg ranks, central charges, and BV anomalies are evaluated by their own traces.

\chapter{The Gritsenko--Cl\'ery atlas}
The eight diagonal-divisor forms are
\[
 \Delta_5,\Delta_2,\Delta_1,\Delta_{1/2},
 \nabla_3,\nabla_2,\nabla_{3/2},Q_1
\]
with weights
\[
 5,2,1,\frac12,3,2,\frac32,1.
\]
Each form is an automorphic product on its own paramodular or Hecke congruence group with its own character or multiplier \cite{GC}. A geometric realization is specified row by row by a charge lattice, Hall problem, trace, and primitive algebra comparison.

\chapter{Scalarization hierarchy}
The constructions descend through four levels:
\[
\begin{aligned}
 \text{factorization or Hall object}
 &\longrightarrow \text{primitive Lie algebra}\\
 &\longrightarrow \text{PBW character}
 \longrightarrow \text{scalar partition function}.
\end{aligned}
\]
Each arrow retains the quotient of structure displayed at its target. The first object contains supports, extension maps, and brackets. The second retains the infinitesimal algebra. The third retains signed graded dimensions. The fourth adds a trace and normalization. Upward reconstruction is supplied by the corresponding geometric operations and comparison morphisms.

\part{The unified theorem system}

\chapter{The foundational theorem}
\begin{definition}[Complete Calabi--Yau quantum-group datum]
A complete comparison datum specifies independent categorical, field,
boundary, Hall, Hopf, and trace objects together with:
\begin{enumerate}
\item PTVV-representable moduli for the smooth proper Calabi--Yau category;
\item a renormalized BV quantization satisfying the quantum master equation and a local compact factorization generator $b$;
\item a boundary reduction or proper locally trivial map to a curve satisfying factorization pushforward, together with the supplied curve-level faces and aligned distributive law;
\item a strong-oriented Hall coefficient theory whose pull--push maps satisfy base change and whose positive algebra is separated and complete;
\item paired negative and Cartan halves, a completed coproduct, a continuous nondegenerate Hopf pairing, and a separated triangular completion on which the canonical tensors converge;
\item compatible open--closed and Hall--boundary morphisms;
\item a parity-refined primitive comparison for every automorphic identification.
\end{enumerate}
\end{definition}

\begin{theorem}[Calabi--Yau quantum-group assembly]
Let a complete datum consist of:
\begin{enumerate}
\item a Calabi--Yau category and a representable moduli-of-objects problem;
\item a quantized cyclic elliptic BV theory and a boundary or defect generator;
\item a Hall extension correspondence with coefficient functor, base-change theorem, strong orientation data, and separated charge completion;
\item paired positive and negative bialgebras, Cartan data, a nondegenerate continuous Hopf pairing, and a completed double;
\item every curve-level distributive, PBW, and analytic comparison used below.
\end{enumerate}
Then the datum determines functorially:
\begin{enumerate}
\item framed $\Etwo$ Hochschild cochains and the shifted-symplectic moduli structure;
\item the three curve-level algebraic faces and their three bars;
\item the positive Hall/BPS algebra and its primitive Lie algebra;
\item the completed quantum double and braided representation category;
\item every supplied comparison among categorical, field-theoretic, Hall, current, quantum, and automorphic realizations;
\item the Weyl-shifted PBW/Chevalley character of any parity-refined primitive Lie-superalgebra identification.
\end{enumerate}
\end{theorem}
\begin{proof}
Cyclic Deligne and shifted symplectic geometry give the first output \cite{BravRozenblyum,PTVV,BravDyckerhoff}.  Boundary factorization gives the curve faces, and Volume~I gives their bars and centres.  The Hall correspondence and its base-change theorem give the associative product; strong orientation supplies the coefficient coherence required in the $3$-Calabi--Yau setting \cite{KPS}.  The paired bialgebra data give the double and universal $R$-matrix.  The comparison maps commute by hypothesis.  PBW and Chevalley antisymmetrization give the final character.  Each construction is natural under a morphism preserving the displayed data.
\end{proof}

The local quiver, equivariant surface, K3 correspondence, and Igusa Borcherds constructions have the carriers stated above.
A geometric $K3\times E$ comparison requires the primitive map and determinant-line map of Problem~\ref{cyiv:compact-comparison}.

\chapter{Landscape of established realizations}
\begin{center}
\footnotesize
\begin{tabularx}{0.99\textwidth}{@{}>{\raggedright\arraybackslash}p{0.15\textwidth}>{\raggedright\arraybackslash}p{0.25\textwidth}>{\raggedright\arraybackslash}X>{\raggedright\arraybackslash}p{0.20\textwidth}@{}}
\toprule
\textbf{Geometry}&\textbf{Positive algebra}&\textbf{Full or centred object}&\textbf{Character}\\
\midrule
$A_2$ quiver&$U_v^+(\mathfrak{sl}_3)$&$U_v(\mathfrak{sl}_3)$&PBW root product\\
$\C^3$&$Y^+(\widehat{\mathfrak{gl}}_1)$&$Y(\widehat{\mathfrak{gl}}_1)$&plane-partition and Fock traces\\
$\widetilde{\C^2/\Gamma}$&$Y^+(\widehat{\mathfrak g}_\Gamma)$&paired completion when the negative half and coproduct are supplied&surface Fock traces\\
$K3$&$U(\mathfrak g_{\BPS}^+)$&Heisenberg and LLV actions&K3 Hall characters\\
$K3\times E$&Borcherds root envelope; geometric reduced character&$U_q(\gDelta)$ and the root-current model and classical BF theory&$-\ChiTen^{-1}$\\
\bottomrule
\end{tabularx}
\end{center}
The $A_2$ and $\C^3$ rows provide complete standard quantum-group models. The ADE row provides an established completed positive half together with its geometric current actions; its paired full double is formed from the negative half, Cartan data, coproduct, continuous Hopf pairing, and triangular completion stated in the table. The K3 row provides global two-Calabi--Yau BPS and cohomological symmetries. The final row records independently constructed root, character, and classical field objects.
Their geometric Hall realization requires the comparison in Problem~\ref{cyiv:compact-comparison}.

\chapter{Constructed comparisons and their geometric extensions}
The Hall flag argument constructs associativity once its coefficient
and orientation maps are defined.
The tensor-action formula constructs representation fusion from a
coassociative counital coproduct.
The matrix-coefficient formulas construct the coend of a supplied
fibre functor.
The primitive presentation gives enveloping and Chevalley maps.
The classical cyclic and K3 field calculations retain their actual
differentials, traces, and brackets.

The quantum field realization requires admissible quantization and
observable completions.
The compact Hall realization requires the brane and orientation maps
of Problem~\ref{cyiv:compact-comparison}.
The Fredholm and automorphic comparisons require their analytic domains
and value-line isomorphisms.
These constructions preserve the full geometric reconstruction problem
beyond its finite algebraic and classical calculations.

\chapter{Actions and comparison maps}
For a specified boundary algebra $A_b$, a coherent bulk action is a map
\[
 \beta:\Obs_{\mathrm{bulk}}|_\partial\longrightarrow Z_{\ord}(A_b).
\]
The target is the derived endomorphism object of the diagonal bimodule
in the chosen factorization coefficient category.
A Hall construction instead starts with
\[
 \M_\alpha\times\M_\beta\xleftarrow{p}\mathcal E_{\alpha,\beta}
 \xrightarrow{q}\M_{\alpha+\beta}.
\]
Critical coefficients, orientations, and admissible pull--push maps
give its product.
Framed moduli supply a representation after their action
correspondences and triple-flag identities have been constructed.
A paired Hopf structure gives $D(H^+,H^0,H^-)$.
Its coproduct defines tensor products of continuous modules, and a
convergent universal $R$-matrix gives exchange.

A strong monoidal equivalence between this module category and a
boundary category transports categorical centres.
It does not identify that categorical centre with the object
$Z_{\ord}(A_b)$ without a realization theorem.
A trace on a specified representation then gives an amplitude.
Its comparison with an automorphic section requires a map of value
lines with the same normalization.

\chapter{Comparison data for Calabi--Yau constructions}
Comparisons among Calabi--Yau categories, holomorphic field theory, chiral algebras, Hall algebras, quantum groups, and automorphic products require the following structures:
\begin{enumerate}
\item for a smooth proper $d$-Calabi--Yau category with its nondegenerate negative-cyclic class, the categorical outputs are framed $\Etwo$ Hochschild cochains and, when the moduli stack of objects is representable in the PTVV setting, $(2-d)$-shifted symplectic moduli; the shift $d$ does not change the Hochschild operadic level to $\En$;
\item a geometric BV theory and boundary category produce factorization observables and transverse ordered chiral algebras; associative chiral convolution is recorded as a separate supplied structure;
\item pushforward to a curve uses an actual map $p:X\to C$; formal cycle-specialization notation alone specifies no pushforward;
\item extension stacks, coefficient sheaves, orientations, and completions define the Hall algebra;
\item coproduct, Cartan, Hopf pairing, and topology define the double;
\item K3 Heisenberg, LLV orthogonal symmetry, local ADE Yangians, and global BPS generalized Kac--Moody algebras occupy distinct strata joined by explicit comparison maps where constructed;
\item the graph of the $K3\times E$ Gram map becomes multiplicative in the polarized central extension;
\item point-local and wrapped sectors form a hybrid factorization problem;
\item the Igusa form is the Weyl-shifted PBW/Chevalley character of $\gDelta$, and compact Hall geometry maps to it through the primitive comparison;
\item anomaly identities are transported through explicit maps between deformation complexes.
\end{enumerate}
The categorical, geometric, Hall, and automorphic operations thus have separate carriers and comparison maps.

\part{Technical appendices to the preceding book}
\chapter{Typed data table}
\begin{center}
\small
\begin{tabularx}{0.98\textwidth}{@{}>{\raggedright\arraybackslash}p{0.20\textwidth}>{\raggedright\arraybackslash}X>{\raggedright\arraybackslash}p{0.28\textwidth}@{}}
\toprule
\textbf{Object}&\textbf{Minimal data}&\textbf{Primary operation}\\
\midrule
$\CY$ category&nondegenerate negative-cyclic trace&Hochschild and shifted-symplectic structures\\
Factorization algebra&local observables and Weiss descent&disjoint-open product\\
Transversely ordered chiral algebra&factorization coefficient with an internal $\Eone$ structure for $\tensor^!$&ordered fusion\\
Associative chiral algebra&an $\Eone$ structure for chiral convolution $\starCh$&nonlocal chiral composition\\
$\CoHA$&extension correspondence and coefficient theory&pull--push product\\
Bialgebra&compatible product and coassociative counital coproduct&fusion of representations\\
Double&positive and negative halves, Cartan data, Hopf pairing&universal $R$-matrix\\
Borcherds product&Jacobi or vector-valued modular input&graded character\\
\bottomrule
\end{tabularx}
\end{center}

\chapter{The three K3 brackets}
The Heisenberg bracket is the excess intersection of two Nakajima incidence kernels.  The LLV bracket is the commutator of diagonal Lefschetz kernels and their adjoints.  The local Yangian bracket is the convolution commutator of ADE Hecke kernels with spectral equivariance.  In $\operatorname{Corr}_{K3}(S)$ all three are compositions of geometric kernels.  Their mixed commutators are the corresponding triple-intersection classes.  The action algebra and Tannaka coend retain these classes while preserving the three geometric constructions.

\chapter{Exact arithmetic checks}
The companion verification calculation establishes:
\begin{enumerate}
\item the low Fourier rows of $\phi_{0,1}$ through $q^3$ by exact theta-series division;
\item the vector-valued cocycle identity for the Gram polarization;
\item the rational Yang--Baxter equation on $(\C^2)^{\tensor3}$;
\item the $A_2$ Hall flag counts at ranks zero and one;
\item the type-II Gram matrix and reflection identities.
\end{enumerate}
Each calculation uses integer or rational arithmetic and appears as a worked formula in the relevant chapter.

\cleardoublepage
\phantomsection
\addcontentsline{toc}{part}{Universal Chiral BV and Einstein Completion}
\chapter*{Coherent boundary actions and BV fields}
\addcontentsline{toc}{chapter}{Coherent boundary actions and BV fields}
\begin{summarybox}
Boundary operations determine a deformation problem once their interface and coherence maps are supplied. A local realization gives a shifted cotangent BV theory. Comparing it with a physical bulk requires a map of fields, actions, and coherent boundary operations.
\end{summarybox}
\part{The full mixed boundary datum}

\chapter{Three products and one interface}
Holomorphic collisions and transverse composition act on different coefficient objects. A common boundary theory must specify their interaction on an interface. This interaction is part of the input to the following constructions.

Let $X$ be a smooth complex curve and let $\Dcal_X=D\text{-}\Mod(\Ran X)$ in a fixed six-functor enhancement.  Write $\tensor^!$ for the pointwise product and $\starCh$ for chiral convolution.

\begin{definition}[Complete mixed boundary datum]
A complete mixed boundary datum on $X$ is a tuple
\[
 \mathbb A=(L,A^{\ch},A^\perp,K,\iota,\lambda)
\]
with the following components.
\begin{enumerate}
\item $L$ is a complete chiral Lie algebra.
\item $A^{\ch}$ is a complete associative algebra for $\starCh$.
\item $A^\perp$ is a complete $\Eone$ algebra in $(\Fact_D(X),\tensor^!)$.
\item $K$ is an $(A^{\ch},A^\perp)$ interface bimodule.
\item $\iota:L\to(A^{\ch})_{\Lie}$ is a chiral Lie morphism.
\item $\lambda$ is the aligned mixed distributive law relating the two module actions on $K$.
\end{enumerate}
Every structure map is continuous for the chosen weight filtrations. The tuple includes the module actions, the mixed coherence homotopies, and their identities in the chosen enhancement. When a reduced bar is used, the relevant augmentations are supplied as additional data. Existence of this tuple for a physical boundary theory is a separate construction.
\end{definition}

The three native bars are
\[
 \BarLieCh(L),\qquad
 \BarAssCh(A^{\ch}),\qquad
 \BarPerp(A^\perp).
\]
The interface has a fourth bar.  Its degree-$r$ terms are alternating words in the two associative colours with one marked interface letter.  Its differential contracts adjacent compatible letters.

\begin{construction}[Interface bar]
Let $\overline A^{\ch}$ and $\overline A^\perp$ be the augmentation ideals.  Set
\[
 \BarMix(\mathbb A;K)
 =\bigoplus_{p,q\ge0}
 (s\overline A^{\ch})^{\tensor p}\tensor sK\tensor
 (s\overline A^\perp)^{\tensor q}.
\]
The differential is the sum of the internal differential, the two bar differentials, the left action on $K$, the right action on $K$, and the mixed terms supplied by $\lambda$.
\end{construction}

\begin{theorem}[Four-colour differential]
The interface differential squares to zero.  Its codimension-two terms are exactly the associative relations in the two colours, the two module relations, and the mixed distributive relation.
\end{theorem}
\begin{proof}
Filter by total word length.  Every component of the square contracts two adjacent edges.  Disjoint contractions cancel by the determinant-line sign.  Nested contractions in one colour cancel by associativity.  A contraction meeting the interface cancels by the corresponding module axiom.  The remaining square has one edge of each colour and is the aligned distributive square.  The datum $\lambda$ identifies its two composites.
\end{proof}

\chapter{Compatible modules and the mixed centre}
An operation on a boundary theory must act on boundary-changing morphisms as well as on objects. Natural transformations retain this compatibility, and derived natural transformations retain its higher homotopies.

A compatible module is a triple $(M^{\ch},M^\perp,\kappa)$ with an $A^{\ch}$ action, an $A^\perp$ action, and an interface comparison over $K$.  Let $\Mod_{\mathbb A}^{\mathrm{mix}}$ be their stable category.

\begin{definition}[Mixed derived centre]
The mixed derived centre is
\[
 Z_{\mathrm{mix}}(\mathbb A)
 =\End\!\left(\id_{\Mod_{\mathbb A}^{\mathrm{mix}}}\right).
\]
It is computed by the endomorphisms of the diagonal compatible bimodule.
\end{definition}

\begin{theorem}[Universal interior action]
$Z_{\mathrm{mix}}(\mathbb A)$ carries an internal $\Etwo$ algebra structure.  Every action already defined as a coherent natural endomorphism of the identity determines a point of this centre. The space of its factorizations through the representing endomorphism object is contractible.  Morita-equivalent complete mixed data have equivalent centres.
\end{theorem}
\begin{proof}
Endomorphisms of the identity functor form the higher centre of a module category.  Composition and insertion of natural transformations give the brace operations and hence the $\Etwo$ structure.  A supplied coherent natural action defines an endomorphism of the identity. A physical bulk insertion determines such an action only after its action on modules, morphisms, and higher compositions has been constructed.  The factorization is the universal property of the endomorphism object.  Morita equivalence identifies the module categories and their identity functors.
\end{proof}

\chapter{Relative induction from a chiral algebra}
Let $V$ be a chiral algebra with conformal vacuum.  Its algebraic collar $\Coll(V)$ is the free transverse $\Eone$ extension generated by one boundary copy of $V$, modulo the chiral relations already present in $V$.  Weight completion gives $\widehat\Coll(V)$.

\begin{theorem}[Universal collar]
For every complete transversely ordered chiral algebra $B$ and every chiral morphism $V\to B$, there is a unique continuous $\Eone$ morphism
\[
 \widehat\Coll(V)\longrightarrow B
\]
extending the chiral map.  The construction is left adjoint to restriction from transverse ordered algebras to chiral coefficients.
\end{theorem}
\begin{proof}
The algebraic collar is the relative operadic induction from the chiral operad to the mixed operad.  Relative induction is left adjoint to restriction.  Completion preserves the adjunction on positive complete objects because every map is determined weight by weight.
\end{proof}

When relative induction supplies a regular interface and a continuous aligned distributive law, these data define a mixed envelope with $L=V_{\Lie}$, $A^{\ch}=U^{\ch}(L)$, and $A^\perp=\widehat\Coll(V)$. Constructing that interface law in the selected geometric coefficient category is part of the envelope construction.

\part{The deformation algebra and its BV envelope}


\chapter{The centre-extended deformation algebra}
Varying a boundary product changes its compatibility with the interface and the central action. The deformation complex must therefore include all these operations. A local field realization is further data needed to turn their deformation equations into a BV action.

Fix a cofibrant coloured operad $\mathcal P_{\mathrm{mix}}$ presenting the structures of a complete mixed datum and its central action.  Let $\mathcal P_{\mathrm{mix}}^{\ash}$ be its bar cooperad.  The convolution complex
\[
 \Def_{\mathbb A}^{\mathrm{full}}
 =\Hom_{\mathbb S}
 \left(\mathcal P_{\mathrm{mix}}^{\ash},\End_{\mathbb A}\right)
\]
contains the chiral operation, both associative products, the interface, and the centre action.

\begin{theorem}[Complete deformation $L_\infty$ algebra]
The convolution complex has a canonical complete $L_\infty$ structure.  Its Maurer--Cartan points are complete mixed boundary structures together with a coherent central action.  Twisting by the point defined by $\mathbb A$ gives the tangent algebra of deformations of the full datum.
\end{theorem}
\begin{proof}
The infinitesimal decomposition maps of the bar cooperad insert collections of operations into one another.  Composition in the endomorphism operad turns these decompositions into the convolution brackets.  The cooperad identities imply the $L_\infty$ identities.  A Maurer--Cartan element is precisely an operad morphism from a cofibrant resolution of $\mathcal P_{\mathrm{mix}}$ to $\End_{\mathbb A}$.  Twisting linearizes this mapping space at the selected structure.
\end{proof}

\section{Field equations and Bianchi identities}
Let $\alpha$ be a degree-one field in the localized deformation algebra.  Its curvature is
\[
 \mathcal F(\alpha)
 =\sum_{n\ge1}\frac1{n!}\ell_n(\alpha^{\tensor n}).
\]
The linearization at $\alpha$ is
\[
 d_\alpha x
 =\sum_{n\ge0}\frac1{n!}
   \ell_{n+1}(\alpha^{\tensor n},x).
\]

\begin{theorem}[Universal Bianchi identity]
For every $\alpha$,
\[
 d_\alpha\mathcal F(\alpha)=0.
\]
At a Maurer--Cartan point, $d_\alpha^2=0$.
\end{theorem}
\begin{proof}
Both identities are the exponential form of the $L_\infty$ Jacobi relations.  Expand the coderivation square on the cofree cocommutative coalgebra and evaluate it on $e^\alpha$ and $e^\alpha x$.
\end{proof}

\chapter{Shifted cotangent transgression}
Let $\flie$ be a complete elliptic $L_\infty$ algebra on a three-manifold $M$.  Its shifted cotangent theory has fields
\[
 \mathcal E=T^*[-1]\flie[1]
 =\flie[1]\oplus\flie^\vee[-2].
\]
The canonical action is
\[
 S(\alpha,\beta)
 =\angles{\beta,\mathcal F(\alpha)}.
\]

\begin{theorem}[Universal BV envelope]
The action satisfies the classical master equation.  Its Euler--Lagrange equations are
\[
 \mathcal F(\alpha)=0,
 \qquad
 d_\alpha^\vee\beta=0.
\]
The tangent complex at a solution is the cotangent complex of the Maurer--Cartan moduli problem.
\end{theorem}
\begin{proof}
The Hamiltonian vector field of $S$ is the cotangent lift of the homological vector field defining the $L_\infty$ structure.  The square of that vector field is zero, hence $\{S,S\}=0$.  Variation in $\beta$ gives the Maurer--Cartan equation.  Variation in $\alpha$ gives the dual linearized equation.
\end{proof}

\begin{definition}[Local realization and chiral BV envelope]
A local realization of the complete mixed datum consists of an elliptic local $L_\infty$ algebra $\flie_{\mathbb A}^{\mathrm{full}}$ on $M=X^{\mathrm{an}}\times\R$, a specified continuous filtered $L_\infty$ quasi-isomorphism
\[
 \Gamma(M,\flie_{\mathbb A}^{\mathrm{full}})
 \longrightarrow \Def_{\mathbb A}^{\mathrm{full}},
\]
and the local pairing needed for cotangent transgression. Here $\Gamma$ denotes the chosen complex of global sections with its declared topology. The comparison includes compatibility with the specified boundary operations. The dual fields use the density-valued continuous dual, with compact support or boundary conditions that make the integrated pairing and integration by parts defined.

For a supplied local realization, define
\[
 \BVEnv_3^{\mathrm{full}}(\mathbb A)
 =T^*[-1]\MC(\flie_{\mathbb A}^{\mathrm{full}}).
\]
Its action is the cotangent action displayed above. Constructing a local realization from an abstract mixed datum remains an additional existence problem.
\end{definition}

\chapter{Intrinsic algebraic quantization}
Let $\mathfrak X$ be a finite perfect derived formal moduli problem with a homological orientation $\mu$: the corresponding Berezinian is invariant under the homological vector field.  Functions on its shifted cotangent are completed polyvectors:
\[
 \Ocal(T^*[-1]\mathfrak X)
 =\widehat{\PV}(\mathfrak X)
 =\prod_{p\ge0}\Gamma(\mathfrak X,\wedge^pT_{\mathfrak X})[p].
\]
The homological orientation defines the divergence operator
\[
 \Delta_\mu:\widehat{\PV}(\mathfrak X)
 \longrightarrow\widehat{\PV}(\mathfrak X)[1].
\]

\begin{theorem}[Oriented BV complex]
The operator $\Delta_\mu$ has order two and satisfies
\[
 \Delta_\mu^2=0,
 \qquad
 [d_{\mathfrak X},\Delta_\mu]=0,
\]
\[
 \Delta_\mu(fg)=\Delta_\mu(f)g+(-1)^{|f|}f\Delta_\mu(g)
 +(-1)^{|f|}\{f,g\}.
\]
Consequently
\[
 \QHdg(\mathfrak X,\mu)
 =\left(\widehat\PV(\mathfrak X)\llbracket\hbar\rrbracket,
 d_{\mathfrak X}+\hbar\Delta_\mu\right)
\]
is a square-zero quantum BV complex.
\end{theorem}
\begin{proof}
In oriented Darboux coordinates, $\Delta_\mu$ is the sum of paired second derivatives.  Their graded commutation gives $\Delta_\mu^2=0$.  Invariance of the Berezinian under the homological vector field gives $[d_{\mathfrak X},\Delta_\mu]=0$.  Expanding the operator on a product gives the bracket as the cross term.
\end{proof}

The orientation identifies polyvectors with differential forms by contraction.  Under this identification,
\[
 \iota_\mu\Delta_\mu=d_{\mathrm{dR}}\iota_\mu.
\]
Hence the BV complex is the Hodge-completed derived de Rham complex with differential
\[
 d_{\mathfrak X}+\hbar d_{\mathrm{dR}}.
\]
The Rees parameter records Hodge degree.

\begin{definition}[Intrinsic quantum complex of a mixed datum]
For a complete mixed datum $\mathbb A$ equipped with compatible homological orientations on its finite perfect charts, define $\QHdg_3^{\mathrm{full}}(\mathbb A)$ by applying the preceding BV complex to finite perfect charts of
\[
 \MC(\flie_{\mathbb A}^{\mathrm{full}})
\]
and taking the separated inverse limit over the chart and weight filtrations.
\end{definition}

\begin{theorem}[Quantum master equation]
A degree-zero action $S=S_0+\hbar S_1+\cdots$ defines a quantum state $e^{S/\hbar}\mu$ precisely when
\[
 dS+\frac12\{S,S\}+\hbar\Delta_\mu S=0.
\]
The equation is equivalent to
\[
 (d+\hbar\Delta_\mu)e^{S/\hbar}=0.
\]
\end{theorem}
\begin{proof}
Apply the second-order product formula to the exponential and collect powers of $\hbar$.
\end{proof}

\chapter{Perturbative factorization realization}
The algebraic quantum differential does not specify a propagator or a space of singular kernels. Fix a local realization, elliptic gauge-fixing data, support and boundary conditions, and a renormalization prescription. A heat-kernel propagator is defined away from the diagonal. Extending its graph integrals to collisions requires counterterms and compatibility with the boundary conditions.

\begin{definition}[Factorization realization data]
A factorization realization consists of renormalized local observable complexes, continuous disjoint-insertion maps, and coherent homotopies for their compositions. It includes extensions of graph weights to the required collision strata, their Stokes identities, and a quantum master equation in the permitted local functional complex. Weiss descent is a further required property of these actual complexes. A boundary realization also includes a comparison with the supplied mixed datum, preserving its two products, Lie operation, interface, and mixed law.
\end{definition}

Under these hypotheses the disjoint-insertion maps and composition homotopies define the stated factorization object. The boundary comparison identifies its boundary operations with those of $\mathbb A$. The construction problem is to produce these data from the physical fields and action. A graphwise identity supplies the corresponding composition equation, whereas descent and boundary identification require their own maps and proofs.

\part{Boundary duality, exchange, and reduction}


\chapter{The Koszul--Morita square}
A change from modules to comodules must transport the interface as well as each algebra. The mixed bar supplies the comparison object. Exchange then asks for motion in an additional topological direction, while BRST reduction asks whether the differential respects every transported operation.

Let $P$ be a dualizable boundary state for $A^\perp$.  Its boundary dual is
\[
 A_P^!=\RHom_{A^\perp}(P,P).
\]
The ordinary bar construction gives explicit twisted tensor functors.
Their extension to the mixed datum must also transport the interface and
its geometric operations.

Let $k$ be a field and let $A_{\mathrm{alg}}=\bigoplus_{w\ge0}A_w$
be an augmented associative algebra in cohomological degree zero.
Assume $A_0=k$, augmentation projects onto $A_0$, and multiplication adds weights.
Set $A=\prod_{w\ge0}A_w$ and $C=\prod_{w\ge0}B(A_{\mathrm{alg}})_w$,
with degreewise products and the weight topology.
The bar suspension lowers cohomological degree by one, with $b_2[a|b]=[ab]$.
The degree-one twisting map $\tau:C\to A$ sends $[a]$ to $a$
and vanishes on the empty word and longer words.

A module or comodule has the form $\prod_{w\ge w_0}M_w$,
with differential preserving weight, for some integer $w_0$.
Actions and coactions preserve weight.
Comodule coactions are counital and coassociative, and use the tensor product
\[
 (C\widehat\otimes M)_w
 =\bigoplus_{u+v=w}B(A_{\mathrm{alg}})_u\otimes M_v.
\]
The same convention defines $A\widehat\otimes M$.
These sums are finite in each weight.
Repeated reduced coactions vanish in weight $w$ after more than $w-w_0$ steps.
Invert quasi-isomorphisms in each weight to obtain $D_{\mathrm{wt}}(A)$
and $D_{\mathrm{wt}}(C)$.

\begin{proposition}[Ordinary twisted tensor equivalence]
The functors
\[
 F(M)=C\widehat\otimes_\tau M,\qquad
 G(N)=A\widehat\otimes_\tau N
\]
induce inverse equivalences $D_{\mathrm{wt}}(A)\simeq D_{\mathrm{wt}}(C)$.
Writing $\Delta c=\sum c_{(1)}\otimes c_{(2)}$
and $\rho n=\sum n_{(-1)}\otimes n_{(0)}$, their differentials are
\[
\begin{aligned}
 d_F(c\otimes m)={}&d_Cc\otimes m+(-1)^{|c|}c\otimes d_Mm
 +\sum(-1)^{|c_{(1)}|}c_{(1)}\otimes\tau(c_{(2)})m,\\
 d_G(a\otimes n)={}&a\otimes d_Nn
 -\sum a\tau(n_{(-1)})\otimes n_{(0)}.
\end{aligned}
\]
The counit and unit are
\[
 \epsilon_M(a\otimes c\otimes m)=a\epsilon_C(c)m,\qquad
 \eta_N(n)=\sum n_{(-1)}\otimes1\otimes n_{(0)}.
\]
\end{proposition}
\begin{proof}
The twisting equation makes both differentials square to zero.
The action, coaction, and counit identities show that the displayed maps
are chain maps and satisfy the triangle identities.
For $GF(M)$, the extra degeneracy of the two-sided bar inserts
$1[\bar a_0|a_1|\cdots|a_r]m$ into $a_0[a_1|\cdots|a_r]m$,
where $\bar a_0=a_0-\epsilon_A(a_0)1$.
With simplicial boundary $\sum_i(-1)^id_i$, it contracts the augmented bar.
Multiplying length-$r$ words by $(-1)^r$ gives the displayed differential convention.
This contraction preserves weight, so the counit is a weightwise quasi-isomorphism.
For $FG(N)$, filter each weight by the weight of the final $N$-factor.
The reduced coaction lowers that weight.
The associated graded unit is the contracted right-sided bar
$k\otimes N_v\to B(k,A_{\mathrm{alg}},A_{\mathrm{alg}})\otimes N_v$.
The filtration is finite in each total weight, so the unit is also a quasi-isomorphism.
The same finite filtrations show that $F,G$ preserve weightwise quasi-isomorphisms,
since tensoring complexes over $k$ preserves them.
Thus the unit and counit give the asserted inverse equivalences.
\end{proof}

For a right module $R$ and a left module $L$ with the same lower-weight
bounds, the two-sided bar has underlying weight-completed complex
\[
 \prod_n\ \bigoplus_{u+v_1+\cdots+v_r+w=n}
 R_u\otimes s\bar A_{v_1}\otimes\cdots\otimes s\bar A_{v_r}\otimes L_w,
 \qquad v_i>0,
\]
where the sum also runs over $r\ge0$ and $\bar A=\ker(A_{\mathrm{alg}}\to k)$.
Internal differentials and the two endpoint actions and adjacent products
form its bar differential. Its realization computes $R\otimes_A^{\mathbb L}L$
in these weight-graded categories.
Comparing it with a derived cotensor product requires a two-sided twisting
comparison on the chosen comodule category and completion.
Weight completion here is not identified with derived augmentation-adic completion.

\begin{proposition}[Centre transport under a mixed equivalence]
Suppose an equivalence between the compatible mixed module categories is supplied,
with coherent compatibility with all four bar objects and the interface operations.
It identifies the derived endomorphism objects of their identity functors.
\end{proposition}
\begin{proof}
Conjugate a natural endomorphism by the equivalence and a quasi-inverse.
The unit and counit homotopies give inverse maps on derived endomorphism objects
and preserve composition and its coherent operations.
\end{proof}

Constructing such a mixed equivalence requires the restriction, collision,
interface, and distributive-law comparisons for the actual geometric coefficients.
Its unit and counit must respect those maps and the chosen inverse limits.
A diagonal bimodule comparison must produce an actual compatible bicomodule
before tensor products can be identified with geometric cotensor products.
Transport of mixed Hochschild chains and traces further requires their
cyclic comparison and the trace domains used in their definition.
These are additional mixed Koszul--Morita constructions.
The ordinary weight-graded equivalence supplies its stated one-colour functors,
without identifying the other mixed carriers or their physical traces.

\chapter{Planar exchange and Tannaka reconstruction}
A second topological direction gives paths exchanging parallel defects.
Their homotopy classes form braid groups.
Let $\mathscr L$ be an essentially small rigid braided linear category
with a $k$-linear strong monoidal functor
$\omega:\mathscr L\to\operatorname{Vect}_{\mathrm{fd}}(k)$.
Its matrix-coefficient coend is
\[
 \Ocal(\mathbb G)=\int^{M\in\mathscr L}\omega(M)^*\tensor\omega(M).
\]

\begin{theorem}[Quantum symmetry from exchange]
The coend is a coquasitriangular Hopf algebra.
Its full linear dual acts naturally on every fiber.
The coefficient form of the braiding lies in the completed dual tensor
defined in Chapter~\ref{ch:vol2-quantum-interfaces}.
Naturality gives an RTT relation on a selected fiber.
A presentation requires a further isomorphism from the algebra
defined by those relations to the coend.
\end{theorem}
\begin{proof}
Coevaluation gives
$\Delta c_M(\lambda,v)=\sum_i c_M(\lambda,e_i)\otimes c_M(e^i,v)$,
and evaluation gives the counit.
Tensor product gives multiplication. Dual objects give the antipode.
These formulas and their identities are proved in
Proposition~\ref{prop:vol2-coend-operations}.

The braiding defines the coefficient form
\eqref{eq:vol2-coend-braiding}, and the hexagon identities imply
its compatibility equations. Naturality gives the RTT relation.
Further morphisms in the category can impose additional relations.
\end{proof}

A spectral version must specify its meromorphic coefficient domains,
expansion maps and continuous tensors. Reconstruction of the whole
category additionally uses the exactness hypotheses of
Theorem~\ref{thm:vol2-tensor-reconstruction}.

\chapter{BRST reduction}
Let $Q$ be a square-zero derivation of a mixed datum.  Extend $Q$ factorwise to every bar word.  Since $Q$ is a derivation for each structure map, it anticommutes with every bar differential.

\begin{theorem}[BRST--bar bicomplex]
The total differential
\[
 D=Q+b_{\ch}+b_\perp+b_{\mathrm{mix}}
\]
satisfies $D^2=0$.  A bounded-below complete filtration gives a convergent spectral sequence
\[
 E_1=H_Q\bigl(\BarMix(\mathbb A;K)\bigr)
 \Longrightarrow H_D\bigl(\BarMix(\mathbb A;K)\bigr).
\]
When $Q$-cohomology is concentrated and tensor-exact, the target is the mixed bar of the reduced datum.
\end{theorem}
\begin{proof}
The square vanishes by the derivation identity and the four-colour bar theorem.  The filtration by BRST degree gives the spectral sequence.  Concentration and tensor exactness identify the $E_1$ differential with the bar differential of the reduced operations.
\end{proof}

\part{Conformal covariance and Einstein completion}


\chapter{Projective connections and the Virasoro tangent complex}
A conformal boundary algebra transforms under changes of coordinates. Projective connections record its central term. Three-dimensional gravity supplies independently defined connection and coframe fields whose boundary transformation law can be compared with this algebraic structure.

Let $\Pcal_c(X)$ be the formal moduli problem of projective connections with central parameter $c$.  At a projective connection $p$, its tangent complex is locally
\[
 \Omega^{0,\bullet}(X,T_X)
 \xrightarrow{\mathcal D_{p,c}}
 \Omega^{0,\bullet}(X,K_X^{\tensor2}),
\]
where
\[
 \mathcal D_{p,c}(\epsilon)
 =\epsilon\partial p+2(\partial\epsilon)p+\frac{c}{12}\partial^3\epsilon.
\]

\begin{theorem}[Projective covariance]
A conformal map $\Vir_c\to A$ defines a formal family of mixed structures over $\Pcal_c(X)$.  Its tangent algebra is the semidirect product of the Virasoro projective complex with the deformation algebra relative to a fixed projective connection.
\end{theorem}
\begin{proof}
The conformal vector acts on fields by infinitesimal coordinate change.  The central term is the affine cocycle in the transformation law of a projective connection.  The action therefore gives a map from the projective tangent algebra to derivations of the mixed datum.  The tangent algebra of the associated family is the semidirect product.
\end{proof}

\chapter{Three-dimensional gravity as a flat Cartan theory}
Let $M$ be an oriented three-manifold, let $\Lambda=-\ell^{-2}$, and let $\mathfrak g_\Lambda$ have basis $J_a,P_a$ with
\[
 [J_a,J_b]=\epsilon_{ab}{}^cJ_c,
 \qquad [J_a,P_b]=\epsilon_{ab}{}^cP_c,
 \qquad [P_a,P_b]=\ell^{-2}\epsilon_{ab}{}^cJ_c.
\]
The invariant pairing is $\angles{J_a,P_b}=\eta_{ab}$ and vanishes on the two diagonal summands.  A Cartan connection is
\[
 \mathcal A=\omega^aJ_a+e^aP_a.
\]
Its curvature is
\[
 F_{\mathcal A}
 =\left(R^a+\frac1{2\ell^2}\epsilon^a{}_{bc}e^b\wedge e^c\right)J_a
 +T^aP_a.
\]

\begin{theorem}[Cartan--Einstein equivalence]
On the locus where $e$ is invertible, $F_{\mathcal A}=0$ is equivalent to
\[
 T^a=0,
 \qquad
 G_{\mu\nu}+\Lambda g_{\mu\nu}=0,
 \qquad
 g_{\mu\nu}=\eta_{ab}e^a_\mu e^b_\nu.
\]
The Chern--Simons action for $\mathcal A$ equals the Einstein--Hilbert action with cosmological constant, up to its boundary term and normalization.
\end{theorem}
\begin{proof}
The $P$ component of the flatness equation is the torsion equation.  Invertibility of $e$ gives the Levi--Civita spin connection.  The $J$ component states that the curvature two-form has constant sectional curvature $-\ell^{-2}$.  In three dimensions the Riemann tensor is determined by the Ricci tensor, so this is the Einstein equation.  Expanding the Chern--Simons three-form with the off-diagonal invariant pairing gives $e\wedge R+\ell^{-2}e^3$ plus an exact form 
\cite{AchucarroTownsend,Witten3DGravity}.
\end{proof}

\chapter{Four equivalent linear complexes}
Fix a constant-curvature solution $(e,\omega)$.  There are four deformation complexes.
\begin{enumerate}
\item The flat adjoint complex $\Omega^\bullet(M,\mathfrak g_\Lambda)$ with differential $d_{\mathcal A}$.
\item The Einstein--Cartan complex on $(\delta e,\delta\omega)$ with linearized torsion and curvature.
\item The metric Einstein complex on $h_{\mu\nu}$ with linearized Einstein operator.
\item The Calabi complex resolving local Killing fields by the Killing, linearized curvature, and Bianchi operators.
\end{enumerate}

\begin{theorem}[Chain equivalence of gravitational presentations]
An invertible torsion-free dreibein gives chain maps among the four complexes.  These maps are quasi-isomorphisms on every geodesically convex open set.  Their cohomology describes infinitesimal isometries, physical linearized fields, and Bianchi identities in the four languages.
\end{theorem}
\begin{proof}
The dreibein identifies internal translations with tangent vectors and internal Lorentz rotations with skew endomorphisms.  Splitting $d_{\mathcal A}$ in this basis gives the linearized torsion and curvature operators.  Eliminating $\delta\omega$ by the linearized torsion equation gives the metric Einstein operator.  The symbol sequence is the Calabi symbol sequence and is exact on constant-curvature space.  The Poincar\'e lemma for the flat adjoint connection and exactness of the symbol homotopies give local quasi-isomorphisms.
\end{proof}

\chapter{Brown--Henneaux and the projective boundary}
Write the AdS connection as
\[
 A^\pm=\omega\pm\ell^{-1}e.
\]
The Einstein action is the difference of two $\mathfrak{sl}_2$ Chern--Simons actions with level
\[
 k=\frac{\ell}{4G}.
\]

\begin{theorem}[Brown--Henneaux central charge]
The classical asymptotic Poisson algebra of each chiral Chern--Simons sector has Virasoro central charge
\[
 c=6k=\frac{3\ell}{2G}.
\]
The projective connection carried by the boundary stress tensor is the projective variable of the preceding conformal family.
\end{theorem}
\begin{proof}
Hamiltonian reduction of the boundary $\mathfrak{sl}_2$ current Poisson algebra gives the classical Virasoro bracket.  Its central term is $6k$.  Substitution of the gravitational level gives the second equality.  The coadjoint transformation of the boundary stress tensor is the projective-connection transformation law 
\cite{BrownHenneaux}.
\end{proof}

\chapter{Einstein completion of a conformal pair}
Let $A_L$ and $A_R$ be conformal mixed boundary data with equal Brown--Henneaux central charge.  Their projective families map to the same $\Pcal_c(X)$.  Define their relative field algebra by the homotopy fibre product over the projective tangent algebra.

\begin{definition}[Candidate Einstein BV completion]
Suppose the relative algebra and a continuous $L_\infty$ action of the Cartan algebra have been supplied, together with their local realization. The associated cotangent candidate is
\[
 \EinBV_3(A_L,A_R)
 =T^*[-1]\MC\!\left(
   \mathfrak g_{\mathrm{grav}}
   \ltimes_{\mathrm{Ham}}
   \mathfrak m_{A_L,A_R}
  \right).
\]
Here $\mathfrak g_{\mathrm{grav}}$ is the Cartan deformation algebra and $\mathfrak m_{A_L,A_R}$ is the relative mixed deformation algebra after projective matching.
\end{definition}

\begin{theorem}[Cotangent action for a supplied coupled algebra]
Suppose projective matching and a continuous $L_\infty$ action have constructed the displayed complete semidirect product. Suppose also that its local realization and integrated dual pairing are defined. Its shifted cotangent has action
\[
 S(\alpha,\beta)=\langle\beta,\mathcal F(\alpha)\rangle
\]
and satisfies the classical master equation. Its Bianchi identity is the $L_\infty$ identity of the coupled algebra.
\end{theorem}
\begin{proof}
The supplied semidirect product has a square-zero homological vector field. Its cotangent lift is Hamiltonian for the displayed action and remains square-zero. This proves the classical master equation and the Bianchi identity by the preceding cotangent construction.
\end{proof}

The gravitational summand of this cotangent action is the cotangent theory of Cartan deformations. It introduces independent dual fields. To identify it with Einstein BV theory requires a comparison preserving the BV pairing and the action, including boundary terms, on the invertible-coframe sector. The flatness equations alone do not provide that comparison. A physical coupled completion must also construct the matter action and its equivariant moment map. The intended Einstein reconstruction therefore retains these field and action comparisons as existence problems.

\chapter{Higher-spin completion}
Replace $\mathfrak{sl}_2$ by $\mathfrak{sl}_N$ and choose its principal $\mathfrak{sl}_2$ embedding.  The adjoint representation decomposes as
\[
 \mathfrak{sl}_N
 \cong\bigoplus_{s=2}^{N}V_{2s-2},
 \qquad
 \sum_{s=2}^{N}(2s-1)=N^2-1.
\]

\begin{theorem}[Principal higher-spin spectrum]
For the specified principal $\mathfrak{sl}_N$ Chern--Simons theory
linearized about its gravitational background, the decomposition gives
one frame-like field of each spin $2,3,\ldots,N$.
With principal Drinfeld--Sokolov boundary conditions, reduction of the
boundary affine Poisson algebra gives currents of these same spins.
\end{theorem}
\begin{proof}
The principal decomposition \eqref{eq:vol2-principal-spectrum}
contains $V_{2s-2}$ once for each $s=2,\ldots,N$.
It is the internal representation of the spin-$s$ frame-like field.
Principal Drinfeld--Sokolov reduction of this theory produces one
generating current in each conformal weight $2,\ldots,N$
\cite{CampoleoniHigherSpin,HenneauxRey}.
\end{proof}

For an independently supplied pair of $W_N$ boundary theories, a bulk
realization still requires a map from this reduction to those theories,
preserving their Poisson operations and coupling.
The spin multiplicities alone do not construct that map.

\part{Exact gravitational complexes and quantum tests}


\chapter{The four complexes in coordinates}
The Cartan flatness equations and the metric Einstein equations describe the same classical solutions on the invertible-coframe locus. Their infinitesimal complexes make the comparison explicit. Quantum determinants additionally depend on gauge fixing, boundary conditions, and determinant-line data.

\section{Flat adjoint complex}
At a flat Cartan connection $\mathcal A$, the first complex is
\[
 \Omega^0(M,\mathfrak g_\Lambda)
 \xrightarrow{d_{\mathcal A}}
 \Omega^1(M,\mathfrak g_\Lambda)
 \xrightarrow{d_{\mathcal A}}
 \Omega^2(M,\mathfrak g_\Lambda)
 \xrightarrow{d_{\mathcal A}}
 \Omega^3(M,\mathfrak g_\Lambda).
\]
Degree zero is the gauge parameter, degree one is the fluctuation, degree two is the linearized curvature, and degree three is the Bianchi identity.

\section{Einstein--Cartan complex}
Write a gauge parameter as $(\xi^a,\lambda^a)$ and a fluctuation as $(u^a,v^a)=(\delta e^a,\delta\omega^a)$.  The first differential is
\[
 \begin{aligned}
 u^a&=D_\omega\xi^a+\epsilon^a{}_{bc}e^b\lambda^c,\\
 v^a&=D_\omega\lambda^a+\ell^{-2}\epsilon^a{}_{bc}e^b\xi^c.
 \end{aligned}
\]
The second differential is
\[
 \begin{aligned}
 \delta T^a&=D_\omega u^a+\epsilon^a{}_{bc}v^b\wedge e^c,\\
 \delta F^a&=D_\omega v^a+\ell^{-2}\epsilon^a{}_{bc}e^b\wedge u^c.
 \end{aligned}
\]
The third differential is the linearized pair of Cartan Bianchi identities.

\begin{proposition}[Chain identity]
The composite of the first two Einstein--Cartan differentials vanishes on a constant-curvature torsion-free background.
\end{proposition}
\begin{proof}
Apply $D_\omega^2$ to $\xi$ and $\lambda$.  Replace the background curvature by $-\ell^{-2}e\wedge e$ and use $D_\omega e=0$.  The curvature terms cancel the explicit $\ell^{-2}$ terms.  The remaining expressions are Jacobi identities for $\epsilon_{ab}{}^c$.
\end{proof}

\section{Metric complex}
The dreibein identifies the translation parameter with a vector field.  Modulo Lorentz rotations,
\[
 h_{\mu\nu}=2\eta_{ab}e^a_{(\mu}u^b_{\nu)}.
\]
The metric differential begins with the Killing operator
\[
 K(\zeta)_{\mu\nu}=\nabla_\mu\zeta_\nu+\nabla_\nu\zeta_\mu.
\]
The linearized cosmological Einstein operator is
\[
 \begin{aligned}
 (E_\Lambda h)_{\mu\nu}
 =\frac12\bigl(&-\nabla^2h_{\mu\nu}-\nabla_\mu\nabla_\nu h
 +\nabla_\mu\nabla^\rho h_{\rho\nu}
 +\nabla_\nu\nabla^\rho h_{\rho\mu}\\
 &-g_{\mu\nu}(\nabla^\rho\nabla^\sigma h_{\rho\sigma}-\nabla^2h)
 \bigr)-\Lambda(h_{\mu\nu}-g_{\mu\nu}h).
 \end{aligned}
\]
It satisfies $E_\Lambda K=0$ and $\nabla^\mu(E_\Lambda h)_{\mu\nu}=0$.

\begin{theorem}[Explicit metric reduction]
Solving the linearized torsion equation for $v=\delta\omega$ and substituting into the linearized curvature equation gives $E_\Lambda h=0$.  The kernel of the map from $(u,v)$ to $h$ is the infinitesimal Lorentz action.
\end{theorem}
\begin{proof}
The torsion equation is algebraic in the antisymmetric part of the covariant derivative of $u$ and determines $v$.  The symmetric part of $u$ is $h/2$.  Substitution into $\delta F=0$ yields the standard variation of the Ricci tensor and the displayed Einstein operator.  An antisymmetric $u_{\mu\nu}$ is exactly a Lorentz rotation of the dreibein.
\end{proof}

\section{Calabi complex}
On a constant-curvature manifold, the Calabi complex begins
\[
 \Gamma(TM)
 \xrightarrow{K}
 \Gamma(S^2T^*M)
 \xrightarrow{R'}
 \Gamma(\mathcal R)
 \xrightarrow{B'}
 \Gamma(\mathcal B),
\]
where $R'$ is the linearized Riemann tensor and $B'$ is the linearized algebraic and differential Bianchi operator.  In three dimensions the Riemann tensor is equivalent to the Einstein tensor, so the metric and Calabi complexes are locally quasi-isomorphic.

\chapter{Boundary BFV theory and the Virasoro moment map}
The variation of Chern--Simons theory on a manifold with boundary gives the boundary symplectic form
\[
 \Omega_{\partial}=\frac{k}{4\pi}\int_{\partial M}\angles{\delta A\wedge\delta A}.
\]
A chiral polarization selects $A_{\bar z}$ as coordinate.  The residual moment map is the flatness constraint.

\begin{theorem}[Boundary current algebra]
Canonical quantization of the boundary symplectic form gives the affine current algebra at level $k$.  Drinfeld--Sokolov reduction of the principal $\mathfrak{sl}_2$ constraint gives the Virasoro moment map
\[
 T(z)=\frac1{2(k+h^\vee)}:J^aJ_a:(z)+\text{improvement},
\]
with the classical Brown--Henneaux limit $c=6k$.
\end{theorem}
\begin{proof}
The inverse of $\Omega_\partial$ is the current Poisson bracket.  Quantization gives the affine central extension.  The BRST complex of the principal constraint eliminates the gauge directions and retains the improved quadratic current.  Its classical central term is $6k$ in the gravitational normalization.
\end{proof}

\chapter{Matter stress tensor and the Einstein moment map}
Let $\mathfrak m$ be the relative field algebra of a conformal pair.  Projective covariance gives a Hamiltonian action of vector fields and therefore a moment map
\[
 \mu_{\mathrm{matt}}:\mathfrak m\longrightarrow
 \Omega^{0,\bullet}(X,K_X^{\tensor2})[-1].
\]
Its classical component is the stress tensor.

\begin{theorem}[Conservation from the master equation]
The coupled BV action
\[
 S_{\mathrm{tot}}=S_{\mathrm{grav}}+S_{\mathrm{matt}}
 +\angles{h,\mu_{\mathrm{matt}}}
\]
satisfies the classical master equation precisely when the matter moment map is equivariant.  The component of that equation linear in the diffeomorphism ghost is
\[
 \nabla^\mu T_{\mu\nu}=0.
\]
\end{theorem}
\begin{proof}
The master bracket of the coupling with itself is the defect of equivariance.  The bracket with the gravitational action is the infinitesimal diffeomorphism of the moment map.  Their cancellation is the Hamiltonian equivariance identity.  Pairing the identity with the diffeomorphism ghost and integrating by parts gives stress-tensor conservation.
\end{proof}

\chapter{One-loop gravity and analytic torsion}
Expand Chern--Simons theory around an acyclic flat connection $\mathcal A$.  Gauge fixing gives the Laplacians of the flat adjoint de Rham complex.

\begin{theorem}[One-loop determinant]
The absolute value of the one-loop partition function is the square root of the Ray--Singer torsion of the flat adjoint bundle:
\[
 |Z_{1\text{-loop}}(\mathcal A)|
 =\tau_{\mathrm{RS}}(M,\ad\mathcal A)^{1/2}.
\]
The phase is the eta-invariant contribution fixed by the framing convention.
\end{theorem}
\begin{proof}
Bosonic and ghost Gaussian integrals give alternating determinants of the Laplacians in the flat de Rham complex.  Their product is the analytic torsion.  The odd signature operator supplies the determinant phase.  This is the standard stationary-phase calculation for Chern--Simons theory \cite{RaySinger1971,Witten3DGravity}.
\end{proof}

The same determinant complex is obtained from the Einstein--Cartan, metric, and Calabi presentations through their chain equivalences.  One-loop equality is therefore a consequence of quasi-isomorphism together with compatible determinant-line orientations.

\chapter{BTZ, Cardy, and the area law}
For a unitary modular boundary theory with left and right central charges $c_L,c_R$, the asymptotic entropy is
\[
 S_{\mathrm{Cardy}}
 =2\pi\sqrt{\frac{c_L}{6}\left(L_0-\frac{c_L}{24}\right)}
 +2\pi\sqrt{\frac{c_R}{6}\left(\bar L_0-\frac{c_R}{24}\right)}.
\]
For the BTZ black hole,
\[
 L_0-\frac c{24}=\frac{(r_++r_-)^2}{16G\ell},
 \qquad
 \bar L_0-\frac c{24}=\frac{(r_+-r_-)^2}{16G\ell},
 \qquad
 c=\frac{3\ell}{2G}.
\]

\begin{theorem}[Cardy--Bekenstein equality]
Substitution gives
\[
 S_{\mathrm{Cardy}}=\frac{2\pi r_+}{4G}.
\]
This is the Bekenstein--Hawking entropy of the BTZ horizon.
\end{theorem}
\begin{proof}
Each square root becomes $\pi(r_+\pm r_-)/(4G)$.  Their sum is $2\pi r_+/(4G)$.  The modular asymptotic is Cardy's theorem, and the geometric charges are the Chern--Simons holonomies of the BTZ solution \cite{Cardy1986,BTZ1992}.
\end{proof}

\chapter{Higher-spin holonomy and boundary charges}
In $\mathfrak{sl}_N$ Chern--Simons gravity, regularity constrains the
conjugacy class of the thermal holonomy.
Such a constraint is imposed together with the boundary conditions.
The functions
\[
 \Tr(\operatorname{Hol}^s),
 \qquad 2\le s\le N,
\]
are functions of the holonomy, whereas the local
Drinfeld--Sokolov currents are functions of the connection.
For a constant connection $A$, holonomy is $\exp(LA)$.

\begin{proposition}[Recovering constant connection coefficients]
For a regular semisimple constant companion connection, its holonomy
conjugacy class determines its characteristic coefficients after
a logarithm branch is fixed on a domain where the exponential is
injective. Without that restriction the coefficients need not be
determined.
\end{proposition}
\begin{proof}
The branch recovers the eigenvalues of $A$ from those of $\exp(LA)$,
and hence its characteristic polynomial.
For $A=\left(\begin{smallmatrix}0&1\\q&0\end{smallmatrix}\right)$,
the trace is $2\cosh(L\sqrt q)$.
At $L=2\pi$, every $q=-n^2$, $n\ge1$, gives holonomy equal to the
identity. Proposition~\ref{prop:vol2-holonomy-logarithm} proves
both assertions.
\end{proof}

Hamiltonian reduction supplies the local $W_N$ Poisson brackets.
Identifying its charges with global holonomy functions requires
the actual Hamiltonian comparison, including the boundary conditions.
The characteristic polynomial of the connection does not give it.

\part{Anomalies, Hall sources, and automorphic completion}


\chapter{Four native anomaly classes}
A local anomaly, a determinant-line curvature, and an arity-zero bar term belong to different complexes. A cancellation in one complex becomes a geometric cancellation only through a comparison preserving the relevant class and its primitive. Hall and automorphic constructions require their coefficient and orientation maps for the same reason.

The construction carries four curvature classes.
\begin{enumerate}
\item A local BV anomaly lies in the degree-one local deformation complex.
\item A projective anomaly is the curvature of the determinant line over parameter space.
\item A framing anomaly is the transgression of a characteristic class to a three-manifold.
\item A curved Koszul term is the arity-zero component of a curved coalgebra.
\end{enumerate}

\begin{theorem}[Curvature-killing extension]
Let $\kappa$ be a closed central degree-two element of a dg Lie
algebra. Adjoining a central degree-one generator $u$ with
$du=\kappa$ makes $\kappa$ exact.
If $\kappa$ is a curved Maurer--Cartan term, twisting by $-u$
makes that term zero.
\end{theorem}
\begin{proof}
The differential squares to zero because $d\kappa=0$.
Its derivation identity on a bracket with $u$ follows from
$[\kappa,x]=0$.
The twisted curvature is
$\kappa+d(-u)+\frac12[-u,-u]=0$.
\end{proof}

This extension can change the deformation problem.
A geometric realization requires an independently constructed
invertible theory with the opposite anomaly.
A comparison must identify both classes in the same complex and normalization.
Under those hypotheses their anomaly classes add to zero.
The central-generator construction alone supplies neither that theory
nor the comparison map.

\chapter{Hall and Calabi--Yau sources}
Let $\Ccal$ be a Calabi--Yau category with a Hall moduli problem.  An oriented extension correspondence gives a Hall algebra $H^+$.  A factorization pushforward to the boundary curve gives a chiral-convolution algebra.  Relative induction supplies the remaining colours.

\begin{proposition}[Transport under a supplied mixed Morita equivalence]
Suppose a Hall construction supplies a complete mixed boundary datum and a Morita equivalence of its compatible module category with that of $\mathbb A$. The equivalence identifies their derived centres. A comparison of local deformation algebras, BV pairings, and actions is additional data for a BV comparison. Quantum line and amplitude comparisons additionally require the monoidal and cyclic structures used to define them.
\end{proposition}
\begin{proof}
Conjugating natural transformations by the equivalence and its inverse identifies the endomorphisms of the two identity functors. The unit and counit homotopies provide the inverse comparison and preserve composition. This proves the centre comparison. It does not construct the additional local, symplectic, or cyclic maps.
\end{proof}

An algebra map alone does not induce the required forward centre map. For example, the homomorphism $\C[x]\to M_2(\C)$ with $x\mapsto\operatorname{diag}(0,1)$ sends the central element $x$ to a noncentral matrix. Its commutator with the matrix unit $E_{12}$ is $-E_{12}$. Thus even the degree-zero condition fails for a centre map extending the given algebra map. Hall multiplication maps must be supplemented by an actual coherent module comparison before centre transport can be applied.

For $X=K3\times E$, suppose a geometric Hall source supplies the complete charge-graded BPS Lie superalgebra and Cartan action required in Volume~IV. Its semidirect product with the common Cartan defines $\mathfrak b_X$. Volume~IV then constructs the completed Lie amalgam with the Igusa Borcherds Borel $\mathfrak b_\Delta$. Applying the completed enveloping functor gives the cocommutative Hopf algebra
\[
 H_{X\boxtimes\Delta}
 =\widehat U(\mathfrak b_X)
  \widehat *_{\widehat U(\mathfrak h_\Pi)}
  \widehat U(\mathfrak b_\Delta).
\]
The two Hopf retractions preserve the geometric and automorphic sectors separately.  A $q$-deformed mixed Hopf algebra is supplied by the quantum amalgam theorem when a geometric Hall coproduct and pairing are present.

\section{Mixed realization of the amalgam}
The Hopf retractions above are maps of their stated algebraic objects. A mixed realization must construct chiral and transverse operations, an interface, and a compatible distributive law. To realize the retractions, it must then supply maps preserving those operations and the chosen completions.

A BV retraction requires further maps between the local deformation algebras, together with compatible pairings and actions. The counterexample to forward centre functoriality prevents deducing these maps from the Hopf retractions alone. The intended geometric--automorphic theory asks for this mixed realization and for a primitive comparison quotient on its actual boundary operators. Its interaction ideal is the kernel of that constructed quotient, rather than an ideal determined by equality of scalar characters.

\chapter{Igusa root states and the mixed-boundary comparison}

The Gritsenko--Nikulin presentation defines the root Lie superalgebra
$\mathfrak g_{\Delta_5}$ and its positive subalgebra $\mathfrak n_+$.
Let $d_\alpha=\operatorname{sdim}(\mathfrak g_{\Delta_5})_\alpha$,
and fix the positive cone, its proper height, and Weyl vector $\rho$.
Put
\[
 P=\prod_{\alpha>0}(1-e^{-\alpha})^{d_\alpha},\qquad
 \DenIg=e^{-\rho}P=D_5(2Z),\qquad \jmath(Z)=Z/2.
\]
The products are first interpreted in the positive-height completion.

\begin{theorem}[Root character and reduced scalar]
The root Fock space $\operatorname{Sym}(\mathfrak n_+)$, formed with
super symmetry, has supercharacter $P^{-1}$. Consequently
\[
 -\jmath^*\!\left(e^{2\rho}
   \operatorname{sch}(\operatorname{Sym}(\mathfrak n_+)^{\otimes2})\right)
 =-(\jmath^*\DenIg)^{-2}=-\chi_{10}^{-1}.
\]
For a nonsingular projective K3 surface $S$ and an elliptic curve $E$, this scalar also has an enumerative interpretation. Let $N_{g,h,d}$ be the reduced disconnected Gromov--Witten invariant in class $(\beta_h,d)$, where $\beta_h$ is nonzero and primitive with square $2h-2$. Use the incidence insertion $\beta_h^\vee\otimes[\mathrm{point}]$, with $\langle\beta_h,\beta_h^\vee\rangle=1$. Then the series
\[
 \mathcal Z(u,q,s)=\sum_{g,h,d\ge0}N_{g,h,d}u^{2g-2}q^{h-1}s^{d-1}
\]
is the Laurent expansion of $-\chi_{10}(r,q,s)^{-1}$ under $r=e^{iu}$, by \cite[Theorem~1]{OP}.
\end{theorem}
\begin{proof}
An even root vector contributes $(1-e^{-\alpha})^{-1}$ to the
supercharacter of its symmetric algebra. An odd root vector contributes
$1-e^{-\alpha}$, because its symmetric algebra in super vector spaces
is an exterior algebra. Multiplying the two contributions gives
$P^{-1}$. Proper height makes every coefficient calculation finite.
The Weyl factor gives $e^{2\rho}P^{-2}=\DenIg^{-2}$.
Pullback by $\jmath$ gives $D_5(Z)^{-2}=\chi_{10}(Z)^{-1}$.
Oberdieck--Pixton, Theorem~1, identifies the separately defined reduced numerical generating series with its negative \cite{OP}. This last theorem
is an enumerative input; the character calculation supplies the root side.
\end{proof}

A primitive realization starts with parity-preserving root classes in a charge-graded Lie superalgebra $\mathfrak p$ satisfying the defining relations. Require the classes to preserve charge. The free Lie presentation gives an even charge-preserving map $f:\mathfrak n_+\to\mathfrak p$.

Applying $f$ to tensor words constructs $U(f)$. With trivial coefficients, use Chevalley chains $C_*^{\mathrm{CE}}(\mathfrak a)=\operatorname{Sym}(\mathfrak a[1])$, where the homological suspension raises degree by one and the differential contracts pairs by the bracket. Applying the suspension of $f$ to every factor gives a chain map because $f$ preserves brackets.

For the level-zero Lie conformal current, define $\operatorname{Cur}(\mathfrak a)=\C[\partial]\otimes\mathfrak a$, with $\partial$ even. Its generator bracket is $[x_\lambda y]=[x,y]$. Extend it by $[\partial x_\lambda y]=-\lambda[x_\lambda y]$ and $[x_\lambda\partial y]=(\partial+\lambda)[x_\lambda y]$. The formula $\partial^j x\mapsto\partial^j f(x)$ defines $\operatorname{Cur}(f)$. Charge preservation makes these maps compatible with the specified proper-height quotients and their inverse limits.

A complete mixed boundary realization requires further geometric maps.
The compact source consists of sheaves and extension flags on a common
elliptic torsor. Its critical realization must supply reduced coefficient
complexes, orientation cross-line maps, unit and flag coherences, and the
admissible pull--push operations. A chiral comparison must preserve actual
support and moving-insertion operations. A quantum-double comparison must
preserve coproducts, pairings and the domain of exchange operators.
The finite Calabi--Yau quiver completions require their specified dg maps;
ordinary full-subquiver inclusions do not in general give those maps.

The mixed datum to which a BV envelope is applied must therefore include
these operations and compatible maps between them. Identifying the
primitive algebra of that envelope requires an actual map of its
deformation complex to the root complex. A determinant comparison requires
a skew perfect family, a square-root line and section, and an isomorphism
to the automorphic line preserving transport and normalization.
An operator-trace interpretation further requires its state space and
trace domain. Constructing this datum and these comparisons is the Igusa
mixed-boundary problem. The root character identity specifies the scalar
that their final evaluation must recover.

\chapter{The carrier lattice}
The complete theory uses the carrier set
\[
 \mathfrak C=
 \{\mathrm{diag},\mathrm{ch},\perp,\mathrm{int},
   \mathrm{ex},\mathrm{cyc},\mathrm{mod},\mathrm{BV},\mathrm{Ein}\}.
\]
For a subset $S\subset\mathfrak C$, let $\operatorname{Str}_S(\mathbb A)$ be the homotopy limit of all mapping spaces attached to incidence cells supported in $S$.

\begin{theorem}[Carrier descent]
When $S$ and $T$ form a relation-excisive pair, the square
\[
\begin{tikzcd}
\operatorname{Str}_{S\cup T}(\mathbb A)\arrow[r]\arrow[d]&
\operatorname{Str}_S(\mathbb A)\arrow[d]\\
\operatorname{Str}_T(\mathbb A)\arrow[r]&
\operatorname{Str}_{S\cap T}(\mathbb A)
\end{tikzcd}
\]
is a homotopy pullback.  Thus a global structure is assembled from its carrier restrictions together with their common interface.
\end{theorem}
\begin{proof}
Relation-excision identifies the incidence category of $S\cup T$ with the pushout of the incidence categories for $S$ and $T$ over their intersection.  Homotopy limits send this pushout to the displayed pullback.
\end{proof}

\begin{definition}[Relation-complete carrier cover]
Let $U\subseteq\mathfrak C$ and let $\mathcal U=\{S_i\}_{i\in I}$ be a finite family with $U=\bigcup_iS_i$.  The family is \emph{relation-complete} when the canonical functor from the Cech homotopy colimit
\[
 \underset{\varnothing\ne J\subseteq I}{\operatorname{hocolim}}
 \mathcal I_{\cap_{j\in J}S_j}
 \longrightarrow \mathcal I_U
\]
is an equivalence.  Equivalently, every incidence cell and every relation among its faces is contained in one member of the cover, with all overlaps recorded by the intersections.
\end{definition}

\begin{theorem}[Power-set coherence]
The assignment
\[
 S\longmapsto\operatorname{Str}_S(\mathbb A)
\]
is a contravariant functor on the full power set of $\mathfrak C$.  For every relation-complete carrier cover $\mathcal U=\{S_i\}_{i\in I}$ of $U$, the restriction map is an equivalence
\[
 \operatorname{Str}_U(\mathbb A)
 \xrightarrow{\sim}
 \underset{\varnothing\ne J\subseteq I}{\operatorname{holim}}
 \operatorname{Str}_{\cap_{j\in J}S_j}(\mathbb A).
\]
Thus every subtheory, every intersection of subtheories, and the complete theory are values of one descent object.
\end{theorem}
\begin{proof}
An inclusion $S\subseteq T$ induces the full inclusion $\mathcal I_S\subseteq\mathcal I_T$ and therefore the restriction map on homotopy limits.  Composition of inclusions gives functoriality.  A relation-complete cover presents $\mathcal I_U$ as the displayed homotopy colimit of its intersection diagram.  Mapping that colimit into the diagram of structure spaces converts it into the stated homotopy limit.  The resulting equivalence is natural under refinement of the carrier cover.
\end{proof}

\begin{definition}[Exact-support deformation complex]
For $S\subseteq\mathfrak C$, let $\Def_{\leq S}$ be the subcomplex generated by incidence cells supported in $S$, and set
\[
 \Def_{=S}
 =\operatorname{cofib}\!\left(
   \underset{T\subsetneq S}{\operatorname{hocolim}}
   \Def_{\leq T}
   \longrightarrow
   \Def_{\leq S}
 \right).
\]
Thus $\Def_{=S}$ retains the deformation cells whose minimal carrier support is exactly $S$.
\end{definition}

\begin{theorem}[Carrier-support spectral sequence]
Let $\Def_{\mathbb A}^{\mathrm{full}}$ be the cellular convolution complex of the complete incidence diagram.  Filter it by the cardinality of minimal carrier support.  The finite filtration gives a convergent spectral sequence
\[
 E_1^{p,q}
 =\bigoplus_{\substack{S\subseteq\mathfrak C\\|S|=p}}
 H^{p+q}(\Def_{=S})
 \Longrightarrow
 H^{p+q}\!\left(\Def_{\mathbb A}^{\mathrm{full}}\right).
\]
A class on the $S$-summand is a deformation or curvature class whose first complete geometric carrier is $S$.  The first differential is the signed incidence map between exact-support quotients.
\end{theorem}
\begin{proof}
The support-cardinality filtration is finite because $\mathfrak C$ is finite.  Its $p$th associated-graded quotient is the direct sum of the exact-support complexes $\Def_{=S}$ with $|S|=p$.  The spectral sequence of a finite filtered complex gives the displayed first page and converges strongly to the cohomology of the total convolution complex.  The cellular determinant orientation gives the incidence signs.
\end{proof}

\chapter{Sixteen structural laws}
The theory is governed by the following laws.
\begin{enumerate}
\item A geometric carrier determines the type of its operation.
\item The diagonal stores the complete principal-part filtration.
\item The interval stores ordered fusion.
\item The plane stores exchange.
\item The circle stores trace.
\item The stable curve stores modular sewing.
\item The transverse, associative-chiral, and chiral-Lie bars are distinct.
\item Chiral PBW compares the last two on enveloping objects.
\item The interface bar records the mixed distributive law.
\item The derived centre is the universal interior action.
\item The deformation algebra is the tangent algebra of the centre-extended datum.
\item Maurer--Cartan curvature is the field equation.
\item The $L_\infty$ identity is the Bianchi identity.
\item Shifted cotangent is the universal classical BV envelope.
\item Hodge-completed derived de Rham theory is the intrinsic algebraic quantum complex of the oriented formal moduli problem.
\item A gravitational realization requires projective matching, a coupled local field algebra, and a comparison preserving the physical BV action and pairing.
\end{enumerate}

\chapter{Coherent actions and gravitational reconstruction}
A supplied mixed boundary datum defines compatible modules and their derived natural endomorphisms. A specified coloured operad organizes deformations of that datum together with its central action. These are constructions relative to the full operations and their coherence maps. The centre alone does not recover the datum from which the deformation algebra is formed.

A supplied local realization of the deformation algebra admits shifted cotangent transgression. With compatible invariant orientations, the finite perfect charts also carry the algebraic BV differential. The physical bulk is defined independently by its fields, action, gauge symmetry, boundary conditions, and quantization. Reconstruction requires a coherent action map from its observables to an independently defined geometric centre, followed by an equivalence theorem on the stated sector.

For gravity, the invertible coframe gives the classical Cartan--Einstein correspondence. The cotangent construction supplies a formal comparison theory with extra dual fields. Its identification with gravitational BV fields requires an action-preserving and pairing-preserving map. Quantum boundary currents, a continuum measure, global holonomy sectors, allowed fillings, and compatible states remain necessary for the quantum and global comparison.

Hall and Calabi--Yau geometry can supply boundary operations when their correspondences, coefficients, orientations, and comparison maps are constructed. Exchange additionally uses a monoidal category and a fibre functor. Cyclic pairings and sewing determine amplitudes only on their stated domains. These comparisons are the mathematical steps by which local boundary operations could reconstruct physical and geometric theories.

\part{Technical appendices to the preceding book}

\chapter{The Chern--Simons normalization}
The gravitational charge calculation uses the action normalization. The mixed bar calculation uses the determinant orientation of its edges. The following formulas specify these conventions.

With invariant pairing $\angles{J_a,P_b}=\eta_{ab}$, the Chern--Simons action is
\[
 S_{\mathrm{CS}}[\mathcal A]
 =\frac{k}{4\pi}\int_M
 \angles{\mathcal A\wedge\dd\mathcal A
 +\frac23\mathcal A\wedge\mathcal A\wedge\mathcal A}.
\]
For $A^\pm=\omega\pm\ell^{-1}e$ one finds
\[
 S_{\mathrm{CS}}[A^+]-S_{\mathrm{CS}}[A^-]
 =\frac{k}{2\pi\ell}\int_M
 \epsilon_{abc}e^a\wedge
 \left(R^{bc}+\frac1{3\ell^2}e^b\wedge e^c\right)
 +\int_{\partial M}(\cdots).
\]
The equality with the Einstein--Hilbert normalization gives $k=\ell/(4G)$.

\chapter{The principal embedding count}
The exponents of $\mathfrak{sl}_N$ are $1,2,\ldots,N-1$.  Under the principal $\mathfrak{sl}_2$ embedding,
\[
 \mathfrak{sl}_N=\bigoplus_{m=1}^{N-1}V_{2m}.
\]
The dimensions satisfy
\[
 \sum_{m=1}^{N-1}(2m+1)
 =(N-1)(N+1)=N^2-1.
\]
Writing $s=m+1$ gives one field at every spin $s=2,\ldots,N$.

\chapter{Low-arity mixed signs}
Let $x|k|y$ be a suspended interface word.  The two length-lowering contractions are
\[
 b_L[x|k|y]=(-1)^{|x|}(xk)|y,
 \qquad
 b_R[x|k|y]=(-1)^{|x|+|k|}x|(ky).
\]
Applying the remaining contraction gives the two composites $(xk)y$ and $x(ky)$ with opposite determinant-line signs.  Their sum vanishes by the interface associativity relation.  A transverse and a chiral contraction acquire the product orientation
\[
 \det(E_{\ch}\oplus E_\perp)
 =\det(E_{\ch})\tensor\det(E_\perp),
\]
so reversing their order contributes the Koszul sign and yields anticommutation.

\backmatter
\bibliographystyle{alpha}
\bibliography{references}
\end{document}
